% ======================================================================
%  Designing AI Serving Systems (at Cloud Scale) — Korean edition
%  Euiseong Seo — Sungkyunkwan University
%
%  Master file.  Compile with XeLaTeX (twice, for the table of contents):
%     xelatex main   &&   xelatex main
%  or simply:  make        (see Makefile)
%
%  Structure: front matter (title, preface, contents), three parts of
%  chapters, and an epilogue.  Chapter 8 is the course midterm and has no
%  chapter of its own, so the numbering skips from 7 to 9 — each chapter
%  file sets its own number, so the gap is intentional and self-contained.
% ======================================================================
\documentclass[11pt,oneside]{book}
\usepackage{style/tbb}

\hypersetup{
  pdftitle={Designing AI Serving Systems (at Cloud Scale)},
  pdfauthor={Euiseong Seo},
  pdfsubject={AI 서빙 시스템, 클라우드 인프라, 추론},
  bookmarksnumbered=true,
}

% figures live in each chapter's own figures/ dir; filenames are unique per
% chapter, so search-path resolution of `figures/<name>.png` is unambiguous.
\graphicspath{%
  {chapters/ch01/}{chapters/ch02/}{chapters/ch03/}{chapters/ch04/}%
  {chapters/ch05/}{chapters/ch06/}{chapters/ch07/}{chapters/ch09/}%
  {chapters/ch10/}{chapters/ch11/}{chapters/ch12/}{chapters/ch13/}%
  {chapters/ch14/}{backmatter/epilogue/}}

\begin{document}
\frontmatter

%% Title page — palette bar, story descriptor,
%% three-part foot line.
\thispagestyle{empty}
\begingroup
\setlength{\parskip}{0pt}\raggedright

% top strip: edition line over a hairline
\noindent{\color{tbbgray}\large A Principles-First Textbook\hfill 한국어판 · 2026}\par
\vspace{3pt}{\color{tbbtableborder}\rule{\linewidth}{0.4pt}}\par

\vspace{2.6cm}

% title
{\titlefont\bfseries\fontsize{36}{40}\selectfont\color{tbbink}Designing AI Serving Systems\par}
\vspace{6pt}
{\titlefont\bfseries\fontsize{28}{32}\selectfont\color{tbbblue}(at Cloud Scale)\par}

\vspace{16pt}
% the five-colour palette bar (the figures' fixed palette, as a signature)
\noindent\textcolor{tbbblue}{\rule{0.2\linewidth}{10pt}}%
\textcolor{tbbteal}{\rule{0.2\linewidth}{10pt}}%
\textcolor{tbbamber}{\rule{0.2\linewidth}{10pt}}%
\textcolor{tbbpurple}{\rule{0.2\linewidth}{10pt}}%
\textcolor{tbbred}{\rule{0.2\linewidth}{10pt}}\par

\vspace{18pt}
% descriptor
{\fontsize{15}{21}\selectfont\color{tbbgray}%
AI 서비스 밑바탕의 원리 --- 가상화, 컨테이너, 쿠버네티스, 서빙 엔진, 그리고 비용 --- 를 하나의 이야기로 엮었다. 21시 04분에 무너진 데모에서 시작해, 의도한 대로 정확히 월 900달러에 운영되는 서비스까지.\par}

\vfill

% author
{\bfseries\fontsize{24}{28}\selectfont\color{tbbink}서의성\par}
\vspace{4pt}
{\fontsize{14}{18}\selectfont\color{tbbgray}성균관대학교\par}

\vspace{1.4cm}
% foot: the three parts over a hairline
{\color{tbbtableborder}\rule{\linewidth}{0.4pt}}\par\vspace{6pt}
\noindent{\fontsize{12}{16}\selectfont%
\textcolor{tbbblue}{\textbf{1부}} \textcolor{tbbgray}{인프라와 원리}\quad
\textcolor{tbbteal}{\textbf{2부}} \textcolor{tbbgray}{AI 서빙 시스템}\quad
\textcolor{tbbamber}{\textbf{3부}} \textcolor{tbbgray}{운영}\par}
\endgroup
\clearpage


%% Preface — 그 위의 계층
\frontchapter{그 위의 계층}{서문}
\chaptersubtitle{운영체제 교수가 운영체제 위의 세계에 관한 책을 쓴 까닭}

나는 17년 동안 학부 운영체제를 가르쳤고, 매번 같은 고백으로 수업을 시작했다. \textbf{여러분 가운데 운영체제를 직접 만들 사람은 거의 없다.} 세상은 손에 꼽을 만큼 적은 커널 위에서 돌아가며, 그 대부분은 강의실의 학생들보다 나이가 많다. 새 커널의 시장 규모는 아무리 후하게 반올림해도 0에 가깝다. 그러면 해마다 용감한 학생 하나가 손을 든다. \textit{그런데 우리는 왜 이 과목을 배우나요?} 이 수업은 애초부터 운영체제라는 결과물 자체를 배우는 자리가 아니었기 때문이다. 그 결과물을 이루는 생각, 곧 스케줄링, 캐싱, 가상 메모리, 격리, 그리고 천 개의 프로세스에 동시에 약속하고도 그중 충분히 많은 약속을 지키는 기술을 배우는 자리였다. 또 학생이 컨텍스트 스위칭은 마법이 아니라 터무니없이 빠른 장부 정리일 뿐임을 깨닫는 바로 그 순간을 지켜보는 자리이기도 했다. 학생은 지적 불꽃을 얻고, 나는 그 불꽃이 튀는 모습을 보았다. 더없이 만족스러운 거래였다.

그런데 땅이 움직였다. 앞으로 학생들이 코드를 직접 쓰는 양이 갈수록 줄어들고, 편집기에 자료구조를 타이핑하는 일은 스스로 \textit{하는} 일보다 누군가에게 \textit{지시하는} 일이 되리라는 전망은 이제 대담한 주장도 아니다. 여전히 프로그래밍 솜씨의 가치가 예전과 다르지 않을 것이라고 말하는 동료도 있다. 나는 그만큼 낭만적이지 않다. 계산기가 나온 뒤에도 산술은 여전히 진짜 기술이고 때로는 결정적이지만, 더는 한 사람의 경력을 다른 사람과 가르는 기준이 아니다. 프로그래밍도 그렇게 될 것이다. 나는 이 말을 기쁜 마음으로 하는 것이 아니다. 나 역시 그 솜씨를 좋아했다. 하지만 머리말이라면 적어도 저자의 진짜 생각만큼은 솔직히 밝혀야 한다.

\subsection*{다시, 1968년}

땅이 움직일 때는 이전에 어디에서 움직였는지 돌아보면 마음을 다잡을 수 있다. 1968년 가을, 뜻밖에도 NATO가 바이에른 가르미슈에 약 50명을 모았다. 일부러 도발적으로 고른 기치가 \textit{소프트웨어 공학(software engineering)}이었다. 당시 많은 사람에게 두 단어의 조합 자체가 모순처럼 들렸다. 모임의 계기는 모두가 사석에서 \textbf{소프트웨어 위기}라고 부르던 상황이었다. 곳곳의 프로젝트가 일정을 넘기고, 예산을 초과하고, 틀린 결과물을 내놓았다. 소프트웨어가 독립된 학문 분야로 태어난 날을 하나 고르라면 바로 그 주만큼 어울리는 때도 없다.

아홉 달 뒤, 그 주장의 의미는 높은 하늘 위에서 입증되었다. 아폴로 11호가 달로 하강하던 중 잘못 설정된 레이더가 무의미한 작업을 쏟아 내자 유도 컴퓨터에 1202와 1201 경보가 번쩍이기 시작했다. 마거릿 해밀턴이 이끈 MIT 팀의 소프트웨어는 우선순위가 낮은 작업을 버리고 착륙 계산을 계속 수행했다. 언론의 제목을 장식한 것은 경보 코드였지만, 착륙을 구한 것은 \textbf{스케줄링 정책}이었다. 그보다 한 달 전 IBM은 소프트웨어를 기계와 분리해 판매하는 이른바 "언번들링"을 발표했고, 소프트웨어는 제품이자 산업이며 돈이 되었다. 보이지 않는 시스템 소프트웨어가 눈앞의 결과를 결정했다. 1960년대를 한 문장으로 요약하면 그렇다.

나는 지금 그 역사가 한 층 위에서 재판되고 있다고 믿는다. 오늘날의 AI 서비스는 운영체제 \textit{위}에 놓인, 그에 못지않게 보이지 않는 시스템 소프트웨어 계층에 서 있다. 하이퍼바이저는 실리콘을 조각내 빌려 주고, 컨테이너 런타임과 쿠버네티스는 수많은 머신을 한 대의 컴퓨터처럼 바꾸며, 서빙 엔진은 GPU 메모리에 천 개의 대화를 밀어 넣는다. 그 위로는 에이전트형 AI의 하니스와 파이프라인이 빠르게 늘어난다. 무대에서 이 계층 자체를 시연하는 사람은 없지만, 무대 위 모든 시연이 이 계층에 기대고 있다.

운영체제 교수에게 가장 즐거운 대목은 이 모든 것이 낯익다는 사실이다. 연구자들이 LLM \textbf{KV 캐시}의 메모리 낭비를 마침내 길들였을 때, 해법이 된 PagedAttention의 핵심은 페이지 테이블이었다. 1962년 Atlas 머신이 도입한 바로 그 기술이다. 제11장의 한 절 제목은 "운영체제가 안으로 들어오다"인데, 이는 비유가 아니다. 스케줄링은 연속 배칭으로, 격리는 네임스페이스와 cgroup으로 돌아왔다. 대기행렬 이론은 p99 지연 시간이라는 이름표를 달고 다시 나타났다. 내가 17년 동안 가르친 내용은 하나도 버려지지 않았다. 이름을 바꾸고 GPU 위에서 실행되는 모습으로 모두 돌아왔다.

\subsection*{두 번 가르치고 폐강한 과목}

이 책에는 실패한 선조가 있다. 독자에게도 그 집안 내력을 들려줄 필요가 있다. 챗봇 덕분에 "토큰"이라는 말이 일상어가 되기 전, 나는 \textit{클라우드와 AI}라는 대학원 과목을 두 차례 가르쳤다. 방향은 옳았지만 변화의 속도를 감당하지 못한 과목이었다. 쿠버네티스가 아직 전쟁에서 승리하기 전이어서, 진지한 사람들도 Mesos와 Swarm에 베팅했다. 9월에 가르친 AI 프레임워크가 12월이면 구식이 되었다. 첫 주에 쓴 강의계획서는 15주 차가 되기도 전에 고쳐야 했다. 내가 원한 교재, 곧 릴리스 노트보다 오래 살아남을 원리를 중심으로 구성한 책은 세상에 없었다. 결국 여기저기서 빌린 슬라이드를 희망으로 간신히 이어 붙여 강의했다.

학생들이 그 과목을 즐겼다고는 생각하지 않는다. 적어도 나는 분명 즐기지 못했다. 두 번째 강의를 마친 뒤 과목을 폐강했다. 더 나은 한 학기를 누릴 자격이 있는 주제를 제대로 가르치지 못했다는 교사 특유의 죄책감이 남았다. 이 책은 무엇보다도 그 빚을 갚는 일이다.

\subsection*{왜 지금인가}

달라진 점은 이제 뼈대가 굳었다는 것이다. 트랜스포머는 아홉 살이 되었고, 이를 송두리째 대체하려던 시도는 아직 왕좌를 빼앗지 못했다. 도전자들은 트랜스포머 곁에서 하이브리드 형태로 살아남았고, 가장 값비싼 부작용인 KV 캐시는 여전히 이 분야의 대표적인 미해결 과제로 남아 이제는 그 자체로 가르칠 만한 주제가 되었다. 서빙 스택도 연속 배칭, 대화를 위한 페이지드 메모리, 지연 시간을 재는 두 시계라는 알아볼 수 있는 형태로 수렴했다. 쿠버네티스는 전쟁에서 완승한 나머지 업계 언론이 이제 \textit{지루하다}고 부른다. 인프라 분야에서 받을 수 있는 최고의 찬사다. 클라우드 네이티브 아키텍처는 조직이 시스템을 만드는 보편적인 방식이 되었다. 앞으로도 하니스, 런타임, 파이프라인 같은 새 모듈은 이 뼈대 위에 계속 등장하겠지만, 모두 이 뼈대 \textit{안으로} 들어올 것이다. 세부 사항이 매주 바뀔 때는 블로그 글을 쓴다. 원리가 10년은 버틸 듯할 때 비로소 조심스럽게 교재를 쓸 수 있다.

\subsection*{이 책을 만든 방법: 공개해야 할 이야기}

교재가 필요하다는 사실을 아는 것과 실제로 쓰는 것은 전혀 다른 종목이다. 도널드 커누스는 1962년 컴파일러에 관한 소박한 책 한 권을 쓰기로 계약했다. 첫 권은 6년 뒤에 나왔고, 60여 년이 지난 지금도 그 장대한 프로젝트는 미완성이다. 프레드 브룩스는 프로젝트가 어떻게 1년씩 늦어지는지 묻고 스스로 답했다. \textit{하루씩 늦어진 결과다.} 교재도 똑같다. 그렇다고 대학원생 한 명에게 한 장씩 맡긴 뒤 교수 이름으로 묶어 내던 학계의 전통적인 해결책을 꺼낼 수도 없었다. 그런 시대는 끝났고, 마땅히 그래야 한다. 대학원생이 세상에 내놓아야 할 것은 연구이지 내 책의 원고가 아니다.

이 책을 구한 것은 다름 아닌 이 책의 주제였다. 나는 AI 에이전트를 공동 저자로 삼는 일에 회의적인 태도로 시작했다. 결과를 있는 그대로 말하자면, 이 책을 손으로 썼다면 지금보다 더 잘 쓰지 못했을 것이며 이번 10년 안에 끝내지도 못했을 것이다. 내가 보탠 것은 뼈대였다. 강의, 강의계획서, 무엇이 근본 원리이고 무엇이 유행인지 가리는 판단이었다. 에이전트는 초안과 그림, 계산, 그리고 지칠 줄 모르는 인내를 보탰다. 나는 동그라미 하나와 선 다섯 개로 사람을 그린 연필 스케치를 건넸고, 영화 같은 결과물을 돌려받았다. 농기계가 한 쌍의 손으로 경작할 수 있는 면적을 몇 배로 늘렸듯, 이제 지적 노동에도 같은 배율의 증폭이 찾아왔다. (전통적으로 머리말은 경건한 문구로 끝난다. 도움을 준 이들에게 감사를 돌리고, 남은 오류는 오롯이 저자의 몫이라고 말한다. 이 시대에 맞게 고쳐 보겠다. 이 책에서 좋은 부분은 모두 내 덕이고, 남은 오류는 전적으로 AI 탓이다. 왜냐하면 사람은 여전히 중요하기 때문이다.)

공개할 이야기 안에 한 가지가 더 있다. 초안의 모든 문장은 빌린 GPU에서 실행되는 모델이 토큰으로 만들어 내게 보낸 것이다. 그 아래에서는 요청이 연속으로 배칭되고, 대화가 KV 캐시를 통해 페이지 단위로 오갔으며, 제9장부터 제12장까지 다루는 장치 전체가 쉼 없이 돌아갔다. 이 책은 자신이 가르치는 바로 그 시스템으로 만들어졌다. 이보다 강의계획서를 더 강하게 뒷받침하는 증거는 떠올리기 어렵다.

\subsection*{이 책의 내용과 읽는 방법}

이 책은 15주 강의를 따라 세 부로 구성된다. \textbf{1부: 인프라와 원리(제1장\textasciitilde{}제6장)}에서는 의도적으로 스택을 내려간다. 클라우드 인스턴스와 컨테이너의 실체(네임스페이스, cgroup, 계층형 파일 시스템이라는 세 커널 기능이 트렌치코트를 걸친 모습), 쿠버네티스가 선언된 의도를 실제로 실행되는 머신 집합으로 바꾸는 방법, 가중치와 데이터가 머무는 곳을 살펴본다. \textbf{2부: AI 서빙 시스템(제7장, 제9장\textasciitilde{}제12장)}에서는 꼬리 지연 시간과 SLO, 서빙 프레임워크, 양자화와 컴파일, LLM 추론만의 물리 법칙, 비용표를 붙인 스케일링을 통해 서빙이라는 분야의 기초를 세운다. \textbf{3부: 운영(제13장\textasciitilde{}제14장)}에서는 레지스트리, 파이프라인, 카나리를 이용해 소란 없이 시스템을 바꾸는 법과 모니터링, 알림, 드리프트로 이루어진 서비스의 신경계를 배운다. 제8장은 없다. 중간고사에서 독자가 직접 쓰는 장이기 때문이다. 마지막에는 맺음말이 이 모든 이야기를 문밖까지 배웅한다.

전체 내용은 하나의 연속된 이야기로 이어진다. 학생 팀의 데모는 21시 04분, 정확히 100명의 사용자가 몰리자 무너진다. 이를 "고치기" 위해 빌린 GPU는 대부분 놀았지만 월말 청구서는 779.66달러다. 15주 동안 엔지니어링을 거친 뒤 같은 서비스는 약속한 범위 안에서 월 900달러에 운영된다. \textbf{정확히} 900달러다. 그때는 비용을 당하는 것이 아니라 스스로 선택하기 때문이다. 각 장은 새벽 2시에도 결정을 내릴 수 있을 만큼 선명한 한 줄짜리 \textit{기념품}을 남긴다. \code{exit 137}은 커널이 남기는 작별 인사이고, \textit{평균은 아무도 설명하지 못한다.} 풀이 예제는 독자 앞에서 계산 과정을 숨김없이 보여 준다. 좋은 세미나라면 질문을 던질 바로 그 지점에서 \textit{생각해 보기} 상자가 흐름을 멈춘다. 모든 연습문제에는 정답과 해설이 딸려 있다. 이 책은 \textbf{혼자} 공부하는 사람도 읽을 수 있어야 하기 때문이다.

이는 의도한 구성이다. 지금은 AI 시대다. 설명은 값싸졌다. 새벽 3시에도 AI 튜터가 무엇이든 기꺼이 설명해 준다. 하지만 \textbf{큐레이션은 여전히 값싸지 않다.} 무엇을 어떤 순서로 배우고, 어떤 원리로 연결할 것인가. 그것이 강의이며, 이 책은 종이 위의 강의가 되려 한다. 학생에게 뼈대와 질문을 건네면 나머지는 학생과 AI가 해낼 것이다. (동료 교수들에게: 사직서는 아니다. 여전히 누군가는 뼈대를 골라야 한다.) 형식적인 선수 요건은 다음과 같다. Python은 필수다. 운영체제 과목은 강력히 권한다. 필요한 내용은 책에서 다시 세우므로 반드시 이수해야 하는 것은 아니다. 기계학습 입문 과목도 도움이 되지만 전제로 삼지는 않는다.

\vspace{1.0pt}

\begin{calloutbox}{tbbamber}{tbbamberfill}
{\normalsize\bfseries\color{tbbamberdark}이 책을 활용하는 방법}\par\vspace{3pt}
\begin{tbbitemize}
  \item \textbf{정규 강의}: 각 장은 15주 학기의 한 주와 일대일로 대응한다. 8주 차에는 중간고사, 15주 차에는 데모 데이가 있으며, 별도 강의계획서에 프로젝트 평가 기준과 실습 일정을 담았다.
  \item \textbf{독학}: 한 장을 읽고 실습한 뒤 AI 튜터와 연습문제의 답을 놓고 논쟁해 보자. 모든 문제에 정답과 해설이 있어 논쟁의 승자를 스스로 판정할 수 있다.
  \item \textbf{하드웨어}: 모든 실습은 소형 GPU 한 개면 충분하다. 무료 Colab T4로도 가능하며 CPU만 사용하는 대안도 제공한다. 보유하지 않은 클러스터가 필요한 실습은 없다.
\end{tbbitemize}
\end{calloutbox}

\vspace{1.5pt}

\subsection*{"시스템SW가 돈이다"}

얼마 전 한 대형 기술 기업의 최고경영자가 이 모든 내용을 네 단어로 요약했다. 이니셜은 S다. 한국 독자라면 누구인지 맞히는 재미가 있을 것이다. \textbf{"시스템SW가 돈이다."} 1969년 IBM도 같은 결론에 도달했으니 훌륭한 선배가 있는 셈이다. 이 글을 쓰는 현재 세계에서 가장 가치 있는 기업은 명함상 반도체 회사다. 하지만 그 플랫폼을 떠나 보려 한 사람에게 물으면 진짜 해자가 CUDA임을 알게 된다. 시스템 소프트웨어는 언제나처럼 보이지 않는 곳에서 조용히 대금을 거둔다.

독자에게 바라는 것은 두 가지다. 첫째, 이 책을 다 읽은 뒤에는 보이지 않는 계층을 \textit{볼 수} 있기를 바란다. 챗봇이 0.5초 만에 답하는 모습을 보면서 그 뒤의 하이퍼바이저, cgroup, 스케줄러, 페이지드 캐시, 청구서를 함께 떠올릴 수 있기를 바란다. 둘째, 여러분 세대의 21시 04분이 찾아왔을 때, 그리고 그 순간은 반드시 찾아온다, 그 방 안에서 어느 구간이 병목인지 이미 아는 사람이 바로 여러분이기를 바란다. 1960년대에는 그런 시야가 수많은 경력을 만들어 냈다. 2020년대에는 한 층 위에서 같은 일이 다시 벌어지고 있다.

아직 답이 없던 시절 이 질문들을 베타 테스트해 준 폐강 과목의 두 기수 학생들에게 감사한다. 내가 졸라맨을 그리는 동안 쉬지 않고 타이핑한 에이전트들에게도 감사를 전한다.

\textcolor{tbbink}{\textbf{여러분의 서비스가 지루할 만큼 평온하기를.}}

{\raggedleft \textcolor{tbbink}{\textbf{의성}}\textcolor{tbbgray}{\textit{  ·  한국 AI 혁명의 심장 판교에서, 2026년 7월}}\par}


\cleardoublepage
\pdfbookmark[0]{Contents}{toc}
{\hypersetup{linkcolor=tbbink}\tableofcontents}

\mainmatter

\part{인프라와 원리}
%% Chapter 1 — 이 책 개요: AI 서비스의 큰 그림
\setcounter{chapter}{0}
\chapter{이 책 개요: AI 서비스의 큰 그림}
\chaptersubtitle{개발 → 학습 → 서빙 → 운영 — 그리고 이 모든 것을 떠받치는 클라우드}

\begin{calloutbox}{tbbblue}{tbbbluefill}
{\normalsize\bfseries\color{tbbblue}학습 목표}\par\vspace{3pt}
\begin{tbbitemize}
  \item 엉망이 된 시연 행사와 한 달 뒤 도착한 클라우드 청구서를 통해, 잘 학습된 모델과 실제 프로덕션 환경에서 작동하는 AI 서비스 사이의 간극을 설명하고, 그 간극에서 생기는 질문 가운데 정확도에 관한 것은 하나도 없는 이유를 설명한다.
  \item 전체 AI 서비스 스택(개발 → 학습 → 서빙 → 운영)을 그리고, 각 단계를 활동·산출물·인프라 요구 사항과 연결한다.
  \item 여섯 가지 이상의 관점에서 학습 워크로드와 추론 워크로드를 비교하고, 그 차이가 인프라 설계를 어떻게 좌우하는지 논의한다.
  \item 지연 시간·처리량·비용으로 이루어진 \textbf{서빙 삼각형}을 이용해 특정 최적화가 무엇을 지불하고 무엇을 얻는지 밝히며, SLO가 이 삼각형 위에서 의도적으로 선택한 한 점인 이유를 설명한다.
  \item 추론 서버의 메모리 사용량이 평탄하지 않고 급격히 출렁이는 이유와, "관측한 유휴 사용량 + 10\%"를 기준으로 용량을 산정하면 문제가 생기는 이유를 설명한다.
  \item 전통적인 클라우드 서비스 모델(IaaS · PaaS · SaaS)을 AI 인프라 관점에서 새롭게 해석하고, GPUaaS 및 모델 API 생태계를 파악하며, 실제 결정을 좌우하는 기준과 과금 단위를 바탕으로 "직접 서빙할 것인가, API를 구매할 것인가"를 논증한다.
\end{tbbitemize}
\end{calloutbox}

\vspace{5.0pt}

\begin{calloutbox}{tbbteal}{tbbtealfill}
{\normalsize\bfseries\color{tbbteal}핵심 용어}\par\vspace{3pt}
AI 서비스 스택 · 서빙 · 추론 · 워크로드 · 지연 시간 · 꼬리 지연 시간(p95/p99) · 처리량 · 서빙 삼각형 · SLO · 사용률 · 레플리카 · IaaS/PaaS/SaaS · GPUaaS · 네오클라우드 · 서버리스 GPU · 모델 API · 토큰, 키, 할당량 · 오픈 웨이트 모델 · 자체 구축 대 외부 도입
\end{calloutbox}

\vspace{5.0pt}

\begin{calloutbox}{tbblightgray}{tbbboxfill}
{\normalsize\bfseries\color{tbbink}장 구성}\par\vspace{3pt}
\textbf{1.1} 모델에서 서비스로: 왜 "서빙 시스템"을 공부하는가    \textbf{1.2} 전체 AI 서비스 스택    \textbf{1.3} 학습 워크로드와 추론 워크로드 — 그리고 서빙 삼각형    \textbf{1.4} 클라우드 서비스 모델과 AI 인프라    \textbf{1.5} GPUaaS와 모델 API 생태계    \textbf{1.6} 이 책의 로드맵과 학습 방법 — 이후 요약, 연습문제, 더 읽을거리로 이어진다.
\end{calloutbox}

\section{모델에서 서비스로: 왜 "서빙 시스템"을 공부하는가}

이 책은 도구 하나를 가르치기에 앞서 그 도구가 왜 존재하는지부터 설명해야 한다. 따라서 정의 대신 두 가지 사례와 세계적으로 유명한 이야기 하나로 시작한다. 두 사례는 여러 경험을 합친 것이다. 어느 공학 과정에서나 해마다 비슷한 일이 벌어지지만, 여기에 제시한 숫자는 모두 산술적으로 정확하다. 이 절을 마칠 때면 직접 계산할 수 있을 것이다. 마지막 이야기는 온전히 실화다.

\subsection*{사례 A: 시연 행사}

수강생 세 명으로 구성된 팀이 한 학기 동안 이미지 캡셔닝 모델을 만들었다. 검증 정확도는 94\%다. 팀 노트북에서는 사진을 넣으면 3분의 1초 만에 캡션이 나온다. 시연 당일 21:04, 팀은 학과 Discord에 링크를 올렸다. 이어진 상황을 터미널 기록으로 재구성하면 다음과 같다.

\begin{terminalbox}
\begin{Verbatim}[breaklines=true,breakanywhere=true,fontsize=\small,formatcom=\color{tbbtermfg},codes=\tbbcodeactive]
$ curl -w "%{time_total}s\n" -o /dev/null -s localhost:8000/caption -d @cat.jpg
0.31s                          # my laptop. instant. ship it.
# 21:04  posted the demo link to the department discord
# 21:11  "it's loading for me"  "same"  "is it down?"
$ hey -z 30s -c 100 http://demo.example.edu/caption   # 100 concurrent users
Summary:
  Requests/sec:  3.2
  Latency   p50: 15.2s    p95: 28.9s    p99: 30.0s
  Error distribution:
  [504] Gateway Timeout: 61 responses
# 21:19  the server never crashed. it queued.
# accuracy is still 94%. nobody stayed to find out.
\end{Verbatim}
\end{terminalbox}
\figcaption{사례 1-A  재구성한 시연 행사. GPU 하나, 한 번에 요청 하나, 사용자 백 명. 모델은 아무 잘못도 하지 않았지만 서비스는 실패했다.}

산술로 부검해 보자. 이 이야기의 전부가 바로 그 계산에 담겨 있다. 서버는 \textbf{한 번에 요청 하나}를 처리하고, 요청마다 0.31초가 걸린다. 따라서 상한은 1 ÷ 0.31 ≈ \textbf{초당 요청 3.2건}이며, 이는 \code{hey}가 측정한 처리량(throughput)과 정확히 같다. 사용자 100명이 연결을 유지한 상태에서는 새 요청이 평균적으로 약 50개 요청 뒤에 선다. 자기 요청을 처리하는 0.31초에 앞서 50 × 0.31 ≈ \textbf{15.5초를 기다려야} 하므로 p50이 설명된다. 가장 운이 나쁜 요청은 99개 요청 뒤에 선다. 99 × 0.31 ≈ 30.7초지만, 앞단 게이트웨이(gateway)가 30초에 포기하므로 이 요청은 \textbf{504} 오류로 끝난다. 어떤 것도 충돌하지 않았고, 일반적인 의미에서 잘못 설정된 것도 없다. 초당 요청 3.2건이라는 단단한 상한을 지닌 서비스에 동시 사용자 100명이 몰리자, 나머지는 대기열 이론이 결정했다. (이 짧은 상한 및 대기열 계산은 7주 차에 리틀의 법칙(Little's Law)이라는 이름으로 정식으로 다시 다루며, 9주 차와 12주 차에는 실제 서버군의 규모를 산정하는 데 사용한다.)

팀이 무엇을 언제 알았는지도 눈여겨보자. 문제가 생겼다는 첫 신호는 \textbf{21:11에 채팅으로 들어온 친구들의 불평}이었다. "200 OK" 줄이 찍히는 노트북 터미널은 모니터링이 아니므로, 서비스는 7분 동안 소리 없이 실패했다. 그리고 21:20에 팀이 한 말에도 주목하자. 이미 짐작했을 것이다. \textit{"하지만 내 컴퓨터에서는 되는데."} 이 말을 기억해 두자. 제4장은 한 계층 더 아래에서 이 상황 전체를 극적으로 재현하며 시작한다.

\subsection*{사례 B: 청구서는 어김없이 도착한다}

쓴맛을 본 팀은 당연한 다음 수를 둔다. 클라우드에서 제대로 된 GPU 머신을 빌려 다시 배포하고, 그 금요일 지도교수 앞에서 시연을 성공적으로 마친다. 이후 모두 다른 마감 일정으로 넘어갔지만 머신은 그러지 않았다. 월말이 되자 클라우드 콘솔에 다음 내역이 나타났다.

\begin{terminalbox}
\begin{Verbatim}[breaklines=true,breakanywhere=true,fontsize=\small,formatcom=\color{tbbtermfg},codes=\tbbcodeactive]
# Billing > Bills > October (excerpt)
Amazon EC2   g5.xlarge (1x NVIDIA A10G, 24 GB)   us-east-1
   744 hrs x $1.006/hr                        $748.46
EBS storage + data transfer                    $31.20
                                     TOTAL    $779.66
# total time the demo was actually used, all month: ~90 minutes.
# utilization: 1.5 hrs / 744 hrs = 0.2%. the meter does not care.
\end{Verbatim}
\end{terminalbox}
\figcaption{사례 1-B  90분 사용한 시연 서비스의 월말 청구서. 상시 실행되는 GPU는 요청 도착 여부와 관계없이 한 달 744시간의 요금이 청구된다.}

이 사례를 부검하는 데 대기열 이론은 필요 없다. 곱셈 한 번이면 된다. GPU 인스턴스에는 \textbf{존재하는 모든 시간에 대해 시간 단위로} 요금이 부과된다. 10월에는 24 × 31 = 744시간이며, 실제로 추론이 일어났는지는 계량기가 전혀 신경 쓰지 않는다. 사용률 0.2\%는 팀이 \textbf{실제 사용 시간 1분당 약 \$8.66}를 냈다는 뜻이다. 두 사례는 서로 거울상이다. 사례 A는 수요가 프로비저닝한 용량을 초과할 때 벌어지는 일이고, 사례 B는 프로비저닝한 용량이 수요를 초과할 때 치르는 비용이다. 트래픽의 호흡에 맞춰 이 둘 사이를 지속적이고 자동으로 조절하는 일을 오토스케일링이라 한다. 12주 차에 다룰 이 작업은 말처럼 쉽지 않다. GPU 서빙 인스턴스가 다시 준비되는 데는 밀리초가 아니라 몇 분이 걸리기 때문이다(\textit{콜드 스타트} 문제). 지금은 이 두 사례를 이 책의 양극으로 기억해 두자. \textbf{너무 느리거나}, \textbf{너무 비싸다}.

\begin{calloutbox}{tbbamber}{tbbamberfill}
{\normalsize\bfseries\color{tbbamberdark}21:04부터 월말까지 생겨난 질문}\par\vspace{3pt}
\begin{tbbitemize}
  \item \textbf{동시성} — 사용자 100명이 동시에 요청을 보내면 어떻게 되는가? 사례 A에서는 요청들이 3.2 req/s라는 상한 뒤에 대기열을 만들었다. 같은 조건의 서비스라면 어떻게 되는가?
  \item \textbf{지연 시간} — 노트북에서 추론에는 0.31 s가 걸렸다. 그런데 중간 사용자는 왜 15 s를 기다렸는가? 왜 \textit{가장 나쁜} 사용자 경험(p99)이 중요한 수치인가?
  \item \textbf{가용성} — 오전 3시에 서버 프로세스가 죽으면 누가 다시 시작하는가? 그동안 사용자에게는 무엇이 보이는가?
  \item \textbf{비용} — "밤새 GPU를 계속 실행할 것인가?"라는 질문에 사례 B는 "기본적으로 그렇다. \$1.006/hr를 내면서"라고 답했다. 어떤 답을 내릴 것인가?
  \item \textbf{배포} — 더 나은 모델이 준비되었다. 서비스를 중단하지 않고 어떻게 교체하며, 모델이 오작동하면 어떻게 롤백하는가?
  \item \textbf{관측 가능성(observability)} — 팀의 모니터링 시스템은 \textit{채팅으로 불평하는 친구들}이었다. 무엇이어야 했는가?
\end{tbbitemize}
\end{calloutbox}

\vspace{3.0pt}

이 여섯 질문의 공통점은 무엇인가? \textbf{모델 정확도에 관한 질문이 하나도 없다.} 모두 시스템에 관한 질문이다. 정확도 94\%인 모델을 96\%로 높이는 일은 머신러닝 문제지만, 위의 여섯 질문에 답하는 일은 운영체제·네트워킹·분산 시스템·클라우드 인프라의 영역이다. 이 책이 다루는 주제가 바로 후자다.

\subsection*{가장 유명한 서빙 이야기}

앞의 사례가 학생 프로젝트 규모로 느껴진다면 크기를 10⁷배로 키워 보자. 그러면 지난 10년간 가장 큰 파장을 일으킨 제품 출시가 된다. \textbf{2022년 11월 30일}, OpenAI는 ChatGPT를 공개하며 겸손하게 "연구 프리뷰"라고 소개했다. 닷새 뒤 사용자는 백만 명이 되었다. 1월에는 월간 사용자가 1억 명으로 추산되었고, 당시 분석가들은 역사상 가장 빠르게 성장한 소비자 애플리케이션이라고 불렀다(널리 퍼진 비교 도표에 따르면 Instagram이 사용자 백만 명에 도달하는 데 약 2.5개월, Netflix는 3.5\textit{년}이 걸렸다). 첫겨울 내내 가장 많이 공유된 화면은 영리한 답변이 아니라 오류 페이지였다. 모델이 자신의 과부하를 소재로 지은 5행시와 함께 \textbf{"ChatGPT의 사용량이 현재 한도에 도달했습니다"}라는 문구가 떴다. 행성 규모로 커진 사례 A였다. 프로비저닝한 용량을 수요가 초과하는 광경이 1억 명 앞에서 펼쳐졌다.

이 이야기가 머신러닝 과정이 아니라 이 장에 속하는 결정적 이유는 \textbf{모델 자체가 뉴스가 아니었다}는 점이다. GPT-3 계열은 이미 2년 넘게 API로 공개되어 있었고, ChatGPT 출시 모델의 기반 역량은 개발자들이 이미 이용할 수 있던 모델과 본질적으로 같았다. 11월 30일에 출시된 것은 \textit{패키징과 서빙}이었다. 채팅 인터페이스, 무료 접근, 그리고 지구상의 모든 사람 앞에 대규모 언어 모델을 동시에 내놓으려는 인프라의 야심이었다. "API 뒤에 그 역량이 존재한다"와 "우리 부모님이 쓴다" 사이의 간극은 이후 기업 가치로 수천억 달러에 이르는 것으로 드러났다. 그 간극은 바로 위 상자 속 여섯 질문을 행성 규모로 던져 얻은 결과다. 사례 B 역시 함께 커졌다. 며칠 만에 OpenAI CEO는 컴퓨팅 비용이 "눈이 튀어나올 정도"라고 공개적으로 인정했고, 두 달 뒤 월 \$20의 Plus 요금제가 등장했다. 서빙 비용은 결코 각주가 아니다. 규모가 커지면 사업 모델 자체를 결정한다.

\subsection*{모델은 서비스의 작은 일부다}

이 패턴은 LLM보다 오래되었으므로 시스템 엔지니어에게는 전혀 놀라운 일이 아니었다. 이를 가장 유명하게 표현한 자료가 2015년 Google 논문 "Hidden Technical Debt in Machine Learning Systems"다. 이 논문의 상징적인 그림에서 전체 프로덕션 ML 시스템을 그리면, 중앙의 "ML code"는 아주 작은 검은 상자이고 그 주위를 훨씬 큰 상자들이 둘러싼다. 데이터 수집, 특징 추출, 서빙 인프라, 모니터링, 자원 관리, 구성이다. 논문의 표현에 따르면 성숙한 시스템에서 ML 코드는 "많아야 5\%"이며, 나머지 95\%는 배관, 즉 시스템이다.

10년 뒤 대규모 언어 모델 시대가 되자 이 비율은 완화되기는커녕 더 극단적으로 변했다. 이제 많은 팀은 모델을 전혀 학습하지 않고 공개된 사전 학습 모델을 가져다 쓴다. 모델 자체가 "주어진 것"이 되면 팀을 가르는 기준은 \textbf{그 모델을 얼마나 빠르고 신뢰성 높고 저렴하게 서비스로 바꾸는가}다. 같은 오픈 소스 모델을 서빙하는 두 팀의 결과가 천양지차일 수 있다. 한 팀은 응답에 4초가 걸리고 한 달에 \$8,000를 쓰지만, 다른 팀은 그 5분의 1 비용으로 0.5초 만에 응답한다. 이 격차를 만드는 지식이 이 책의 내용이다.

\subsection*{이 책의 관점: 원리에서 시작한다}

서빙 시스템을 가르치는 방법은 두 가지다. 하나는 도구부터 시작하는 방식이다. Docker 명령과 Kubernetes YAML을 빠르게 익혀 무언가를 배포한다. 다른 하나는 원리부터 시작하는 방식이다. 도구를 만지기 전에 그 도구를 떠받치는 커널 메커니즘을 이해한다. 이 책은 두 번째 길을 택한다. 첫 6주 동안 하이퍼바이저, Linux 네임스페이스, cgroup, 컨테이너 런타임을 커널 수준까지 파고드는 데에는 두 가지 이유가 있다.

첫째, \textbf{도구는 바뀌지만 원리는 남는다.} 졸업할 때쯤이면 Kubernetes보다 더 유행하는 오케스트레이터(orchestrator)가 등장할 수도 있다. 하지만 그 도구 역시 네임스페이스로 격리를 구현하고 cgroup으로 자원을 제한할 것이다. 원리를 알면 새로운 도구는 모두 "이미 아는 개념을 감싼 새로운 래퍼"일 뿐이다.

둘째, \textbf{추상화가 새는 곳에서 장애가 발생한다.} 평소에는 컨테이너를 "경량 VM"처럼 생각해도 문제없다. 그러나 추론 서버가 눈에 보이는 원인 없이 계속 죽고 로그에는 수수께끼 같은 \code{exit code 137}만 남는 날에는 이야기가 달라진다. 컨테이너의 메모리 한도가 커널 cgroup이며, 프로세스가 그 한도를 넘는 순간 커널의 OOM 킬러가 신호 9인 SIGKILL을 보내고, 이것이 128 + 9 = 137이 된다는 사실을 아는 사람만이 로그 한 줄 남기지 않고 서버가 사라진 이유를 설명하고 고칠 수 있다. 3주 차 실습에서는 의도적으로 자기 프로세스를 정확히 이런 방식으로 죽인다. 이후 이 책 내내 숫자 137이 유용하게 따라다닐 것이다.

요컨대 이 책은 AI를 \textit{만드는} 법이 아니라 \textbf{이미 만들어진 AI를 실행하는 법}을 다룬다. 학습 알고리즘은 선수 과목인 머신러닝과 딥러닝의 영역이다. 누군가 학습이 끝난 모델 파일을 건네는 순간부터 이야기를 이어받아, 그 파일이 수만 명이 사용하는 안정적인 서비스가 되기까지의 모든 것을 책임진다.

\subsection*{책 전반에서 사용할 용어}

본격적으로 시작하기 전에 이 책에서 되풀이해 사용할 용어를 정리하자. 지금은 대략적인 감만 잡으면 충분하며, 각 용어는 해당 주차에서 자세히 다룬다.

\begin{tbbtablewrap}
  \begin{tabular}{@{}>{\raggedright\arraybackslash}m{\dimexpr 0.2447\tbltextwidth-2\tabcolsep\relax} >{\raggedright\arraybackslash}m{\dimexpr 0.7553\tbltextwidth-2\tabcolsep\relax}@{}}
  \arrayrulecolor{tbbnavy}\specialrule{1pt}{0pt}{0pt}
  \rowcolor{tbbtablehead}{\centering \bfseries\color{tbbnavy}용어} & {\centering \bfseries\color{tbbnavy}의미} \\
  \arrayrulecolor{tbbnavy}\specialrule{0.7pt}{0pt}{2pt}
  {\textbf{모델}} & {학습의 산출물이다. 물리적으로는 네트워크 아키텍처와 수억 개에서 수천억 개에 이르는 파라미터(가중치)를 담은 파일이다.} \\
  \arrayrulecolor{tbbtableborder}\hline
  {\textbf{추론}} & {학습된 모델에 입력을 통과시켜 출력을 계산하는 일이다. 학습과 달리 가중치를 절대 바꾸지 않으며, 순전파만 수행한다.} \\
  \arrayrulecolor{tbbtableborder}\hline
  {\textbf{서빙}} & {네트워크를 통해 추론을 서비스 형태로 사용자에게 제공하는 일이다. 추론 + 요청 처리 + 스케일링 + 운영으로 이루어진 시스템 문제다.} \\
  \arrayrulecolor{tbbtableborder}\hline
  {\textbf{엔드포인트}} & {클라이언트가 추론을 요청할 때 호출하는 네트워크 주소(URL)다. 예: \code{POST /\allowbreak v1/\allowbreak predict}.} \\
  \arrayrulecolor{tbbtableborder}\hline
  {\textbf{워크로드(workload)}} & {인프라에서 실행되는 작업의 성격과 부하 패턴을 두루 가리키는 말이다. "학습 워크로드"와 "추론 워크로드"처럼 쓴다.} \\
  \arrayrulecolor{tbbtableborder}\hline
  {\textbf{지연 시간}} & {요청이 도착한 시점부터 응답이 나가는 시점까지 걸린 시간이다. 사용자 경험을 좌우하며, 평균이 아니라 꼬리(p95/p99)를 기준으로 관리한다.} \\
  \arrayrulecolor{tbbtableborder}\hline
  {\textbf{처리량}} & {단위 시간에 처리하는 요청(또는 토큰)의 수다. 사례 A의 상한인 3.2 req/s가 처리량이다. 자원 효율과 비용을 좌우한다.} \\
  \arrayrulecolor{tbbtableborder}\hline
  {\textbf{레플리카(replica)}} & {로드 밸런서(load balancer) 뒤에서 실행되는 동일한 서빙 인스턴스 사본 여러 개 중 하나다. 레플리카가 많을수록 처리량과 가용성, 비용이 모두 늘어난다.} \\
  \arrayrulecolor{tbbtableborder}\hline
  {\textbf{사용률}} & {비용을 지불한 자원이 유용한 작업을 수행하는 비율이다. 사례 B에서는 0.2\%였다. 서빙 경제성을 좌우하는 핵심 지렛대다.} \\
  \arrayrulecolor{tbbtableborder}\hline
  {\textbf{SLO}} & {서비스 수준 목표(Service Level Objective)로, 스스로 설정하는 정량적 품질 목표다. 예: "p99 지연 시간 ≤ 500 ms, 가용성 ≥ 99.9\%." 7주 차에서 자세히 다룬다.} \\
  \arrayrulecolor{tbbnavy}\specialrule{1pt}{2pt}{0pt}
  \end{tabular}\par
  \figcaption{표 1-1  이 책 전반에서 사용하는 기본 용어}
\end{tbbtablewrap}

\begin{calloutbox}{tbbpurple}{tbbpurplefill}
{\normalsize\bfseries\color{tbbpurple}생각해 보기}\par\vspace{3pt}
\begin{tbbitemize}
  \item 두 가지 조건을 바꾸어 사례 A를 다시 계산해 보자. 팀은 로드 밸런서 뒤에 레플리카 4개를 실행하며, 약간 느린 클라우드 GPU에서는 요청 하나에 0.4 s가 걸린다. 새로운 처리량 상한은 얼마인가? 동시 사용자 100명일 때 p50은 대략 얼마가 되는가? 충분한가?
  \item 사례 B에서 팀이 90분 사용에 \$748를 내지 않을 수 있었던 서로 다른 방법 세 가지를 나열해 보자. 이미 최소 세 가지는 생각할 수 있다(끄기, 어떤 방식으로든 초 단위로 빌리기, 다른 팀과 공유하기). 각 답은 이후 어느 주차의 내용을 미리 보여 주는가?
  \item ChatGPT의 기반 모델 역량은 ChatGPT 출시 2년 전부터 공개되어 있었다. 과거 기술 가운데 같은 형태의 사례를 하나 들어 보자. 역량은 제품보다 훨씬 먼저 존재했고, 제품의 성취는 패키징과 인프라에 있었다. (정답은 많다.)
\end{tbbitemize}
\end{calloutbox}

\section{전체 AI 서비스 스택: 개발 → 학습 → 서빙 → 운영}

아이디어에서 시작해 일상적인 운영에 이르는 AI 서비스의 생애는 네 단계로 나눌 수 있다. 아이디어와 데이터에서 출발하는 \textbf{개발}, 모델을 만들어 내는 \textbf{학습}, 모델을 사용자 앞에 내놓는 \textbf{서빙}, 서비스가 살아 있는 동안 건강하게 유지하는 \textbf{운영}이다. 그림 1-1은 네 단계와 그 아래의 공통 기반인 클라우드 인프라를 함께 보여 준다.

\begin{tbbfigure}
  \adjustbox{max width=0.94\linewidth,max totalheight=0.46\textheight}{\includegraphics{figures/fig1-1_stack.png}}\\[3pt]
  \figcaption{그림 1-1  전체 AI 서비스 스택. 네 단계는 직선이 아니라 피드백으로 연결된 순환 고리를 이루며, 모든 단계가 클라우드 인프라 위에서 실행된다.}
\end{tbbfigure}

\subsection*{네 단계를 빠르게 둘러보기}

\textbf{① 개발} 단계에서는 문제를 정의하고, 데이터를 수집·정제하며, 모델 아키텍처와 학습 레시피를 실험한다. 주로 Jupyter 노트북과 실험 추적기를 사용한다. 인프라 관점에서 보면 작은 GPU를 짧고 빈번하게 사용하는 대화형 패턴이다. 이 단계의 산출물은 정제된 데이터셋과 재현 가능한 학습 파이프라인이다.

\textbf{② 학습} 단계에서는 그 파이프라인을 실제로 실행한다. 대형 모델의 경우 수십 개에서 수천 개 GPU에 걸친 분산 학습이 며칠 동안 이어지며, 장애에 대비해 중간중간 체크포인트를 저장한다. 이 단계의 산출물이자 이 책의 출발점은 \textbf{학습된 모델 가중치}다. 실제로는 체크포인트나 \code{safetensors} 같은 형식으로, 몇 기가바이트에서 수백 기가바이트에 이르는 파일이다. 모델 형식은 7주 차에 다룬다.

\textbf{③ 서빙} 단계에서는 모델 파일이 실제 트래픽을 처리하는 서비스가 된다. 모델을 추론 서버와 컨테이너로 감싸 클라우드에 배포하고 엔드포인트로 공개한다. 양자화와 배칭 같은 최적화로 지연 시간과 비용을 관리하며, 트래픽 변화에 맞춰 인스턴스 수를 조절한다. 이 책의 중반부 전체(7\textasciitilde{}12주 차)가 이 단계에 해당한다. 사례 A의 팀에 도움이 필요했던 곳도 바로 여기다.

\textbf{④ 운영} 단계는 서비스가 살아 있는 한 끝나지 않는다. 지연 시간, 처리량, 오류율, GPU 사용률을 계속 관찰하고, 새 모델 버전을 안전하게 배포하는 자동화 파이프라인을 구축한다. 또한 입력 분포가 모델의 학습 분포에서 벗어날 때 조용히 성능이 저하되는 \textbf{드리프트}를 감시한다. 13\textasciitilde{}15주 차의 주제다.

중요한 점은 이 네 단계가 \textbf{폭포형태가 아니라 순환고리(cycle)형태}라는 사실이다. 운영 과정에서 수집한 사용자 데이터와 실패 사례가 다시 개발과 학습으로 흘러 들어가 다음 모델 버전을 만들고, 새 모델은 다시 서빙 단계를 거친다. 조직이 성숙할수록 이 고리는 더 자동화되고 더 빠르게 돈다. 이 고리를 원활하게 돌리는 문화와 도구를 \textbf{MLOps}라 하며, 13주 차에 다룬다.

\subsection*{서빙 단계 확대하기}

이 책의 중심인 서빙 단계를 확대해 보자. 그림 1-2는 사용자 요청 하나가 도착해 응답으로 돌아오기까지의 경로를 추적한다. 요청은 먼저 API 게이트웨이(API gateway)와 로드 밸런서를 통과한다. 사례 A에서 뒤쪽 대기열이 인내 시간을 넘기자 504를 반환한 구성 요소가 바로 이들이다. 이곳에서 인증을 거친 요청은 여러 서빙 \textbf{레플리카} 중 하나로 전달된다. 레플리카 안에서 모델 실행은 전처리와 후처리 사이의 한 구간일 뿐이다. 모델 가중치는 시작할 때 \textbf{객체 스토리지(object storage) 기반 모델 저장소}에서 불러오며, 모니터링 시스템은 이 경로 곳곳에서 메트릭을 수집한다. 오토스케일러(autoscaler)는 이 근거를 이용해 레플리카 수를 조절한다.

\begin{tbbfigure}
  \adjustbox{max width=0.94\linewidth,max totalheight=0.46\textheight}{\includegraphics{figures/fig1-2_request_path.png}}\\[3pt]
  \figcaption{그림 1-2  온라인 추론 요청의 여정. 사용자가 체감하는 지연 시간은 모델 실행 시간이 아니라 전체 경로의 합이다.}
\end{tbbfigure}

이 그림에서 두 가지는 반드시 기억해야 한다. 첫째, \textbf{사용자가 체감하는 지연 시간은 전체 경로의 합이다.} GPU에서 모델 실행이 50 ms밖에 걸리지 않아도 요청이 대기열에서 2초를 기다렸다면 사용자는 느린 서비스를 경험한 것이다. 다른 해석의 여지가 없다. 사례 A는 극단적인 예였다. 30초의 경험 가운데 모델 시간이 0.31 s에 불과했다. 따라서 서빙 시스템의 성능 분석은 언제나 "어느 구간이 병목인가?"라는 질문이다. 둘째, 그림의 모든 구성 요소가 이 책 어딘가에서 다시 등장한다. 로드 밸런서와 객체 스토리지는 6주 차, 서빙 레플리카 내부의 배칭은 9\textasciitilde{}11주 차, 오토스케일러는 12주 차, 모니터링은 14주 차에 다룬다. 그림 1-2를 이 책의 지도라고 생각해도 좋다.

\begin{calloutbox}{tbbpurple}{tbbpurplefill}
{\normalsize\bfseries\color{tbbpurple}생각해 보기}\par\vspace{3pt}
\begin{tbbitemize}
  \item 그림 1-2의 지연 시간 구간을 직접 제어할 수 있는 것과 대부분 제어할 수 없는 것으로 나누어 보자. (힌트: 사용자의 네트워크는 누가 제어하는가?)
  \item 모델 가중치를 컨테이너 이미지에 넣지 않고 객체 스토리지에서 불러오는 이유는 무엇인가? 모델이 50 GB라면 어떤 문제가 생기는가? (6주 차에 답을 확인한다.)
\end{tbbitemize}
\end{calloutbox}

\section{학습 워크로드와 추론 워크로드}

같은 GPU와 같은 모델을 사용하더라도 인프라 관점에서 "학습"과 "추론 서빙"은 완전히 다른 존재다. 이 차이를 정확히 이해하는 것이 이 장의 가장 중요한 목표다. 서빙에서 작은 인스턴스를 여러 개 두어 수평확장(horizontal scale-out)하는 이유, 오토스케일링(autoscaling)이 필수인 이유, GPU 사용률이 골칫거리인 이유 등 이 책 전반에 등장하는 설계 결정이 모두 여기서 비롯되기 때문이다. 그림 1-3은 여섯 가지 관점에서 둘의 차이를 펼쳐 보인다.

\begin{tbbfigure}
  \adjustbox{max width=0.94\linewidth,max totalheight=0.46\textheight}{\includegraphics{figures/fig1-3_train_vs_infer.png}}\\[3pt]
  \figcaption{그림 1-3  학습 워크로드와 추론 워크로드. 아래쪽의 GPU 사용률 대비는 이 책 후반의 오토스케일링과 GPU 공유 주제를 예고한다.}
\end{tbbfigure}

\subsection*{작업과 서비스}

가장 근본적인 차이는 실행 형태다. 학습은 \textbf{시작과 끝이 있는 배치 작업}이다. 3일이 걸리든 3주가 걸리든, 끝나면 자원을 반납한다. 서빙은 \textbf{끝없이 상시 실행되는 서비스}다. 언제 도착할지 모르는 요청에 늘 응답할 준비가 되어 있어야 한다. 이 차이가 인프라를 바꾼다. 학습 플랫폼의 핵심 구성 요소는 작업 대기열과 배치 스케줄러다. 반면 서빙 플랫폼의 핵심은 로드 밸런서, 서비스가 계속 실행되도록 보장하는 오케스트레이터(Kubernetes Deployment가 하는 일이며 5주 차에 다룬다), 장애 자동 복구다. 사례 B가 벌어진 이유도 여기에 있다. \textit{작업}은 끝나면 머신을 반납하지만, \textit{서비스}에는 끝이 없으므로 머신을 기본적으로 끄는 스위치도 없다.

\subsection*{처리량 중심과 지연 시간 중심}

학습의 성능 지표는 처리량이다. 에포크(epoch)당 몇 시간이 걸리는지, 초당 몇 개 샘플을 처리하는지를 본다. 전체 실행이 더 일찍 끝난다면 개별 샘플에 0.1 s가 걸리는지 1 s가 걸리는지는 중요하지 않다. 따라서 학습은 GPU가 한 번에 처리하는 샘플 묶음인 \textbf{배치}를 \textbf{가능한 한 크게} 만들어 하드웨어를 포화시킨다.

서빙에서는 성능지표가 요청별 지연 시간으로 바뀌면서 근본적인 고민이 생긴다. GPU 관점에서는 요청을 더 큰 배치로 모을수록 처리량이 올라가고 비용이 낮아진다. 그러나 먼저 도착한 요청은 배치가 찰 때까지 \textbf{기다려야} 하므로 지연 시간이 나빠진다. 이 줄다리기는 서빙 시스템 설계 전체를 관통한다. 이를 영리하게 해결하는 동적 배칭(10주 차)과 연속 배칭(11주 차)은 이 책의 핵심 내용이다.

한 가지 더 기억하자. 서빙의 지연 시간은 평균이 아니라 \textbf{꼬리(tail)}로 관리한다. 평균 응답 시간이 100 ms여도 사용자 100명 중 한 명이 5초를 기다린다면, 하루 요청 100,000건을 처리하는 서비스는 매일 1,000번의 나쁜 경험을 제공한다. 따라서 실무에서는 모든 요청을 빠른 순서부터 느린 순서로 늘어놓았을 때 95번째와 99번째 백분위에 해당하는 지연 시간인 p95와 p99를 목표로 삼는다. "p99 ≤ 500 ms"는 요청의 99\%에 0.5초 안에 응답하겠다는 약속이다. 사례 A의 p99는 30초였으며, 모델이 아니라 대기열이 결정했다. 서비스 품질이 평균이 아니라 꼬리에 존재하는 이유가 바로 이것이다. 7주 차에는 이를 SLO로 발전시킨다.

\subsection*{서빙 삼각형: 지연 시간, 처리량, 비용}

배칭의 줄다리기는 이 책의 중심을 이루는 패턴의 한 사례이며, 별도의 이름과 그림을 붙일 만하다. 모든 서빙 시스템에는 세 가지가 요구된다. \textbf{낮은 지연 시간}(각 사용자에게 빠르게 응답), \textbf{높은 처리량}(초당 많은 사용자에게 응답), \textbf{낮은 비용}(응답 한 번에 적은 비용)이다. 엔지니어링의 현실은 \textit{어느 하나를 지불하면 나머지 둘을 살 수 있지만}, 어떤 묘수로도 셋을 한꺼번에 얻을 수는 없다는 것이다. 그림 1-7은 이 책의 나침반이다.

\begin{tbbfigure}
  \adjustbox{max width=0.94\linewidth,max totalheight=0.46\textheight}{\includegraphics{figures/fig1-7_triangle.png}}\\[3pt]
  \figcaption{그림 1-7  서빙 삼각형. 이 책의 모든 최적화 기법은 이 삼각형 위에서 움직이는 일이며, 얻을 것을 취하기 전에 무엇을 지불하는지 밝히는 태도가 중요하다.}
\end{tbbfigure}

세 변을 한 바퀴 돌아보자. \textbf{빠르면서 대용량}인 서비스는 쉽게 만들 수 있다. 레플리카를 많이 실행하고 배치는 작게 유지하며 모든 곳에 여유 용량을 둔다. 대신 \textbf{비용}을 치른다(사례 B는 이 꼭짓점의 경고담이다. 아무도 쓰지 않는 용량이었다). \textbf{대용량이면서 저렴한} 서비스도 쉽다. 공격적으로 배칭하고 모든 GPU를 포화 상태에 가깝게 실행한다. 대신 \textbf{지연 시간}을 치른다. 파이프가 가득 차면 대기열이 생기고, 대기열은 꼬리에 자리하기 때문이다(사례 A가 의도치 않게 이 경우였다). \textbf{빠르면서 저렴한} 서비스도 트래픽이 늘기 전까지는 가능하다. 작고 준비된 인스턴스 하나가 적은 요청에 즉시 저렴하게 응답하다가 트래픽이 증가하는 순간 무너진다. 주어진 엔지니어링 역량으로 세 꼭짓점을 동시에 얻는 선택지는 \textit{없다}. 이 책이 가르치는 바를 한 문장으로 표현하면 \textbf{삼각형 자체를 바꾸는 방법}이다. 양자화(10주 차)는 요청 하나를 더 저렴하면서도 \textit{더} 빠르게 만들고, 연속 배칭(11주 차)은 단순 배칭보다 적은 지연 시간 비용으로 처리량을 얻으며, 오토스케일링(12주 차)은 트래픽이 많을 때만 비용의 꼭짓점으로 잠시 위치하게 한다.

이 삼각형에서 두 가지 규칙이 따라 나오며, 반드시 지켜야 한다. 첫째, \textbf{SLO는 바람이 아니라 선택한 점}이다. "200 req/s에서 p99 ≤ 500 ms, 월 \$900 이하"는 삼각형 안의 한 지점을 구체적으로 지목하고 지키겠다고 약속한다. 팀 프로젝트 제안서(7주 차)에서 바로 이 작업을 하며, 15주 차 시연은 그 약속을 기준으로 측정한다. 둘째, \textbf{모든 최적화 주장은 대가를 밝혀야 한다.} "더 빠르게 만들었다"는 말은 "GPU-hours가 X\% 늘어나는 대가로" 또는 "8보다 큰 배치 크기를 포기하여" 같은 설명이 뒤따라야 비로소 엔지니어링이 된다. 뒤 장에서 \textit{제1장의 비용 관점으로 읽으라}고 할 때 말하는 관점이 바로 이것이다. 어느 꼭짓점을 사기 위해 어느 꼭짓점으로 값을 지불하며, 그 교환은 그만한 가치가 있는가?

\subsection*{일정한 부하와 급증하는 부하}

학습 작업의 부하는 예측할 수 있다. 직접 시작한 작업이 직접 선택한 배치 크기로 실행된다. 서빙 워크로드의 부하는 사용자가 결정한다. 점심시간에 급증하고 밤에는 조용해지며, 마케팅 행사나 입소문을 탄 게시물 때문에 10배로 뛸 수도 있다. 사례 A의 Discord 메시지는 소규모 바이럴 게시물이었고, ChatGPT 출시는 그 완전한 규모의 사례였다. 이 특성은 서빙 인프라의 고전적인 딜레마를 만든다. \textbf{최대치에 맞춰 프로비저닝하면 한산할 때 돈을 낭비하고}(트래픽이 10분의 1인 오전 4시에도 GPU 10개가 돌아간다. 서버군 규모의 사례 B다), \textbf{평균에 맞춰 프로비저닝하면 최대치에서 서비스가 무너진다}(사례 A). 해결책은 부하에 맞춰 자원을 늘리고 줄이는 오토스케일링이며, 논리적 극단에는 요청이 없을 때 인스턴스를 0개로 만드는 0으로 스케일링(scale-to-zero)이 있다. 그러나 GPU 서빙 인스턴스가 준비되려면 수십 초에서 몇 분이 걸린다. 모델 로딩이 느린 부분이다. 이것이 \textbf{콜드 스타트(cold start)} 문제이며, 12주 차에 이 문제를 만날 때 2주 차의 Firecracker microVM에서 해결책을 찾는다.

\subsection*{장애가 지니는 무게}

학습 중 GPU가 죽으면 성가시지만 재앙은 아니다. 마지막 체크포인트부터 재개하면 된다. 잃는 것은 시간과 돈이다. 서빙 장애는 성격부터 다르다. \textbf{바로 지금 서비스를 사용하는 사람들이 즉시 피해를 본다.} 오류 페이지를 보거나 로딩 표시만 바라보다 떠난다. 따라서 서빙 시스템은 단일 장애점(single point of failure)을 없애는 데 집착한다. 항상 레플리카를 둘 이상 실행하고, 죽은 인스턴스를 자동으로 교체하며(Kubernetes의 자가 치유, 5주 차), 새 버전 출시 때는 트래픽의 작은 비율만 새 버전에서 처리하기 시작하여 점진적으로 새 버전으로 모두 이전한다(카나리 배포, 13주 차).

\subsection*{자원 사용의 형태 — 시간에 따른 형태까지}

메모리 요구량은 근본부터 다르다. 학습은 가중치뿐 아니라 그래디언트, 옵티마이저 상태(Adam에서는 가중치의 두 배), 중간 활성값(activations)까지 GPU 메모리에 보관해야 한다. 추론에는 가중치와 현재 요청의 활성값만 필요하다. 경험칙상 같은 모델에서 추론에 필요한 메모리는 학습의 약 3분의 1에서 4분의 1이다. 가중치의 수치 정밀도를 낮추는 \textbf{양자화}(10주 차)를 적용하면 더 줄어든다. 따라서 최신 대형 GPU가 없어도 추론을 실행할 수 있으며, 이 책의 실습도 T4 한 대나 Colab 세션 하나에 맞도록 설계했다.

그러나 "메모리를 덜 쓴다"를 "메모리 사용량이 일정하다"로 받아들이면 안 된다. 추론 서버의 메모리 사용량에는 \textbf{시간에 따른 형태}가 있고, 그 형태는 뾰족하게 출렁인다. 가중치는 한 번 로드되어 계속 상주하는 평평한 기준선이다. 그 위에 처리 중인 요청마다 활성값용 작업 메모리를 할당하며, 이 할당량은 \textbf{동시에 처리하는 요청 수(배치)와 각 입력의 크기}에 비례한다(LLM에서는 시퀀스 길이와 요청마다 커지는 KV 캐시라는 구조가 중요하며, 11주 차의 주인공이다). 유휴 상태에서 6 GB를 쓰는 서버도 긴 입력이 한꺼번에 몰리는 순간 8 GB를 훌쩍 넘길 수 있다. 사례 A를 떠올려 보자. 프로덕션 트래픽은 바로 이렇게 급증한다. 따라서 \textit{서버를 하루 동안 유휴 상태로 관찰하고 10\%를 더한 값을 메모리 한도로 설정한다}는 흔히 쓰는 용량 산정 규칙은 신중함이 아니라 예정된 장애다. 이 한도는 조용한 모든 시간에는 유지되다가 처음으로 실제 요청이 몰리는 순간 서버를 다운시키고 로그에 \code{137}을 남긴다(요청당 몇 킬로바이트의 문자열만 사용하는 많은 무상태(stateless) 웹 앱에는 같은 규칙이 잘 통한다. 그래서 이런 습관이 생겼고, 추론에 적용하면 위험한 이유도 여기에 있다). 3주 차에는 이 서버다운 메커니즘의 실제 이름인 cgroup과 OOM 킬러를 배운다.

하드웨어 선택도 갈린다. 학습은 그래디언트 동기화로 통신량이 많으므로 NVLink와 InfiniBand로 긴밀하게 연결한 고성능 GPU(H100, B200급) 클러스터를 선호한다. 추론 레플리카끼리는 거의 통신하지 않으므로 저렴한 추론 GPU(L4, T4)를 여러 곳에 분산하는 편이 보통 더 싸다. 소형 모델은 \code{llama.cpp} 같은 런타임으로 CPU에서도 서빙할 수 있으며, GPU가 없을 때 사용하는 대안이기도 하다.

\begin{calloutbox}{tbbteal}{tbbtealfill}
{\normalsize\bfseries\color{tbbx115E59}상자 1-1  GPU를 사용하는 이유 — 유휴 GPU가 문제인 이유}\par\vspace{3pt}
딥러닝 계산의 본질은 거대한 행렬 곱셈이며, 행렬 곱셈은 동시에 실행할 수 있는 수천 개의 작고 독립적인 계산으로 나뉜다. GPU는 수천 개 코어로 정확히 이 작업을 수행하도록 만든 하드웨어다. CPU가 뛰어난 일꾼 몇 명이라면, GPU는 단순한 일꾼 수천 명으로 이루어진 여단이다.

문제는 가격이다. 데이터 센터 GPU 서버는 시간당 몇 달러가 든다(1.5절). 학습은 GPU를 95\% 이상 바쁘게 유지한 뒤 반납하지만, 사용자 트래픽을 따르는 추론 서비스는 평균 사용률이 흔히 수십 퍼센트 이하로 떨어진다. 사례 B의 0.2\%는 정도가 극단적일 뿐이다. \textbf{비싼 자원이 놀고 있는 상태}가 AI 서빙 비용 문제의 본질이며, 배칭(10주 차), GPU 공유와 MIG(2·12주 차), 오토스케일링(12주 차), 스팟 인스턴스(12주 차)는 모두 이 문제를 공략한다.
\end{calloutbox}

\vspace{3.0pt}

\subsection*{비용의 형태: 일회성 지출과 지속적 지출}

마지막은 비용 구조다. 학습 비용은 프로젝트 형태다. "이번 학습 실행에 400 GPU-hours, 약 X달러가 들었다"처럼 한 번 지불하면 끝난다. 서빙 비용은 \textbf{서비스가 살아 있는 동안 매달 반복되는 운영 비용}이다. 트래픽과 함께 늘어나며, 방치하면 조용히 부풀어 오른다. 따라서 서빙 비용은 총액이 아니라 \textbf{요청당 비용}(또는 토큰당 비용)으로 관리한다. 이 지표는 최적화의 성공 기준(10주 차)이자 팀 프로젝트의 필수 산출물인 월간 운영 비용 추정치다.

요청당 비용 계산에서 따라오는 한 가지 결과는 별도의 문장으로 강조할 만하다. 제5장이 이 개념에 크게 의존하기 때문이다. 청구서는 머신 단위로 오지만 가치는 요청 단위로 생기므로, 서빙 경제성은 \textbf{밀도}라는 게임을 끊임없이 한다. 서로 해치지 않으면서 머신 한 대를 레플리카 몇 개가 공유할 수 있는가? 각 레플리카에 넉넉한 전용 여유 용량을 주면 머신당 실행하는 레플리카 수가 적어진다. 안전하지만 비싸다. 촘촘히 넣으면 레플리카 하나의 부하 급증이 이웃 레플리카의 지연 시간 급증으로 이어진다. 따라서 레플리카당 여유 용량 1 GB는 모두 \textbf{의도적으로 선택한 밀도와 비용의 절충}이다. 자원 요청량과 한도라는 도구를 이용하면 이 절충을 강제할 수 있으며, 제3장과 제5장이 바로 이를 향해 나아간다.

정리하면, 이 책 후반부의 모든 주제는 이 절에서 살펴본 대비의 귀결이다. 서빙은 상시 실행 서비스이므로 오케스트레이션(orchestration)과 자가 치유(self-healing)가 필요하고(5주 차), 지연 시간 중심이므로 배칭과 최적화가 어렵고도 흥미로우며(10\textasciitilde{}11주 차), 부하가 급증하므로 오토스케일링과 콜드 스타트가 중요하고(12주 차), 장애의 무게가 크므로 배포 전략과 모니터링이 정교해야 한다(13\textasciitilde{}14주 차). 이 문장과 삼각형을 이 장의 요약으로 기억해 두자.

\begin{calloutbox}{tbbpurple}{tbbpurplefill}
{\normalsize\bfseries\color{tbbpurple}생각해 보기}\par\vspace{3pt}
\begin{tbbitemize}
  \item 온라인 상점의 상품 추천 모델을 생각해 보자. 이 모델의 "학습"과 "서빙"에 대해 그림 1-3의 여섯 차원별로 구체적인 수치를 그려 보자. (예를 들어 최대 트래픽은 언제인가?)
  \item 무료 취미 프로젝트, 주식 거래 신호 API, 사진 보관소를 밤새 처리하는 배치 캡셔너를 서빙 삼각형(serving triangle) 위에 배치해 보자. 각 서비스는 어느 꼭짓점을 희생하며, 어떤 SLO 수치를 적을 것인가?
  \item 어떤 추론은 실시간일 필요가 없다. 예를 들어 어제 찍은 사진을 밤새 분류할 수 있다. 이런 "배치 추론"은 그림 1-3의 어느 열과 닮았는가? 인프라에서 무엇이 단순해지는가? (7주 차 예고.)
\end{tbbitemize}
\end{calloutbox}

\section{클라우드 서비스 모델과 AI 인프라}

앞 절에서 살펴본 추론 워크로드의 특성, 즉 급증하는 부하, 상시 운영, 비싼 GPU는 자연스럽게 한 가지 질문으로 이어진다. \textbf{이 인프라는 어디에 두어야 하며, 그중 얼마만큼을 직접 관리해야 하는가?} 사례 B의 팀은 반사적으로 "VM 하나를 빌려 모든 것을 직접 설치한다"고 답했고 그 대가로 \$779를 냈다. 이 절에서는 팀이 생각하지 못했던 선택지를 정리한다.

\subsection*{클라우드를 사용하는 이유}

GPU 서버를 직접 구매해 온프레미스에서 운영하는 방식과 비교하면, AI 워크로드에서 클라우드의 장점은 특히 두드러진다. 첫째, \textbf{자본 지출이 운영 지출로 바뀐다.} 데이터 센터급 GPU 서버는 수만에서 수십만 달러가 들지만, 클라우드에서는 시간당 몇 달러에 빌리고 작업이 끝나면 떠날 수 있다. 둘째는 \textbf{탄력성}이다. 임대한 인프라만이 1.3절에서 본 급증하는 추론 부하에 맞춰 늘고 줄 수 있다. 셋째, \textbf{하드웨어 세대교체의 위험을 이전한다.} 1\textasciitilde{}2년마다 새 GPU가 출시되고 구형 GPU의 가치는 빠르게 떨어지지만, 클라우드에서는 인스턴스 유형만 바꾸면 된다. 물론 클라우드가 항상 정답은 아니다. 대규모 자원을 계속 높은 사용률로 쓰는 조직은 결국 소유하는 편이 더 저렴하며, 어떤 데이터는 건물 밖으로 반출할 수 없다. 그러므로 정확한 질문은 "클라우드를 쓸 것인가?"가 아니라 "\textbf{어느 수준까지 빌릴 것인가?}"다.

\subsection*{고전적인 세 계층: IaaS, PaaS, SaaS}

전통적으로 클라우드 서비스는 공급자가 얼마나 많이 관리해 주는지에 따라 분류한다. \textbf{IaaS}(Infrastructure as a Service)는 가상 머신, 스토리지, 네트워크 같은 원재료를 빌려 주며, OS 위의 모든 것은 사용자의 책임이다. \textbf{PaaS}(Platform as a Service)는 런타임 플랫폼까지 관리하므로 애플리케이션만 가져오면 된다. \textbf{SaaS}(Software as a Service)는 완성된 소프트웨어를 그대로 사용하는 형태다. 아래로 갈수록 제어권이 많아지고, 위로 갈수록 운영할 대상이 줄어든다.

AI 인프라는 이 세 계층을 이어받으면서 그 사이에 더 세밀한 단계를 넣는다. 그림 1-4가 이 스펙트럼을 보여 준다.

\begin{tbbfigure}
  \adjustbox{max width=0.94\linewidth,max totalheight=0.46\textheight}{\includegraphics{figures/fig1-4_abstraction.png}}\\[3pt]
  \figcaption{그림 1-4  AI 인프라의 추상화 스펙트럼. 오른쪽으로 갈수록 운영할 대상이 줄지만 제어권도 줄어든다. 이 책은 의도적으로 IaaS부터 관리형 K8s까지의 범위를 다룬다.}
\end{tbbfigure}

각 단계를 서빙 관점에서 읽어 보자. \textbf{GPU VM (IaaS)}은 GPU가 장착된 머신 전체를 빌려 준다. 무엇이든 할 수 있지만 OS 패치부터 장애 복구까지 모든 일을 직접 해야 한다. \textbf{관리형 Kubernetes (CaaS)}는 컨테이너 오케스트레이션 클러스터를 빌려 준다. 컨테이너의 배포, 스케일링, 복구는 플랫폼이 처리한다. 오늘날 자체 관리형 서빙의 사실상 표준이며 이 책의 주 무대다(5주 차). \textbf{관리형 ML 플랫폼}(SageMaker, Vertex AI 등)은 "모델 파일을 건네면 엔드포인트를 받는" 수준까지 올라간다(13주 차에 소개한다). \textbf{서버리스 GPU} 플랫폼은 컨테이너나 함수만 받아 요청을 처리하는 동안에만 GPU를 할당하고 초 단위로 과금한다. 콜드 스타트(12주 차)를 감수하는 대신 사례 B가 수학적으로 일어날 수 없다. 가장 오른쪽의 \textbf{모델 API}는 다른 사람의 인프라에서 다른 사람의 모델을 사용하고 토큰 단위로 비용을 내는 방식, 즉 AI를 SaaS처럼 소비하는 형태다.

\begin{tbbtablewrap}
  \begin{tabular}{@{}>{\raggedright\arraybackslash}m{\dimexpr 0.2021\tbltextwidth-2\tabcolsep\relax} >{\raggedright\arraybackslash}m{\dimexpr 0.2660\tbltextwidth-2\tabcolsep\relax} >{\raggedright\arraybackslash}m{\dimexpr 0.2766\tbltextwidth-2\tabcolsep\relax} >{\raggedright\arraybackslash}m{\dimexpr 0.2553\tbltextwidth-2\tabcolsep\relax}@{}}
  \arrayrulecolor{tbbnavy}\specialrule{1pt}{0pt}{0pt}
  \rowcolor{tbbtablehead}{\centering \bfseries\color{tbbnavy}계층} & {\centering \bfseries\color{tbbnavy}빌리는 것} & {\centering \bfseries\color{tbbnavy}직접 관리할 것} & {\centering \bfseries\color{tbbnavy}서빙 예시} \\
  \arrayrulecolor{tbbnavy}\specialrule{0.7pt}{0pt}{2pt}
  {\textbf{IaaS}} & {GPU VM, 스토리지, 네트워크} & {OS, 런타임, 서빙 소프트웨어, 스케일링, 복구} & {EC2/GCE GPU 인스턴스에 vLLM을 직접 설치} \\
  \arrayrulecolor{tbbtableborder}\hline
  {\textbf{CaaS}} & {+ K8s 클러스터 운영} & {컨테이너 이미지, 배포 명세, 서빙 구성} & {EKS · GKE의 모델 서버(이 책)} \\
  \arrayrulecolor{tbbtableborder}\hline
  {\textbf{PaaS / ML 플랫폼}} & {+ 서빙 런타임, 스케일링, 배포} & {모델 아티팩트, 구성} & {SageMaker · Vertex AI 엔드포인트} \\
  \arrayrulecolor{tbbtableborder}\hline
  {\textbf{서버리스 GPU}} & {+ 전체 인스턴스 수명 주기} & {컨테이너/함수 코드} & {Modal, RunPod Serverless, ...} \\
  \arrayrulecolor{tbbtableborder}\hline
  {\textbf{SaaS (모델 API)}} & {모델과 인프라 전체} & {프롬프트, 호출 코드} & {OpenAI, Anthropic 등의 모델 API} \\
  \arrayrulecolor{tbbnavy}\specialrule{1pt}{2pt}{0pt}
  \end{tabular}\par
  \figcaption{표 1-2  서비스 계층별 서빙: 빌리는 것과 직접 관리하는 것}
\end{tbbtablewrap}

이 책이 택한 교육 방식도 분명히 해 두자. 업계에서는 표의 아래쪽 행, 즉 추상화 수준이 높은 선택지를 고르는 것이 합리적일 때가 많다. 그럼에도 이 책은 IaaS에서 CaaS까지의 길을 따라 \textbf{직접 구축한다}. 높은 추상화 수준에서 무언가 고장 났을 때 아래에서 어떤 일이 일어나는지 아는 사람과 모르는 사람은 전혀 다르게 대응하기 때문이다. 관리형 엔드포인트의 청구서를 해석하거나 서버리스 플랫폼의 콜드 스타트를 줄이는 일도 결국 아래 계층에 대한 이해를 요구한다. 처음부터 한 번 직접 만들어 보면 어떤 추상화도 다시는 두렵지 않을 것이다.

\begin{calloutbox}{tbbpurple}{tbbpurplefill}
{\normalsize\bfseries\color{tbbpurple}생각해 보기}\par\vspace{3pt}
\begin{tbbitemize}
  \item 수강 신청 기간에만 트래픽이 급증하는 교내 챗봇을 만든다고 하자. 그림 1-4에서 어느 지점을 선택하겠는가? 유휴 비용과 가용한 운영 인력을 근거로 논증해 보자.
  \item "공동 책임 모델(shared responsibility model)"이라는 표현을 찾아보자. 관리형 K8s에서 보안 사고가 발생하면 어느 부분이 공급자의 책임이고 어느 부분이 사용자의 책임인가?
\end{tbbitemize}
\end{calloutbox}

\section{GPUaaS와 모델 API 생태계}

이제 실제 시장을 살펴보자. 지난 몇 년 동안 두 극점 사이에서 생태계가 폭발적으로 성장했다. 하나는 GPU 계산 자원 자체를 판매하는 \textbf{GPUaaS}(GPU-as-a-Service), 다른 하나는 모델의 역량 자체를 판매하는 \textbf{모델 API}다. 시장은 빠르게 변하므로 이 절의 목표는 회사 이름을 외우는 것이 아니라 \textbf{지형을 읽는 법}을 익히는 데 있다.

\subsection*{하드웨어 감각: 데이터 센터 GPU}

시장보다 먼저 상품부터 살펴보자. NVIDIA는 데이터 센터 GPU의 사실상 표준이며, 세대가 바뀔 때마다 속도뿐 아니라 \textbf{메모리}도 커진다. A100 (40/80 GB) → H100 (80 GB) → H200 (141 GB) → Blackwell 세대 B200 (192 GB)로 이어지며, 여러 CPU와 GPU를 하나의 장치로 결합한 GB200 NVL72 같은 랙 규모 시스템도 있다. 서빙에서 메모리가 이토록 중요한 이유는 무엇인가? \textbf{추론을 실행하려면 모델이 GPU 메모리에 들어가야 하기} 때문이다. 16-bit 정밀도의 파라미터가 700억 개인 모델은 가중치만 약 140 GB가 필요하므로 80 GB H100 한 대에 들어가지 않는다. 양자화(10주 차)로 줄이거나 여러 GPU에 샤딩해야 한다. 한편 추론용 GPU인 L4 (24 GB), T4 (16 GB)는 저렴하고 전력 효율이 높은 일꾼이며, 이 책의 실습은 T4 한 대를 기준으로 설계했다.

여기서 이 책 내내 본능이 될 때까지 반복할 용량 산정 원칙이 처음 등장한다. \textbf{먼저 결정적인 제약 차원을 찾아라.} 대부분의 추론에서 실제 제약은 순수 FLOPs가 아니라 메모리 용량(모델이 들어가는가?)이나 메모리 대역폭(가중치가 계산 장치를 얼마나 빠르게 통과할 수 있는가?)이다. 그래서 소형 모델은 적당한 L4에서도 훨씬 저렴한 가격으로 H100과 거의 같은 사용자 경험을 제공할 때가 많다. "안전을 위해" 가장 큰 GPU를 고르는 일은 대개 애초에 제약이 아니었던 차원에 돈을 내는 셈이다. 인스턴스는 결정적인 제약 차원에 따라 선택해야 하며, 이후의 모든 비용 연습(9·10·12주 차)은 여기서 시작한다.

\begin{tbbtablewrap}
  \begin{tabular}{@{}>{\raggedright\arraybackslash}m{\dimexpr 0.3830\tbltextwidth-2\tabcolsep\relax} >{\raggedright\arraybackslash}m{\dimexpr 0.2553\tbltextwidth-2\tabcolsep\relax} >{\raggedright\arraybackslash}m{\dimexpr 0.3617\tbltextwidth-2\tabcolsep\relax}@{}}
  \arrayrulecolor{tbbnavy}\specialrule{1pt}{0pt}{0pt}
  \rowcolor{tbbtablehead}{\centering \bfseries\color{tbbnavy}임대처} & {\centering \bfseries\color{tbbnavy}대략적인 가격대} & {\centering \bfseries\color{tbbnavy}비고} \\
  \arrayrulecolor{tbbnavy}\specialrule{0.7pt}{0pt}{2pt}
  {하이퍼스케일러(AWS · GCP · Azure)} & {\$3–7 / hour} & {생태계 통합의 대가로 프리미엄 부과} \\
  \arrayrulecolor{tbbtableborder}\hline
  {GPU 전문 클라우드(Lambda, CoreWeave, ...)} & {\$2.0–2.5 / hour} & {전문화로 단가 절감} \\
  \arrayrulecolor{tbbtableborder}\hline
  {마켓플레이스(Vast.ai, ...)} & {\$1.9–3.5 / hour} & {유휴 GPU 중개, 저렴하지만 변동성 있음} \\
  \arrayrulecolor{tbbtableborder}\hline
  {\textcolor{tbbgray}{(참고) B200}} & {\textcolor{tbbgray}{\$4–6 / hour}} & {\textcolor{tbbgray}{차세대, 더 높은 처리량}} \\
  \arrayrulecolor{tbbnavy}\specialrule{1pt}{2pt}{0pt}
  \end{tabular}\par
  \figcaption{표 1-3  2026년 중반 기준 H100 (80 GB) 온디맨드 시간당 대략적인 가격}
\end{tbbtablewrap}

\begin{calloutbox}{tbbred}{tbbxFEF2F2}
{\normalsize\bfseries\color{tbbx991B1B}주의: 숫자가 아니라 구조를 기억하라}\par\vspace{3pt}
이 표의 숫자는 독자가 읽을 때쯤 이미 틀렸을 수 있다. GPU 가격은 세대와 공급 상황에 따라 움직인다. 기억할 것은 구조다. \textbf{① 전문 클라우드는 GPU 가격으로 하이퍼스케일러보다 싸게 경쟁한다. ② 같은 GPU도 약정 사용, 스팟, 마켓플레이스 채널에 따라 가격이 몇 배씩 달라진다. ③ 시간 단위로 계량되므로 유휴 시간은 곧 돈이다.} 그래서 첫 실습 과제도 과금 알림 설정이다(2주 차).
\end{calloutbox}

\vspace{3.0pt}

\subsection*{생태계 지도: 여섯 개의 집단}

그림 1-5는 생태계를 "인프라 판매"와 "모델 역량 판매"라는 하나의 축을 따라 여섯 개 군집으로 배치한다.

\begin{tbbfigure}
  \adjustbox{max width=0.94\linewidth,max totalheight=0.46\textheight}{\includegraphics{figures/fig1-5_ecosystem.png}}\\[3pt]
  \figcaption{그림 1-5  GPUaaS와 모델 API 생태계. 왼쪽은 계산 자원을, 오른쪽은 역량을 판매한다. 회사 이름은 2026년의 예시일 뿐이며 중요한 것은 군집의 경계다.}
\end{tbbfigure}

\textbf{하이퍼스케일러}(AWS, Google Cloud, Azure)는 범용 클라우드의 일부로 GPU를 판매하며 관리형 ML 플랫폼도 함께 제공한다. 운영 중인 다른 모든 서비스와 통합하기 쉽고 보안과 규정 준수가 성숙했다는 장점이 있지만, GPU 가격에는 프리미엄이 붙는다. 시장에서 \textit{네오클라우드}라 부르는 \textbf{GPU 전문 클라우드}는 AI 워크로드를 위해 특별히 구축한 클라우드다. CoreWeave, Lambda, Nebius, Crusoe가 대표적이다. 새 GPU를 빠르게 도입하고 공격적으로 가격을 책정한다. Vast.ai 같은 마켓플레이스형 사업자는 유휴 GPU를 중개한다. 공급 변동성을 감수한다면 가장 저렴하다.

\textbf{서버리스 GPU 플랫폼}(Modal, RunPod Serverless, Baseten, Replicate, ...)은 서버라는 개념 자체를 숨기고 코드나 컨테이너 호출 단위로 GPU를 빌려 준다. 요청이 실행되는 동안에만 과금하므로 트래픽이 급증하거나 적은 서비스에 매력적이지만, 1.3절에서 만난 콜드 스타트가 약점이다. \textbf{오픈 모델 추론 플랫폼}(Together AI, Fireworks AI, Hugging Face, ...)은 공개된 오픈 웨이트 모델을 API 뒤에서 서빙한다. 직접 서빙할 때의 모델 자유도와 API의 운영 편의성을 함께 제공하는 중간 지대다. 이 군집에서는 Groq (LPU), Cerebras (WSE) 같은 기업이 GPU가 아닌 맞춤형 실리콘으로 속도 경쟁을 벌인다. 마지막으로 \textbf{프런티어 모델 API}(OpenAI, Anthropic, Google)는 폐쇄형 최첨단 모델을 토큰 단위로 판매하며, \textbf{라우터/애그리게이터}(OpenRouter 등)는 앞의 모든 서비스를 단일 API 뒤에 묶는다.

최근 몇 년 사이 이 지형을 가장 크게 바꾼 흐름은 \textbf{오픈 웨이트 모델의 부상}이다. 가중치를 공개적으로 내려받을 수 있는 Llama, Qwen, DeepSeek, Mistral 계열의 품질이 폐쇄형 모델을 따라잡으면서 "좋은 모델을 내 인프라에서 서빙한다"는 선택이 현실이 되었다. 자체 서빙 수요가 늘면 지도 왼쪽의 GPU 인프라와 가운데의 서빙 기술도 함께 성장한다. 이 책이 가르치는 내용이 바로 그 서빙 기술이다.

\subsection*{AI 서비스 경제의 장치: 키, 계량(meters), 할당량}

어느 군집에서 구매하든 거래의 형태는 공통적이다. 이 책에서 결국 직접 만들게 될 것이므로 이름을 붙여 둘 가치가 있다. 모든 모델 API와 GPU 플랫폼은 \textbf{API 키}(신원: \textit{누가} 호출하는가)를 발급하고, \textbf{계량기}(과금: 토큰당, 초당, 요청당)를 돌리며, 사용자를 \textbf{할당량과 속도 제한}(보호: 등급별 분당 요청 또는 토큰 상한) 안에 둔다. 상한을 넘으면 HTTP \textbf{429 Too Many Requests}를 받는다. 서비스 장애나 코드 오류가 아니라, 사용자의 트래픽이 다른 테넌트(tenant)에 피해를 주지 않도록 공급자의 로드 밸런서가 의도한 대로 정중히 거절한 것이다. 무료 등급, 사용량 대시보드, 지출 알림, 단계별 가격은 모두 키, 계량기, 울타리로 이루어진 같은 세 부분 장치다.

이처럼 평범한 일을 왜 장황하게 강조할까? 이 장치가 곧 제품이기 때문이다. 모델은 요청에 답하지만, 키·계량·할당량 장치가 요청에 답하는 행위를 \textit{사업}으로 만든다. 다른 사람이 운영하는 GPU 서버군을 예산에 넣을 수 있는 비용 항목으로 바꾸는 방법이다. 그리고 이 책이 닫으려는 고리가 바로 여기에 있다. 직접 서빙하는 순간, 당신이 곧 제공자가 된다. 그러면 모든 조각이 이번엔 당신의 책임으로 되돌아온다. 키를 확인하고 한도를 강제하던 API 게이트웨이는 6주차, 토큰당 원가를 알려주던 계량은 12주차, 그리고 팀 프로젝트는 바로 이런 정문 뒤에 자신의 모델을 내걸어야 한다. 6주차가 되면, 당신은 이 계산대의 반대편에 서 있게 된다.

\subsection*{자체 구축 대 외부 도입: 모델 API와 자체 서빙}

실제로 생태계의 두 끝은 하나의 결정으로 모인다. \textbf{모델 API를 구매할 것인가, 오픈 모델을 직접 서빙할 것인가?} 표 1-4에 판단 기준을 모았다.

\begin{tbbtablewrap}
  \begin{tabular}{@{}>{\raggedright\arraybackslash}m{\dimexpr 0.1915\tbltextwidth-2\tabcolsep\relax} >{\raggedright\arraybackslash}m{\dimexpr 0.4043\tbltextwidth-2\tabcolsep\relax} >{\raggedright\arraybackslash}m{\dimexpr 0.4043\tbltextwidth-2\tabcolsep\relax}@{}}
  \arrayrulecolor{tbbnavy}\specialrule{1pt}{0pt}{0pt}
  \rowcolor{tbbtablehead}{\centering \bfseries\color{tbbnavy}기준} & {\centering \bfseries\color{tbbnavy}모델 API에 유리(외부 도입)} & {\centering \bfseries\color{tbbnavy}자체 서빙에 유리(자체 구축)} \\
  \arrayrulecolor{tbbnavy}\specialrule{0.7pt}{0pt}{2pt}
  {\textbf{규모 / 트래픽}} & {작거나 불확실함 — 사용한 만큼만 지불} & {크고 일정함 — GPU를 포화 상태로 유지 가능} \\
  \arrayrulecolor{tbbtableborder}\hline
  {\textbf{데이터 / 규제}} & {데이터의 경계 외부 반출에 제약 없음} & {민감한 데이터, 온프레미스 또는 데이터 주권 요구 사항} \\
  \arrayrulecolor{tbbtableborder}\hline
  {\textbf{맞춤화}} & {범용 역량으로 충분함} & {미세 조정 모델, 특수 도메인, 출력 완전 제어} \\
  \arrayrulecolor{tbbtableborder}\hline
  {\textbf{지연 시간 제어}} & {공급자 SLA로 충분함} & {꼬리 지연 시간, 배칭, 하드웨어를 직접 조정해야 함} \\
  \arrayrulecolor{tbbtableborder}\hline
  {\textbf{팀 역량}} & {가용한 인프라 엔지니어가 없음} & {이 책으로 공부한 엔지니어가 있음} \\
  \arrayrulecolor{tbbtableborder}\hline
  {\textbf{최고 역량}} & {프런티어급 역량이 반드시 필요함} & {작업에 오픈 모델 품질로 충분함} \\
  \arrayrulecolor{tbbnavy}\specialrule{1pt}{2pt}{0pt}
  \end{tabular}\par
  \figcaption{표 1-4  모델 API와 자체 서빙의 선택 기준}
\end{tbbtablewrap}

비용 축은 빠른 어림 계산으로 살펴볼 수 있다. 아래 숫자는 \textbf{가상 값}이다. 소형 오픈 모델로 하루 토큰 2,000만 개를 생성하는 서비스를 가정하자. 토큰 100만 개당 \$0.5인 API를 쓰면 하루 \$10, 한 달에 약 \textbf{\$300}다. 같은 부하를 시간당 \$0.8인 L4 GPU 한 대로 직접 서빙하면서 상시 실행하면 한 달 약 \textbf{\$580}다. \textbf{이 규모에서는 API가 더 저렴하다.} 이제 트래픽이 20배 증가한다고 하자. API 청구액은 선형으로 늘어 \$6,000가 되지만, 자체 서빙 클러스터는 GPU 몇 대만 더 추가하면 따라갈 수 있다. 배칭이 처리량을 초선형으로 높이기 때문이다(10주 차). 중간 어딘가에서 두 곡선이 교차한다. 손익분기점을 제대로 찾으려면 GPU 사용률, 최적화 수준, 운영 인원까지 고려해야 한다. 팀 프로젝트의 "월간 운영 비용 추정"이 요구하는 작업이 바로 이것이다.

마지막으로 강조할 점은 이것이 양자택일이 아니라는 사실이다. 실제 조직은 흔히 \textbf{하이브리드 전략}을 쓴다. 프런티어 API로 프로토타입을 만들고 검증한 뒤, 성능이 입증되고 트래픽이 많은 핵심 부분을 자체 서빙 오픈 모델로 옮긴다. 어느 쪽을 선택하든 이 절의 목표는 \textbf{무엇을 포기하고 무엇을 얻는지 아는 상태에서} 선택하는 것이다.

\begin{calloutbox}{tbbpurple}{tbbpurplefill}
{\normalsize\bfseries\color{tbbpurple}생각해 보기}\par\vspace{3pt}
\begin{tbbitemize}
  \item 병원에서 임상 기록을 요약하는 AI를 도입하려 한다. 표 1-4의 기준을 적용해 보자. 어느 기준이 결정적인 장애물인가?
  \item 어림 계산 예시에서 자체 서빙 GPU의 사용률이 30\%에 불과해 하루 대부분을 유휴 상태로 보낸다고 하자. 실제 토큰당 비용은 어떻게 되는가? 사용률을 높이는 방법에는 무엇이 있는가? (12주 차 예고.)
  \item 자체 서비스의 트래픽이 급증한 동안 모델 API에서 429를 받았다. 이 역설을 추적해 보자. 공급자의 키·계량·할당량 장치 중 어느 부분이 \textit{공급자의} p99를 지키는 대가로 \textit{자체 서비스의} p99를 희생시켰는가? 공급자 서비스에 의존하는 동안 자체 서비스가 약속할 수 있는 SLO에 관해 무엇을 알려 주는가?
\end{tbbitemize}
\end{calloutbox}

\section{이 책의 로드맵과 학습 방법}

이 장에서 그린 큰 그림이 앞으로 15주 동안 어떻게 펼쳐지는지 살펴보자. 그림 1-6은 전체 로드맵이다.

\begin{tbbfigure}
  \adjustbox{max width=0.94\linewidth,max totalheight=0.46\textheight}{\includegraphics{figures/fig1-6_roadmap.png}}\\[3pt]
  \figcaption{그림 1-6  15주 로드맵. 빨간 점선은 제1부의 원리가 뒤의 주제에서 다시 호출되는 대표적인 경로를 표시한다.}
\end{tbbfigure}

\textbf{제1부(1\textasciitilde{}6주 차)}는 모든 AI 서비스가 서 있는 지반 아래의 커널까지 파고든다. 가상화(virtualization)가 CPU 명령어 수준에서 작동하는 방식(2주 차), 실제로 컨테이너를 구성하는 커널 기능(3주 차), Docker와 Kubernetes가 그 기능을 패키징하는 방식(4\textasciitilde{}5주 차), 모델 파일과 트래픽을 운반하는 스토리지와 네트워크의 모습(6주 차)을 배운다. \textbf{제2부(7\textasciitilde{}12주 차)}는 그 지반 위에 서빙 시스템을 구축한다. SLO와 p99라는 서빙 성능 언어(7주 차), 프레임워크로 실제 모델을 서빙하고 부하 테스트하기(9주 차), 양자화·컴파일·배칭으로 속도 높이기(10주 차), LLM 서빙의 독특한 세계(11주 차), 스케일링과 비용(12주 차)을 다룬다. \textbf{제3부(13\textasciitilde{}15주 차)}는 운영이다. 배포 자동화(13주 차), 모니터링과 신뢰성(14주 차), 팀 프로젝트 마무리(15주 차)로 이어진다.

로드맵이 강조하는 것은 빨간 점선, 즉 \textbf{콜백 구조}다. 이 책 전반부는 후반부를 예고하며, 의도적으로 되풀이되는 기념물을 여럿 만나게 된다. 나중에 다시 등장할 때 반갑게 알아볼 수 있도록 일부를 미리 소개한다.

\begin{tbbitemize}
  \item \textbf{숫자 137} — 이 장에서 수수께끼 같은 종료 코드로 처음 만나고, 3주 차에는 자기 프로세스에 직접 일으킨다(커널의 OOM 킬러다. 128 + SIGKILL의 9). 5주 차에는 Kubernetes가 이를 강제하며, 이후 내리는 모든 메모리 용량 결정에 따라붙는다.
  \item \textbf{수수께끼 같은 부모 프로세스 PID 41288} — 3주 차의 첫 번째 기록에서 컨테이너를 시작할 때 설명 없이 등장하고, 4주 차에 마침내 심문을 거쳐 정체가 밝혀진다. (미리 밝히지는 않겠다.)
  \item \textbf{1956년의 화물 컨테이너} — 4주 차는 강철 상자 하나가 화물 운송 비용을 \textasciitilde{}37× 줄인 이야기와 Docker 이미지가 소프트웨어에 같은 묘수를 적용한 이유로 시작한다. 이후에는 "컨테이너"라는 단어가 전과 달리 보일 것이다.
  \item \textbf{2주 차의 Firecracker microVM}(밀리초 만에 부팅되는 VM)은 12주 차의 \textbf{GPU 콜드 스타트}에서 돌아온다. "서버리스 GPU에서 첫 요청이 왜 느린가?"에 답하려면 머신이 부팅된다는 것이 물리적으로 무엇을 뜻하는지 알아야 한다.
  \item \textbf{3주 차에 직접 작성하는 cgroup 파일}(\code{memory.max}, \code{cpu.max})은 5주 차에 Kubernetes 자원 요청량/한도가 도달하는 바로 그곳이다. 하나의 메커니즘이 두 인터페이스를 쓴 것이며, 이 장의 밀도와 비용 절충에 대한 직접적인 답이다.
  \item \textbf{2주 차의 GPU 가상화}(PCIe 패스스루, vGPU, MIG)는 4주 차에서 커널 자체가 GPU 시간을 중재하지 못하는 이유를 설명한 뒤, 12주 차에 여러 서비스가 비싼 GPU 하나를 나누는 \textbf{GPU 공유} 전략으로 다시 등장한다.
\end{tbbitemize}

\subsection*{팀 프로젝트: 이 책의 종합 시험}

성적에서 가장 큰 비중(30\%)을 차지하는 팀 프로젝트의 과제는 한 문장으로 정리된다. \textbf{"오픈 소스 AI 모델을 선택해 클라우드에 배포하고, 목표 SLO를 충족함을 부하 테스트로 입증하라."} Kubernetes 기반 컨테이너 서빙, 부하 상황에서 SLO를 검증하는 오토스케일링, 모니터링 대시보드, 최적화 전후 비교, 월간 운영 비용 추정까지 이 장에서 둘러본 모든 주제가 프로젝트 하나에 들어간다. 제안서는 7주 차, 중간 점검은 12주 차, 최종 시연은 15주 차에 진행한다. 이번 주에는 팀을 구성하기만 하면 된다. 매주 배우는 내용을 차례로 조립하도록 설계한 프로젝트이므로 지금부터 겁먹을 이유는 없다.

\begin{tbbtablewrap}
  \begin{tabular}{@{}>{\raggedright\arraybackslash}m{\dimexpr 0.2553\tbltextwidth-2\tabcolsep\relax} >{\raggedright\arraybackslash}m{\dimexpr 0.1489\tbltextwidth-2\tabcolsep\relax} >{\raggedright\arraybackslash}m{\dimexpr 0.5957\tbltextwidth-2\tabcolsep\relax}@{}}
  \arrayrulecolor{tbbnavy}\specialrule{1pt}{0pt}{0pt}
  \rowcolor{tbbtablehead}{\centering \bfseries\color{tbbnavy}평가 요소} & {\centering \bfseries\color{tbbnavy}비중} & {\centering \bfseries\color{tbbnavy}평가 범위} \\
  \arrayrulecolor{tbbnavy}\specialrule{0.7pt}{0pt}{2pt}
  {중간고사} & {\centering 20\%} & {1\textasciitilde{}7주 차: 클라우드 인프라, 가상화 원리, 서빙 기초} \\
  \arrayrulecolor{tbbtableborder}\hline
  {기말고사} & {\centering 20\%} & {9\textasciitilde{}14주 차: AI 서빙 시스템 구축, 최적화, 운영} \\
  \arrayrulecolor{tbbtableborder}\hline
  {실습 과제} & {\centering 25\%} & {주별 실습 산출물과 벤치마크 보고서(5\textasciitilde{}6회 제출)} \\
  \arrayrulecolor{tbbtableborder}\hline
  {팀 프로젝트} & {\centering 30\%} & {제안서 · 중간 점검 · 최종 시연, 요구 사항과 SLO 달성 여부 평가} \\
  \arrayrulecolor{tbbtableborder}\hline
  {출석} & {\centering 5\%} & {대학 규정에 따름} \\
  \arrayrulecolor{tbbnavy}\specialrule{1pt}{2pt}{0pt}
  \end{tabular}\par
  \figcaption{표 1-5  평가 방식 한눈에 보기}
\end{tbbtablewrap}

\subsection*{실습 환경과 이번 주에 할 일}

실습은 AWS Academy 또는 Google Cloud for Education의 교육용 크레딧으로 진행한다. 모든 실습은 파라미터 1–3 B 규모의 소형 모델에 양자화를 적용한다는 가정 아래 T4 GPU 한 대나 Colab 세션 하나에서 실행되도록 설계했다. GPU가 전혀 없다면 \code{llama.cpp}를 통한 CPU 추론 경로로 대체할 수 있다. Kubernetes 실습은 비용을 아끼기 위해 기본적으로 로컬 환경(minikube/kind)을 사용한다. 이번 주에 할 일은 두 가지다. \textbf{① 환경 준비}(Python과 Docker를 실행할 수 있는 개인 노트북, 클라우드 교육용 크레딧 신청), \textbf{② 팀 구성}이다. 자세한 내용은 별도의 실습 안내서에서 설명한다.

\subsection*{이 책에서 좋은 성과를 내는 법}

네 가지 학습 조언이 있다. 첫째, \textbf{지금 선수 지식을 점검하라.} 2\textasciitilde{}3주 차부터 커널 수준의 내용을 다루므로 프로세스, 가상 메모리, 시스템 호출이 희미하게 느껴진다면 이번 주에 운영체제 과목을 복습하는 것이 가장 수익률 높은 투자다. Python 숙련도는 이 책 전체에서 전제로 삼는다. 둘째, \textbf{실습을 구경만 하지 마라.} 이 책의 지식은 남이 하는 모습을 지켜본 눈이 아니라 직접 \code{unshare}를 실행하고 네임스페이스를 만든 손에 쌓인다. 셋째, \textbf{어두운 상자를 재현하라.} 이 책에 실린 모든 터미널 기록(사례 1-A와 1-B 등)은 직접 실행하거나 산술적으로 검증할 수 있다. 제2장부터는 대부분 2주 차 실습에서 직접 빌릴 VM에서 실행된다. 직접 입력하라. 재현한 기록은 지식이고, 읽기만 한 기록은 잡학이다. 넷째, \textbf{계속 "왜"라고 물어라.} 모든 도구 뒤에는 그 도구가 해결하려고 탄생한 문제가 있다. "Kubernetes를 배웠다"보다 "Kubernetes 같은 것이 왜 존재해야 했는지 이해한다"가 더 가치 있다. 이 책이 추구하는 바가 후자이며, 이 책의 모든 장이 그 주의 기술을 필요하게 만든 재난, 청구서, 이야기로 시작하는 이유이기도 하다.

\namedsection{요약}{열두 가지로 정리한 제1장}

\qitem{1}{\textbf{모델은 서비스가 아니다.} 정확도는 모델의 속성이지만 동시성, 지연 시간, 가용성, 비용, 배포, 관측 가능성은 시스템의 속성이다. 사례 A는 정확도 94\%인 상태에서 실패했고, ChatGPT 출시는 모델링이 아니라 서빙 사건이었다. 프로덕션 ML 시스템에서 ML 코드는 작은 부분(Hidden Technical Debt)에 불과하며, 이 책은 나머지 전부를 다룬다.}

\qitem{2}{\textbf{사례 A의 계산은 일반화할 수 있다.} 서빙 시스템에는 처리량 상한(레플리카당 1 ÷ 요청별 시간)이 있다. 상한을 넘는 수요는 대기열이 되고, 대기열은 꼬리 지연 시간이 되며, 서비스 앞단의 구성 요소는 지나치게 긴 대기를 오류로 바꾼다. 서비스가 실패하는 데 어떤 것도 "충돌"할 필요는 없다.}

\qitem{3}{\textbf{사례 B의 계산도 일반화할 수 있다.} 상시 실행 용량은 트래픽 유무와 관계없이 매시간 과금된다. 기록해야 할 계량 지표는 사용률이며, 서빙 비용의 정직한 단위는 월이 아니라 요청당 비용(또는 토큰당 비용)이다.}

\qitem{4}{AI 서비스 스택은 \textbf{개발 → 학습 → 서빙 → 운영의 피드백 고리}이며, 모든 단계가 클라우드 인프라에서 실행된다. 이 책은 학습된 모델을 넘겨받는 시점부터 시작해 서빙과 운영을 책임진다.}

\qitem{5}{사용자가 체감하는 지연 시간은 모델 실행 시간이 아니라 \textbf{요청 경로의 합}이다. 성능 분석은 "어느 구간이 병목인가?"를 묻는 일이다.}

\qitem{6}{학습 워크로드는 \textbf{유한한 배치 작업으로, 처리량 중심이고 예측 가능하며 재시작할 수 있다.} 추론 워크로드는 \textbf{상시 실행 서비스로, 꼬리 지연 시간(p95/p99) 중심이고 부하가 급증하며 장애가 사용자에게 즉시 피해를 준다.} 후반부의 모든 주제(배칭, 오토스케일링, 콜드 스타트, 카나리 배포, 모니터링)는 이 대비에서 나온다.}

\qitem{7}{\textbf{지연 시간·처리량·비용으로 이루어진 서빙 삼각형은 이 책의 나침반이다.} 어느 하나를 지불하면 나머지 둘을 살 수 있다. SLO는 삼각형 위에서 의도적으로 선택한 점이며, 모든 최적화 주장은 무엇을 지불하는지 밝혀야 한다. 배칭, 양자화, 오토스케일링 같은 기법은 삼각형을 없애는 것이 아니라 움직인다.}

\qitem{8}{같은 모델에서 추론은 학습보다 훨씬 적은 메모리가 필요하지만, 메모리 사용량은 \textbf{시간에 따라 뾰족하게 출렁인다}(기준선인 가중치 + 부하 급증에 비례하는 활성값). 따라서 "관측한 유휴 사용량 + 10\%"를 메모리 한도로 삼는 것은 안전 여유가 아니라 예정된 OOM 처형(exit 137, 3주 차의 메커니즘)이다.}

\qitem{9}{클라우드 서비스 모델은 \textbf{"얼마만큼 빌리는가"}의 스펙트럼이다. GPU VM (IaaS) → 관리형 K8s (CaaS) → 관리형 엔드포인트 → 서버리스 GPU → 모델 API (SaaS)로 이어진다. 오른쪽으로 갈수록 제어권을 편의성과 교환한다. 이 책은 원리를 배우기 위해 의도적으로 왼쪽에서 직접 구축한다.}

\qitem{10}{GPUaaS 생태계는 \textbf{하이퍼스케일러 / GPU 전문 클라우드(네오클라우드) / 마켓플레이스 / 서버리스 GPU / 오픈 모델 추론 플랫폼 / 프런티어 모델 API / 라우터}로 이루어진다. 회사 이름은 바뀌지만 군집 구조는 남는다. 어느 군집에서든 인스턴스 선택은 \textbf{결정적인 제약 차원}, 보통 FLOPs가 아니라 메모리를 따라야 한다.}

\qitem{11}{AI 경제의 배관은 \textbf{키, 계량기, 할당량}이다(429는 장애가 아니라 울타리다). 자체 서빙을 하면 공급자가 되어 세 가지 모두 다시 구축한다. 6주 차의 게이트웨이, 12주 차의 계량기, 7주 차의 그 뒤에 있는 SLO다.}

\qitem{12}{이 책의 구조는 \textbf{원리(1\textasciitilde{}6주 차) → 서빙(7\textasciitilde{}12주 차) → 운영(13\textasciitilde{}15주 차)}이며, exit 137, PID 41288, 1956년의 화물 컨테이너, Firecracker, 모든 Kubernetes 한도 아래에 있는 cgroup 파일이 전반부와 후반부를 의도적으로 연결한다.}

\namedsection{연습문제}{}

\subsection*{A. 객관식}

\qitem{1}{AI 서비스 스택의 순서로 옳은 것은?}

\begin{tbboptions}
  \item ① 학습 → 개발 → 운영 → 서빙        ② 개발 → 학습 → 서빙 → 운영
  \item ③ 개발 → 서빙 → 학습 → 운영        ④ 학습 → 서빙 → 개발 → 운영
\end{tbboptions}

\qitem{2}{추론(서빙) 워크로드를 올바르게 설명한 것은?  ㉠ 시작과 끝이 있는 유한한 작업이다  ㉡ 지연 시간이 핵심 성능 지표다  ㉢ 사용자 트래픽에 따라 부하가 변동한다  ㉣ 체크포인트에서 재시작할 수 있으므로 장애 비용이 작다}

\begin{tbboptions}
  \item ① ㉠, ㉡        ② ㉡, ㉢        ③ ㉠, ㉣        ④ ㉢, ㉣
\end{tbboptions}

\qitem{3}{"p99 지연 시간 500 ms"의 정확한 의미는?}

\begin{tbboptions}
  \item ① 평균 응답 시간이 500 ms다
  \item ② 가장 느린 요청에 500 ms가 걸렸다
  \item ③ 전체 요청의 99\%에 500 ms 안에 응답했다
  \item ④ 요청의 99\%가 500 ms보다 오래 걸렸다
\end{tbboptions}

\qitem{4}{한 팀이 컨테이너 이미지를 만들고 관리형 Kubernetes(EKS/GKE)에 모델 서버를 배포한다. 이 선택을 가장 잘 설명한 것은?}

\begin{tbboptions}
  \item ① SaaS를 사용하는 것이므로 인프라 지식이 필요 없다
  \item ② CaaS 수준의 선택이다. 클러스터 운영은 위임하지만 컨테이너와 배포 명세는 계속 팀이 담당한다
  \item ③ GPU VM 전체를 빌리는 IaaS이므로 팀이 OS를 패치해야 한다
  \item ④ 서버리스이므로 유휴 상태에는 요금이 전혀 부과되지 않는다
\end{tbboptions}

\qitem{5}{자체 서빙(자체 구축)보다 모델 API(외부 도입)가 가장 분명히 유리한 상황은?}

\begin{tbboptions}
  \item ① 환자 기록을 구내 밖으로 반출할 수 없는 병원 시스템
  \item ② GPU를 포화 상태로 유지할 수 있을 만큼 트래픽이 크고 일정한 서비스
  \item ③ 트래픽이 거의 없는 아이디어 검증 단계의 MVP
  \item ④ 미세 조정 모델의 출력 형식을 완전히 제어해야 하는 서비스
\end{tbboptions}

\qitem{6}{파라미터 70 B 규모의 모델을 FP16(파라미터당 2 bytes)으로 로드하려 한다. 가중치만 고려할 때 80 GB H100 한 대에 대한 올바른 메모리 추정과 판단은?}

\begin{tbboptions}
  \item ① ≈ 70 GB — 들어간다        ② ≈ 140 GB — 들어가지 않는다
  \item ③ ≈ 35 GB — 들어간다        ④ ≈ 280 GB — 들어가지 않는다
\end{tbboptions}

\qitem{7}{한 팀이 GPU 서비스에 공격적인 배칭을 적용하자 GPU당 처리량이 60\% 늘고 토큰당 비용이 35\% 줄었다. 서빙 삼각형에 따르면 가장 가능성 높은 대가는 무엇이며, 어디에서 먼저 나타나는가?}

\begin{tbboptions}
  \item ① 모델 정확도 — 검증 점수
  \item ② 꼬리 지연 시간 — 요청이 배치가 찰 때까지 기다리면서 p95/p99에 나타남
  \item ③ 가용성 — 더 잦은 충돌
  \item ④ 없음 — 배칭은 대가 없는 개선임
\end{tbboptions}

\qitem{8}{추론 서버가 유휴 상태에서 6.0 GB를 쓴다. "관측한 유휴 사용량 + 10\%" 규칙에 따라 메모리 한도를 6.6 GB로 설정했다. 앞으로 벌어질 일을 가장 잘 예측한 것은?}

\begin{tbboptions}
  \item ① 안전하다. 10\% 여유가 현실적인 모든 변동을 감당한다
  \item ② 모델을 더 큰 크기로 다시 학습할 때만 위험하다
  \item ③ 위험하다. 활성값 메모리는 동시 요청 수와 입력 길이에 비례하므로, 처음으로 실제 부하가 몰리면 6.6 GB를 넘어 서버가 죽을 수 있다(exit 137)
  \item ④ 커널은 프로세스를 죽이지 않고 메모리를 스로틀링하므로 안전하다
\end{tbboptions}

\subsection*{B. 단답형}

\qitem{1}{"p99 지연 시간 ≤ 500 ms, 가용성 ≥ 99.9\%"처럼 표현하며, "서비스가 스스로 설정하는 정량적 품질 목표"를 무엇이라 하는가?}

\qitem{2}{0으로 스케일링은 유휴 상태일 때 서비스의 인스턴스를 0개로 줄인다. 조용한 시간이 지난 뒤 첫 요청이 인스턴스 시작과 모델 로딩을 기다리느라 매우 느려지는 문제를 무엇이라 하는가?}

\qitem{3}{빈칸을 채워라. 학습 워크로드의 주요 성능 지표는 ( ㉠ )이며, 온라인 추론 워크로드의 주요 성능 지표는 ( ㉡ )이다.}

\qitem{4}{파라미터 13 B 규모 모델의 가중치는 FP16에서 GPU 메모리를 약 얼마나 차지하는가? (계산 과정을 제시하라.)}

\qitem{5}{코드가 모델 API에서 HTTP 429를 받았다. 어떤 메커니즘이 이를 만들었으며, 이것이 코드의 버그나 공급자 서비스의 장애가 아닌 이유를 한 문장으로 설명하라.}

\subsection*{C. 서술형}

\qitem{1}{신입 팀원이 말한다. "모델의 검증 정확도가 94\%이므로 내일 서비스로 출시하겠습니다." 동시성, 지연 시간, 가용성, 비용, 배포, 관측 가능성 중 최소 세 가지 관점에서 시연과 프로덕션 서비스 사이의 간극을 구체적인 예로 설명하라.}

\qitem{2}{학습 워크로드와 추론 워크로드의 차이 중 두 가지를 선택하고, 각각이 만들어 내는 인프라 요구 사항(예: 스케줄러와 로드 밸런서, 수직 확장과 수평 확장)에 명시적인 논리 사슬로 연결하라.}

\qitem{3}{초기 스타트업이 법률 문서 요약 서비스를 만든다. 고객은 법률 회사 다섯 곳이고, 문서에는 민감한 계약 세부 사항이 담겨 있으며, 내년에 도메인 미세 조정 모델을 만들 계획이다. 모델 API(외부 도입), 자체 서빙(자체 구축), 단계적 전략 중 무엇을 권하겠는가? 현재 단계와 성장 단계를 구분해 표 1-4의 기준으로 논증하라.}

\qitem{4}{당직 엔지니어가 되어 사례 1-A를 다시 실행해 보자. 레플리카 하나는 한 번에 요청 하나를 0.25 s에 처리하고, 게이트웨이 타임아웃은 20 s이며, 사용자 80명이 도착해 연결을 유지한다. (a) 처리량 상한, (b) 대략적인 중앙값 대기 시간, (c) 504가 발생하는지 계산하라. 그런 다음 서빙 삼각형의 서로 다른 꼭짓점을 공략하는 해결책 두 가지를 제안하고 각각 무엇을 지불하는지 밝혀라.}

\namedsection{정답 및 해설}{}

\subsection*{A. 객관식}

\qitem{1}{\textbf{②} 개발(데이터, 실험) → 학습(모델 생성) → 서빙(서비스로 제공) → 운영(건강하게 유지)이며, 운영에서 개발과 학습으로 돌아가는 피드백 경로가 이들을 순환 고리로 연결한다. (1.2절)}

\qitem{2}{\textbf{②} ㉡과 ㉢이 옳다. ㉠은 학습을 설명하며(추론은 끝없이 상시 실행되는 서비스다), ㉣도 마찬가지다. 서빙 장애는 현재 사용자를 즉시 해치므로 비용이 크다. (1.3절)}

\qitem{3}{\textbf{③} 요청을 빠른 순서에서 느린 순서로 정렬했을 때 p99는 99번째 백분위의 지연 시간이다. 평균(①)은 느린 꼬리를 감추므로 서빙에서는 백분위 지표를 사용한다. (1.3절)}

\qitem{4}{\textbf{②} 관리형 K8s에서는 클러스터(컨트롤 플레인) 운영을 공급자에게 위임하지만, 이미지·배포 명세·서빙 구성은 사용자가 계속 책임진다. CaaS 수준이다. ③은 IaaS, ④는 서버리스 GPU를 설명한다. (1.4절, 표 1-2)}

\qitem{5}{\textbf{③} 트래픽이 적고 불확실하면 사용량 기반 API가 거의 언제나 더 저렴하며 더 빨리 출시할 수 있다. ①은 데이터 규제, ②는 규모, ④는 맞춤화에 관한 논거로, 모두 자체 서빙에 유리하다. (1.5절, 표 1-4)}

\qitem{6}{\textbf{②} 70 B × 2 bytes = 140 GB로, H100 한 대의 80 GB를 초과한다. 양자화하거나(예: 4-bit이면 약 35 GB) 여러 GPU에 샤딩해야 한다. 이런 추정에서 서빙 하드웨어 선택이 시작된다. (1.5절)}

\qitem{7}{\textbf{②} 배칭은 요청을 다른 요청과 함께 처리할 때까지 기다리게 하여 처리량과 비용을 얻는다. 가장 운이 나쁜 요청은 배치가 가장 많이 찰 때까지 기다리므로 대기 시간은 꼬리에 먼저 나타난다. 삼각형의 규율을 기억하자. 60\%/35\%의 이득도 실재하고 대가도 실재하므로, 엔지니어의 주장에는 둘 다 들어가야 한다. (1.3절)}

\qitem{8}{\textbf{③} 유휴 메모리는 기준선(가중치 + 런타임)이다. 요청별 활성값 메모리는 동시성과 입력 크기에 따라 늘어나므로 부하가 몰리면 유휴 상태보다 훨씬 높게 치솟는다. 6.6 GB 한도는 조용한 하루 내내 버티다가 처음으로 부하가 몰릴 때 서버를 처형한다. OOM 킬러가 보낸 SIGKILL이며 종료 코드는 137이다. 이 메커니즘(cgroup \code{memory.max})은 3주 차의 주제다. (1.3절)}

\subsection*{B. 단답형}

\qitem{1}{\textbf{SLO}(Service Level Objective, 서비스 수준 목표)다. 7주 차에는 SLI 및 SLA와 함께 발전시킨다.}

\qitem{2}{\textbf{콜드 스타트}다. GPU 서빙에서는 인스턴스 부팅 + 컨테이너 이미지 가져오기 + 모델 가중치 로딩이 겹쳐 수십 초에서 몇 분에 이른다. 12주 차에는 2주 차의 Firecracker와 함께 해결책을 다룬다.}

\qitem{3}{㉠ \textbf{처리량}, ㉡ \textbf{지연 시간, 특히 꼬리 지연 시간(p95/p99)}이다.}

\qitem{4}{13 B × 2 bytes/parameter = \textbf{≈ 26 GB}다. 24 GB L4에는 FP16으로 로드할 수 없어 양자화가 필요하다는 함의까지 덧붙이면 만점이다.}

\qitem{5}{\textbf{속도 제한 / 할당량}이다. 공급자의 허용 제어가 다른 테넌트와 공급자 자신의 p99를 보호하기 위해 사용자 등급의 상한을 넘는 요청을 거절한다. 의도대로 작동한 울타리이며, 해결책은 백오프(backoff) 후 재시도, 상위 등급 사용, 또는 직접 같은 울타리를 구축하는 자체 서빙(6주 차)이다.}

\subsection*{C. 서술형(채점 기준)}

\qitem{1}{핵심 명제는 정확도가 시스템 품질을 전혀 보장하지 않는다는 것이다. 동시성(동시 요청, 대기열), 지연 시간(경로의 합, 꼬리), 가용성(자동 복구, 무중단), 비용(유휴 GPU), 배포(롤링 교체, 롤백), 관측 가능성(메트릭이 없으면 고장 사실조차 알 수 없음) 중 세 가지 이상을 각각 서비스 맥락의 예와 연결하면 만점이다. (1.1절)}

\qitem{2}{모범 답안 예시: ① "작업과 서비스" → 학습에는 배치 스케줄러와 작업 대기열이 필요하고, 서빙에는 로드 밸런서, 오케스트레이터, 자가 치유가 필요하다. ② "처리량과 지연 시간" → 학습은 큰 배치로 GPU를 포화시키지만, 배칭은 지연 시간을 늘리므로 서빙은 동적 배칭으로 절충해야 한다. 부하 패턴 → 오토스케일링, 장애 비용 → 레플리카와 카나리 배포, 통신 패턴 → 수직 확장과 수평 확장도 옳다. 차이를 단순히 나열하지 않고 인프라 결과에 \textit{연결했는지}가 채점의 핵심이다. (1.3절)}

\qitem{3}{유일한 정답은 없으며 논증을 평가한다. 좋은 논리는 다음과 같다. 민감한 데이터(규제)와 미세 조정 계획(맞춤화)은 자체 구축을 가리키지만, 고객 다섯 곳이라는 작은 규모는 외부 도입을 가리킨다. 따라서 단계적 전략을 쓴다. (a) 현재는 개인정보를 보호하는 계약 아래 API로 빠르게 검증한다. (b) 규제로 불가능하다면 소형 오픈 모델의 자체 서빙으로 시작한다. (c) 미세 조정/성장 단계에는 자체 서빙으로 확실히 전환한다. 표 1-4의 기준을 명시적으로 사용하고 손익분기점을 어느 정도 추론해야 한다. (1.5절)}

\qitem{4}{(a) 레플리카당 1 ÷ 0.25 s = \textbf{4 req/s}다. (b) \textasciitilde{}80개 연결이 유지되므로 중앙값 요청은 \textasciitilde{}40개 요청 뒤에 선다. 40 × 0.25 ≈ \textbf{10 s}다. (c) 가장 긴 대기열은 \textasciitilde{}79개이므로 79 × 0.25 ≈ 19.8 s다. 20 s 타임아웃 바로 아래이므로 \textbf{아직은 504가 거의 발생하지 않는다}. 다만 서비스 시간의 작은 변동이나 사용자 한 차례만 더 몰려도 꼬리가 절벽을 넘는다. 계산을 제시하고 어느 쪽이든 논증하면 만점이다. 해결책은 다음 중 두 가지를 제시하고 꼭짓점을 밝혀야 한다. 레플리카 추가(지연 시간과 처리량을 얻고 비용을 지불), GPU 배칭(처리량/비용을 얻고 꼬리 지연 시간을 지불), 동시성 제한과 빠른 부하 차단(허용된 요청의 지연 시간을 보호하고 처리량 형태의 가용성을 지불), 더 작거나 양자화된 모델로 요청을 빠르게 처리(삼각형을 이동하며 엔지니어링 시간과 경우에 따라 품질을 지불)다. (1.1절, 1.3절)}

\namedsection{더 읽을거리}{}

더 깊이 파고들고 싶은 독자를 위한 자료다. ◆는 과정 교재를, ▶는 특정 주차에 다시 만날 자료를 표시한다.

\begin{tbbitemize}
  \item ◆ \textbf{Chip Huyen, Designing Machine Learning Systems (O'Reilly, 2022)} — 주 교재다. 이 장과 관련해서는 제1장(큰 관점의 ML 시스템)과 제7장(모델 배포와 예측 서비스)을 먼저 읽는다. ML을 시스템 관점에서 보는 능력을 기르기에 가장 좋은 단 한 권이다.
  \item ◆ \textbf{Chip Huyen, AI Engineering (O'Reilly, 2025)} — 파운데이션 모델 시대에 맞게 그림을 갱신한 부 교재다. 제9장(추론 최적화)은 이 책의 10\textasciitilde{}11주 차와 짝을 이룬다.
  \item \textbf{D. Sculley et al., "Hidden Technical Debt in Machine Learning Systems," NeurIPS 2015} — 1.1절에서 인용한 "ML 코드는 작은 상자" 그림의 출처다. 10쪽 남짓이므로 이번 주에 원문을 읽어 보자.
  \item ▶ \textbf{A. Agache et al., "Firecracker: Lightweight Virtualization for Serverless Applications," NSDI 2020} — 2주 차 필독 자료다. AWS Lambda 아래의 microVM 설계이자 이 장에서 예고한 콜드 스타트 이야기의 기술적 뿌리다.
  \item ▶ \textbf{W. Kwon et al., "Efficient Memory Management for Large Language Model Serving with PagedAttention," SOSP 2023} — 11주 차 필독 자료다. 운영체제 페이징을 LLM 서빙으로 이식한 vLLM의 핵심 아이디어를 다룬다. 이 책의 정신에 이보다 더 잘 맞는 논문은 없다.
  \item ▶ \textbf{Liz Rice, "Containers from Scratch" (GOTO 2018, video)} — 3주 차 실습의 원전이다. 수십 줄의 Go 코드로 아무것도 없는 상태에서 컨테이너를 만드는 30분짜리 영상이다. 보고 나면 컨테이너를 둘러싼 신비감이 사라진다.
  \item ▶ \textbf{Google SRE, Site Reliability Engineering (O'Reilly, 2016), Chapter 4 "Service Level Objectives"} — 온라인에서 무료로 읽을 수 있으며, 이 장이 서빙 삼각형으로 예고한 SLO 규율을 정석적으로 다룬다. 이 용어 체계를 그대로 도입하는 7주 차 전에 읽는다.
  \item \textbf{[이번 주의 흥미로운 읽을거리] ChatGPT 출시 후 첫 100일을 다룬 당시 보도} — 예를 들어 1억 사용자 추정치를 보도한 Reuters의 2023년 2월 1일 기사와 OpenAI 상태 페이지의 과거 기록이 있다. 팬이 아니라 서빙 엔지니어의 눈으로 읽어 보자. 모든 일화가 사상 최대 규모로 대중 앞에서 이루어진 용량, 대기열, 비용 결정이다.
  \item \textbf{기술 문서} — Kubernetes Concepts 페이지(5주 차 전에 훑어본다), OCI Runtime Specification, Linux cgroups v2 커널 문서, vLLM · NVIDIA Triton · TensorRT 공식 문서가 있다. 소설이 아니라 사전처럼 활용하라.
  \item \textbf{생태계 관찰} — GPU 가격과 공급자 지형은 고정된 참고 자료가 따라잡기 어려울 만큼 빠르게 변한다. 중요한 것은 확인하는 습관이다. 클라우드 가격 페이지, 서빙 엔진 릴리스 노트, NSDI · OSDI · SOSP · MLSys의 서빙 논문은 주기적으로 관찰하기 좋은 지점이다.
\end{tbbitemize}

\vspace{5.0pt}

\begin{calloutbox}{tbbblue}{tbbbluefill}
{\normalsize\bfseries\color{tbbblue}다음 장 — 제2장: 가상화 원리와 클라우드 컴퓨팅}\par\vspace{3pt}
이 장에서는 "클라우드에서 GPU VM을 빌린다"고 자연스럽게 말했다. 다음 장은 작은 충격으로 끝나는 심문에서 시작한다. 엄밀히 말해 빌린 머신은 존재하지 않는다. 임대의 실제 의미를 열어 본다. 물리 서버 하나를 여러 "가상" 머신으로 나누는 방식(하이퍼바이저, Intel VT-x / AMD-V), Linux 자체가 클라우드의 하이퍼바이저가 된 과정(KVM/QEMU, virtio), 밀리초 만에 부팅되는 VM을 가능하게 하는 기술(Firecracker microVM), 실재하는 GPU를 가상의 컴퓨터 안에 넣는 방식(PCIe 패스스루, vGPU, MIG)을 배운다. 실습에서는 클라우드 계정을 만들고, 사례 1-B를 기억하며 가장 먼저 과금 알림을 설정한다. 그런 다음 첫 VM을 시작해 같은 추론 작업에서 CPU와 GPU의 속도를 겨룬다.
\end{calloutbox}

%% Chapter 2 — 가상화 원리와 클라우드 컴퓨팅
\setcounter{chapter}{1}
\chapter{가상화 원리와 클라우드 컴퓨팅}
\chaptersubtitle{"GPU VM을 빌릴" 때 실제로 일어나는 일 — CPU 모드부터 microVM과 MIG까지}

\begin{calloutbox}{tbbblue}{tbbbluefill}
{\normalsize\bfseries\color{tbbblue}학습 목표}\par\vspace{3pt}
\begin{tbbitemize}
  \item 임대한 클라우드 VM을 \code{dmidecode}, \code{systemd-detect-virt}, \code{kvm-clock}로 심문해 물리적으로 존재하지 않는다는 증거를 읽고, 산업이 이 환상을 계속 다시 만드는 이유를 설명하는 60년의 사용률 역사(수백만 달러짜리 메인프레임 → 2000년대의 빈 서버 → EC2의 시간당 \$0.10)를 재구성한다.
  \item Popek과 Goldberg가 제시한 가상 머신 모니터의 세 요구 사항을 말하고, 고전적 x86이 이를 충족하지 못한 이유와 2005년까지 산업을 지탱한 두 소프트웨어 우회책인 바이너리 변환과 반가상화를 설명한다.
  \item VMX root/non-root 모드, VMCS, VM exit, EPT/NPT를 포함한 하드웨어 지원 가상화를 설명하고, "exit의 횟수와 비용"이라는 관점으로 가상화 성능을 분석한다.
  \item KVM, QEMU, virtio가 Linux에서 VM 실행 작업을 어떻게 나누는지, 반가상 장치가 에뮬레이션 장치보다 빠른 이유를 설명한다.
  \item 프로세스, 컨테이너, gVisor, microVM, 완전한 VM을 격리 스펙트럼 위에 놓고, AWS가 Lambda를 위해 Firecracker를 만든 이유를 포함해 주어진 멀티테넌트 워크로드에 어느 지점이 맞는지 논증한다.
  \item GPU가 워크로드에 도달하는 세 방식인 PCIe 패스스루, vGPU/시간 분할, MIG를 비교하고 성능 격리와 비용에 미치는 영향을 설명한다.
  \item 과금 안전장치를 가장 먼저 마련한 뒤 첫 클라우드 GPU VM을 시작하고, CPU와 GPU 추론을 벤치마크하며, 안전하게 종료한다.
\end{tbbitemize}
\end{calloutbox}

\vspace{5.0pt}

\begin{calloutbox}{tbbteal}{tbbtealfill}
{\normalsize\bfseries\color{tbbteal}핵심 용어}\par\vspace{3pt}
시분할(CP-67 / VM) · 서버 난립 · 사용률 · 하이퍼바이저 / VMM · Type 1 대 Type 2 · 트랩 후 에뮬레이션 · 민감 명령어 대 특권 명령어 · 바이너리 변환 · 반가상화 · VT-x / AMD-V · VMX root/non-root · VMCS · VM exit · EPT / NPT · KVM · QEMU · vCPU · 스틸 타임 · virtio · virtqueue · vhost · SR-IOV · IOMMU (VT-d) · microVM · Firecracker · gVisor · Kata Containers · PCIe 패스스루(VFIO) · vGPU · 시간 분할 · MPS · MIG
\end{calloutbox}

\vspace{5.0pt}

\begin{calloutbox}{tbblightgray}{tbbboxfill}
{\normalsize\bfseries\color{tbbink}장 구성}\par\vspace{3pt}
\textbf{2.1} 빌린 머신은 존재하지 않는다    \textbf{2.2} 하이퍼바이저와 트래핑 문제    \textbf{2.3} 하드웨어 지원 가상화: VT-x, VMCS, EPT    \textbf{2.4} KVM, QEMU, virtio: 클라우드의 하이퍼바이저가 된 Linux    \textbf{2.5} VM은 얼마나 작아질 수 있는가? microVM과 격리 스펙트럼    \textbf{2.6} 클라우드의 GPU: 인스턴스, 패스스루, 공유    \textbf{2.7} 실습: 첫 클라우드 GPU를 안전하게 사용하기 — 이후 요약, 연습문제, 더 읽을거리로 이어진다.
\end{calloutbox}

\section{빌린 머신은 존재하지 않는다}

이번 주 어느 시점에 \textbf{Launch instance}라고 쓰인 버튼을 누르고 40초가량 기다린 뒤 "자신의" 머신에 ssh로 접속할 것이다. 그 머신에는 호스트 이름, CPU, 메모리, 디스크, 그리고 이 과정에서 처음으로 GPU까지 있다. 재부팅할 수도 있고 망가뜨릴 수도 있다. 모든 면에서 자기 소유의 컴퓨터처럼 동작한다.

이 장에서 풀어낼 불편한 진실은 \textbf{그 컴퓨터가 존재하지 않는다}는 것이다. 버튼을 눌렀을 때 누군가 서버를 랙에 꽂은 것이 아니다. 데이터 센터 어딘가에는 128개 코어와 2테라바이트 RAM을 갖췄을 법한 거대한 물리 머신이 있다. 그 위에서 이와 같은 게스트 수십 개가 나란히 실행되며, 각자 하드웨어를 소유한다고 믿는다. 이 환상을 유지하는 소프트웨어가 \textbf{하이퍼바이저}, 또는 \textbf{가상 머신 모니터(VMM)}다. 환상 자체는 \textbf{가상 머신(VM)}으로, 물리 컴퓨터를 충실하고 격리된 소프트웨어 정의 사본으로 만든 것이다.

\subsection*{사례 A: 머신을 심문하다}

강한 주장에는 증거가 필요하다. 이 주장은 명령 여섯 개만으로 스스로 유죄를 입증한다. 머신은 파일에 자신의 본성을 기록해 두고 있으므로 물어보기만 하면 된다.

\begin{terminalbox}
\begin{Verbatim}[breaklines=true,breakanywhere=true,fontsize=\small,formatcom=\color{tbbtermfg},codes=\tbbcodeactive]
$ ssh ubuntu@203.0.113.7            # "your" machine, forty seconds old
$ sudo dmidecode -s system-manufacturer
Amazon EC2                          # the "manufacturer" is a retailer
$ sudo dmidecode -s system-product-name
g4dn.xlarge                         # the "model number" is a price tag
$ systemd-detect-virt
kvm                                 # the OS knows. it has always known.
$ lscpu | grep Hypervisor
Hypervisor vendor:   KVM
$ cat /sys/devices/system/clocksource/clocksource0/current_clocksource
kvm-clock                           # even TIME arrives via the hypervisor
$ nvidia-smi -L
GPU 0: Tesla T4 (UUID: GPU-8f61c81e...)   # and yet THIS is physically real
\end{Verbatim}
\end{terminalbox}
\figcaption{사례 2-A  임대한 머신을 심문한 결과. 펌웨어 문자열, 가상화 감지기, CPU 플래그, 시계까지 모든 계층이 이 컴퓨터가 소프트웨어임을 증언한다. 마지막 줄만은 예외이며, 2.6절에서 설명한다.}

자백을 한 줄씩 읽어 보자. 노트북이라면 \textit{LENOVO}와 모델 번호가 들어갈 DMI 펌웨어 문자열 필드에 소매업체 이름과 과금 SKU가 들어 있다. 하이퍼바이저가 작성했고 \textit{그 부분}까지 거짓말할 이유는 없었기 때문이다. \code{systemd-detect-virt}가 특별한 묘기를 부리는 것도 아니다. 게스트는 단순히 물어볼 수 있다(하이퍼바이저용으로 예약된 \code{CPUID} 리프이며, 2.3절에서 누가 답하는지 설명한다). 스케줄에서 밀려난 게스트는 혼자 정확한 시간을 유지할 수 없으므로 시계조차 반가상 서비스다. \code{kvm-clock}과 사촌인 \textit{스틸 타임}은 p99를 달고 2.4절에 다시 등장한다. 그런데 마지막 줄은 이 패턴을 거부한다. 디스크도 가상이고 NIC도 가상이며 시간 자체도 가상이지만, T4는 \textbf{통째로 패스스루된 실제 카드}다. 가장 비싼 구성 요소 하나만 빼고 전부 가짜인 머신은 기묘한 산물이다. 2.6절을 마치면 산업이 정확히 왜 이런 식으로 만드는지 알게 된다.

\subsection*{사례 B: 텅 빈 머신들}

애초에 이 환상은 왜 존재할까? 우아함 때문이 아니라 \textbf{돈} 때문이다. 50년 간격으로 두 번이나 그랬다. 이 과정이 이야기의 3막 안에 자리하기 때문에 세 문단을 들여 살펴볼 가치가 있다.

\textbf{1막: 머신 가격이 수백만 달러다(1960년대).} 당시 메인프레임은 수백만 달러였으며, 제작자가 처음 의도한 대로 한 번에 작업 하나만 실행하면 사람이 테이프를 바꾸는 작업 사이에 놀고 있었다. \textit{귀중한 머신 하나를 여러 사용자가 나눠 써야 한다}는 경제적 압력이 시분할을 낳았다. IBM Cambridge 연구소에서는 더 기묘하고 오래가는 산물인 \textbf{CP-67 (1967)}이 탄생했다. 사용자에게 로그인 세션이 아니라 \textit{각자의 완전한 System/360 시뮬레이션}을 제공하고, 그 안에서 사용자별 단일 사용자 OS를 실행하는 소프트웨어였다. 가상화는 클라우드를 위해 발명된 것이 아니다. 하드웨어가 비싸면서도 놀고 있었기 때문에 탄생했으며, UNIX보다도 오래되었다.

\textbf{2막: 머신이 텅 비어 있다(2000년대).} 시간을 빠르게 돌려 보자. 서버가 저렴해지자 기업은 애플리케이션마다 한 대씩 구매했다. 관리 편의도 이유였고, "메일 서버와 데이터베이스가 \textit{반드시} 싸울 것"이라는 미신도 한몫했다. 그 결과 \textbf{서버 난립}이 벌어졌다. 베이지색 상자가 랙을 가득 채웠고, 당시 조사에서 일반적인 사용률은 \textbf{5–15\%}였다. 사실상 기업 데이터 센터는 각각 전력과 냉각 비용을 태우는 유휴 CPU 창고였다. 2.2절에서 설명하듯 기술적 이유로 가상화를 거부하던 x86에 VMware가 1960년대의 묘수를 되살렸다(1999–2001). 대부분 놀고 있는 머신 10대를 한 대로 접는 \textit{통합}은 스스로 가치를 입증했다. 투자 수익률 논리는 말 그대로 "전기 요금 청구서 아홉 장을 없앤다"였다.

\textbf{3막: 계량기(2006년).} 2006년 8월 25일, Amazon은 Elastic Compute Cloud라는 서비스의 베타를 열었다. API로 빌리는 가상 머신이며 소형 인스턴스 가격은 \textbf{\$0.10 per hour}였다. 중요한 단어는 \textit{클라우드}가 아니라 \textit{시간}이었다. 할당이 소프트웨어 작업(이 장의 주제)이 되었으므로 계산 자원을 전기처럼 계량할 수 있었다. 조달도, 랙 설치도, 최소 사용량도 필요 없었다. 제1장에서 배운 모든 경제 구조, 즉 GPU 시간당 가격, 토큰당 API, 서버리스 과금, 사례 1-B의 \$779 교훈까지 모두 그 계량기에서 내려왔다. 그리고 이 이야기를 우리 이야기로 만드는 대칭이 있다. \textbf{GPU가 지금 다시 1막을 살고 있다.} H100은 이번 10년의 수백만 달러짜리 메인프레임이다. 터무니없이 비싸고 만성적으로 충분히 공유되지 않는다. 2.6절의 패스스루 → 시간 분할 → MIG는 CPU가 세 막에 걸쳐 밟은 여정(독점 사용 → 소프트웨어 공유 → 하드웨어 파티션)을 실리콘에서 다시 재생한다. 아키텍처에는 운율이 있고, 이 장은 그 운율 체계를 가르친다.

제1장에서는 "GPU VM을 빌린다"는 말을 믿고 받아들였다. 이 장에서는 상자를 연다. 안의 메커니즘은 세부 사항이 아니라 이 과정의 모든 주제가 서 있는 기반이기 때문이다. 먼저 이유 네 가지를 살펴보자.

\subsection*{클라우드가 가상화 위에 구축된 이유}

\textbf{첫째: 사용률의 경제성.} 제1장에서 AI 인프라 비용의 핵심 문제는 비싼 하드웨어가 놀고 있는 것이라고 배웠다. 고객 한 명에게 통째로 판매한 물리 서버는 그 고객이 쉴 때 함께 논다. 물리 서버를 가상 머신으로 나누면 부하의 고점과 저점이 엇갈리는 여러 테넌트를 수용할 수 있다. 공급자는 같은 실리콘을 여러 번 판매하고, 고객이 원하는 만큼 작은 조각을 원하는 시간만큼 제공할 수 있다. 제1장 표 전체의 경제적 전제인 시간당 과금은 할당이 소프트웨어 작업이기 때문에 가능하다.

\textbf{둘째: 낯선 사용자 사이의 격리.} VM은 절대 만날 일이 없는, 어쩌면 경쟁사일 수도 있는 테넌트와 물리 머신을 공유한다. 하이퍼바이저는 다른 테넌트가 내 메모리를 읽거나 CPU 몫을 소진하거나 커널을 충돌시키지 못하도록 보장하는 심판이다. 그 심판이 얼마나 강력하며 성능 비용은 얼마인지가 이 장을 관통하는 주제다. 워크로드가 "빌린 서버"에서 "200 ms 동안만 사는 서버리스 함수"로 작아지면 답이 달라지므로 2.5절 전체가 이 문제를 다룬다.

\textbf{셋째: 탄력성은 소프트웨어 호출 한 번이면 된다.} 제1장에서 추론 부하는 급증하며 해결책은 탄력적 스케일링이라고 배웠다. "머신 생성"이 하드웨어를 랙에 꽂는 일이 아니라 프로세스와 비슷한 소프트웨어 객체를 만드는 일이므로 탄력성이 물리적으로 가능하다. 12주 차에 오토스케일링을 배울 때 모든 수평 확장 이벤트의 바닥에는 이 장에서 설명하는 할당, 부팅, 로드, 서빙의 순서가 있다. 이 순서에 걸리는 시간이 \textbf{콜드 스타트}다. 하한은 가상화 기술이 정하며, AWS가 Lambda를 실용화하려고 새로운 VM인 Firecracker(2.5절)를 발명해야 했던 이유다.

\textbf{넷째: 서버군은 이기종이지만 VM은 그렇지 않다.} 공급자는 여러 세대의 하드웨어를 나란히 운영한다. 가상화 계층은 이 동물원을 안정적인 인스턴스 유형으로 정규화한다. 공급자는 게스트가 실행되는 동안에도 호스트를 마이그레이션하고 비우고 패치할 수 있으며, 사용자는 머신 전체를 파일로 스냅숏할 수 있다. 하드웨어 프로젝트가 될 작업이 API 호출로 바뀐다.

\subsection*{AI 서빙 엔지니어가 내부 구조를 알아야 하는 이유}

VMCS가 무엇인지 모르고도 수년 동안 클라우드를 사용할 수 있다. 그러나 이 과정은 의존하는 추상화 아래로 내려가겠다고 약속했고, 이 계층을 방문하면 특히 큰 보상을 얻는다.

\begin{tbbitemize}
  \item \textbf{p99의 뿌리가 이곳까지 내려온다.} vCPU는 호스트가 스케줄링하는 스레드다(2.4절). 과도하게 판매된 호스트에서 게스트는 "스틸 타임"을 겪는다. 애플리케이션을 아무리 튜닝해도 고칠 수 없는 보이지 않는 일시 중지가 꼬리 지연 시간 급증으로 드러난다. 이 사실을 알면 수수께끼가 진단으로 바뀐다.
  \item \textbf{콜드 스타트는 ML 문제이기 전에 가상화 문제다.} 0으로 스케일링(12주 차)의 성패는 격리된 실행 환경을 얼마나 빨리 만들어 내는지에 달렸다. 125 ms 만에 부팅하는 Firecracker microVM이라는 최신 기술을 분석하게 된다.
  \item \textbf{GPU 공유(12주 차)는 GPU 가상화의 응용이다.} H100 한 대가 테넌트 세 곳을 안전하게 서빙할 수 있는지는 패스스루, vGPU, MIG 중 무엇을 쓰느냐에 달렸다. 2.6절의 주제다.
  \item \textbf{이번 주 실습 청구서는 진짜 돈이다.} 인스턴스가 무엇인지, 그리고 stop, terminate, 연결된 스토리지가 아래에서 무엇을 의미하는지 이해하는 것이 \$2짜리 실습과 뜻밖의 \$200 청구서를 가른다(2.7절).
\end{tbbitemize}

\subsection*{이 장의 용어}

\begin{tbbtablewrap}
  \begin{tabular}{@{}>{\raggedright\arraybackslash}m{\dimexpr 0.2660\tbltextwidth-2\tabcolsep\relax} >{\raggedright\arraybackslash}m{\dimexpr 0.7340\tbltextwidth-2\tabcolsep\relax}@{}}
  \arrayrulecolor{tbbnavy}\specialrule{1pt}{0pt}{0pt}
  \rowcolor{tbbtablehead}{\centering \bfseries\color{tbbnavy}용어} & {\centering \bfseries\color{tbbnavy}의미} \\
  \arrayrulecolor{tbbnavy}\specialrule{0.7pt}{0pt}{2pt}
  {\textbf{가상 머신(VM)}} & {수정하지 않은 운영체제를 실행할 수 있을 만큼 물리 컴퓨터를 충실하게 소프트웨어로 복제한 것이다.} \\
  \arrayrulecolor{tbbtableborder}\hline
  {\textbf{하이퍼바이저 / VMM}} & {VM을 만들고 중재하는 소프트웨어다. 실제 하드웨어를 게스트 사이에서 다중화하면서 서로 격리한다.} \\
  \arrayrulecolor{tbbtableborder}\hline
  {\textbf{호스트 / 게스트}} & {물리 머신과 그 운영 환경(호스트), 그리고 VM 안에서 실행되는 가상화 환경과 OS(게스트)를 대비하는 말이다.} \\
  \arrayrulecolor{tbbtableborder}\hline
  {\textbf{가상화 대 에뮬레이션}} & {가상화는 안전할 때 게스트 명령어를 실제 CPU에서 직접 실행한다. 에뮬레이션은 소프트웨어로 해석하므로 유연하지만 느리다. 전통적으로 하이퍼바이저는 CPU를 가상화하고 장치를 에뮬레이션한다.} \\
  \arrayrulecolor{tbbtableborder}\hline
  {\textbf{트랩}} & {코드가 허용되지 않은 작업을 시도할 때 하드웨어가 제어권을 더 높은 특권 계층으로 강제로 넘기는 일이다. VMM이 기반으로 삼는 기본 메커니즘이다.} \\
  \arrayrulecolor{tbbtableborder}\hline
  {\textbf{VM exit / entry}} & {VT-x/AMD-V에서 게스트 실행이 하이퍼바이저로 전환되는 것(exit)과 다시 돌아가는 것(entry)이다. 가상화 오버헤드의 단위다.} \\
  \arrayrulecolor{tbbtableborder}\hline
  {\textbf{반가상화}} & {완벽한 환상을 유지하는 대신 게스트가 가상화 사실을 알고 효율적인 인터페이스로 협력하게 하는 방식이다. virtio가 현대의 후손이다.} \\
  \arrayrulecolor{tbbtableborder}\hline
  {\textbf{장치 모델}} & {게스트에 NIC나 디스크 같은 하드웨어 장치를 가장하는 하이퍼바이저 구성 요소다. QEMU의 주된 역할이다.} \\
  \arrayrulecolor{tbbtableborder}\hline
  {\textbf{microVM}} & {최소한의 장치 모델과 직접 커널 부팅을 사용해 밀리초 만에 시작하고 호스트당 수천 개를 실행하도록 설계한 VM이다(Firecracker, Cloud Hypervisor).} \\
  \arrayrulecolor{tbbtableborder}\hline
  {\textbf{IOMMU (VT-d / AMD-Vi)}} & {장치 DMA를 변환하고 통제하여 GPU 같은 물리 장치를 게스트에 직접 안전하게 넘길 수 있게 하는 하드웨어다.} \\
  \arrayrulecolor{tbbnavy}\specialrule{1pt}{2pt}{0pt}
  \end{tabular}\par
  \figcaption{표 2-1  이 장에서 사용할 용어}
\end{tbbtablewrap}

\begin{calloutbox}{tbbpurple}{tbbpurplefill}
{\normalsize\bfseries\color{tbbpurple}생각해 보기}\par\vspace{3pt}
\begin{tbbitemize}
  \item 사례 2-A의 명령을 물리 Linux 머신에서 실행해 보자. WSL2도 좋지만 먼저 WSL2가 무엇을 보고할지, 그리고 그 답 자체가 왜 흥미로운지 예측한다. 어느 줄이 다르며 각 차이가 무엇을 증명하는가?
  \item 클라우드 VM이 4 vCPUs를 보고한다. 호스트에서 이를 물리적으로 구현할 수 있는 방법 세 가지와 각 방법이 추론 지연 시간에 미치는 영향을 나열해 보자. (2.4절 뒤에 다시 살펴본다.)
  \item 공급자는 가상화를 건너뛰고 물리 서버 전체를 빌려 줄 수도 있다("베어 메탈" 상품이 실제로 존재한다). 공급자에게는 무엇이 더 어려워지고 고객에게는 무엇이 더 좋아지는가? 누가 어떤 이유로 베어 메탈을 선택하는가?
  \item 사례 2-B의 "GPU가 다시 1막을 살고 있다"는 비유를 확장해 보자. GPU의 3막, 즉 GPU의 2006년과 계량기는 어떤 모습일까? 오늘날의 서버리스 GPU 플랫폼(제1장)이 이에 해당하는가, 아니면 아직 부족한 것이 있는가?
\end{tbbitemize}
\end{calloutbox}

\section{하이퍼바이저와 트래핑 문제}

하이퍼바이저의 직무 설명은 처음 보면 불가능해 보인다. 머신을 소유하도록 작성된 소프트웨어인 완전하고 수정하지 않은 운영체제를 특권 없는 테넌트로 실행하면서, 스스로 그 사실을 눈치채거나 탈출하지 못하게 해야 한다. 이 일이 가능한 것은 오래되고 우아한 아이디어 하나 덕분이며, \textit{빠르게} 가능한 것은 2.3절에서 만날 하드웨어 기능 덕분이다. 먼저 아이디어부터 세워 보자.

\subsection*{중간 개입 아이디어와 Popek \& Goldberg의 규칙}

운영체제 커널은 페이지 테이블 전환, 인터럽트 비활성화, 장치 통신 같은 소수의 위험한 작업으로 하드웨어를 관리한다. 실제 머신에서 CPU의 특권 체계(x86에서는 \textbf{링}, 커널은 ring 0, 애플리케이션은 ring 3)는 커널만 이 작업을 수행하도록 보장한다. 고전적인 가상화 묘수인 \textbf{트랩 후 에뮬레이션(trap-and-emulate)}은 게스트 커널의 \textit{특권을 낮춰} 실행하는 방식이다. 게스트 코드는 실제 CPU에서 최고 속도로 직접 실행되지만, 위험한 작업을 시도하는 순간 하드웨어가 하이퍼바이저로 \textbf{트랩}한다. 하이퍼바이저는 게스트가 보는 가상 머신 상태, 즉 가상 인터럽트 상태, 가상 페이지 테이블, 가상 장치에 대해 동등하지만 안전한 작업을 수행한 뒤 게스트를 재개한다. 게스트는 거의 모든 시간에 실제 실리콘 위에서 실행되지만 실제 머신의 제어 장치에는 손대지 못한다.

클라우드보다 수십 년 앞선 1974년, Gerald Popek과 Robert Goldberg는 오늘날도 사용하는 용어 체계로 이 묘수가 언제 가능한지를 논문에서 형식화했다.

\begin{calloutbox}{tbbblue}{tbbbluefill}
{\normalsize\bfseries\color{tbbblue}상자 2-1  Popek \& Goldberg (1974): 진정한 VMM이 제공해야 할 것}\par\vspace{3pt}
\begin{tbbitemize}
  \item \textbf{동등성(충실성).} VM에서 실행되는 소프트웨어는 실제 하드웨어에서 실행될 때와 본질적으로 똑같이 동작한다. 게스트 OS에 가상화 사실을 알릴 필요가 없어야 한다.
  \item \textbf{효율성(성능).} 압도적으로 많은 게스트 명령어가 VMM 개입 없이 실제 CPU에서 직접 실행된다. "모든 것을 그냥 해석한다"는 방식은 배제된다. 그것은 가상화가 아니라 에뮬레이션이다.
  \item \textbf{자원 제어(안전성).} VMM이 물리 자원을 완전히 통제하며, 어떤 게스트 작업도 할당 범위 밖의 자원에 영향을 줄 수 없다.
\end{tbbitemize}

그리고 정리는 다음과 같다. \textbf{모든 민감 명령어(특권 머신 상태를 읽거나 바꾸는 명령어)가 특권 명령어(특권 없이 실행하면 트랩하는 명령어)이기도 할 때, 그리고 그때에만} 머신을 이 방식으로 가상화할 수 있다. 다시 말해 위험한 것은 모두 트랩할 수 있어야 한다. — "Formal Requirements for Virtualizable Third Generation Architectures," CACM 1974.
\end{calloutbox}

\vspace{3.0pt}

이 정리는 깔끔한 정신 모형을 제공한다. 게스트의 특권을 낮춰 실행하고, 위험한 작업은 트랩하게 하며, VMM에서 에뮬레이션한다. 동등성과 자원 제어는 트래핑에서 나오고, 효율성은 나머지 모든 것이 네이티브로 실행되는 데서 나온다.

\subsection*{x86의 난처함: 민감하지만 특권은 아니다}

문제가 하나 있었다. 세계에서 가장 중요한 CPU 아키텍처가 이 정리를 통과하지 못했다. Robin과 Irvine의 2000년 분석은 \textbf{고전적 x86에서 민감하지만 특권은 아닌 명령어 17개}를 목록화했다. 특권이 낮아진 게스트 커널에서 실행하면 올바르게 작동하지도, 트랩하지도 않는다. 그저 조용히 잘못된 일을 한다. 두 가지 예를 살펴보자.

\begin{tbbitemize}
  \item \code{POPF}는 인터럽트 활성화 플래그를 포함한 CPU 플래그 레지스터를 복원한다. 게스트 커널은 인터럽트를 켜거나 끌 생각으로 이를 사용한다. 특권이 낮은 상태에서는 인터럽트 비트가 \textit{조용히 무시된다}. 트랩도 오류도 없고, 게스트 커널은 자기 인터럽트 상태를 잘못 믿게 된다. 디버깅은 악몽이다.
  \item \code{SGDT}/\code{SIDT} 같은 명령어는 \textit{특권 검사 없이} 특권 테이블 레지스터를 읽어 실제 머신의 상태를 게스트에 누출하고 동등성의 환상을 깨뜨린다.
\end{tbbitemize}

따라서 x86에서는 순수한 트랩 후 에뮬레이션이 불가능했다. 하필 닷컴 시대가 서버 통합을 상업적으로 절실하게 만들던 때였다. 산업은 영웅적인 우회책 두 가지를 만들어 냈다. 그 아이디어가 오늘날의 스택 안에 살아 있으므로 둘 다 알아둘 가치가 있다.

\subsection*{우회책 1: 바이너리 변환(VMware, 1999)}

VMware의 통찰은 CPU가 문제의 명령어 17개를 트랩하지 않는다면 \textit{애초에 실행되지 않도록 게스트를 다시 쓰자}는 것이었다. VMM은 게스트 커널 코드가 실행되기 직전에 스캔하고 변환한다. 대부분 명령어는 그대로 복사하지만, 민감 명령어는 모두 VMM으로 진입하는 안전한 시퀀스로 바꾼다. 변환한 블록은 캐시하므로 자주 실행되는 코드는 한 번만 변환되고 이후에는 네이티브에 가까운 속도로 실행된다. 걱정할 민감 명령어가 없는 사용자 모드 게스트 코드는 손대지 않는다. 그 결과 정리의 문언을 위반한 하드웨어에서 Popek \& Goldberg의 정신을 충족했다. 명령어 스트림을 적시 수술하여 수정하지 않은 게스트를 실행하는 \textbf{완전 가상화}다. 자기 수정 코드, 정확한 예외, 변환 캐시 무효화 같은 공학 문제는 엄청나게 어려웠다. 하드웨어가 마침내 문제를 해결했을 때 산업이 그토록 고마워한 이유는 이 기억에 있다.

\subsection*{우회책 2: 반가상화(Xen, 2003)}

Xen 프로젝트는 정반대의 길을 택했다. \textbf{환상을 포기한다.} 동등성을 유지하는 부분이 비싸다면 게스트가 게스트라는 진실을 알려 주고, 하이퍼바이저와 통신할 정직하고 효율적인 인터페이스를 제공한다. 게스트 커널은 \textit{소스 수준에서 수정}된다. 특권 작업을 실행할 곳에서 대신 \textbf{하이퍼콜}을 수행한다. 시스템 호출과 같은 형태로 하이퍼바이저를 의도적으로 호출하는 것이다. 배치된 하이퍼콜이 트랩 폭풍을 대체하면서 성능은 뛰어났다. 대가는 게스트 커널 소스가 필요하고(Linux에는 괜찮지만 Windows에는 곤란하다), 수정 사항을 영원히 유지해야 한다는 점이다.

하드웨어 지원이 등장하자 순수한 형태의 반가상화는 사라졌지만, \textit{가상화 사실을 알고 협력하는 게스트가 속여야 하는 게스트보다 빠르다}는 핵심 아이디어는 오늘날 모든 클라우드 VM의 중심에 살아 있다. \textbf{반가상 장치}인 virtio로 이어졌으며 2.4절에서 다룬다. Xen 자체도 초기 클라우드를 떠받쳤다. EC2는 2006년 Xen에서 출시되었고, AWS는 2017년에야 KVM 기반 Nitro 시스템으로 교체하기 시작했다.

\subsection*{Type 1과 Type 2 — KVM의 위치}

1970년대 문헌에서 나온 또 하나의 고전적 구분은 하이퍼바이저가 \textit{어디서 실행되는지}에 따라 나눈다(그림 2-1). \textbf{Type 1}(베어 메탈) 하이퍼바이저는 머신에서 가장 낮은 소프트웨어 계층 자체다. VMware ESXi, Microsoft Hyper-V, Xen이 있다. 데이터 센터는 공격 표면이 가장 작고 아래에서 경쟁적으로 스케줄링하는 호스트 OS가 없는 Type 1을 사용한다. \textbf{Type 2}(호스팅형) 하이퍼바이저는 일반 OS 위의 애플리케이션으로 실행된다. 노트북의 VirtualBox나 VMware Workstation처럼 개발에 편리하지만, VM 아래에 호스트 OS 스케줄링 계층이 하나 더 있다.

\begin{tbbfigure}
  \adjustbox{max width=0.94\linewidth,max totalheight=0.46\textheight}{\includegraphics{figures/fig2-1_hypervisors.png}}\\[3pt]
  \figcaption{그림 2-1  Type 1(베어 메탈)과 Type 2(호스팅형) 하이퍼바이저. KVM은 경계를 흐린다. 이미 베어 메탈에서 실행되는 Linux 커널을 하이퍼바이저로 바꾸는 커널 모듈이다.}
\end{tbbfigure}

세계에서 가장 중요한 하이퍼바이저는 이 이분법을 거부한다. \textbf{KVM}은 Linux 커널 모듈이므로 "호스트 OS"와 "하이퍼바이저"가 같은 것이 된다. 이미 베어 메탈에서 실행 중인 Linux 자체가 VMM 능력을 얻는다. 실용적으로는 \textit{Type 2의 사용 편의성을 갖춘 Type 1}로 분류할 수 있다. 아키텍처는 하이퍼바이저급이지만, VM은 일반 Linux 도구로 일반 Linux 프로세스처럼 관리한다. 이 조합에 오픈 소스라는 장점까지 더해져 KVM은 Google Cloud, 현대 AWS (Nitro), OpenStack 클라우드, 그리고 이 장의 2.4절을 떠받친다.

\begin{calloutbox}{tbbpurple}{tbbpurplefill}
{\normalsize\bfseries\color{tbbpurple}생각해 보기}\par\vspace{3pt}
\begin{tbbitemize}
  \item Docker Desktop의 Linux VM을 실행하는 MacBook, GCP 호스트, WSL2를 실행하는 Windows를 분류해 보자. (힌트: WSL2는 Hyper-V에서 실행되며, Hyper-V는 실행 중인 Windows \textit{아래로} 들어간다.) 각각 어느 유형이며, 이 질문에 깔끔한 답이 있기는 한가?
  \item 바이너리 변환은 복잡성을 대가로 동등성을 보존했고, 반가상화는 속도를 위해 동등성을 희생했다. 오늘날에도 각 절충이 옳은 상황을 하나씩 제시해 보자.
\end{tbbitemize}
\end{calloutbox}

\section{하드웨어 지원 가상화: VT-x, VMCS, EPT}

2005\textasciitilde{}2006년, Intel과 AMD는 소프트웨어가 몸을 비틀며 피하려 했던 일을 했다. 프로세서를 바꿨다. Intel의 \textbf{VT-x}와 AMD의 \textbf{AMD-V}는 명령어 17개를 하나씩 고치는 대신 CPU에 완전히 새로운 특권 차원을 추가하여 x86이 마침내 Popek \& Goldberg의 조건을 충족하게 했다. 이론에 앞서 실제 머신에서 해결책을 찾아보자. \code{grep} 한 번이면 된다.

\begin{terminalbox}
\begin{Verbatim}[breaklines=true,breakanywhere=true,fontsize=\small,formatcom=\color{tbbtermfg},codes=\tbbcodeactive]
# on any modern laptop / workstation running Linux:
$ grep -m1 -c -E "vmx|svm" /proc/cpuinfo
1                          # vmx = Intel VT-x, svm = AMD-V. 2005's fix.
$ ls -l /dev/kvm
crw-rw---- 1 root kvm 10, 232 Jul  6 09:14 /dev/kvm
#   ^ the kernel offering hypervisor powers as a device file (sec. 2.4)

# now the same two commands on this week's rented EC2 guest:
$ grep -m1 -c -E "vmx|svm" /proc/cpuinfo
0                          # the flag is gone. and /dev/kvm does not exist.
# the hypervisor HID it: your guest may not host guests of its own.
# the illusion, by default, refuses to nest.
\end{Verbatim}
\end{terminalbox}
\figcaption{기록 2-1  VT-x 찾기. 물리 노트북은 vmx/svm을 알리고 /dev/kvm을 공개한다. 임대한 VM 안에서는 하이퍼바이저가 이 기능을 마스킹한다. 곧 만날 VMCS가 선택한 CPUID 비트 하나다.}

\subsection*{두 세계: VMX root 모드와 non-root 모드}

VT-x는 링과 \textit{직교하는} 모드 전환을 추가한다(그림 2-2). 하이퍼바이저는 실제 머신을 완전히 제어하며 자체 rings 0–3을 지닌 \textbf{VMX root 모드}에서 실행된다. 게스트는 하드웨어의 감독을 받으면서 \textit{역시} rings 0–3을 지닌 \textbf{VMX non-root 모드}에서 실행된다. 이 문장을 두 번 읽어 보자. 오래된 딜레마가 사라진다. \textbf{게스트 커널은 다시 ring 0에서}, 작성된 그대로 실행되므로 보기 흉한 링 압축 묘수가 필요 없다. non-root ring 0은 진짜 ring 0이 아닐 뿐이다. CPU 자체가 게스트 실행을 감시하며, 하이퍼바이저가 개입할 가치가 있다고 표시한 작업이 일어나면 하드웨어가 \textbf{VM exit}를 수행해 자동으로 세계 전체를 root 모드로 전환한다. 처리가 끝나면 하이퍼바이저가 \textbf{VM entry}를 실행하고 게스트는 아무것도 모른 채 재개한다. 마침내 \textit{하드웨어}가 트래핑하는 트랩 후 에뮬레이션이 완성되고, 문제의 명령어 17개도 다른 것처럼 트랩한다.

\begin{tbbfigure}
  \adjustbox{max width=0.94\linewidth,max totalheight=0.46\textheight}{\includegraphics{figures/fig2-2_vtx.png}}\\[3pt]
  \figcaption{그림 2-2  한 장으로 보는 VT-x. 왼쪽은 2005년 이전의 문제로, 특권이 낮아진 게스트 커널과 조용히 실패하는 명령어를 보여 준다. 오른쪽은 root/non-root 모드, VM exit와 entry, 규칙과 상태를 기록하는 VMCS다.}
\end{tbbfigure}

\subsection*{VMCS: 두 세계 사이의 계약}

세계를 전환할 때는 두 세계의 상태와 교전 규칙을 보관할 곳이 필요하다. 그것이 가상 CPU마다 하나씩 존재하는 \textbf{VMCS (Virtual Machine Control Structure)}다. (a) entry 때 로드하고 exit 때 저장하는 게스트 레지스터 상태, (b) exit 때 복원하는 호스트 상태, (c) 하이퍼바이저가 하드웨어를 프로그래밍하는 \textbf{제어 필드}를 담는다. 어떤 이벤트가 exit를 일으키는지(\code{CPUID}가 exit해야 하는가? 모든 I/O 포트 접근은? 이 제어 레지스터에 쓰기는?), 특정 기능을 게스트에 어떻게 보일지, exit 핸들러가 어디에 있는지를 정한다. exit가 발생하면 하드웨어는 \textbf{exit 이유}도 VMCS에 기록한다. 따라서 하이퍼바이저의 핵심은 문자 그대로 exit 이유를 처리하는 거대한 \code{switch} 문이다. 2.4절에서 KVM 실행 루프를 볼 때 바로 이 구조를 구동한다.

\subsection*{Exit는 가상화 성능의 통화다}

VM exit는 아키텍처상 무거운 작업이다. 한 세계를 저장하고 다른 세계를 로드한다. 캐시 효과까지 세면 왕복에 대략 마이크로초가 걸리며, 최선의 경우에도 CPU 수천 사이클이 든다. exit가 드문 게스트는 사실상 네이티브 속도로 실행되지만, 계속 exit하는 게스트는 기어간다. 이 한 문장이 20년의 가상화 엔지니어링을 정리한다.

\begin{tbbitemize}
  \item \textbf{Exit 횟수를 줄인다.} 과거 하이퍼바이저가 처리하던 작업을 하드웨어가 non-root 모드에서 처리하게 한다(아래의 EPT, 인터럽트용 APICv). 레지스터 쓰기마다 exit하지 않고 exit 한 번에 작업을 배치하는 장치 인터페이스를 설계한다(virtio, 2.4절).
  \item \textbf{남은 exit를 저렴하게 만든다.} CPU 세대가 거듭되면서 세계 전환 지연 시간이 줄었다. 하이퍼바이저는 빈번한 exit를 가벼운 커널 경로로 보내고 드물고 복잡한 것만 사용자 공간에 맡긴다.
\end{tbbitemize}

실습 이후에도 이 관점을 유지하자. 가상화된 워크로드의 성능이 수수께끼처럼 낮으면 실무자가 가장 먼저 묻는 질문은 \textit{무엇이 얼마나 자주 exit하는가}다. 호스트의 \code{perf kvm stat} 같은 도구는 exit를 이유별로 표로 만들어 정확히 답한다. VM의 "exit 폭풍"은 프로세스의 시스템 호출 폭풍과 같다.

\subsection*{두 번째 문제: 섀도 페이지에서 EPT까지}

CPU 가상화는 전투의 절반에 불과했다. 나머지 절반은 메모리다. 이제 주소 변환에 \textbf{두 겹의 거짓말}이 있기 때문이다(그림 2-3). 게스트 OS는 게스트 가상 주소를 \textbf{게스트 물리} 주소로 변환하는 페이지 테이블을 만든다. 그러나 게스트 물리 주소 자체도 허구이며, 게스트 페이지가 \textbf{호스트 물리} 메모리의 어디에 실제로 위치하는지는 하이퍼바이저가 정한다. 2008년 이전의 하드웨어 페이지 테이블 워커는 테이블 집합 하나만 따라갈 수 있었다. 두 변환을 합칠 무언가가 필요했다.

\begin{tbbfigure}
  \adjustbox{max width=0.94\linewidth,max totalheight=0.46\textheight}{\includegraphics{figures/fig2-3_ept.png}}\\[3pt]
  \figcaption{그림 2-3  메모리 가상화. 게스트 페이지 테이블(게스트 관리)과 EPT/NPT(하이퍼바이저 관리)의 두 변환을 모두 하드웨어가 순회하며, TLB가 합성된 결과를 캐시한다.}
\end{tbbfigure}

소프트웨어 시대의 답은 \textbf{섀도 페이지 테이블}이었다. 하이퍼바이저가 게스트 주소 공간마다 게스트 가상→호스트 물리 변환을 합친 테이블을 유지하고 실제 하드웨어가 이를 순회했다. 게스트 페이지 테이블을 쓰기 보호하고 \textbf{모든 게스트 페이지 테이블 갱신을 트랩}하여 일관성을 유지했다. 프로세스를 포크하거나 메모리의 넓은 영역에 접근하는 워크로드에서는 exit 폭풍과 고통스러운 성능 저하가 발생했다. 섀도 페이징은 가상화 오버헤드와 복잡성의 가장 큰 단일 원인이었다.

하드웨어의 답은 Intel의 \textbf{EPT}(Extended Page Tables, \textasciitilde{}2008)와 AMD의 \textbf{NPT}로 등장했다. CPU가 게스트 자체의 페이지 테이블에 \textit{더해} 순회하는, 하이퍼바이저 소유의 \textit{두 번째} 페이지 테이블 집합이다. 게스트는 트랩, 섀도, exit 없이 자유롭게 테이블을 갱신한다. 대가는 정직하고 눈에 보인다. 이제 TLB 미스에는 2차원 순회가 필요하다. 최악에는 게스트의 4–5개 레벨마다 자체 EPT 순회가 필요해 최대 \textasciitilde{}24번의 메모리 참조가 일어난다. 따라서 \textbf{VM에서는 TLB 미스가 더 아프다.} 일반적인 완화책은 4 KB 대신 2 MB / 1 GB인 \textbf{대형 페이지}를 사용해 TLB 압력을 크게 줄이는 것이다. 추론 서버는 수 기가바이트 가중치 텐서를 보관하는 메모리 집약적 워크로드라 이 차이를 알아차리므로 이 과정에서도 기억할 가치가 있다. EPT와 exit가 줄어든 인터럽트 처리를 함께 사용하면 잘 구성한 VM의 계산 및 메모리 성능은 베어 메탈과 몇 퍼센트 이내 차이에 그친다. 클라우드를 그 위에 세울 수 있는 이유다.

CPU와 메모리는 해결했다. 세 번째 기둥인 \textbf{I/O}에서 아키텍처 이야기가 특히 흥미로워진다. Linux, QEMU, virtio라는 반가상 아이디어가 등장하기 때문이다. 다음 절의 주제다.

\begin{calloutbox}{tbbpurple}{tbbpurplefill}
{\normalsize\bfseries\color{tbbpurple}생각해 보기}\par\vspace{3pt}
\begin{tbbitemize}
  \item 다음 게스트 작업을 "계속 exit함"부터 "전혀 exit하지 않음"까지 순서대로 놓고 근거를 설명해 보자. 두 레지스터 더하기, 에뮬레이션 직렬 포트에서 한 바이트 읽기, TLB 히트인 메모리 로드, CPUID 실행, 자체 페이지 테이블 갱신(EPT 시대).
  \item EPT는 게스트 페이지 테이블 갱신에서 exit를 없앴지만 TLB 미스의 비용을 높였다. 30 GB LLM 가중치 버퍼를 4 KB 페이지와 1 GB 대형 페이지로 접근할 때 각각 대략 몇 개의 TLB 엔트리가 필요한가? 무엇을 시사하는가?
\end{tbbitemize}
\end{calloutbox}

\section{KVM, QEMU, virtio: 클라우드의 하이퍼바이저가 된 Linux}

VT-x는 모든 운영체제에 하이퍼바이저의 원재료를 제공했다. Linux는 거의 즉시 이를 차지했다. \textbf{KVM (Kernel-based Virtual Machine)}은 2007년 초 메인라인 커널에 병합되었으며, 재사용 설계의 작은 걸작이다. 하이퍼바이저 역할을 위해 스케줄러, 메모리 관리자, 드라이버를 모두 갖춘 새 운영체제를 만드는 대신, KVM은 Linux에 이미 세 기능의 세계적 구현이 \textit{모두} 있다는 점에 주목하고 VT-x를 구동하는 법만 가르쳤다. 그 결과가 앞으로 사용할 클라우드 대부분의 아래에 있는 소프트웨어 스택이며, 작업을 세 부분으로 나눈다(그림 2-4).

\begin{tbbfigure}
  \adjustbox{max width=0.94\linewidth,max totalheight=0.46\textheight}{\includegraphics{figures/fig2-4_kvm_virtio.png}}\\[3pt]
  \figcaption{그림 2-4  KVM/QEMU/virtio 스택. KVM(커널)은 VT-x를 구동하고, QEMU(VM당 사용자 공간 프로세스 하나)는 장치를 제공하며, virtio는 하드웨어의 환상을 정직하고 빠른 인터페이스로 바꾼다. vCPU는 그저 호스트 스레드다.}
\end{tbbfigure}

\subsection*{작업 분담}

\begin{tbbitemize}
  \item \textbf{KVM(커널 안)}은 작은 ioctl API인 \code{/\allowbreak dev/\allowbreak kvm}을 공개한다. VM 생성, 메모리 제공, vCPU 생성, 그리고 핵심인 \code{KVM\_RUN}이다. VMCS를 프로그래밍하고 VM entry를 실행하며 VM exit를 받아 처리한다. 빈번하고 단순한 exit(EPT 폴트, 대부분의 MSR 접근)는 커널에서 바로 처리한다.
  \item \textbf{QEMU(사용자 공간, VM당 프로세스 하나)}는 나머지 모든 일을 한다. 게스트 메모리(그저 프로세스 메모리다!)를 할당하고, 펌웨어나 커널을 로드하며, 대표적인 역할인 \textbf{장치 모델}을 제공한다. NIC, 디스크, 직렬 포트와 마더보드의 나머지 장치를 소프트웨어로 가장한다. 장치 I/O 때문에 VM exit가 발생하면 KVM이 QEMU로 넘긴다. QEMU 장치 코드는 게스트가 "하드웨어"에 무슨 일을 하려 했는지 파악하고 일반 Linux 시스템 호출을 통해 현실의 동등한 작업을 한다. 실제 소켓에 패킷을 보내거나 실제 파일에 블록을 쓴다.
  \item \textbf{호스트 Linux 커널}은 둘 다 다시 구현하지 않는 기능을 제공한다. 스케줄러가 vCPU를 실행하고, 메모리 관리자가 게스트 RAM을 뒷받침하며, 드라이버와 파일 시스템이 실제 I/O를 운반한다.
\end{tbbitemize}

게스트를 움직이는 실행 루프를 한 번은 문장으로 살펴볼 가치가 있다. vCPU 스레드가 \code{ioctl(KVM\_RUN)}을 호출하고, 커널이 VM entry를 수행한다. 게스트는 무언가 exit할 때까지 non-root 모드에서 네이티브로 계산한다. KVM은 exit를 처리하고 즉시 다시 entry하거나, 장치 작업을 위해 QEMU로 돌아갔다가 다시 entry한다. VM의 수명 내내 이 과정이 반복된다. 2.3절에서 배운 exit 이유와 exit 비용은 모두 이 루프 내부의 경제성이다.

\subsection*{vCPU는 그저 스레드다 — p99가 신경 써야 하는 이유}

그림 2-4에서 vCPU가 어디에 있는지 다시 보자. \textbf{일반 호스트 스레드}다. 그림의 말만 믿지 않아도 된다. KVM 게스트를 실행하는 Linux 머신이라면 어느 것이든 호스트가 직접 보여 준다(virt-manager를 사용하는 노트북이면 충분하다).

\begin{terminalbox}
\begin{Verbatim}[breaklines=true,breakanywhere=true,fontsize=\small,formatcom=\color{tbbtermfg},codes=\tbbcodeactive]
# on the HOST, with one 4-vCPU guest running:
$ pstree -p $(pgrep -of qemu) | head -5
qemu-system-x86(9114)-+-{CPU 0/KVM}(9131)
                      |-{CPU 1/KVM}(9132)
                      |-{CPU 2/KVM}(9133)
                      |-{CPU 3/KVM}(9134)
$ ps -o pid,psr,pcpu,comm -Lp 9114 | grep CPU  # psr = core hosting it
 9131  3  99.2  CPU 0/KVM     # "processor 0" = thread 9131, on core 3
 9132 11   1.0  CPU 1/KVM     # idle vCPU, parked on core 11
# renice it, pin it, cgroup it - everything Linux does to threads applies.
\end{Verbatim}
\end{terminalbox}
\figcaption{기록 2-2  호스트에서 본 게스트의 "프로세서" 네 개. QEMU 프로세스 하나 안에 있는 스케줄 가능한 스레드 네 개다. 이후의 오버커밋, 스틸 타임, 전용 코어는 모두 평범한 스레드 스케줄링이 의상을 입은 것이다.}

이 한 가지 사실에서 운영 중 만나게 될 여러 클라우드 현상이 나온다.

\begin{tbbitemize}
  \item \textbf{오버커밋.} 호스트는 테넌트의 최대 부하가 동시에 오지 않을 것에 베팅하며 물리 코어보다 많은 vCPU를 만들 수 있다. 항공 좌석과 같은 경제성이다. 저가 인스턴스 계열은 이 베팅을 바탕으로 가격이 정해진다.
  \item \textbf{스틸 타임.} vCPU가 실행 가능할 때 호스트 스케줄러가 \textit{다른 누군가}를 실행하면 게스트는 사이클을 얻지 못한 채 시간이 흐르는 경험을 한다. 게스트 안의 \code{top}에서는 \code{\%st}로 보인다. 지연 시간에 민감한 서빙에서 스틸 타임은 순수한 꼬리 지연 시간 독극물이다. 빼앗긴 슬라이스에 도착한 요청은 보이지 않게 기다린다. 진지한 서빙 서버군이 전용 코어 인스턴스 유형을 구매하는 이유이며, "CPU 그래프는 멀쩡한데 p99가 급증한다"는 질문의 답이 VM 한 계층 아래에 있을 수 있는 이유다.
  \item \textbf{시끄러운 이웃, 그러나 경계가 있다.} 스케줄러와 cgroup 가중치(제3장!)는 테넌트가 서로 해칠 수 있는 정도를 제한한다. 2.1절의 심판은 다음 주에 직접 다룰 바로 그 커널 메커니즘으로 구현된다.
\end{tbbitemize}

스틸 타임은 별도 사례로 살펴볼 가치가 있다. 프로덕션 텔레메트리에서 만날 가능성이 가장 높은 가상화 산물이며, 애플리케이션 문제로 오진하기도 가장 쉽다.

\begin{terminalbox}
\begin{Verbatim}[breaklines=true,breakanywhere=true,fontsize=\small,formatcom=\color{tbbtermfg},codes=\tbbcodeactive]
# inside a guest on a busy (oversold) host, during a load test:
$ top -bn1 | grep "%Cpu"
%Cpu(s): 41.2 us, 3.1 sy, 0.0 ni, 30.5 id, 0.2 wa, 0.0 hi, 0.6 si, 24.4 st
#                                                                  ^^^^^^^
# st = steal: 24.4% of this guest's runnable time went to... someone else.
# no process inside the guest will ever show this in its profile.
$ grep -A1 steal /proc/stat | head -2      # the raw counter, if you script it
cpu  184023 340 44522 830112 902 0 4211 267001 0 0
#                                        ^^^^^^ jiffies stolen since boot
\end{Verbatim}
\end{terminalbox}
\figcaption{기록 2-3  현장에서 붙잡은 스틸 타임. 게스트는 실행 가능했지만 호스트 스케줄러가 다른 테넌트를 실행했고, 요청은 볼 수 없는 슬라이스를 기다렸다. 깨끗한 CPU 그래프와 함께 p99가 급증한다면 한 계층 아래에 답이 있다.}

\subsection*{장치 문제와 virtio의 정직한 답}

이제 가상화의 세 번째 기둥인 I/O를 살펴보자. 완전한 환상을 제공하는 방식은 \textbf{장치 에뮬레이션}이다. QEMU가 오래되고 유명한 특정 실제 NIC를 가장하고(고전적인 예는 Intel e1000), 게스트는 기본 드라이버를 실행한다. e1000 드라이버가 있는 모든 OS가 바로 작동하므로 동등성은 완벽하다. 그러나 대가를 생각해 보자. 이 드라이버는 레지스터 쓰기가 나노초밖에 걸리지 않는 하드웨어용으로 작성되었다. 에뮬레이션 장치에서는 \textbf{모든 레지스터 접근이 사용자 공간으로 향하는 VM exit}다. 패킷 하나를 보내면서 장치를 여러 번 건드린다. exit, exit, exit다. 에뮬레이션 I/O는 세계 전환에 빠져 네이티브보다 10배 이상 느릴 수 있다.

해결책은 2.2절의 반가상 아이디어를 가장 잘 맞는 곳에서 되살리는 것이다. \textbf{virtio}(Rusty Russell, 2008, 현재 표준)는 가장을 멈추라고 말한다. 게스트에 \textit{공공연히 가상인} 장치를 제공하되, 소프트웨어 구성 요소 두 개가 공유 메모리로 버퍼를 전달하는 실제 상황에 맞게 인터페이스를 설계한다. 게스트는 \textbf{프런트엔드 드라이버}(\code{virtio-net}, \code{virtio-blk}, ...)를 로드하고, 하이퍼바이저 쪽은 짝이 되는 \textbf{백엔드}를 실행한다. 둘 사이에는 \textbf{virtqueue}가 있다. 게스트가 버퍼 디스크립터를 배치로 올리고 백엔드가 소비하는 게스트 메모리의 링 버퍼다. 게스트는 레지스터마다가 아니라 \textit{배치마다} 가벼운 exit 한 번으로 백엔드를 "깨운다". 완료 신호는 인터럽트로 돌아오거나 폴링 중에 수거한다. 부하가 걸리면 링이 계속 바쁘게 돌아가므로 exit가 거의 없다. 한 계층 아래에서 10주 차 동적 배칭과 같은 물리 법칙이 작동한다. \textbf{비싼 경계 통과 비용을 여러 작업 단위에 분산한다.}

두 가지 개선으로 그림이 완성된다. \textbf{vhost}는 백엔드를 QEMU에서 호스트 커널로 옮겨 데이터 경로의 사용자 공간 홉을 없앤다(vhost-user는 DPDK 기반 네트워크 스택처럼 전용 사용자 공간 스위치로 옮긴다). 최고 등급에서는 하드웨어가 다시 등장한다. \textbf{SR-IOV}는 물리 NIC 하나가 가벼운 \textit{가상 함수} 여러 개를 제공하게 하며, 각 함수를 IOMMU를 통해 게스트에 곧바로 안전하게 할당할 수 있다. 실리콘 자체가 장치 "에뮬레이션"을 하는 셈이다. 이 생각을 붙들어 두자. 실제 하드웨어를 게스트에 할당하는 방식은 2.6절에서 GPU가 VM에 도달하는 방법과 정확히 같다.

\begin{calloutbox}{tbbteal}{tbbtealfill}
{\normalsize\bfseries\color{tbbx115E59}상자 2-2  AWS Nitro: 가상화 비용을 0에 가깝게}\par\vspace{3pt}
이 절의 논리를 끝까지 따라가면 현대 EC2 아래의 시스템인 Nitro에 도달한다. AWS는 장치 모델을 \textit{소프트웨어 밖으로 완전히} 옮겼다. 네트워킹, 스토리지, 관리는 전용 \textbf{Nitro 카드}(맞춤형 실리콘)에서 실행되고 게스트에는 직접 할당된 하드웨어로 제공된다. 호스트 CPU에는 의도적으로 최소화한 KVM 기반 하이퍼바이저만 남으며 데이터 경로에 QEMU가 없다. 게스트는 베어 메탈에 가까운 성능을 얻고, 호스트는 배관에 코어를 거의 쓰지 않는다. 같은 아키텍처가 진정한 베어 메탈 인스턴스도 제공한다(Nitro 카드가 마더보드의 I/O \textit{자체}다). 역사의 흐름은 깔끔하다. EC2는 반가상화된 Xen에서 시작했고(2006), 20년에 걸친 exit 제거는 exit가 하던 일을 하드웨어로 옮기며 끝났다(2017–).
\end{calloutbox}

\vspace{3.0pt}

\begin{calloutbox}{tbbpurple}{tbbpurplefill}
{\normalsize\bfseries\color{tbbpurple}생각해 보기}\par\vspace{3pt}
\begin{tbbitemize}
  \item 게스트의 모델 서버에서 클라이언트로 나가는 HTTPS 응답 하나를 추적해 보자. (a) 에뮬레이션 e1000, (b) virtio + vhost, (c) SR-IOV에서 각 바이트가 그림 2-4의 어느 구성 요소를 통과하는가? 세계 전환 횟수를 정성적으로 세어 보자.
  \item 부하 테스트 중 지연 시간이 급증할 때 게스트 안의 \code{top}에 25\% steal이 나타난다. 클라우드 공급자를 바꾸지 않고 선택할 수 있는 방법은 무엇인가? (인스턴스 계열과 "전용"의 비용을 생각해 보자.)
\end{tbbitemize}
\end{calloutbox}

\section{VM은 얼마나 작아질 수 있는가? microVM과 격리 스펙트럼}

지금까지는 몇 시간에서 몇 년 동안 사는 임대 서버라는 고전적인 클라우드 단위를 가정했다. 2010년대에는 새로운 극단이 등장했다. \textbf{밀리초에서 몇 분} 동안 살고, \textbf{서로 신뢰하지 않는 낯선 사용자}에게 속하며, 경제성을 확보하려면 \textbf{호스트 한 대에 수천 개}를 넣어야 하는 서버리스 함수와 요청별 샌드박스다. 환경 생성에 40 s와 500 MB가 든다면 100 ms 단위로 과금할 수 없으므로, 이 규모에서는 고전적 VM의 형태가 맞지 않는다. 산업의 대응은 가상화 자체를 바꾸었다. 이를 이해하려면 \textit{테넌트 코드와 호스트 커널 사이에 무엇이 있으며 각 계층의 비용은 얼마인가?}라는 스펙트럼으로 보아야 한다(그림 2-5).

\begin{tbbfigure}
  \adjustbox{max width=0.94\linewidth,max totalheight=0.46\textheight}{\includegraphics{figures/fig2-5_spectrum.png}}\\[3pt]
  \figcaption{그림 2-5  격리 스펙트럼. 오른쪽으로 갈수록 테넌트와 호스트 커널 사이의 계층이 많아지고 격리는 강해지지만 시작은 느려지고 밀도는 낮아진다. 서버리스의 최적 지점은 microVM이다.}
\end{tbbfigure}

\subsection*{컨테이너만으로 해결되지 않는 이유}

다음 주에는 커널 기본 요소로 컨테이너를 만들 것이므로 여기서는 보안에 관련된 윤곽만 보면 된다. 컨테이너는 일반 프로세스 묶음에 비공개 \textit{뷰}(네임스페이스)와 제한된 \textit{몫}(cgroup)을 제공한다. 그러나 \textbf{호스트의 다른 모든 사용자와 여전히 같은 커널을 공유한다.} 시작에는 \textasciitilde{}100 ms가 걸리고 오버헤드는 거의 0이므로 컨테이너가 패키징 전쟁에서 승리했다. 하지만 격리 경계는 시스템 호출 인터페이스다. 수백 개 진입점이 하나의 커널로 이어지며, 커널 취약점 하나가 모든 테넌트의 취약점이 된다. 컨테이너 탈출은 가설이 아니라 반복해서 등장하는 CVE 유형이다. \textit{자신의} 클러스터에서 \textit{자신의} 워크로드를 공유할 때는 보통 이 위험을 받아들일 수 있으며, Kubernetes도 바로 그 판단 위에 세워졌다. 그러나 공개 FaaS, CI 서비스, 모델이 생성한 프로그램을 실행하는 AI 에이전트 샌드박스처럼 \textbf{임의의 낯선 코드}를 실행할 때는 대부분 운영자가 공유 커널로 부족하다고 판단한다. (마지막 예시는 이 과정의 시대에 특히 중요하다. "LLM이 방금 작성한 코드를 실행한다"는 것은 멀티테넌트 격리 문제의 전형이다.)

\subsection*{컨테이너와 VM 사이의 세 가지 답}

\textbf{gVisor (Google, 2018)}는 컨테이너 형태를 유지하되 경호원을 넣는다. Go로 작성한 사용자 공간 커널인 \textit{Sentry}가 애플리케이션의 시스템 호출을 가로채 직접 처리하며, 대폭 줄이고 강화한 호출 집합으로 실제 호스트 커널에 접근한다. 따라서 테넌트 코드는 호스트 커널과 직접 통신하지 않는다. 커널 익스플로잇은 먼저 메모리 안전성을 갖춘 두 번째 커널을 뚫어야 한다. 대가는 시스템 호출 비용(I/O가 많은 워크로드가 체감한다)과 Linux 동작의 경계에서 발생하는 불완전한 호환성이다. Google은 App Engine, Cloud Run, Cloud Functions의 대규모 서버군에서 이를 사용한다.

\textbf{Kata Containers}는 반대 방향으로 간다. 컨테이너 \textit{인터페이스}는 유지하지만 \textit{구현}을 실제 하드웨어 VM으로 바꾼다. 각 pod는 경량 VM 안에 자체 게스트 커널을 가지며, 표준 런타임 배관(제4장)을 통해 Kubernetes에 들어간다. 컨테이너 형태의 워크플로에 VM급 격리를 얻는 대신 pod마다 VM 형태의 오버헤드를 치른다. Kata 아래에 최소 VMM을 결합하면 이를 줄일 수 있으며, 이 절의 주인공으로 이어진다.

\textbf{Firecracker (AWS, 2018년 오픈 소스 공개, NSDI 2020)}는 잔인할 만큼 구체적인 질문에 답한다. Lambda가 실제로 필요로 하는 \textit{최소} 가상 머신은 무엇인가? 답은 KVM과 하드웨어 격리를 유지하되 범용 장치 박물관을 없앤 VMM이다(그림 2-6). BIOS 없이 게스트 커널을 직접 로드한다. PCI 열거도, VGA·USB·사운드도, 라이브 마이그레이션도 없다. virtio-net, virtio-block, vsock 채널, 직렬 콘솔, 이를 구동하는 REST API만 있다. QEMU의 수백만 줄 C 코드에 비해 Rust 코드 약 5만 줄이며, 각 microVM은 seccomp로 제한된 "jailer"가 한 번 더 가둔다. 공개된 수치는 VM 비용에 대한 산업의 직관을 다시 배선했다. \textbf{125 ms 안에 게스트 사용자 공간까지 부팅, microVM당 5 MB 미만의 VMM 오버헤드, 호스트당 초당 microVM 150개 생성}이다. 운영 관점에서 프로세스처럼 동작하는 VM이다. 실제로 모든 AWS Lambda 호출과 Fargate 작업이 이 위에서 초당 수백만 개 규모로 실행된다.

\begin{tbbfigure}
  \adjustbox{max width=0.94\linewidth,max totalheight=0.46\textheight}{\includegraphics{figures/fig2-6_firecracker.png}}\\[3pt]
  \figcaption{그림 2-6  Firecracker: KVM은 유지하고 장치 박물관은 없앤다. 125 ms, 5 MB, 150/s라는 수치는 이번 주 필독 자료인 NSDI 2020 논문에서 가져왔다.}
\end{tbbfigure}

\begin{tbbtablewrap}
  \begin{tabular}{@{}>{\raggedright\arraybackslash}m{\dimexpr 0.2021\tbltextwidth-2\tabcolsep\relax} >{\raggedright\arraybackslash}m{\dimexpr 0.1968\tbltextwidth-2\tabcolsep\relax} >{\raggedright\arraybackslash}m{\dimexpr 0.1968\tbltextwidth-2\tabcolsep\relax} >{\raggedright\arraybackslash}m{\dimexpr 0.2021\tbltextwidth-2\tabcolsep\relax} >{\raggedright\arraybackslash}m{\dimexpr 0.2021\tbltextwidth-2\tabcolsep\relax}@{}}
  \arrayrulecolor{tbbnavy}\specialrule{1pt}{0pt}{0pt}
  \rowcolor{tbbtablehead}{\centering \bfseries\color{tbbnavy} } & {\centering \bfseries\color{tbbnavy}컨테이너} & {\centering \bfseries\color{tbbnavy}gVisor} & {\centering \bfseries\color{tbbnavy}microVM (FC/Kata)} & {\centering \bfseries\color{tbbnavy}완전한 VM} \\
  \arrayrulecolor{tbbnavy}\specialrule{0.7pt}{0pt}{2pt}
  {\textbf{커널}} & {공유하는 호스트 커널} & {사용자 공간 프록시 커널} & {자체 게스트 커널} & {자체 게스트 커널} \\
  \arrayrulecolor{tbbtableborder}\hline
  {\textbf{격리 경계}} & {시스템 호출 인터페이스} & {Sentry + 축소한 시스템 호출} & {하드웨어(VT-x)} & {하드웨어(VT-x)} \\
  \arrayrulecolor{tbbtableborder}\hline
  {\textbf{시작 시간}} & {\textasciitilde{}100 ms} & {\textasciitilde{}100 ms} & {\textasciitilde{}125 ms–1 s} & {초–분} \\
  \arrayrulecolor{tbbtableborder}\hline
  {\textbf{인스턴스당 오버헤드}} & {≈ 0} & {시스템 호출 비용} & {\textasciitilde{}5 MB + 게스트 커널} & {수백 MB} \\
  \arrayrulecolor{tbbtableborder}\hline
  {\textbf{일반적인 용도}} & {자체 워크로드, K8s 기본값} & {멀티테넌트 PaaS (Google)} & {FaaS, 신뢰할 수 없는 코드(AWS)} & {임대한 클라우드 서버} \\
  \arrayrulecolor{tbbnavy}\specialrule{1pt}{2pt}{0pt}
  \end{tabular}\par
  \figcaption{표 2-2  격리 선택지 비교(벤치마크가 아닌 규모 차수)}
\end{tbbtablewrap}

\subsection*{서빙에 주는 의미(12주 차로 이어지는 실마리)}

이를 제1장의 0으로 스케일링 이야기와 연결하면 그림이 유용하게 선명해진다. Firecracker는 \textit{환경}이 \textasciitilde{}125 ms 안에 나타날 수 있음을 입증한다. 따라서 GPU 기반 모델 엔드포인트의 콜드 스타트가 30–60 s라면 microVM이 범인이 아니다. GPU 호스트를 스케줄링하고, 수 기가바이트 컨테이너 이미지를 가져오며, 수십 기가바이트 가중치를 GPU 메모리에 로드하고, 런타임 캐시를 워밍업하는 데 시간이 쓰였다. 이 분해 자체가 12주 차 수업의 씨앗이다. 가상화는 제 몫을 다했고, 남은 콜드 스타트 엔지니어링은 이미지, 가중치, GPU 상태에 관한 것이다. 그래서 \textit{로드된} VM의 스냅숏/복원 같은 기법(Firecracker 스냅숏은 현재 서버리스 GPU 연구와 실무의 한 흐름이다)이 최전선에 있다. 이제 서버리스 GPU 플랫폼이 콜드 스타트 수치를 제시하면 그 수치를 검증할 방법을 안다.

\begin{calloutbox}{tbbpurple}{tbbpurplefill}
{\normalsize\bfseries\color{tbbpurple}생각해 보기}\par\vspace{3pt}
\begin{tbbitemize}
  \item 팀이 AI 에이전트 제품에서 모델이 생성한 신뢰할 수 없는 Python을 실행한다고 하자. 커널 0-day의 피해 범위, 시작 지연 시간 예산(대화형이다!), 호스트당 밀도, GPU 접근(2.6절을 미리 보자)을 기준으로 이 샌드박스에 컨테이너, gVisor, Firecracker 중 무엇이 적합한지 논증해 보자.
  \item 밀도를 계산해 보자. 호스트에 테넌트용 RAM 384 GB가 있다. 유휴 상태인 128 MB 힙을 쓰는 함수가 (a) 컨테이너(\textasciitilde{}0 오버헤드), (b) Firecracker microVM(각각 +5 MB + \textasciitilde{}20 MB 게스트 커널), (c) 완전한 VM(각각 +300 MB)에 몇 개씩 들어가는지 추정하라. 그 결과는 호출당 가격에 어떤 영향을 주는가?
\end{tbbitemize}
\end{calloutbox}

\section{클라우드의 GPU: 인스턴스, 패스스루, 공유}

지금까지 이 장에서는 CPU, 메모리, I/O를 가상화했다. 이제 이 과정의 실제 중심인 구성 요소를 더해 보자. GPU는 앞의 모든 방법론에 완강히 저항한다. "GPU용 VT-x"는 없다. 아키텍처의 non-root 모드도, CUDA용 하드웨어 트랩 후 에뮬레이션도 없다. GPU는 마이크로초 만에 세계 전환할 수 없는 수 기가바이트의 컨텍스트를 보관하고, 거대한 독점 드라이버 스택으로 구동되며, 정중한 시분할이 아니라 까다로운 소유자 한 명을 위해 설계되었다. 따라서 "테넌트가 GPU를 어떻게 얻는가?"에 대한 클라우드의 답은 앞의 모든 것과 다르다. 세 가지 답(그림 2-7)은 12주 차까지 따라다닌다. \textbf{GPU를 얼마나 세밀하게 공유할 수 있는가는 가상화 의상을 입은 서빙 경제성 문제}이기 때문이다.

\subsection*{먼저 가격표의 형태: GPU 인스턴스 계열}

공급자는 제1장의 학습과 추론 구분을 반영한 인스턴스 계열로 GPU를 패키징한다. \textbf{학습형 인스턴스}는 플래그십 가속기(H100/H200/B200급) 8개를 한 섀시에 넣고 내부는 NVLink, 섀시 사이는 매우 빠른 패브릭으로 연결한다. AWS P 계열, GCP A3, Azure ND 계열이 있다. \textbf{추론형 인스턴스}는 적당한 CPU 수에 소형 GPU(T4, L4, A10급) 하나에서 몇 개를 연결한다. AWS G 계열, GCP G2가 있으며 시간당 비용이 약 10배 낮다. 이 과정의 실습도 이 범위에서 진행한다(AWS g4dn 같은 T4 인스턴스나 GCP N1+T4 조합). 초보자를 놀라게 하는 실용적인 주의 사항이 두 가지 있다. GPU 인스턴스는 흔히 \textbf{할당량으로 잠겨} 있어 시작 전에 접근을 요청해야 하므로 실습 초반에 신청해야 한다. 또한 제1장의 메모리 적합성 계산에 따라 이 과정에서는 GPU 사양표보다 \textbf{VRAM}이 더 중요하다.

\begin{tbbtablewrap}
  \begin{tabular}{@{}>{\raggedright\arraybackslash}m{\dimexpr 0.1862\tbltextwidth-2\tabcolsep\relax} >{\raggedright\arraybackslash}m{\dimexpr 0.2660\tbltextwidth-2\tabcolsep\relax} >{\raggedright\arraybackslash}m{\dimexpr 0.1915\tbltextwidth-2\tabcolsep\relax} >{\raggedright\arraybackslash}m{\dimexpr 0.3564\tbltextwidth-2\tabcolsep\relax}@{}}
  \arrayrulecolor{tbbnavy}\specialrule{1pt}{0pt}{0pt}
  \rowcolor{tbbtablehead}{\centering \bfseries\color{tbbnavy}GPU (VRAM)} & {\centering \bfseries\color{tbbnavy}대표 인스턴스} & {\centering \bfseries\color{tbbnavy}대략적인 \$/hr} & {\centering \bfseries\color{tbbnavy}서빙 최적 용도} \\
  \arrayrulecolor{tbbnavy}\specialrule{0.7pt}{0pt}{2pt}
  {\textbf{T4 (16 GB)}} & {AWS g4dn.xlarge · GCP N1+T4} & {≈ \$0.5} & {이 과정의 실습 기준, 소형 모델, INT8 추론} \\
  \arrayrulecolor{tbbtableborder}\hline
  {\textbf{L4 (24 GB)}} & {GCP g2-standard-4 · AWS g6} & {≈ \$0.7–1.0} & {현대 추론의 주력, 토큰당 낮은 전력} \\
  \arrayrulecolor{tbbtableborder}\hline
  {\textbf{A10G (24 GB)}} & {AWS g5.xlarge} & {≈ \$1.0} & {중형 모델 — 제1장 사례 1-B에서 사용} \\
  \arrayrulecolor{tbbtableborder}\hline
  {\textbf{A100 (80 GB)}} & {AWS p4d (÷8) · GCP a2} & {GPU당 ≈ \$3–5} & {대형 모델(양자화 70B), 분할용 MIG 호스트(아래 절)} \\
  \arrayrulecolor{tbbtableborder}\hline
  {\textbf{H100 (80 GB)}} & {AWS p5 (÷8) · GCP a3} & {GPU당 \$2–7(채널에 따라 다름)} & {프런티어 서빙 및 학습, NVLink 연결} \\
  \arrayrulecolor{tbbtableborder}\hline
  {\textcolor{tbbgray}{(참고) B200 (192 GB)}} & {\textcolor{tbbgray}{차세대}} & {\textcolor{tbbgray}{\$4–6+}} & {\textcolor{tbbgray}{GPU당 상한을 다시 높임}} \\
  \arrayrulecolor{tbbnavy}\specialrule{1pt}{2pt}{0pt}
  \end{tabular}\par
  \figcaption{표 2-3  서빙 서버군의 가격(2026년 중반의 온디맨드 예시 — 숫자보다 구조가 오래간다)}
\end{tbbtablewrap}

제1장의 두 관점으로 이 표를 읽으면 목록이 아니라 의사 결정 도구가 된다. \textbf{결정적인 제약 차원:} 양자화 모델에 11 GB가 필요하다면 T4와 그 위의 모든 GPU가 \textit{작동}하지만, T4에서 L4 사이만 \textit{경제적}이다. T4 경제성으로 모델을 검증하고 "안전을 위해" A100에 배포하면 제약이 되지 않는 VRAM과 FLOPs에 시간당 6–10×를 지불한다. \textbf{사용률:} 서빙하든 놀든 가격 열의 요금이 시간당 부과되므로, 서버군에 관한 질문은 "어느 GPU가 가장 좋은가?"가 아니라 "실제로 유지할 수 있는 사용률에서 SLO를 충족하는 가장 저렴한 카드는 무엇인가?"다. 12주 차에는 오토스케일링과 GPU 공유가 이 표를 최적화 문제로 바꾼다. \textit{서버군 표}로 기억하자.

\subsection*{VM이 GPU 전체를 얻는 방법: PCIe 패스스루}

\begin{terminalbox}
\begin{Verbatim}[breaklines=true,breakanywhere=true,fontsize=\small,formatcom=\color{tbbtermfg},codes=\tbbcodeactive]
# on this week's rented VM - recall Exhibit 2-A: everything so far was fake
$ lspci | grep -i nvidia
00:1e.0 3D controller: NVIDIA Corporation TU104GL [Tesla T4]
$ nvidia-smi --query-gpu=name,driver_version,memory.total --format=csv
name, driver_version, memory.total [MiB]
Tesla T4, 550.54.15, 16384 MiB
# a REAL PCI device, running the REAL driver, inside the fake machine.
# the one component the hypervisor did not fake. how? and why not?
\end{Verbatim}
\end{terminalbox}
\figcaption{기록 2-4  환상 속의 예외. 디스크는 virtio, 시계는 kvm-clock, CPU 모델은 메뉴에서 고른 것이지만 T4는 패스스루된 물리 실리콘이다. 이 하위 절에서는 그 방법(VFIO + IOMMU)과 대가(카드당 테넌트 하나)를 설명한다.}

이번 주 g4dn을 시작했을 때 \code{nvidia-smi}에 보이는 T4는 에뮬레이션이나 프록시가 아니라 \textbf{물리 카드}다. 메커니즘은 \textbf{PCIe 패스스루}다. 호스트가 Linux VFIO 프레임워크를 통해 장치를 자신에게서 분리해 게스트에 매핑한다. 게스트는 실제 NVIDIA 드라이버를 로드하고 실제 실리콘과 통신한다. 이를 가능하게 하는 하드웨어는 표 2-1의 \textbf{IOMMU}(Intel VT-d / AMD-Vi)다. GPU는 메모리를 직접 읽고 쓰는 DMA 엔진이므로, 통제되지 않은 DMA 엔진을 테넌트에 넘기면 호스트 전체를 읽을 수 있다. IOMMU는 장치에 자체 주소 변환을 제공한다. 게스트 메모리 접근이 EPT를 통과하듯 DMA 주소도 하이퍼바이저가 제어하는 테이블을 통과하므로, GPU는 자기 게스트의 페이지만 건드릴 수 있다. 패스스루 덕분에 클라우드 GPU는 네이티브에 가까운 성능을 낸다. 설정 후 데이터 경로에는 하이퍼바이저가 전혀 관여하지 않는다. 한계도 마찬가지로 구조적이다. \textbf{할당 단위가 물리 GPU 하나 전체}다. 테넌트 하나, 카드 하나이므로 중재할 것이 없다. (VM이 호스트에 고정되기도 한다. 패스스루 장치를 사용하면 라이브 마이그레이션할 수 없어서 GPU 인스턴스가 유지보수 때 이동하지 않고 중단되는 이유 중 하나다.)

\subsection*{GPU 하나 공유하기: 시간 분할, vGPU, MPS, MIG}

모델이 작아지는 순간 GPU 전체 할당은 제1장의 사용률 문제와 충돌한다. 80 GB H100에서 3 GB 모델을 실행하면 랙에서 가장 비싼 부품의 96\%를 낭비한다. 공유 선택지는 \textit{무엇이 공정성을 강제하는지}에 따라 나뉜다. 소프트웨어 스케줄인지 실리콘 파티션인지가 차이다(그림 2-7).

\begin{tbbfigure}
  \adjustbox{max width=0.94\linewidth,max totalheight=0.46\textheight}{\includegraphics{figures/fig2-7_gpu_virt.png}}\\[3pt]
  \figcaption{그림 2-7  GPU에 접근하는 세 가지 방식. 인스턴스가 판매하는 독점 패스스루, 시간적 공유(vGPU/시간 분할/MPS), 공간적 하드웨어 파티셔닝(MIG)이다.}
\end{tbbfigure}

\begin{tbbitemize}
  \item \textbf{단순 시간 분할}(Kubernetes GPU 시간 분할 등)은 여러 워크로드가 번갈아 같은 GPU에 작업을 제출하게 한다. 메모리를 나누는 것도, 차례의 공정성을 보장하는 것도 없다. 테넌트 하나가 VRAM 전체를 할당하거나 긴 커널을 독점할 수 있다. 신뢰하고 협력하는 작업 사이에서는 쓸 만하지만, 낯선 사용자나 꼬리 지연 시간 약속에는 적합하지 않다.
  \item \textbf{NVIDIA vGPU}는 VM의 시간적 공유를 제품화한다. 라이선스가 필요한 호스트 소프트웨어와 \textit{VRAM을 정적으로 나누는} 프로필을 사용해 각 vGPU가 자기 조각을 소유하며, 스케줄러가 거의 일반적인 드라이버를 실행하는 게스트 사이에서 계산 시간을 분할한다. 가상 데스크톱용으로 탄생했지만 추론에도 쓸 수 있다. 계산 QoS는 여전히 하드웨어급이 아니라 스케줄러급이다.
  \item \textbf{MPS (Multi-Process Service)}는 한 소유자의 프로세스 \textit{내에서} CUDA 수준으로 공유한다. 프로세스 여러 개의 커널이 GPU 하나에서 동시에 실행되며 선택적으로 자원 한도를 둘 수 있다. 자체 모델을 함께 배치하기에는 훌륭하지만 보안 경계는 전혀 제공하지 않는다.
  \item \textbf{MIG (Multi-Instance GPU}, Ampere/A100 이후\textbf{)}는 마침내 CPU가 2005년에 한 일을 수행한다. 격리를 실리콘 안으로 옮긴다. 다이는 최대 \textbf{7개 인스턴스}로 분할되며 각 인스턴스가 자체 SM 조각, L2 조각, DRAM 컨트롤러/용량을 소유한다. 예를 들어 H100-80GB 한 대에서 1g.10gb 인스턴스 7개를 만들거나 3g.40gb + 2g.20gb + 2g.20gb처럼 크기를 섞을 수 있다. 각 인스턴스에는 하드웨어가 강제하는 메모리 대역폭과 용량이 있으므로 이웃이 난동을 부려도 지연 시간이 움직이지 않는다. 서빙 서버군이 갈망하는 \textbf{조각별 예측 가능한 p99}를 제공한다. 대가는 고정된 프로필 메뉴, 하나의 작업에 협력하지 못하는 인스턴스, 중단을 일으키는 재구성이다. 보장을 얻기 위해 경직성을 감수하는 절충이며 12주 차에 가격을 계산한다.
\end{tbbitemize}

\begin{tbbtablewrap}
  \begin{tabular}{@{}>{\raggedright\arraybackslash}m{\dimexpr 0.2181\tbltextwidth-2\tabcolsep\relax} >{\raggedright\arraybackslash}m{\dimexpr 0.1915\tbltextwidth-2\tabcolsep\relax} >{\raggedright\arraybackslash}m{\dimexpr 0.1862\tbltextwidth-2\tabcolsep\relax} >{\raggedright\arraybackslash}m{\dimexpr 0.1809\tbltextwidth-2\tabcolsep\relax} >{\raggedright\arraybackslash}m{\dimexpr 0.2234\tbltextwidth-2\tabcolsep\relax}@{}}
  \arrayrulecolor{tbbnavy}\specialrule{1pt}{0pt}{0pt}
  \rowcolor{tbbtablehead}{\centering \bfseries\color{tbbnavy} } & {\centering \bfseries\color{tbbnavy}시간 분할(K8s)} & {\centering \bfseries\color{tbbnavy}NVIDIA vGPU} & {\centering \bfseries\color{tbbnavy}MPS} & {\centering \bfseries\color{tbbnavy}MIG} \\
  \arrayrulecolor{tbbnavy}\specialrule{0.7pt}{0pt}{2pt}
  {\textbf{공유 유형}} & {시간적} & {시간적 + 정적 VRAM} & {동시 실행(CUDA)} & {공간적(하드웨어)} \\
  \arrayrulecolor{tbbtableborder}\hline
  {\textbf{메모리 격리}} & {없음} & {정적 파티션} & {없음} & {하드웨어 강제} \\
  \arrayrulecolor{tbbtableborder}\hline
  {\textbf{QoS / p99}} & {없음} & {스케줄러급} & {한도, 경계 없음} & {하드웨어급} \\
  \arrayrulecolor{tbbtableborder}\hline
  {\textbf{보안 경계}} & {없음} & {VM급(라이선스)} & {없음} & {있음(인스턴스별)} \\
  \arrayrulecolor{tbbtableborder}\hline
  {\textbf{최적 용도}} & {신뢰하는 배치 작업} & {VDI, 혼합 VM} & {함께 배치한 자체 모델} & {멀티테넌트 추론} \\
  \arrayrulecolor{tbbnavy}\specialrule{1pt}{2pt}{0pt}
  \end{tabular}\par
  \figcaption{표 2-4  GPU 공유 선택지 한눈에 보기}
\end{tbbtablewrap}

마지막 연결점을 보자. GPU는 이 장에서 CPU의 역사 전체를 되풀이했다. 가상화 이전 서버 같은 독점 소유, 바이너리 변환 시대처럼 보장이 약한 소프트웨어 시분할, 그리고 하드웨어가 강제하는 파티션으로 이어졌다. MIG는 GPU의 VT-x 순간이다. 아키텍처에는 운율이 있다.

\begin{calloutbox}{tbbpurple}{tbbpurplefill}
{\normalsize\bfseries\color{tbbpurple}생각해 보기}\par\vspace{3pt}
\begin{tbbitemize}
  \item 내부 추론 서비스 세 개가 각각 \textasciitilde{}15 GB VRAM과 p99 SLO를 필요로 한다. H100-80GB 한 대가 있다. MIG 프로필, MPS, 별도의 L4 인스턴스 세 대를 격리, 사용률, 비용, 장애 피해 범위 관점에서 비교해 논증하라.
  \item 패스스루는 라이브 마이그레이션을 막지만 virtio 장치는 막지 않는 이유는 무엇인가? (각 경우에 상태가 어디에 존재하는지 생각해 보자.)
\end{tbbitemize}
\end{calloutbox}

\section{실습: 첫 클라우드 GPU를 안전하게 사용하기}

실습은 이 장을 구체화한다. 클라우드 계정을 만들고 \textbf{무엇보다 먼저 과금 안전장치를 설정}한 뒤 GPU VM을 시작한다. \code{nvidia-smi}로 하드웨어를 만나고, 같은 추론 작업에서 CPU와 GPU의 속도를 겨루며, 계량기를 켜 둔 채 남기지 않고 모두 종료한다. 클릭 단위의 자세한 절차는 실습 안내서에 있다. 콘솔은 교재보다 빨리 바뀌기 때문이다. 여기서는 각 단계의 \textit{이유}와 예상할 수치를 다룬다. 산출물은 측정한 지연 시간과 이 장에 연결한 관찰 세 가지를 담은 1쪽 보고서다.

\subsection*{0단계 — 계산 자원보다 먼저 비용 안전장치}

제1장은 "시간당 과금 = 유휴 시간은 돈"이라는 말로 끝났다. 여기서 이를 몸에 익힌다. 무엇이든 시작하기 전에 \textbf{(a)} 공급자 과금 콘솔에서 알림 임계값이 있는 예산(예: \$10의 50/80/100\%에서 이메일)을 만들고, \textbf{(b)} 교육용 크레딧을 사용 중이며 계량값이 어디에 표시되는지 확인하고, \textbf{(c)} 두 시계를 이해해야 한다. \textit{중지한} 인스턴스는 계산 요금이 멈추지만 \textbf{연결된 디스크 스토리지 요금은 계속 부과}된다. \textit{종료한} 인스턴스는 모든 자원과 별도 지시가 없다면 디스크도 해제한다. 초보자의 전형적인 청구서는 GPU 시간 때문이 아니다. 오래 중지한 인스턴스에 연결한 채 잊은 200 GB 볼륨이나, 버려진 브라우저 탭에서만 "중지"한 인스턴스 때문이다. 실습 체크리스트의 마지막은 콘솔에 \textbf{실행 중인 인스턴스 0개, 연결되지 않은 볼륨 0개}가 표시되는지 확인하는 것이다. 보고서용으로 화면을 캡처한다.

\subsection*{1\textasciitilde{}3단계 — 할당량, 시작, GPU 만나기}

GPU 할당량은 일찍 요청한다. 승인에 몇 시간이 걸릴 수 있다. 가까운 리전에서 가장 저렴한 T4 인스턴스 유형(AWS \code{g4dn.xlarge} 또는 GCP N1 + T4급, 시간당 약 1달러)를 고른다. 측정 자체에는 계정 없이 쓸 수 있는 대안으로 Colab이 있다. 드라이버가 사전 설치된 딥러닝 OS 이미지를 선택하고 ssh로 접속한다. 그런 다음 시스템 관점에서 \code{nvidia-smi}를 읽어 보자. 모든 필드가 이 장의 내용을 되풀이한다. \textbf{드라이버/CUDA 버전}은 \textit{게스트}에 속한다. T4가 \textbf{패스스루}되었으므로 VM이 실제 드라이버를 실행한다(§2.6). \textbf{메모리 표시값(0/16384 MiB)}은 제1장의 계산과 10주 차의 양자화가 관리할 VRAM 예산이다. \textbf{사용률과 프로세스 목록}은 12주 차 오토스케일러와 14주 차 대시보드가 소비할 메트릭이다. \code{lspci | grep -i nvidia}도 실행한다. GPU가 물리 카드처럼 게스트 PCI 버스에 있으며, 패스스루가 눈에 보인다. 40초 동안의 시작이 기계적으로 무엇이었는지도 주목하자. 컨트롤 플레인이 호스트를 골랐고, 하이퍼바이저(이 장)가 게스트를 구체화했으며, 펌웨어가 없거나 빠른 부팅이 커널을 올렸고, cloud-init이 개인화했다. 이제 콜드 스타트를 직접 눈으로 보았다. 12주 차에는 이 관찰을 산업 규모로 확장한다.

\subsection*{4단계 — 측정: CPU 추론 대 GPU 추론}

안내서의 벤치마크 스크립트를 실행한다. 소형 트랜스포머(DistilBERT급)나 소형 CNN에 같은 입력 배치를 넣고 CPU와 GPU에서 시간을 잰다. PyTorch에서는 \code{.to("cuda")} 하나가 차이다. 코드보다 프로토콜이 중요하다. \textbf{먼저 워밍업하라.} 첫 GPU 호출에는 CUDA 컨텍스트, 메모리 풀, 커널 선택 같은 일회성 비용이 들어 CPU보다 \textit{느릴 수도 있다}. 개인적으로 얻은 첫 콜드 스타트 데이터로 따로 보고하고, 이후 약 100회의 정상 상태 시간을 측정한다. 배치 크기(1, 8, 32, 64)를 바꾸며 \textbf{배치당 지연 시간과 처리량}을 모두 기록한다. 결과의 대략적인 규모는 다음과 같을 것이다. CPU는 배치 크기 1에서 수십 ms이며 배치에 따라 거의 선형으로 증가한다. GPU는 배치 크기 1에서 수 ms이고, 배치가 커져도 \textit{거의 움직이지 않다가} 배치 크기 32–64에서 처리량이 CPU의 10–50×가 된다. GPU의 배치 크기 1 성능이 기대에 못 미친다면 이유를 다시 발견한 것이다. 작은 작업에서는 시작 오버헤드와 PCIe 홉이 지배하며, 병렬 하드웨어는 폭넓게 일을 줄 때 값을 한다. 제1장의 처리량과 지연 시간 사이 긴장, 그리고 10주 차 배칭의 전체 동기를 직접 측정한 셈이다. 보고서에 바로 이런 관찰을 적어야 한다.

\begin{calloutbox}{tbbpurple}{tbbpurplefill}
{\normalsize\bfseries\color{tbbpurple}실습 체크리스트(크레딧을 지키는 버전)}\par\vspace{3pt}
\begin{tbbitemize}
  \item 첫 시작 \textbf{전에} 예산 + 알림이 존재하며, 과금 페이지 위치를 안다.
  \item GPU 할당량을 일찍 요청하고, 가장 저렴한 T4 인스턴스 유형, 가까운 리전, 드라이버가 포함된 딥러닝 이미지를 선택한다.
  \item \code{nvidia-smi}와 \code{lspci} 출력을 저장하고, 워밍업과 정상 상태를 따로 측정하며, 배치 변화(1/8/32/64)를 기록한다.
  \item 이유가 없다면 인스턴스를 단순히 중지하지 않고 \textbf{종료}한다. 볼륨을 확인하고, 실행 중이거나 고아가 된 자원이 없는 최종 콘솔 화면을 캡처한다.
  \item 보고서에는 수치와 §2.3–§2.6에 연결한 관찰 세 가지를 담는다(예: 스틸 타임을 보았는가? 워밍업 비용은? 배치 효과는?).
\end{tbbitemize}
\end{calloutbox}

\vspace{3.0pt}

\begin{calloutbox}{tbbpurple}{tbbpurplefill}
{\normalsize\bfseries\color{tbbpurple}생각해 보기}\par\vspace{3pt}
\begin{tbbitemize}
  \item 배치 크기 1의 GPU 지연 시간은 CPU보다 거의 빠르지 않지만, 배치 크기 64의 GPU 처리량은 CPU의 40×다. 프로덕션 트래픽은 한 번에 요청 하나씩 도착한다. 어떤 선택지가 있는가? (방금 10\textasciitilde{}12주 차의 의제를 다시 발명했다.)
  \item g4dn.xlarge가 \$0.53/hr일 때 효율적으로 2시간 작업하면 실습 전체 비용이 얼마인지 추정해 보자. 인스턴스를 한 달 동안 잊으면 얼마인가? 안전장치는 어느 숫자를 막으려 하는가?
\end{tbbitemize}
\end{calloutbox}

\namedsection{요약}{열두 가지로 정리한 제2장}

\qitem{1}{\textbf{임대한 머신은 존재하지 않으며, 스스로 그렇게 말한다.} 하이퍼바이저(VMM)는 물리 호스트 하나를 격리된 가상 머신 여러 개로 다중화한다. 게스트는 소매업체 이름이 들어간 DMI 문자열, \code{kvm}이라고 답하는 \code{systemd-detect-virt}, \code{kvm-clock}으로 도착하는 시계에서 직접 자백을 읽을 수 있다(사례 2-A). 이 소프트웨어 환상이 시간당 과금, 멀티테넌시, 탄력성의 기반이다.}

\qitem{2}{\textbf{가상화는 "비싼 하드웨어가 놀고 있다"는 질문에 대한 60년 된 답이다.} 1막: 수백만 달러짜리 메인프레임 → 시분할과 CP-67의 사용자별 가상 머신(1967). 2막: 5–15\% 사용률로 난립한 2000년대 서버 → VMware 통합. 3막: EC2의 \textit{시간당 \$0.10}(2006). 할당이 소프트웨어가 되면서 계산 자원에 계량기가 생겼다. GPU는 다시 1막을 살고 있으며, §2.6은 진행 중인 GPU의 2막과 3막이다.}

\qitem{3}{\textbf{Popek \& Goldberg (1974)}는 동등성, 효율성, 자원 제어라는 계약과 정리를 정의했다. 모든 민감 명령어가 특권 명령어, 즉 트랩하는 명령어일 때만 트랩 후 에뮬레이션이 작동한다. 고전적 x86에는 민감하지만 특권은 아닌 명령어 17개가 있어 이 조건을 깨뜨렸다.}

\qitem{4}{두 소프트웨어 시대의 기법이 간극을 메웠다. \textbf{바이너리 변환}(VMware)은 게스트를 실행 중에 다시 써서 수정하지 않은 게스트를 완전 가상화했고, \textbf{반가상화}(Xen)는 게스트를 수정해 하이퍼콜을 사용함으로써 소스가 필요하지만 더 빠르게 실행했다. 협력하는 게스트라는 아이디어는 오늘날 virtio로 살아 있다.}

\qitem{5}{\textbf{VT-x/AMD-V (2005–06)}는 VMCS를 계약으로 삼아 VMX root/non-root 모드를 추가했다. 게스트 커널은 다시 ring 0에서 실행되고 위험한 작업은 하드웨어 VM exit를 일으킨다. \textbf{Exit는 가상화 성능의 통화}이므로 횟수를 세고 비용을 낮춰야 한다.}

\qitem{6}{\textbf{EPT/NPT}는 섀도 페이지 테이블을 끝냈다. 하드웨어가 두 변환(guest VA → guest PA → host PA)을 순회한다. 비용은 더 무거운 TLB 미스로 드러나므로 대형 페이지가 중요하며, 메모리를 많이 쓰는 추론에서는 특히 그렇다.}

\qitem{7}{\textbf{KVM + QEMU + virtio}는 하이퍼바이저가 된 Linux다. KVM은 커널에서 VT-x를 구동하고, QEMU는 VM별 장치 모델을 제공하며, virtio는 장치의 환상을 exit마다 작업을 배치하는 정직한 공유 메모리 virtqueue로 대체한다. \textbf{vCPU는 그저 호스트 스레드}이므로 오버커밋과 스틸 타임, 그리고 꼬리 지연 시간에 미치는 영향이 따라온다.}

\qitem{8}{\textbf{격리 스펙트럼}은 테넌트 코드와 호스트 커널 사이에 무엇이 있는지에 따라 선택지를 나열한다. 프로세스 < 컨테이너(공유 커널!) < gVisor(사용자 공간 커널) < microVM < 완전한 VM 순이며, 시작 시간과 밀도를 격리 강도와 교환한다.}

\qitem{9}{\textbf{Firecracker}는 서버리스 시대의 답이다. KVM은 유지하고 장치 박물관은 없앤다. <125 ms 부팅, <5 MB 오버헤드, 호스트당 150 microVMs/s로, 프로세스처럼 동작하며 Lambda와 Fargate를 실행하는 VM이다. 12주 차에 주는 교훈은 환경이 더는 콜드 스타트 병목이 아니며 이미지, 가중치, GPU 상태가 병목이라는 것이다.}

\qitem{10}{\textbf{GPU는 고전적 가상화에 저항하며} VT-x에 해당하는 기능이 없다. 따라서 클라우드는 카드 전체를 \textbf{PCIe 패스스루}한다(IOMMU/VFIO 사용, 네이티브에 가까운 성능, 공유 불가, 기록 2-4의 "가짜 머신 안의 진짜 카드"). 공유에는 \textbf{시간 분할/vGPU}(소프트웨어 공정성), \textbf{MPS}(함께 배치한 자체 워크로드), \textbf{MIG}(하드웨어 파티션, 최대 7개 격리 인스턴스와 조각별 메모리 및 예측 가능한 p99, GPU의 VT-x 순간)를 사용한다.}

\qitem{11}{\textbf{서버군 표(표 2-3)는 의사 결정 도구다.} 추론용 T4/L4/A10G의 \textasciitilde{}\$0.5–1.0/hr부터 GPU당 수 달러인 A100/H100까지 제시한다. "클수록 안전하다"는 목록으로 읽지 말고, 제1장의 결정적인 제약 차원(들어가면서 SLO를 충족하는 \textit{가장 저렴한} 카드는 무엇인가?)과 사용률 관점으로 읽어야 한다.}

\qitem{12}{\textbf{과금 안전장치는 엔지니어링 관행이다.} 처음 시작하기 전에 예산과 알림을 설정하고, 중지해도 무료가 아니며(디스크는 계속 과금된다), 종료하면 자원을 해제한다는 점을 이해한다. 실습은 빈 콘솔, 그리고 직접 측정한 첫 콜드 스타트와 CPU 대 GPU 배칭 곡선으로 끝난다.}

\namedsection{연습문제}{}

\subsection*{A. 객관식}

\qitem{1}{다음 중 Popek \& Goldberg가 제시한 가상 머신 모니터의 요구 사항이 아닌 것은?}

\begin{tbboptions}
  \item ① 동등성 — 게스트가 실제 하드웨어에서처럼 동작한다
  \item ② 효율성 — 대부분의 명령어가 CPU에서 직접 실행된다
  \item ③ 자원 제어 — VMM이 물리 자원을 완전히 통제한다
  \item ④ 이식성 — 어느 CPU 아키텍처에서든 게스트가 수정 없이 실행된다
\end{tbboptions}

\qitem{2}{VT-x에서 VMX non-root 모드로 실행되는 게스트 커널이 하이퍼바이저가 개입 대상으로 설정한 작업(예: I/O 포트 접근)을 실행하면 어떻게 되는가?}

\begin{tbboptions}
  \item ① 고전적 x86처럼 명령어가 조용히 아무 일도 하지 않는다
  \item ② 하드웨어가 VM exit를 수행한다. 게스트 상태를 VMCS에 저장하고 제어권을 root 모드의 하이퍼바이저로 넘긴다
  \item ③ 게스트 커널이 특권 폴트로 충돌한다
  \item ④ 명령어가 실행 중에 바이너리 변환된다
\end{tbboptions}

\qitem{3}{EPT/NPT가 해결하는 문제는?}

\begin{tbboptions}
  \item ① 게스트 물리 주소를 호스트 물리 주소로 하드웨어에서 변환하여 소프트웨어 섀도 페이지 테이블을 없앤다
  \item ② 하이퍼바이저로부터 게스트 메모리를 암호화한다
  \item ③ 게스트가 페이지 테이블을 공유하여 메모리를 절약하게 한다
  \item ④ 페이지 크기 전송으로 디스크 I/O를 가속한다
\end{tbboptions}

\qitem{4}{virtio 장치가 에뮬레이션 e1000 NIC보다 빠른 이유는?}

\begin{tbboptions}
  \item ① virtio 장치 코드가 더 빠른 언어로 작성되었다
  \item ② virtio 하드웨어가 e1000 하드웨어보다 물리적으로 새롭다
  \item ③ 게스트가 가상화 사실을 알고 협력한다. 공유 메모리 virtqueue를 통해 작업을 배치하므로 레지스터 접근마다 VM exit하는 대신 배치당 \textasciitilde{}한 번 exit한다
  \item ④ virtio는 모든 구성에서 커널을 완전히 우회한다
\end{tbboptions}

\qitem{5}{컨테이너와 (micro)VM 사이의 핵심 격리 차이는?}

\begin{tbboptions}
  \item ① 컨테이너는 CPU나 메모리 사용량을 제한할 수 없다
  \item ② 호스트의 컨테이너는 호스트 커널을 공유하므로 시스템 호출 인터페이스가 보안 경계다. 각 VM은 하드웨어 경계 뒤에 자체 커널을 둔다
  \item ③ VM에서는 Linux를 실행할 수 없다
  \item ④ 차이가 없으며 두 용어는 서로 바꿔 쓸 수 있다
\end{tbboptions}

\qitem{6}{테넌트별로 하드웨어가 강제하는 메모리와 대역폭 격리를 제공하는 GPU 공유 메커니즘은?}

\begin{tbboptions}
  \item ① Kubernetes 시간 분할        ② NVIDIA MPS
  \item ③ MIG (Multi-Instance GPU)      ④ CUDA 스트림
\end{tbboptions}

\qitem{7}{팀원이 말한다. "nvidia-smi에 실제 NVIDIA 드라이버를 사용하는 진짜 Tesla T4가 보이므로 EC2 인스턴스는 물리 서버임이 틀림없다." 올바른 반박은?}

\begin{tbboptions}
  \item ① 맞는 말이다. 실제 GPU가 있으면 실제 머신이다
  \item ② GPU는 디스크와 NIC처럼 QEMU가 에뮬레이션하므로 아무것도 증명하지 못한다
  \item ③ 머신은 KVM 게스트이며 dmidecode, systemd-detect-virt, kvm-clock이 모두 이를 증명한다. T4만 물리 장치이며 PCIe 패스스루로 통째로 매핑한, 환상에 대한 의도적 예외다
  \item ④ 가상 머신 안에서는 nvidia-smi를 전혀 실행할 수 없다
\end{tbboptions}

\subsection*{B. 단답형}

\qitem{1}{VT-x가 도입한 두 CPU 모드의 이름과 하이퍼바이저가 사용하는 모드를 쓰라.}

\qitem{2}{게스트 내부 모니터링에서 p99 지연 시간이 급증할 때 정확히 높은 \code{\%st}(스틸 타임)가 나타난다. 물리적으로 무슨 일이 일어나며 해결책 하나는 무엇인지 한 문장으로 설명하라.}

\qitem{3}{빈칸을 채워라. 물리 GPU를 게스트에 넘기는 것을 ( ㉠ )라 하며, 장치 DMA를 변환하고 통제하는 ( ㉡ ) 덕분에 안전하다.}

\qitem{4}{Firecracker는 microVM을 \textasciitilde{}125 ms에 부팅하지만 서버리스 GPU 모델 엔드포인트의 콜드 스타트에는 30초 이상이 걸릴 수 있다. 남은 시간이 쓰이는 곳 두 가지를 쓰라.}

\qitem{5}{양자화 모델에 VRAM 11 GB가 필요하고 L4(\textasciitilde{}\$0.8/hr)에서 SLO를 충족한다. 팀원이 "안전을 위해" A100-80GB(\textasciitilde{}\$4/hr)에 배포하자고 한다. 서버군 표가 이 제안에 적용하는 제1장의 두 관점을 각각 한 문장으로 설명하라.}

\subsection*{C. 서술형}

\qitem{1}{고전적 x86이 Popek \& Goldberg 정리에 따라 가상화할 수 없었던 이유를 설명하라. 바이너리 변환과 반가상화가 각각 어떻게 우회했는지 강점과 비용을 하나씩 들고, VT-x가 아키텍처 수준에서 문제를 해결한 방법을 설명하라.}

\qitem{2}{익명 사용자를 대신해 LLM이 작성한 코드를 실행하는 공개 "AI 에이전트" 서비스의 샌드박스를 설계한다. 격리 스펙트럼(컨테이너 / gVisor / Firecracker microVM / 완전한 VM)에서 한 지점을 선택하고, 커널 장애 피해 범위, 시작 지연 시간, 호스트당 밀도, 운영 복잡성을 고려해 양옆의 두 선택지보다 나은 이유를 논증하라.}

\qitem{3}{팀이 H100-80GB 한 대를 보유하고 있으며 p99 SLO가 있는 내부 모델 세 개(각각 \textasciitilde{}15 GB VRAM)와 카드 전체를 기꺼이 사용할 야간 배치 임베딩 작업을 실행해야 한다. MIG, MPS, 시간 분할, 시간대별 스케줄링을 자유롭게 조합해 GPU 공유 계획을 제안하고 격리, 사용률, 운영 단순성을 근거로 정당화하라. 계획 검증을 위해 무엇을 측정할지도 밝혀라.}

\namedsection{정답 및 해설}{}

\subsection*{A. 객관식}

\qitem{1}{\textbf{④} 세 요구 사항은 동등성, 효율성, 자원 제어다. 아키텍처 간 이식성은 포함되지 않는다. VMM은 자신이 실행되는 것과 \textit{같은} 아키텍처를 제공한다. (§2.2, 상자 2-1)}

\qitem{2}{\textbf{②} 정확히 VM exit다. 하드웨어가 게스트 상태를 VMCS에 저장하고 exit 이유를 기록한 뒤 root 모드로 진입한다. ①은 VT-x \textit{이전}의 병리 현상이고, ④는 소프트웨어 시대 VMware의 우회책이었다. (§2.3)}

\qitem{3}{\textbf{①} EPT/NPT는 CPU가 순회하는 두 번째 하이퍼바이저 소유 변환을 추가한다. 게스트는 exit 없이 자체 페이지 테이블을 관리하며 섀도 테이블은 사라진다. 비용은 TLB 미스 지연 시간으로 이동하므로 대형 페이지가 필요하다. (§2.3)}

\qitem{4}{\textbf{③} e1000 환상에서는 레지스터 접근마다 exit가 발생하지만, virtio의 정직한 인터페이스는 공유 링에 디스크립터를 배치하고 배치당 \textasciitilde{}한 번만 깨운다. 협력이 환상을 이긴다는 반가상화의 오래가는 교훈이다. (§2.4)}

\qitem{5}{\textbf{②} 컨테이너는 공유 커널 하나의 \textit{뷰}와 \textit{몫}을 격리한다(네임스페이스/cgroup, 제3장). VM은 테넌트와 호스트 사이에 하드웨어 경계와 비공개 커널을 둔다. 그림 2-5의 전체 축이 바로 이 차이다. (§2.5)}

\qitem{6}{\textbf{③} MIG만이 SM, L2, DRAM 등 실리콘 자체를 나누어 인스턴스별 하드웨어급 QoS를 제공한다. 시간 분할과 MPS는 하드웨어 경계 없이 시간적으로 또는 동시에 공유한다. (§2.6)}

\qitem{7}{\textbf{③} 사례 2-A의 증거(DMI 문자열, \code{systemd-detect-virt}, \code{kvm-clock})는 머신이 가상임을 입증한다. 기록 2-4는 그럼에도 GPU는 실제임을 보여 준다. VFIO/IOMMU를 통한 패스스루가 물리 카드를 게스트에 매핑한다. 가상 머신 안의 실제 장치는 모순이 아니라 의도한 결과다. (§2.1, §2.6)}

\subsection*{B. 단답형}

\qitem{1}{\textbf{VMX root 모드}(하이퍼바이저)와 \textbf{VMX non-root 모드}(게스트)다. 둘 다 rings 0–3을 포함하며, 게스트 커널은 하드웨어의 감독 아래 ring 0을 되찾는다.}

\qitem{2}{vCPU가 실행 가능할 때 호스트 스케줄러가 다른 테넌트의 스레드를 실행해 "CPU"가 보이지 않게 멈춘다. 전용 코어 인스턴스 유형을 쓰거나, 호스트당 vCPU를 줄여 오버커밋을 낮추거나, 지연 시간에 중요한 서빙을 공유 계열에서 옮길 수 있다.}

\qitem{3}{㉠ \textbf{PCIe 패스스루}(VFIO 사용), ㉡ \textbf{IOMMU}(Intel VT-d / AMD-Vi)다.}

\qitem{4}{GPU 호스트 스케줄링/탐색, 수 GB 컨테이너 이미지 가져오기, 수십 GB 모델 가중치를 GPU 메모리에 로드하기, CUDA/런타임 초기화와 캐시 워밍업 중 두 가지를 쓰면 된다. microVM 자체는 병목이 아니게 되었으며, 이것이 Firecracker의 성취이자 12주 차의 출발점이다.}

\qitem{5}{\textbf{결정적인 제약 차원:} 제약은 VRAM 11 GB이며 L4의 24 GB로 이미 충족한다. A100의 추가 56 GB와 FLOPs는 제약이 되지 않으므로 순수한 비용이다. \textbf{사용률/비용:} 계량기는 항상 시간당 돌아간다. 같은 SLO를 제공하면서 가격이 \textasciitilde{}5×면 요청당 비용은 \textasciitilde{}5× 늘고 비용을 낸 카드의 측정 사용률은 떨어진다. 여기서 "안전"은 "안전하게 낭비한다"는 뜻이다.}

\subsection*{C. 서술형(채점 기준)}

\qitem{1}{만점 답안에는 다음이 필요하다. 정리(민감 ⊆ 특권)와 이를 위반해 \textit{조용히} 실패하는 x86 명령어 17개(POPF가 대표 예시), 바이너리 변환 = 실행 중 게스트 커널 코드를 다시 쓰고 변환을 캐시(강점: 수정하지 않은 게스트, 비용: 엄청난 복잡성과 경계 사례), 반가상화 = 게스트 소스를 수정해 하이퍼콜(강점: 속도와 배칭, 비용: 게스트 변경이 필요하고 영구히 유지), VT-x = 직교하는 root/non-root 모드 + VMCS로 게스트에 ring 0을 돌려주고 모든 민감 작업이 하드웨어에서 확실히 트랩하게 함. virtio가 반가상화의 후손임을 언급하면 가산점이 있다. (§2.2–2.3)}

\qitem{2}{좋은 답은 \textbf{Firecracker형 microVM}을 선택하고 다음과 같이 논증한다. 커널 0-day에 대한 하드웨어 경계(공유 커널이나 Go 프록시가 상한인 컨테이너/gVisor와 대비), 대화형 예산 안의 \textasciitilde{}125 ms 부팅, 수천 개 밀도를 유지하는 \textasciitilde{}5 MB 오버헤드, 프로덕션 선례(Lambda)다. 컨테이너보다 구성 요소가 많고, 시스템 호출 호환성은 gVisor보다 단순하지만 장치/GPU 접근은 컨테이너보다 어렵다는 점도 솔직히 인정해야 한다. GPU 접근과 폭넓은 Linux 호환성이 중요하면 gVisor도 방어 가능한 대안이다. 선택한 브랜드가 아니라 \textit{절충에 대한 추론}을 평가한다. (§2.5)}

\qitem{3}{깔끔한 계획 하나는 H100을 SLO 서비스용으로 크기를 맞춘 격리 조각 세 개(예: 각각 2g.20gb ≈ 60 GB, p99용 하드웨어 QoS)로 MIG 파티셔닝하고, \textasciitilde{}1g.10gb는 오버플로/카나리용으로 남기는 것이다. 야간 배치 작업은 남은 조각에서 계속 실행하거나 한산한 시간에 \textit{재구성}해 실행한다. MIG 재구성은 중단을 일으키므로, 조각에서 MPS를 통해 배치 작업을 속도 제한하여 실행하는 방식과 비교해야 한다. 함께 배치했을 때와 격리했을 때의 서비스별 p99, 조각 사용률, 구성별 배치 작업 처리량을 반드시 측정해야 한다. 사용률과 격리의 대립이 하드웨어 조절 장치를 얻은 제1장의 비용 긴장임을 알아차린 답안에 가산점을 준다. (§2.6)}

\namedsection{더 읽을거리}{}

\textbf{[이번 주]} 표시는 지금 읽어야 할 필독 자료이고, \textbf{[N주 차]} 표시는 해당 주차에 도움이 될 자료다.

\begin{tbbitemize}
  \item \textbf{[이번 주] A. Agache et al., "Firecracker: Lightweight Virtualization for Serverless Applications," NSDI 2020} — 필독 자료다. 짧고 구체적이며 이 장과 서빙을 연결하는 가장 좋은 단일 자료다. 설계 제약(§2–3)과 밀도/부팅 수치(§5)를 읽고, 전반에서 12주 차의 콜드 스타트 경제성을 염두에 둔다.
  \item \textbf{G. Popek \& R. Goldberg, "Formal Requirements for Virtualizable Third Generation Architectures," CACM 1974} — 이 장에서 사용한 동등성, 효율성, 자원 제어, 민감 명령어와 특권 명령어라는 용어를 제시한 1974년 논문이다. 50년 된 형식 논문치고 놀라울 만큼 읽기 쉽다.
  \item \textbf{E. Bugnion, J. Nieh, D. Tsafrir, "Hardware and Software Support for Virtualization" (Morgan \& Claypool, 2017)} — x86의 17가지 문제, 바이너리 변환 내부 구조, VT-x/EPT 설계, KVM 사례 연구까지 §2.2–2.4의 모든 내용을 책 한 권 분량으로 가장 잘 다룬다. 이 장의 깊이가 잘 맞았다면 자연스러운 다음 단계다.
  \item \textbf{P. Barham et al., "Xen and the Art of Virtualization," SOSP 2003} — 하이퍼콜, 분할 드라이버, virtio로 이어진 정직한 인터페이스 철학을 다룬 반가상화의 고전이다. 역사적으로 초기 클라우드가 이 논문 위에서 실행되었다.
  \item \textbf{R. Russell, "virtio: Towards a De-Facto Standard for Virtual I/O Devices," ACM OSR 2008} — 짧고 유쾌하게 비형식적이며, 설계자가 virtqueue를 설명한다. 현대적인 세부 사항이 필요하면 현재 OASIS virtio 명세와 함께 읽는다.
  \item \textbf{[제3장 뒤] gVisor 문서("Architecture Guide")와 Kata Containers 문서} — 네임스페이스/cgroup으로 스펙트럼의 컨테이너 쪽이 구체화되는 제3장 뒤에 읽는다. 둘 다 제4장의 런타임 이야기에서 OCI 런타임인 \code{runsc}, \code{kata-runtime}으로 들어간다.
  \item \textbf{[12주 차] NVIDIA MIG User Guide와 NVIDIA vGPU 문서} — §2.6의 운영 세부 사항인 프로필, 재구성, K8s 장치 플러그인 통합을 다룬다. 12주 차 GPU 공유 실습에서 직접 사용할 자료가 된다.
  \item \textbf{KVM API 문서(kernel.org, `api.rst`)} — ioctl 목록을 한 번 훑어보자. \code{KVM\_CREATE\_VM} / \code{KVM\_RUN}이 \textasciitilde{}40개 호출로 이루어진 것을 보면 어떤 그림보다 "하이퍼바이저가 된 Linux"의 신비가 잘 풀린다. 작은 교육용 VMM인 \code{kvmtool}도 읽기 좋은 동반 자료다.
  \item \textbf{Nitro System에 관한 AWS re:Invent 발표} — 장치 모델을 하드웨어로 옮기고 호스트 하이퍼바이저를 거의 아무것도 남지 않을 정도로 줄인 과정을 가장 명확하게 공개한다. §2.3–2.4의 exit 제거 흐름을 제품 이야기로 들려주는 부분에 주목한다.
  \item \textbf{[이번 주의 흥미로운 읽을거리] Melinda Varian, "VM and the VM Community: Past, Present, and Future" (1997)} — 온라인에서 무료로 읽을 수 있는 CP-67과 VM/370의 결정적인 내부 역사다. Cambridge의 비공식 프로젝트가 1967년 가상 머신을 발명한 과정과 IBM이 이를 거의 없앨 뻔한 여러 순간을 다룬다. Amazon의 원래 EC2 베타 발표(2006년 8월 25일, 지금도 AWS 블로그에 있다)와 함께 읽고, 40년 간격으로 작성된 두 문서를 사례 2-B 이야기의 1막과 3막으로 보자.
\end{tbbitemize}

\vspace{5.0pt}

\begin{calloutbox}{tbbblue}{tbbbluefill}
{\normalsize\bfseries\color{tbbblue}다음 장 — 제3장: 컨테이너 내부}\par\vspace{3pt}
이 장에서는 하드웨어 경계와 완전한 두 번째 커널이라는 무거운 방식으로 격리를 구축했다. 다음 주에는 가벼운 방식으로 만들며, 컨테이너라는 \textit{물건}은 없고 레시피만 있다는 사실을 발견한다. Linux \textbf{네임스페이스}는 프로세스 그룹에 PID, 네트워크, 마운트, 호스트 이름, 사용자에 관한 자체 세계관을 제공한다. \textbf{cgroups v2}는 CPU와 메모리 몫을 제한하며, 언젠가 추론 서버를 죽일 OOM 킬러와 그 종료 코드 137을 소개한다. 계층형 \textbf{파일 시스템}(overlayfs)은 이식 가능한 이미지를 제공한다. Liz Rice의 유명한 라이브 코딩 시연을 본보기로 \code{unshare}를 사용해 처음부터 직접 만들고, 5주 차에 Kubernetes가 오케스트레이션할 바로 그 메커니즘을 만난다. 한 가지 경고가 있다. 다음 주의 첫 기록에는 PID 41288인 프로세스가 등장하지만 장은 그 정체를 설명해 주지 않는다. 제2장의 스펙트럼과 인내심을 가져오자. 4주 차에 보상을 받는다.
\end{calloutbox}

%% Chapter 3 — 컨테이너 내부
\setcounter{chapter}{2}
\chapter{컨테이너 내부}
\chaptersubtitle{네임스페이스, cgroup, overlayfs — 모델 서버를 담아 보내는 상자 만들기}

\begin{calloutbox}{tbbblue}{tbbbluefill}
{\normalsize\bfseries\color{tbbblue}학습 목표}\par\vspace{3pt}
\begin{tbbitemize}
  \item "컨테이너는 일반 프로세스 주위를 네임스페이스 + cgroup + 계층형 파일 시스템으로 감싼 것"이라는 문장을 설명하고, 커널에는 컨테이너 객체 자체가 없다는 더 강한 주장을 논증한다.
  \item 고전적인 네임스페이스 여섯 개가 각각 무엇을 가상화하는지 설명하고, PID 네임스페이스 의미론(PID 1의 의무, 신호, 좀비)과 NET 네임스페이스 배선(veth, 브리지, 포트)의 결과를 추론한다.
  \item 파일 시스템 인터페이스를 통해 cgroups v2를 사용한다. CPU 할당량과 메모리 한도를 설정하고, 사용량과 이벤트를 읽으며, CPU 스로틀링이 요청 지연 시간에 미치는 영향을 직접 계산한다.
  \item 프로세스가 memory.max를 넘을 때 메모리 회수, OOM 강제 종료, SIGKILL, 종료 코드 137로 이어지는 정확한 과정을 설명하고, 실제 메모리 사용 형태를 수치와 함께 따져 추론 서버의 메모리 한도를 적절히 산정한다.
  \item overlayfs가 이미지 계층을 루트 파일 시스템으로 조립하는 방식을 설명하고, 셸에서 copy-up과 whiteout을 시연하며, 계층 위생을 이미지 크기·빌드 시간·콜드 스타트 지연 시간과 연결한다.
  \item unshare, pivot\_root, /proc, cgroup으로 작동하는 컨테이너를 직접 만들면서 모든 명령의 출력을 읽고, Docker가 그 위에 정확히 무엇을 더하는지 설명한다(다음 주 주제).
\end{tbbitemize}
\end{calloutbox}

\vspace{5.0pt}

\begin{calloutbox}{tbbteal}{tbbtealfill}
{\normalsize\bfseries\color{tbbteal}핵심 용어}\par\vspace{3pt}
네임스페이스(PID · NET · MNT · UTS · IPC · USER) · unshare / setns · PID 1 · 좀비 회수 · veth 쌍 · 브리지 · 루트리스 컨테이너 · cgroups v2 · 컨트롤러 · cpu.max · 스로틀링 · cpu.weight · memory.max · memory.high · memory.events · OOM 킬러 · 종료 코드 137 · oom\_score\_adj · overlayfs · lowerdir / upperdir / merged · copy-up · whiteout · 이미지 계층 · pivot\_root · chroot
\end{calloutbox}

\vspace{5.0pt}

\begin{calloutbox}{tbblightgray}{tbbboxfill}
{\normalsize\bfseries\color{tbbink}장 구성}\par\vspace{3pt}
\textbf{3.1} 컨테이너라는 것은 없다    \textbf{3.2} 네임스페이스: 세계를 보는 비공개 관점    \textbf{3.3} cgroups v2: 머신의 제한된 몫    \textbf{3.4} overlayfs: 계층으로 이루어진 파일 시스템    \textbf{3.5} 직접 만들기 — Docker가 실제로 더하는 것    \textbf{3.6} 심판이 모델을 죽일 때: OOM 대응 지침    \textbf{3.7} 실습: 미니 컨테이너와 관찰한 OOM — 이후 요약, 연습문제, 더 읽을거리로 이어진다.
\end{calloutbox}

\vspace{3.0pt}

\begin{calloutbox}{tbbamber}{tbbamberfill}
{\normalsize\bfseries\color{tbbamberdark}이 장을 읽는 방법}\par\vspace{3pt}
이 장은 \textbf{터미널 옆에서 읽도록} 썼다. 아래의 어두운 블록은 모두 어느 Linux VM에서든 재현할 수 있는 실제 명령 기록이다. 제2장의 인스턴스나 로컬 Linux를 사용하면 되며, root 권한이 필요한 명령은 sudo를 사용한다고 가정한다. 첫 시연을 제외하면 Docker를 설치할 필요가 전혀 없다. 바로 그것이 핵심이다.
\end{calloutbox}

\section{컨테이너라는 것은 없다}

제2장은 실제 GPU까지 포함한 임대 VM이라는 커널 바로 위 계층에서 끝났다. 이 장은 마지막 한 계단을 내려가 모든 서빙 시스템이 실제로 담겨 배포되는 패키징에 도달한다. 일부러 약간 도발적으로 벼린 주장으로 시작한다. \textbf{컨테이너라는 것은 없다.} Linux 소스에서 컨테이너 객체인 \code{struct container}나 \code{container\_create()} 시스템 호출을 찾아도 나오지 않는다. \code{struct task\_struct}(프로세스), \code{struct pid\_namespace}, \code{struct mem\_cgroup}은 있지만 컨테이너는 없다. 커널은 컨테이너라는 말을 들어본 적이 없다.

커널에 \textit{실제로} 있는 것은 프로세스와 그 주위를 감쌀 수 있는 독립적인 메커니즘 세 가지다. 프로세스가 \textit{볼 수 있는 것}을 바꾸는 \textbf{네임스페이스}, \textit{사용할 수 있는 것}을 바꾸는 \textbf{cgroup}, \textit{그 세계가 어디에서 오는지}를 바꾸는 마운트 메커니즘과 계층형 파일 시스템(\textbf{overlayfs})이다. "컨테이너"는 그저 세 가지를 한꺼번에 적용하는 관례다(그림 3-1). \code{docker run}을 입력한다고 신기한 객체가 튀어나오는 것이 아니다. 평범한 Linux 프로세스가 세 벌의 의상을 입고 시작된다.

\begin{tbbfigure}
  \adjustbox{max width=0.94\linewidth,max totalheight=0.46\textheight}{\includegraphics{figures/fig3-1_recipe.png}}\\[3pt]
  \figcaption{그림 3-1  docker run 분해하기. 컨테이너는 공유 커널 위의 일반 프로세스에 네임스페이스 + cgroup + 계층형 파일 시스템을 적용하는 레시피다.}
\end{tbbfigure}

\subsection*{직접 눈으로 확인하기}

이런 주장은 믿지 말고 확인해야 한다. Docker가 설치된 머신에서 컨테이너를 시작한 뒤 \textit{호스트에서} 심문해 보자. Docker 명령이 아니라 첫 운영체제 수업 때부터 써 온 평범한 프로세스 도구를 쓴다.

\begin{terminalbox}
\begin{Verbatim}[breaklines=true,breakanywhere=true,fontsize=\small,formatcom=\color{tbbtermfg},codes=\tbbcodeactive]
$ docker run -d --name demo -m 512m --cpus 1 python:3.11 sleep 9999
3f9a1c2b77e1...
$ docker inspect --format '{{.State.Pid}}' demo     # ask Docker for the host PID
41320
$ ps -o pid,ppid,rss,cmd -p 41320                  # plain ps, on the host
  PID   PPID   RSS  CMD
41320  41288  9184  sleep 9999
$ cat /proc/41320/cgroup                           # which budget group is it in?
0::/system.slice/docker-3f9a1c2b77e1....scope
$ cat /sys/fs/cgroup/system.slice/docker-3f9a*.scope/memory.max
536870912                                          # there is your -m 512m
$ ls -l /proc/41320/ns/pid                         # and its private view
... /proc/41320/ns/pid -> pid:[4026532871]
\end{Verbatim}
\end{terminalbox}
\figcaption{기록 3-1  일반 프로세스 도구를 사용해 호스트에서 실행 중인 컨테이너를 심문한 결과. cgroup과 네임스페이스가 연결된 프로세스(PID 41320)다.}

기록을 아래에서 위로 읽으면 다섯 줄에 이 장 전체가 예고되어 있다. 컨테이너는 \textbf{호스트 ps에 보인다}. 런타임 shim의 자식인 프로세스 41320이며, 다른 프로세스처럼 메모리 9 MB를 사용한다. \code{-m 512m} 플래그는 \textbf{파일 속 숫자}가 되었다(\code{memory.max = 536870912}, 3.3절). "격리"는 \code{/\allowbreak proc/\allowbreak 41320/\allowbreak ns/\allowbreak } 아래의 \textbf{네임스페이스 링크 집합}이다(3.2절). 호스트 커널은 이 프로세스를 스케줄링하고 사용량을 계산하며, 3.6절에서 생생히 보여 주듯 죽일 수도 있다. Docker의 기여는 설정을 오케스트레이션하는 것이지 새로운 종류의 무언가를 만들어 내는 것이 아니다.

\subsection*{신비를 벗기는 일이 중요한 이유(철학적인 이유만은 아니다)}

제2장의 격리 스펙트럼에서 이미 이 사실을 한 번 사용했다. 컨테이너가 호스트 커널을 공유하므로 멀티테넌트 플랫폼은 gVisor나 microVM을 선택한다. 그러나 레시피 관점은 이 책 내내 매주 운영 현장에서 값을 한다.

\begin{tbbitemize}
  \item \textbf{이 책의 나머지 동안 모델 서버는 컨테이너다.} 4주 차에는 이미지로 패키징하고, 5주 차에는 Kubernetes가 스케줄링하며, 9\textasciitilde{}12주 차에는 그 안에서 서빙한다. 자원 한도, 상태 검사, 네트워킹, GPU 접근 등 시스템이 컨테이너에 하는 모든 일은 결국 이 장의 세 메커니즘으로 귀결된다. 오작동하면 이 계층에서 디버깅한다.
  \item \textbf{이 장의 대미를 장식하는 것은 실전에서 가장 흔한 서빙 장애다.} \code{OOMKilled}, 종료 코드 137과 함께 애플리케이션 로그 없이 죽는 추론 pod도 memory.max가 cgroup 파일이고 OOM 킬러가 이를 강제한다는 사실을 알면 수수께끼가 아니다. 3.6절은 언젠가 오전 2시에 읽게 될 실제 \code{dmesg} 줄과 \code{kubectl} 출력을 통해 이 사고를 체크리스트로 바꾼다.
  \item \textbf{이미지 계층은 지연 시간 예산이다.} PyTorch/CUDA 이미지는 5–10 GB에 이른다. 새 노드로 이미지를 가져오는 시간이 12주 차 콜드 스타트 청구서에서 단일 항목으로는 최대일 때가 많다. \code{RUN rm}을 실행해도 아무것도 작아지지 않는 이유를 포함해 overlayfs(3.4절)를 이해하면 이미지 축소가 맹목적인 관습이 아니라 합리적인 작업이 된다.
  \item \textbf{디버깅의 성격이 바뀐다.} 컨테이너가 프로세스라는 사실을 알면 기록 3-1의 도구 상자를 어디서든 사용할 수 있다. \code{ps}, \code{/\allowbreak proc/\allowbreak <pid>/\allowbreak ns}, cgroup 파일, 내부로 들어가는 \code{nsenter}다. 상자가 더는 불투명하지 않고, "컨테이너에서는 작동합니다"가 대화의 끝이 아니게 된다.
\end{tbbitemize}

이 장의 구성은 레시피를 반영한다. 3.2\textasciitilde{}3.4절은 보는 것, 사용하는 것, 존재하는 것에 해당하는 재료를 하나씩 다루며, 각각 명령 기록과 계산 예시를 제시한다. 3.5절에서는 이번 주 시청 자료인 Liz Rice의 유명한 라이브 코딩 시연 구조를 따라 셸만으로 재료를 조립해 작동하는 컨테이너를 만든다. 3.6\textasciitilde{}3.7절에서는 이렇게 얻은 이해를 이 장에서 가장 가치 있는 운영 기술 하나에 쓴다. OOM 킬러를 예측하고 관찰하며 살아남는 법이다.

\begin{calloutbox}{tbbpurple}{tbbpurplefill}
{\normalsize\bfseries\color{tbbpurple}생각해 보기}\par\vspace{3pt}
\begin{tbbitemize}
  \item Docker for Mac/Windows는 내부에서 Linux VM을 실행한다. 이 절의 주장과 제2장을 이용해 왜 그래야 하는지 설명해 보자. 이 사실은 노트북 위에서 "컨테이너는 VM보다 가볍다"는 말에 어떤 함의를 갖는가?
  \item 기록 3-1에서 sleep의 부모는 입력한 docker CLI가 아니라 "shim" 프로세스(PID 41288)다. 수명이 긴 중간자가 유용한 이유를 추측해 보자. (힌트: Docker 데몬이 다시 시작되면 컨테이너에는 무슨 일이 일어나야 하는가?) 4주 차에 확인한다.
  \item 컨테이너가 그저 프로세스라면 \code{docker exec}는 무엇을 해야 하는가? 네임스페이스를 힌트 삼아 답을 그려 본 뒤, §3.2의 setns와 맞춰 보자.
\end{tbbitemize}
\end{calloutbox}

\section{네임스페이스: 세계를 보는 비공개 관점}

네임스페이스는 프로세스 테이블, 네트워크 스택, 마운트 테이블 같은 커널의 전역 자원 하나를 가져와 프로세스 그룹에 \textbf{그 관점의 비공개 사본}을 제공한다. 아무것도 에뮬레이션하지 않고 두 번째 커널도 나타나지 않는다(두 번째 커널을 띄우는 것은 제2장의 방식이었다). 커널은 네임스페이스마다 장부를 따로 관리하고 각 프로세스에 자신이 속한 장부를 보여 줄 뿐이다. PID 네임스페이스의 프로세스가 \code{getpid()}를 호출하면 커널은 호출자가 어느 네임스페이스에 속하는지 찾아 \textit{그곳에서 보이는} ID를 반환한다. 같은 커널과 같은 코드 경로지만 보는 사람마다 답이 다르다. 그림 3-2는 호스트와 컨테이너에 보이는 것을 비교하면서 고전적인 네임스페이스 여섯 개를 둘러본다.

\begin{tbbfigure}
  \adjustbox{max width=0.94\linewidth,max totalheight=0.46\textheight}{\includegraphics{figures/fig3-2_namespaces.png}}\\[3pt]
  \figcaption{그림 3-2  여섯 네임스페이스와 여섯 가지 비공개 차원: PID, 네트워크, 마운트, 호스트 이름, IPC, 사용자 ID. 원하는 것만 골라 조합할 수 있다.}
\end{tbbfigure}

\subsection*{API는 호출 세 개와 디렉터리 하나다}

그 결과에 비해 네임스페이스 API는 작다. \code{unshare()}는 \textit{호출한} 프로세스를 새로 만든 네임스페이스로 옮긴다. \code{unshare} 명령은 이를 감싸며 실습의 주력 도구다. \code{clone()}은 새 네임스페이스 안에 \textit{자식}을 바로 만들며 컨테이너 런타임이 사용한다. \code{setns()}는 핸들을 받아 기존 네임스페이스에 \textit{합류}한다. 3.1절의 예고에 답하므로 잠시 주목할 만하다. \code{docker exec}는 겉모습만 바꾼 \code{setns()}다. 대상 컨테이너의 네임스페이스 파일을 열고 합류한 뒤 그곳에서 셸을 실행한다. 마법의 문은 없다. 프로세스가 읽을 장부를 바꿀 뿐이다.

핸들은 검사할 수 있는 디렉터리에 있다. 모든 프로세스는 \code{/\allowbreak proc/\allowbreak <pid>/\allowbreak ns/\allowbreak }에 심볼릭 링크 형태로 네임스페이스 소속을 공개하며, 대괄호 속 숫자가 네임스페이스의 신원이다. \textbf{두 프로세스의 숫자가 같으면 같은 네임스페이스이고, 다르면 다른 네임스페이스다. 예외는 없다.}

\begin{terminalbox}
\begin{Verbatim}[breaklines=true,breakanywhere=true,fontsize=\small,formatcom=\color{tbbtermfg},codes=\tbbcodeactive]
$ ls -l /proc/self/ns/            # my shell, on the host
mnt -> mnt:[4026531841]   net -> net:[4026531840]
pid -> pid:[4026531836]   uts -> uts:[4026531838]  ...
$ ls -l /proc/41320/ns/           # the container process from Transcript 3-1
mnt -> mnt:[4026532868]   net -> net:[4026532866]
pid -> pid:[4026532871]   uts -> uts:[4026532864]  ...
# every number differs -> different views on every axis.
# same number would mean: shared namespace. Free debugging test.
\end{Verbatim}
\end{terminalbox}
\figcaption{기록 3-2  네임스페이스 신원은 검사할 수 있다. inode 번호가 같으면 네임스페이스를 공유한다. "컨테이너에서 X가 보이지 않는다"면 가장 먼저 확인할 항목이다.}

네임스페이스는 원하는 것만 골라 조합할 수 있다. 프로세스가 NET만 분리할 수도 있고(네트워크 테스트에서 흔하다), 모든 것을 한꺼번에 분리할 수도 있다(컨테이너). 더 늦게 등장한 cgroup 네임스페이스와 시간 네임스페이스까지 더하면 현대의 여덟 개가 되지만, 이 장은 고전적인 여섯 개를 다룬다.

\subsection*{PID: 프로세스 하나, 두 가지 진실}

PID 네임스페이스는 컨테이너에 자체 PID 1을 뿌리로 하는 프로세스 트리를 제공한다(그림 3-3 왼쪽). 따라서 같은 커널 태스크가 호스트에는 PID 41320, 자신에게는 PID 1이라는 두 신원을 갖는다. 두 번째 신원이 생기는 모습을 직접 볼 수 있다.

\begin{terminalbox}
\begin{Verbatim}[breaklines=true,breakanywhere=true,fontsize=\small,formatcom=\color{tbbtermfg},codes=\tbbcodeactive]
$ sudo unshare --pid --fork --mount-proc bash    # new PID ns + fresh /proc
$ echo $$                                        # "what is my PID?"
1
$ ps aux
USER  PID  %CPU  %MEM  ... CMD
root    1   0.0   0.0  ... bash
root    8   0.0   0.0  ... ps aux
# two processes exist in this whole "machine" - and I am init.
\end{Verbatim}
\end{terminalbox}
\figcaption{기록 3-3  새 PID 네임스페이스 내부. 셸은 PID 1이고 프로세스 테이블에는 항목 두 개가 있다. (--fork가 중요하다. PID 소속은 자식에게 적용되므로 unshare가 셸을 새 세계로 fork해야 한다.)}

\code{--mount-proc} 플래그가 함께 따라오는 이유는 무엇인가? \code{ps}는 커널에 직접 묻지 않고 \textit{/proc 파일 시스템을 읽는다}. \code{/\allowbreak proc}가 여전히 호스트의 마운트를 보여 주면 새 네임스페이스 안에서도 \code{ps}는 호스트의 수천 개 프로세스를 거리낌 없이 보고한다. \code{/\allowbreak proc}를 다시 마운트하면 도구가 새로운 진실을 말한다. 조금 철학적인 이 교훈을 기억하자. \textbf{Linux 시스템에 대한 관점은 마운트가 정직한 만큼만 정직하다.} 실습의 4단계와 14주 차 모니터링 각주에서 다시 등장한다.

\subsection*{PID 1은 숫자가 아니라 직무다}

Unix는 PID 1(init)에 의무 두 가지를 부여한다. 대부분의 애플리케이션 코드는 이를 생각해 볼 일이 없다 — 컨테이너에 담겨 갑자기 \textit{자신이} PID 1이 되기 전까지는. 첫째는 \textbf{신호}다. PID 1에 대해 커널은 보통 SIGTERM을 치명적으로 만드는 기본 핸들러를 비활성화한다. 명시적인 SIGTERM 핸들러가 없는 프로세스가 PID 1로 실행되면 SIGTERM을 그냥 \textit{무시한다}. 둘째는 \textbf{회수}다. 프로세스가 죽으면 부모가 \code{wait()}를 호출해 종료 상태를 수집해야 한다. 그때까지 프로세스 테이블에 죽은 항목인 \textbf{좀비}로 남는다. 고아 프로세스는 PID 1에 다시 연결되므로, PID 1은 다른 프로세스의 자식을 영원히 회수해야 한다. 그렇지 않으면 \code{pids.max}나 노드가 질식할 때까지 좀비가 쌓인다.

이 두 의무를 무시하는 모델 서버에서 Kubernetes 롤링 배포(13주 차)가 어떻게 흘러가는지 돌려 보자. t=0에 kubelet이 컨테이너의 PID 1, 즉 핸들러를 설치하지 않은 Python 서버에 SIGTERM을 보낸다. 아무 일도 일어나지 않고, 서버는 끝내지 못할 요청을 계속 받는다. 기본 유예 기간인 t=30 s가 되면 kubelet이 인내심을 잃고 SIGKILL을 보낸다. 서버는 요청 처리 중간에 죽고, 처리 중인 모든 요청이 오류를 낸다. 이 과정이 \textbf{모든 배포마다} 반복된다. 스스로 만들어 낸 영구 장애 제조기인데, 겉보기에는 "K8s가 불안정하다"로 보인다. 해결책은 작고 표준적이다.

\begin{tbbitemize}
  \item \textbf{서버에서 SIGTERM을 처리한다.} 새 요청을 받지 않고, 처리 중인 요청을 비운 뒤 종료한다. FastAPI/uvicorn과 대부분의 제대로 된 서버는 신호를 받기만 하면 이 작업을 한다. 다음 항목을 보자.
  \item \textbf{신호가 서버에 도달하는지 확인한다.} 고전적인 Dockerfile 버그가 있다. \textit{셸 형식}의 \code{CMD python serve.py}는 \code{/\allowbreak bin/\allowbreak sh -c "python serve.py"}를 실행한다. 셸이 PID 1이며 자식에게 SIGTERM을 전달하지 않는다. exec 형식(\code{CMD ["python", "serve.py"]})을 사용해 서버 \textit{자체가} PID 1이 되고 실제 신호를 받게 한다.
  \item \textbf{또는 진짜 init을 설치한다.} \code{tini}(플래그 하나: \code{docker run --init})는 신호를 전달하고 좀비를 회수하는 수백 줄짜리 PID 1이다. 애플리케이션이 다시 평범한 자식이 될 수 있다.
\end{tbbitemize}

\begin{tbbfigure}
  \adjustbox{max width=0.94\linewidth,max totalheight=0.46\textheight}{\includegraphics{figures/fig3-3_pidnet.png}}\\[3pt]
  \figcaption{그림 3-3  중요한 결과를 만드는 두 네임스페이스. 왼쪽 PID는 프로세스 하나의 두 신원과 PID 1의 의무를 보여 준다. 오른쪽 NET은 veth 쌍을 통해 호스트 브리지에 연결된 비공개 스택이다.}
\end{tbbfigure}

\subsection*{NET: 비공개 스택과 요청 하나의 생애}

NET 네임스페이스는 여섯 개 중 가장 물리적으로 느껴진다. 컨테이너가 인터페이스, IP 주소, 라우팅 테이블, 방화벽 규칙, 그리고 무엇보다 \textbf{자체 포트 공간}을 포함한 완전한 비공개 네트워크 스택을 얻는다. 호스트 하나의 모델 서버 레플리카 20개가 충돌 없이 모두 \code{:8000}에 바인딩할 수 있다. 이 덕분에 레플리카 서버군을 관리할 수 있다. 1998년처럼 포트를 하나씩 배정할 필요가 없다.

갓 만든 비공개 스택은 아직 \textit{아무 데도 연결되지 않은} 스택이기도 하다. 새 NET 네임스페이스에는 루프백 장치만 있고 그마저 내려가 있다. 표준 배선(그림 3-3 오른쪽)은 패치 케이블처럼 연결한 가상 NIC 두 개, 즉 \textbf{veth 쌍}을 만든다. 제2장의 virtio 정신을 네트워킹에 적용한 것이다. 한쪽 끝은 컨테이너 안에 두고 \code{eth0}으로 이름을 바꾸며, 다른 끝은 호스트의 소프트웨어 \textbf{브리지}(Docker의 \code{docker0})에 연결한다. 구체적으로, 노트북에서 출발한 요청 하나가 \code{-p 8080:8000}으로 공개된 컨테이너의 8000번 포트에 닿기까지를 따라가 보자.

\begin{tbbitemize}
  \item \textbf{도착:} \code{host:8080}으로 주소가 지정된 패킷이 호스트의 실제 NIC에 도착한다.
  \item \textbf{변환:} \code{-p 8080:8000}을 지정할 때 Docker가 설치한 iptables/nftables DNAT 규칙이 목적지를 \code{172.17.0.2:8000}으로 고쳐 쓴다. \textit{공개한 포트는 문자 그대로 NAT 규칙 하나다.}
  \item \textbf{브리지와 통과:} 호스트가 패킷을 \code{docker0}으로 라우팅하고, 브리지가 컨테이너의 veth로 전달한다. 패킷은 NET 네임스페이스 안의 \code{eth0}에 나타나며 서버가 \code{:8000}에서 수신 대기한다.
  \item \textbf{반환:} 응답은 경로를 거슬러 간다. 연결 추적이 소스의 NAT를 되돌리므로 노트북은 아무것도 모른 채 \code{host:8080}의 응답을 본다.
\end{tbbitemize}

5주 차를 위해 미리 기록해 둘 내용이 있다. Kubernetes \textbf{pod}는 하나의 NET 네임스페이스를 소유하는 작은 "pause" 자리 표시자 프로세스를 통해 \textbf{NET 네임스페이스 하나를 공유하는} 컨테이너로 정의된다. 이 한 가지 설계 결정 때문에 사이드카(sidecar) 프록시가 \code{localhost}로 모델 서버에 도달하고, "pod당 IP 하나"가 가능하며, pod 안의 컨테이너 둘은 포트가 충돌할 \textit{수 있지만} pod 둘은 충돌하지 않는다.

\subsection*{USER: 내부의 root는 외부의 root가 아니다}

USER 네임스페이스는 명시적인 테이블 \code{/\allowbreak proc/\allowbreak <pid>/\allowbreak uid\_map}으로 사용자 ID를 다시 매핑한다. \code{0 100000 65536}이라는 매핑은 "이 네임스페이스 내부의 uid 0이 외부의 uid 100000에 대응하며 65536개 ID 범위에 적용된다"는 뜻이다. 따라서 프로세스는 내부에서 비공개 파일 시스템에 패키지를 설치하고 비공개 80번 포트에 바인딩하는 \textit{root}이면서, 호스트 커널에는 특권 없는 uid 100000일 수 있다. 내부에서 파일을 만들면 호스트에서는 특별할 것 없는 100000이 소유한다. 이것이 \textbf{루트리스 컨테이너(rootless container)}(Podman의 대표 기능, Docker의 루트리스 모드)의 기반이다. 런타임까지 특권 없이 실행하여 컨테이너 탈출의 결과를 크게 줄인다. 루트리스 컨테이너를 탈출한 공격자는 호스트에서 아무 권한 없는 사용자로 나온다. (완화책일 뿐 제2장의 microVM 경계는 아니다. 공유 커널의 시스템 호출 표면이 남아 있다.)

\subsection*{UTS와 IPC: 짧지만 가볍지 않게}

UTS(자체 호스트 이름)는 가장 작지만 가장 많이 가르쳐 주는 네임스페이스다. 내부 호스트 이름을 바꾸면서 호스트가 그대로인 모습을 보는 것은 관점이 실제로 갈라졌음을 입증하는 가장 빠른 방법이며, 그래서 실습 2단계에 둔다. IPC는 System V와 POSIX 공유 메모리를 비공개로 만든다. 호랑이 담배 피우던 시절의 각주처럼 들리지만 \textbf{PyTorch DataLoader 워커가 `/dev/shm`에 마운트된 공유 메모리를 통해 프로세스 사이에서 텐서를 전달한다}는 사실을 떠올리면 달라진다. Docker의 기본 \code{/\allowbreak dev/\allowbreak shm}은 64 MB다. 워커가 여러 개인 DataLoader가 실제 배치를 처음 옮기면 이 상한에 부딪혀, 아무 도움도 안 되기로 악명 높은 \code{ERROR: Unexpected bus error encountered in worker. This might be caused by insufficient shared memory (shm)}와 함께 죽는다. 모든 ML 엔지니어가 결국 외우는 해결책은 \code{--shm-size=8g}(또는 K8s의 동등한 볼륨)이다. 이제 그 플래그가 어느 네임스페이스를 향한 것인지 안다.

\begin{calloutbox}{tbbpurple}{tbbpurplefill}
{\normalsize\bfseries\color{tbbpurple}생각해 보기}\par\vspace{3pt}
\begin{tbbitemize}
  \item Kubernetes pod 하나에 컨테이너 두 개가 있다. 그림 3-2의 여섯 축을 사용해 어느 네임스페이스를 공유하고 어느 것을 비공개로 유지하는지 생각해 보자. (localhost로 통신하지만 파일 시스템은 다르다는 점에서 추론할 수 있다.)
  \item 로컬 테스트(python serve.py, Ctrl-C 작동)에서는 서버가 SIGTERM을 잘 처리하지만 컨테이너에서는 무시한다. 이 절을 이용해 가장 가능성 높은 원인 두 가지와 각각의 한 줄 해결책을 제시하라.
  \item USER 네임스페이스 안에서 uid 0으로 실행하는 것이 네임스페이스 없이 uid 0으로 실행하는 것보다 훨씬 안전하지만, 제2장의 microVM 경계만큼 강하지는 않은 이유는 무엇인가? 남아 있는 공유 표면의 이름을 쓰라.
\end{tbbitemize}
\end{calloutbox}

\section{cgroups v2: 머신의 제한된 몫}

네임스페이스는 프로세스에 보이는 것을 바꾸지만 소비할 수 있는 양은 제한하지 않는다. 아름답게 비공개인 PID 트리를 가진 컨테이너화된 모델 서버도 노드의 모든 코어와 RAM을 마지막 바이트까지 먹어 치우고 이웃 레플리카 19개를 함께 쓰러뜨릴 수 있다. 두 번째 재료인 \textbf{컨트롤 그룹, cgroups v2}는 예산 시스템이다. 프로세스를 그룹에 넣고 그룹별 \textbf{컨트롤러}가 CPU, 메모리, I/O, 프로세스 수를 계량하고 제한한다. 네임스페이스가 "어떤 세계가 보이는가?"에 답한다면 cgroup은 "머신을 얼마만큼 사용할 수 있는가?"에 답한다(그림 3-4).

\begin{tbbfigure}
  \adjustbox{max width=0.94\linewidth,max totalheight=0.46\textheight}{\includegraphics{figures/fig3-4_cgroups.png}}\\[3pt]
  \figcaption{그림 3-4  cgroups v2는 파일 시스템으로 공개된 단일 계층이다. cpu.max는 뭉텅이로 스로틀링해 p99를 위협하며, memory.max는 OOM 킬러가 강제하는 단단한 벽이다.}
\end{tbbfigure}

\subsection*{인터페이스는 파일 시스템이다 — 라이브러리가 필요 없다}

cgroups v2는 \code{/\allowbreak sys/\allowbreak fs/\allowbreak cgroup} 아래에 통합 트리 하나를 제공하며, 전체 API가 파일 작업이다. 완전한 예산을 만들고 적용한 뒤 다시 읽는 과정을 살펴보자. 실습 5단계의 예고이기도 하다.

\begin{terminalbox}
\begin{Verbatim}[breaklines=true,breakanywhere=true,fontsize=\small,formatcom=\color{tbbtermfg},codes=\tbbcodeactive]
$ cd /sys/fs/cgroup && sudo mkdir demo        # a new group = a new directory
$ echo "50000 100000" | sudo tee demo/cpu.max # 50ms CPU per 100ms window (0.5 CPU)
$ echo $((256*1024*1024)) | sudo tee demo/memory.max   # 256 MiB hard wall
$ echo $$ | sudo tee demo/cgroup.procs        # enroll THIS shell (children follow)
$ cat demo/memory.current                     # live usage, in bytes
3862528
$ cat demo/cpu.stat | head -4
usage_usec 184220
user_usec  132009
system_usec 52211
nr_throttled 0                                # (watch this one in a moment)
\end{Verbatim}
\end{terminalbox}
\figcaption{기록 3-4  다섯 번의 파일 작업으로 만든 완전한 cgroup 예산. mkdir, 한도를 쓰는 echo 두 번, 등록하는 echo 한 번, 관찰하는 cat이다. Docker와 kubelet이 사용자를 대신해 정확히 이 일을 한다.}

\code{mkdir}, \code{echo}, \code{cat}이 메커니즘의 전부다. 대필자를 알아보자(아래 표 3-2). Docker 플래그나 Kubernetes 한도는 결국 프로그램이 이 파일에 숫자를 쓰는 것이다. v2의 단일 트리는 v1의 컨트롤러별로 뒤얽힌 계층을 대체했다. 현대 배포판과 Kubernetes는 v2를 기본으로 사용하며 이 책도 줄곧 v2를 다룬다.

\begin{tbbtablewrap}
  \begin{tabular}{@{}>{\raggedright\arraybackslash}m{\dimexpr 0.2447\tbltextwidth-2\tabcolsep\relax} >{\raggedright\arraybackslash}m{\dimexpr 0.4574\tbltextwidth-2\tabcolsep\relax} >{\raggedright\arraybackslash}m{\dimexpr 0.2979\tbltextwidth-2\tabcolsep\relax}@{}}
  \arrayrulecolor{tbbnavy}\specialrule{1pt}{0pt}{0pt}
  \rowcolor{tbbtablehead}{\centering \bfseries\color{tbbnavy}파일} & {\centering \bfseries\color{tbbnavy}의미} & {\centering \bfseries\color{tbbnavy}일반적인 작성자} \\
  \arrayrulecolor{tbbnavy}\specialrule{0.7pt}{0pt}{2pt}
  {\textbf{cgroup.procs}} & {이 그룹에 등록한 PID(프로세스를 옮기려면 씀)} & {런타임 / 사용자(실습)} \\
  \arrayrulecolor{tbbtableborder}\hline
  {\textbf{cpu.max}} & {하드 할당량: µs 단위 "<quota> <period>", 예: "50000 100000" = 0.5 CPU} & {docker --cpus, kubelet(한도)} \\
  \arrayrulecolor{tbbtableborder}\hline
  {\textbf{cpu.weight}} & {경합 시 비례 몫(1–10000), 소프트하고 남는 자원을 활용함} & {kubelet(요청량에서 계산)} \\
  \arrayrulecolor{tbbtableborder}\hline
  {\textbf{cpu.stat}} & {사용량 + nr\_throttled / throttled\_usec — 스로틀링(throttling) 진단} & {(읽기 전용 텔레메트리)} \\
  \arrayrulecolor{tbbtableborder}\hline
  {\textbf{memory.max}} & {단단한 메모리 벽(memory wall), 회수 실패 후 초과 → OOM 강제 종료} & {docker -m, kubelet(한도)} \\
  \arrayrulecolor{tbbtableborder}\hline
  {\textbf{memory.high}} & {소프트 상한: 스로틀링 + 회수, 강제 종료 없음 — 경보선(§3.6)} & {사용자 / 플랫폼 정책} \\
  \arrayrulecolor{tbbtableborder}\hline
  {\textbf{memory.current}} & {그룹의 실시간 사용량 — 실습에서 관찰} & {(읽기 전용 텔레메트리)} \\
  \arrayrulecolor{tbbtableborder}\hline
  {\textbf{memory.events}} & {oom\_kill을 포함한 카운터(counter) — OOM 증거 장부} & {(읽기 전용 텔레메트리)} \\
  \arrayrulecolor{tbbtableborder}\hline
  {\textbf{pids.max}} & {최대 프로세스/스레드 수 — 포크 폭탄 울타리} & {docker --pids-limit} \\
  \arrayrulecolor{tbbnavy}\specialrule{1pt}{2pt}{0pt}
  \end{tabular}\par
  \figcaption{표 3-1  실제로 다룰 cgroup v2 파일}
\end{tbbtablewrap}

\subsection*{CPU: 할당량은 뭉텅이로 스로틀링한다 — 지연 시간 계산 예시}

CPU 컨트롤러에는 성격이 \textit{매우} 다른 조절 장치 두 개가 있다. \code{cpu.weight}는 \textbf{경합할 때만} CPU를 비례 배분한다. 바쁜 4코어 머신에서 가중치 300 대 100이면 대략 코어 3개 대 1개를 얻지만, 이웃이 쉬면 어느 그룹이든 전부 사용할 수 있다. 노는 자원을 놀리지 않는 방식이며, 이를 강제하느라 사이클을 낭비하는 일도 없다. \code{cpu.max}는 주기가 있는 하드 할당량이다. \code{"50000 100000"}은 이웃이 쉬든 말든 \textit{100 ms 창마다 이 그룹이 50 ms 동안 실행될 수 있다}는 뜻이다. 강제 방식이 중요한 미묘함이다. 부드럽게 2× 느려지는 것이 아니라 \textbf{스로틀링}한다. 그룹은 창의 할당량을 소진할 때까지 제 속도 그대로 내달리다가, 다음 창이 열릴 때까지 멈춘다.

산술로 체감해 보자. 추론 요청 하나에 \textbf{CPU 시간 30 ms}가 필요하고 컨테이너 할당량이 \code{"20000 100000"}이라고 하자. 실제 클러스터가 "소형" 서비스에 설정하는 0.2 CPU다. 최선의 경우에는 새 창이 시작될 때 요청이 도착한다. 20 ms 동안 실행하고 할당량에 도달해 \textbf{80 ms 동안 멈춘} 뒤 다음 창에서 마지막 10 ms를 실행한다. 완료까지 \textbf{110 ms}가 걸리며 그중 80 ms는 강제 일시 중지다. 최악의 경우에는 앞선 요청이 할당량을 막 소진한 뒤 도착한다. 실행을 시작하기도 전에 최대 80 ms를 더 기다려 CPU 작업 30 ms에 \textbf{190 ms} 가까이 걸린다. 한편 부하가 적은 p50 요청은 창 하나 안에 완전히 들어가 아무 영향도 받지 않을 수 있다. 특유의 징후는 \textbf{p50은 멀쩡한데 p99는 창 잔여 시간의 대략 배수 단위로 뾰족하게 튀고, 애플리케이션 프로파일러에는 아무것도 보이지 않는다}는 것이다. 프로세스가 실행되지 않았으므로 프로파일링할 것도 없다.

\begin{terminalbox}
\begin{Verbatim}[breaklines=true,breakanywhere=true,fontsize=\small,formatcom=\color{tbbtermfg},codes=\tbbcodeactive]
# after driving load through the throttled cgroup from Transcript 3-4:
$ cat demo/cpu.stat
usage_usec     9184220
nr_periods     412
nr_throttled   118        # windows in which we hit the quota
throttled_usec 8221340    # 8.2 s of enforced pause, total
\end{Verbatim}
\end{terminalbox}
\figcaption{기록 3-5  스로틀링 진단. 부하가 걸릴 때 nr\_throttled가 증가하면 p99 급증은 코드가 아니라 여기서 만들어진다.}

표준 해결책은 메커니즘에서 바로 나온다. 할당량을 높이거나, 지연 시간이 중요한 pod에서는 하드 상한보다 \textbf{가중치}(요청량)를 선호해 유휴 용량을 계속 사용할 수 있게 한다. 또는 플랫폼이 무엇을 설정하는지 이해하고 실질적인 일시 중지를 줄인다. 5주 차에는 이 절을 근거로 "K8s CPU 한도를 설정해야 하는가?"라는 논쟁에 돌아온다. 이 이야기는 제2장의 스틸 타임과 나란히 기억해 두자. \textbf{서로 다른 두 메커니즘이 주는 교훈은 하나다. 코드 아래 계층이 꼬리 지연 시간을 만들어 낼 수 있다.}

\subsection*{메모리: 집행자가 지키는 단단한 벽 — 메모리 회수의 미묘함}

메모리 컨트롤러의 \code{memory.max}는 진정한 하드 한도지만, 벽에 이르는 길에는 순진한 그림이 놓치는 중간 단계가 하나 있다. \textbf{메모리 회수}다. 그룹 사용량이 한도에 도달하면 커널은 먼저 다시 만들 수 있는 메모리를 해제하려 한다. 대표적으로 RAM에 캐시한 파일 내용인 \textbf{페이지 캐시(page cache)}가 있다. 대시보드를 읽을 때 중요하다. \code{memory.current}는 페이지 캐시를 포함하므로, 대형 파일을 스트리밍하는 컨테이너(예: 모델 가중치 로딩)가 몇 시간 동안 한도에 붙어 있는 \textit{것처럼 보여도} 완벽하게 건강할 수 있다. 캐시는 회수 가능한 짐이지 필수 메모리가 아니다. 회수할 수 없는 것은 Python 힙과 텐서 같은 \textbf{익명 메모리}다. 벽에 닿은 사용량이 익명 메모리라면 회수에 실패하고, 커널은 \textbf{이 cgroup 범위의 OOM 킬러}를 호출한다. 그룹에서 가장 큰 프로세스(\code{oom\_score\_adj}로 조정 가능)를 골라 항소할 수도 잡을 수도 없고 피해자에게 보이지도 않는 SIGKILL을 보낸다. 애플리케이션 로그에는 아무것도 없고 셸에는 \textbf{종료 코드 137}(128 + 9)이 보이며 Kubernetes는 \code{OOMKilled}를 찍는다. 더 온화한 형제인 \code{memory.high}는 죽이지 \textit{않고} 스로틀링하고 공격적으로 회수한다. 3.6절에서 설정할 경보선이며, 여기서부터의 자세한 이야기는 그 절이 도맡는다.

\subsection*{cgroup이 심판 볼 수 있는 것과 없는 것}

나머지 컨트롤러도 살펴보자. \code{io}는 블록 장치 대역폭과 IOPS를 계량하며, 가중치 로딩이 공유 디스크를 두들길 때 중요하다. \code{pids}는 프로세스 수를 제한하는 Liz Rice 시연의 포크 폭탄 울타리다. \code{cpuset}은 그룹을 특정 코어에 고정하며, 제2장 전용 코어 인스턴스의 컨테이너 사촌이다. 이 책의 방향을 결정하므로 없는 것에도 주목해야 한다. \textbf{cgroup에는 GPU 컨트롤러가 없다.} 커널은 \code{/\allowbreak dev/\allowbreak nvidia*} 장치에 대한 \textit{접근}을 통제할 수 있으며, 권한 수준에서 \code{--gpus}가 작동하는 방식이다. 4주 차에 툴킷의 내부 배관을 들여다본다. 그러나 cgroup 트리에는 \code{cpu.max}와 \code{memory.max}처럼 GPU 계산 시간을 나누거나 VRAM을 분할하는 것이 없다. 구체적으로 GPU 하나에 프로세스 둘을 놓으면 어느 쪽이든 VRAM 전체를 \code{cudaMalloc}하거나 SM을 독점할 수 있다. 이 부재 때문에 GPU 공유에는 제2장의 공급자 메커니즘인 시간 분할, MPS, MIG가 필요하며 12주 차 GPU 공유 도구 상자는 이 절과 전혀 다르게 생겼다.

\begin{tbbtablewrap}
  \begin{tabular}{@{}>{\raggedright\arraybackslash}m{\dimexpr 0.3511\tbltextwidth-2\tabcolsep\relax} >{\raggedright\arraybackslash}m{\dimexpr 0.3298\tbltextwidth-2\tabcolsep\relax} >{\raggedright\arraybackslash}m{\dimexpr 0.3191\tbltextwidth-2\tabcolsep\relax}@{}}
  \arrayrulecolor{tbbnavy}\specialrule{1pt}{0pt}{0pt}
  \rowcolor{tbbtablehead}{\centering \bfseries\color{tbbnavy}사용자가 쓰는 것} & {\centering \bfseries\color{tbbnavy}cgroup v2에 기록되는 것} & {\centering \bfseries\color{tbbnavy}비고} \\
  \arrayrulecolor{tbbnavy}\specialrule{0.7pt}{0pt}{2pt}
  {docker run -m 512m} & {memory.max = 536870912} & {하드 벽, §3.6의 주제} \\
  \arrayrulecolor{tbbtableborder}\hline
  {docker run --cpus 1.5} & {cpu.max = "150000 100000"} & {부드럽지 않고 뭉텅이로 스로틀링} \\
  \arrayrulecolor{tbbtableborder}\hline
  {docker run --pids-limit 256} & {pids.max = 256} & {포크 폭탄 울타리} \\
  \arrayrulecolor{tbbtableborder}\hline
  {K8s resources.requests (5주 차)} & {cpu.weight (+ 스케줄러 계산)} & {계획 + 소프트 몫} \\
  \arrayrulecolor{tbbtableborder}\hline
  {K8s resources.limits (5주 차)} & {cpu.max, memory.max} & {정확히 이 절의 벽} \\
  \arrayrulecolor{tbbtableborder}\hline
  {docker run --gpus all (4주 차)} & {— cgroup 컨트롤러 없음 —} & {장치 접근만 가능, 공유에는 제2장의 vGPU/MIG 필요(12주 차)} \\
  \arrayrulecolor{tbbnavy}\specialrule{1pt}{2pt}{0pt}
  \end{tabular}\par
  \figcaption{표 3-2  사용하는 도구가 이 파일을 대필하는 방식}
\end{tbbtablewrap}

\begin{calloutbox}{tbbpurple}{tbbpurplefill}
{\normalsize\bfseries\color{tbbpurple}생각해 보기}\par\vspace{3pt}
\begin{tbbitemize}
  \item 지연 시간에 민감한 추론 pod에 cpu.max = "100000 100000"(1 CPU)을 설정했다. 부하가 걸리면 p50은 멀쩡하지만 p99가 40–90 ms씩 급증하고 nr\_throttled가 증가한다. 계산 예시의 산술로 급증 크기를 설명하고, 절충이 서로 다른 해결책 두 가지를 제시하라.
  \item 대시보드에서 가중치를 로드하는 컨테이너가 한 시간 동안 memory.max의 98\%에 붙어 있지만 OOM 강제 종료는 발생하지 않는다. 메모리 회수의 미묘함을 이용해 이유를 설명하고, 실제 위험에 얼마나 가까운지 알려 줄 파일의 이름을 쓰라.
  \item "메모리 한도 = 관측한 유휴 메모리 + 10\%" 규칙이 많은 무상태 웹 앱에는 잘 통하면서 추론 서버에서는 OOM 강제 종료를 부르는 이유는 무엇인가? (제1장에서 이미 답할 수 있으며 §3.6에서 수치로 확인한다.)
\end{tbbitemize}
\end{calloutbox}

\section{overlayfs: 계층으로 이루어진 파일 시스템}

세 번째 재료는 "컨테이너의 세계가 어디에서 오는가?"에 답한다. 컨테이너에는 자체 \code{/\allowbreak usr/\allowbreak bin/\allowbreak python}과 자체 라이브러리가 있는 비공개 루트 파일 시스템이 필요하다. (a) 아티팩트로 운반할 수 있고, (b) 컨테이너마다 전체 사본 비용이 들지 않으며, (c) 밀리초 만에 나타나야 한다. 세 조건의 기반 메커니즘은 디렉터리를 \textit{쌓는} 유니온 파일 시스템 \textbf{overlayfs}다. 읽기 전용 \textbf{lower} 디렉터리 여러 개(이미지 계층)와 쓰기 가능한 \textbf{upper} 디렉터리 하나(컨테이너의 임시 공간)를 겹쳐, 컨테이너가 \code{/\allowbreak }로 마운트하는 하나의 \textbf{merged} 관점으로 제공한다(그림 3-5).

\begin{tbbfigure}
  \adjustbox{max width=0.94\linewidth,max totalheight=0.46\textheight}{\includegraphics{figures/fig3-5_overlayfs.png}}\\[3pt]
  \figcaption{그림 3-5  overlayfs: 읽기 전용 이미지 계층 + 얇은 쓰기 가능 계층 하나 = 컨테이너의 루트. 쓰기는 copy-up하고 삭제는 whiteout을 남기며 계층은 깨끗한 공유 상태를 유지한다.}
\end{tbbfigure}

\subsection*{90초 만에 시연하는 세 가지 규칙}

overlayfs는 모든 기본 Linux 커널에 있는 마운트 유형이므로 Docker 없이 임시 디렉터리에서 전체 과정을 실험할 수 있다. 아래 기록은 실제로 실행해 볼 가치가 있다. 각 명령이 세 규칙 중 하나를 가르친다.

\begin{terminalbox}
\begin{Verbatim}[breaklines=true,breakanywhere=true,fontsize=\small,formatcom=\color{tbbtermfg},codes=\tbbcodeactive]
$ mkdir lower upper work merged
$ echo "from the image"  > lower/config.txt
$ echo "also the image"  > lower/keep.txt
$ sudo mount -t overlay overlay \
    -o lowerdir=lower,upperdir=upper,workdir=work merged
$ cat merged/config.txt                # RULE 1 - reads fall through
from the image
$ echo "changed!" > merged/config.txt  # RULE 2 - writes copy-up
$ cat lower/config.txt                 #   the layer is untouched...
from the image
$ ls upper/                            #   ...the copy lives in upper
config.txt
$ rm merged/keep.txt                   # RULE 3 - deletes leave whiteouts
$ ls -l upper/keep.txt
c--------- 1 root root 0, 0 ... upper/keep.txt   # a marker, not a file
$ ls lower/                            #   the bytes still exist below!
config.txt  keep.txt
\end{Verbatim}
\end{terminalbox}
\figcaption{기록 3-6  한 번에 보는 overlayfs. 읽기는 스택을 통과하고, 쓰기는 파일을 위로 복사한 뒤 수정하며, 삭제는 whiteout 표식으로 파일을 숨긴다. lower 계층은 깨끗하고 공유 가능한 상태를 유지한다.}

세 규칙의 결과를 분명히 해 보자. \textbf{읽기는 위에서 아래로 찾으므로} merged 관점에는 각 경로의 가장 위 버전이 보인다. 뒤 계층이 앞 계층을 수정하지 않고 파일을 "갱신"하는 방식이다. \textbf{쓰기는 copy-up을 일으킨다.} 먼저 파일 전체를 upper 디렉터리로 복사하고 그 사본을 바꾼다. 컨테이너마다 자체 upper가 있으므로 컨테이너 30개가 같은 기본 이미지의 모든 읽기 전용 바이트를 공유하면서도 그 이미지를 각각 "수정"할 수 있다. \textbf{삭제는 whiteout을 남긴다.} 문자 장치 표식(\code{c---------}, 장치 0:0)이 merged 관점에서 \code{keep.txt}를 숨기지만 바이트는 lower 계층에 남는다. whiteout 규칙에는 유명한 비용상 결과가 있다. \textbf{뒤 Dockerfile 계층의 RUN rm은 이미지를 줄이는 것이 아니라 늘린다.} "삭제한" 바이트가 앞 계층에 실려 가고 표식까지 더해진다. 파일을 만든 \textit{바로 그 계층에서} 삭제해야 하며, Dockerfile에서는 같은 \code{RUN} 줄을 뜻한다.

\begin{terminalbox}
\begin{Verbatim}[breaklines=true,breakanywhere=true,fontsize=\small,formatcom=\color{tbbtermfg},codes=\tbbcodeactive]
# WRONG - image contains 900MB of downloads AND a whiteout hiding them:
RUN apt-get install -y build-essential
RUN rm -rf /var/lib/apt/lists /tmp/downloads

# RIGHT - the layer never contains the junk at all:
RUN apt-get install -y build-essential && \
    rm -rf /var/lib/apt/lists /tmp/downloads
\end{Verbatim}
\end{terminalbox}
\figcaption{기록 3-7  Dockerfile 실무의 whiteout 규칙. 정리 작업은 치울 대상을 만든 RUN과 같아야 하며, 그렇지 않으면 바이트가 그대로 운반된다.}

\subsection*{계층이 경제적 설계인 이유 — 서빙 이미지를 장부로 읽기}

4주 차 실습에서 실제로 만들 이미지를 이용해 경제학자의 눈으로 스택을 읽어 보자. 서빙 이미지는 \textit{변경 속도}에 따라 자연스럽게 나뉘며, 순서를 잘 잡은 Dockerfile은 가장 천천히 바뀌는 내용을 가장 아래에 둔다.

\begin{tbbtablewrap}
  \begin{tabular}{@{}>{\raggedright\arraybackslash}m{\dimexpr 0.3617\tbltextwidth-2\tabcolsep\relax} >{\raggedright\arraybackslash}m{\dimexpr 0.1809\tbltextwidth-2\tabcolsep\relax} >{\raggedright\arraybackslash}m{\dimexpr 0.2234\tbltextwidth-2\tabcolsep\relax} >{\raggedright\arraybackslash}m{\dimexpr 0.2340\tbltextwidth-2\tabcolsep\relax}@{}}
  \arrayrulecolor{tbbnavy}\specialrule{1pt}{0pt}{0pt}
  \rowcolor{tbbtablehead}{\centering \bfseries\color{tbbnavy}계층(Dockerfile 단계)} & {\centering \bfseries\color{tbbnavy}일반적 크기} & {\centering \bfseries\color{tbbnavy}변경 주기} & {\centering \bfseries\color{tbbnavy}재빌드/가져오기 비용} \\
  \arrayrulecolor{tbbnavy}\specialrule{0.7pt}{0pt}{2pt}
  {FROM ubuntu / python-slim 기본 이미지} & {\textasciitilde{}80–150 MB} & {매년} & {거의 영구 캐시} \\
  \arrayrulecolor{tbbtableborder}\hline
  {CUDA 런타임 + cuDNN 라이브러리} & {\textasciitilde{}2–3 GB} & {CUDA 업그레이드마다} & {수개월 캐시} \\
  \arrayrulecolor{tbbtableborder}\hline
  {pip install torch + 종속성} & {\textasciitilde{}2–3 GB} & {requirements.txt 변경 시} & {수주 캐시} \\
  \arrayrulecolor{tbbtableborder}\hline
  {\textbf{COPY serve.py /app/}} & {\textbf{\textasciitilde{}40 KB}} & {\textbf{배포마다}} & {\textbf{수초}} \\
  \arrayrulecolor{tbbnavy}\specialrule{1pt}{2pt}{0pt}
  \end{tabular}\par
  \figcaption{표 3-3  계층별 모델 서버 이미지(대표 크기)}
\end{tbbtablewrap}

이 순서 덕분에 이 책 내내 두고두고 덕을 보는 세 가지가 생긴다. \textbf{중복 제거:} 노드 하나에 \textasciitilde{}5–6 GB 이미지의 레플리카 30개를 두어도 계층은 \textit{한 번만} 저장하고 거의 비어 있는 upper 디렉터리 30개만 더한다. ML 이미지로 컨테이너 밀도를 확보할 수 있는 이유다. \textbf{캐시:} \code{serve.py}를 편집하고 재빌드하면 아래의 모든 계층을 재사용한다. 재빌드는 수초, 푸시는 수 킬로바이트에 그친다. 반면 \code{COPY . /\allowbreak app}이 \code{pip install} \textit{위에} 있는 초보자의 고전적인 순서 역전에서는 모든 코드 편집이 종속성 계층을 무효화하고 torch를 다시 빌드하므로 20분과 수 기가바이트가 드는 재앙이 된다. \textbf{시작:} 컨테이너 파일 시스템의 "시작"은 이미 존재하는 계층을 overlay로 한 번 마운트하는 일이므로 밀리초밖에 걸리지 않는다. 파일 시스템 자체는 콜드 스타트 문제가 아니다. \textit{없는 계층을 네트워크로 가져오는 일이 문제다.} 12주 차를 미리 보는 계산을 해 보자. 새 노드에 6 GB 이미지를 유효 100–250 MB/s로 가져오면 레지스트리, 네트워크, 압축 해제가 모두 파이프에 부담을 주어 \textbf{코드 한 줄 실행 전에 콜드 스타트 25–60초}가 걸린다. 대개 단일 항목으로는 최대이며, 계층 위생이 정리 정돈이 아니라 지연 시간 엔지니어링인 이유다.

서빙에 특화된 한 가지 주의 사항으로 그림을 완성하자. \textbf{모델 가중치는 보통 이미지 계층에 넣지 않는다.} 가중치는 거대하고 코드와 다른 주기로 바뀌며, 제1장에서 배웠듯 객체 스토리지의 모델 저장소에 두고 시작할 때 로드하는 편이 맞다. 6주 차 실습에서 바로 그 경로를 만들며, 13주 차 레지스트리/버전 관리 이야기도 이에 의존한다. 30 GB 체크포인트를 계층에 구우면 모델 버전과 이미지 버전이 결합되고, 모든 노드의 디스크가 부풀며, 배포할 때마다 바뀌지 않은 수 기가바이트를 다시 보낸다. 불변 데모 이미지나 엣지 장비에서는 때때로 옳지만, 기본값이 아니라 의식적인 결정이어야 한다.

\begin{calloutbox}{tbbpurple}{tbbpurplefill}
{\normalsize\bfseries\color{tbbpurple}생각해 보기}\par\vspace{3pt}
\begin{tbbitemize}
  \item 팀원의 Dockerfile이 FROM python-slim → COPY . /app → pip install -r /app/requirements.txt → CMD ... 순서다. (a) serve.py 변경, (b) requirements.txt 변경 시 재빌드 동작을 예측하라. 올바르게 순서를 바꾸고 각 경우의 새 비용을 설명하라.
  \item 기록 3-6의 규칙을 이용해 보자. 같은 이미지의 컨테이너 두 개가 내부의 /var/log/app.log에 각각 내용을 덧붙인다. 서로의 줄이 보이는가? 두 파일은 물리적으로 어디에 존재하는가? 컨테이너를 제거하면 둘은 어떻게 되는가?
  \item 팀이 serve.py (40 KB)만 바뀌는 5.5 GB 이미지를 하루 6번 배포한다고 하자. 올바르게 계층화하면 배포마다 몇 바이트가 이동하는가? COPY가 pip 위에 있는 순서 역전에서는 어떠한가? (레지스트리 송신에도 요금이 부과되므로 제1장의 비용 관점을 적용한다.)
\end{tbbitemize}
\end{calloutbox}

\section{직접 만들기 — Docker가 실제로 더하는 것}

맨손으로 레시피를 조립할 시간이다. 아래 순서는 이번 주 실습의 지적 핵심이며 Liz Rice의 유명한 "Containers from Scratch" 시연 구조를 따른다. Rice는 Go로 라이브 코딩하지만, 우리는 신비를 더 벗겨 주는 기본 셸 명령으로 실행한다. 이 장의 각 절을 활용하는 다섯 단계다(그림 3-6). 순서대로 실행하고 모든 출력을 읽자.

\begin{tbbfigure}
  \adjustbox{max width=0.94\linewidth,max totalheight=0.46\textheight}{\includegraphics{figures/fig3-6_byhand.png}}\\[3pt]
  \figcaption{그림 3-6  다섯 단계로 만드는 컨테이너. 네임스페이스를 unshare하고, 신원을 얻고, 비공개 루트로 pivot하며, /proc를 마운트하고, cgroup에 등록한다. Docker가 더하는 것은 편의성과 표준이지 커널 마법이 아니다.}
\end{tbbfigure}

\subsection*{1\textasciitilde{}2단계: 새로운 관점과 비공개임을 보여 주는 증거}

\begin{terminalbox}
\begin{Verbatim}[breaklines=true,breakanywhere=true,fontsize=\small,formatcom=\color{tbbtermfg},codes=\tbbcodeactive]
$ sudo unshare --uts --pid --mount --net --fork bash   # MOVE 1
# (prompt looks identical - but check your namespace links vs Transcript 3-2)
$ hostname minibox                                     # MOVE 2
$ hostname
minibox
# ...meanwhile, in a second terminal on the host:
$ hostname
gpu-node-17                    # untouched. The view really is private.
\end{Verbatim}
\end{terminalbox}
\figcaption{기록 3-8  1\textasciitilde{}2단계. PID 소속은 자식에게 적용되므로 --fork가 중요하다. 터미널 두 개의 호스트 이름 대비가 첫 실습 아티팩트다.}

\subsection*{3단계: 비공개 루트(pivot\_root를 쓰고 chroot를 쓰지 않는 이유)}

alpine이나 ubuntu-base tar 아카이브 같은 소형 rootfs를 풀어 \code{/\allowbreak }로 만든다. rootfs는 문자 그대로 이미지 계층을 디렉터리 하나로 평탄화한 것이다. 시연에서는 \code{chroot}라는 지름길을 쓰며 실습에는 충분하다. 실제 런타임은 \code{pivot\_root}를 사용한다. 둘의 차이는 보안 소양을 위해 한 문장으로 정리할 가치가 있다. \code{chroot}는 프로세스가 생각하는 \code{/\allowbreak }의 위치만 \textit{다시 가리키고} 기존 루트는 마운트 네임스페이스에 남기므로 잘 알려진 탈출 묘수로 다시 접근할 수 있다. 감옥 밖 파일 디스크립터를 가진 프로세스는 되돌아 나갈 수 있다. 반면 \code{pivot\_root}는 마운트 네임스페이스의 루트를 \textit{교체}하며 기존 세계를 완전히 언마운트할 수 있다. 다시 가리키기냐 통째로 갈아 끼우기냐 — 고전적인 컨테이너 보안 상식 하나가 이제 당신 것이 되었다.

\begin{terminalbox}
\begin{Verbatim}[breaklines=true,breakanywhere=true,fontsize=\small,formatcom=\color{tbbtermfg},codes=\tbbcodeactive]
$ curl -LO https://.../alpine-minirootfs-3.20-x86_64.tar.gz   # ~3 MB(!)
$ mkdir rootfs && tar -xzf alpine-minirootfs-*.tar.gz -C rootfs
$ ls rootfs
bin  dev  etc  home  lib  ...  usr  var        # a whole tiny Linux userland
$ chroot rootfs /bin/sh                        # MOVE 3 (lab shortcut)
/ # ls /                                       # your / is now the rootfs
bin  dev  etc  home  lib  ...  usr  var
\end{Verbatim}
\end{terminalbox}
\figcaption{기록 3-9  3단계. rootfs는 그저 디렉터리이고 이미지는 계층으로 운반되는 rootfs다. Alpine은 3 MB이며 6 GB CUDA 이미지를 바라볼 기준을 제공한다.}

\subsection*{4단계: 도구가 새로운 진실을 말하게 하기}

\begin{terminalbox}
\begin{Verbatim}[breaklines=true,breakanywhere=true,fontsize=\small,formatcom=\color{tbbtermfg},codes=\tbbcodeactive]
/ # ps aux                       # BEFORE mounting /proc: confusion
ps: can't open /proc: No such file or directory
/ # mount -t proc proc /proc     # MOVE 4
/ # ps aux
PID   USER   COMMAND
  1   root   /bin/sh             # there it is: you are init
  9   root   ps aux
\end{Verbatim}
\end{terminalbox}
\figcaption{기록 3-10  4단계. ps는 /proc를 읽는다. 네임스페이스를 인식하는 proc를 마운트하면 도구가 프로세스 두 개와 PID 1인 셸이라는 새로운 세계를 보고한다.}

\subsection*{5단계: 예산과 전체 목록}

호스트의 두 번째 터미널에서 기록 3-4와 똑같이 cgroup을 만든다. \code{mkdir /\allowbreak sys/\allowbreak fs/\allowbreak cgroup/\allowbreak minibox}를 실행하고 \code{memory.max}에 \code{268435456}을 쓰며, 셸의 호스트 PID를 \code{cgroup.procs}에 쓴다. 이제 직접 만든 컨테이너에는 256 MiB 하드 벽이 생겼으며, 실습 B부에서 의도적으로 부딪힌다. 이 순간 존재하는 것을 목록으로 정리해 보자. 자체 호스트 이름(UTS), 자신이 init인 자체 PID 트리(PID), 자체 마운트 테이블과 루트 파일 시스템(MNT + rootfs), 아직 비어 있는 자체 네트워크 스택(NET), 단단한 메모리 예산(cgroup)을 가진 프로세스다. \textbf{이것이 컨테이너다.} 아무것도 설치하지 않았고 데몬도 실행되지 않았으며 Docker는 관여하지 않았다. 사용한 도구는 \code{unshare}, \code{chroot}, \code{mount}, \code{mkdir}, \code{echo}, 3 MB tar 아카이브가 전부다.

\subsection*{그렇다면 Docker가 실제로 더하는 것은 무엇인가?}

레시피 주위의 모든 것을 더하지만 커널 안에는 아무것도 더하지 않는다. \textbf{이미지와 레지스트리:} rootfs를 내용 주소화 계층으로 패키징하는 표준 방식이다. 각 계층은 바이트의 SHA-256으로 이름을 붙이므로 계층 캐시를 신뢰할 수 있다. 네트워크로 옮기고 어느 머신에서든 바이트까지 똑같이 재현하여 "내 노트북에서 작동한다"를 "GPU 노드에서도 작동한다"로 바꾼다. \textbf{네트워킹 기본값:} §3.2의 veth/브리지/NAT 배선을 컨테이너마다 만들고 제거하므로 \code{ip link add}를 직접 입력할 필요가 없다. \textbf{수명 주기와 API:} tmux 창에 unshare를 잔뜩 쌓는 대신 재시작 정책과 로그 캡처를 갖춘 데몬 기반 API로 create/start/stop/exec/logs를 제공한다. 그리고 \textbf{런타임 분리:} Docker 자체는 실제 레시피 실행을 아래로 위임한다. Docker → \code{containerd}(수명 주기 관리자) → \code{runc}로 이어진다. runc는 다섯 단계를 엄격하게 수행하는 것이 \textit{유일한 임무}인 작은 독립 실행 바이너리다. 어느 네임스페이스를 unshare할지, 어떤 cgroup 값을 쓸지, 어느 rootfs로 pivot할지를 지정하는 JSON 파일(OCI 런타임 명세)이 구동하며, 이 장을 설정 파일로 옮겨 적은 것처럼 읽힌다. 이 표준화된 바닥 계층 덕분에 제2장의 Kata와 gVisor가 Kubernetes 아래에서 대체재로 끼어들 수 있다. 바로 다음 주의 주제이며, 이 책이 아껴 둔 질문인 GPU를 상자 안에 주입하는 방식도 다룬다.

\begin{calloutbox}{tbbpurple}{tbbpurplefill}
{\normalsize\bfseries\color{tbbpurple}생각해 보기}\par\vspace{3pt}
\begin{tbbitemize}
  \item 직접 만든 컨테이너에서는 인터페이스를 올리기 전까지 localhost에도 ping이 실패한다. 호스트까지 작동하는 네트워킹을 제공하는 데 필요한 최소 추가 단계를 나열해 보자. (§3.2의 패킷 순회에 모든 부품이 있다. 쌍, 주소, up이 필요하다.)
  \item 직접 만든 컨테이너와 docker run alpine sh는 모두 같은 rootfs를 제공한다. Docker 버전에는 있지만 직접 만든 것에는 없는 운영상 중요한 요소 두 가지와 바이트 단위로 동일한 요소 하나를 제시하라.
  \item 해시(sha256)로 계층을 내용 주소화하면 서로 만난 적 없는 머신 사이에서도 계층 캐싱이 안전한 이유는 무엇인가? 부수적으로 어떤 공격을 막는가?
\end{tbbitemize}
\end{calloutbox}

\section{심판이 모델을 죽일 때: OOM 대응 지침}

이 장의 모든 메커니즘은 프로덕션의 한순간으로 모인다. 이 책에서 절 하나를 온전히 할애할 만큼 흔한 상황이다. 추론 서버가 실행되고 트래픽이 흐르다가 프로세스가 사라진다. 스택 추적도 마지막 로그 줄도 없으며 애플리케이션 텔레메트리에는 침묵만 남는다. 어딘가에서 재시작 루프가 시작된다. 이것이 컨테이너화된 모델 서버가 죽는 단연 가장 흔한 방식인 \textbf{OOM 강제 종료}다. 이 절을 마치면 진단하기 가장 쉬운 문제 중 하나가 된다(그림 3-7).

\begin{tbbfigure}
  \adjustbox{max width=0.94\linewidth,max totalheight=0.46\textheight}{\includegraphics{figures/fig3-7_oom.png}}\\[3pt]
  \figcaption{그림 3-7  memory.max 아래에서 살아가는 추론 서버의 메모리 생애. 가중치 계단, 요청별 활성값 급증, 치명적인 경계 초과, 포렌식 흔적(137, dmesg, memory.events)을 보여 준다.}
\end{tbbfigure}

\subsection*{처음부터 끝까지 보는 메커니즘}

3.3절의 규칙을 시간순으로 되짚어 보자. 그룹의 \code{memory.current}가 \code{memory.max}를 향해 올라간다. 커널은 회수할 수 있는 것부터 회수하며 페이지 캐시가 먼저다. 그래서 파일 I/O가 많은 단계는 한도에 붙어도 해롭지 않을 수 있다. 메모리 회수 후에도 새 \textbf{익명} 할당(힙, 텐서)을 충족할 수 없으면 커널이 \textbf{cgroup 범위의 OOM 킬러}를 호출한다. 그룹의 프로세스에 본질적으로 메모리 사용량을 기준으로 점수를 매기고(\code{oom\_score\_adj}로 조정할 수 있다), 가장 큰 것을 고른다. 모델 서버 pod에서는 사실상 항상 모델 서버다. 그리고 \textbf{SIGKILL}을 보낸다. SIGKILL은 잡거나 막거나 처리할 수 없다. 정리도 플러시도 작별 인사도 없다. 따라서 증거는 피해자의 \textit{바깥에만} 남는다. 오전 2시에 필요해지기 전에 각 아티팩트를 눈으로 알아볼 수 있어야 한다.

\begin{terminalbox}
\begin{Verbatim}[breaklines=true,breakanywhere=true,fontsize=\small,formatcom=\color{tbbtermfg},codes=\tbbcodeactive]
$ echo $?                          # 1) the exit code, wherever the process ran
137                                #    = 128 + 9 (killed by signal 9, SIGKILL)
$ sudo dmesg | grep -i -A1 "out of memory"     # 2) the kernel's confession
Memory cgroup out of memory: Killed process 41320 (python)
  total-vm:1892456kB, anon-rss:498112kB, file-rss:1024kB, ...
$ cat /sys/fs/cgroup/.../memory.events          # 3) the cgroup's ledger
oom 3
oom_kill 1
# 4) and the same story, as Kubernetes tells it (Week 5):
$ kubectl describe pod model-server-7d4b9 | grep -A3 "Last State"
Last State:   Terminated
  Reason:     OOMKilled
  Exit Code:  137
\end{Verbatim}
\end{terminalbox}
\figcaption{기록 3-11  완전한 OOM 증거 사슬. 종료 코드 137, 한도에 가까운 회수 불가능 메모리 anon-rss가 보이는 dmesg 줄, oom\_kill 카운터, 그리고 세 가지를 되풀이하는 Kubernetes다.}

\code{dmesg} 줄에는 자세히 읽은 보람이 있는 세부 사항이 하나 있다. 512 MiB 한도에 \code{anon-rss:498112kB}다. 회수할 수 없는 \textit{익명} 메모리가 예산을 채워 kill이 일어났다. 캐시와 비슷한 \code{file-rss}는 메모리 회수가 이미 벗겨 냈기 때문에 작다. 커널이 자신의 자백에 인쇄한 3.3절의 메모리 회수 미묘함이다.

\subsection*{특히 추론 서버가 벽으로 걸어 들어가는 이유}

OOM 킬러가 형을 집행하지만 범죄는 거의 언제나 잘못 산정한 예산이다. 추론 메모리에는 무상태 웹 앱만 다루며 길러진 직관이 잘못 짚는 \textit{형태}가 있기 때문이다. 웹 핸들러는 평평한 기준선 주위에서 요청당 수 킬로바이트를 할당하므로 관측한 사용량 ≈ 필요한 사용량이고 "유휴 + 10\%"도 괜찮은 한도다. 추론 서버는 세 곡선이 겹친다. 시작할 때의 크고 고정된 \textbf{가중치} 계단(예측 가능하며 한 번 측정), 적당하고 일정한 평탄 구간, 그리고 트래픽에 의존해 \textbf{배치 크기와 시퀀스 길이에 비례하여 높아지는 요청별 활성값 급증}이다. 그 위에 할당자 오버헤드도 더해진다. PyTorch 캐싱 할당자는 해제된 블록을 풀에 보관하고 단편화가 더 늘린다. 평탄 구간만으로는 최대치에 관해 거의 아무것도 알 수 없다.

\begin{calloutbox}{tbbblue}{tbbbluefill}
{\normalsize\bfseries\color{tbbblue}계산 예시 — 소형 BERT 서버의 memory.max 산정(실습 규모 수치)}\par\vspace{3pt}
CPU에서 서빙하는 BERT-base 분류기를 생각해 보자. 실습 B부 워크로드와 비슷하다. 가중치는 110 M parameters × 4 bytes (fp32) ≈ \textbf{440 MB}이며 시작할 때 한 번 로드한다. 그림 3-7의 계단이다. 부하 상태에서 측정하면 다음과 같은 대표 수치가 나온다.

\begin{tbbitemize}
  \item 로드 후 유휴 상태: \textbf{\textasciitilde{}620 MB}(가중치 + 런타임 + 할당자 풀)
  \item batch 8 × seq 128: \textbf{\textasciitilde{}900 MB}까지 급증  ·  batch 32 × seq 256: \textbf{\textasciitilde{}1.4 GB}까지 급증
\end{tbbitemize}

정책을 비교해 보자. \textbf{단순 정책:} limit = 유휴 + 10\% ≈ \textbf{680 MiB} → \textit{첫} batch-8 요청에서 죽고 재시작 루프에 들어간다. 재시작할 때마다 440 MB 로드 비용을 다시 낸다. \textbf{최대치 기준:} limit = 측정한 최악 급증 × 1.25(여유) ≈ \textbf{1.75 GiB} → 부하 테스트한 모든 상황에서 살아남지만 레플리카당 유휴 상태보다 \textasciitilde{}1.1 GB를 더 예약한다. 제1장의 밀도와 비용 절충이다. \textbf{경계 설정:} 서버 구성에서 batch ≤ 8, seq ≤ 128로 서빙을 제한하고 limit ≈ 900 MB × 1.25 ≈ \textbf{1.1 GiB}로 둔다. 예약량은 줄고 최악 조건에는 경계가 생기지만 처리량 상한을 대가로 치른다. 지연 시간·처리량·비용 삼각형이 다시 등장한다. 단순하지 않은 두 정책은 모두 방어할 수 있지만, 단순 정책은 트래픽이라는 도화선을 단 시한폭탄이다.
\end{calloutbox}

\vspace{3.0pt}

정직성을 위해 두 가지를 분명히 하자. 첫째, 위 산술은 \textit{방법}이며 실제 수치는 달라진다. 대응 지침의 첫 동사가 추측이 아니라 \textbf{측정}인 이유다. 둘째, GPU 서빙에서는 같은 이야기가 두 번 반복된다. 커널 OOM은 이 절에서 설명한 그대로 호스트 RAM을 관리하지만, \textbf{GPU VRAM에는 다른 심판}이 있다. VRAM이 고갈되면 CUDA 할당자의 \code{torch.cuda.OutOfMemoryError}가 발생한다. SIGKILL이 아니라 할당자 오류이므로 Python에서 \textit{잡을 수 있고}, 해당 요청만 떨궈 낼 수 있다. 종료 코드 137 + dmesg와 Python 트레이스백이라는 징후로 두 OOM을 구분하는 것은 실제로 유용한 진단 기술이다. 11주 차에는 VRAM 쪽에서 계속 커지는 세 번째 등장인물인 KV 캐시를 더한다.

\subsection*{대응 지침}

\begin{tbbitemize}
  \item \textbf{평탄 구간이 아니라 최대치에 맞춰 산정한다.} 구성 가능한 최대 배치와 가장 긴 시퀀스를 사용하는 \textit{최악 조건 부하 테스트}에서 \code{memory.current}를 측정한다. 9주 차 Locust 실행이 정확히 이를 만든다. 명시적인 여유 용량을 더해 그 최대치보다 높게 \code{memory.max}를 설정한다(예시는 1.25× 사용). 유휴 관찰로 한도를 도출해서는 안 된다.
  \item \textbf{가변 부분에 경계를 둔다.} 급증은 배치와 시퀀스 길이에 비례한다. 서버 구성에서 둘을 제한하면 상한 없던 메모리 동작에 상한이 생기고, 앞 단계의 "최악 조건"이 실제로 존재하는 숫자가 된다. 제1장의 SLO 사고를 메모리에 적용한 것이다.
  \item \textbf{memory.high를 경보선으로 사용한다.} max보다 아래에 설정하고(예: 1.75 GiB max 아래 high = 1.5 GiB) \code{memory.events} / 메모리 압력 신호에 경보를 건다. 조용한 처형 대신 회수, 스로틀링, 증가하는 카운터가 있는 경고 단계를 그래프로 얻는다. 14주 차에 대시보드로 연결한다.
  \item \textbf{137을 수수께끼가 아니라 진단으로 읽는다.} 종료 코드 137 + memory.events의 \code{oom\_kill} = 버그가 아니라 예산 문제다. 해결을 논의할 때 나와야 할 말은 "크기를 조정하거나 경계를 둔다"이지 "try/except를 추가한다"가 아니다. SIGKILL은 잡을 수 없고, 계산 예시처럼 치명적인 할당이 자기 코드의 것이 아닐 때도 많다.
  \item \textbf{5주 차보다 먼저 Kubernetes에서 무엇이 무엇에 대응하는지 알아둔다.} requests는 스케줄러의 계획 수치이고 limits는 이 장의 cgroup 벽이다. QoS 클래스는 OOM \textit{우선순위} 정책이다. Guaranteed pod(requests = limits)는 가장 보호적인 \code{oom\_score\_adj}를 받고, \textit{노드}의 메모리가 부족하면 BestEffort pod가 가장 먼저 희생된다. 서빙 클러스터에서는 보통 swap을 끈다. 벽은 그대로 벽이다.
\end{tbbitemize}

\begin{calloutbox}{tbbpurple}{tbbpurplefill}
{\normalsize\bfseries\color{tbbpurple}생각해 보기}\par\vspace{3pt}
\begin{tbbitemize}
  \item 동료가 우아하게 성능을 낮추기 위해 Python에서 OOM을 잡자고 제안한다(try/except MemoryError). 메커니즘과 두 OOM의 차이를 이용해 (a) cgroup 벽에서는 작동할 수 없는 이유, (b) GPU 쪽에서 매우 비슷한 아이디어가 작동하는 곳을 설명하라.
  \item 서버에서 측정한 급증은 batch 32에서 1.9 GB, batch 64에서 3.4 GB다. 노드 용량상 레플리카당 한도는 2.5 GiB다. 계산 예시의 방법으로 구성을 제안하고, 지연 시간·처리량·비용 삼각형의 어느 꼭짓점을 희생했는지 밝혀라.
  \item 기록 3-11에서 dmesg가 대신 file-rss:400000kB와 anon-rss:90000kB를 보여 주었다고 하자. "OOM"에 관해 무엇을 알려 주며 해결책이 완전히 달라야 하는 이유는 무엇인가?
\end{tbbitemize}
\end{calloutbox}

\section{실습: 미니 컨테이너와 관찰한 OOM}

이번 주 실습은 상자를 만들고 그 벽에 부딪히는 이 장의 흐름을 반영한 두 부분으로 이루어진다. 어느 Linux VM이나 이 책의 클라우드 인스턴스에서도 실행할 수 있으며 GPU는 필요 없다. OOM 시연은 의도적으로 CPU만 사용하므로 비싼 인스턴스를 걸 필요가 없다. 안내서에는 정확한 명령이 있고, 이 절에서는 의도·프로토콜·예상 관찰 결과를 정해 실습에서 맹목적으로 입력하지 않고 예측을 \textit{확인}하게 한다. 산출물은 명령 기록의 핵심과 이 장의 절에 각각 연결한 관찰 세 가지를 담은 1쪽 보고서다.

\subsection*{A부 — 다섯 단계 실행하기}

기록 3-8부터 3-10까지 처음부터 끝까지 실행하며 아티팩트 네 가지를 수집한다. 첫째, 터미널 두 개의 \textbf{호스트 이름 대비}로 §3.2의 비공개 관점을 입증한다. 둘째, \code{/\allowbreak proc} 마운트 뒤 \textbf{PID 1 = 자기 셸이라고 표시하는 ps}로 §3.2의 두 가지 진실을 확인한다(마운트 \textit{전}의 \code{ps} 출력도 기록하자 — 그 오류 메시지가 바로 실전에서 만나는 "도구는 /proc를 읽는다"는 교훈이다). 셋째, 내부와 호스트 셸의 \code{ls -l /\allowbreak proc/\allowbreak \$\$/\allowbreak ns}를 비교해 모든 숫자가 다른지 확인하는 \textbf{네임스페이스 inode 비교}다(기록 3-2의 테스트를 직접 만든 컨테이너에 적용한다). 넷째, \code{cat /\allowbreak sys/\allowbreak fs/\allowbreak cgroup/\allowbreak minibox/\allowbreak cgroup.procs}가 컨테이너가 1이라고 믿는 프로세스의 \textit{호스트} PID를 출력하는 \textbf{호스트 쪽 예산 관점}이다. 마지막 아티팩트는 출력 한 줄에 이 장의 명제를 담는다. 같은 프로세스, 두 가지 진실, 하나의 예산이다.

\subsection*{B부 — 처형 관찰하기}

이제 서빙의 교훈을 살펴본다. 예산을 설정한 컨테이너 또는 \code{memory.max} = 256 MiB인 새 cgroup 안에서 안내서의 장난감 추론 스크립트를 실행한다. 의도적으로 큰 모델 객체를 로드한 뒤(가중치 계단) 점점 커지는 배치를 서빙한다(활성값 급증). 두 번째 터미널의 호스트에서 관찰한다.

\begin{terminalbox}
\begin{Verbatim}[breaklines=true,breakanywhere=true,fontsize=\small,formatcom=\color{tbbtermfg},codes=\tbbcodeactive]
# terminal 2 (host): watch the two files that matter, side by side
$ watch -n 0.5 'cat /sys/fs/cgroup/minibox/memory.current \
                 /sys/fs/cgroup/minibox/memory.events'
# you will watch Figure 3-7 draw itself:
#   ~8 MB (idle shell) -> ~120 MB (model loaded) -> spikes with each batch
#   -> at the fatal batch: memory.current snaps DOWN, oom_kill: 0 -> 1
# terminal 1 (container): the victim's-eye view
$ python toy_serve.py
loaded model (112.4 MB)
batch  4 ok   peak ~158 MB
batch 16 ok   peak ~214 MB
Killed                          # <- the shell says this, not Python
$ echo $?
137
\end{Verbatim}
\end{terminalbox}
\figcaption{기록 3-12  예상대로 펼쳐지는 B부. 신호를 보고하는 셸이 "Killed"라고 말하며 Python에는 마지막 말할 기회가 없다는 점에 주목한다. 이후 기록 3-11처럼 dmesg와 memory.events를 수집한다.}

관찰을 엔지니어링으로 바꾸는 후속 실험 두 가지를 실행한다. \textbf{(1)} \code{memory.high}를 200 MiB로 설정하고 반복한다. 강제 종료가 일어나기 전에 벽 앞 단계의 성격이 바뀐다. 진행이 느려지고 회수가 요동치며 \code{memory.events}의 메모리 압력 카운터가 증가한다. 경보선을 직접 경험하고 14주 차의 좋은 경보가 무엇에 반응하는지 미리 본다. \textbf{(2)} 스크립트 배치 크기를 살아남은 가장 큰 값으로 제한하고 같은 256 MiB 한도에서 전체 실행이 살아남는지 보인다. "가변 부분에 경계를 둔다"는 해결책을 직접 입증하며 계산 예시의 방법을 자체 수치로 뒷받침한다.

\begin{calloutbox}{tbbpurple}{tbbpurplefill}
{\normalsize\bfseries\color{tbbpurple}실습 체크리스트}\par\vspace{3pt}
\begin{tbbitemize}
  \item A부 아티팩트: 호스트 이름 대비 · /proc 마운트 전후 ps · 내부와 외부의 /proc/\$\$/ns · 호스트 쪽 cgroup.procs.
  \item B부 아티팩트: memory.current 시간 흐름(수기로 적어도 됨) · 종료 코드 137 · dmesg 줄(anon-rss 수치 인용) · memory.events 카운터.
  \item 후속 실험: memory.high 경고 단계의 동작 · 같은 한도에서 살아남는 배치 상한 재실행.
  \item 보고서: 각각 해당 절을 밝힌 관찰 세 가지(예: "proc 마운트 전에 ps 실패 — §3.2의 도구는 /proc를 읽는다는 교훈").
  \item 정리: 셸 종료, cgroup rmdir, rootfs 삭제. (이번 주에는 과금될 클라우드 자원이 없지만 제2장의 습관을 유지한다.)
\end{tbbitemize}
\end{calloutbox}

\vspace{3.0pt}

\begin{calloutbox}{tbbpurple}{tbbpurplefill}
{\normalsize\bfseries\color{tbbpurple}생각해 보기}\par\vspace{3pt}
\begin{tbbitemize}
  \item B부에서 강제 종료 전에 memory.current가 저절로 떨어질 때가 있다. 커널은 무엇을 하고 있으며 §3.3의 어떤 미묘함을 보여 주는가? 여기서 OOM을 늦추지만 막지는 못하는 이유는 무엇인가?
  \item 실행 전에 예측해 보자. 장난감 서버가 워커 두 개를 포크하고 치명적인 급증이 워커에서 발생하면 OOM 킬러는 어느 프로세스를 선택하며 이후 서버는 어떤 모습인가? (oom\_score가 힌트다. 실습에서 검증하고 tini의 회수 의무가 중요한 이유와 연결하라.)
\end{tbbitemize}
\end{calloutbox}

\namedsection{요약}{열 가지로 정리한 제3장}

\qitem{1}{\textbf{커널에는 컨테이너 객체가 없다.} 컨테이너는 공유 커널의 일반 프로세스를 네임스페이스(프로세스에 보이는 것) + cgroup(사용할 수 있는 것) + 계층형 파일 시스템(세계가 오는 곳)으로 감싸는 관례다. 제2장 격리 스펙트럼의 주의 사항을 이제 메커니즘으로 이해한다.}

\qitem{2}{\textbf{네임스페이스는 머신이 아니라 관점을 비공개로 만든다.} PID, NET, MNT, UTS, IPC, USER와 cgroup, 시간 네임스페이스가 있다. API는 unshare/clone/setns이며, \code{docker exec}는 겉모습만 바꾼 setns일 뿐이다.}

\qitem{3}{\textbf{PID 1은 직무다.} 고아를 회수하고 신호를 처리해야 한다. 그렇지 않으면 롤아웃(rollout)에서 SIGTERM이 무시되고 끝내 SIGKILL을 받는다. 서버에서 SIGTERM을 처리하거나 tini를 실행해야 한다. 이 한 문단이 실제 사고를 막는다.}

\qitem{4}{\textbf{NET 네임스페이스는 모든 컨테이너에 자체 포트 공간을 제공}하고, veth 쌍과 NAT를 통해 호스트 브리지에 연결한다. Kubernetes pod는 NET 네임스페이스 하나를 공유하는 컨테이너다. 그래서 localhost 사이드카와 pod당 IP 하나가 가능하다(5주 차).}

\qitem{5}{\textbf{cgroups v2는 예산의 파일 시스템이다.} 파일에 한도를 echo하고 PID를 등록한 뒤 사용량을 읽는다. Docker 플래그와 K8s 한도가 이 파일을 대필한다.}

\qitem{6}{\textbf{cpu.max는 뭉텅이로 스로틀링한다.} 프로세스가 창의 할당량을 소진하면 통째로 멈춘다. 프로파일러에는 보이지 않고 cpu.stat의 nr\_throttled에는 보이며 p99에 치명적이다. 제2장 스틸 타임의 사촌이다.}

\qitem{7}{\textbf{memory.max는 집행자가 지키는 벽이다.} 회수 실패 후 벽을 넘으면 OOM 킬러가 cgroup에서 가장 큰 프로세스에 SIGKILL을 보낸다. 종료 코드 137, dmesg 자백, memory.events oom\_kill, K8s "OOMKilled"가 남는다. 오류 경로가 아니라 직접 선택한 예산의 집행이다.}

\qitem{8}{\textbf{추론 메모리는 계단 + 평탄 구간 + 트래픽에 비례하는 급증}이다. 가중치 뒤에 배치/시퀀스에 의존하는 활성값이 온다. 측정한 최대치로 한도를 산정하고, 배치/길이에 경계를 두며, memory.high를 경보선으로 쓴다. 유휴 상태로 산정해서는 안 된다.}

\qitem{9}{\textbf{overlayfs는 이미지에 경제성을 제공한다.} 호스트 전체가 공유하는 읽기 전용 계층, 컨테이너당 하나의 얇은 쓰기 가능 upper, 쓰기 시 copy-up, 삭제 시 whiteout이다. 변경 속도에 따른 계층화 ⇒ 중복 제거, 빌드 캐시, 즉각적인 파일 시스템 시작이며, 이미지 가져오기가 실제 콜드 스타트 비용이다(12주 차). 가중치는 계층이 아니라 모델 저장소에 둔다.}

\qitem{10}{\textbf{다섯 단계로 컨테이너를 직접 만들 수 있다.} unshare, hostname, pivot\_root, /proc 마운트, cgroup이다. 따라서 Docker가 더하는 이미지/레지스트리, 네트워크 기본값, 수명 주기 API, 표준화된 런타임 분리(containerd/runc/OCI)가 무엇인지 정확히 안다. GPU 접근과 함께 4주 차의 주제다.}

\namedsection{연습문제}{}

\subsection*{A. 객관식}

\qitem{1}{컨테이너와 Linux 커널에 관한 설명으로 옳은 것은?}

\begin{tbboptions}
  \item ① 커널은 전용 시스템 호출로 만드는 컨테이너 객체를 제공한다
  \item ② 컨테이너는 커널이 스케줄링하는 일반 프로세스를 네임스페이스, cgroup, 마운트한 계층형 파일 시스템으로 감싼 것이다
  \item ③ 컨테이너는 경량 VM처럼 두 번째 커널을 실행한다
  \item ④ 컨테이너는 Docker 독점 커널 기능이다
\end{tbboptions}

\qitem{2}{컨테이너화 서버가 롤아웃 중 SIGTERM을 무시하고 유예 기간 뒤 SIGKILL을 받는다. 가장 가능성 높은 원인은?}

\begin{tbboptions}
  \item ① NET 네임스페이스가 외부 신호를 버린다
  \item ② CPU 스로틀링 시 cgroup이 신호 전달을 막는다
  \item ③ 서버가 PID 네임스페이스의 PID 1로 실행되면서 SIGTERM 핸들러를 설치하지 않았다. PID 1에는 기본 처리가 적용되지 않는다
  \item ④ overlayfs가 프로세스를 읽기 전용으로 만들었다
\end{tbboptions}

\qitem{3}{추론 컨테이너의 cpu.max를 "25000 100000"으로 설정했다. 부하 상태에서 가장 일치하는 관찰 결과는?}

\begin{tbboptions}
  \item ① 모든 요청이 균일하게 4× 느려진다
  \item ② p50은 거의 움직이지 않지만 cpu.stat의 nr\_throttled가 증가하면서 p99가 최대 \textasciitilde{}75 ms씩 급증한다
  \item ③ 컨테이너가 종료 코드 137로 죽는다
  \item ④ 아무 영향이 없다. cpu.max는 배치 작업에만 영향을 준다
\end{tbboptions}

\qitem{4}{pod가 애플리케이션 로그 없이 종료 코드 137로 죽고 memory.events에는 oom\_kill: 1이 나타난다. 가장 먼저 해야 할 올바른 해석은?}

\begin{tbboptions}
  \item ① 모델 런타임의 세그멘테이션 폴트이므로 상위 프로젝트에 버그를 보고한다
  \item ② 커널이 memory.max를 강제했다. 회수에 실패하자 OOM 킬러가 cgroup에서 가장 큰 프로세스에 SIGKILL을 보냈다
  \item ③ 가져오는 동안 이미지가 손상되었다
  \item ④ 디스크 부족으로 Kubernetes가 pod를 퇴거시켰다
\end{tbboptions}

\qitem{5}{노드 하나의 컨테이너 30개가 모두 같은 8 GB PyTorch 기본 이미지를 쓴다. 30번째 컨테이너의 파일 시스템에 추가되는 대략적인 디스크 용량은?}

\begin{tbboptions}
  \item ① \textasciitilde{}8 GB — 컨테이너마다 이미지를 복사한다
  \item ② \textasciitilde{}4 GB — 쓰기 시 복사(copy-on-write)가 절반으로 줄인다
  \item ③ \textasciitilde{}0 — 계층은 읽기 전용으로 공유하며 컨테이너마다 얇은 쓰기 가능 upper 디렉터리만 추가한다
  \item ④ memory.max에 따라 다르다
\end{tbboptions}

\qitem{6}{실제 컨테이너 런타임이 비공개 루트에 chroot 대신 pivot\_root를 사용하는 이유는?}

\begin{tbboptions}
  \item ① pivot\_root의 파일 I/O가 더 빠르다
  \item ② 네임스페이스 안에서는 chroot를 사용할 수 없다
  \item ③ pivot\_root는 실제로 마운트 네임스페이스의 루트를 교체하고 기존 루트를 분리한다. chroot는 경로만 다시 가리키며 설계상 탈출할 수 있다
  \item ④ pivot\_root가 rootfs를 압축한다
\end{tbboptions}

\subsection*{B. 단답형}

\qitem{1}{컨테이너 레시피의 세 재료와 각각의 역할을 한 구절로 쓰라.}

\qitem{2}{종료 코드 137을 산술적으로 해석하고 OOM 강제 종료가 프로세스 밖에서 증거를 남기는 두 위치를 쓰라.}

\qitem{3}{Kubernetes pod의 컨테이너 두 개가 localhost로 통신한다. 이를 위해 반드시 공유해야 하는 네임스페이스와 그 결과로 생기는 pod 속성은 무엇인가?}

\qitem{4}{Dockerfile의 \code{RUN rm -rf /\allowbreak opt/\allowbreak samples} 단계가 별도 계층에 있고 이미지 크기를 줄이지 못했다. 원인이 되는 overlayfs 메커니즘을 한 문장으로 설명하라.}

\subsection*{C. 서술형}

\qitem{1}{컨테이너를 직접 만드는 다섯 단계(§3.5)를 설명하라. 각 단계에서 이 장의 어느 메커니즘을 사용하며 작동을 입증하기 위해 어떤 관찰 가능한 증거를 수집할지 밝혀라. 직접 만든 컨테이너에는 없지만 Docker가 더하는 것 두 가지로 마무리하라.}

\qitem{2}{한도 2 GiB인 추론 pod가 트래픽 최대치 부근에서 거의 매일 OOM 강제 종료를 겪는다. 유휴 사용량은 1.1 GiB다. §3.6의 대응 지침으로 진단과 해결 계획을 작성하라. 무엇을 측정하고, 어떻게 크기를 조정하거나 경계를 두며, 어느 경보선을 설정할지, 각 단계가 제1장의 용어로 어떤 절충을 하는지 설명하라.}

\qitem{3}{"컨테이너가 승리한 것은 네임스페이스가 아니라 overlayfs 덕분이다." §3.4의 경제성(중복 제거, 캐시, 시작, 배포)과 §3.2–3.3의 격리 및 운영 속성을 이용해 이 주장에 찬성과 반대 논거를 제시하라. 최종 입장을 정하고 서빙 시스템과 관련된 관점에서 논증하라.}

\namedsection{정답 및 해설}{}

\subsection*{A. 객관식}

\qitem{1}{\textbf{②} 레시피 관점이다. 커널에는 컨테이너 객체가 없고(①), 제2장의 VM과 달리 두 번째 커널이 관여하지 않으며(③), 메커니즘은 공급자 중립적인 커널 기능이다(④). (§3.1)}

\qitem{2}{\textbf{③} PID 1 의미론 때문이다. 커널은 init의 기본 신호 처리를 비활성화하므로 처리되지 않은 SIGTERM은 치명적인 대신 무시된다. 코드 또는 tini/--init으로 고친다. (§3.2)}

\qitem{3}{\textbf{②} "25000 100000" = 100 ms 창마다 25 ms다. 최고 속도로 뭉텅이 실행한 뒤 최대 75 ms 동안 강제로 멈춘다. 고전적인 스로틀링 꼬리 지연 시간이며 nr\_throttled로 확인한다. 137(③)은 CPU가 아니라 메모리의 징후다. (§3.3)}

\qitem{4}{\textbf{②} 137, 조용한 애플리케이션, oom\_kill 카운터라는 세 가지가 그대로 OOM 증거 사슬이다. 대화의 주제는 "충돌 디버깅"이 아니라 "예산 크기 조정 또는 경계 설정"이다. (§3.6)}

\qitem{5}{\textbf{③} 계층은 호스트 전체에서 읽기 전용으로 공유되고, 컨테이너별 비용은 거의 비어 있는 upperdir뿐이다. 수 GB ML 이미지를 실용화하는 중복 제거의 경제성이다. (§3.4)}

\qitem{6}{\textbf{③} pivot\_root는 마운트 네임스페이스의 루트를 교체하고 기존 세계를 언마운트할 수 있게 한다. chroot는 기존 루트에 접근할 수 있어 고전적인 탈출이 가능하다. 시연용 지름길과 프로덕션 메커니즘의 차이다. (§3.5)}

\subsection*{B. 단답형}

\qitem{1}{\textbf{네임스페이스}는 프로세스에 보이는 것을 비공개로 만들고, \textbf{cgroups v2}는 사용할 수 있는 것을 제한하며, \textbf{overlayfs (+ pivot\_root)}는 공유 읽기 전용 계층에서 루트 파일 시스템을 제공한다.}

\qitem{2}{137 = 128 + 9, 즉 신호 9(SIGKILL)에 의한 죽음이다. 증거는 호스트 \code{dmesg}("Out of memory: Killed process ...")와 cgroup의 \code{memory.events}(\code{oom\_kill} 카운터)에 있다. 플랫폼은 K8s의 \code{OOMKilled} 상태로 이를 되풀이한다.}

\qitem{3}{\textbf{NET 네임스페이스}다. 공유하면 모든 pod 컨테이너가 네트워크 스택 하나를 사용하므로 pod당 IP 하나와 localhost의 사이드카가 가능하다. MNT 네임스페이스는 별도로 유지해 파일 시스템은 다르다.}

\qitem{4}{\textbf{Whiteout 의미론} 때문이다. 뒤 계층은 lower 읽기 전용 계층의 파일을 \textit{숨길} 수만 있고 바이트를 제거할 수 없으므로 삭제한 내용도 앞 계층에 그대로 실려 간다.}

\subsection*{C. 서술형(채점 기준)}

\qitem{1}{다음 요소를 확인한다. unshare(네임스페이스 §3.2, 증거: /proc/\$\$/ns inode가 다름) → hostname(UTS, 증거: 내부/외부 대비) → 푼 rootfs 위에서 pivot\_root(MNT + §3.4, 증거: 비공개 /) → /proc 마운트(증거: ps가 PID 1을 표시, 그리고 도구가 /proc를 읽는다는 상위 교훈) → cgroup 등록(§3.3, 증거: cgroup.procs가 호스트 PID를 표시). Docker가 더하는 것 중 두 가지는 이미지/레지스트리, 네트워크 배선, 수명 주기 API/데몬, OCI 런타임 분리다. 명령만이 아니라 증거가 있어야 만점이다. (§3.5)}

\qitem{2}{진단: 한도를 평탄 구간(유휴 1.1 GiB) 가까이 설정했지만 급증은 배치/시퀀스에 비례한다. 최대 트래픽 때 죽는 것은 §3.6의 전형적인 징후다. 계획: memory.current 최대치를 기록하면서 최대 배치/길이로 부하 테스트한다(9주 차 도구). 한도를 측정한 최대치 + 여유 용량보다 높이거나(절충: 노드당 레플리카 감소, 제1장의 비용/사용률 절충), 배치/시퀀스에 상한을 둔다(절충: 처리량 상한, 지연 시간/처리량 긴장). memory.high 경보선을 설정하고 memory.events에 알림을 걸며, 같은 부하 테스트로 재검증한다. 재시작 루프가 모델 로드 콜드 스타트까지 반복해 사용자 영향을 키운다고 언급하면 가산점이 있다. (§3.6)}

\qitem{3}{좋은 찬성 논거는 배포 경제성을 말한다. 계층은 수 GB 환경을 운반·캐시할 수 있고 밀도 높고 빠르게 시작하게 했다. 이것이 없었다면 네임스페이스 기술은 시스템 관리자만의 신기한 기능으로 남았을 수 있다(cgroup 역시 Docker가 뜨기 전부터 있었다). 좋은 반대 논거는 포트, PID 트리, 자원 벽이 공동 테넌시를 \textit{원할 만큼 안전하게} 만든다고 말한다. 격리 없는 밀도는 공유 상자일 뿐이다. 서빙의 구체적 요소를 다루면 어느 최종 입장도 만점이다. 찬성에서는 이미지 가져오기 콜드 스타트와 계층 위생, 반대에서는 OOM 벽과 레플리카 서버군의 포트 비공개성을 다룬다. (§3.2–3.4)}

\namedsection{더 읽을거리}{}

\textbf{[이번 주]} 표시는 지금 보거나 읽을 필수 자료이며, \textbf{[N주 차]} 표시는 해당 주차에 도움이 될 자료다.

\begin{tbbitemize}
  \item \textbf{[이번 주] Liz Rice, "Containers from Scratch" (GOTO 2018, video)} — 이번 주 필수 시청 자료이자 실습의 원전이다. 수십 줄의 Go로 컨테이너를 실시간으로 만든다. A부를 직접 시도한 \textit{뒤}에 보자. 기시감 자체가 목적이다.
  \item \textbf{[이번 주] man 페이지: namespaces(7), cgroups(7), unshare(1), pivot\_root(2)} — 커널 문서치고는 유난히 잘 쓰인 글로, §3.2–3.3의 권위 있는 세부 사항을 제공한다. namespaces(7)은 한 번 완전히 읽고 나머지는 실습 지원 자료로 훑어본다.
  \item \textbf{Liz Rice, "Container Security" (O'Reilly, 2020)} — 이 장의 세계관을 책 한 권으로 확장했다. 제2\textasciitilde{}6장은 이 책의 서빙 관점을 보완하는 보안 관점으로 네임스페이스, cgroup, 이미지 계층을 다룬다.
  \item \textbf{cgroup v2 커널 문서(admin-guide/cgroup-v2)} — 컨트롤러 의미론의 1차 자료다. 메모리 컨트롤러 절은 §3.6의 해설처럼 읽힌다. 설계 역사를 보려면 LWN의 cgroup v2 연재와 함께 읽는다.
  \item \textbf{Docker 문서: "About storage drivers"와 "Use the OverlayFS storage driver"} — 계층, copy-up, whiteout을 구체적인 그림과 함께 다룬다. §3.4의 실용적 동반 자료다.
  \item \textbf{[4주 차] OCI Runtime Specification} (opencontainers.org) — 지금 config.json을 한 번 훑어보면 모든 필드(네임스페이스, cgroup, 마운트, 루트)가 이 장을 JSON으로 옮긴 것임을 알아볼 수 있다. 4주 차의 중심이 된다.
  \item \textbf{[5주 차] Kubernetes 문서: "Resource Management for Pods and Containers"} — 5주 차 장을 읽은 뒤 살펴본다. requests/limits가 스케줄러 입력과 cgroup 쓰기로 나뉘고, QoS 클래스가 OOM 우선순위 정책임을 볼 수 있다.
  \item \textbf{[4주 차] NVIDIA Container Toolkit 문서} — 다음 주로 이어지는 다리다. GPU 장치, 드라이버 라이브러리, CUDA 런타임이 이번 주에 만든 상자 안에 들어가는 방식을 설명한다.
\end{tbbitemize}

\vspace{5.0pt}

\begin{calloutbox}{tbbblue}{tbbbluefill}
{\normalsize\bfseries\color{tbbblue}다음 장 — 제4장: Docker와 컨테이너 런타임 스택}\par\vspace{3pt}
이제 맨손으로 컨테이너를 만들 수 있다. 다음 주에는 이 일을 하루에 10억 번 수행하는 산업용 기계를 만난다. 구성 요소가 Docker → containerd → runc로 조직되는 방식, OCI 명세가 런타임을 교체 가능하게 하는 이유(Kubernetes 아래에 들어간 Kata와 gVisor를 떠올리자), 이미지를 잘 빌드하는 법(Dockerfile 계층 규율, 다단계 빌드 — §3.4를 엔지니어링 실무로 전환), 이미지를 옮기는 법(레지스트리, 다이제스트, 가져오기 — 콜드 스타트 항목)을 배운다. 그리고 이 책이 아껴 둔 질문을 다룬다. \textbf{GPU는 상자 안에 어떻게 들어가는가?} NVIDIA Container Toolkit, 장치 주입, "cgroup은 GPU 계산을 스케줄링할 수 없다"는 사실(§3.3)이 ML 컨테이너의 모든 것을 좌우하는 이유를 설명한다. 실습에서는 1주 차 스택 그림의 FastAPI 모델 서버를 컨테이너화해 레지스트리에 푸시한다. 5주 차에 Kubernetes에 배포할 이미지다.
\end{calloutbox}

%% Chapter 4 — Docker와 컨테이너 런타임 스택
\setcounter{chapter}{3}
\chapter{Docker와 컨테이너 런타임 스택}
\chaptersubtitle{다섯 단계에서 산업으로: 이미지, 레지스트리, OCI 분리, 그리고 GPU를 상자에 넣는 법}

\begin{calloutbox}{tbbblue}{tbbbluefill}
{\normalsize\bfseries\color{tbbblue}학습 목표}\par\vspace{3pt}
\begin{tbbitemize}
  \item "내 머신에서는 되는데"라는 인수인계 실패와 1956년 해상 컨테이너 이야기를 바탕으로 업계가 이미지를 채택한 이유, 곧 환경 자체가 배포 결과물이 된다는 원리를 설명하고 이미지 안에 있는 것과 없는 것을 정확히 구분한다.
  \item 하나의 \code{docker run}이 CLI → dockerd → containerd → shim → runc로 이어지는 전체 스택을 추적하고, 각 구성 요소의 단 하나뿐인 역할과 데몬 재시작 뒤에도 컨테이너가 살아남는 이유를 설명한다(제3장의 shim 예고를 마무리한다).
  \item OCI의 세 규격(이미지, 배포, 런타임)을 빌더·레지스트리·런타임을 교체할 수 있게 하는 접합부로 설명하고, Kata와 gVisor가 어디에 들어가며 kubelet이 Docker와 직접 통신하지 않는 이유를 밝힌다.
  \item 이미지를 매니페스트, 설정, 계층, 태그와 다이제스트로 이루어진 콘텐츠 주소 지정 결과물로 읽고, 콘텐츠 주소 지정이 중복 제거·신뢰할 수 있는 캐시·변조 탐지에 주는 이점을 설명한다.
  \item 변경 빈도에 따른 계층 순서, 다단계 빌드, exec 형식 CMD(PID 1), 비루트 USER를 적용해 서빙용 Dockerfile을 설계하고, 단순 이미지와 개선한 이미지의 풀 시간과 용량 차이를 계산한다.
  \item 계층 중복 제거 push/pull, 다이제스트 고정, 인증, 속도 제한, 배포 시 풀 폭풍 실패 모드를 이해하며 레지스트리를 의식적으로 운영한다.
  \item 호스트 드라이버, 주입되는 사용자 공간, 이미지의 CUDA, 툴킷 훅, 디바이스 플러그인을 거쳐 GPU가 실제로 컨테이너 안에 들어오는 과정을 설명하고, cgroup은 장치 \textit{접근}만 통제할 뿐 GPU \textit{연산}을 스케줄링할 수 없는 이유를 밝힌다.
\end{tbbitemize}
\end{calloutbox}

\vspace{5.0pt}

\begin{calloutbox}{tbbteal}{tbbtealfill}
{\normalsize\bfseries\color{tbbteal}핵심 용어}\par\vspace{3pt}
의존성 지옥 · "내 머신에서는 되는데" · 사용자 공간 · 재래식 하역과 복합 운송 컨테이너 · docker CLI · dockerd · containerd · containerd-shim · runc · OCI(image-spec / distribution-spec / runtime-spec) · config.json · CRI · dockershim(제거됨) · 이미지 매니페스트 · 이미지 설정 · 계층 · 다이제스트(sha256) · 태그 · 콘텐츠 주소 지정 · Dockerfile · 빌드 캐시 · .dockerignore · 다단계 빌드 · -slim / -runtime / -devel 베이스 · exec 형식 CMD · USER · 레지스트리 · push / pull · 풀스루 캐시 · imagePullSecrets · 속도 제한 · 풀 폭풍 · buildx / 다중 아키텍처 · NVIDIA Container Toolkit · 드라이버/CUDA 런타임 분리 · NVIDIA\_VISIBLE\_DEVICES · 디바이스 플러그인 · docker inspect · OOMKilled · exec format error
\end{calloutbox}

\vspace{5.0pt}

\begin{calloutbox}{tbblightgray}{tbbboxfill}
{\normalsize\bfseries\color{tbbink}장 구성}\par\vspace{3pt}
\textbf{4.1} 다섯 단계에서 산업으로    \textbf{4.2} 스택: Docker → containerd → runc와 OCI 접합부    \textbf{4.3} 엔지니어링 결과물로서의 이미지: 계층, 다이제스트, Dockerfile    \textbf{4.4} 레지스트리: 콘텐츠 주소 기반 배포    \textbf{4.5} GPU를 상자 안에 넣는 법    \textbf{4.6} 검시관처럼 컨테이너 읽기: inspect 현장 안내서    \textbf{4.7} 실습: 모델 서버를 컨테이너화하고 push한 뒤 검증하기. 이어서 요약, 연습문제, 더 읽을거리가 나온다.
\end{calloutbox}

\vspace{3.0pt}

\begin{calloutbox}{tbbamber}{tbbamberfill}
{\normalsize\bfseries\color{tbbamberdark}이 장을 읽는 방법}\par\vspace{3pt}
이번에도 터미널에서 시작한다. §4.4까지의 모든 내용은 Docker가 설치된 머신이라면 어디서나 재현할 수 있다. 제2장에서 만든 VM이나 노트북이면 충분하고, 실습 레지스트리는 로컬 \code{registry:2} 컨테이너이므로 계정도 비용도 필요 없다. §4.5의 GPU 터미널 기록은 GPU 노드에서 실행한 것으로 표시한다. 하드웨어가 없는 팀은 내용을 읽은 뒤 5주 차 관리형 클러스터 부록에서 재현한다. 가능하면 강의 전에 터미널 기록 4-1을 직접 실행해 보자. 프로세스 트리 한 장에 이 장의 논지가 모두 담겨 있다.
\end{calloutbox}

\section{다섯 단계에서 산업으로}

제3장은 맨손으로 컨테이너를 만들었다는 자랑으로 끝났다. 이 장에서 그 기술을 산업화하기 전에 더 근본적인 질문에 답해야 한다. \textit{사람들은 왜 굳이 이런 수고를 하는가?} "컨테이너가 VM보다 가볍기 때문"이라는 답은 제2장에서 이미 다루었고, "Kubernetes가 필요로 하기 때문"이라는 답은 제5장의 주제다. 출발점은 훨씬 인간적인 이유다. 세상의 모든 엔지니어링 팀이 한 번쯤 겪었을 만큼 보편적인 작은 참사다. 수많은 실화를 바탕으로 재구성해 보자.

\subsection*{사례 A: 인수인계}

\begin{terminalbox}
\begin{Verbatim}[breaklines=true,breakanywhere=true,fontsize=\small,formatcom=\color{tbbtermfg},codes=\tbbcodeactive]
[mon 09:12] jun   model server done! just clone the repo and run:
                  pip install -r requirements.txt && python serve.py
[mon 09:58] mina  ModuleNotFoundError: No module named 'sentencepiece'
[mon 10:20] jun   ah, install that too (added to requirements... i think)
[mon 11:03] mina  OSError: libgomp.so.1: cannot open shared object file
[mon 11:30] jun   right - sudo apt install libgomp1. which OS are you on?
[mon 13:44] mina  RuntimeError: CUDA error: no kernel image is available
                  for execution on the device
[mon 14:02] jun   hm. works on my machine :/
[tue 16:55] mina  finally running locally! deploying to the prod VM now
[tue 17:40] mina  prod python is 3.8. we built on 3.11. rolling back.
\end{Verbatim}
\end{terminalbox}
\figcaption{사례 4-A  인수인계. 각 줄은 같은 병이 서로 다른 계층에서 드러난 모습이다. 실행 환경은 저장소가 아니라 준의 노트북에 살고 있었다.}

각 실패는 서로 다른 깊이에 있으므로 한 줄씩 검시해 보자. \code{sentencepiece}는 requirements.txt에 들어가지 못한 \textbf{Python 계층} 의존성이다. 이 파일에는 코드에 \textit{필요한 것}이 아니라 준이 \textit{기억한 것}만 적혀 있었다. \code{libgomp.so.1}은 pip 세계 전체의 아래에 있는 \textbf{시스템 라이브러리}다. 아무리 꼼꼼한 requirements.txt도 이를 표현할 수 없고, virtualenv는 Python 계층만 둘러쌀 뿐 그 아래에는 미치지 못한다. CUDA 오류는 세 계층 아래의 \textbf{사용자 공간과 드라이버} 불일치다. 화요일의 마지막 문제는 노트북과 프로덕션의 \textbf{인터프리터 자체}가 달랐다는 점이다. 패턴은 분명하다. \code{requirements.txt}가 실제 환경의 5\%쯤을 운반했다면, 나머지는 누구도 기록하지 않은 채 한 노트북에 수개월 동안 \code{pip install}과 \code{apt install}로 쌓인 \textbf{상태}였다. 차이를 비교할 수도 없고, "온보딩"이라는 제목의 위키 문서 속에서 서서히 썩어 갔다. 이 병의 이름은 \textbf{의존성 지옥}이고 구호는 \textbf{"내 머신에서는 되는데"}다. 구호가 정확한 까닭은 머신 자체가 결과물이 되었기 때문이다. 컨테이너 이전의 처방은 모두 이 사실을 피했다. 위키는 낡고, 영웅 한 명에게 의존하는 방식은 확장되지 않는다. VM 전체 스냅숏은 머신을 고정할 수 있지만, 차이를 보거나 소스에서 다시 만들거나 CI로 운반할 수 없는 불투명한 10 GB 덩어리가 된다.

\subsection*{사례 B: 세상을 작게 만든 상자}

해법은 더 오래된 산업에서 왔다. 그 계보(lineage)는 어떤 설명서보다 Docker의 설계를 잘 보여 주므로 2분쯤 살펴볼 가치가 있다. 1956년 이전 해상 화물은 \textbf{재래식 하역(break-bulk)} 방식으로 움직였다. 부두 노동자 무리가 배마다 통, 포대, 상자를 손으로 실었고, 배는 바다보다 항구에서 더 오래 머물기도 했다. 하역 비용은 톤당 약 \textbf{5.86달러}로 유명했다. 1956년 4월 26일, 운송 사업가 \textbf{맬컴 맥린(Malcolm McLean)}은 개조한 유조선 \textit{Ideal-X}에 똑같은 강철 상자 58개를 싣고 뉴어크에서 휴스턴으로 보냈다. 하역 비용은 톤당 약 \textbf{0.16달러}, 약 37분의 1로 떨어졌다. 천재적인 부분은 상자 자체가 아니라 \textbf{표준화된 모서리 고정부}였다. 어떤 크레인과 배, 트럭, 열차도 상자 안에 무엇이 있는지 \textit{알거나 신경 쓰지 않고} 모든 컨테이너를 집을 수 있었다. 항구는 상자를 열지 않고 그대로 옮기기 시작했다. 취급 방식이 통일되자 자동화할 수 있었고, 20년 안에 세계 무역이 이 상자를 중심으로 재편되었다. 이 역사를 다룬 마크 레빈슨(Marc Levinson)의 \textit{The Box}는 이번 주의 재미있는 읽을거리다(더 읽을거리 참고).

Docker가 내세운 비유는 장식이 아니다. 로고부터 컨테이너를 쌓은 고래다. 대응 관계도 거의 기계적이다. \textbf{이미지}는 밀봉된 표준 모양의 상자이며 내용물은 무엇이든 될 수 있다. \textbf{OCI 규격}(§4.2)은 표준 손잡이인 모서리 고정부다. \textbf{런타임과 오케스트레이터}, 곧 runc, containerd, Kubernetes는 크레인과 배와 열차다. 모든 상자를 같은 방식으로 다루기 때문에 하나의 스케줄러가 모든 팀의 워크로드를 실행할 수 있다. \textbf{레지스트리}(§4.4)는 항구와 물류 기지다. 경제성도 닮았다. 강철 상자가 싸서 성공한 것이 아니라 \textit{취급 방식}이 통일되었기 때문이다. 소프트웨어에서 "인프라는 상자를 열어 보지 않는다"는 무지가 바로 제5장과 제6장에서 계속 "플랫폼"이라고 부른 성질이다. 클러스터는 이미지 안에 Python이 있는지 Rust가 있는지, 13 GB 모델 서버가 있는지 모른다. 그것이 기능이며 전부다.

\subsection*{그렇다면 상자 안에는 정확히 무엇이 있을까?}

비유는 충분하다. 이제 상자를 열어 보자. 실습에서 사용할 베이스 이미지면 된다.

\begin{terminalbox}
\begin{Verbatim}[breaklines=true,breakanywhere=true,fontsize=\small,formatcom=\color{tbbtermfg},codes=\tbbcodeactive]
$ docker export $(docker create python:3.11-slim) | tar -t | wc -l
28414                    # ~28,000 files (representative count)
$ docker export $(docker create python:3.11-slim) | tar -t \
    | grep -m4 -E "libc.so|ca-certificates.crt|Seoul|bin/python3.11$"
lib/x86_64-linux-gnu/libc.so.6            # glibc: the userland floor
etc/ssl/certs/ca-certificates.crt         # trust store: HTTPS works
usr/share/zoneinfo/Asia/Seoul             # timezones. all of them.
usr/local/bin/python3.11                  # the interpreter itself
$ docker export $(docker create python:3.11-slim) | tar -t \
    | grep -cE "^boot/vmlinuz|^lib/modules"
0                        # and NO kernel. none. (Ch. 3 explained why)
\end{Verbatim}
\end{terminalbox}
\figcaption{사례 4-B  python:3.11-slim의 내부. 인터프리터, glibc, 인증서, 시간대 데이터를 갖춘 완전한 Linux 사용자 공간이 들어 있지만, 주목할 점은 어디에도 커널이 없다는 사실이다.}

이제 정직한 목록을 만들 수 있다. 이미지는 \textbf{프로세스의 사용자 공간 전체를 얼린 것}이다. 인터프리터와 모든 site-package에서 C 라이브러리와 glibc에 이르기까지 프로세스가 \code{open()}할 수 있는 모든 파일이 들어 있다. 프로세스가 당연히 존재한다고 여기는 조용한 문명도 함께다. 컨테이너에서 HTTPS가 작동하게 하는 CA 인증서, 시간대 데이터베이스, \code{/\allowbreak etc/\allowbreak passwd}, 셸이 그 예다. 바로 이 점 때문에 미나의 월요일을 다시 겪을 필요가 없다. "먼저 X를 설치하라"는 절차가 없고, 모든 X가 바이트 단위로 똑같이 따라온다. 상자에 \textbf{없는 것}을 적은 반대 목록도 중요하다. \textbf{커널은 없다.} 어떤 컨테이너에서든 \code{uname -r}을 실행하면 \textit{호스트} 커널 버전이 나온다. 제3장의 핵심이 컨테이너가 커널을 공유한다는 사실이었기 때문이다. "이미지는 작은 VM 아닌가?"라는 질문에도 답이 된다. VM 이미지는 부팅해야 하지만 컨테이너 이미지는 부팅을 시도할 수도 없다. \textbf{GPU 드라이버도 없다.} 호스트의 속성이며 시작할 때 주입된다(§4.5). \textbf{모델 가중치도 없다.} 모델 저장소의 속성이며 부팅할 때 가져온다(제6장). \textbf{init 시스템, sshd, cron도 없다.} 머신의 세계가 아니라 프로세스 하나의 세계다. 극단적으로 단순한 경우에도 이 정의는 그대로다.

\begin{terminalbox}
\begin{Verbatim}[breaklines=true,breakanywhere=true,fontsize=\small,formatcom=\color{tbbtermfg},codes=\tbbcodeactive]
$ cat Dockerfile.scratch
FROM scratch             # the empty image: zero bytes, zero files
COPY hello /hello        # one static binary, 1.2 MB
CMD ["/hello"]
$ docker build -q -t tiny -f Dockerfile.scratch . && \
    docker export $(docker create tiny) | tar -t
hello                    # the ENTIRE filesystem. one file.
\end{Verbatim}
\end{terminalbox}
\figcaption{사례 4-C  직접 만들어 확인한 증명. 이미지에는 넣은 것만 정확히 들어 있다. "베이스 이미지"는 아래에서 실행되는 OS가 아니라, 다른 사람이 만든 첫 몇 개 계층을 바닥으로 채택한 것이다.}

이제 \textbf{이미지 = 얼린 사용자 공간 + 시작 방법을 적은 메모}라는 모델을 기억하자. 여기서 메모는 §4.3의 설정, 곧 환경 변수·명령·사용자다. 그러면 인수인계 이야기가 제대로 끝난다. 준은 금요일 오후에 \code{docker build}를 실행하고, 상자는 인터프리터를 포함한 환경 전체를 싣는다. 미나는 월요일 09시 13분에 컨테이너 런타임이 있는 어떤 머신에서든 이를 실행하며, 프로덕션도 \textit{같은 다이제스트}를 실행한다. 위키는 은퇴한다. 머신은 더는 결과물이 아니다. 결과물이 진짜 결과물이 된다.

\subsection*{노트북 한 대의 치료에서 산업으로}

이 치료법을 한 번 적용하면 편리한 도구에 그친다. 이 장의 주제는 \textbf{서로 만난 적도 없고 누구도 직접 조립하지 않은 머신들에서 하루 10억 번} 적용하는 데 무엇이 필요했는가다. 더 큰 셸 스크립트만으로는 그 간격을 건널 수 없다. 문제는 운송 이야기와 대응하는 세 산업적 과제로 나뉜다. \textbf{패키징}은 파일 시스템과 설정을 운반하고 검증할 수 있는 결과물로 만드는 일, 곧 상자 자체다(§4.3). \textbf{배포}는 수 GB짜리 상자를 머신 사이에서 몇 초 만에 옮기는 일, 곧 항구다(§4.4). \textbf{수명 주기}는 사람 없이 수천 개를 시작·감시·재시작하는 일, 곧 크레인이다(§4.2). Docker의 역사적 성취는 세 문제를 하나의 일관된 도구 모음으로 해결한 것이다. 더 오래 남은 성취, 그래서 이 장이 특정 도구보다 오래가는 이유는 그 해법을 누구나 구현할 수 있는 인터페이스로 표준화했다는 점이다(§4.2).

이 장은 제3장, 제5장, 제6장이 적어 온 빚 장부도 정리한다. 제3장에서 컨테이너의 부모로 나타난 수수께끼의 shim 프로세스(PID 41288)는 설명을 약속했고 §4.2에서 갚는다. \code{RUN rm}으로 이미지를 줄일 수 없음을 증명하며 "계층 규율"을 예고한 내용은 §4.3에서 가격표가 붙은 엔지니어링 실천으로 바뀐다. 이미지 풀에 25\textasciitilde{}60초가 든다는 제3장의 콜드 스타트 계산은 §4.4에서 대부분의 풀을 2초로 줄이는 장치를 만난다. 제3장의 가장 큰 부정적 결론, 곧 \textit{GPU용 cgroup 컨트롤러는 없다}는 사실은 §4.5의 중심 질문이 된다. 커널이 GPU를 중재할 수 없다면 컨테이너는 어떻게 GPU를 받고, 받았다는 것은 정확히 무엇을 약속하는가?

\subsection*{영역에 들어가기 전에 지도를 보자}

움직이는 부품의 이름이 많으므로 먼저 지도를 고정하자. \code{docker run}을 입력하면 네 프로그램이 협력한다. \textbf{CLI}는 REST 클라이언트일 뿐이다. 제5장의 kubectl과 같은 방식으로 신비를 벗긴다. 이 책에서는 "플랫폼"을 열어 볼 때마다 HTTP 클라이언트가 튀어나온다. \textbf{데몬}(\code{dockerd})은 개발자용 API와 편의 기능을 맡는다. \textbf{수명 주기 엔진}(\code{containerd})은 이미지와 실행 중인 작업을 관리하며, Kubernetes는 Docker를 건너뛰고 이를 \textit{직접} 사용한다. \textbf{런타임}(\code{runc})은 JSON 규격을 엄격히 따라 제3장의 다섯 단계을 수행한 뒤 \textit{종료}하는 것이 전부다. containerd와 프로세스 사이에는 작지만 오래 살아남는 \textbf{shim}이 있어, 이후 그 프로세스를 책임진다. §4.2에서는 터미널 기록으로 이 사슬을 따라간다. 한 문장으로 요약하면 다음과 같다. \textbf{Docker 스택의 어떤 것도 모델을 직접 실행하지 않는다. 실행은 커널이 한다. 나머지는 네임스페이스, cgroup, overlayfs 주변의 장부 정리이며, 바로 그 때문에 모든 부품을 교체할 수 있다.}

두 결과물이 이 장치 안을 흐르며 장 중반의 중심을 잡는다. \textbf{이미지}(§4.3)는 제3장의 overlayfs 계층 스택에 운반 가능한 콘텐츠 주소 기반 파일 형식을 입히고, 런타임이 무엇을 해야 하는지 알려 주는 설정을 더한 것이다. \textbf{레지스트리}(§4.4)는 이 결과물이 이동하는 배포 시스템이다. 모든 대화가 다이제스트로 이루어지므로 코드 한 줄을 바꿔 push할 때는 KB만 움직이지만 CUDA 베이스를 처음 pull할 때는 GB가 움직인다. 제6장에서 모델 저장소를 만들 때 불변의 버전 결과물, 체크섬, "새 버전 = 새 접두사"가 익숙했다면 이미 몇 주 동안 이 장의 레지스트리 규율을 \textit{사용}했기 때문이다. 이미지는 이 책이 콘텐츠 주소 기반 저장소로 운반한 첫 결과물이었다. 책의 순서로는 여기서 한 번 배우고 다른 곳에서 계속 알아보게 된다.

\subsection*{서빙 엔지니어가 관심을 가져야 하는 이유 — "바로 아래 계층이니까"를 넘어}

\begin{tbbitemize}
  \item \textbf{이미지는 이 책의 나머지 전체가 쓰는 배포 단위다.} 이번 주 실습에서 컨테이너화한 FastAPI 서버는 바이트 하나까지 똑같이 제5장이 스케줄링하고(\code{registry.example.com/\allowbreak team/\allowbreak model-server}), 제6장이 모델 저장소에 연결하며, 9\textasciitilde{}14주 차가 서빙·오토스케일링·모니터링하는 바로 그 결과물이다. 지금 정한 크기, 계층 순서, PID 1 정확성, 아키텍처를 이후 모든 주가 물려받는다.
  \item \textbf{이미지 크기는 지연 시간이자 청구 금액이다.} 제3장은 6 GB 풀을 25–60 s의 콜드 스타트로 계산했다. §4.3의 계산 예시는 계층 순서를 개선해 같은 서버를 3분의 1 크기로 운반하며, 12주 차에는 이를 오토스케일링 반응 속도로 곧바로 환산한다.
  \item \textbf{GPU 절이 이후 모든 내용을 떠받친다.} 제2장, 제3장, 제5장은 공급업체 공유 기법, cgroup 컨트롤러의 부재, 정수 단위 스케줄링이라는 각자의 관점에서 "GPU는 특별하다"는 사실을 다루었다. §4.5에서 세 관점이 기계적으로 만난다. 이후에는 어느 것도 전설처럼 외울 필요가 없다.
  \item \textbf{이 계층의 실패는 표의 한 행으로 읽기 전까지 마법처럼 보인다.} \code{exec format error}, \code{toomanyrequests}, \code{CUDA driver version is insufficient}를 §4.6 현장 안내서에서 제3장의 exit 137, 제5장의 Pending 옆에 놓고 징후·원인·첫 명령으로 정리한다.
\end{tbbitemize}

\begin{calloutbox}{tbbpurple}{tbbpurplefill}
{\normalsize\bfseries\color{tbbpurple}생각해 보기}\par\vspace{3pt}
\begin{tbbitemize}
  \item 사례 4-A의 검시 결과를 표로 다시 작성해 보자. 각 실패가 속한 계층(Python / 시스템 라이브러리 / 드라이버 / 인터프리터)을 적고, virtualenv로 해결되는 것, conda로 해결되는 것, 사용자 공간 전체를 얼려야 해결되는 것을 나눈다. Docker를 써도 여전히 실패하는 행이 정확히 하나 있다. 찾아서 §4.5가 어떻게 다루는지 설명하라.
  \item 운송 비유가 깨질 때까지 확장해 보자. 세관 검사는 무엇에 해당하는가? 선하증권과 냉장 컨테이너는 무엇인가? 그런 다음 소프트웨어 상자는 할 수 있지만 강철 상자는 할 수 없는 행동을 찾아 비유가 실제로 실패하는 지점을 하나 제시하라. §4.3의 계층 공유가 좋은 출발점이다.
  \item 제3장에서 손으로 만든 컨테이너는 "한 머신에서 한 번" 작동했다. rootfs tarball과 설정 스크립트를 다른 배포판을 쓰는 동료에게 이메일로 보낼 때 무엇이 구체적으로 깨지는지 나열하고, 각 실패가 세 산업적 문제(패키징, 배포, 수명 주기) 중 어디에 속하는지 분류하라.
\end{tbbitemize}
\end{calloutbox}

\section{스택: Docker → containerd → runc와 OCI 접합부}

제3장의 첫 터미널 기록은 눈앞에 단서를 남겨 두었다. 컨테이너 프로세스(PID 41320)의 부모는 예상했던 Docker 데몬이 아니라 정체불명의 중간자 PID 41288이었다. 이제 가족 전체를 불러 조사해 보자(그림 4-1).

\begin{terminalbox}
\begin{Verbatim}[breaklines=true,breakanywhere=true,fontsize=\small,formatcom=\color{tbbtermfg},codes=\tbbcodeactive]
$ docker run -d --name demo -m 512m --cpus 1 python:3.11 sleep 9999
3f9a1c2b77e1...                      # Ch. 3, Transcript 3-1 - same run
$ ps -eo pid,ppid,cmd | grep -E "dockerd|containerd|sleep" | grep -v grep
  912      1  /usr/bin/dockerd -H fd://
  931      1  /usr/bin/containerd
41288      1  /usr/bin/containerd-shim-runc-v2 -id 3f9a1c2b77e1...
41320  41288  sleep 9999
# note what is MISSING: runc. it started the container and exited.
# and note the shim's parent: PID 1, not containerd, not dockerd.
$ pstree -p 41288
containerd-shim(41288)---sleep(41320)
\end{Verbatim}
\end{terminalbox}
\figcaption{터미널 기록 4-1  가족 관계를 조사한 결과. 컨테이너의 조상에는 dockerd도 runc도 없다. 평범한 프로세스를 붙들고 있는 작은 shim, 곧 제3장의 PID 41288만 남아 있다.}

각 구성 요소의 단 하나뿐인 역할과 함께 명단을 읽어 보자. \textbf{docker CLI}는 REST 클라이언트다. \code{docker run}은 Unix 소켓에 보내는 몇 번의 HTTP 호출이며, \code{docker --debug}가 그대로 보여 준다. \textbf{dockerd}는 그 소켓 뒤의 데몬으로 개발자용 API, \code{docker build}, 네트워크와 볼륨 편의 기능을 맡는다. 컨테이너를 직접 실행하지 않고 gRPC를 통해 이미지·스냅숏(제3장의 overlayfs 스택)·실행 중인 작업을 관리하는 수명 주기 엔진 \textbf{containerd}에 위임한다. 작업을 시작할 때 containerd는 \textbf{containerd-shim}을 만들고, shim은 JSON 규격을 읽어 제3장의 다섯 단계을 프로덕션 수준으로 수행하는 작은 독립 바이너리 \textbf{runc}를 호출한다. runc는 프로세스를 시작한 뒤 \textit{종료}한다. 실행 중에 남는 것은 터미널 기록 그대로다. 프로세스 하나, shim 하나, 제3장의 커널 메커니즘이다. "컨테이너 런타임"은 상주하는 하이퍼바이저가 아니라 설치 작업반이다. 컨테이너 하나당 오버헤드는 작은 프로세스 하나뿐이며, 제2장의 격리 스펙트럼에서 컨테이너가 그 위치에 놓인 큰 이유다.

\subsection*{shim이 제 밥값을 하는 모습 — 실시간 확인}

왜 굳이 shim을 끼워 넣을까? \textit{누군가}는 컨테이너의 표준 입출력을 열어 두고, 종료 코드를 수집하고, 종료한 프로세스를 회수해야 한다. 제3장의 좀비 교훈이 한 계층 위로 올라온 셈이다. 그 누군가가 데몬이어서는 안 된다. 그렇지 않으면 데몬 업그레이드가 매번 플릿 전체의 장애가 된다. shim은 의도적으로 작고 지루하며 오래 살아남는다. 덕분에 그 위의 구성 요소는 마음껏 크고 복잡하며 자주 바뀔 수 있다.

\begin{terminalbox}
\begin{Verbatim}[breaklines=true,breakanywhere=true,fontsize=\small,formatcom=\color{tbbtermfg},codes=\tbbcodeactive]
$ docker exec demo python -c "print('serving...')" && \
    sudo systemctl restart docker           # kill the whole daemon
serving...
$ sleep 3; docker ps --format "{{.Names}}: {{.Status}}"
demo: Up 9 minutes                          # never blinked
$ ps -o pid,ppid,cmd -p 41320
  PID   PPID  CMD
41320  41288  sleep 9999                    # same PID, same shim
# the daemon died and returned; the container never noticed.
# upgrades, crashes, restarts of dockerd/containerd: non-events.
\end{Verbatim}
\end{terminalbox}
\figcaption{터미널 기록 4-2  shim이 존재하는 이유를 직접 확인한다. Docker 데몬 전체를 재시작해도 컨테이너는 같은 shim 아래에서 계속 서비스한다. 제3장의 예고를 마무리했다.}

\subsection*{OCI 접합부: 세 규격과 세 교체 지점}

스택의 구성 요소는 \textbf{오픈 컨테이너 이니셔티브(Open Container Initiative, OCI)} 규격을 통해 서로의 약속을 지킨다. Docker가 runc를 기부하며 시작된 2015년의 표준화는 한 회사의 제품을 업계 공통 인터페이스로 바꾸었다. §4.1의 세 산업적 문제마다 규격 하나가 대응한다. \textbf{image-spec}은 이미지가 \textit{무엇인지} 정의한다. 매니페스트(manifest), 설정, 계층이 모두 콘텐츠 주소로 식별된다(§4.3). \textbf{distribution-spec}은 이미지가 \textit{어떻게 이동하는지} 정의한다. 모든 레지스트리가 구현하는 push/pull HTTP API다(§4.4). \textbf{runtime-spec}은 컨테이너가 \textit{어떻게 시작하는지} 정의한다. rootfs와 \code{config.json}의 묶음이며, JSON은 제3장을 설정 형식으로 옮겨 적은 것처럼 읽힌다.

\begin{terminalbox}
\begin{Verbatim}[breaklines=true,breakanywhere=true,fontsize=\small,formatcom=\color{tbbtermfg},codes=\tbbcodeactive]
$ runc spec && python3 -c "import json;c=json.load(open('config.json'));\
print(json.dumps({k:c[k] for k in ['process','root','linux']}, indent=1)[:600])"
{
 "process": { "args": ["sh"], "user": {"uid": 0}, ... },
 "root":    { "path": "rootfs", "readonly": true },
 "linux": {
  "namespaces": [ {"type":"pid"}, {"type":"network"}, {"type":"ipc"},
                  {"type":"uts"}, {"type":"mount"} ],
  "resources": { "memory": {"limit": ...}, "cpu": {...} },
  ...
# read it aloud: which namespaces to unshare (sec. 3.2), which cgroup
# values to write (sec. 3.3), which rootfs to pivot into (sec. 3.5).
\end{Verbatim}
\end{terminalbox}
\figcaption{터미널 기록 4-3  OCI 런타임 규격의 config.json은 제3장의 내용을 데이터로 옮긴 것이다. 네임스페이스, cgroup 리소스, 루트 파일 시스템이 들어 있다. runc는 이 파일의 인터프리터이며 Kata와 gVisor도 마찬가지다.}

표준화된 접합부에서 교체가 일어나며, 이 책에서는 두 지점이 중요하다(그림 4-2). \textbf{런타임 접합부}에서는 runtime-spec 묶음을 받아들이는 무엇이든 runc를 대신할 수 있다. \textbf{Kata Containers}는 하드웨어 경계를 제공하는 microVM을 시작하지만 위쪽 호출자는 그대로다. \textbf{gVisor의 runsc}는 시스템 호출 경계에 사용자 공간 커널을 끼워 넣는다. 제2장의 "더 강한 격리, 같은 워크플로" 선택지가 실제 클러스터에 들어가는 방식이 바로 이것이다. 런타임 설정 한 줄만 바꾸며 이미지와 매니페스트는 손대지 않는다. \textbf{오케스트레이터 접합부}에서는 Kubernetes의 kubelet이 gRPC 인터페이스인 \textbf{CRI}로 containerd와 직접 통신한다. 제5장의 그림 5-2가 그린 사슬이다. Docker 데몬은 애초에 클러스터 구성 요소가 아니었다. 2022년 Kubernetes가 dockershim 어댑터를 제거했을 때도 업계의 이미지는 수정 없이 계속 작동했다. \textit{이미지는 Docker 결과물이 아니라 OCI 결과물}이기 때문이다. 이 싱거운 결말이 표준의 존재 이유 전체다. \textbf{Docker가 오래 남긴 공헌은 데몬이 아니라 이제 모두가 구현하는 규격이다.}

\begin{tbbfigure}
  \adjustbox{max width=0.94\linewidth,max totalheight=0.46\textheight}{\includegraphics{figures/fig4-1_flow.png}}\\[3pt]
  \figcaption{그림 4-1  docker run의 흐름과 이후 남는 구성 요소. CLI와 데몬은 장부를 정리하고, runc는 제3장의 다섯 단계을 수행한 뒤 종료한다. 컨테이너마다 shim 하나가 남아 꼭 붙들어야 할 것을 맡는다.}
\end{tbbfigure}

\begin{tbbfigure}
  \adjustbox{max width=0.94\linewidth,max totalheight=0.46\textheight}{\includegraphics{figures/fig4-2_oci.png}}\\[3pt]
  \figcaption{그림 4-2  OCI 접합부. 빌드, 배포, 실행이 서로 다른 규격이므로 빌더·레지스트리·런타임을 독립적으로 교체할 수 있고, kubelet에는 Docker가 전혀 필요 없다.}
\end{tbbfigure}

\begin{calloutbox}{tbbpurple}{tbbpurplefill}
{\normalsize\bfseries\color{tbbpurple}생각해 보기}\par\vspace{3pt}
\begin{tbbitemize}
  \item 스택의 구성 요소(dockerd, containerd, shim, runc)를 장애가 서빙 플릿에 미치는 피해가 큰 순서대로 정렬하고, 각 구성 요소가 보유한 상태나 책임을 근거로 설명하라. 왜 아키텍처가 가장 작은 프로그램을 프로세스에 가장 가까이 두는지 답에서 드러나야 한다.
  \item Kata는 런타임 접합부에서 microVM으로 교체되며 "위쪽은 하나도 바뀌지 않는다." 제2장과 제3장을 바탕으로 워크로드에서 조용히 바뀌는 두 가지를 찾아라. 커널의 정체, /proc 내용, 장치 접근을 생각해 보자. 접합부는 계약을 표준화할 뿐 물리 법칙을 같게 만들지는 않는다.
  \item dockershim 제거는 거의 아무것도 깨뜨리지 않았지만 일부 팀에서는 \textit{무언가}가 깨졌다. 모든 노드에 Docker 데몬이 존재한다고 가정한 워크플로가 무엇인지 추론하라. CI 파드(Pod) 안의 docker build와 docker.sock 마운트를 단서로 삼고 OCI 네이티브 대안을 찾아보자.
\end{tbbitemize}
\end{calloutbox}

\section{엔지니어링 결과물로서의 이미지: 계층, 다이제스트, Dockerfile}

제3장에서는 overlayfs의 동작, 곧 읽기는 아래 계층으로 내려가고 쓰기는 위로 복사하며 삭제는 whiteout을 남긴다는 원리와 중복 제거·캐시·콜드 스타트 풀이라는 경제성을 배웠다. 이 절에서는 메커니즘을 \textbf{결과물}로 바꾼다. 정확히 무엇을 push하고, 어떻게 이름을 붙이며, 프로덕션에서 잘 동작하도록 어떻게 작성할까? 실제 이미지를 분해하는 데서 시작하자(그림 4-3).

\begin{terminalbox}
\begin{Verbatim}[breaklines=true,breakanywhere=true,fontsize=\small,formatcom=\color{tbbtermfg},codes=\tbbcodeactive]
$ docker build -t model-server:v1 . && docker images model-server
REPOSITORY     TAG   IMAGE ID       SIZE
model-server   v1    9f2c4a01be77   2.4GB
$ docker image inspect model-server:v1 --format \
    '{{.Architecture}} {{len .RootFS.Layers}} layers'
amd64 9 layers
$ docker history model-server:v1 --format "{{.Size}}\t{{.CreatedBy}}" | head -5
40kB    COPY serve.py /app/serve.py
1.9GB   RUN pip install -r requirements.txt
0B      WORKDIR /app
120MB   FROM python:3.11-slim
# every Dockerfile step that touched the filesystem = one layer,
# stackable and shippable exactly as Chapter 3 described.
\end{Verbatim}
\end{terminalbox}
\figcaption{터미널 기록 4-4  이미지를 분해했다. 파일 시스템을 건드리는 Dockerfile 단계는 계층과 일대일로 대응하며, 결과물의 정체는 친숙한 이름이 아니라 다이제스트다.}

구조적으로 OCI 이미지는 세 종류의 객체로 이루어지며, 각 이름은 바이트의 sha256이다. \textbf{매니페스트}는 작은 루트 문서로 \textbf{설정} 하나와 N개의 \textbf{계층} 다이제스트를 나열한다. 설정에는 runc에 전달할 env, \code{Cmd}, \code{User}, 노출 포트 같은 런타임 정보와 순서가 있는 계층 목록이 들어 있다. 계층은 제3장에서 본 읽기 전용 overlayfs 층을 압축한 tarball이다. 이 구조의 무게를 떠받치는 설계 선택은 \textbf{콘텐츠 주소 지정}이다. 모든 객체의 이름이 곧 해시이므로 몇 개의 이미지가 공유하든 동일한 계층은 한 번만 저장하고 전송한다(중복 제거). \code{sha256:22...}를 캐시한 머신은 누가 만들었는지 묻지 않고 절대적으로 신뢰할 수 있다. 서로 만난 적 없는 머신 사이에서도 계층 캐시가 작동하는 이유다. 무엇이든 변조하면 이름이 바뀐다(무결성). 반대로 \code{:v7} 같은 \textbf{태그}는 불변 세계를 가리키는 \textit{가변 포인터}, 즉 누군가 옮길 수 있는 별명이다. 운영 규칙은 자연스럽다. 사람은 태그로 말하지만 \textbf{배포는 다이제스트를 고정}한다(\code{image@sha256:aa41f9...}). 또는 레지스트리가 강제하는 불변 태그를 쓴다. 제6장에서 모델 버전에 같은 규율을 적용했는데, 이 책이 그 규율을 처음 배운 곳이 여기다.

\begin{tbbfigure}
  \adjustbox{max width=0.94\linewidth,max totalheight=0.46\textheight}{\includegraphics{figures/fig4-3_image.png}}\\[3pt]
  \figcaption{그림 4-3  이미지의 구조. 매니페스트가 설정과 계층을 가리키며 모든 것이 콘텐츠 주소로 식별된다. 태그는 바뀔 수 있는 별명이고 다이제스트는 정체다. 프로덕션 배포는 후자를 고정한다.}
\end{tbbfigure}

\subsection*{한 줄씩 설계하는 Dockerfile}

이미지를 작성하는 순간 이 책에서 쌓은 교훈이 파일 하나로 모인다. 아래는 실습 모델 서버의 Dockerfile이며, 모든 줄이 앞선 교훈을 증거와 함께 다시 불러온다. 그림 4-4에 전체 주석을 달았다.

\begin{terminalbox}
\begin{Verbatim}[breaklines=true,breakanywhere=true,fontsize=\small,formatcom=\color{tbbtermfg},codes=\tbbcodeactive]
# ---- stage 1: builder - compilers live here, and die here ----
FROM python:3.11 AS builder
COPY requirements.txt .
RUN pip install --prefix=/venv -r requirements.txt

# ---- stage 2: runtime - what actually ships ----
FROM python:3.11-slim
COPY --from=builder /venv /venv         # deps: change monthly
ENV PATH=/venv/bin:$PATH
RUN useradd -u 10001 serve
USER 10001                              # non-root (Ch. 3)
COPY serve.py /app/serve.py             # code: changes DAILY - last!
EXPOSE 8000
CMD ["uvicorn", "serve:app", "--host", "0.0.0.0", "--port", "8000"]
#   ^^ EXEC form: the server IS PID 1; SIGTERM lands (Ch. 3's bug)
\end{Verbatim}
\end{terminalbox}
\figcaption{터미널 기록 4-5  실습 Dockerfile. 순서로 제3장의 캐시 경제성을 구현하고, 다단계 빌드로 빌드 시점과 실행 시점 의존성을 분리하며, 마지막 줄로 제3장의 CMD 형식 버그를 영구히 고친다.}

가치의 대부분은 세 결정에서 나온다. \textbf{변경 빈도에 따른 순서}는 제3장의 표 3-3을 실행 가능하게 만든다. 천천히 바뀌는 베이스와 의존성은 아래에, 매일 바뀌는 \code{serve.py}는 마지막에 둔다. 그러면 매일 빌드할 때 40 KB짜리 마지막 계층만 새로 만들고 나머지는 모두 재사용한다. §4.4에 따라 매일의 \textit{push와 pull}도 KB만 움직인다. 여기서 중요한 결론이 하나 나온다. \code{COPY requirements.txt}와 \code{pip install}은 모든 \code{COPY . .}보다 \textit{앞에} 와야 한다. 전체 트리를 먼저 복사하면 코드 편집 하나가 의존성 계층을 무효화해 빌드마다 20분을 쓰게 된다. \textbf{.dockerignore}도 짝을 이룬다. 이미지에 들어오지 않는 빌드 컨텍스트는 캐시를 무효화할 수 없다. \code{.git}, 데이터셋, 제6장에서 다룰 로컬 모델 가중치를 제외하자. \textbf{다단계 빌드}는 실재하는 딜레마를 해결한다. 과학 계산 Python 패키지를 \textit{설치}하려면 약 1.9 GB의 컴파일러와 헤더가 필요하지만 실행할 때는 전혀 필요 없다. 1단계는 전체 도구 모음으로 \code{/\allowbreak venv}에 설치하고 2단계는 슬림 베이스에 \code{/\allowbreak venv}만 복사한다. 빌드 의존성이 실행 페이로드에서 빠져 CPU 실습 이미지는 2.4 GB에서 약 610 MB로 줄어든다. 마지막은 \textbf{PID 1 정확성과 비루트 실행}이다. exec 형식 \code{CMD}는 서버를 PID 1로 만들어 SIGTERM이 실제로 도달하게 한다. 제3장의 롤아웃, 제5장의 드레이닝, 제6장의 502 정렬이 한 줄에서 결실을 본다. \code{USER 10001}은 애플리케이션이 침해되어도 공격자가 비특권 사용자에 머물게 한다. 컨테이너 탈출의 피해를 줄이는 제3장의 교훈을 빌드 시점에 적용한 것이다.

\begin{tbbfigure}
  \adjustbox{max width=0.94\linewidth,max totalheight=0.46\textheight}{\includegraphics{figures/fig4-4_dockerfile.png}}\\[3pt]
  \figcaption{그림 4-4  설계된 Dockerfile. 다단계 빌드는 1.9 GB 도구 모음을 610 MB 결과물과 분리하고, 순서는 캐시를 보호하며, CMD와 USER 줄은 제3장의 교훈을 빌드에 얼려 넣는다.}
\end{tbbfigure}

\subsection*{GPU판: 가격표가 붙은 CUDA 태그 동물원}

GPU 서빙 이미지도 같은 규율을 따르지만 중요한 선택이 하나 더 있다. \textbf{어떤 CUDA 베이스를 쓸 것인가?} NVIDIA는 CUDA 버전마다 세 형태를 제공한다. \code{base}에는 드라이버와 맞닿는 라이브러리만 있고, \code{runtime}에는 CUDA 런타임과 cuDNN처럼 프레임워크를 \textit{실행}할 때 필요한 라이브러리가 더해진다. \code{devel}에는 nvcc, 헤더, 정적 라이브러리처럼 \textit{컴파일}할 때 필요한 약 8 GB의 도구가 추가된다. "개발 환경에서 됐으니까"라는 이유로 \code{devel}을 프로덕션에 보내는 실수가 반복되어 GB 단위 낭비를 만든다. 다단계 패턴은 올바른 선택을 기계적으로 만든다. \code{devel} 빌더 단계에서 커스텀 커널을 만들고, 최종 단계에는 \code{runtime}을 보낸다. 다음 계산 예시에서 차이에 가격표를 붙여 보자.

\begin{calloutbox}{tbbamber}{tbbamberfill}
{\normalsize\bfseries\color{tbbamberdark}계산 예시 — pull 시간과 비용으로 따진 이미지 다이어트}\par\vspace{3pt}
GPU 서빙 이미지를 세 방식으로 만든다. 용량은 대표값이며 제3장에서 계산한 유효 속도 100–250 MB/s로 pull한다.

\begin{tbbitemize}
  \item \textbf{단순형:} cuda:-devel 위의 단일 단계, pip보다 COPY가 먼저 오는 역전 구조로 \textbf{9.8 GB}다. 새 노드의 pull은 40–100 s가 걸리고, 역전 때문에 의존성 계층이 무효화되어 \textit{모든 코드 배포}마다 4 GB를 다시 보낸다.
  \item \textbf{순서 개선형:} 같은 베이스에 계층 순서만 바로잡는다. 첫 pull은 여전히 9.8 GB지만 코드 배포는 약 40 KB만 움직인다. 두 줄을 옮겨 문제의 절반을 해결했다.
  \item \textbf{설계형:} 다단계 빌드의 최종 단계를 cuda:-runtime으로 만든 \textbf{3.1 GB} 이미지다. 새 노드 pull은 12–31 s다. 12주 차의 콜드 스타트 항목은 약 3분의 1이 되고, 이미지 버전마다 노드 디스크를 약 7 GB 절약하며, 레지스트리 저장소와 송신 비용도 비례해 줄어든다.
\end{tbbitemize}

플릿 규모에서는 효과가 곱해진다. 새 노드 20대 × (9.8 − 3.1) GB = \textbf{플릿 전체의 콜드 pull 한 번당 레지스트리 트래픽 134 GB 절감}이다. 오토스케일링으로 노드를 늘릴 때마다 약 3배 빠르게 끝난다. 서버도 모델도 하드웨어도 같고 Dockerfile 하나만 다르다.
\end{calloutbox}

\vspace{3.0pt}

마지막으로 경계를 하나 그어야 한다. 결과물에 무엇을 넣을지 결정하는 경계다. 이미지는 runtime과 cuDNN 같은 CUDA \textbf{사용자 공간}을 담지만 \textbf{드라이버}는 절대 담지 않는다. 드라이버는 호스트의 속성이며 시작할 때 주입된다(§4.5). 제3장과 제6장에서 보았듯 \textbf{모델 가중치도 담지 않는다.} 모델 저장소에 두고 부팅할 때 불러온다. 그러면 이미지에는 코드, 의존성, 설정만 남는다. 코드와 같은 주기로 바뀌고 코드와 같은 리뷰 절차를 거쳐야 하는 부분, 바로 이미지가 담아야 할 내용이다.

\begin{calloutbox}{tbbpurple}{tbbpurplefill}
{\normalsize\bfseries\color{tbbpurple}생각해 보기}\par\vspace{3pt}
\begin{tbbitemize}
  \item 동료가 "모든 파일을 사용할 수 있게" Dockerfile 세 번째 줄에 COPY . /app을 넣는 빠른 수정을 제안했다. 빌드 시간, push 크기, 노드 pull, 12주 차 오토스케일러에 미치는 결과를 차례로 추적한 뒤, 두 줄을 옮기고 .dockerignore를 추가해 고쳐라.
  \item exec 형식 CMD 한 줄이 제3장, 제5장, 제6장에서 모두 보상받는 이유는 무엇인가? 각 장의 메커니즘을 이름 붙이고, 셸 형식을 쓰면 롤아웃·드레인(drain)·로드 밸런서에서 구체적으로 무엇이 깨지는지 설명하라.
  \item 다단계 빌드는 빌드 시점과 실행 시점 의존성을 분리한다. 한 계층 위의 모델 서빙에서 대응하는 사례를 들어라. §6.3의 모델 저장소는 무엇을 분리하며, "학습 컨테이너가 트래픽을 서빙한다"가 "devel 이미지를 배포한다"와 같은 실수인 이유는 무엇인가?
\end{tbbitemize}
\end{calloutbox}

\section{레지스트리: 콘텐츠 주소 기반 배포}

레지스트리는 겉보기에는 맥 빠질 만큼 단순하다. OCI distribution-spec을 구현한 HTTP 서비스로, 계층·설정·매니페스트 블롭을 다이제스트별로 저장했다가 돌려준다. Docker Hub, ECR, GCR, GHCR, 실습에서 실행할 \code{registry:2} 컨테이너는 \textit{같은 API}를 노출한다. 그래서 어느 곳에서든 \code{docker push}와 kubelet의 이미지 pull이 똑같이 작동한다. 블롭 자체는 보통 객체 스토리지에 놓이며 제6장에서 이 사실을 하나의 계약으로 다룬다. 흥미로운 엔지니어링은 저장소가 아니라 \textbf{대화}에 있다. push와 pull은 방향만 반대인 같은 협상이다. \textit{"여기 매니페스트가 있다. 이 중 없는 다이제스트가 무엇인가?"} 없는 블롭만 이동한다. 중복 제거는 최적화가 아니라 프로토콜 자체다(그림 4-5).

\begin{terminalbox}
\begin{Verbatim}[breaklines=true,breakanywhere=true,fontsize=\small,formatcom=\color{tbbtermfg},codes=\tbbcodeactive]
$ docker tag model-server:v1 localhost:5000/team/model-server:v1
$ docker push localhost:5000/team/model-server:v1
a1b2c3d4e5f6: Pushed        # 120MB  base (first push: everything moves)
b2c3d4e5f6a1: Pushed        # 1.9GB  deps
c3d4e5f6a1b2: Pushed        # 40kB   serve.py
v1: digest: sha256:aa41f9... size: 2415
# now change one line of serve.py, rebuild as :v2, push again:
$ docker build -q -t localhost:5000/team/model-server:v2 . && \
    docker push localhost:5000/team/model-server:v2
a1b2c3d4e5f6: Layer already exists
b2c3d4e5f6a1: Layer already exists
f6a1b2c3d4e5: Pushed        # 40kB - the only bytes that moved
v2: digest: sha256:bc52a0... size: 2415
\end{Verbatim}
\end{terminalbox}
\figcaption{터미널 기록 4-6  프로토콜이 작동하는 모습. 다른 다이제스트가 모두 "이미 존재"하므로 두 번째 push는 40 KB만 움직인다. pull도 대칭이며, 따라서 일상적인 배포 비용은 낮지만 콜드 노드의 비용은 높다.}

pull 쪽을 보면 경제성이 완성된다. 노드는 containerd 스냅숏 형태의 \textbf{로컬 계층 캐시}를 유지한다. \code{v1}을 서비스하던 노드에 \code{v2}를 pull하면 40 KB 계층 하나만 가져온다. 그래서 제5장의 롤링 업데이트(rolling update)는 이미 따뜻한 플릿에서 조용히 지나간다. 비싼 경우는 \textbf{새 노드}다. CUDA 크기의 이미지는 모든 것을 한 번 내려받아야 하므로 제3장에서 계산한 25–60 s가 든다. 병적인 경우는 \textbf{풀 폭풍}이다. 확장이나 플릿 전체 배포가 여러 새 노드를 동시에 건드려 GB에 노드 수를 곱한 트래픽을 하나의 레지스트리 경로로 보낸다. 가장 급한 순간에 일어난다. 12주 차에는 오토스케일링의 콜드 스타트 세금으로 다시 만난다. 대응은 §4.3의 이미지 다이어트, 노드 이미지 캐시와 사전 pull, 같은 지역의 풀스루 미러, 더 나아가 지연 로딩(lazy loading) 스냅쇼터 순이다. 터미널 기록이 push 끝마다 출력하는 \textbf{매니페스트 다이제스트}도 보자. 이 문자열이 배포 가능한 정체다. 태그는 움직이므로 제5장의 매니페스트는 다이제스트(\code{image@sha256:...})나 불변 태그를 참조해야 한다. "v2라는 이름의 \textit{무언가}를 배포했다"는 감사 기록이 아니다. 13주 차 파이프라인이 환경 사이에서 다이제스트를 승격하는 이유다.

\begin{tbbfigure}
  \adjustbox{max width=0.94\linewidth,max totalheight=0.46\textheight}{\includegraphics{figures/fig4-5_registry.png}}\\[3pt]
  \figcaption{그림 4-5  콘텐츠 주소 기반 배포. push와 pull은 없는 다이제스트만 움직이고, 노드 캐시는 코드 배포를 KB 단위로 만든다. 새 노드와 풀 폭풍은 별도로 설계해 피해야 할 비싼 예외다.}
\end{tbbfigure}

\subsection*{장애로 이어지는 운영상의 세부 사항}

\begin{tbbitemize}
  \item \textbf{인증.} \code{docker login}은 자격 증명을 저장하고, Kubernetes에는 같은 자격 증명을 \code{imagePullSecrets} 참조로 건네야 한다. 제5장의 ImagePullBackOff 행에 적힌 첫 번째 원인이다. 클라우드 레지스트리는 인스턴스/파드 IAM과 통합한다. 잘 구성한 클러스터의 노드는 정적 시크릿 없이도 비공개 이미지를 pull한다. 제6장의 키 없는 패턴이 다시 등장한다.
  \item \textbf{속도 제한.} Docker Hub의 익명 pull 제한은 작은 클러스터도 배포 도중 할당량을 다 쓸 만큼 엄격하다. \code{toomanyrequests: You have reached your pull rate limit}는 실제 장애 유형이다(§4.6). 구조적 해법은 "더 열심히 로그인"하는 것이 아니다. 의존하는 공개 베이스를 자체 레지스트리에 미러링하자. \code{registry:2}의 플래그 하나로 풀스루 캐시를 만들 수 있고, 업스트림 장애와 삭제에도 견딘다.
  \item \textbf{위생.} 태그는 쌓이고, 태그 없는 매니페스트와 주인 없는 블롭은 가비지 컬렉션 전까지 남는다. 레지스트리에 보존 정책이 있는 데는 이유가 있다. "GC를 한 번도 하지 않아 레지스트리 비용이 두 배가 됐다"는 분기마다 실제로 마주치는 황당한 비용 증가다. 보존 정책에 push 시 CVE를 검사하는 \textbf{스캐닝}을 함께 적용하자. 13주 차 파이프라인은 둘 모두를 관문으로 삼는다.
  \item \textbf{아키텍처.} 저장소는 같은 태그가 amd64와 arm64 이미지로 각각 해석되는 \textbf{다중 아키텍처 매니페스트 목록}을 담을 수 있다. Apple Silicon 노트북에서 \code{--platform} 없이 빌드하면 클라우드 노드가 arm64 전용 이미지를 pull하는 순간 §4.6의 \code{exec format error}를 만난다. \code{docker buildx --platform linux/\allowbreak amd64,linux/\allowbreak arm64}로 둘 다 게시하자. 노트북이 아니라 CI가 맡아야 할 일이다.
\end{tbbitemize}

\begin{calloutbox}{tbbpurple}{tbbpurplefill}
{\normalsize\bfseries\color{tbbpurple}생각해 보기}\par\vspace{3pt}
\begin{tbbitemize}
  \item 터미널 기록 4-6의 동작을 이용해 계산해 보자. 올바르게 계층화한 이미지를 따뜻한 노드 20대에 하루 6회 배포한다. 하루 레지스트리 트래픽을 추정한 뒤 §4.3의 COPY 순서 역전 버그가 있을 때 다시 계산하라. 제3장의 생각해 보기 문제에 배포 비용까지 더한 셈이다.
  \item 풀스루 미러는 같은 지역에 공개 이미지를 캐시한다. 서로 다른 세 가지 실패/비용 유형 가운데 무엇을 막아 주는지 나열하고, 새로 생기는 실패 하나도 찾아라.
  \item 13주 차 파이프라인은 환경마다 이미지를 다시 태깅("staging" → "prod")하지 않고 다이제스트를 승격해야 한다. 각 접근에서 생기는 구체적인 경합과 롤백(rollback) 시나리오를 제시하라.
\end{tbbitemize}
\end{calloutbox}

\section{GPU를 상자 안에 넣는 법}

이제 제2장부터 아껴 둔 질문, 강의계획서의 괄호 안 질문을 다룬다. \textit{cgroup은 왜 GPU 연산을 스케줄링할 수 없으며, 그렇다면 컨테이너는 어떻게 GPU를 사용하는가?} 먼저 컨테이너의 본질을 떠올리자(제3장). 프로세스 주위에 네임스페이스와 cgroup을 두지만 \textbf{아무것도 가상화하지 않는다.} 따라서 제2장의 virtio가 디스크를 VM 안에 에뮬레이션한 것처럼 GPU를 컨테이너 안에 흉내 낼 수 없다. 진짜 장치를 건네야 한다. 그 "진짜 장치"도 단일 부품이 아니라 단층선이 있는 \textit{스택}이다. 호스트에는 \textbf{커널 모듈}(\code{nvidia.ko})과 버전이 정확히 일치하는 \textbf{드라이버 사용자 공간 라이브러리}가 있고, 그 위에는 애플리케이션이 링크하는 런타임·cuDNN·프레임워크 등의 \textbf{CUDA 사용자 공간}이 있다. 이 단층선이 ML 이미지를 운반할 수 있게 하는 패키징 규칙을 결정한다. \textbf{이미지는 CUDA 사용자 공간을 담고, 호스트는 드라이버를 담으며, 이미지 안의 어떤 것도 특정 드라이버 빌드에 의존해서는 안 된다}(그림 4-6).

\begin{tbbfigure}
  \adjustbox{max width=0.94\linewidth,max totalheight=0.46\textheight}{\includegraphics{figures/fig4-6_gpu.png}}\\[3pt]
  \figcaption{그림 4-6  GPU는 생성 시점에 들어온다. 툴킷의 OCI 훅이 장치 노드와 호스트의 드라이버 사용자 공간을 주입하고, 이미지에는 CUDA 사용자 공간만 들어 있다. cgroup은 문을 지키고 드라이버는 연산을 맡는다.}
\end{tbbfigure}

\subsection*{메커니즘: OCI 훅과 세 가지 주입}

\textbf{NVIDIA Container Toolkit}은 이 결합을 OCI 계층의 훅으로 구현한다. shim 런타임인 \code{nvidia-container-runtime}이 runtime-spec 묶음을 받아 수정한 뒤 평범한 runc에 넘긴다. 컨테이너 생성 시점, PID 1이 존재하기 전에 세 가지를 주입한다. 첫째는 \textbf{장치 노드}다. \code{/\allowbreak dev/\allowbreak nvidia0}, \code{/\allowbreak dev/\allowbreak nvidiactl}, \code{/\allowbreak dev/\allowbreak nvidia-uvm}이 컨테이너에 나타나고 cgroup의 \textit{장치 규칙}이 이를 허용한다. GPU 이야기 전체에서 cgroup이 관여하는 부분은 이것뿐이다. 파일을 열 수 있는지 허용/거부하는 일로, 제3장의 "접근 관문"이 구체화된 모습이다. 둘째, \textbf{드라이버 사용자 공간 라이브러리}를 호스트에서 바인드 마운트한다. \code{nvidia.ko}와 버전이 맞아야 하는 절반을 \code{nvidia.ko}와 같은 머신에서 가져오는 것이 핵심이다. 셋째, 환경 변수 \code{NVIDIA\_VISIBLE\_DEVICES}가 컨테이너에 보일 GPU를 선택한다. GPU 4개가 있는 노드가 서로 다른 컨테이너에 서로 다른 장치를 건네는 방식이다.

\begin{terminalbox}
\begin{Verbatim}[breaklines=true,breakanywhere=true,fontsize=\small,formatcom=\color{tbbtermfg},codes=\tbbcodeactive]
$ docker run --rm --gpus '"device=0"' nvidia/cuda:12.4-runtime nvidia-smi \
    --query-gpu=index,name,memory.total --format=csv,noheader
0, NVIDIA A10G, 23028 MiB
$ docker inspect $(docker ps -lq) --format '{{json .HostConfig.DeviceRequests}}'
[{"Driver":"nvidia","Count":0,"DeviceIDs":["0"],"Capabilities":[["gpu"]]}]
# and the injected mounts, visible like any other config:
$ docker inspect $(docker ps -lq) | grep -c "nvidia"
23        # device nodes + driver libs + env - all inspectable
# the failure mode, staged: an image whose CUDA is newer than the host driver:
$ docker run --rm --gpus all nvidia/cuda:99.9-runtime python3 -c "..."
CUDA driver version is insufficient for CUDA runtime version
# the version contract of Figure 4-6, breached on purpose (sec. 4.6).
\end{Verbatim}
\end{terminalbox}
\figcaption{터미널 기록 4-7  GPU 접근은 마법이 아니라 설정이다. HostConfig의 장치 요청, 주입된 노드와 마운트를 모두 inspect할 수 있다. 이미지의 CUDA가 호스트 드라이버보다 앞서면 고전적인 오류가 발생한다.}

\subsection*{Kubernetes 쪽 절반과 질문의 완전한 답}

클러스터에서는 발견과 할당이 한 계층 위로 올라간다. GPU 노드의 DaemonSet인 \textbf{디바이스 플러그인(device plugin)}은 kubelet과 gRPC로 통신한다. 노드 용량에 \code{nvidia.com/\allowbreak gpu: 4}를 알리고, GPU를 요청한 파드가 배치되면 UUID로 어떤 물리 장치를 허용할지 kubelet에 알려 준다. 이는 제5장의 스케줄러가 예산으로 사용한 바로 그 숫자다. 이어 kubelet의 CRI 호출이 방금 Docker에서 본 것과 같은 툴킷 주입을 촉발한다. 아래쪽에서는 새로운 일이 없다. Kubernetes는 변하지 않은 메커니즘 위에 재고 관리와 배치를 더한다. 제5장 §5.5의 "확장 리소스는 정수다"라는 문장의 기계실을 이제 들여다보았다.

이제 강의계획서가 물었던 \textit{왜}가 남는다. 커널은 자신이 볼 수 있는 것을 스케줄링한다. 실행 대기열을 소유하므로 CPU 시간을 보고 \code{cpu.max}로 버스트 중인 작업도 멈출 수 있다. 페이지 테이블을 소유하므로 메모리를 보고 \code{memory.max}로 할당을 거부할 수 있다. GPU 연산은 설계상 커널에 보이지 않는다. 사용자 프로세스는 \textbf{드라이버가 관리하는 명령 대기열}로 작업을 제출한다. 이 대기열은 사용자 공간에 직접 매핑되어 제출할 때 시스템 호출조차 하지 않는 경우가 많다. 다음에 실행할 작업은 GPU 자체의 엔진과 펌웨어가 고른다. 커널이 보는 GPU 전체는 열기를 허용하거나 거부할 수 있는 장치 파일일 뿐이다. cgroup 컨트롤러가 측정할 만한, 커널에 보이는 "GPU 시간" 단위가 없으므로 GPU 컨트롤러도 없다. 이 한 문장에서 이후 모든 내용이 따라온다. 공유는 대기열을 \textit{볼 수 있는} 드라이버와 하드웨어, 곧 \textbf{공급업체 계층}에서 제공해야 했다(제2장의 시간 분할, MPS, MIG). Kubernetes는 아래에서 강제할 수 없는 것을 약속하지 않으므로 GPU를 \textbf{나눌 수 없는 정수}로 스케줄링한다(제5장). 메모리에서 누린 OOM식 안전망은 VRAM에는 없다. 전부 쓰면 커널과 무관하게 제3장의 \textit{다른} OOM, 곧 잡아서 처리할 수 있는 CUDA 할당자 오류를 받는다.

\begin{calloutbox}{tbbblue}{tbbbluefill}
{\normalsize\bfseries\color{tbbblue}네 장에 걸친 GPU 이야기를 한곳에 모으기}\par\vspace{3pt}
\begin{tbbitemize}
  \item \textbf{제2장:} 하드웨어를 공유할 수 있지만 PCIe 패스스루, vGPU, 시간 분할, MPS, MIG 같은 공급업체 메커니즘만 가능하다.
  \item \textbf{제3장:} 커널은 GPU를 중재할 수 없다. cgroup GPU 컨트롤러는 없고 장치 파일 접근만 통제한다.
  \item \textbf{제4장(현재):} 따라서 접근은 \textit{주입}이다. OCI 훅이 생성 시점에 호스트 드라이버와 장치 노드를 이미지의 CUDA 사용자 공간에 결합하고, 디바이스 플러그인이 발견과 할당을 맡는다.
  \item \textbf{제5장:} 오케스트레이션은 GPU를 정수 확장 리소스로 다룬다. 분수도 오버커밋도 없고 requests = limits다. 분할 수요는 제2장의 메커니즘으로 돌아가며 12주 차에 운영 방법을 다룬다.
\end{tbbitemize}
\end{calloutbox}

\vspace{3.0pt}

\begin{calloutbox}{tbbpurple}{tbbpurplefill}
{\normalsize\bfseries\color{tbbpurple}생각해 보기}\par\vspace{3pt}
\begin{tbbitemize}
  \item 호스트의 드라이버 라이브러리를 이미지에 넣는 방식, 곧 "내 노드에서는 됐는데"가 왜 이식성 시한폭탄이 되는가? 이미지는 고정된 채 플릿의 노드 풀이 드라이버 버전을 올릴 때 벌어지는 일을 추적하라.
  \item NVIDIA\_VISIBLE\_DEVICES는 주입으로 강제한다. 허용받지 않은 장치는 컨테이너에 아예 보이지 않는다. 이를 제3장의 memory.max 강제 모델과 비교하라. 각 방식에서 가능한 실패는 무엇인가? 두 컨테이너에 \textit{같은} GPU를 허용하면 어떤 일이 생길지 생각해 보자.
  \item 이 절과 제2장을 바탕으로 "GPU 시간용 cgroup 컨트롤러를 만들자. 얼마나 어렵겠어?"라는 제안을 두 문단으로 기술적으로 반박하라. 커널이 보아야 하지만 볼 수 없는 것은 무엇이며, 측정은 대신 어느 계층에 있어야 하는가?
\end{tbbitemize}
\end{calloutbox}

\section{검시관처럼 컨테이너 읽기: inspect 현장 안내서}

징후, 원인, 첫 명령으로 정리하는 이 책의 현장 안내서 습관은 제3장의 종료 코드 하나에서 시작해 장마다 한 계층씩 자란다. Docker 계층에서는 도구 하나가 중심이다. \textbf{docker inspect}는 무엇을 \textit{실제로} 설정했고 무엇이 \textit{실제로} 일어났는지 플랫폼이 JSON으로 제출하는 선서 증언이다. 선언한 이미지·환경·상한·장치와 스택이 기록한 상태·종료 코드·OOM 플래그·재시작 횟수가 모두 들어 있다. 어떤 질문에 어느 필드가 답하는지 아는 것이 기술이다. 실습의 필수 연습이 가장 좋은 도입이다. 제3장의 첫 터미널 기록에서 시작한 고리를 닫기 때문이다.

\begin{terminalbox}
\begin{Verbatim}[breaklines=true,breakanywhere=true,fontsize=\small,formatcom=\color{tbbtermfg},codes=\tbbcodeactive]
$ docker run -d --name demo -m 512m --cpus 1 model-server:v1
$ docker inspect demo --format \
    'mem={{.HostConfig.Memory}} nanocpus={{.HostConfig.NanoCpus}}'
mem=536870912 nanocpus=1000000000     # your flags, recorded verbatim
$ PID=$(docker inspect demo --format '{{.State.Pid}}') && echo $PID
41320
$ CG=$(cat /proc/$PID/cgroup | cut -d: -f3) && \
    cat /sys/fs/cgroup$CG/memory.max /sys/fs/cgroup$CG/cpu.max
536870912                             # the SAME number, in the kernel
100000 100000                         # --cpus 1, in Ch. 3 grammar
# intent (inspect) == enforcement (cgroupfs): the ghostwriter, audited.
# this is Transcript 3-1 closed from both ends - and the lab's Part B.
\end{Verbatim}
\end{terminalbox}
\figcaption{터미널 기록 4-8  강의계획서가 요구한 감사. docker inspect는 기록된 의도를, cgroup 파일 시스템은 커널의 강제를 보여 준다. 제3장의 대필자가 이 장의 스택이므로 숫자가 일치한다.}

여섯 필드가 질문의 90\%에 답하므로 작은 JSON 지도를 기억하자. \textbf{.State}는 실행 중인지, 아니라면 \textit{어떻게 끝났는지} 알려 준다. \code{ExitCode}, \code{OOMKilled}, \code{Error}, 타임스탬프가 있다. 죽은 컨테이너는 서류가 딸린 시신이며 \code{docker logs}도 여전히 작동한다. 제5장의 \code{--previous} 비행 기록 장치 습관이 한 계층 아래에 있다. \textbf{.Config}는 \textit{이미지}가 말한 env, \code{Cmd}, \code{User}, 노출 포트다. \textbf{.HostConfig}는 \textit{운영자}가 말한 상한, 재시작 정책, 장치, 마운트, 포트 게시다. 의도와 현실이 다를 때 의도는 여기에 적혀 있다. \textbf{.GraphDriver}는 제3장의 lower/upper/merged overlayfs 디렉터리를 실제 \code{ls}할 수 있는 경로로 보여 준다. \textbf{.NetworkSettings}에는 IP, 포트, §3.2의 veth 이야기가 있다. \textbf{.RootFS.Layers}는 다이제스트 목록이며 태그가 배신했을 때 \textit{실제로 어떤 바이트}를 실행하는지 확인한다. 지도를 들면 반복되는 실패가 스스로 분류된다(그림 4-7).

\begin{tbbfigure}
  \adjustbox{max width=0.94\linewidth,max totalheight=0.46\textheight}{\includegraphics{figures/fig4-7_fieldguide.png}}\\[3pt]
  \figcaption{그림 4-7  빌드와 실행 현장 안내서. 아키텍처 불일치, 속도 제한, 사라진 태그, CUDA 버전 계약, Docker식 OOM, runc의 규격 거부를 각각 첫 명령과 함께 정리한다.}
\end{tbbfigure}

\begin{tbbitemize}
  \item \textbf{즉시 나타나는 exec format error}: 이 CPU에서 바이너리를 실행할 수 없다. Apple Silicon 노트북이 기본으로 만든 arm64 이미지를 amd64 노드에서 실행한 경우다. 2021년 이후 노트북에서 클라우드로 옮길 때 가장 흔한 놀라움이다. \code{docker image inspect --format '\{\{.Architecture\}\}'}로 확인하고 CI의 \code{buildx --platform linux/\allowbreak amd64}로 고치거나 다중 아키텍처를 게시한다. 옆에 \code{no matching manifest for linux/\allowbreak amd64}도 분류하자. 같은 병을 실행 시점이 아니라 pull 시점에 잡은 것이다.
  \item \textbf{pull 중 toomanyrequests}: 플릿이 배포 도중 Docker Hub의 익명 속도 제한을 소진했다. 우선 인증으로 급한 불을 끄고, 장기적으로는 의존하는 공개 베이스를 미러링한다(§4.4). 배포 주기를 제3자의 속도 제한과 결합해서는 안 된다.
  \item \textbf{manifest unknown}: 태그가 없다. 오타, 삭제되었거나 GC된 태그, 빌드가 push하기 전에 배포한 파이프라인이 원인이다. 지속 가능한 해결책은 §4.3의 다이제스트 배포다. 다이제스트는 존재하거나 처음부터 존재하지 않았을 뿐 절반만 push될 수 없다.
  \item \textbf{CUDA driver version is insufficient for CUDA runtime version}: 그림 4-6의 버전 계약을 어겼다. 이미지의 CUDA 사용자 공간이 호스트 드라이버보다 새롭다. 오래된 노드 풀에 "베이스 이미지만 올린" 변경을 배포했을 때 흔하다. 호스트에서 \code{nvidia-smi}가 보여 주는 드라이버/CUDA 상한과 이미지 태그를 비교하고, 우연이 아니라 계획에 따라 드라이버를 올린다.
  \item \textbf{조용한 죽음, exit 137}: 제3장 §3.6을 Docker로 포장한 모습이다. \code{.State.OOMKilled: true}는 \code{-m} 장벽이 강제되었다고 말한다. 측정한 피크에 맞춰 용량을 조정하거나 입력을 제한하는 해법은 같다. 이 플래그는 버그가 아니라 예산 문제였음을 한 명령으로 증명한다.
  \item \textbf{OCI runtime create failed}: 프로세스가 생기기도 전에 runc가 규격을 거부했다. 존재하지 않는 바인드 마운트 원본, 이미지에 없는 엔트리포인트 경로, 잘못된 cgroup 값이 원인일 수 있다. 진짜 원인은 중첩된 오류 가운데 \textit{가장 안쪽}에 있다. 전체 줄을 읽고 \code{docker inspect}로 실제 제출된 Config/HostConfig를 확인해 의도와 비교한다.
\end{tbbitemize}

한 습관이 이 절을 책 전체와 잇는다. \textbf{플랫폼은 자신이 한 일을 기록한다.} 제3장은 커널 장부(\code{/\allowbreak proc}, cgroupfs, dmesg)를 읽었고, 이 절은 런타임 장부(\code{inspect}, \code{events}, \code{logs})를 읽는다. 제5장은 오케스트레이터 장부(describe, events), 제6장은 관문의 장부(상태 코드, 대상 상태)를 읽는다. 네 계층에 규율은 하나다. 장애 대응은 각 계층에 올바른 순서로 증언을 요청하는 연습일 뿐이다.

\begin{calloutbox}{tbbpurple}{tbbpurplefill}
{\normalsize\bfseries\color{tbbpurple}생각해 보기}\par\vspace{3pt}
\begin{tbbitemize}
  \item 컨테이너가 docker run -it에서는 완벽히 실행되지만 Kubernetes에서는 exit 127("not found")로 즉시 죽는다. .Config/.HostConfig 구분을 이용해 대화형 실행과 매니페스트 제출 사이에서 가장 가능성 높은 차이 두 가지를 추론하라. 엔트리포인트 덮어쓰기와 환경 변수 의존 경로를 생각해 보자.
  \item 그림 4-7의 행 가운데 배포 전에 막을 수 있는 모든 실패를 예방하는 다섯 줄짜리 CI 관문을 설계하라. 원칙상 CI에서 잡을 수 없는 두 행은 무엇이며, 각각 어떤 런타임 검사가 대신 다루는가?
  \item docker inspect에는 .HostConfig.Memory = 536870912라고 나오지만 컨테이너가 800 MB를 문제없이 사용한다. 버그라고 부르기 전에 제3장에서 배운 세 가지 설명을 확인 순서대로 나열하라. 실제로 강제하는 cgroup, 페이지 캐시와 익명 메모리의 차이, 실행 중 cgroup을 바꾼 주체를 생각해 보자.
\end{tbbitemize}
\end{calloutbox}

\section{실습: 모델 서버를 컨테이너화하고 push한 뒤 검증하기}

이번 주 실습에서는 이 책의 중심 결과물을 만든다. FastAPI 모델 서버를 잘 설계한 이미지로 만들고 레지스트리에 저장한 뒤 cgroup 파일까지 검증한다. 제5장이 Kubernetes에 배포하고 제6장이 모델 저장소에 연결할 바로 그 이미지다. 모든 작업은 제2장의 VM이나 노트북에서 실행한다. 레지스트리는 무료이고 즉시 실행되며 실제 서비스와 같은 API를 쓰는 로컬 \code{registry:2} 컨테이너다. GPU 부분은 하드웨어가 있는 팀을 위한 선택 부록이다. 제출물은 아래 결과물 체크리스트와 관찰마다 관련 절을 표시한 한 페이지 보고서다.

\subsection*{A부: 제대로 빌드하고 계층의 동작을 증명하기}

\code{serve.py}를 작성한다. 유인물에는 제5장과 제6장에서 프로브할 \code{/\allowbreak predict}, \code{/\allowbreak healthz}, \code{/\allowbreak ready} 엔드포인트를 가진 FastAPI 앱 뼈대가 있다. 두 번 컨테이너화한다. 먼저 단일 단계, 전체 베이스, \code{pip install}보다 앞선 \code{COPY . .}를 쓰는 \textbf{단순 빌드}를 만든다. 다음은 터미널 기록 4-5의 다단계 Dockerfile과 \code{.dockerignore}를 쓴 \textbf{설계 빌드}다. 각각 이미지 크기(\code{docker images}), 계층 지도(\code{docker history}), 순서의 효과를 체감하게 할 측정값인 \textbf{serve.py 한 줄 변경 뒤 재빌드 시간과 "이동할 바이트"}를 기록한다. 둘 다 다시 빌드해 단순 빌드는 pip를 재실행하지만 설계 빌드는 캐시 계층을 모두 재사용해 몇 초 만에 끝나는 모습을 확인한다. \code{docker image inspect --format '\{\{.Architecture\}\}'}로 아키텍처도 명시적으로 확인한다. Apple Silicon 팀은 \code{--platform linux/\allowbreak amd64}를 추가하고 에뮬레이션 비용을 보고서에 적는다.

\subsection*{B부: 실행하고 제한하고 감사하기(강의계획서의 inspect 연습)}

설계 이미지를 실제 예산 \code{-m 512m --cpus 1 -p 8080:8000}으로 실행하고, 자신의 컨테이너에서 터미널 기록 4-8 전체를 재현한다. \code{inspect}로 기록된 의도를 확인하고 \code{/\allowbreak proc/\allowbreak <pid>/\allowbreak cgroup}에서 cgroup 디렉터리로 이동해 \code{memory.max}와 \code{cpu.max}가 플래그와 일치하는 화면을 캡처한다. 이제 제3장이 연 플래그 → 데몬 → containerd → runc → 커널 파일의 고리를 직접 닫았다. 이어 두 가지를 실시간으로 검증한다. \textbf{PID 1 정확성}은 \code{docker stop} 시간을 재서 확인한다. exec 형식 CMD는 SIGTERM을 받아 1초 안에 드레인하고 끝나지만, 의도적으로 만든 셸 형식 버전은 10초의 유예 시간을 꽉 채워 버티다가 SIGKILL된다. 두 시간을 모두 기록한다. \textbf{Docker식 OOM}은 \code{-m 128m}으로 다시 실행해 장난감 추론 루프를 137이 날 때까지 구동한 뒤 \code{.State.OOMKilled: true}를 수집한다. 제3장의 증거 사슬을 이 장의 도구로 확인한 것이다.

\subsection*{C부: 자체 레지스트리에서 프로토콜 관찰하기}

\code{docker run -d -p 5000:5000 --name registry registry:2}를 실행하고 설계 이미지에 태그를 붙여 push한 뒤 계층별 출력을 읽는다. 이어 프로토콜을 가르치는 두 실험을 한다. \textbf{중복 제거}에서는 serve.py 한 줄을 바꾸고 \code{:v2}로 다시 빌드해 push한다. \code{Layer already exists}가 쌓인 가운데 작은 \code{Pushed} 하나만 나오는 화면을 캡처한다. 자신의 터미널 기록 4-6이다. \textbf{다이제스트 규율}에서는 push가 반환한 다이제스트를 각각 기록하고, 다이제스트로 pull한다(\code{docker pull localhost:5000/\allowbreak team/\allowbreak model-server@sha256:...}). 그런 다음 일부러 \code{:v2}를 \textit{다른} 이미지에 다시 붙여 태그는 거짓말할 수 있지만 다이제스트는 그렇지 못함을 증명하고 이를 한 문장 정책으로 쓴다. §6.3의 모델 저장소에서 이미 아는 규칙이 어디서 왔는지 확인하는 순간이다. 마지막으로 현장 안내서의 두 행을 일부러 재현한다. 없는 태그를 pull해 \code{manifest unknown}을 만들고, 팀에 ARM 머신이 있다면 플랫폼을 지정하지 않은 빌드를 x86에서 실행해 \code{exec format error}를 만든다. 화면을 보고서에 넣는다.

\begin{calloutbox}{tbbpurple}{tbbpurplefill}
{\normalsize\bfseries\color{tbbpurple}실습 체크리스트}\par\vspace{3pt}
\begin{tbbitemize}
  \item A부 결과물: 단순형과 설계형의 크기·계층 표, 한 줄 변경 뒤 재빌드 시간 두 개, 아키텍처 확인.
  \item B부 결과물: inspect와 cgroupfs의 일치(자신의 터미널 기록 4-8), exec 형식과 셸 형식의 docker stop 시간, -m 128m에서 OOMKilled: true 캡처.
  \item C부 결과물: 첫 push와 두 번째 push의 기록("Layer already exists"), 다이제스트 pull 증명과 태그 배신 시연, 일부러 재현한 현장 안내서 두 행.
  \item 보고서: 한 페이지. 각 관찰에 관련 절을 적는다. 예: "두 번째 push는 40 KB를 옮겼다. §4.4의 프로토콜이다. 단순 이미지는 1.9 GB를 옮겼을 것이다. §4.3의 순서 역전 때문이다."
  \item 이미지 보관: 최종 빌드에 model-server:v1 태그를 붙이고 디스크가 허용하면 레지스트리를 계속 실행한다. 5주 차 실습에서 바로 이 결과물을 배포한다. 정리해야 한다면 docker rm -f와 docker system prune을 사용하되 먼저 삭제 대상을 읽는다.
\end{tbbitemize}
\end{calloutbox}

\vspace{3.0pt}

\begin{calloutbox}{tbbpurple}{tbbpurplefill}
{\normalsize\bfseries\color{tbbpurple}생각해 보기}\par\vspace{3pt}
\begin{tbbitemize}
  \item B부에서 정지 시간을 재니 exec 형식은 약 0.3 s, 셸 형식은 약 10 s 만에 끝났다. 이 두 숫자를 제5장의 롤링 업데이트와 제6장의 로드 밸런서 드레이닝에 적용하라. Dockerfile 한 줄이 이후 모든 배포에서 사용자에게 어떤 차이를 만드는가?
  \item 실습 레지스트리는 인증도 TLS도 없는 컨테이너 하나다. 프로덕션이라 부르기 전에 더해야 할 항목을 우선순위대로 나열하고, 13주 차 파이프라인이 가장 크게 의존하는 항목도 고르라.
\end{tbbitemize}
\end{calloutbox}

\namedsection{요약}{열 가지로 정리한 제4장}

\qitem{1}{\textbf{프로덕션 컨테이너화는 더 큰 셸 스크립트가 아니라 세 산업적 문제}, 곧 패키징·배포·수명 주기를 푸는 일이다. Docker는 셋을 모두 해결했고 OCI는 그 해법을 누구나 구현할 수 있는 인터페이스로 표준화했다.}

\qitem{2}{\textbf{스택의 각 계층에는 역할이 하나뿐이다.} docker CLI(REST 클라이언트) → dockerd(API와 편의 기능) → containerd(수명 주기 엔진이며 kubelet의 직접 통신 상대) → containerd-shim(컨테이너마다 하나씩 stdio·종료 코드·회수를 담당) → runc(JSON 규격에 따라 제3장의 다섯 단계을 수행한 뒤 종료)로 이어진다. 스택이 모델을 실행하는 것이 아니라 커널이 실행한다.}

\qitem{3}{\textbf{shim은 제3장의 수수께끼 같은 부모 PID 41288을 설명한다.} 가장 작고 지루한 프로그램이 프로세스에 가장 가까이 있으므로 그 위의 모든 구성 요소가 장애·업그레이드·재시작을 겪어도 컨테이너에는 아무 일도 일어나지 않는다.}

\qitem{4}{\textbf{OCI 접합부는 모든 구성 요소를 교체 가능하게 한다.} image-spec은 빌드, distribution-spec은 배포, runtime-spec은 실행을 정의하며 config.json은 제3장을 데이터로 옮긴 것이다. Kata와 gVisor는 런타임 접합부에 들어가고 kubelet은 CRI로 containerd와 직접 통신한다. 2022년 dockershim이 제거되어도 이미지는 신경 쓰지 않았다. Docker 결과물이 아니라 OCI 결과물이기 때문이다.}

\qitem{5}{\textbf{이미지는 콘텐츠 주소 기반 결과물이다.} 매니페스트가 설정과 계층을 가리키며 모든 이름은 sha256이다. 중복 제거, 머신 간 캐시 신뢰, 변조 탐지가 이름에서 자연스럽게 나온다. 태그는 움직일 수 있는 별명이므로 배포는 다이제스트를 고정한다. §6.3 모델 저장소의 불변성 규율도 여기서 배웠다.}

\qitem{6}{\textbf{Dockerfile은 이 책의 누적된 교훈을 빌드에 얼려 넣은 파일이다.} 변경 빈도에 따른 순서(제3장의 캐시), 코드보다 앞선 requirements(순서가 뒤집히면 재빌드에 20분), .dockerignore, 다단계 빌드(빌드 의존성은 실행 페이로드가 아니다: 2.4 GB→610 MB, cuda-devel→cuda-runtime: 9.8 GB→3.1 GB), exec 형식 CMD(PID 1, 제3·5·6장에서 보상), 비루트 USER를 적용한다.}

\qitem{7}{\textbf{레지스트리는 없는 다이제스트만 움직인다.} 중복 제거는 최적화가 아니라 프로토콜이다. 따뜻한 플릿의 코드 배포는 KB만 움직이지만 새 노드와 풀 폭풍은 GB를 움직인다(12주 차의 콜드 스타트 세금). 풀 시크릿/IAM으로 인증하고, 속도 제한이 장애가 되지 않도록 공개 베이스를 미러링하며, GC와 스캐닝을 적용한다(13주 차의 관문).}

\qitem{8}{\textbf{GPU는 패키징하지 않고 주입한다.} 이미지는 CUDA 사용자 공간을, 호스트는 커널 모듈과 버전이 맞는 드라이버 사용자 공간을 담는다. NVIDIA 툴킷의 OCI 훅이 생성 시점에 장치 노드·바인드 마운트 라이브러리·NVIDIA\_VISIBLE\_DEVICES로 둘을 결합한다. 디바이스 플러그인은 같은 메커니즘 위에 클러스터 재고와 할당을 더한다(제5장의 정수 리소스).}

\qitem{9}{\textbf{커널이 GPU 연산을 볼 수 없으므로 cgroup은 이를 스케줄링할 수 없다.} 작업은 드라이버가 관리하는 명령 대기열로 들어가며 사용자 공간에 매핑되어 시스템 호출조차 없을 수 있다. 커널에는 장치 파일만 있고 연산 단위는 없다. cgroup의 역할은 접근 통제뿐이고 공유는 공급업체 메커니즘에 있다(제2장). VRAM 고갈은 커널 OOM이 아니라 잡을 수 있는 할당자 OOM이다.}

\qitem{10}{\textbf{docker inspect는 런타임의 선서 증언이다.} State는 OOMKilled와 ExitCode 등 종료 방식을, Config와 HostConfig는 이미지의 의도와 운영자의 의도를, GraphDriver는 제3장의 overlay 디렉터리를, RootFS.Layers는 실제 실행 중인 바이트를 보여 준다. 현장 안내서의 여섯 징후인 exec format error, toomanyrequests, manifest unknown, CUDA driver insufficient, 조용한 137, OCI runtime create failed에는 각각 첫 명령이 있고 실습에서 직접 재현한다.}

\namedsection{연습문제}{}

\subsection*{A. 객관식}

\qitem{1}{model-server 컨테이너 30개가 실행 중인 노드에서 Docker 데몬을 재시작했다. 컨테이너에는 무슨 일이 일어나며, 그 이유는 무엇인가?}

\begin{tbboptions}
  \item ① dockerd가 부모이므로 모두 재시작한다.
  \item ② PID 1의 자식인 장기 실행 containerd-shim이 각각 붙들고 있으므로 아무 영향 없이 계속 실행한다.
  \item ③ 데몬이 돌아올 때까지 일시 정지한다.
  \item ④ 계속 실행하지만 네트워크 네임스페이스를 잃는다.
\end{tbboptions}

\qitem{2}{runc에 관한 설명으로 옳은 것은?}

\begin{tbboptions}
  \item ① 각 컨테이너 옆에서 상주 감독자로 실행된다.
  \item ② Kubernetes가 CRI로 통신하는 데몬이다.
  \item ③ OCI runtime-spec의 config.json을 읽고 네임스페이스/cgroup/rootfs를 설정해 프로세스를 시작한 뒤 종료한다.
  \item ④ 서로 다른 아키텍처 사이에서 이미지를 변환한다.
\end{tbboptions}

\qitem{3}{serve.py 한 줄을 바꿔 다시 빌드하고 push했다. 계층 순서가 올바르면 약 40 KB만 움직이는 이유는?}

\begin{tbboptions}
  \item ① 레지스트리가 다른 계층을 압축해 없애기 때문이다.
  \item ② push가 다이제스트로 협상하며 바뀌지 않은 모든 계층은 레지스트리에 "이미 존재"하므로 새 코드 계층만 업로드하기 때문이다.
  \item ③ Docker가 소스 트리의 차이를 계산해 패치를 보내기 때문이다.
  \item ④ 태그가 포인터이므로 어떤 데이터도 두 번 움직이지 않기 때문이다.
\end{tbboptions}

\qitem{4}{Apple Silicon 노트북에서 만든 이미지가 클라우드 노드에서 "exec format error"와 함께 즉시 죽었다. 원인과 해결책은?}

\begin{tbboptions}
  \item ① push가 손상되었다. 다시 push한다.
  \item ② CUDA 버전이 맞지 않는다. 드라이버 버전을 맞춘다.
  \item ③ amd64 노드에 arm64 이미지를 배포했다. CI에서 --platform linux/amd64로 빌드하거나 다중 아키텍처를 만든다.
  \item ④ 실행 권한이 없다. Dockerfile에서 chmod +x를 실행한다.
\end{tbboptions}

\qitem{5}{이미지는 CUDA 런타임을 담지만 NVIDIA 드라이버 라이브러리는 절대 담지 않는 이유는?}

\begin{tbboptions}
  \item ① 라이선스가 재배포를 금지하기 때문이다.
  \item ② 드라이버 사용자 공간은 호스트 커널 모듈과 정확히 맞아야 하므로 생성 시 호스트에서 주입한다. 이미지에 넣으면 드라이버 업그레이드 때 깨진다.
  \item ③ 드라이버가 이미지 계층에 넣기에는 너무 크기 때문이다.
  \item ④ 드라이버는 C로 작성되었고 이미지는 Python만 담기 때문이다.
\end{tbboptions}

\qitem{6}{로그 없이 컨테이너가 죽었다. docker inspect에 ExitCode: 137, OOMKilled: true가 나온다. 올바른 해석은?}

\begin{tbboptions}
  \item ① 커널 OOM 킬러가 -m 상한을 강제했다(제3장 §3.6). 측정한 피크에 맞춰 늘리거나 입력을 제한한다.
  \item ② 애플리케이션이 세그멘테이션 오류로 죽었다.
  \item ③ 이미지가 비정상이어서 Docker가 종료했다.
  \item ④ 레지스트리가 이미지를 회수했다.
\end{tbboptions}

\subsection*{B. 단답형}

\qitem{1}{docker CLI → dockerd → containerd → shim → runc 사슬에서 각 구성 요소의 단 하나뿐인 역할을 쓰고, 실행 중인 컨테이너의 프로세스 트리에 없는 두 구성 요소와 그 이유를 설명하라.}

\qitem{2}{태그와 다이제스트를 각각 정의하고, 태그로 배포할 때의 실패 모드를 제시하며, 이 책이 한 계층 위에서 같은 불변성 규율을 적용하는 두 곳을 쓰라.}

\qitem{3}{NVIDIA Container Toolkit이 컨테이너 생성 시 수행하는 세 가지 주입을 나열하고, 그중 cgroup이 관여하는 항목과 제한된 역할을 설명하라.}

\qitem{4}{팀 이미지가 9.8 GB(cuda-devel, 단일 단계, COPY 순서 역전)다. 약 3.1 GB이면서 코드 배포는 KB 단위가 되게 하는 세 변경과 각각 구현하는 장의 교훈을 쓰라.}

\subsection*{C. 서술형}

\qitem{1}{"Docker가 오래 남긴 공헌은 데몬이 아니라 규격이다." shim 아키텍처, dockershim 제거, 런타임 접합부(Kata/gVisor)를 근거로 이 주장에 찬성하거나 반박하라. 이미지와 런타임이 한 공급업체의 독점 기술로 남았다면 ML 서빙이 어떻게 달라졌을지도 논하라.}

\qitem{2}{팀 엔지니어링 안내서의 "이미지 표준" 절을 작성하라. 베이스 이미지 정책(slim/runtime 대 devel), 계층 순서, 다단계 빌드, USER/CMD, 아키텍처, 태그/다이제스트 정책과 각각을 강제하는 CI 관문을 포함하고, 모든 규칙을 이 장 또는 제3·5·6장의 실패나 비용으로 정당화하라.}

\qitem{3}{GPU 이야기는 네 장에 걸쳐 이어진다(§4.5의 상자). 새 동료를 위해 하나의 일관된 글로 다시 설명하라. 커널이 GPU 연산을 스케줄링하지 못하는 이유, 그 결과 공급업체가 맡게 된 일(제2장), 패키징에 생긴 요구(이 장), 오케스트레이션에 생긴 요구(제5장)를 인과 관계로 연결하고, 12주 차의 공유 결정에 미칠 영향으로 마무리하라.}

\namedsection{정답과 해설}{}

\subsection*{A. 객관식}

\qitem{1}{\textbf{②} shim은 실행 중에 데몬이 하중을 떠받치는 부품이 되지 않게 하려고 존재한다. PID 1의 자식으로 stdio와 종료 코드를 붙든다. 터미널 기록 4-2에서 데몬 재시작으로 직접 증명했다. (§4.2)}

\qitem{2}{\textbf{③} runc는 상주 프로그램이 아니라 설치 작업반이다. config.json에 따라 다섯 단계을 수행한 뒤 종료한다. ②는 CRI 통신 상대인 containerd와 혼동했다. (§4.2)}

\qitem{3}{\textbf{②} 중복 제거가 곧 프로토콜이다. push는 "없는 다이제스트가 무엇인가?"라고 묻는다. 순서가 올바르면 40 KB 코드 계층만 새롭다. ③처럼 소스 차이를 보내는 것이 아니며 단위는 차이가 아니라 계층이다. (§4.3–4.4)}

\qitem{4}{\textbf{③} 아키텍처 불일치는 노트북에서 클라우드로 옮길 때 가장 흔한 문제다. 바이너리를 실행할 수 없어 애플리케이션 코드에 도달하기 전에 죽는다. --platform은 CI가 책임져야 한다. (§4.4, §4.6)}

\qitem{5}{\textbf{②} 호스트마다 커널 모듈과 드라이버 사용자 공간의 버전이 서로 고정되어 있다. 툴킷이 생성 시 호스트의 라이브러리를 바인드 마운트하므로 드라이버를 올려도 이미지는 이식성을 유지한다. (§4.5)}

\qitem{6}{\textbf{①} 137 + OOMKilled: true는 제3장의 장벽에 Docker의 서류가 붙은 것이다. 버그가 아니라 예산 문제다. 피크를 측정해 크기를 조정하거나 입력을 제한하는 절차는 같다. (§3.6, §4.6)}

\subsection*{B. 단답형}

\qitem{1}{CLI는 REST 클라이언트다. dockerd는 개발자 API, 빌드, 편의 기능을 맡는다. containerd는 이미지·스냅숏·작업의 수명 주기를 관리하며 CRI를 통한 kubelet의 통신 상대다. shim은 컨테이너별 stdio·종료 코드·회수를 장기간 맡는다. runc는 config.json의 다섯 단계을 수행하고 종료한다. 프로세스 트리에 없는 것은 dockerd와 runc다. dockerd는 처음부터 조상이 아니며 shim은 PID 1의 자식이 되고, runc는 설정 뒤 종료한다. (§4.2)}

\qitem{2}{다이제스트는 매니페스트의 sha256 정체로 구조상 불변이다. 태그는 이름→다이제스트를 잇는 가변 포인터다. 태그가 이동하거나 다시 push되면 노드와 시간에 따라 "v2"가 다른 바이트를 뜻해 감사 기록이 사라지고 캐시가 오염된다. 한 계층 위의 같은 규율은 §6.3 모델 저장소 버전(새 버전 = 새 접두사, sha256 매니페스트)과 13주 차 다이제스트 승격 파이프라인에 적용된다. (§4.3, §6.3)}

\qitem{3}{(1) \code{/dev/nvidia*} 장치 노드를 컨테이너에 넣고 cgroup \textit{장치 허용 규칙}으로 접근을 통제한다. 이것이 유일한 cgroup 역할이다. (2) nvidia.ko와 맞아야 하는 호스트 드라이버 사용자 공간 라이브러리를 바인드 마운트한다. (3) NVIDIA\_VISIBLE\_DEVICES로 컨테이너에 보일 GPU를 선택한다. 목록 어디에도 스케줄링은 없다. 바로 그 점이 핵심이다. (§4.5)}

\qitem{4}{(1) 최종 단계를 cuda:-runtime으로 한 다단계 빌드로 약 7 GB devel 도구 모음을 뺀다. 빌드 의존성은 실행 페이로드가 아니다. (2) requirements+설치를 코드 COPY보다 앞에 두어 제3장의 캐시 경제성을 적용한다. 코드 배포가 의존성 계층을 더는 무효화하지 않는다. (3) .dockerignore로 컨텍스트를 정리하고 제6장에 따라 가중치를 이미지에서 뺀다. 결과는 약 3.1 GB 이미지와 약 40 KB 배포다. (§4.3)}

\subsection*{C. 서술형(채점 기준)}

\qitem{1}{강한 찬성 답안은 shim/데몬 분리로 데몬을 교체 가능한 존재로 만든 점(Docker 자체 구성 요소도 특권적이지 않음), dockershim 제거라는 자연 실험(생태계가 데몬을 버려도 결과물은 그대로), Kata/gVisor를 받아들이는 런타임 접합부(제2장의 스펙트럼이 설정 한 줄이 됨)를 든다. 반사실적 상황에서는 공급업체 종속 이미지가 캐시/중복 제거의 경제성을 분열시키고 §4.4의 레지스트리 상호운용성을 막아 Kubernetes 생태계 자체를 늦추었을 가능성을 논한다. 반박은 데몬 시대의 사용성이 보급을 이끌었고 표준은 성공 뒤에 왔다고 지적할 수 있다. 어느 입장이든 논증을 평가한다. (§4.2)}

\qitem{2}{근거가 붙은 규칙을 찾는다. devel이 아닌 runtime(9.8 GB → 3.1 GB, §4.3), 순서+.dockerignore(20분 재빌드, KB 배포), 컴파일 의존성의 다단계 빌드 의무화, exec 형식 CMD+비루트 USER(제3·5·6장의 보상), CI의 --platform(exec format error), 배포의 다이제스트 또는 불변 태그(manifest unknown, 감사)가 포함되어야 한다. CI 관문으로 두 아키텍처 빌드, 크기 예산, 스캔, push 후 배포 순서, 다이제스트 승격을 제시한다. 공개 베이스 미러 정책(속도 제한)은 가산점이다. (§4.3–4.6)}

\qitem{3}{점수를 줄 핵심 흐름은 다음과 같다. 커널은 볼 수 있는 것만 스케줄링한다. GPU 작업은 커널에 보이지 않는 드라이버 관리 대기열로 들어간다 → cgroup 컨트롤러가 존재할 수 없다(제3장의 공백을 메커니즘으로 설명) → 공유는 공급업체 계층의 시간 분할/MPS/MIG가 맡아야 했다(제2장) → 패키징은 드라이버(호스트)와 CUDA 사용자 공간(이미지)을 나누고 주입으로 결합해야 했다(이 장) → 오케스트레이션은 정수만 약속할 수 있다(제5장) → 12주 차의 공유 결정은 공급업체 메커니즘 하나를 고르고 격리 절충을 받아들이는 일이다. 목록만 아니라 인과 화살표가 있어야 만점이다. (§4.5)}

\namedsection{더 읽을거리}{}

\textbf{[이번 주]} 표시는 지금 읽어야 할 필수 자료이고, \textbf{[N주 차]} 표시는 해당 주에 도움이 될 자료다.

\begin{tbbitemize}
  \item \textbf{[이번 주] Marc Levinson, "The Box: How the Shipping Container Made the World Smaller and the World Economy Bigger" (2006), 1장}: §4.1 사례 B의 실제 역사다. 재래식 하역의 경제성, 맥린, Ideal-X, 표준화의 정치를 다룬다. §4.1은 단순한 비유가 아니라 계보이며, 이 책에서 가장 재미있는 필수 읽을거리다.
  \item \textbf{[이번 주] OCI runtime-spec과 image-spec (opencontainers.org)}: 제3장에서는 config.json을 한 번 훑었다. 이제 모국어처럼 읽어 보자. image-spec의 매니페스트 절은 10분이면 읽고 §4.3의 내용을 오래 남게 한다.
  \item \textbf{[이번 주] Docker 문서: "Dockerfile best practices"와 BuildKit 문서}: 공급업체가 설명하는 §4.3이다. 캐시 마운트, 다단계 패턴, 빌드 비밀을 다룬다. 자격 증명을 이미지에 굽지 않는 규칙은 제6장의 키 없는 규칙을 빌드 시점에 적용한 것이다.
  \item \textbf{[이번 주] NVIDIA Container Toolkit 문서("How it works")}: §4.5의 주입 메커니즘을 원 출처에서 설명하며 생태계가 표준화 중인 CDI(Container Device Interface) 방향도 포함한다.
  \item \textbf{containerd 문서: 아키텍처와 스냅쇼터}: 수명 주기 엔진이 이미지와 제3장의 overlay 스택을 모델링하는 법을 다룬다. 지연 로딩 스냅쇼터 절은 12주 차 콜드 스타트의 최전선이다.
  \item \textbf{Kubernetes 디바이스 플러그인 문서}: 제5장의 정수 리소스 뒤에 있는 gRPC 계약이다. 짧으면서 §4.5와 §5.5의 고리를 닫는다.
  \item \textbf{[5주 차] Kubernetes 문서: "Images"(imagePullPolicy, 풀 시크릿)}: 이 장의 레지스트리 규율이 연결되는 kubelet 쪽 설정이다.
  \item \textbf{[12주 차] 이미지 스트리밍/사전 pull에 관한 클라우드 공급업체 문서와 "690 ms cold starts" 같은 엔지니어링 글}: §4.4의 풀 폭풍 문제를 산업 규모로 확장한 사례다.
  \item \textbf{[13주 차] Sigstore/cosign과 SLSA 개요}: 파이프라인에 공급망 보안이 들어올 때 "다이제스트를 신뢰하라"가 서명과 출처 증명으로 자라는 모습을 다룬다.
\end{tbbitemize}

\vspace{5.0pt}

\begin{calloutbox}{tbbblue}{tbbbluefill}
{\normalsize\bfseries\color{tbbblue}다음 장: 제5장 Kubernetes}\par\vspace{3pt}
이제 레지스트리에 저장된 잘 만든 이미지를 손에 넣었고 GPU 접근도 주입 훅까지 이해했다. 한 번의 \code{docker run}으로 머신 한 대에서 서비스할 수 있다. 다음 주에는 플릿으로 간다. 원하는 상태를 영원히 조정하므로 다시는 서버를 직접 "시작"하지 않는다는 Kubernetes의 핵심 생각, 이를 구현하는 구조와 CRI를 통해 이 장의 containerd를 구동하는 kubelet, 서빙의 삼위일체인 Pod / Deployment / Service, 플릿 규모에서 제3장의 cgroup 파일을 대필하는 리소스 requests와 limits, 비싼 GPU를 정수·테인트(taint)·빈 패킹(bin packing)으로 배치하는 스케줄러를 배운다. 실습에서는 이번 주에 push한 바로 그 이미지를 minikube 클러스터에 배포하고 노출한 뒤 여섯 방식으로 망가뜨려 각 징후로 진단한다. 인프라 이야기가 완성되기까지 이제 한 장 남았다.
\end{calloutbox}

%% Chapter 5 — Kubernetes
\setcounter{chapter}{4}
\chapter{Kubernetes}
\chaptersubtitle{서빙 플릿을 움직이는 제어 루프 — 원하는 상태, 정직한 자원 예산, 값비싼 하드웨어를 배치하는 스케줄러}

\begin{calloutbox}{tbbblue}{tbbbluefill}
{\normalsize\bfseries\color{tbbblue}학습 목표}\par\vspace{3pt}
\begin{tbbitemize}
  \item Kubernetes의 핵심 아이디어 하나, 즉 데이터베이스에 원하는 상태를 저장하고 제어 루프로 조정한다는 개념을 설명하고, 파드가 죽거나 노드가 사라지거나 배포가 명세를 바꾸는 등 현실이 원하는 상태에서 벗어날 때 클러스터가 어떻게 움직일지 예측한다.
  \item 컨트롤 플레인과 노드의 구성 요소를 열거하고, 한 번의 \code{kubectl apply}가 API 서버 → etcd → 컨트롤러 → 스케줄러 → kubelet → containerd/runc → 제3장의 네임스페이스와 cgroup으로 이어지는 전 과정을 추적한다.
  \item 모델 서버를 Pod + Deployment + Service로 배포하고, 느린 모델 로딩에 맞춰 시작·준비·활성 프로브를 올바르게 연결하며, 준비 상태가 트래픽과 롤링 업데이트를 어떻게 제어하는지 설명한다.
  \item 리소스 requests와 limits가 실제로 무엇인지, 즉 스케줄러의 산술과 cgroup 파일이라는 차이로 해석하고, 노드 패킹을 직접 계산하며, 서빙 파드에 알맞은 QoS 클래스를 의도적으로 선택한다.
  \item 확장 리소스와 정수 단위 할당, 테인트와 어피니티, 빈 패킹과 분산 배치, 그리고 플릿에서 가장 값비싼 하드웨어를 조용히 놀리는 단편화까지 GPU 스케줄링을 논리적으로 분석한다.
  \item 대표적인 여섯 가지 파드 실패 상태인 Pending, ImagePullBackOff, CrashLoopBackOff, OOMKilled, Evicted, 준비됐지만 라우팅되지 않는 상태(ready-but-unrouted)를 각각의 징후로 진단하고, 상황별로 가장 먼저 실행할 명령을 고른다.
\end{tbbitemize}
\end{calloutbox}

\vspace{5.0pt}

\begin{calloutbox}{tbbteal}{tbbtealfill}
{\normalsize\bfseries\color{tbbteal}핵심 용어}\par\vspace{3pt}
컨트롤 플레인 · API 서버 · etcd · 스케줄러 · 컨트롤러 매니저 · 조정 · kubelet · CRI · 모놀리스 · 마이크로서비스 아키텍처(MSA) · 파드 · pause 컨테이너 · 레이블 / 셀렉터 · Deployment · ReplicaSet · 롤링 업데이트 · Service · ClusterIP / NodePort / LoadBalancer · 엔드포인트 · kube-proxy · CoreDNS · 시작 / 준비 / 활성 프로브 · 요청량(requests) / 상한(limits) · 할당 가능 용량 · 오버커밋 · QoS 클래스 (Guaranteed / Burstable / BestEffort) · 노드 압박 퇴거 · 테인트 / 톨러레이션 · nodeSelector / 어피니티 · 확장 리소스 · nvidia.com/gpu · 디바이스 플러그인 · Pending · CrashLoopBackOff · OOMKilled · Evicted
\end{calloutbox}

\vspace{5.0pt}

\begin{calloutbox}{tbblightgray}{tbbboxfill}
{\normalsize\bfseries\color{tbbink}장 구성}\par\vspace{3pt}
\textbf{5.1} 제어 루프로서의 클러스터    \textbf{5.2} 구조: 자기주장이 있는 데이터베이스    \textbf{5.3} Pods, Deployments, Services — 서빙의 세 축    \textbf{5.4} requests, limits, QoS: 플릿 규모로 확장한 제3장    \textbf{5.5} GPU와 스케줄러    \textbf{5.6} 죽어 가는 파드 현장 안내서    \textbf{5.7} 실습: 모델 서버 오케스트레이션 — 이어서 요약, 연습문제, 더 읽을거리.
\end{calloutbox}

\vspace{3.0pt}

\begin{calloutbox}{tbbamber}{tbbamberfill}
{\normalsize\bfseries\color{tbbamberdark}이 장을 읽는 방법}\par\vspace{3pt}
제3장과 마찬가지로 이 장은 \textbf{터미널을 옆에 두고 읽도록} 썼다. GPU 예시만 관리형 GPU 노드 풀에서 실행했을 뿐, 나머지 모든 터미널 기록은 노트북 클러스터에서 그대로 재현할 수 있다. 유일한 사전 준비는 \code{minikube start}(또는 kind)이며, GPU에서도 명령은 같고 하드웨어만 다르다. YAML 파일은 강의 저장소에 있지만, 한 번 직접 입력해 보는 편이 복제만 하는 것보다 훨씬 유익하다.
\end{calloutbox}

\section{제어 루프로서의 클러스터}

제4장은 만족스러운 결과물로 끝났다. 모델 서버를 이미지로 패키징해 레지스트리에 푸시했고, GPU 접근 방식도 디바이스 플러그인까지 이해했다. VM 한 대에서 \code{docker run} 한 번이면 서빙이 시작된다. 데모라면 이야기는 여기서 끝이다. 그러나 제1장의 실제 요구사항이 한꺼번에 밀려온다. 처리량을 높이고 한 번의 장애가 전체 중단으로 번지지 않게 하려면 \textbf{많은 레플리카}가 필요하다. \textbf{요청을 하나도 놓치지 않는 배포}도 필요하다(제3장의 SIGTERM 이야기가 이제 매주 되풀이된다). 사람이 깨어 있지 않은 \textbf{새벽 3시에 노드가 죽어도} 견뎌야 한다. 하드웨어 요구사항이 서로 다른 \textbf{20개 모델}을 하나의 비용 항목 아래에서 운영해야 한다. 그리고 그 비용에서 가장 큰 몫을 차지하는 GPU가 단지 누군가 작업 배치를 잊었다는 이유로 놀아서는 안 된다.

이 모든 일을 스크립트로 처리할 수도 있다. 2014년 무렵에는 모두가 그랬다. SSH를 순회하는 bash 루프, 헬스 엔드포인트에 curl을 보내 실패한 프로세스를 재시작하는 cron 작업, 어느 모델이 어디서 실행되는지 적어 둔 위키 페이지가 있었다. 이런 스크립트를 현실의 요구에 맞춰 키우다 보면 결국 한 형태로 수렴한다. \textbf{무엇이 실행되어야 하는지 기술한 파일, 그 파일과 현재 실행 상태를 비교하는 루프, 둘 사이의 차이를 없애는 동작}이다. Kubernetes는 이 형태를 산업 수준으로 다듬은 시스템이다. Kubernetes가 이를 외부 세계에 공개하기 전, Google의 내부 시스템 Borg는 바로 이 아이디어로 10년간 운영되었다. 이 형태를 이해하면 나머지는 세부사항일 뿐이다.

\subsection*{시스템 전체를 가르쳐 주는 실험}

실행 중인 클러스터에서 시작하자(실습에서는 \code{minikube start}부터 안내한다). 이 클러스터의 Deployment는 모델 서버 레플리카 세 개를 요구한다. 이제 다른 곳이라면 방해 행위로 여길 일을 해 보자. 하나를 죽인 뒤 곧바로 다시 확인한다.

\begin{terminalbox}
\begin{Verbatim}[breaklines=true,breakanywhere=true,fontsize=\small,formatcom=\color{tbbtermfg},codes=\tbbcodeactive]
$ kubectl get pods
NAME                            READY   STATUS    RESTARTS   AGE
model-server-7d4b9c6f5-8xk2p    1/1     Running   0          2d
model-server-7d4b9c6f5-tq9wz    1/1     Running   0          2d
model-server-7d4b9c6f5-zl6vn    1/1     Running   0          2d
$ kubectl delete pod model-server-7d4b9c6f5-zl6vn   # sabotage
pod "model-server-7d4b9c6f5-zl6vn" deleted
$ kubectl get pods
NAME                            READY   STATUS    RESTARTS   AGE
model-server-7d4b9c6f5-8xk2p    1/1     Running   0          2d
model-server-7d4b9c6f5-tq9wz    1/1     Running   0          2d
model-server-7d4b9c6f5-fh5m8    1/1     Running   0          4s
\end{Verbatim}
\end{terminalbox}
\figcaption{터미널 기록 5-1  파드를 삭제하자 4초 뒤 새 파드가 생겼다(새 이름 접미사와 AGE에 주목하자). 누구도 재시작하지 않았다. 컨트롤러가 2와 3이 다르다는 사실을 알아챘을 뿐이다.}

마지막 줄을 자세히 보자. 이전 파드가 되살아난 것이 아니라, 접미사가 다르고 생성된 지 4초밖에 되지 않은 \textbf{새} 파드다. 아무것도 재시작하지 않았다. 더 정확히 말하면 \textbf{애초에 파드를 요청한 적도 없다.} 단지 명세에 \code{replicas: 3}이라고 적힌 Deployment를 선언했다. 원하는 상태와 실제 상태를 끊임없이 비교하는 루프인 컨트롤러가 2 ≠ 3을 관찰하고 하나를 만들었다. 파드 삭제는 누군가 처리한 이벤트가 아니라 \textbf{편차(drift)}였고, 편차는 교정된다(그림 5-1). 그래서 Kubernetes의 문법은 처음에 낯설다. 동사는 쓰지 않고 \textbf{명사}만 쓴다(객체: "레플리카가 세 개인 Deployment"). 생성, 재시작, 재스케줄링 같은 동사는 명사를 대신해 컨트롤러가 수행하면서 자연스럽게 나타난다.

\begin{tbbfigure}
  \adjustbox{max width=0.94\linewidth,max totalheight=0.46\textheight}{\includegraphics{figures/fig5-1_loop.png}}\\[3pt]
  \figcaption{그림 5-1  조정 루프. 원하는 상태는 API 서버 뒤의 etcd에 저장된다. 컨트롤러는 양쪽을 지켜보며 차이가 사라질 때까지 영원히 동작한다.}
\end{tbbfigure}

\subsection*{새벽 3시에는 선언형이 명령형을 이긴다}

이 차이에는 이름이 있다. \textbf{명령형(imperative)} 시스템은 "서버 세 대를 시작하라" 같은 지시를 받는다. 지시는 한 번 실행되고 나면 시간이 갈수록 효력이 사라진다. 나중에 서버가 죽어도 지시는 이미 끝났으므로, 누군가 이를 발견해 새 지시를 내려야 한다. \textbf{선언형(declarative)} 시스템은 "서버 세 대가 존재한다"라는 \textit{불변 조건}을 유지한다. 불변 조건은 효력을 잃지 않는다. 계속해서 다시 확인하고 강제하므로 "X가 실패하면 어떻게 되는가?"라는 질문의 답은 언제나 같다. \textit{루프가 알아채고 다시 수렴한다.} 가정용 온도조절기도 같은 원리다. 21°C를 선언하면 누군가 창문을 열 때마다 선언을 반복하지 않아도 된다. 서빙 플릿에 꼭 필요한 성질이 바로 이것이다. 새벽 3시의 노드 장애, 실수로 누른 삭제, 플릿을 절반으로 줄여 버린 배포는 루프의 눈에는 모두 똑같은 편차다.

선언형 습관은 이 강의 뒤쪽에서도 복리처럼 이익을 낸다. 12주 차에 오토스케일링을 추가해도 오토스케일러는 전혀 새로운 종류의 장치가 아니다. 부하가 커지면 \code{replicas: 3}을 \code{replicas: 7}로 바꾸고 나머지는 이 절의 루프에 맡기는 \textbf{또 하나의 원하는 상태 작성자}일 뿐이다. 13주 차에 CI/CD로 배포를 자동화할 때도 파이프라인의 마지막 동작은 객체 한 줄을 바꾸는 것이다. 이 강의의 운영 파트에서 일어나는 모든 일은 사람이나 시스템이 명사를 편집하는 과정이며, 유일한 엔진은 그림 5-1의 루프다.

\subsection*{이 장이 실제로 다루는 것}

서빙 엔지니어에게 Kubernetes는 네 가지 높이에서 중요하며, 이 장은 각 층을 차례로 살펴본다. 가장 위의 \textbf{제어 루프}(이 절)는 앞으로 마주칠 모든 실패 상황에서 시스템이 어떻게 움직이는지 설명한다. 한 층 아래의 \textbf{객체}, 즉 Pod, Deployment, Service(§5.3)는 모델 서버를 기술하고 상태를 확인하며 안정적인 주소를 부여하는 어휘다. 그 아래의 \textbf{리소스와 QoS}(§5.4)는 제3장의 cgroup 벽을 플릿 규모로 대신 작성하는 곳이다. 이제 자다가도 설명할 수 있는 OOM 동작도 여기에 포함된다. 가장 아래에서는 \textbf{스케줄러}(§5.5)가 어떤 값비싼 GPU에서 어떤 모델을 실행할지 결정한다. 매달 청구서로 결과가 드러나는 배치 문제다. 마지막으로 §5.6에서는 이 전체 스택을 실제 업무에서 쓰이는 형태, 즉 체계적인 진단법으로 바꾼다.

\begin{calloutbox}{tbbpurple}{tbbpurplefill}
{\normalsize\bfseries\color{tbbpurple}생각해 보기}\par\vspace{3pt}
\begin{tbbitemize}
  \item 컨트롤러는 2 ≠ 3을 보고 파드를 하나 만들었다. 이번에는 ReplicaSet이 선택하는 것과 같은 레이블(label)을 단 네 번째 파드를 직접 만들었다고 하자. 컨트롤러는 무엇을 해야 하며, 이는 구성원 자격을 결정하는 방식에 관해 무엇을 뜻하는가?
  \item 컨트롤러 자체가 중단되었을 때 조정 루프(reconciliation loop)가 다시 복구되는 능력이, 명령형 스크립트를 실행하던 호스트가 중단되었을 때보다 뛰어난 이유는 무엇인가? (힌트: 올바르게 재개하려면 각각 어떤 상태가 필요한가?)
  \item 12주 차에는 Horizontal Pod Autoscaler가 부하에 반응해 한 객체의 정수 하나를 바꾼다. 그 뒤 자동으로 일어나야 한다고 이 절에서 설명한 모든 일을 나열해 보자. 오토스케일러 자체가 알아야 하는 것은 놀라울 만큼 적다.
\end{tbbitemize}
\end{calloutbox}

\section{구조: 자기주장이 있는 데이터베이스}

신비감을 걷어 내면 Kubernetes의 아키텍처는 한 문장으로 정리된다. \textbf{앞에는 API(API 서버)가 있고 뒤에는 수많은 제어 루프(컨트롤러, 스케줄러, 노드마다 하나씩 있는 kubelet)가 붙은 데이터베이스(etcd)}다. 나머지는 모두 부연 설명이다. 그림 5-2는 구성 요소를 한눈에 펼쳐 보인다. 이 절에서는 각 요소를 하나씩 설명한 뒤, 한 번의 \code{kubectl apply}가 모든 요소를 통과하는 과정을 추적한다. 실제로 디버깅하는 대상은 개별 부품이 아니라 이 연결 고리이기 때문이다.

\begin{tbbfigure}
  \adjustbox{max width=0.94\linewidth,max totalheight=0.46\textheight}{\includegraphics{figures/fig5-2_architecture.png}}\\[3pt]
  \figcaption{그림 5-2  클러스터 구조. 컨트롤 플레인은 원하는 상태를 저장하고 조정한다. 각 노드의 kubelet은 제4장의 런타임 스택을 거쳐 파드를 현실로 만들며, 그 끝에는 제3장의 네임스페이스와 cgroup이 있다.}
\end{tbbfigure}

\subsection*{컨트롤 플레인: 구성 요소별 한 문단}

\textbf{API 서버}는 유일한 출입구다. kubectl을 사용하는 사람, 모든 컨트롤러, 모든 kubelet은 오직 API 서버하고만 통신한다. 그 외에는 아무것도 etcd에 접근할 수 없다. API 서버는 인증하고 권한을 확인하며 유효성을 검사한 뒤, 다소 싱거울 만큼 단순한 일을 한다. 객체에 CRUD를 수행한다. 내부에 숨겨진 오케스트레이션 두뇌는 없다. 출력 상세도를 높이면 \code{kubectl}조차 얇은 REST 클라이언트일 뿐이라는 사실이 드러난다.

\begin{terminalbox}
\begin{Verbatim}[breaklines=true,breakanywhere=true,fontsize=\small,formatcom=\color{tbbtermfg},codes=\tbbcodeactive]
$ kubectl get pods -v=8 2>&1 | grep -E "GET http"
GET https://192.168.49.2:8443/api/v1/namespaces/default/pods?limit=500
# kubectl get pods = one HTTP GET. kubectl is a REST client with manners.
$ kubectl get pods -n kube-system
NAME                               READY   STATUS    RESTARTS      AGE
etcd-minikube                      1/1     Running   0             9d
kube-apiserver-minikube            1/1     Running   0             9d
kube-controller-manager-minikube   1/1     Running   0             9d
kube-scheduler-minikube            1/1     Running   0             9d
coredns-5dd5756b68-mfqpx           1/1     Running   0             9d
kube-proxy-b47rl                   1/1     Running   0             9d
\end{Verbatim}
\end{terminalbox}
\figcaption{터미널 기록 5-2  두 가지 오해를 걷어 낸다. kubectl은 HTTP이며, 컨트롤 플레인 자체도 kube-system 네임스페이스의 파드로 실행된다. 시스템이 자기 자신을 관리하는 셈이다.}

\textbf{etcd}는 클러스터의 기억 장치다. 원하는 상태와 마지막으로 관찰한 상태를 포함한 모든 객체를 보관하는 작고 일관된(Raft 합의 기반) 키-값 저장소다. 이것이 \textit{유일한} 상태 저장 구성 요소다. 다른 구성 요소는 언제든 중단해도 재시작한 뒤 API 서버를 다시 지켜보는 것만으로 작업에 필요한 지식을 복구한다. 잠시 감탄할 만한 설계 결정이다. 상태는 복제된 저장소 하나에 집중하고, 나머지는 모두 상태가 없으며 교체 가능하게 만든다. 나중에 직접 설계할 서빙 시스템에도 똑같이 적용하라고 권할 형태다.

\textbf{스케줄러}는 역할이 매우 좁은 중매자다. 아직 노드가 배정되지 않은 파드를 지켜보다가, 각 파드에 가장 적합한 노드를 고르고(먼저 필터링하고 점수를 매기는 과정은 GPU 때문에 흥미로워지는 §5.5에서 해부한다), 그 선택을 파드 객체에 다시 기록한다. 하는 일은 이것뿐이다. 아무것도 시작하지 않고 어떤 노드도 직접 건드리지 않는다. \textbf{컨트롤러 매니저}는 §5.1의 조정 루프들을 실행한다. Deployment 컨트롤러, ReplicaSet 컨트롤러, 죽은 노드를 알아채는 노드 수명 주기 컨트롤러를 비롯한 수십 개의 루프가 있으며, 각 루프는 한 객체 유형을 대상으로 "관찰하고, 차이를 찾고, 동작하는" 작은 과정이다.

모든 노드에서 \textbf{kubelet}은 손발 역할을 한다. API 서버를 지켜보다가 \textit{자기} 노드에 바인딩된 파드를 현실로 만든다. 제4장의 containerd → runc 스택 앞에 놓인 gRPC 인터페이스인 CRI를 통해 컨테이너 런타임에 지시하고, 볼륨을 마운트하며, §5.3의 프로브를 실행하고, 상태를 다시 보고한다. 메모리 압박이 생기면 §5.4의 퇴거도 수행한다. 이 설명에서 kubelet이 누구인지 눈여겨보자. 바로 \textbf{제3장이 계속 예고했던 대필자}다. 파드가 \code{limits.memory: 1.75Gi}를 선언하면, 런타임을 거쳐 \code{memory.max} 파일에 \code{1879048192}가 나타나게 만드는 주체가 kubelet이다. 그 옆에는 §5.3의 Service NAT를 프로그래밍하는 \textbf{kube-proxy}가 실행된다. GPU 노드에서는 제4장의 \textbf{디바이스 플러그인}도 함께 실행되어 GPU가 몇 개 있고 어느 것이 비어 있는지 kubelet에 알려 준다.

\subsection*{한 번의 apply를 처음부터 끝까지}

이제 연결 고리를 따라가 보자. 레플리카가 세 개인 새 모델 서버를 위해 \code{kubectl apply -f deployment.yaml}을 입력한다. 전체를 지휘하는 주체 없이 다음 일이 순서대로 일어난다.

\begin{tbbitemize}
  \item \textbf{1 — 저장.} API 서버가 Deployment 객체의 유효성을 검사하고 etcd에 기록한다. 아직 아무것도 실행되지 않았다. 데이터베이스의 한 행만 바뀌었다.
  \item \textbf{2 — 확장.} Deployment 컨트롤러는 일치하는 ReplicaSet이 없는 Deployment를 발견하고 하나를 만든다(이번에도 객체다). ReplicaSet 컨트롤러는 원하는 수가 3인데 실제 수가 0임을 알아채고, \code{nodeName}이 비어 있는 Pod 객체 세 개를 만든다.
  \item \textbf{3 — 배치.} 스케줄러는 노드에 바인딩되지 않은 파드 세 개를 발견하고, 파드마다 필터링과 점수 계산을 거쳐 노드 이름을 기록한다. 여전히 어느 곳에서도 아무것도 실행되지 않았다.
  \item \textbf{4 — 실체화.} 선택된 각 노드의 kubelet은 자신에게 바인딩된 파드를 발견한다. 이미지가 없으면 가져오고(제3장에서 계산한 25–60 s의 콜드 풀 시간이 노드마다 적용된다), containerd → runc를 구동해 pause 컨테이너와 서버 컨테이너를 만든다. 분리된 네임스페이스, 작성된 cgroup 파일, 마운트된 overlayfs까지 제3장의 다섯 동작을 시스템이 수행한다.
  \item \textbf{5 — 보고.} kubelet이 상태를 다시 게시하고 프로브가 시작된다. 준비 상태 검사를 통과하면 파드의 IP가 Service 엔드포인트(§5.3)에 들어가고 트래픽이 흐른다.
\end{tbbitemize}

이 연결 고리에서 두 가지는 특히 강조할 가치가 있다. 첫째, \textit{Kubernetes라는 이름의 무언가가 컨테이너를 실행한 것이 아니다.} 실행은 runc가 했고, Kubernetes는 위치를 결정하고 상태를 기록했다. 제3장이 컨테이너의 신비를 걷어 냈다면, 이 절은 오케스트레이터의 신비를 걷어 낸다. 같은 커널 메커니즘 주위에서 장부를 관리할 뿐이다. 둘째, \textit{연결 고리의 모든 화살표는 API 서버에 대한 watch다.} 구성 요소끼리는 서로 호출하지 않는다. 이처럼 철저히 분리했기 때문에 어느 한 구성 요소가 멈춰도 시스템이 견딘다. 작업은 etcd에서 기다리고, 재개한 구성 요소가 가져가면 된다. (또한 컨트롤 플레인에 문제가 생겨도 데이터 플레인이 계속 서빙한다는 뜻이다. 이미 실행 중인 파드, 이미 프로그래밍된 Service 경로, 이미 작성된 cgroup 파일은 계속 동작하는 데 허가가 필요 없다. 컨트롤 플레인이 돌아올 때까지 아무것도 \textit{변경}하지 못할 뿐이다.)

\begin{calloutbox}{tbbpurple}{tbbpurplefill}
{\normalsize\bfseries\color{tbbpurple}생각해 보기}\par\vspace{3pt}
\begin{tbbitemize}
  \item etcd가 10분 동안 쿼럼을 잃었다. 다음 가운데 계속 동작하는 것은 무엇이며, 그 이유는 무엇인가? (a) 기존 파드의 트래픽 서빙, (b) 파드로 향하는 Service 라우팅, (c) 롤링 업데이트, (d) 파드 중단 후 자가 복구
  \item 구성 요소끼리 직접 통신하게 하지 않고(예: 스케줄러 → kubelet) API 서버에 대한 watch를 거치게 한 이유는 무엇인가? 이 간접 계층이 흡수하는 실패 시나리오 두 가지를 들어 보자.
  \item 컨트롤러가 중단되었다가 5분 뒤 재시작했다. 관찰한 상태를 기준으로 동작하는 "레벨 트리거" 조정이 이벤트를 기준으로 동작하는 "에지 트리거" 방식보다 이 상황을 자연스럽게 견디는 이유를 설명하자.
\end{tbbitemize}
\end{calloutbox}

\section{Pods, Deployments, Services — 서빙의 세 축}

이 강의에서 실행할 거의 모든 서빙 워크로드는 세 객체가 떠받친다. \textbf{Pod}는 서버와 보조 컨테이너를 묶는 실행 단위다. \textbf{Deployment}는 개수, 버전, 업그레이드 방식을 관리하는 플릿 관리자다. \textbf{Service}는 언젠가 사라질 레플리카 앞에 하나의 이름을 제공하는 안정적인 주소다. 이 세 가지를 깊이 이해하면 앞으로 만날 Kubernetes YAML 대부분이 그 변형으로 보인다.

\subsection*{먼저 워크로드의 형태부터: 마이크로서비스 개요}

세 객체를 해부하기 전에 이들이 암묵적으로 가정하는 아키텍처 스타일에 이름을 붙여 보자. 전통적인 대안은 \textbf{모놀리스(monolith)}다. UI, 비즈니스 로직, 데이터 접근 같은 모든 기능을 하나의 배포 단위에 담고, 하나의 코드베이스에서 빌드해 하나의 결과물로 출시하며, 전체 복사본을 더 실행해 확장한다. \textbf{마이크로서비스(microservice) 아키텍처(MSA)}는 애플리케이션을 작고 독립적으로 배포 가능한 서비스로 나눈다. 각 서비스는 하나의 기능과 이상적으로는 자체 데이터를 소유하고, HTTP나 gRPC 같은 네트워크 API를 통해서만 다른 서비스에 기능을 제공한다. 보통 전담 팀 아래에서 독자적인 출시 주기로 발전한다.

업계가 유행을 좇아 이 구조로 옮겨 간 것은 아니다. 이점은 운영에서 드러난다. \textbf{독립 배포}: 결제 서비스를 다시 출시하지 않고도 추천 서비스를 매일 배포할 수 있다. \textbf{독립 확장}: 부하가 몰리는 곳에만, ML 시스템이라면 GPU가 필요한 곳에만 레플리카를 추가한다. \textbf{장애 격리}: 추천 서비스가 중단되면 페이지 품질은 낮아져도 결제까지 멈추지는 않는다. \textbf{조직과의 정합성}: 작은 팀이 작은 서비스를 처음부터 끝까지 소유한다. 물론 대가도 분명하며, 이것이 솔직한 나머지 절반이다. 프로세스 내부 함수 호출이 모두 \textbf{네트워크 호출}로 바뀌어 지연 시간, 타임아웃, 부분 실패, 분산 디버깅을 동반한다. 하나의 배포 대상은 \textbf{N개의 배포 대상}이 되고, 각각 빌드, 버전 관리, 서비스 발견(service discovery), 상태 검사, 확장, 모니터링이 필요하다.

\begin{calloutbox}{tbbteal}{tbbtealfill}
{\normalsize\bfseries\color{tbbteal}모놀리스 대 마이크로서비스 — 한눈에 보는 절충 관계}\par\vspace{3pt}
\begin{tbbitemize}
  \item \textbf{배포} — 모놀리스: 결과물 하나를 한 번에 전체 출시. MSA: 서비스마다 독립적인 주기로 출시.
  \item \textbf{확장} — 모놀리스: 전체를 복제. MSA: 부하가 큰 경로만 확장(서빙에서는 GPU를 많이 쓰는 부분만 확장).
  \item \textbf{실패} — 모놀리스: 운명을 공유. MSA: 장애 영향 범위를 격리. 단, 타임아웃과 대체 경로를 실제로 설계해야 하며 저절로 얻어지지는 않는다.
  \item \textbf{호출} — 모놀리스: 프로세스 내부에서 나노초 단위. MSA: 네트워크를 거치므로 지연 시간과 부분 실패가 생기고, 무슨 일이 있었는지 재구성하려면 추적이 필요하다.
  \item \textbf{조직} — 모놀리스: 하나의 코드베이스, 회의로 조율. MSA: 서비스별 팀(Conway의 법칙을 의도적으로 활용).
\end{tbbitemize}

한 문장으로 요약하면 \textbf{코드 복잡성을 운영 복잡성과 맞바꾸며, Kubernetes는 업계가 운영 측면을 표준화한 결과물이다.}
\end{calloutbox}

\vspace{3.0pt}

마지막 문장이 이 장의 존재 이유를 설명한다. MSA가 만드는 운영 부담은 도입 조직마다 거의 같다. 그래서 산업화할 수 있었고, 그 산업화의 결과가 Kubernetes다. 아래의 세 축을 \textbf{마이크로서비스 입문 키트}로 보자. 배포 단위인 Pod, 서비스별 플릿 관리자인 Deployment, 그리고 서비스 발견과 로드 밸런싱(load balancing)을 제공하는 Service(서비스마다 DNS 이름 하나)다. 사실 이 스타일을 직접 선택한 것도 아니다. \textbf{AI 서빙 시스템은 태어날 때부터 작은 MSA다.} 앞에는 CPU 비용이 낮은 API 게이트웨이가 있고(작은 레플리카가 많으며 매주 출시), 뒤에는 모델마다 GPU 비용이 큰 모델 서버가 있으며(레플리카가 적고 모델 버전이 바뀔 때마다 재배포 — 13주 차), 각자 병목에 맞춰 확장하는 전처리기, 후처리기, 익스포터(exporter)가 따라붙는다. 서로 다른 하드웨어, 확장 방식, 출시 주기는 마이크로서비스가 존재하는 바로 그 조건이다. 강의 나머지 전체에 통하는 경험 법칙은 \textbf{마이크로서비스 하나 = Deployment 하나 + Service 하나}다. 9주 차의 서빙 프레임워크는 이를 "model-as-a-service"로 가정하고, 12주 차의 KServe는 이를 형식화하며, 14주 차의 관측 가능성 스택은 바로 이 구조 때문에 존재한다.

\subsection*{Pod: 스케줄링과 공유의 단위}

파드는 하나의 단위로 스케줄링되며 \textbf{네트워크 정체성}과 \textbf{볼륨} 두 가지를 공유하는 하나 이상의 컨테이너 묶음이다. Kubernetes는 제3장에서 이미 본 기법으로 공유를 구현한다. 먼저 작은 \textbf{pause 컨테이너}를 시작해 \textit{파드의 NET 네임스페이스}(그리고 IPC)를 소유하게 하고, 나머지 컨테이너를 여기에 합류시킨다(그림 5-3). 파드 하나에는 IP가 하나 있다. 내부 컨테이너는 \code{localhost}로 서로 통신하며, MNT와 PID 네임스페이스는 §3.2에서 설명한 그대로 필요에 따라 각자 분리된다. 애플리케이션 컨테이너가 재시작되어도 파드가 살아남는 이유도 pause 컨테이너에 있다. 아무 일도 하지 않아서 중단될 일이 없는 프로세스가 네트워크 정체성을 유지한다.

컨테이너를 그대로 쓰지 않고 파드로 묶는 이유는 무엇일까? 실제 서버에는 메트릭 익스포터, 로그 수집기, 나중에는 서비스 메시(service mesh) 프록시 같은 수행원이 따라붙기 때문이다(14주 차에 이 가운데 두 가지를 만난다). 파드는 "함께 태어나고 죽으며 스케줄링되어야 하는 것들"의 경계다. 서빙의 경험 법칙은 \textbf{파드마다 주 컨테이너는 하나만 두고, 횡단 관심사를 처리하는 기반 기능에만 사이드카를 사용한다}는 것이다. 모델 서버 두 개를 한 파드에 넣지 않는다. 그러면 서로 독립적으로 확장하거나 업그레이드하거나 OOM을 격리할 수 없다. 제3장에서 보았듯 memory.max는 어차피 컨테이너별이므로 이 격리는 실제로 작동한다.

\begin{tbbfigure}
  \adjustbox{max width=0.94\linewidth,max totalheight=0.46\textheight}{\includegraphics{figures/fig5-3_pod.png}}\\[3pt]
  \figcaption{그림 5-3  파드 구조(왼쪽): pause가 공유 네임스페이스를 소유하고 컨테이너가 합류한다. 세 가지 프로브(오른쪽): 시작 프로브는 모델 로딩을 보호하고, 준비 상태는 트래픽을 제어하며, 활성 상태는 실제로 죽은 프로세스를 재시작한다.}
\end{tbbfigure}

\subsection*{실제로 작성할 매니페스트}

강의에서 사용할 모델 서버의 전체 Deployment 매니페스트다. §5.4와 §5.6이 각각 여기의 한 구절을 바탕으로 절 하나를 전개하므로, 한 줄씩 읽어 볼 가치가 있다.

\begin{terminalbox}
\begin{Verbatim}[breaklines=true,breakanywhere=true,fontsize=\small,formatcom=\color{tbbtermfg},codes=\tbbcodeactive]
apiVersion: apps/v1
kind: Deployment
metadata:
  name: model-server
spec:
  replicas: 3
  selector:
    matchLabels: { app: model }
  template:
    metadata:
      labels: { app: model }        # must match selector - membership tag
    spec:
      terminationGracePeriodSeconds: 30   # Ch. 3's SIGTERM tape, the knob
      containers:
      - name: server
        image: registry.example.com/team/model-server:v7
        ports: [ { containerPort: 8000 } ]
        resources:                  # sec. 5.4 lives here
          requests: { cpu: "1", memory: 1.2Gi }
          limits:   { cpu: "2", memory: 1.75Gi }
        startupProbe:               # covers the model load
          httpGet: { path: /healthz, port: 8000 }
          periodSeconds: 5
          failureThreshold: 36      # 36 x 5s = 3 min of grace
        readinessProbe:             # gates traffic
          httpGet: { path: /ready, port: 8000 }
          periodSeconds: 5
        livenessProbe:              # restarts the truly dead
          httpGet: { path: /healthz, port: 8000 }
          periodSeconds: 10
          failureThreshold: 3
\end{Verbatim}
\end{terminalbox}
\figcaption{터미널 기록 5-3  강의에서 사용할 Deployment 매니페스트. 세 구절이 핵심 역할을 한다. 레이블(구성원 자격), 리소스(§5.4), 프로브(아래와 §5.6)다.}

\begin{terminalbox}
\begin{Verbatim}[breaklines=true,breakanywhere=true,fontsize=\small,formatcom=\color{tbbtermfg},codes=\tbbcodeactive]
$ kubectl apply -f model-server.yaml
deployment.apps/model-server created
$ kubectl get pods -o wide
NAME                            READY  STATUS    IP            NODE
model-server-7d4b9c6f5-8xk2p    1/1    Running   10.244.1.17   node-a
model-server-7d4b9c6f5-tq9wz    1/1    Running   10.244.2.31   node-b
model-server-7d4b9c6f5-fh5m8    0/1    Running   10.244.1.44   node-a
# 0/1 = running but not READY: the model is still loading.
# readiness is failing on purpose - and no traffic arrives yet.
\end{Verbatim}
\end{terminalbox}
\figcaption{터미널 기록 5-4  적용되어 여러 노드에 배치된 모습. 세 번째 파드에 주목하자. Running은 프로세스가 실행 중이라는 뜻이고, READY는 서빙할 준비가 되었다는 뜻이다. 둘 사이의 간격이 모델 로딩 시간이다.}

\subsection*{프로브: 세 질문과 서로 다른 세 결과}

kubelet은 모든 컨테이너에 세 가지 질문을 계속 던진다. 서빙 시스템의 생사는 이 답을 구별하는 데 달려 있다. \textbf{startupProbe}는 "부팅을 마쳤는가?"를 묻는다. 모델 때문에 시작이 느려지는 상황을 위해 존재하며, 실행되는 동안 다른 두 프로브(probe)는 중지된다. 예산(\code{failureThreshold × periodSeconds}; 터미널 기록 5-3에서는 3분)은 최악의 가중치 로딩 시간을 포괄해야 한다. \textbf{readinessProbe}는 "트래픽을 보내도 되는가?"를 묻고, 오직 Service 엔드포인트의 구성원 자격만 제어한다. 준비 상태 검사 실패는 아무것도 재시작하지 않는다. 조용히 파드를 로테이션에서 제외하며, 이는 로딩 중이거나 과부하 상태이거나 처리 중인 요청을 비우는 동안 원하는 동작이다. \textbf{livenessProbe}는 "살아 있기는 한가?"를 묻는 둔중한 도구다. 실패하면 컨테이너를 종료하고 재시작한다. 의존성이 아니라 프로세스 자체의 상태(이벤트 루프의 응답 여부)를 검사해야 하며, 충분히 긴 주기와 넉넉한 임계값을 두어 \textit{느긋하게} 동작해야 한다.

이제 대표적인 서빙 사고를 한 문장으로 설명할 수 있다. 기본 대기 시간으로 설정된 활성 상태 프로브가 모델 로딩에 90초가 걸리는 서버를 검사하면, 로딩 중인 컨테이너를 영원히 죽인다. \textbf{아무 문제가 없는 이미지에서 CrashLoopBackOff가 발생한다.} 실제 현장에서 틀림없이 만날 문제이므로 §5.6에서 터미널 기록과 함께 전체 사고를 재현한다. 예방책은 위의 startupProbe 설정이고, 진단 습관은 터미널 기록 5-4의 구분이다. Running은 Ready가 아니며, 어느 쪽도 죽었다는 뜻은 아니다.

\subsection*{Deployment: 버전이 있는 플릿과 무중단 롤아웃}

Deployment는 파드 템플릿 버전마다 하나씩 존재하는 ReplicaSet을 관리한다(템플릿 해시는 터미널 기록 5-1의 파드 이름에 그대로 나타난다). 새 이미지 태그나 새 메모리 상한처럼 템플릿의 무엇이든 바꾸면 Deployment는 \textit{새} ReplicaSet을 만들고 정해진 속도로 기존 파드와 새 파드를 교대한다. 속도는 전환 중 추가로 존재할 수 있는 수를 뜻하는 \code{maxSurge}와 비어 있어도 되는 수를 뜻하는 \code{maxUnavailable}이 결정한다. 서빙에 적합한 설정은 \code{maxSurge: 1, maxUnavailable: 0}이다. 용량을 절대 떨어뜨리지 않고, 하나를 추가해 \textbf{Ready}가 될 때까지 기다린 뒤 하나를 제거한다. 여기서 "Ready가 될 때까지 기다린다"는 문구가 핵심이다. \textbf{롤아웃의 속도는 준비 상태가 결정한다.} 새 버전이 끝내 준비 상태가 되지 않으면 추가 파드 하나에서 롤아웃이 멈춘다. 기존 버전은 계속 서빙하고 장애 영향 범위는 파드 하나로 제한된다. 시스템이 조용히 사용자의 실수로부터 사용자를 구해 주는 셈이다.

\begin{terminalbox}
\begin{Verbatim}[breaklines=true,breakanywhere=true,fontsize=\small,formatcom=\color{tbbtermfg},codes=\tbbcodeactive]
$ kubectl set image deployment/model-server server=.../model-server:v8
deployment.apps/model-server image updated
$ kubectl rollout status deployment/model-server
Waiting for rollout: 1 out of 3 new replicas are available...
$ kubectl get pods
NAME                            READY   STATUS        RESTARTS   AGE
model-server-7d4b9c6f5-8xk2p    1/1     Running       0          2d
model-server-7d4b9c6f5-tq9wz    1/1     Terminating   0          2d
model-server-6f8d5b9c44-2mwpx   1/1     Running       0          38s
model-server-6f8d5b9c44-9qk7d   0/1     Running       0          6s
# surge pod loading; an old pod drains only after a new one is Ready.
# regret it? one command points back at the old ReplicaSet:
$ kubectl rollout undo deployment/model-server
\end{Verbatim}
\end{terminalbox}
\figcaption{터미널 기록 5-5  진행 중인 롤링 업데이트. v8 파드는 Ready가 되는 대로 합류하고, 그제야 v7 파드가 종료된다. 롤백은 규모 0으로 유지해 둔 이전 ReplicaSet을 다시 가리키는 일이다.}

제3장의 관점으로 \code{Terminating} 파드를 보자. kubelet은 \textbf{PID 1에 SIGTERM}을 보내고 \code{terminationGracePeriodSeconds}만큼 기다린 뒤(매니페스트에서는 30으로 설정했다) SIGKILL을 보낸다. exec 형식 CMD, 새 요청 수락을 멈추고 처리 중인 요청을 비우는 핸들러, 필요하다면 tini까지 §3.2에서 배운 모든 것이 \code{Terminating} 시간을 클라이언트에게 보이지 않게 한다. Kubernetes가 우아한 종료(graceful shutdown)를 해결한 것은 아니다. 단지 종료를 \textit{스케줄링}한다. 이 강의의 나머지 기간 동안 모든 롤아웃, 노드 비우기, 축소 때마다 실제 작업은 여전히 핸들러가 수행한다.

\subsection*{Service: 사라질 파드 앞에 놓는 안정적인 이름 하나}

파드는 가축과 같다. 태어날 때 IP를 받고 죽을 때 잃는다. 터미널 기록 5-5에서 죽음이 얼마나 일상적인지 보았다. 클라이언트에게 이를 일일이 따라가라고 할 수는 없다. 해결책은 \textbf{Service}다. 현재 \textbf{레이블 셀렉터(selector)}와 일치하면서 준비 상태 검사를 통과한 파드 앞에 안정적인 가상 IP와 DNS 이름(\code{model-svc.ml.svc.cluster.local}, CoreDNS가 제공)을 둔다. 일치하는 파드의 IP가 \textbf{엔드포인트} 목록을 구성하고 컨트롤러가 이를 실시간으로 갱신한다. \textbf{kube-proxy}는 가상 IP로 들어온 연결을 정상 백엔드 하나로 DNAT 변환하도록 노드의 패킷 경로(iptables 또는 IPVS)를 프로그래밍한다. 제3장의 \code{-p 8080:8000} 패킷 경로에서 정확히 같은 메커니즘을 만났다. 공개 포트는 NAT 규칙 하나였고, Service는 같은 기법을 클러스터 전체에 건강 상태를 반영해 적용한다(그림 5-4). MSA 관점에서는 \textbf{플랫폼이 해결한 서비스 디스커버리}다. 소비자는 이름만 알 뿐 토폴로지는 알 필요가 없다.

\begin{tbbfigure}
  \adjustbox{max width=0.94\linewidth,max totalheight=0.46\textheight}{\includegraphics{figures/fig5-4_deploy_svc.png}}\\[3pt]
  \figcaption{그림 5-4  Deployment → ReplicaSet → 파드 앞에 Service가 놓인 구조. 엔드포인트 목록에는 READY 파드만 들어간다. 그래서 로딩 중인 파드와 처리 중인 요청을 비우는 파드가 트래픽에서 자동으로 빠진다.}
\end{tbbfigure}

\begin{terminalbox}
\begin{Verbatim}[breaklines=true,breakanywhere=true,fontsize=\small,formatcom=\color{tbbtermfg},codes=\tbbcodeactive]
$ kubectl expose deployment model-server --name=model-svc \
    --port 80 --target-port 8000
service/model-svc exposed
$ kubectl get svc model-svc
NAME        TYPE        CLUSTER-IP     EXTERNAL-IP   PORT(S)   AGE
model-svc   ClusterIP   10.96.44.7     <none>        80/TCP    5s
$ kubectl get endpoints model-svc
NAME        ENDPOINTS                             AGE
model-svc   10.244.1.17:8000,10.244.2.31:8000     5s
# two backends - the third pod is still loading (not Ready).
# from any pod in the cluster:
$ curl -s http://model-svc.default/predict -d '{"x": [1,2]}'
{"y": 0.87, "model": "v7"}
\end{Verbatim}
\end{terminalbox}
\figcaption{터미널 기록 5-6  플릿 앞에 놓인 Service. 가상 IP, DNS 이름, 준비 상태 프로브에 응답하는 파드만 담는 엔드포인트 목록을 제공한다.}

세 가지 Service 유형은 노출 범위를 단계적으로 넓힌다. \textbf{ClusterIP}(위 예시의 기본값)는 클러스터 내부에서만 접근할 수 있어 다른 서비스 뒤에 놓인 모델 서버에 적합하다. \textbf{NodePort}는 모든 노드에서 높은 번호의 포트(30000–32767)를 연다. minikube 실습에서 외부로 나가는 방법이다. \textbf{LoadBalancer}는 해당 노드 포트에 연결된 실제 외부 로드 밸런서를 클라우드에 요청한다. 관리형 클러스터에서 쓰는 방식이며, 6주 차에는 앞단에 API 게이트웨이, TLS, 인증을 추가한다. 지금 기억할 점이 하나 있다. Service는 요청이 아니라 \textit{연결}을 로드 밸런싱한다. 스트리밍이나 gRPC 멀티플렉싱처럼 오래 유지되는 연결은 이를 무력화할 수 있으며, 이 문제는 11주 차 LLM 스트리밍에서 다시 등장한다.

세 축이 아직 기억에 생생할 때 습관을 하나 더 들이자. 바로 \textbf{네임스페이스}다. 이름, 할당량, 접근 권한의 범위를 나눈다. 실습에서는 모든 것을 \code{default}가 아니라 \code{ml}이라는 네임스페이스 안에 만든다. 실제로 접할 모든 클러스터가 이런 식으로 나뉘며, 13주 차의 CI/CD 파이프라인도 환경별 네임스페이스에 배포하기 때문이다.

\begin{calloutbox}{tbbpurple}{tbbpurplefill}
{\normalsize\bfseries\color{tbbpurple}생각해 보기}\par\vspace{3pt}
\begin{tbbitemize}
  \item model-server-v7과 실험용 v8이라는 두 Deployment에 실수로 모두 app=model 레이블이 붙었고, 하나의 Service가 app=model을 선택한다. 클라이언트가 정확히 무엇을 경험하는지, 어떤 구성 요소도 오류를 알리지 않는 이유는 무엇인지, 어디에서 이 상황을 가장 먼저 확인할 수 있는지 설명하자. (이 구성을 의도적으로 만들면 카나리라 부른다. 13주 차에서 다룬다.)
  \item 롤링 업데이트는 컨테이너가 시작되기만을 기다리는 것이 아니라 준비 상태를 기다린다. 전체 트래픽을 받는 상황에서 로딩에 90초가 걸리는 모델에 maxSurge: 1, maxUnavailable: 0을 적용하되 컨테이너 시작 여부만 기준으로 삼으면 어떤 구체적인 재난이 생기는지 구성해 보자.
  \item 파드는 살아 있고(활성 상태 검사 통과) 준비는 끝내 되지 않으며(준비 상태 검사 실패) 재시작도 되지 않은 채 무기한 머물 수 있다. 이 끝나지 않는 상태가 기능인 이유를 논증한 뒤, 의도하지 않았을 때 §5.6의 어떤 실패가 되는지 말해 보자.
\end{tbbitemize}
\end{calloutbox}

\section{requests, limits, QoS: 플릿 규모로 확장한 제3장}

터미널 기록 5-3의 \code{resources:} 구절은 여덟 줄에 불과하지만, 나머지 매니페스트 전체를 합친 것보다 운영상 중요하다. 이해하기 까다로운 핵심은 \textbf{서로 다른 두 가지 일을 하는 두 시스템이 이 구절을 읽는다}는 점이다(그림 5-5). 배치 시점에 \textit{스케줄러}는 \textbf{requests}를 산술 값으로 읽는다. 노드의 \textbf{allocatable} 용량에서 이미 배치된 requests의 합을 빼고도 새 파드의 requests를 감당할 수 있으면 파드가 들어간다. requests는 그저 계획을 위한 숫자다. \textbf{런타임에는 절대 강제되지 않는다.} 노드에 여유가 있으면 1 CPU를 요청한 파드가 3을 사용해도 된다. 반면 \textit{kubelet}이 읽는 \textbf{limits}는 제4장의 런타임 스택을 거쳐 제3장의 cgroup 파일로 정확히 변환되며, 커널이 매 밀리초마다 강제한다.

\begin{tbbfigure}
  \adjustbox{max width=0.94\linewidth,max totalheight=0.46\textheight}{\includegraphics{figures/fig5-5_requests_limits.png}}\\[3pt]
  \figcaption{그림 5-5  하나의 구절을 읽는 두 주체. requests는 스케줄러의 빈 패킹에 쓰이고, limits는 cpu.max와 memory.max가 된다. 둘 사이의 간격은 의도적인 오버커밋이며, 노드 압박이 시작되는 지점이기도 하다.}
\end{tbbfigure}

\subsection*{핵심을 증명하는 기록: YAML을 넣으면 cgroup 파일이 나온다}

제3장은 호스트에서 Docker 컨테이너를 조사해 cgroup 파일이 연결된 프로세스를 찾아내며 시작했다. Kubernetes에서 같은 조사를 수행하면 이 장의 핵심 주장이 완성된다.

\begin{terminalbox}
\begin{Verbatim}[breaklines=true,breakanywhere=true,fontsize=\small,formatcom=\color{tbbtermfg},codes=\tbbcodeactive]
$ kubectl get pod model-server-7d4b9c6f5-8xk2p \
    -o jsonpath='{.spec.nodeName} {.status.qosClass}'
node-a Burstable
$ minikube ssh                        # onto the node itself
node-a $ ps aux | grep uvicorn        # the server is just a process (Ch. 3)
root  41320  2.1  7.8  ...  uvicorn serve:app --port 8000
node-a $ cat /proc/41320/cgroup
0::/kubepods.slice/kubepods-burstable.slice/kubepods-...-pod9f31.slice/...
node-a $ cd /sys/fs/cgroup/kubepods.slice/kubepods-burstable.slice/...
node-a $ cat memory.max
1879048192                            # = 1.75Gi, your YAML, in bytes
node-a $ cat cpu.max
200000 100000                         # = limits.cpu "2", Ch. 3 grammar
node-a $ cat cpu.weight
39                                    # from requests.cpu "1" - soft share
\end{Verbatim}
\end{terminalbox}
\figcaption{터미널 기록 5-7  연결 고리를 직접 완성했다. YAML의 limits는 노드의 memory.max와 cpu.max가 되었고, CPU 요청량은 cpu.weight가 되었다. Kubernetes는 제3장의 앞에 데이터베이스를 붙인 시스템이다.}

세 파일을 제3장과 연결해 읽으면 모든 동작이 자연스럽게 따라 나온다. \code{memory.max = 1879048192}는 파드가 \textbf{같은 OOM 킬러, 같은 exit 137, 같은 dmesg의 자백, 같은 크기 산정법}을 만난다는 뜻이다. §3.6이 그대로 적용되고 \code{kubectl describe}에는 \code{Last State: OOMKilled}로 나타난다. \code{cpu.max = "200000 100000"}는 파드가 \textbf{버스트 중에 스로틀링될 수 있음}을 뜻한다. §3.3에서 인위적으로 만든 p99가 이제 YAML 한 줄로 설정된다. requests에서 파생된 \code{cpu.weight}는 경합이 있을 때만 영향을 주는, 정중하고 작업 보존적인 지분이다. 기억해야 할 표는 다음과 같다.

\begin{tbbtablewrap}
  \begin{tabular}{@{}>{\raggedright\arraybackslash}m{\dimexpr 0.3085\tbltextwidth-2\tabcolsep\relax} >{\raggedright\arraybackslash}m{\dimexpr 0.3511\tbltextwidth-2\tabcolsep\relax} >{\raggedright\arraybackslash}m{\dimexpr 0.3404\tbltextwidth-2\tabcolsep\relax}@{}}
  \arrayrulecolor{tbbnavy}\specialrule{1pt}{0pt}{0pt}
  \rowcolor{tbbtablehead}{\centering \bfseries\color{tbbnavy}작성 내용(YAML)} & {\centering \bfseries\color{tbbnavy}변환 결과} & {\centering \bfseries\color{tbbnavy}성격} \\
  \arrayrulecolor{tbbnavy}\specialrule{0.7pt}{0pt}{2pt}
  {requests.cpu: "1"} & {스케줄러 산술 + cpu.weight} & {계획 + 유연한 지분, 작업 보존적} \\
  \arrayrulecolor{tbbtableborder}\hline
  {requests.memory: 1.2Gi} & {스케줄러 산술(+ 퇴거 순위)} & {계획에만 사용, 커널 객체 없음} \\
  \arrayrulecolor{tbbtableborder}\hline
  {limits.cpu: "2"} & {\textcolor{tbbamberdark}{cpu.max = "200000 100000"}} & {엄격한 할당량, 버스트에서 스로틀링(§3.3)} \\
  \arrayrulecolor{tbbtableborder}\hline
  {limits.memory: 1.75Gi} & {\textcolor{tbbred}{memory.max = 1879048192}} & {엄격한 벽, OOM killer(§3.6)} \\
  \arrayrulecolor{tbbtableborder}\hline
  {nvidia.com/gpu: 1} & {디바이스 플러그인 할당(§5.5)} & {정수, cgroup이 없어 중재할 커널 기능 없음} \\
  \arrayrulecolor{tbbnavy}\specialrule{1pt}{2pt}{0pt}
  \end{tabular}\par
  \figcaption{표 5-1  한 줄씩 살펴보는 변환 관계}
\end{tbbtablewrap}

\subsection*{CPU: request는 넉넉하게, limit는 신중하게}

CPU에서 두 설정값은 제3장에서 본 서로 다른 두 성격을 이어받으며, 서빙 지침도 기계적으로 도출된다. request는 \textit{예약}이다. 배치 위치를 결정하고, \code{cpu.weight}를 통해 경합 중에는 비례 배분의 하한을 보장하면서도 이웃이 놀 때는 자원을 낭비하지 않는다. limit는 \code{cpu.max}라는 \textit{상한}이다. §3.3의 계산에서 상한이 꼬리 지연 시간에 미치는 영향을 보았다. CPU 시간 30 ms가 필요한 요청은 0.2-CPU 할당량 아래에서 110–190 ms가 걸리지만 프로파일러에는 원인이 보이지 않는다. 따라서 \textbf{CPU requests는 정직하게 설정}해야 한다. 이웃 워크로드 간섭이 있을 때 p99를 지키는 보험이기 때문이다. 반면 지연 시간에 민감한 파드에서는 \textbf{CPU limits를 신중하게 사용}해야 한다. 버스트를 흡수하는 데 쓰인 유휴 코어가 p99 구간에서 스로틀링되는 것보다 저렴하다. 다만 현실에는 두 가지 복잡한 점이 있다. 플랫폼 팀이 비용 관리와 폭주 방지를 위해 limits를 \textit{요구}할 때가 많고, 다음 문단에서 볼 Guaranteed QoS는 requests = limits를 요구한다. 일정한 성능과 QoS가 모두 필요한 팀은 두 값을 똑같이 넉넉하게 잡거나 kubelet의 정적 CPU 매니저 정책을 사용한다. 이 정책은 Guaranteed 등급의 정수-CPU 파드에 \textit{전용 코어}를 준다. 제2장의 전용 코어 개념이 kubelet에서 다시 등장한 것이다.

\subsection*{메모리: 간격, 압박, 두 번째 처형자}

메모리에서는 requests와 limits의 간격을 설명하는 별도 어휘가 필요하다. 스케줄러는 노드마다 Σrequests ≤ allocatable을 보장하지만 limits에는 관여하지 않으므로 \textbf{Σlimits가 노드 용량을 넘는 일은 흔하다.} 이것이 오버커밋이며 기능이다. 모두가 동시에 버스트하지는 않으리라는 가정 아래 누구나 버스트할 수 있게 한다. 이 가정이 틀리면, 예를 들어 여러 레플리카의 사용량이 함께 치솟으면(§3.6에서 보았듯 추론 메모리는 \textit{본질적으로 급격히 변한다}) 개별 파드가 자기 벽에 닿기 전에 \textbf{노드 자체}의 메모리가 부족해진다. 커널의 노드 수준 OOM 킬러에 맡기면 혼란스러우므로 Kubernetes는 질서를 택한다. \textbf{kubelet이 퇴거}를 수행한다. QoS 클래스와 requests 초과 사용량을 기준으로 희생 대상을 고르고 정상적으로 종료한 뒤 \code{Evicted}로 표시한다. 이제 서로 다른 두 처형자를 알게 되었다. 둘을 구별하는 능력은 §5.6의 핵심이다.

\begin{tbbtablewrap}
  \begin{tabular}{@{}>{\raggedright\arraybackslash}m{\dimexpr 0.2340\tbltextwidth-2\tabcolsep\relax} >{\raggedright\arraybackslash}m{\dimexpr 0.3830\tbltextwidth-2\tabcolsep\relax} >{\raggedright\arraybackslash}m{\dimexpr 0.3830\tbltextwidth-2\tabcolsep\relax}@{}}
  \arrayrulecolor{tbbnavy}\specialrule{1pt}{0pt}{0pt}
  \rowcolor{tbbtablehead}{\centering \bfseries\color{tbbnavy}} & {\centering \bfseries\color{tbbnavy}cgroup OOM 종료(제3장)} & {\centering \bfseries\color{tbbnavy}노드 압박(node pressure) 퇴거} \\
  \arrayrulecolor{tbbnavy}\specialrule{0.7pt}{0pt}{2pt}
  {\textbf{발동 조건}} & {파드가 자기 memory.max를 초과} & {노드 메모리 부족(오버커밋 가정 실패)} \\
  \arrayrulecolor{tbbtableborder}\hline
  {\textbf{처형자}} & {커널 OOM 킬러 — SIGKILL} & {kubelet — 파드를 정상 절차로 종료} \\
  \arrayrulecolor{tbbtableborder}\hline
  {\textbf{영향 범위}} & {컨테이너 하나} & {QoS 순서에 따라 고른 희생 대상, 노드 전체} \\
  \arrayrulecolor{tbbtableborder}\hline
  {\textbf{징후}} & {\textcolor{tbbred}{OOMKilled · exit 137 · dmesg}} & {\textcolor{tbbpurple}{상태 Evicted · "노드 메모리가 부족함"}} \\
  \arrayrulecolor{tbbtableborder}\hline
  {\textbf{우선 해결책}} & {최댓값에 맞춰 크기 조정 또는 배치 제한(§3.6)} & {정직한 requests, 서빙에 Guaranteed QoS 적용} \\
  \arrayrulecolor{tbbnavy}\specialrule{1pt}{2pt}{0pt}
  \end{tabular}\par
  \figcaption{표 5-2  두 처형자와 두 가지 징후}
\end{tbbtablewrap}

\subsection*{QoS 클래스: 누가 먼저 죽을지도 설정이다}

kubelet의 퇴거 순서와, 그 단계까지 간다면 노드 OOM 킬러의 점수는 추측으로 정해지지 않는다. 플랫폼이 requests/limits 패턴에 따라 할당하는 제3장의 \code{oom\_score\_adj}다. 모든 컨테이너의 CPU와 메모리에서 requests = limits인 \textbf{Guaranteed}는 ≈ −997을 받아 거의 건드릴 수 없다. 아무것도 설정하지 않은 \textbf{BestEffort}는 +1000을 받아 언제나 가장 먼저 희생된다. 그 사이에 있는 \textbf{Burstable}은 사용량이 requests를 얼마나 초과했는지에 따라 점수가 달라진다. 서빙에 적용할 원칙은 명확하다. \textbf{모델 서버는 Guaranteed로 실행}해야 한다. 최소한 memory requests = limits로 맞춘다. 그러면 오버커밋된 노드가 압박을 받을 때 p99를 책임지는 파드가 아니라 배치 작업과 BestEffort 실험이 죽는다. (Kubernetes에는 이와 독립적인 두 번째 설정인 PriorityClass도 있다. 이는 \textit{워크로드 사이의 선점과 퇴거} 순서를 정하고, QoS는 \textit{한 노드에서 종료할 우선순위}를 정한다. 12주 차 오토스케일링에서는 둘 다 사용한다.)

\begin{calloutbox}{tbbblue}{tbbbluefill}
{\normalsize\bfseries\color{tbbblue}계산 예제 — 직접 해 보는 노드 패킹}\par\vspace{3pt}
범용 8-vCPU / 32 GiB 노드를 생각해 보자. 시스템과 kubelet 예약분을 빼면 \textbf{allocatable(할당 가능 용량)은 ≈ 7.5 vCPU 및 29 GiB}다(터미널 기록 5-8에서 확인 위치를 보여 준다). 노드마다 실행되는 데몬(로그 수집기, 모니터링 에이전트, kube-proxy)은 ≈ 0.3 vCPU 및 0.5 GiB를 요청한다. 강의 모델 서버의 requests는 \code{cpu: 1, memory: 1.2Gi}이고 limits는 2 / 1.75Gi다.

\begin{tbbitemize}
  \item \textbf{CPU 계산:} (7.5 − 0.3) / 1 = CPU requests 기준으로 \textbf{레플리카 7개}가 들어간다.
  \item \textbf{메모리 계산:} (29 − 0.5) / 1.2 ≈ 메모리 requests 기준으로 \textbf{레플리카 23개}가 들어간다.
\end{tbbitemize}

이 워크로드에서 노드는 \textbf{CPU 병목}이다. 노드마다 레플리카 7개가 들어가고 메모리는 남는다. 따라서 메모리가 많은 인스턴스 유형에 비용을 더 내면 낭비다. 제1장의 비용 관점처럼 인스턴스는 제약이 되는 차원에 맞춰 골라야 한다. 이제 오버커밋 노출을 확인하자. 레플리카 7개 × limits = 14 vCPU 및 12.25 GiB다. CPU limits(14)는 allocatable(7.5)을 넘지만 합법이다. 동시에 버스트하면 requests에 비례해 cpu.weight가 분배하며, p99는 requests가 얼마나 정직한지에 달렸다. 메모리 limits(12.25 GiB)는 29 GiB보다 훨씬 작다. 이 패킹은 원리상 노드 메모리 압박을 일으킬 수조차 없으며, 서빙 노드에 의도적으로 적용한 보수적인 자세다.

\code{cpu: 2, memory: 9Gi}를 요청하는 7B 매개변수 INT8 LLM 서버에 같은 계산을 적용해 보자. 이제 메모리가 제약이 되어 (29 − 0.5) / 9 = \textbf{레플리카 3개}가 들어가고, 적합한 노드 계열은 컴퓨팅 최적화에서 메모리 최적화로 바뀐다. 2분간 나눗셈만 해도 인스턴스 계열을 결정할 수 있다. 12주 차 FinOps 절은 바로 이 습관을 바탕으로 전체 비용 모델을 세운다.
\end{calloutbox}

\vspace{3.0pt}

\begin{calloutbox}{tbbpurple}{tbbpurplefill}
{\normalsize\bfseries\color{tbbpurple}생각해 보기}\par\vspace{3pt}
\begin{tbbitemize}
  \item 파드가 requests(1 CPU / 1 Gi)는 설정했지만 limits는 설정하지 않았다. 이 파드의 QoS 클래스는 무엇인가? 다른 파드가 놀고 있는 노드에서 이 파드가 메모리를 누수하면 어떻게 되는가? 반대로 \textit{이웃}이 누수하면 누가 먼저 위험해지는가?
  \item 스케줄러가 limits를 완전히 무시하는 이유는 무엇인가? 오버커밋이 기능이라는 관점에서 논리를 구성한 뒤, limits를 기준으로 스케줄링하는 편이 분명히 나은 워크로드 조합 하나를 들어 보자.
  \item 계산 예제의 레플리카 7개 모두 피크 트래픽에서 nr\_throttled가 증가한다(limits.cpu 2 × 7 = 할당 가능 용량 7.5에서 14). §3.3과 이 절을 바탕으로 p99 징후를 설명하고, 비용이 서로 다른 해결책 두 가지를 제안하자.
\end{tbbitemize}
\end{calloutbox}

\section{GPU와 스케줄러}

지금까지는 노드를 서로 바꿔 쓸 수 있다고 가정했다. GPU가 등장하는 순간 배치는 비용이 큰 결정으로 바뀐다. GPU 플릿을 잘 배치했을 때와 잘못 배치했을 때의 차이는 매달 수천 달러에 이른다. CPU와 달리 \textbf{GPU에는 오버커밋이 없으므로} 실수를 오버커밋으로 흡수할 수도 없다. 이 절에서는 가장 까다로운 사례인 GPU를 중심으로 스케줄러를 살펴본다.

\subsection*{필터링, 점수화, 바인딩 — 배치 과정 자체}

스케줄러는 바인딩되지 않은 파드마다 두 단계의 과정을 수행한다(그림 5-6). \textbf{필터링}은 파드를 실행할 수 \textit{없는} 모든 노드를 제거한다. requests를 감당할 할당 가능 리소스가 부족하거나, 필요한 확장 리소스가 없거나, nodeSelector 또는 어피니티(affinity)가 맞지 않거나, 테인트를 톨러레이션하지 않는 노드가 탈락한다. 남은 노드가 실행 가능 집합이다. 이 집합이 비면 파드는 \textbf{Pending}에 머문다. 특히 \textit{이벤트 메시지는 어떤 필터에서 몇 개 노드가 실패했는지 정확히 열거한다}는 점이 중요하다(§5.6에서 실제 메시지를 읽는다). 이어서 \textbf{점수화}는 가용 영역(availability zone) 분산, 부하가 더 적거나 더 많은 노드 선호, 이미지 존재 여부 같은 유연한 선호 조건으로 실행 가능 집합의 순위를 매긴다. 마지막으로 \textbf{바인딩}이 승자의 이름을 파드 객체에 기록한다. 스케줄러의 출력은 객체 하나의 필드 하나이며, 이후 과정은 모두 §5.2의 연결 고리다.

\begin{tbbfigure}
  \adjustbox{max width=0.94\linewidth,max totalheight=0.46\textheight}{\includegraphics{figures/fig5-6_scheduler.png}}\\[3pt]
  \figcaption{그림 5-6  GPU 파드의 스케줄링 과정. 필터가 노드를 제거하고 실패 이유는 Pending 이벤트가 된다. 점수가 순위를 매기고 바인딩이 이름을 기록한다. 패킹과 분산 중 무엇을 택할지는 GPU만큼 값비싼 정책 결정이다.}
\end{tbbfigure}

\subsection*{GPU는 확장 리소스다: 정수, 오버커밋 없음, requests = limits}

GPU는 \textbf{확장 리소스}로 Kubernetes에 들어온다. 각 GPU 노드에서 실행되는 제4장의 디바이스 플러그인이 kubelet에 \code{nvidia.com/\allowbreak gpu: 4}를 알리고, kubelet이 이를 노드 용량에 보고한다. 파드는 구절 하나로 GPU를 소비한다. 이때 CPU나 메모리보다 엄격한 세 규칙이 적용된다. 수량은 \textbf{정수}여야 하고, \textbf{requests와 limits가 같아야 하며}, \textbf{어떤 경우에도 오버커밋은 없다.} 파드가 시작되면 디바이스 플러그인이 인덱스/UUID로 \textit{특정} GPU를 할당해 컨테이너에 주입한다. 제4장에서 본 \code{/\allowbreak dev/\allowbreak nvidia*} 배관이다.

\begin{terminalbox}
\begin{Verbatim}[breaklines=true,breakanywhere=true,fontsize=\small,formatcom=\color{tbbtermfg},codes=\tbbcodeactive]
# in the pod spec - the entire GPU ask:
        resources:
          limits:
            nvidia.com/gpu: 1        # integer; requests implied equal

$ kubectl describe node gpu-node-3 | sed -n '/Capacity/,/pods/p'
Capacity:
  cpu:                48
  memory:             196Gi
  nvidia.com/gpu:     4
Allocatable:
  cpu:                47
  memory:             186Gi
  nvidia.com/gpu:     4
$ kubectl describe node gpu-node-3 | grep -A4 "Allocated resources"
Allocated resources:
  Resource           Requests      Limits
  cpu                18 (38%)      36 (76%)
  memory             88Gi (47%)    120Gi (64%)
  nvidia.com/gpu     2             2          # 2 of 4 in use
\end{Verbatim}
\end{terminalbox}
\figcaption{터미널 기록 5-8  스케줄러가 보는 GPU 노드. 디바이스 플러그인이 Capacity/Allocatable에 알린 값과 현재 GPU 점유 상태가 보인다. CPU limits는 오버커밋되어 있지만 GPU는 그렇지 않다.}

왜 이렇게 엄격할까? 제3장에서 이미 답을 보았다. \textbf{cgroup에는 GPU 컨트롤러가 없다.} CPU는 cpu.weight가 경합을 작업 보존적으로 중재하므로 오버커밋이 안전하다. 메모리는 memory.max와 OOM 킬러가 벽을 강제한다. GPU에는 두 컨테이너의 연산을 중재하거나 VRAM을 분할할 커널 메커니즘이 없다. 따라서 Kubernetes는 아래 계층에서 강제할 수 없는 것을 약속하지 않는다. GPU 하나에는 소유자 하나만 둔다. 작은 INT8 모델이 A100 전체를 채우는 일은 드물기 때문에 GPU의 일부만 필요한 \textit{수요}는 실제로 존재한다. 그러나 해답은 이 절의 스케줄러 밖에 있다. \textbf{시간 분할}은 디바이스 플러그인이 GPU마다 N개의 가상 복사본을 알리는 방식으로, 격리 없이 배치만 하므로 이웃 워크로드가 보호 장치 없이 VRAM을 공유한다. \textbf{MPS}는 공간을 공유하지만 여전히 협력적인 방식이다. \textbf{MIG}는 제2장의 하드웨어 파티션으로, 각각 \code{nvidia.com/\allowbreak mig-1g.5gb} 같은 독립 확장 리소스로 알려져 실제 격리를 제공하며 다른 리소스처럼 스케줄링할 수 있다. 12주 차에는 이 선택지를 실제 설정으로 바꾼다. 이 장에서는 이 선택지가 \textit{왜} 존재하는지 이해하면 된다.

\subsection*{배치 제어: 테인트는 밀어내고 어피니티는 끌어당긴다}

두 가지 메커니즘이 특수 노드를 기준으로 파드의 위치를 유도하며, GPU 플릿에는 양방향 모두가 필요하다. 노드의 \textbf{테인트}는 이를 명시적으로 톨러레이션하지 않는 파드를 밀어낸다. 관리형 GPU 노드 풀에 \code{nvidia.com/\allowbreak gpu=present:NoSchedule} 테인트가 기본으로 붙는 이유는, CPU가 남았다는 이유만으로 무관한 로그 처리 Deployment가 시간당 \$3짜리 노드를 차지하지 못하게 하기 위해서다. \textbf{톨러레이션(toleration)}은 문을 열 뿐 파드를 특정 방향으로 보내지는 않는다. 방향은 레이블을 기준으로 한 \textbf{nodeSelector / 노드 어피니티}가 정한다. \code{gpu=t4}와 \code{gpu=a100}의 차이는 중요하다. 제2장의 플릿 표가 실제 비용으로 이어지기 때문이다. T4의 경제성을 기준으로 검증한 모델을 A100에서 실행하면 낭비일 수 있고, 10주 차의 양자화 모델은 모델마다 적합한 답을 바꾼다. 따라서 GPU 파드에는 톨러레이션(허가)과 어피니티(의도)가 모두 들어간다.

\subsection*{패킹 대 분산: 스케줄링 정책도 비용 항목이다}

상태 없는 CPU 플릿에서는 기본적인 점수 계산 원칙, 즉 파드를 분산하고 노드 부하를 고르게 유지하는 전략이 건전하다. 이웃 워크로드 간섭을 완화하고 노드 손실도 자연스럽게 견딘다. 하지만 GPU 플릿에서는 두 가지 방식으로 조용히 비용이 커진다. 첫째는 \textbf{단편화}다. 1-GPU 파드를 여러 노드에 흩으면 모든 노드가 일부만 사용된다. 나중에 다중 GPU 파드가 와도 여유 GPU가 충분한 단일 노드를 찾지 못해, 전체 여유 용량은 많은데도 Pending이 된다. 둘째는 \textbf{유휴 고립}이다. GPU가 없는 파드가 CPU나 메모리를 차지한 GPU 노드에서는 GPU가 비어 있고 비용도 내고 있지만 사용할 수 없다. 해결책은 코드가 아니라 정책이다. GPU 노드에 \textbf{빈 패킹} 점수(스케줄러의 MostAllocated 전략)를 적용해 사용 중인 노드는 채우고 빈 노드는 계속 비워 둔다. 그러면 12주 차의 클러스터 오토스케일러(cluster autoscaler)가 완전히 유휴 상태인 노드를 삭제할 수도 있다. GPU 노드에는 테인트를 붙여 CPU 위주의 작업이 들어오지 못하게 하고, GPU 네 개 가운데 두 개만 쓰고도 CPU가 바닥나지 않도록 GPU 파드의 CPU requests도 정직하게 설정한다.

\begin{calloutbox}{tbbamber}{tbbamberfill}
{\normalsize\bfseries\color{tbbamberdark}계산 예제 — 단편화의 비용}\par\vspace{3pt}
클러스터에 A100 4개를 갖춘 노드 4대가 있다(온디맨드 ≈ \textbf{\$3/GPU-hour}, 대표값). GPU 1개를 쓰는 서빙 파드 16개가 들어오면 분산 점수 전략을 쓰는 스케줄러가 노드마다 4개씩 배치한다. 깔끔하게 전부 사용되어 \textit{아직은} 낭비가 없다. 밤이 되어 트래픽이 줄고 플릿이 파드 6개로 축소된다. 분산 전략은 노드마다 파드를 약 1–2개씩 유지한다. \textbf{노드 4대가 계속 실행되고 유휴 GPU 10개에도 계속 비용이 청구된다.} ≈ 10 × \$3 × 24 ≈ \textbf{\$720/day의 유휴 비용}이며, 모든 노드에 파드가 있어 오토스케일러는 아무것도 삭제할 수 없다.

같은 클러스터를 빈 패킹하면 파드 6개가 노드 1.5대를 채운다. 노드 2대만 실행하고 나머지 2대는 비운 뒤 삭제할 수 있어 유휴 비용이 ≈ \$144/day로 줄어든다. 새벽에 GPU 4개를 쓰는 학습 작업이 들어오면 패킹 배치에는 통째로 비어 있는 노드 하나를 줄 수 있다. 반면 분산 배치는 GPU 10개가 조각난 채 비어 있는데도 \textbf{Pending — 노드 4대 가운데 여유 GPU 4개를 갖춘 노드는 0대}라고 보고한다. 단편화는 비용표가 붙은 Pending이다.
\end{calloutbox}

\vspace{3.0pt}

\begin{calloutbox}{tbbpurple}{tbbpurplefill}
{\normalsize\bfseries\color{tbbpurple}생각해 보기}\par\vspace{3pt}
\begin{tbbitemize}
  \item GPU requests는 limits와 같아야 하지만 CPU requests와 limits는 달라도 되는 이유를 제3장의 컨트롤러 목록을 근거로 정확히 설명하자.
  \item 로그 수집용 DaemonSet(노드마다 파드 하나)이 CPU 2개를 요청한다. 터미널 기록 5-8의 48-CPU GPU 노드에서는 별문제 없어 보인다. 노드 20개로 구성된 GPU 플릿에서 vCPU-hour당 \textasciitilde{}\$0.05를 기준으로 연간 비용을 계산한 뒤 더 저렴한 설계를 제시하자.
  \item 모델 A(지연 시간에 민감하고 T4에 맞게 양자화됨)와 모델 B(배치 작업, A100-80GB 필요), 범용 CPU 마이크로서비스를 서빙하는 클러스터의 레이블, 테인트, 톨러레이션을 설계하자. 어디에는 중복 배치를 허용하고 어디에는 금지할 것인가?
\end{tbbitemize}
\end{calloutbox}

\section{죽어 가는 파드 현장 안내서}

단일 머신에서 "고장 났다"는 말이 나오면 살펴볼 곳도 하나다. 클러스터에서는 같은 소식이 파드 \textbf{상태}, 객체 \textbf{이벤트}, 컨테이너 \textbf{로그}, Service \textbf{엔드포인트}처럼 서로 다른 여러 표면으로 들어온다. 어떤 질문에 어느 표면이 답하는지 아는 것이 기술이다. 다행히 실제로는 \textbf{여섯 가지 상태가 거의 모든 서빙 사고를 포괄하며}(그림 5-7), 각각 뚜렷한 징후와 가장 먼저 실행할 표준 명령이 있다. 이 절에서는 가장 교훈적인 사례 하나를 처음부터 끝까지 재현하고, 나머지 다섯 가지를 정리한다. 전체에 걸쳐 한 가지 반사 행동을 익히자. \textbf{먼저 kubectl describe pod를 실행한다. 보통 Events 절에 판결문이 있다.}

\begin{tbbfigure}
  \adjustbox{max width=0.94\linewidth,max totalheight=0.46\textheight}{\includegraphics{figures/fig5-7_fieldguide.png}}\\[3pt]
  \figcaption{그림 5-7  여섯 가지 대표 실패 상태. 무엇이 보이고, 무엇을 뜻하며, 어디부터 살펴봐야 하는지 정리했다. 이번 주 실습에서 여섯 가지를 모두 의도적으로 재현한다.}
\end{tbbfigure}

\subsection*{사고의 전 과정: 정상 이미지에서 발생한 CrashLoopBackOff}

상황은 이렇다. 팀원이 새 모델을 로딩하는 데 \textasciitilde{}90초가 걸리는 \code{model-server:v9}를 배포했다. \code{docker run}에서는 잘 동작했다. 하지만 startupProbe는 설정하지 않고 성급한 livenessProbe를 그대로 둔 Deployment를 클러스터에서 실행하자 다음 결과가 나온다.

\begin{terminalbox}
\begin{Verbatim}[breaklines=true,breakanywhere=true,fontsize=\small,formatcom=\color{tbbtermfg},codes=\tbbcodeactive]
$ kubectl get pods
NAME                            READY   STATUS             RESTARTS   AGE
model-server-59f6c9d8b7-qv4tn   0/1     CrashLoopBackOff   5          8m
# five restarts in eight minutes. First reflex: describe, read Events.
$ kubectl describe pod model-server-59f6c9d8b7-qv4tn
...
Events:
  Warning  Unhealthy  Liveness probe failed: Get "http://10.244.1.61:8000
           /healthz": context deadline exceeded
  Normal   Killing    Container server failed liveness probe, will be
           restarted
  Normal   Pulled     Container image already present on machine
  Normal   Started    Started container server
# the kubelet is the killer - and it says so. Second reflex:
$ kubectl logs model-server-59f6c9d8b7-qv4tn --previous | tail -3
INFO  loading weights shard 5/8 ...
INFO  loading weights shard 6/8 ...
INFO  loading weights shard 7/8 ...
\end{Verbatim}
\end{terminalbox}
\figcaption{터미널 기록 5-9  포렌식 연결 고리. CrashLoopBackOff → describe가 활성 상태 프로브를 처형자로 지목 → logs --previous에는 느리게 로딩했을 뿐인 프로세스가 보인다. 잘못을 저지른 것은 아무것도 없으므로 어디에도 오류는 나타나지 않는다.}

제3장에서 OOM 연결 고리를 분석한 방식으로 읽어 보자. \textbf{상태}는 컨테이너가 계속 죽는다고 말한다. \code{CrashLoopBackOff}는 원인이 아니라 재시작 사이의 대기 시간을 두 배씩 늘리는 재시작 \textit{정책}이다. \textbf{이벤트}는 처형자를 지목한다. 단일 스레드로 가중치를 읽느라 서버가 답할 수 없을 때 응답을 요구한 활성 상태 프로브를 대신해 kubelet이 죽였다. \textbf{이전 로그}는 비행 기록 장치다. 일반 \code{logs}는 막 재시작한 \textit{현재} 시도를 보여 주지만 이전 로그는 오류 하나 없이 로딩 도중에 끝난다. 종료하고, 다시 로딩하고, 다시 종료한다. 재시작할 때마다 전체 로딩 비용을 다시 치르고(제3장의 재시작 루프 경제학), 파드는 한 번도 Ready가 되지 못했다. 해결책은 최악의 로딩 시간을 포괄하는 예산을 둔 startupProbe, 즉 터미널 기록 5-3의 구절이다. 적용하면 같은 이미지가 조용히 90초를 보낸 뒤 Ready가 된다. 이 사고보다 오래 남을 일반 교훈이 있다. \textbf{Kubernetes에서는 플랫폼 자체도 실패 원인 가운데 하나다.} 프로브, 상한, 퇴거는 정상 프로세스를 죽일 수 있으며, 증거는 애플리케이션 로그가 아니라 이벤트에 남는다.

\subsection*{Pending: 스케줄러가 풀이 과정을 보여 준다}

\begin{terminalbox}
\begin{Verbatim}[breaklines=true,breakanywhere=true,fontsize=\small,formatcom=\color{tbbtermfg},codes=\tbbcodeactive]
$ kubectl get pods
NAME                            READY   STATUS    RESTARTS   AGE
llm-server-7b9d5f4c66-x2jkp     0/1     Pending   0          3m
$ kubectl describe pod llm-server-7b9d5f4c66-x2jkp | tail -4
Events:
  Warning  FailedScheduling  0/6 nodes are available:
    4 Insufficient nvidia.com/gpu,
    2 node(s) had untolerated taint {team: batch}.
\end{Verbatim}
\end{terminalbox}
\figcaption{터미널 기록 5-10  Pending은 스케줄러가 필터 보고서를 게시한 상태다(§5.5). 노드 여섯 개가 각각 어떤 필터에서 실패했는지 정확히 보여 준다. 파드에는 아무 문제가 없고, 클러스터가 파드의 제약 조건을 충족하지 못할 뿐이다.}

Pending은 시스템에서 가장 충실하게 스스로 설명하는 상태다. §5.5의 필터 단계가 실패 집계를 이벤트에 그대로 기록하기 때문이다. \textit{"각 노드가 탈락한 이유는 다음과 같다"}라고 읽으면 된다. 해결책은 단계적으로 나뉜다. requests를 줄이고(정말 GPU 4개가 필요한가?), 용량을 비우거나 추가하고(12주 차 오토스케일러가 바로 이 일을 자동화한다), 주소 지정 문제(톨러레이션과 레이블)를 고친다. 해결책이 아닌 것도 하나 기억하자. 파드를 삭제하고 다시 만들면 같은 클러스터에 같은 필터를 다시 실행하므로 또 Pending이 된다. 증상이 아니라 제약 조건을 고쳐야 한다.

\subsection*{나머지 네 가지 상태}

\begin{tbbitemize}
  \item \textbf{ImagePullBackOff} — 노드가 이미지를 가져올 수 없다. 태그가 틀렸거나, 비공개 레지스트리에 필요한 \code{imagePullSecrets}가 없거나, 노드에서 레지스트리에 접근할 수 없는 경우다. 제4장의 레지스트리 배관 문제이며, \code{describe} events에 원본 pull 오류가 들어 있다. (미묘한 서빙 변형도 있다. tag는 존재하지만 이미지 크기가 8 GB이고 노드 디스크가 가득 찬 경우로, 이벤트가 그대로 알려 준다.)
  \item \textbf{OOMKilled / exit 137} — 제3장 §3.6의 상황을 플랫폼이 다시 보여 준다. \code{describe}에는 \code{Last State: Terminated, Reason: OOMKilled, Exit Code: 137}이 나타난다. Kubernetes에는 한 가지 변형이 있다. \code{restartPolicy}가 \textit{같은} limit로 컨테이너를 되살리므로, 너무 작은 limit는 매번 모델 로딩 콜드 스타트 비용을 치르는 충돌 루프가 된다. §3.6의 순진한 정책에 관성이 붙은 셈이다. 해결 방법은 똑같다. 측정한 최댓값에 맞게 크기를 조정하거나 변동 부분을 제한해야 하며, try/except로 처리해서는 안 된다.
  \item \textbf{Evicted} — \textit{노드}의 자원이 부족해졌다(표 5-2). 상태는 \code{Evicted}, 메시지는 "노드의 메모리 리소스가 부족함"으로 나온다. 파드는 잘못이 없었을 수도 있으며, 희생 대상은 QoS가 결정했다. 퇴거가 반복되면 노드 어딘가의 requests가 정직하지 않다는 뜻이다. 네임스페이스를 감사하고 서빙 파드를 Guaranteed로 옮긴다.
  \item \textbf{실행 중이지만 라우팅되지 않음} — 조용한 상태이자 목록에서 어디에도 오류가 없는 유일한 상태다. 파드가 실행되고 활성 상태 검사를 통과하지만 준비 상태가 끝내 true가 되지 않거나 레이블이 Service 셀렉터와 일치하지 않는다. 따라서 \textbf{엔드포인트 목록이 비고} Service는 트래픽을 보낼 곳이 없다. 클라이언트에는 연결 오류가 발생하지만 모든 파드 대시보드는 초록색이다. 이를 드러내는 명령은 \code{kubectl get endpoints model-svc} 하나다. 저자의 경험상 이 명령을 확인하는 반사 행동이 이 장에서 얻을 수 있는 가장 가치 있는 습관이다.
\end{tbbitemize}

\begin{calloutbox}{tbbpurple}{tbbpurplefill}
{\normalsize\bfseries\color{tbbpurple}다섯 가지 반사 행동(모니터에 붙여 두자)}\par\vspace{3pt}
\begin{tbbitemize}
  \item \code{kubectl describe pod <p>} — 언제나 가장 먼저 실행한다. Events에 판결문이 있다(프로브, 스케줄링, pull, OOM, 퇴거).
  \item \code{kubectl logs <p> --previous} — 비행 기록 장치다. 새 시도가 아니라 죽은 시도가 남긴 마지막 말을 보여 준다.
  \item \code{kubectl get endpoints <svc>} — 대시보드는 초록색인데 아무 응답이 없다면, 반증될 때까지 엔드포인트가 비었다고 본다.
  \item \code{kubectl exec -it <p> -- sh} — 내부로 들어가(제3장에서 이것이 setns라고 배웠다) /ready를 직접 시험하고 환경 변수와 경로를 확인한다.
  \item \code{kubectl port-forward pod/\allowbreak <p> 8000:8000} — Service를 우회해 특정 파드 하나와 통신함으로써 "파드 고장"과 "라우팅 고장"을 구분한다.
\end{tbbitemize}
\end{calloutbox}

\vspace{3.0pt}

\begin{calloutbox}{tbbpurple}{tbbpurplefill}
{\normalsize\bfseries\color{tbbpurple}생각해 보기}\par\vspace{3pt}
\begin{tbbitemize}
  \item 컨테이너가 CrashLoopBackOff를 보이고 logs --previous는 오류 없이 끝나며 exit code는 137이다. 서로 다른 두 처형자가 137을 만든다(제3장의 SIGKILL 계산). 둘을 말하고, 구별할 수 있는 정확한 describe 필드를 제시하자.
  \item 팀의 readinessProbe가 데이터베이스 의존성을 검사한다. 데이터베이스가 30초 동안 잠깐 멈추자 모든 파드가 동시에 준비되지 않은 상태가 되고 → 엔드포인트가 비며 → 캐시된 응답을 제공할 수 있었던 서비스까지 완전히 중단된다. 준비 상태 프로브(readiness probe)에서 의존성을 검사하는 방식의 장단점을 논하고 절충안을 제안하자.
  \item "모델 서버는 정상이고 Service 셀렉터가 잘못되었다"는 사실을 1분 안에 증명할 명령 순서를 설계하자. (port-forward와 get endpoints가 유용하다.)
\end{tbbitemize}
\end{calloutbox}

\section{실습: 모델 서버 오케스트레이션}

이번 주 실습에서는 제4장의 이미지를 \code{docker run}에서 관리형 플릿으로 승격한 뒤, §5.6의 모든 방식으로 의도적으로 망가뜨린다. 모든 과정은 노트북에서 실행된다. \code{minikube start --cpus 4 --memory 8g}를 사용하며 kind도 똑같이 동작한다. 이미지는 강의 레지스트리에서 제공하고, 공개 대체 태그는 유인물에 있다. 이번 주에는 클라우드 비용이 필요 없다. §5.5의 GPU 터미널 기록은 시연 전용이며, 이미 크레딧이 있는 팀을 위해 관리형 클러스터에서 재현하는 선택 부록을 제공한다. \code{kubectl create namespace ml}로 만든 네임스페이스 안에서 작업한다.

\subsection*{파트 A — 배포하고 노출한 뒤 루프를 신뢰하기}

터미널 기록 5-3의 Deployment와 터미널 기록 5-6의 Service를 적용한다(minikube에서는 NodePort 유형). \code{get pods -o wide}로 플릿을 확인한 뒤, \textbf{자기 플릿에서 터미널 기록 5-1의 실험을 다시 실행}한다. curl 루프 도중 파드를 삭제하고 조정 과정을 기록한다. 이어서 터미널에서 Service에 초당 요청 하나를 보내는 curl 루프를 실행한 채 \code{:v8}로 롤링 업데이트한다. 제출물의 핵심 장면은 전체 롤아웃 동안의 루프 출력이다. 준비 상태가 속도를 제어하고 제3장의 SIGTERM 핸들러가 기존 파드의 처리 중인 요청을 비우므로 \textbf{실패한 요청은 0개}여야 한다. (그런 다음 의도적으로 망가뜨린다. readinessProbe를 제거하고 \code{:v9-slow}로 롤아웃해 같은 curl 루프에서 오류가 나는 모습을 확인한다. 준비 상태의 효과를 증명하는 음성 대조군이다.)

\subsection*{파트 B — 직접 검증하는 변환 관계}

터미널 기록 5-7을 직접 재현한다. \code{minikube ssh}로 접속해 \code{/\allowbreak sys/\allowbreak fs/\allowbreak cgroup/\allowbreak kubepods.slice/\allowbreak } 아래의 cgroup 트리를 따라가고, YAML의 숫자가 \code{memory.max}와 \code{cpu.max}에 들어 있는 모습을 캡처한다. 자기 터미널에서 이 장의 핵심을 증명하는 것이다. 그런 다음 OOM을 재현한다. 표준 모델에 \code{limits.memory: 256Mi}를 설정해 적용하고, Kubernetes식 증거 연결 고리를 수집한다. \code{describe}의 \code{OOMKilled /\allowbreak  137}과 함께 충돌 루프가 매번 모델 로딩 콜드 스타트 비용을 치르면서 재시작 횟수가 증가하는 모습을 담는다. 정직한 상한으로 복구하고, §3.6의 어휘를 사용해 전후 차이를 보고서에 기록한다.

\subsection*{파트 C — 망가뜨리고, 진단하고, 고치기}

세 가지 사고를 재현한다. 각 사고는 해결책을 적용하고 어떤 절에서 이를 예측했는지 한 문장으로 밝히며 끝낸다. \textbf{(1) Pending:} GPU가 없는 minikube에서 \code{nvidia.com/\allowbreak gpu: 1}을 요청한다. FailedScheduling 이벤트를 캡처하고 필터 보고서를 소리 내어 읽은 뒤(§5.5), 요청을 제거해 고친다. \textbf{(2) 프로브 사고:} 환경 변수 \code{SLOW\_LOAD=90}을 설정해 강의 이미지가 느린 모델 로딩을 모의하게 하고, startupProbe를 삭제하며, 활성 상태 프로브를 \code{periodSeconds: 5, failureThreshold: 2}로 강화한다. 터미널 기록 5-9의 전체 포렌식 연결 고리인 상태, 이벤트, \code{logs --previous}를 재현한다. startupProbe 설정으로 고치고 같은 이미지가 Ready가 되는 모습을 보여 준다. \textbf{(3) 조용한 실패:} 준비 상태 경로를 서버가 제공하지 않는 오타 \code{/\allowbreak readyz}로 바꾼다. Running 상태 파드, 정상인 활성 상태 검사, 비어 있는 \code{get endpoints}, 실패하는 클라이언트 curl 요청을 관찰한다. 수정한 뒤 엔드포인트가 다시 채워지는 모습을 확인한다.

\begin{calloutbox}{tbbpurple}{tbbpurplefill}
{\normalsize\bfseries\color{tbbpurple}실습 체크리스트}\par\vspace{3pt}
\begin{tbbitemize}
  \item 파트 A 결과물: 플릿의 \code{get pods -o wide} · 삭제와 조정 캡처 · 롤링 업데이트 내내 중단 없이 실행된 curl 루프 · 음성 대조군(준비 상태 검사 없음 → 오류).
  \item 파트 B 결과물: YAML과 일치하는 노드 쪽 \code{memory.max}와 \code{cpu.max} · 재시작 횟수를 포함한 OOMKilled describe 블록 · 복원한 설정.
  \item 파트 C 결과물: FailedScheduling 필터 보고서 · 터미널 기록 5-9 재현(이벤트 + logs --previous)과 수정 후 Ready 상태 · 비어 있는 엔드포인트라는 징후와 다시 채워진 모습.
  \item 보고서: 한 페이지, 관찰 세 가지, 각각 관련 절을 명시(예: "활성 상태 프로브 종료는 애플리케이션 로그를 남기지 않았다 — §5.6의 실패 원인으로서의 플랫폼").
  \item 정리: \code{kubectl delete namespace ml}; \code{minikube stop}. (이번 주에는 클라우드 비용이 없지만 습관은 유지한다.)
\end{tbbitemize}
\end{calloutbox}

\vspace{3.0pt}

\begin{calloutbox}{tbbpurple}{tbbpurplefill}
{\normalsize\bfseries\color{tbbpurple}생각해 보기}\par\vspace{3pt}
\begin{tbbitemize}
  \item 파트 A의 롤링 업데이트 동안 curl 루프는 한 번도 실패하지 않았다. 준비 상태 관문과 엔드포인트 갱신부터 제3장의 SIGTERM 핸들러까지 협력해야 했던 모든 메커니즘을 열거하고, 어느 하나를 제거했을 때 가장 손쉽게 보장이 깨지는지 찾자.
  \item minikube는 노드가 하나다. 실제 노드 다섯 개 클러스터라면 달라질 실습 동작 세 가지와, 단일 노드 클러스터에서는 제대로 재현할 수 없는 §5.6의 상태 하나를 말해 보자.
\end{tbbitemize}
\end{calloutbox}

\namedsection{요약}{열 가지로 정리한 이 장}

\qitem{1}{\textbf{Kubernetes는 조정 루프다.} 원하는 상태는 etcd에 있고, 컨트롤러는 영원히 관찰하고 차이를 찾고 동작한다. 명사를 선언하면 동사가 자연스럽게 나타난다. 삭제된 파드는 누군가 처리할 이벤트가 아니라 편차이며, 편차는 교정된다.}

\qitem{2}{\textbf{선언형 불변 조건은 시간이 지나도 효력을 잃지 않으므로 명령형 지시보다 강하다.} 노드 장애, 실수로 누른 삭제, 배포는 루프의 관점에서 모두 같다. 오토스케일링(12주 차)과 CI/CD(13주 차)는 원하는 상태를 작성하는 주체가 하나씩 더 늘어나는 것뿐이다.}

\qitem{3}{\textbf{구조는 자기주장이 있는 데이터베이스다.} API 서버(유일한 출입구), etcd(유일한 상태), 스케줄러(중매자), 컨트롤러 매니저(루프), kubelet(손발), kube-proxy(NAT)로 구성된다. 구성 요소끼리는 직접 호출하지 않고 모든 통신은 API 서버를 감시하는 방식로 이루어진다. 그래서 어느 한 구성 요소가 죽어도 시스템 전체가 중단되지 않는다.}

\qitem{4}{\textbf{Kubernetes라는 이름의 무언가가 컨테이너를 실행하는 것이 아니다.} 연결 고리는 apply → 객체 생성 → 스케줄러 바인딩 → kubelet → CRI → containerd → runc → 제3장의 네임스페이스와 cgroup이다. 마지막에 memory.max를 작성하는 대필자는 kubelet이다.}

\qitem{5}{\textbf{파드는 스케줄링과 공유의 단위다.} 컨테이너는 pause 컨테이너의 NET 네임스페이스(하나의 IP, localhost로 연결된 사이드카)에 합류하고 볼륨을 공유하며, MNT와 PID는 각자 분리된다. 파드는 언젠가 사라지고 레이블로 주소를 지정하므로 IP를 직접 대상으로 삼아서는 안 된다. 이 모든 것이 가정하는 워크로드 형태는 \textbf{마이크로서비스}다. 서비스마다 기능 하나를 두고 코드 복잡성을 운영 복잡성과 맞바꾼다. 세 축은 이 절충 관계의 입문 키트이며, 마이크로서비스 하나 = Deployment 하나 + Service 하나다.}

\qitem{6}{\textbf{프로브는 서로 다른 세 질문을 던진다.} 시작 프로브("부팅을 마쳤는가?")는 모델 로딩을 보호하고, 준비 상태 프로브("트래픽을 보내도 되는가?")는 Service 엔드포인트와 롤아웃 속도를 제어하며, 활성 상태 프로브("살아 있기는 한가?")는 재시작한다. 로딩에 90초가 걸리는 모델에서 활성 상태와 준비 상태를 혼동하면 정상 이미지가 영원히 충돌 루프를 돈다.}

\qitem{7}{\textbf{Deployment는 기계적인 방식으로 무중단 롤아웃을 만든다.} 템플릿 버전마다 ReplicaSet 하나를 두고, maxSurge/maxUnavailable로 속도를 조절하며, 준비 상태를 관문으로 사용하고, 롤백할 때는 이전 대상을 다시 가리킨다. Terminating 파드마다 제3장의 SIGTERM 과정이 되풀이된다. 처리 중인 요청을 비우는 일은 핸들러가 하고 Kubernetes는 그 작업을 스케줄링할 뿐이다.}

\qitem{8}{\textbf{requests는 스케줄러의 산술 값이고 limits는 cgroup의 벽이다.} requests.cpu는 유연하고 작업 보존적인 cpu.weight의 바탕도 된다. limits는 cpu.max(버스트 스로틀링 — §3.3에서 만든 p99)와 memory.max(§3.6의 OOM 벽)가 된다. requests와 limits 사이의 간격은 의도적인 오버커밋이다. 가정이 틀리면 kubelet이 퇴거를 수행하며, 이는 Evicted 대 OOMKilled/137이라는 서로 다른 징후를 보이는 두 번째 처형자다. QoS 클래스(Guaranteed / Burstable / BestEffort)는 oom\_score\_adj를 설정하며, 서빙 파드는 Guaranteed로 실행한다.}

\qitem{9}{\textbf{GPU는 정수 단위 확장 리소스로 스케줄링된다. 일부만 요청할 수 없고 오버커밋도 없으며 requests = limits다. GPU 공유를 중재할 cgroup 컨트롤러가 없기 때문이다(제3장에서 비어 있던 부분).} 테인트는 밀어내고 어피니티는 방향을 정하며 디바이스 플러그인은 리소스를 알리고 할당한다. GPU 플릿은 빈 패킹해야 한다. 단편화는 비용이 붙은 Pending이고 고립된 GPU는 유휴 상태로도 비용이 청구된다. 일부만 나누어 쓰는 방식은 제2장의 MIG/시간 분할에서 온다(12주 차).}

\qitem{10}{\textbf{여섯 가지 상태가 거의 모든 서빙 사고를 포괄한다.} Pending(필터 보고서 확인), ImagePullBackOff(제4장의 배관), CrashLoopBackOff(애플리케이션 충돌 또는 프로브 종료 — logs --previous), OOMKilled/137(제3장과 동일), Evicted(노드 압박), 실행 중이지만 라우팅되지 않음(빈 엔드포인트 — 조용한 실패)다. describe → 이벤트가 진단의 출입구이며 플랫폼 자체도 실패 원인 가운데 하나다.}

\namedsection{연습문제}{}

\subsection*{A. 객관식}

\qitem{1}{레플리카가 세 개인 Deployment에서 파드 하나를 삭제했다. 어떤 메커니즘으로 무슨 일이 일어나는가?}

\begin{tbboptions}
  \item ① 아무 일도 일어나지 않는다. 파드는 다음 배포 때만 다시 생성된다.
  \item ② restartPolicy가 Always이므로 kubelet이 같은 파드를 재시작한다.
  \item ③ ReplicaSet 컨트롤러가 원하는 수 3 ≠ 실제 수 2를 관찰하고 새 파드를 만든다(새 이름, 경우에 따라 새 노드).
  \item ④ kube-proxy가 Service의 엔드포인트 캐시에서 파드를 다시 만든다.
\end{tbboptions}

\qitem{2}{limits.memory: 1.75Gi를 노드의 memory.max 파일로 바꾸는 구성 요소는 무엇인가?}

\begin{tbboptions}
  \item ① 매니페스트를 검증할 때의 API 서버
  \item ② 바인딩 시점의 스케줄러
  \item ③ 저장소 훅을 거치는 etcd
  \item ④ CRI를 통해 컨테이너 런타임을 구동하는 노드의 kubelet
\end{tbboptions}

\qitem{3}{requests와 limits에 관한 설명으로 옳은 것은?}

\begin{tbboptions}
  \item ① requests는 cpu.max가 강제하고 limits는 권고값이다.
  \item ② requests는 스케줄러의 배치와 cpu.weight에 쓰이고, limits는 런타임에 강제되는 cgroup 벽이 된다.
  \item ③ 둘 다 똑같이 강제되며 requests는 단지 최솟값이다.
  \item ④ limits는 GPU에만 영향을 준다.
\end{tbboptions}

\qitem{4}{파드가 Running이고 liveness도 통과하지만 클라이언트가 Service를 통해 연결 오류를 겪는다. 가장 먼저 확인할 가치가 큰 항목은?}

\begin{tbboptions}
  \item ① kubectl logs — 스택 추적를 찾는다.
  \item ② Deployment를 재시작한다.
  \item ③ kubectl get 엔드포인트 — 목록이 비었다면 준비 상태 또는 셀렉터가 문제다.
  \item ④ 노드를 재부팅한다.
\end{tbboptions}

\qitem{5}{컨테이너가 메모리에서만 requests = limits로 설정했고 CPU request는 CPU limit보다 작다. 이 컨테이너의 QoS 클래스는?}

\begin{tbboptions}
  \item ① Guaranteed — 중요한 것은 메모리다.
  \item ② Burstable — Guaranteed가 되려면 모든 컨테이너의 CPU와 메모리에서 requests = limits여야 한다.
  \item ③ BestEffort — 클래스는 CPU가 결정한다.
  \item ④ 노드의 allocatable에 따라 달라진다.
\end{tbboptions}

\qitem{6}{1-GPU 파드가 계속 Pending에 머문다. describe에는 "노드 0/6개를 사용할 수 없음: 4개는 nvidia.com/gpu 부족, 2개는 허용되지 않은 테인트 \{nvidia.com/gpu: present\}가 있음."라고 나온다. 실제로 노드 두 개에는 여유 GPU가 있다. 올바른 해결책은?}

\begin{tbboptions}
  \item ① 파드의 CPU requests를 늘려 다른 파드보다 높은 값을 제시한다.
  \item ② 파드에 일치하는 톨러레이션을 추가해 테인트가 붙은 GPU 노드 두 개가 필터링을 통과하게 한다.
  \item ③ 더 쉽게 들어가도록 nvidia.com/gpu: 0.5로 설정한다.
  \item ④ 파드를 삭제하고 다시 만들어 스케줄링을 다시 실행한다.
\end{tbboptions}

\subsection*{B. 단답형}

\qitem{1}{레플리카가 3개인 새 Deployment에 kubectl apply를 실행한 순간부터 컨테이너가 실행될 때까지를 다섯 단계로 추적하고, 단계마다 담당 구성 요소를 밝히자.}

\qitem{2}{노드에서 Σ limits > allocatable은 허용되지만 Σ requests > allocatable은 불가능한 이유는 무엇인가? 오버커밋 가정이 틀렸을 때 대응하는 두 가지 안전장치도 말하자.}

\qitem{3}{모델 로딩에 최대 120초가 걸리는 서버에 적합한 프로브 설정(프로브 유형, 주기, failureThreshold)을 작성하고, 각 숫자의 근거를 한 구절씩 제시하자.}

\qitem{4}{OOMKilled와 Evicted를 비교해 각각 발동 조건, 처형자, 한 줄짜리 징후를 제시하자.}

\subsection*{C. 논술형}

\qitem{1}{"지연 시간에 민감한 서빙 파드에는 CPU limits를 절대 설정하지 말라." §3.3의 스로틀링 계산, cpu.weight의 작업 보존적 동작, Guaranteed-QoS 요구사항, 플랫폼 거버넌스의 현실을 근거로 찬반을 논하자. 강의 프로젝트의 모델 서버에 적용할 권고안을 밝히고 옹호하자.}

\qitem{2}{하나의 클러스터에서 (a) T4에서 실행되는 지연 시간에 민감한 양자화 모델, (b) A100을 원하는 야간 배치 스코어링 작업, (c) 일반 CPU 마이크로서비스를 서빙하기 위한 노드 풀, 레이블, 테인트, QoS 선택, 패킹 정책을 설계하자. 각 선택이 방지하는 실패 또는 비용 시나리오로 근거를 제시하고, 12주 차 오토스케일러가 설계에 무엇을 추가하는지 밝히자.}

\qitem{3}{"자가 복구는 근본 원인을 숨긴다. 충돌하는 파드를 조용히 교체하는 조정 루프는 사고 해결 시스템이 아니라 사고 억제 시스템이다." 모델 로딩의 재시작 루프 경제학(§3.6, §5.6), §5.6의 관측 표면, 수렴과 진단을 양립시키기 위해 성숙한 팀이 추가해야 할 것(14주 차)을 바탕으로 논하자.}

\namedsection{정답과 해설}{}

\subsection*{A. 객관식}

\qitem{1}{\textbf{③} 되살리는 것이 아니라 조정한다. ReplicaSet 조정 루프가 2 ≠ 3을 세고 새 파드(새 접미사, 새 IP, 경우에 따라 새 노드)를 만든다. ②는 파드 내부의 컨테이너 재시작과 파드 교체를 혼동했다. (§5.1)}

\qitem{2}{\textbf{④} kubelet은 손발이자 대필자다. CRI를 통해 containerd/runc를 구동하고, 이들이 cgroup 파일을 작성한다. API 서버는 저장하고 스케줄러는 배치할 뿐, 어느 쪽도 노드를 건드리지 않는다. (§5.2, §5.4)}

\qitem{3}{\textbf{②} 이 장의 핵심 비대칭이다. requests는 계획용 숫자로 런타임에 강제되지 않으며 CPU 요청량은 cpu.weight의 바탕도 된다. limits는 매 밀리초 강제되는 제3장의 cgroup 벽이다. (§5.4)}

\qitem{4}{\textbf{③} 실행 중이지만 라우팅되지 않는 상태는 조용한 실패다. 활성 상태 검사는 초록색이지만 준비 상태 검사가 실패하거나 레이블이 셀렉터와 맞지 않아 → 엔드포인트가 비고 → Service에는 보낼 대상이 없다. 잘못을 저지른 것이 없으므로 로그(①)에는 아무것도 나타나지 않는다. (§5.3, §5.6)}

\qitem{5}{\textbf{②} Guaranteed는 엄격하다. \textit{모든} 컨테이너의 \textit{두} 리소스에서 requests = limits여야 한다. 일부만 같으면 Burstable이다. 메모리만 Guaranteed처럼 설정하는 것은 흔하고 합리적인 서빙 절충안이지만 클래스 이름은 아니다. (§5.4)}

\qitem{6}{\textbf{②} 필터 보고서를 읽으면 된다. 노드 네 개에는 여유 GPU가 없고, 두 개에는 있지만 테인트가 붙었으며 파드에는 톨러레이션이 없다. 톨러레이션이 테인트가 붙은 두 노드의 문을 연다. ③은 정수만 허용되므로 불법이고, ④는 같은 필터를 다시 실행해 같은 판결을 낸다. (§5.5, §5.6)}

\subsection*{B. 단답형}

\qitem{1}{(1) API 서버가 Deployment를 검증하고 etcd에 저장한다. (2) Deployment 컨트롤러가 ReplicaSet을 만들고, ReplicaSet 컨트롤러가 nodeName이 비어 있는 Pod 객체 3개를 만든다. (3) 스케줄러가 각 파드를 필터링/점수화/바인딩해 노드에 배치한다. (4) 각 노드의 kubelet이 이미지를 가져오고 containerd → runc를 구동해 네임스페이스, cgroup, overlayfs를 만든다. (5) kubelet이 상태를 보고하고 프로브가 통과하면 엔드포인트가 파드를 받아들인다. (§5.2)}

\qitem{2}{스케줄러는 requests만 계산하므로 배치 과정에서 이를 초과 예약할 수 없다. limits는 합계가 노드 용량을 넘도록 의도적으로 허용하는 버스트 상한이다. 모두가 버스트할 수 있지만 동시에 그러는 일은 드물다는, 기능으로서의 오버커밋이다. 가정이 틀릴 때를 위한 안전장치는 컨테이너별 cgroup 벽(memory.max → OOMKilled/137)과 kubelet의 노드 압박 퇴거(QoS 순서, 상태 Evicted)다. (§5.4)}

\qitem{3}{startupProbe: httpGet /healthz, periodSeconds 5, failureThreshold ≥ 30이다. 5 × 30 = 150 s > 120 s로 최악의 로딩 시간에 여유를 더해 다른 프로브를 보류한다. readinessProbe: /ready, 주기 5, 임계값 2–3으로 트래픽을 빠르게 제어한다. livenessProbe: /healthz, 주기 10–15, 임계값 3으로 의도적으로 느긋하게 잡아 실제 교착 상태만 재시작한다. (§5.3, §5.6)}

\qitem{4}{OOMKilled: 컨테이너가 자기 memory.max를 넘었고 커널 OOM 킬러가 SIGKILL을 보낸다. 징후는 Last State: OOMKilled, exit 137(+ dmesg)이다. Evicted: 오버커밋 상태에서 노드 자원이 부족해져 kubelet이 QoS/requests 초과 사용량에 따라 희생 대상을 종료한다. 징후는 상태 Evicted, "노드의 메모리 리소스가 부족함."다. (§5.4, §5.6)}

\subsection*{C. 논술형(채점 기준)}

\qitem{1}{"limit를 설정하지 말라"는 쪽은 §3.3의 계산(30 ms 작업 → 0.2-CPU 상한에서 110–190 ms, p99 급증, 프로파일러에는 보이지 않음), cpu.weight가 자원을 낭비하지 않으면서 이미 경합의 공정성을 보장한다는 점, 유휴 여유분은 저렴한 보험이라는 점을 들어야 한다. 반대쪽은 플랫폼 규모의 폭주 방지와 비용 관리, 정직하지 않은 requests를 설정한 파드로부터 멀티테넌트 공정성을 지키는 문제, Guaranteed QoS가 limits를 요구하고 오버커밋 노드에서 Guaranteed의 OOM 보호가 중요하다는 점, 정적 CPU 매니저 고정에 정수 requests = limits가 필요하다는 점을 들어야 한다. 좋은 답은 이 긴장을 다룬 뒤 입장을 정한다. 예를 들어 측정한 최대 CPU보다 넉넉하게 requests = limits(Guaranteed)를 설정하거나, CPU limit 없이 메모리만 Guaranteed인 Burstable과 PriorityClass를 사용하고 각각의 절충 관계를 밝힌다. (§5.4)}

\qitem{2}{다음 요소를 확인한다. gpu=t4/a100 레이블이 붙은 세 노드 풀(t4-serving, a100-batch, cpu-general), 일반 CPU 작업이 들어오지 않도록 두 GPU 노드 풀에 적용한 테인트(§5.5), 톨러레이션 + gpu=t4 어피니티와 Guaranteed QoS를 적용한 모델 A 파드, 낮은 PriorityClass로 a100 노드 풀을 허용하고 CPU limits는 두지 않은 배치 작업 B, 축소와 다중 GPU 여유 공간을 위한 GPU 노드 풀의 빈 패킹(MostAllocated) 점수, 유휴 고립을 피하도록 GPU 파드에 설정한 정직한 CPU requests다. 비용 시나리오는 단편화로 인한 강제 노드 추가, CPU를 많이 쓰는 DaemonSet 뒤에 고립된 GPU, 배치 작업이 지연 시간에 민감한 파드를 선점하는 문제(우선순위 + 테인트로 방지)다. 12주 차에는 오토스케일러가 빈 패킹으로 비워진 노드를 삭제하고 배치 노드 풀을 0까지 축소한다. (§5.4–5.5)}

\qitem{3}{좋은 답은 먼저 전제의 근거를 인정한다. 90 s마다 자가 복구하는 충돌 루프의 모델 서버는 매번 콜드 스타트 지연 시간을 치르면서도 너무 작은 limit를 몇 주 동안 숨길 수 있고(§3.6의 관성), 터미널 기록 5-9의 사고는 애플리케이션 로그 증거를 전혀 남기지 않는다. 그런 다음 반론한다. 수렴이 새벽 3시의 장애를 견디게 하며, 루프는 증거를 숨기기보다 재시작 횟수, 이벤트, Last State, memory.events에 \textit{국소화}한다. 애플리케이션이 아니라 플랫폼의 표면을 보면 된다. 성숙한 운영(14주 차)은 재시작 횟수, oom\_kill 카운터, 비어 있는 엔드포인트, 멈춘 롤아웃에 경보를 건다. 즉 가용성뿐 아니라 편차 교정의 \textit{빈도}에 경보를 건다. 관측 가능성 없는 조정은 사고 억제지만, 관측 가능성과 결합하면 진단 시간을 벌어 주는 사고 \textit{완충}이라는 결론을 높이 평가한다. (§5.1, §5.6)}

\namedsection{더 읽을거리}{}

\textbf{[이번 주]} 표시가 붙은 항목은 이번 주 필독 자료다. \textbf{[N주 차]} 표시는 해당 주에 도움이 될 자료를 뜻한다.

\begin{tbbitemize}
  \item \textbf{[이번 주] Kubernetes 문서: Concepts → "Pods", "Deployments", "Service"} — §5.3의 바탕이 되는 공식 설명이다. 작업 페이지가 아니라 Concepts 페이지를 읽자. 개념 문서는 튜토리얼이 생략하는 \textit{이유}를 설명한다.
  \item \textbf{[이번 주] Kubernetes 문서: "Pod Lifecycle"} (프로브 포함) — §5.3의 세 질문, 컨테이너 상태, 재시작 백오프에 관한 권위 있는 자료다. §5.6의 CrashLoopBackOff 메커니즘도 모두 여기에 있다.
  \item \textbf{[이번 주] Kubernetes 문서: "Resource Management for Pods and Containers" + "Configure Quality of Service for Pods"} — §5.4의 1차 자료다. requests와 limits, allocatable, QoS 클래스, 그리고 노드 문서에 나오는 노드 압박 퇴거 순서를 다룬다.
  \item \textbf{Kelsey Hightower, "Kubernetes the Hard Way" (GitHub)} — 설치 도구 없이 클러스터를 한 번 만들어 보면 §5.2의 구조가 더 이상 그림에 머물지 않는다. 비 오는 주말에 평생 한 번만 해 보자. 그 정신 모형은 경력 내내 남는다.
  \item \textbf{Marko Lukša, "Kubernetes in Action" (2nd ed., Manning)} — 책 한 권 분량으로 다룬 자료 가운데 가장 좋다. Pod, Deployment, Service, 리소스 관리 장은 더 넓은 API 표면을 다루면서도 §5.3–5.4와 일대일로 대응한다.
  \item \textbf{Verma et al., "Large-scale cluster management at Google with Borg" (EuroSys 2015)} — 선언형 작업, 조정, 빈 패킹, §5.5가 물려받은 사용률 경제학의 기원을 보여 준다. §5(utilization)를 제1장의 비용 관점으로 읽자.
  \item \textbf{[6주 차] Kubernetes 문서: "Ingress" 및 클라우드 공급업체 로드 밸런서 문서} — 6주 차에 Service 앞에 추가할 계층(API 게이트웨이, TLS, 라우팅)을 다룬다.
  \item \textbf{[12주 차] NVIDIA GPU Operator 및 디바이스 플러그인 문서(시간 분할, MIG 전략); Kubernetes Cluster Autoscaler FAQ; HPA/KEDA 문서} — §5.5의 운영상 후속 내용인 분할 GPU, Pending에 따른 확장, 빈 패킹으로 비운 노드의 축소, 이벤트 기반 오토스케일링을 다룬다.
  \item \textbf{[13주 차] Argo Rollouts 문서(카나리, 블루-그린)} — 13주 차에 점진적 전달(progressive delivery)과 자동 분석을 추가하면 Deployment의 롤링 업데이트가 어떤 형태로 발전하는지 보여 준다.
\end{tbbitemize}

\vspace{5.0pt}

\begin{calloutbox}{tbbblue}{tbbbluefill}
{\normalsize\bfseries\color{tbbblue}다음 장 — 제6장: 스토리지와 네트워킹}\par\vspace{3pt}
이제 모델 서버는 안정적인 이름을 가진 플릿으로 실행된다. 하지만 가중치는 여전히 이미지에 구워 넣은 상태이고(제3장에서 경고했다), Service도 여전히 클러스터 내부 배관에 머문다. 다음 주에는 서빙의 골격을 완성한다. \textbf{블록 스토리지(block storage)와 객체 스토리지}의 차이와 모델 가중치를 객체 스토어(object store)에 보관하는 이유, S3 호환 스토리지에 가중치와 데이터셋을 두고 시작할 때 로딩하며 이미지와 독립적으로 버전을 관리하는 \textbf{모델 저장소 패턴}(6주 차 실습에서 바로 이 경로를 만들고 13주 차 레지스트리 이야기의 바탕이 된다), Service 앞의 출입구인 \textbf{로드 밸런서와 API 게이트웨이}, 값비싼 GPU 플릿을 공개 인터넷에서 격리하는 데 필요한 \textbf{VPC와 보안 그룹(security group)} 기본 지식을 다룬다. 6주 차가 끝나면 강의의 인프라 절반이 완성되고, 파트 2에서는 그 위에서 모델 서빙을 시작한다.
\end{calloutbox}

%% Chapter 6 — 스토리지와 네트워킹
\setcounter{chapter}{5}
\chapter{스토리지와 네트워킹}
\chaptersubtitle{가중치가 사는 곳과 플릿의 출입구 — 객체 스토어, 로드 밸런서, 의도적으로 설계한 경계}

\begin{calloutbox}{tbbblue}{tbbbluefill}
{\normalsize\bfseries\color{tbbblue}학습 목표}\par\vspace{3pt}
\begin{tbbitemize}
  \item 파드의 모든 바이트는 이미지에서 왔거나, 저장소에서 가져왔거나, 버려도 된다는 상태 없는 서버 원칙을 설명하고 제5장의 가축 메커니즘을 근거로 정당화하며, 각 서빙 상태를 알맞은 위치에 둔다.
  \item 블록, 객체, 파일 스토리지를 의미 체계, 지연 시간, 처리량, 공유, 내구성, 가격이라는 \textit{계약}으로 비교하고, 유행이 아니라 계산으로 워크로드마다 알맞은 방식을 고른다.
  \item 변경 불가능한 버전별 배치, 무결성 매니페스트, 병렬 범위 다운로드, 노드 캐싱을 갖춘 객체 스토리지 기반 모델 스토어를 설계하고, 콜드 스타트 비용에서 차지하는 몫을 계산한다.
  \item 이미지를 다시 빌드하지 않고 환경 변수 하나만 바꿔 새 모델 버전을 배포하며, 코드와 모델의 출시 주기를 분리하는 일이 13주 차 자동화에서 중요한 이유를 설명한다.
  \item Service에 실제 출입구를 둔다. L4와 L7, Ingress와 TLS, API 게이트웨이의 제품 기능, 꼬리가 긴 추론과 스트리밍을 견디는 타임아웃·재시도 설정을 이해한다.
  \item 공개 출입구, 사설 하드웨어, 가중치로 가는 저렴한 경로, 키 없는 IAM을 갖춘 최소 서빙 VPC를 그리고, 상태 코드만으로 출입구 장애(502/503/504/429, 정확히 N초에서 끊기는 스트림)를 진단한다.
\end{tbbitemize}
\end{calloutbox}

\vspace{5.0pt}

\begin{calloutbox}{tbbteal}{tbbtealfill}
{\normalsize\bfseries\color{tbbteal}핵심 용어}\par\vspace{3pt}
상태 없는 서비스 · 상태 외부화 · 블록 스토리지(EBS) · IOPS · 객체 스토리지(S3) · 버킷 / 키 / 접두사 · 범위 GET · 멀티파트 · 강한 일관성 · 내구성(일레븐 나인) · 스토리지 클래스 / 수명 주기 · 외부 전송 · 파일 스토리지(NFS / EFS) · PersistentVolume / RWO / RWX · 모델 스토어 · 변경 불가능한 버전 관리 · 매니페스트 · safetensors · mmap · initContainer · emptyDir · 콜드 스타트 예산 · 워밍업 · L4 / L7 · 로드 밸런서 · TLS 종료 · Ingress · API 게이트웨이 · 속도 제한(429) · 유휴 타임아웃 · 연결 비우기 · VPC · 서브넷(공개 / 사설) · 라우팅 테이블 · 보안 그룹 · NAT 게이트웨이 · VPC 게이트웨이 엔드포인트 · IAM 역할 · IRSA / 워크로드 아이덴티티 · 502 / 503 / 504
\end{calloutbox}

\vspace{5.0pt}

\begin{calloutbox}{tbblightgray}{tbbboxfill}
{\normalsize\bfseries\color{tbbink}장 구성}\par\vspace{3pt}
\textbf{6.1} 모델 서버는 아무것도 소유하지 않는다    \textbf{6.2} 두 가지 계약: 블록과 객체    \textbf{6.3} 모델 스토어    \textbf{6.4} 출입구: 로드 밸런서, Ingress, 게이트웨이    \textbf{6.5} 경계: VPC, 보안 그룹, 저렴한 경로    \textbf{6.6} 출입구 읽기: 상태 코드 현장 안내서    \textbf{6.7} 실습: S3에서 로딩하는 서빙 구조 완성 — 이어서 요약, 연습문제, 더 읽을거리.
\end{calloutbox}

\vspace{3.0pt}

\begin{calloutbox}{tbbamber}{tbbamberfill}
{\normalsize\bfseries\color{tbbamberdark}이 장을 읽는 방법}\par\vspace{3pt}
이번 장의 터미널 기록도 노트북에서 재현할 수 있다. 실습에서는 minikube 클러스터 안에서 S3 호환 저장소인 \textbf{MinIO}를 실행하므로 스토리지 부분에는 클라우드 계정이나 비용이 필요 없다. 로드 밸런서, VPC, IAM처럼 클라우드가 필요한 예시는 AWS에서 대표 가격으로 실행했음을 표시했다. 개념은 물론 대부분의 플래그 이름까지 GCS/Azure의 대응 요소와 일대일로 연결된다. 가격이 나오면 \textbf{절대값이 아니라 형태}로 받아들이자. 가격은 바뀌어도 비율은 오래 남는다.
\end{calloutbox}

\section{모델 서버는 아무것도 소유하지 않는다}

제5장은 조용한 믿음 하나에 기대어 끝났다. 파드를 삭제하고도 대수롭지 않게 여겼고, 플릿을 롤아웃하며 대낮에 기존 레플리카가 죽는 모습을 지켜봤으며, 스케줄러가 여유 있는 노드로 워크로드를 옮기게 두었다. 이 모든 동작에는 한 번도 명시하지 않은 가정이 있다. \textbf{가치 있는 것은 아무것도 파드와 함께 죽지 않는다.} 파드가 모델 파일, 사용자의 업로드, 세션 같은 무언가의 유일한 복사본을 가지고 있다면 제5장의 모든 동작은 데이터 손실 엔진이 된다. 따라서 이 책이 플릿 위에 무엇을 더 쌓기 전에 이 장에서 밀린 숙제를 해결한다. 가치 있는 것은 실제로 어디에 살며, 바이트와 패킷은 그곳과 파드 사이를 어떻게 오가는가?

먼저 서빙 파드 안에 무엇이 있는지 점검해 보자. 아래 터미널 기록은 13 GB 모델 서버의 실행 중인 레플리카에 들어가 파일 시스템 목록을 조사한다. 보이는 모든 마운트가 제3장의 메커니즘이므로 제3장의 관점으로 읽자.

\begin{terminalbox}
\begin{Verbatim}[breaklines=true,breakanywhere=true,fontsize=\small,formatcom=\color{tbbtermfg},codes=\tbbcodeactive]
$ kubectl exec -it model-server-6f8d5b9c44-2mwpx -- sh
/app $ df -h
Filesystem      Size  Used  Avail  Mounted on
overlay          98G  6.2G    87G  /            # the image (Ch. 3)
tmpfs            64M     0    64M  /dev/shm     # Ch. 3's --shm-size
/dev/nvme1n1    200G   13G   187G  /models      # emptyDir on node NVMe
/app $ ls -lh /models
-rw-r--r--  4.9G  model-00001.safetensors
-rw-r--r--  4.9G  model-00002.safetensors
-rw-r--r--  3.4G  model-00003.safetensors
-rw-r--r--   612  manifest.json
/app $ env | grep MODEL
MODEL_URI=s3://ml-models/llama-8b-int8/2026-07-01/
# 13 GB of weights - which were NOT in the image. they arrived at
# boot, from the URI above. kill this pod: every byte reappears.
\end{Verbatim}
\end{terminalbox}
\figcaption{터미널 기록 6-1  점검 결과. 이미지 기반 루트(대부분 읽기 전용이며 버려도 됨), 임시 마운트, 부팅 때 다른 곳에서 가져온 13 GB 가중치가 있다. 여기에는 유일한 것이 아무것도 없다.}

세 종류의 바이트와 세 가지 출처가 원칙의 전부다. 정상적인 서빙 파드 안의 모든 것은 \textbf{(a) 이미지에서 왔거나}(코드와 의존성 — 레지스트리에서 다시 만들 수 있는 제4장의 산출물), \textbf{(b) 부팅할 때 저장소에서 가져왔거나}(이 장에서 다룰 가중치), \textbf{(c) 버려도 되는 임시 데이터}(캐시, 임시 파일, \code{/\allowbreak dev/\allowbreak shm})다. 네 번째 종류는 없다. 다른 곳에는 존재하지 않는 바이트 하나라도 파드가 가지는 순간 가축이기를 멈춘다. 그러면 재스케줄링, 롤아웃, 오토스케일링, 축출 같은 제5장의 모든 메커니즘이 데이터를 잃는 방법으로 변한다(그림 6-1).

\begin{tbbfigure}
  \adjustbox{max width=0.94\linewidth,max totalheight=0.46\textheight}{\includegraphics{figures/fig6-1_ownnothing.png}}\\[3pt]
  \figcaption{그림 6-1  아무것도 소유하지 않는 원칙. 코드는 레지스트리에서, 가중치는 모델 스토어에서 오고, 트래픽은 출입구로 들어오며, 로그는 밖으로 내보낸다. 파드를 죽여도 모든 바이트가 다시 나타난다. 이 원칙이 제5장의 모든 동작을 가능하게 했다.}
\end{tbbfigure}

\subsection*{상태 목록}

"상태 없음"은 서빙 시스템에 상태가 없다는 뜻이 아니다. 오히려 상태가 넘쳐난다. \textit{서버 프로세스}가 어떤 상태도 유일하게 보유하지 않으며, 각 상태를 형태에 맞는 계약을 제공하는 저장소로 밀어낸다는 뜻이다. 한 번은 빠짐없이 정리할 가치가 있다. 이 책에서 앞으로 나올 모든 아키텍처 논의는 사실 아래 표의 한 행에 해당하기 때문이다.

\begin{tbbtablewrap}
  \begin{tabular}{@{}>{\raggedright\arraybackslash}m{\dimexpr 0.2660\tbltextwidth-2\tabcolsep\relax} >{\raggedright\arraybackslash}m{\dimexpr 0.3723\tbltextwidth-2\tabcolsep\relax} >{\raggedright\arraybackslash}m{\dimexpr 0.3617\tbltextwidth-2\tabcolsep\relax}@{}}
  \arrayrulecolor{tbbnavy}\specialrule{1pt}{0pt}{0pt}
  \rowcolor{tbbtablehead}{\centering \bfseries\color{tbbnavy}상태} & {\centering \bfseries\color{tbbnavy}저장 위치} & {\centering \bfseries\color{tbbnavy}선택 이유} \\
  \arrayrulecolor{tbbnavy}\specialrule{0.7pt}{0pt}{2pt}
  {코드 + 의존성} & {레지스트리의 이미지(제4장)} & {버전 관리, 내용 주소 지정, 노드별 캐시} \\
  \arrayrulecolor{tbbtableborder}\hline
  {\textbf{모델 가중치}} & {\textbf{모델 스토어(model store) — 객체 스토리지(§6.3)}} & {매우 크고 변경 불가능하며 버전이 있고 모든 레플리카가 읽음} \\
  \arrayrulecolor{tbbtableborder}\hline
  {요청 상태} & {별도 저장 안 함 — 요청 자체에 존재} & {상태 없는 HTTP: 어느 레플리카든 어느 요청이든 응답 가능} \\
  \arrayrulecolor{tbbtableborder}\hline
  {캐시(토크나이저, JIT, 워밍업)} & {로컬 임시 공간(emptyDir) — 버려도 됨} & {다시 만드는 비용이 저렴하며 파드의 수명과 묶을 가치가 없음} \\
  \arrayrulecolor{tbbtableborder}\hline
  {로그 · 지표 · 트레이스} & {지속해서 외부로 전송(14주 차)} & {읽기 전에 파드가 죽을 수 있음} \\
  \arrayrulecolor{tbbtableborder}\hline
  {사용자 데이터 · 특성 · 세션} & {데이터베이스와 특성 스토어} & {트랜잭션을 지원하고 변경 가능 — 블록 기반 엔진(§6.2)} \\
  \arrayrulecolor{tbbnavy}\specialrule{1pt}{2pt}{0pt}
  \end{tabular}\par
  \figcaption{표 6-1  서빙 상태별로 알맞은 위치}
\end{tbbtablewrap}

두 행은 다시 볼 가치가 있다. \textbf{요청 상태}가 없기 때문에 제5장의 Service는 "연결마다 정상 백엔드 하나"를 가볍게 고를 수 있었다. 어느 레플리카든 어느 요청에나 응답할 수 있다면 로드 밸런싱은 자유로운 선택이다. 세션 어피니티("고정 세션")가 끌리기 시작했다면 대개 상태를 서버 안으로 몰래 들여온 것이다. 그 대가로 부하가 고르지 않고, 롤아웃마다 처리 중 요청을 비우기 어려우며, 장애 대응도 나빠진다. 그래서 서빙 시스템은 어피니티를 기능이라기보다 나쁜 냄새로 본다. \textbf{캐시}는 더 미묘하다. 워밍업된 파드는 콜드 파드와 실제로 다르다. 제3장의 페이지 캐시에는 가중치가 차 있고, 커널은 JIT로 컴파일되어 있으며, 할당기도 준비되어 있다. 상태 없음은 \textit{성능}이 아니라 \textit{정확성}에 관한 원칙이다. 캐시를 잃어도 정답을 잃어서는 안 되지만, 지연 시간은 분명 나빠질 수 있다. 바로 그래서 이 장은 콜드 스타트 비용을 항목별로 계산하고(§6.3), 12주 차에는 이를 줄이는 방법을 다룬다.

\subsection*{이 장이 제 몫을 하는 이유}

\begin{tbbitemize}
  \item \textbf{가중치 문제가 곧 서빙 문제다.} 440 MB BERT는 실수에 너그럽지만 13 GB 양자화 LLM은 그렇지 않으며, 11주 차 모델은 더 크다. 가중치가 어디에 살고 얼마나 빨리 도착하며 어떻게 버전을 관리하는지가 콜드 스타트 비용(12주 차), 배포 방식(13주 차), 그리고 청구서에서 놀랄 만큼 큰 몫(§6.5의 NAT 함정)을 결정한다.
  \item \textbf{출입구가 곧 제품이다.} 제5장의 Service에는 클러스터 내부 이름만 있다. ClusterIP로는 누구에게도 비용을 청구할 수 없다. TLS, 라우팅, 할당량, 사용량 계측(§6.4)이 Deployment를 제1장 생태계의 모델 API로 바꾼다. 강의 프로젝트도 바로 이 노출을 올바르게 구현해야 한다.
  \item \textbf{경계는 선택 사항이 아니다.} GPU 플릿은 대다수 팀이 운영하는 시스템 가운데 가장 비싸고 공격하기도 쉽다. §6.5의 VPC 패턴인 공개 출입구 하나, 사설 하드웨어, 키 없는 아이덴티티가 제대로 된 아키텍처와 사고 보고서를 가른다.
  \item \textbf{이 장에서 인프라 절반을 완성한다.} 이 장이 끝나면 커널(제2–3장), 패키징(제4장), 오케스트레이션(제5장), 상태와 경계(제6장)를 갖춘다. 이 책의 파트 2에서는 이 스택 위에서 모델을 서빙하며 아래를 다시 내려다보지 않는다. 다만 디버깅할 때는 예외이므로 §6.6이 필요하다.
\end{tbbitemize}

\begin{calloutbox}{tbbpurple}{tbbpurplefill}
{\normalsize\bfseries\color{tbbpurple}생각해 보기}\par\vspace{3pt}
\begin{tbbitemize}
  \item 대부분의 "상태 없는" 서버에 숨어 있는 상태인 처리 중 요청을 찾아보자. 종료할 때 이 상태를 비우기 위해 제3장과 제5장에서 마련한 두 메커니즘은 무엇이며, 생략하면 무엇이 깨지는가?
  \item 팀원이 사용자마다 요청이 같은 레플리카에 가도록 고정 세션(sticky session)을 제안한다("KV 캐시가 워밍업 상태를 유지할 것이다"). 롤아웃, 오토스케일 축소, 노드 장애 때 치러야 할 비용을 열거한 뒤, 11주 차 LLM 세션에서 이 긴장이 실제로 다시 나타나는 지점을 살펴보자.
  \item 대신 가중치를 이미지에 구워 넣었다고 하자(제3장 §3.4의 안티패턴). 표 6-1의 이미지 행을 따라가며 조용히 깨지는 성질을 나열하자. 이제 "버전 관리"는 무엇을 뜻하고, 코드 긴급 수정의 비용은 얼마이며, 13 GB × 매주 배포할 때 레지스트리 청구서는 어떻게 되는가?
\end{tbbitemize}
\end{calloutbox}

\section{두 가지 계약: 블록과 객체}

클라우드 스토리지는 세 계열로 나뉜다. 이들을 생산적으로 비교하는 방법은 기능 목록이 아니라 서비스가 무엇을 약속하고 그 약속에 얼마를 받는지 나타내는 \textbf{계약}을 보는 것이다. 계약을 한 번 정확히 이해하면 이 책의 모든 스토리지 결정은 패턴 맞추기로 단순해진다. 접근 패턴은 무엇이고, 누가 공유하며, 어느 정도의 내구성(durability)이 필요하고, 정확히 그 조건을 파는 계약은 무엇인가(그림 6-2)?

\subsection*{블록: 연결해서 쓰는 디스크}

\textbf{블록 스토리지}(AWS EBS, GCP Persistent Disk)는 컴퓨팅에서 가장 오래된 약속인 \textit{디스크}를 판매한다. 원하는 크기의 원시 장치를 받으며 OS는 이를 로컬 하드웨어와 구별하지 못한다. 사실 제2장의 virtio 이야기가 다시 등장한 것이다. 블록은 네트워크 너머의 복제 스토리지 서버에 있지만 반가상화 디스크로 보인다. 여기에 파일 시스템을 올리면 임의 읽기·쓰기, 제자리 갱신, 바이트 단위 탐색, 밀리초 미만에서 밀리초 단위의 지연 시간까지 완전한 POSIX 의미 체계를 얻는다. 데이터베이스에 필요한 계약이 바로 이것이다. B-tree는 4,381번 페이지를 초당 만 번 제자리에서 다시 쓰고 싶어 한다. 가격도 그에 맞게 책정된다. 용량과 성능(GB 단위 크기, IOPS, 처리량)을 \textbf{프로비저닝(provisioning)}하고 실제 사용 여부와 관계없이 예약한 만큼 비용을 낸다(gp3급: 대표값 월 GB당 약 \$0.08). 운영상 중요한 작은 조항이 두 가지 있다. 블록 볼륨은 \textbf{한 번에 인스턴스 하나에만} 연결되며(공유 공간이 아니라 디스크다), \textbf{하나의 가용 영역}에 존재한다. 그 이상의 내구성은 스냅샷으로 직접 책임져야 하며, 흥미롭게도 스냅샷은 객체 스토리지에 저장된다.

\begin{terminalbox}
\begin{Verbatim}[breaklines=true,breakanywhere=true,fontsize=\small,formatcom=\color{tbbtermfg},codes=\tbbcodeactive]
$ lsblk
NAME     SIZE  TYPE  MOUNTPOINT
nvme0n1  100G  disk  /              # boot volume (also block!)
nvme1n1  200G  disk                 # the EBS volume we attached
$ sudo mkfs.ext4 /dev/nvme1n1 && sudo mount /dev/nvme1n1 /scratch
$ dd if=/dev/zero of=/scratch/t bs=4k count=1 oflag=dsync
4096 bytes copied, 0.00089 s      # ~0.9 ms: a disk-shaped latency
# the OS cannot tell this from local hardware. that is the product.
\end{Verbatim}
\end{terminalbox}
\figcaption{터미널 기록 6-2  직접 체감하는 블록 계약. 연결하고, 포맷하고, 마운트한 뒤 밀리초 미만으로 쓴다. 모든 줄 아래에서 제2장의 virtio 메커니즘이 작동한다.}

\subsection*{객체: 호출해서 쓰는 API}

\textbf{객체 스토리지}(S3, GCS)는 전혀 다른 약속, 즉 \textit{HTTP API 뒤에 놓인 내구성 높은 블롭}을 판매한다. 장치도, 파일 시스템도, 탐색도 없다. \code{bert-classifier/\allowbreak v7/\allowbreak model.safetensors} 같은 키 아래에 최대 5 TB인 객체를 \code{PUT}하고, 전체 또는 원하는 바이트 범위를 \code{GET}한다. 제자리 편집은 할 수 없다. "객체의 100–200번 바이트를 덮어쓴다"는 동작이 없으므로 객체 전체를 교체하거나 아예 건드리지 않아야 한다. 디스크 의미 체계를 포기하는 대신 디스크가 제공할 수 없는 성질을 얻는다. \textbf{내구성}은 일레븐 나인(99.999999999\%)이다. 모든 객체를 여러 장치와 가용 영역에 삭제 코딩해 달성한다. 객체 하나가 사라져도 뉴스가 될 정도이며 누구도 S3를 스냅샷하지 않는다. \textbf{규모와 공유} 측면에서는 용량이 사실상 무제한이고 리전 안의 \textit{모든} 읽기 주체가 동시에 가져갈 수 있다. 파드 천 개가 같은 가중치를 끌어오는 일도 평범한 화요일일 뿐이다. \textbf{가격}은 실제 저장량(표준 클래스 월 GB당 약 \$0.023), 요청 수(GET 1,000회당 약 \$0.0004), 그리고 모두가 잊는 조항인 \textbf{외부로 나가는 바이트 수}(인터넷 외부 전송 GB당 약 \$0.09, §6.5에서 NAT 변형을 더한다)에 따라 낸다. 지연 시간에 관한 작은 조항도 있다. \textbf{첫 바이트 도착 시간은 수십 밀리초}로 디스크보다 100배 느리지만, 이후 각 스트림은 50–100 MB/s를 내고 스트림은 병렬로 늘릴 수 있다. 많은 객체의 여러 범위를 병렬로 요청하면 객체 스토리지가 10 Gbps NIC를 기꺼이 포화시킨다. 시작은 느리지만 합계에는 사실상 상한이 없는 형태를 기억하자. §6.3 전체가 이를 활용한다.

\begin{terminalbox}
\begin{Verbatim}[breaklines=true,breakanywhere=true,fontsize=\small,formatcom=\color{tbbtermfg},codes=\tbbcodeactive]
$ aws s3 cp model.safetensors s3://ml-models/bert-classifier/v7/
upload: ./model.safetensors to s3://ml-models/.../model.safetensors
$ curl -sI "$(aws s3 presign s3://ml-models/.../model.safetensors)" \
    -H "Range: bytes=0-1048575" | grep -E "HTTP|Content-Range"
HTTP/1.1 206 Partial Content
Content-Range: bytes 0-1048575/461301120
# a ranged GET: any slice of any object, over plain HTTPS.
# no mount, no filesystem, no single owner - and 206 is the
# status code your parallel loader (sec. 6.3) lives on.
\end{Verbatim}
\end{terminalbox}
\figcaption{터미널 기록 6-3  직접 체감하는 객체 계약. 업로드는 HTTP PUT이고 사전 서명 URL과 Range 헤더를 사용하면 440 MB 객체의 원하는 1 MB를 가져올 수 있다. 병렬 처리는 이 요청을 여러 개 보내는 것뿐이다.}

두 조항을 더하면 계약이 완성된다. \textbf{일관성}: S3는 2020년부터 강한 일관성을 제공한다. \code{PUT} 뒤의 \code{GET}은 새 객체를 반환하며 예외가 없다. 덕분에 10년간 전승되던 우회책이 사라졌다. 따라서 §6.3에서 세울 규율인 변경 불가능한 버전은 일관성 버그가 아니라 \textit{운영}의 제정신을 지키기 위한 것이다. \textbf{계층}: 같은 API 뒤에 저장 가격은 낮아지고 검색 비용은 높아지는 여러 \textbf{스토리지 클래스}(표준 → 비정기 접근 → 아카이브)가 있으며, \textbf{수명 주기 규칙}이 객체 나이에 따라 계층 사이를 이동시킨다. 모델 스토어에서는 플릿 규모에서 중요하다. 작년의 모델 버전 400개에 표준 클래스 가격을 낼 필요는 없다. "현재 버전이 아니면 30일 뒤 비정기 접근으로 이동" 같은 한 줄짜리 수명 주기 규칙이 깔끔한 청구서와 계속 불어나는 청구서를 가른다.

\subsection*{파일, PersistentVolumes, 그리고 솔직한 비교표}

세 번째 계열인 \textbf{파일 스토리지(file storage)}(서비스형 NFS: EFS, Filestore)는 절충 계약이다. 블록보다 GB당 약 네 배 비싼 가격과 네트워크 파일 시스템 지연 시간을 감수하고 POSIX 의미 체계와 많은 동시 읽기·쓰기 주체를 함께 얻는다. 여러 학습 작업이 이어 쓰는 공유 캐시처럼 여러 파드가 \textit{변경 가능한} 파일을 공유해야 할 때 제값을 한다. 하지만 객체 스토어와 로컬 캐시의 조합이 더 저렴하고 빠른 곳에서도 흔히 잘못 구매한다. Kubernetes에서 블록과 파일은 \textbf{PersistentVolume}으로 드러난다. \code{ReadWriteOnce}(RWO) 접근을 사용하는 PVC는 한 번에 노드 하나만 쓸 수 있는, 정확히 디스크 계약인 블록이다. \code{ReadWriteMany}(RWX)는 파일이다. 의미심장하게도 객체 스토리지는 이 목록에 아예 없다. 마운트하는 볼륨이 아니라 호출하는 API이기 때문에 파일 시스템의 방해 없이 읽기 주체 천 개가 공유할 수 있다. (버킷을 파일 시스템으로 마운트하는 s3fs, mountpoint 같은 FUSE 어댑터도 있으며 탐색에는 편리하다. 다만 편의용 연결 계층으로만 보고, 객체 스토리지에 애초에 하지 않은 디스크 약속을 지키게 하는 수단으로 여겨서는 안 된다.)

\begin{tbbtablewrap}
  \begin{tabular}{@{}>{\raggedright\arraybackslash}m{\dimexpr 0.2181\tbltextwidth-2\tabcolsep\relax} >{\raggedright\arraybackslash}m{\dimexpr 0.2660\tbltextwidth-2\tabcolsep\relax} >{\raggedright\arraybackslash}m{\dimexpr 0.2766\tbltextwidth-2\tabcolsep\relax} >{\raggedright\arraybackslash}m{\dimexpr 0.2394\tbltextwidth-2\tabcolsep\relax}@{}}
  \arrayrulecolor{tbbnavy}\specialrule{1pt}{0pt}{0pt}
  \rowcolor{tbbtablehead}{\centering \bfseries\color{tbbnavy}} & {\centering \bfseries\color{tbbnavy}블록(EBS)} & {\centering \bfseries\color{tbbnavy}객체(S3)} & {\centering \bfseries\color{tbbnavy}파일(EFS)} \\
  \arrayrulecolor{tbbnavy}\specialrule{0.7pt}{0pt}{2pt}
  {\textbf{제공 형태}} & {원시 디스크} & {블롭을 다루는 HTTP API} & {공유 POSIX 트리} \\
  \arrayrulecolor{tbbtableborder}\hline
  {\textbf{의미 체계}} & {제자리 임의 읽기·쓰기} & {전체 또는 범위 PUT/GET, 편집 불가} & {POSIX, 동시 접근} \\
  \arrayrulecolor{tbbtableborder}\hline
  {\textbf{지연 시간}} & {밀리초 미만–밀리초} & {첫 바이트까지 수십 밀리초} & {밀리초(네트워크 홉)} \\
  \arrayrulecolor{tbbtableborder}\hline
  {\textbf{처리량}} & {볼륨별 프로비저닝} & {\textbf{병렬: NIC 포화까지}} & {프로비저닝·버스트} \\
  \arrayrulecolor{tbbtableborder}\hline
  {\textbf{공유}} & {인스턴스 하나(RWO)} & {읽기 주체 제한 없음} & {다수 읽기·쓰기(RWX)} \\
  \arrayrulecolor{tbbtableborder}\hline
  {\textbf{내구성}} & {AZ 하나 + 직접 만든 스냅샷} & {\textbf{일레븐 나인(eleven nines), AZ 간 복제}} & {다중 AZ, 관리형} \\
  \arrayrulecolor{tbbtableborder}\hline
  {\textbf{가격 형태}} & {프로비저닝한 용량: 월 GB당 약 \$0.08} & {월 GB당 약 \$0.023 + 요청 + 외부 전송} & {사용량: 월 GB당 약 \$0.30} \\
  \arrayrulecolor{tbbtableborder}\hline
  {\textbf{서빙 용도}} & {데이터베이스 · 노드 임시 공간} & {\textbf{가중치 · 데이터셋 · 이미지 · 로그}} & {공유 변경 가능 캐시} \\
  \arrayrulecolor{tbbnavy}\specialrule{1pt}{2pt}{0pt}
  \end{tabular}\par
  \figcaption{표 6-2  세 가지 계약 비교(대표 가격)}
\end{tbbtablewrap}

\begin{calloutbox}{tbbblue}{tbbbluefill}
{\normalsize\bfseries\color{tbbblue}계산 예시 — 하나의 13 GB 모델을 저장하는 네 가지 방법}\par\vspace{3pt}
13 GB 모델을 노드 5개에 걸친 레플리카 20개에서 사용할 수 있어야 한다. 월 스토리지 비용과 동작을 비교해 보자(대표 가격).

\begin{tbbitemize}
  \item \textbf{객체 스토어(권장 패턴):} 복사본 하나에 13 GB × \$0.023 ≈ \textbf{월 \$0.30}이며, 배포마다 GET 요청 비용이 몇 센트 더 든다. 각 노드는 처음 가져올 때 이미 있는 블록 임시 공간에 로컬 캐시한다. 버전 관리는 새 접두사를 만드는 것으로 거의 무료다.
  \item \textbf{노드별 블록 볼륨:} 5 × 13 GB × \$0.08 ≈ \textbf{월 \$5.20}에 더해 새 버전을 디스크 5개에 원자적으로 배포하고 롤백하는 오케스트레이션까지 직접 책임져야 한다. 더 나쁜 모델 스토어를 다시 만든 셈이다.
  \item \textbf{파일 공유:} 13 GB × \$0.30 ≈ \textbf{월 \$3.90}이며 공유 마운트라 간편하다. 다만 레플리카 20개가 동시에 콜드 스타트해 하나의 공유 처리량을 두고 경합할 때까지다. 버전 관리 규율도 여전히 필요하다.
  \item \textbf{이미지에 구워 넣기(제3장의 안티패턴):} 레지스트리도 어차피 객체 스토리지 위에 있으므로 저장 비용은 저렴하다. 하지만 코드를 배포할 때마다 레지스트리 경로를 통해 13 GB × 노드 5개를 다시 보내고, 모델 버전이 이미지 버전에 용접된다. 실제 비용은 GB가 아니라 워크플로다.
\end{tbbitemize}

승리하는 패턴은 \textbf{객체 스토어를 원본으로, 블록을 노드별 캐시로} 쓰는 조합이다. 한 번 저장하고 로컬에 캐시하며 접두사로 버전을 관리한다. 이 조합에는 이름이 있으며 다음 절에서 다룬다.
\end{calloutbox}

\vspace{3.0pt}

\begin{tbbfigure}
  \adjustbox{max width=0.94\linewidth,max totalheight=0.46\textheight}{\includegraphics{figures/fig6-2_contracts.png}}\\[3pt]
  \figcaption{그림 6-2  두 가지 계약. 블록은 디스크(POSIX, 마이크로초 단위, 소유자 하나)를 팔고, 객체는 HTTP 뒤의 내구성 높은 블롭(범위 읽기, 제한 없는 읽기 주체, 병렬 처리량)을 판다. 파일은 비싼 절충안이다.}
\end{tbbfigure}

\begin{calloutbox}{tbbpurple}{tbbpurplefill}
{\normalsize\bfseries\color{tbbpurple}생각해 보기}\par\vspace{3pt}
\begin{tbbitemize}
  \item PostgreSQL 데이터베이스를 객체 스토리지에서 실행하는 것은 끔찍한 생각이지만 분석 엔진은 하루 종일 S3의 Parquet 파일을 즐겁게 질의한다. 두 워크로드 사이에서 무엇이 달라졌는가? 이를 결정하는 계약 조항 두 가지를 말하자.
  \item S3는 일레븐 나인의 내구성을 약속하지만 가용성 SLA는 약 99.9\%로 훨씬 낮다. 둘의 차이는 무엇이며, 이 간격 때문에 로더(§6.3)는 무엇을 해야 하는가?
  \item 팀이 "현재 모델"을 보관하는 RWX 파일 공유 하나를 두고, 출시 때마다 제자리에서 업데이트하려 한다. 데모에서는 동작한다. 한 분기 안에 문제를 일으킬 세 가지 방식을 나열하자. 다운로드 중 갱신, 롤백, 콜드 스타트 경합을 생각해 보자.
\end{tbbitemize}
\end{calloutbox}

\section{모델 스토어}

제3장 §3.4는 \textit{"모델 가중치는 보통 이미지 레이어에 둘 것이 아니라 모델 스토어에 보관하고 시작할 때 로딩해야 한다"}는 약속을 남기며 끝났다. 이 절에서 그 약속을 지킨다. 논리는 워크로드와 계약이 완벽하게 맞는다는 데 있다. 가중치는 \textbf{매우 크고}(수백 MB에서 수백 GB), \textbf{변경 불가능하며}(학습된 산출물은 바뀌지 않고 새 학습은 새 산출물이다), \textbf{버전이 있고}(여러 버전을 보관하고 전환한다), \textbf{한 번 쓰고 모든 레플리카가 읽으며}, 병렬화하기 좋은 \textbf{대규모 순차 읽기}로 소비된다. §6.2의 객체 계약을 다시 읽어 보면 정확히 같은 목록이다. 가중치는 시스템 전체에서 가장 객체다운 바이트다(그림 6-3).

\begin{tbbfigure}
  \adjustbox{max width=0.94\linewidth,max totalheight=0.46\textheight}{\includegraphics{figures/fig6-3_modelstore.png}}\\[3pt]
  \figcaption{그림 6-3  모델 스토어. 버킷 하나, 이름 규칙 하나, 변경 불가능한 버전, 그리고 URI를 서빙 모델로 바꾸는 시작 경로로 구성된다. 환경 변수가 이미지와 가중치 사이의 유일한 결합 지점이다.}
\end{tbbfigure}

\subsection*{레이아웃: 하나의 규칙과 변경 불가능성}

\begin{terminalbox}
\begin{Verbatim}[breaklines=true,breakanywhere=true,fontsize=\small,formatcom=\color{tbbtermfg},codes=\tbbcodeactive]
$ aws s3 ls s3://ml-models/ --recursive --human-readable | head -12
bert-classifier/v7/model.safetensors        440.0 MiB
bert-classifier/v7/config.json                1.2 KiB
bert-classifier/v7/tokenizer.json           695.4 KiB
bert-classifier/v7/manifest.json              0.6 KiB
bert-classifier/v8/model.safetensors        440.1 MiB
bert-classifier/v8/...
llama-8b-int8/2026-07-01/model-00001.safetensors  4.9 GiB
llama-8b-int8/2026-07-01/model-00002.safetensors  4.9 GiB
llama-8b-int8/2026-07-01/model-00003.safetensors  3.4 GiB
llama-8b-int8/2026-07-01/manifest.json            0.7 KiB
$ cat manifest.json   # (from v7)
{ "model": "bert-classifier", "version": "v7",
  "files": { "model.safetensors":
    "sha256:9f2c4a01..." },
  "trained": "2026-06-14", "framework": "torch-2.6" }
\end{Verbatim}
\end{terminalbox}
\figcaption{터미널 기록 6-4  스토어의 전체 스키마. 모델·버전·파일과 모든 다운로드를 검증할 수 있게 하는 sha256 합계가 든 매니페스트다. 제4장의 내용 주소 지정 이미지 레이어와 의도적으로 운을 맞췄다.}

배치는 민망할 만큼 단순한 \code{<model>/\allowbreak <version>/\allowbreak <files>}이며, 전체 패턴은 한 가지 규율에 기대어 있다. \textbf{버전은 변경할 수 없다.} 새 학습 실행, 새 양자화, 토크나이저(tokenizer) 파일 수정까지 모두 새 접두사를 사용한다. 아무리 마감이 급해도 \code{v7/\allowbreak }을 제자리에서 덮어써서는 안 된다. 이유는 계속 쌓인다. 롤링 업데이트(제5장)에서는 몇 분 동안 기존 파드와 새 파드가 나란히 실행된다. 그 아래의 \code{v7}이 바뀔 수 있으면 \textit{선언한 버전이 같은} 레플리카끼리도 서로 다른 결과를 낸다. 아래에서 볼 노드 캐시는 URI를 키로 사용하므로 내용을 바꾸면 모든 캐시가 조용히 오염된다. 롤백(제5장의 \code{rollout undo}, 13주 차 자동화)은 돌아갈 대상이 변경되지 않은 채 그대로 존재할 때만 의미가 있다. \code{manifest.json}의 체크섬은 모든 소비자가 다운로드를 스스로 검증하게 한다. 제4장의 이미지 다이제스트에 적용한 무결성 아이디어를 스택에서 한 단계 위의 산출물에 적용한 것이다. 변경 불가능성은 "어느 모델이 실행 중인가?"라는 질문을 조사 작업에서 문자열 비교로 바꾼다. URI \textit{자체}가 답이다.

\subsection*{빠르게 로딩하기: 병렬 로더}

이제 성능 측면을 보자. 첫 바이트는 느리지만 합계에는 사실상 상한이 없다는 객체 계약의 형태를 떠올리고, 순진한 로더의 동작을 살펴보자. \code{GET} 하나, 스트림 하나로 운이 좋아야 약 80 MB/s다. 13 GB라면 GPU 노드가 비용만 내며 아무 일도 하지 않는 \textbf{162초}다. 해결책은 \textit{계약을 활용하는 것}이다. 객체를 여러 범위로 나누고 8–16개 범위를 동시에 가져와(터미널 기록 6-3의 \code{206 Partial Content}를 여러 번 실행) 제자리에 쓴다. 총 처리량은 \textbf{NIC}(10 Gbps ≈ 1.25 GB/s 상한)나 로컬 디스크에 닿을 때까지 증가한다. 이 규모에서는 S3 자체가 병목이 되지 않으며 접두사마다 대규모 요청 병렬 처리를 명시적으로 지원한다. 전용 도구인 \code{s5cmd}, 동시성을 설정할 수 있는 AWS CLI, \code{hf\_transfer}도 정확히 이 방식을 구현한다. 실습에서는 최소한의 로더를 직접 작성해 메커니즘을 익힌다.

\begin{terminalbox}
\begin{Verbatim}[breaklines=true,breakanywhere=true,fontsize=\small,formatcom=\color{tbbtermfg},codes=\tbbcodeactive]
# the experiment the lab reproduces (13 GB, same node, same bucket):
$ time python loader.py --streams 1 $MODEL_URI /models
real  2m42.1s                       # ~80 MB/s: one stream
$ time python loader.py --streams 16 $MODEL_URI /models
real  0m16.8s                       # ~790 MB/s: sixteen ranged streams
# 9.6x from the same bucket, same NIC - the contract, exploited.
$ sha256sum -c <(python manifest_sums.py /models/manifest.json)
model-00001.safetensors: OK
model-00002.safetensors: OK
model-00003.safetensors: OK        # verify before you serve. always.
\end{Verbatim}
\end{terminalbox}
\figcaption{터미널 기록 6-5  병렬 로더. 하나의 GET을 16개의 범위 GET으로 바꿔 162 s → 17 s로 줄인 뒤, 첫 요청을 서빙하기 전에 매니페스트를 기준으로 체크섬을 검증한다.}

다운로드 다음에는 \textbf{로드}가 오며, 형식을 하나 잘 고르면 비용이 거의 사라진다. 바로 \textbf{safetensors}다. JSON 헤더(텐서 이름, 자료형, 형상, 오프셋) 뒤에 원시 텐서 바이트가 이어지는 형식이다. 로더가 역직렬화하지 않고 파일을 \code{mmap}해 텐서를 무복사(zero-copy)로 \textit{제자리에서} 쓰도록 설계된 형식이다. pickle 기반 체크포인트와 달리 열어도 코드가 실행될 수 없다는 점은 7주 차 형식 논의의 보안 측면이다. 제3장과 연결되는 결과가 인상적이다. mmap한 가중치 파일은 \textbf{파일 기반 페이지 캐시}다. 커널은 텐서를 처음 건드릴 때 페이지를 들이고 압박이 생기면 회수할 수 있다. 정상 모델 서버의 \code{memory.current}는 제한에 가깝지만 \code{anon}은 적당히 유지되는 이유다. 이제 §3.3의 대시보드 교훈에 직접 설계한 원인이 생겼다. GPU 서빙에서는 호스트 페이지를 VRAM으로 복사하는 단계가 이어지며, 10주 차 최적화 장이 그 지점부터 시작한다.

\subsection*{파드에 연결하기: initContainer, emptyDir, 계산된 프로브}

\begin{terminalbox}
\begin{Verbatim}[breaklines=true,breakanywhere=true,fontsize=\small,formatcom=\color{tbbtermfg},codes=\tbbcodeactive]
    spec:
      volumes:
      - name: models
        emptyDir: {}              # node-local scratch; dies with pod
      initContainers:
      - name: fetch-model
        image: registry.example.com/team/model-loader:v3
        args: ["$(MODEL_URI)", "/models", "--streams", "16"]
        env:
        - { name: MODEL_URI,
            value: "s3://ml-models/bert-classifier/v7/" }
        volumeMounts: [ { name: models, mountPath: /models } ]
      containers:
      - name: server
        image: registry.example.com/team/model-server:v12
        env:
        - { name: MODEL_DIR, value: "/models" }
        volumeMounts: [ { name: models, mountPath: /models } ]
        startupProbe:             # budget = download + load + margin
          httpGet: { path: /healthz, port: 8000 }
          periodSeconds: 5
          failureThreshold: 24    # 24 x 5s = 2 min for the 440 MB case
\end{Verbatim}
\end{terminalbox}
\figcaption{터미널 기록 6-6  연결 방식. 서버가 시작되기 전에 initContainer가 emptyDir에 다운로드한다. 이제 startupProbe 예산(제5장)은 터미널 기록 6-5의 측정값으로 계산하는 숫자다.}

각 요소는 의도된 역할 하나씩을 맡는다. \textbf{initContainer}는 서버 컨테이너가 시작되기 \textit{전에} 완료될 때까지 실행된다. 따라서 서버 이미지는 로더 의존성에서 자유롭고, 다운로드 실패는 별도의 \code{Init:Error} 상태로 나타나며(제5장의 현장 안내서에서 명확히 구별 가능), 로더도 독립적으로 업그레이드할 수 있다. \textbf{emptyDir}은 정직한 임시 공간이다. 노드 디스크에서 파드와 정확히 같은 기간만 존재하며(tmpfs가 필요하면 \code{medium: Memory}), 파드와 함께 삭제되는 것은 버그가 아니라 플랫폼이 아무것도 소유하지 않는 원칙을 강제한 결과다. \textbf{startupProbe} 예산은 더 이상 구전 지식이 아니라 산술이다. 측정한 다운로드(17 s) + 로드(10 s) + 여유분 → 임계값으로 정한다. 그리고 \textbf{환경 변수가 유일한 결합 지점}이다. 서버는 \code{/\allowbreak models}를 읽고 initContainer는 \code{MODEL\_URI}에서 파일을 가져와 채운다. \textit{두 이미지 어디에도 모델 버전은 나오지 않는다.} 이제 이 장에서 가장 멋진 터미널 기록을 볼 준비가 되었다.

\begin{terminalbox}
\begin{Verbatim}[breaklines=true,breakanywhere=true,fontsize=\small,formatcom=\color{tbbtermfg},codes=\tbbcodeactive]
$ kubectl set env deploy/model-server MODEL_URI=s3://ml-models/bert-classifier/v8/
deployment.apps/model-server env updated
$ kubectl rollout status deploy/model-server && kubectl get pods \
    -o custom-columns=NAME:.metadata.name,IMAGE:.spec.containers[0].image
deployment "model-server" successfully rolled out
NAME                            IMAGE
model-server-59d8c7f9b6-2xk4p   registry.example.com/team/model-server:v12
model-server-59d8c7f9b6-m7qvw   registry.example.com/team/model-server:v12
model-server-59d8c7f9b6-tt8zn   registry.example.com/team/model-server:v12
# same image. new model. a rolling, health-gated, rollback-able
# model deploy (Ch. 5 machinery) - and no software was built.
\end{Verbatim}
\end{terminalbox}
\figcaption{터미널 기록 6-7  결실. 이제 모델 배포는 환경 변수 변경이다. IMAGE 열은 그대로 둔 채 제5장의 다른 롤아웃처럼 여분 레플리카를 만들고 준비 상태로 제어하며 되돌릴 수 있다. 13주 차에는 바로 이 한 줄을 자동화한다.}

방금 일어난 일을 잠시 되새겨 보자. 이 장 전체의 아키텍처 핵심이기 때문이다. 두 출시 주기를, 즉 코드는 이미지에 두고 가중치는 스토어에 두는 방식으로 분리하자 \textbf{모델 출시가 소프트웨어 출시이기를 멈췄다.} Dockerfile도, CI 빌드도, 레지스트리 푸시도 없다. "배포"는 필드 하나를 바꾸는 일이며, 제5장에서 만든 모든 안전 메커니즘(여분 레플리카, 준비 상태를 통한 속도 조절, 실패 시 중단, 한 명령으로 롤백)이 그대로 적용된다. 이것 \textit{자체가} 제5장의 롤아웃이기 때문이다. 13주 차 모델 레지스트리는 어느 URI가 "스테이징"이고 어느 것이 "프로덕션"인지, 누가 언제 승격했는지를 형식적으로 기록한다. 하지만 레지스트리는 메타데이터일 뿐이며 실제 메커니즘은 이 터미널 기록에 있다.

\subsection*{콜드 스타트 청구서와 캐싱 단계}

\begin{tbbfigure}
  \adjustbox{max width=0.94\linewidth,max totalheight=0.46\textheight}{\includegraphics{figures/fig6-4_coldstart.png}}\\[3pt]
  \figcaption{그림 6-4  새 레플리카 하나의 콜드 스타트 비용을 항목별로 나눈 모습. 순진한 막대는 단일 스트림 다운로드가 지배하고, 설계한 막대는 같은 레플리카를 4.6배 빠르게 준비 상태로 만든다. 12주 차에는 이 막대를 두고 씨름한다.}
\end{tbbfigure}

\begin{calloutbox}{tbbamber}{tbbamberfill}
{\normalsize\bfseries\color{tbbamberdark}계산 예시 — 레플리카 하나의 콜드 스타트 항목별 계산(13 GB 모델)}\par\vspace{3pt}
항목을 두 번 더해 보자. 대표값이며 실습에서는 직접 측정한다.

\begin{tbbitemize}
  \item \textbf{순진한 방식:} 스케줄링 약 2 s + 이미지 가져오기 60 s(6 GB, 콜드 노드 — 제3장의 계산) + 가중치 162 s(단일 스트림) + 로드 25 s + 워밍업 13 s ≈ \textbf{준비 상태까지 260 s}다.
  \item \textbf{설계한 방식:} 스케줄링 약 2 s + 이미지 약 0 s(노드에 캐시됨) + 가중치 17 s(스트림 16개) + 로드 25 s + 워밍업 13 s ≈ \textbf{준비 상태까지 56 s}다.
\end{tbbitemize}

같은 하드웨어에서 4.6배 차이가 난다. 이제 모든 항목의 책임자를 안다. 이미지 → 레이어 위생과 노드 캐시(제3–4장), 다운로드 → 이 절의 로더, 로드 → safetensors + mmap, 워밍업 → 준비 상태 전 합성 요청(제5장의 프로브, §6.4의 상태 확인 경로)이다. 남은 수단인 파드가 공유하는 노드별 가중치 캐시, 사전 워밍업 풀(pre-warmed pool), 이 막대가 좌우하는 제로 스케일링 논의는 12주 차에서 다룬다. 그때는 이 계산을 숙지했다고 가정한다.
\end{calloutbox}

\vspace{3.0pt}

캐싱 단계를 한 문단으로 미리 살펴볼 가치가 있다. emptyDir 패턴은 \textbf{파드}마다 다시 다운로드한다. URI를 키로 사용하는 \code{hostPath}/PVC 캐시는 \textbf{노드}마다 다시 다운로드한다. 버전은 오래될 수 없고 없을 수만 있으므로 변경 불가능성이 이를 안전하게 만든다. 미리 가져오는 DaemonSet은 파드가 도착하기 \textbf{전에} 노드를 워밍업한다. 가중치를 이미지에 구워 넣는 방식은 망 분리 배포나 데모 장비 같은 정당한 \textit{극단 사례}로 남는다. 제3장에서 이미 "기본값이 아닌 결정"으로 분류했다. 각 단계는 스토리지 중복과 시작 지연 시간을 맞바꾼다. 이 책에서 반복해서 등장하며 GB 단위로 구매할 수 있는 절충 관계다.

스토어에는 마지막 거주자인 \textbf{데이터셋}도 있다. 학습 데이터, 평가 세트, 팀이 회귀 테스트에 사용할 골든 요청(golden request) 픽스처(13주 차 CI 관문)는 보통 Parquet 또는 JSONL 형식으로 같은 버킷(bucket) 규율, 즉 버전이 있는 접두사, 매니페스트, 변경 불가능성을 따른다. 서빙 경로의 지연 시간 문제를 공유하지 않으므로 새 메커니즘도 필요 없다. 스토어는 신경 쓰지 않으며 바로 그것이 핵심이다.

\begin{calloutbox}{tbbpurple}{tbbpurplefill}
{\normalsize\bfseries\color{tbbpurple}생각해 보기}\par\vspace{3pt}
\begin{tbbitemize}
  \item 당직 중 사고를 넘겨받았다. 누군가 "긴급 수정"이라며 v7의 tokenizer.json을 제자리에서 다시 업로드했고, 그 뒤 레플리카 절반이 재시작되었다. 같은 URI가 두 가지 동작을 보이는 다양한 증상을 설명하고, 캐시가 문제를 악화시키는 이유와 이를 막았을 한 줄짜리 정책을 제시하자.
  \item manifest의 sha256은 TLS가 막지 못하는 무언가를 막고, TLS는 manifest가 막지 못하는 무언가를 막는다. 두 위협을 정확히 말하자.
  \item 레플리카 100개로 확장하는 사건(12주 차)에서 모든 새 파드가 동시에 13 GB를 가져온다. (a) 파드별 emptyDir, (b) 노드 캐시, (c) 사전 워밍업 노드일 때 병목은 실제로 S3, 노드별 NIC, 노드 캐시 가운데 어디에 생기는가? 파드 4개를 실행하는 10 Gbps 노드 하나의 계산을 보여 주자.
\end{tbbitemize}
\end{calloutbox}

\section{출입구: 로드 밸런서, Ingress, 게이트웨이}

제5장의 노출 사다리는 Service에 클라우드 통로를 직접 연결하는 \code{type: LoadBalancer}에서 끝났다. 데모에는 실제로 그것만 있으면 된다. 프로덕션의 출입구에는 더 긴 목록이 필요하다. 누구도 평문으로 추론 입력을 보내지 않도록 \textbf{TLS}, 한 출입구에서 여러 서비스로 \code{/\allowbreak predict}, \code{/\allowbreak embed}, \code{staging.api.example.com}을 나누는 \textbf{라우팅}, 제1장의 모델 API 경제를 움직이는 키와 계량기를 위한 \textbf{인증과 할당량}, LB 뒤의 다른 무엇보다 긴 추론 꼬리를 감당하도록 \textbf{의도적으로 선택한 타임아웃과 재시도}, §6.6에서 상태 코드를 읽기 위한 \textbf{엣지에서의 관측 가능성}이 필요하다. 이 절에서는 로드 밸런서, Ingress, 게이트웨이라는 세 계층으로 출입구를 만든 뒤, 특히 모델 서빙을 망가뜨리는 설정을 충분히 살펴본다(그림 6-5).

\begin{tbbfigure}
  \adjustbox{max width=0.94\linewidth,max totalheight=0.46\textheight}{\includegraphics{figures/fig6-5_frontdoor.png}}\\[3pt]
  \figcaption{그림 6-5  요청의 여정. DNS → L7 로드 밸런서(TLS, 타임아웃) → Ingress(클러스터 내부 라우팅) → Service → 파드로 이어진다. 파드 전까지는 모두 합해 약 5 ms이고 p99는 파드에서 결정된다. §6.6에서는 이 사실을 진단 도구로 바꾼다.}
\end{tbbfigure}

\subsection*{L4 대 L7: 연결 대 요청}

로드 밸런서는 \textit{무엇을 읽는지}에 따라 두 종류로 나뉜다. \textbf{L4 로드 밸런서}는 전송 계층에서 동작한다. TCP 연결을 전달할 뿐 HTTP 바이트는 하나도 해석하지 않는다. 눈부시게 빠르고 저렴하며 프로토콜에 구애받지 않지만 지능은 딱 그 정도다. \textbf{L7 로드 밸런서}는 연결을 종료하고 \textit{요청을 읽어} 요청마다 새 결정을 내린다. 경로나 호스트로 라우팅하고, 멱등인 실패를 재시도하며, 헤더를 주입하고, TLS를 종료하고, 클라이언트와는 HTTP/2로, 파드와는 HTTP/1.1로 통신한다. 이 차이는 제5장이 주의사항을 남긴 바로 그 지점에서 빛을 발한다. Service(kube-proxy DNAT)는 \textbf{연결}을 로드 밸런싱한다. 따라서 다중화한 요청 만 개를 운반하는 장기 gRPC 채널 하나는 모두 단일 파드로 간다. 연결 유지 클라이언트 뒤의 플릿은 모든 대시보드가 "정상"이라고 말하는데도 부하가 극단적으로 불균형할 수 있다. L7 홉은 \textbf{요청} 수준에서 다시 로드 밸런싱해 정확히 이 문제를 고친다. 그래서 gRPC와 스트리밍 워크로드(11주 차)에는 L7 프록시가 부속품이 아니라 필수 요소다. 실용적으로는 필요할 때 원시 패킷 배관의 앞단에 L4를 두거나 클라우드에서 클러스터 내부 L7로 들어가는 입구로 쓰고, 요청이 중요한 단위인 곳에는 L7을 둔다. 서빙에서는 거의 모든 곳이 여기에 해당한다.

\subsection*{Ingress: 하나의 출입구, 여러 Service}

Service마다 클라우드 로드 밸런서를 하나씩 프로비저닝하는 것은 초보자가 튜토리얼에서 배우는 비싼 습관이다. 각 LB는 자체 IP와 유휴 비용이 있는 청구 대상 리소스다. Kubernetes의 답은 \textbf{Ingress}다. 여러 Service에 대응하는 호스트와 경로 규칙을 담은 라우팅 객체이며, \textbf{Ingress 컨트롤러}(nginx, Envoy 기반 컨트롤러, 또는 ALB를 직접 프로그래밍하는 클라우드 컨트롤러)가 실제로 구현한다. 로드 밸런서 하나가 클러스터에 들어오면 컨트롤러가 원하는 수의 Service로 요청을 분배한다. TLS 인증서는 호스트마다 붙이며 \code{cert-manager}가 발급과 갱신을 자동화한다. 실습에서 한 번 실행한 뒤에는 인증서를 다시 수동으로 관리하지 말자.

\begin{terminalbox}
\begin{Verbatim}[breaklines=true,breakanywhere=true,fontsize=\small,formatcom=\color{tbbtermfg},codes=\tbbcodeactive]
apiVersion: networking.k8s.io/v1
kind: Ingress
metadata:
  name: serving-ingress
  annotations:
    nginx.ingress.kubernetes.io/proxy-read-timeout: "120"
spec:
  tls:
  - hosts: [ api.example.com ]
    secretName: api-tls            # cert-manager fills this
  rules:
  - host: api.example.com
    http:
      paths:
      - path: /predict
        pathType: Prefix
        backend: { service: { name: model-svc, port: { number: 80 } } }
      - path: /embed
        pathType: Prefix
        backend: { service: { name: embed-svc, port: { number: 80 } } }

$ curl -sv https://api.example.com/predict -d '{"x":[1,2]}' 2>&1 \
    | grep -E "subject:|HTTP/"
*  subject: CN=api.example.com     # real TLS, terminated at the door
< HTTP/2 200
\end{Verbatim}
\end{terminalbox}
\figcaption{터미널 기록 6-8  Ingress 하나가 호스트의 TLS, 두 모델 Service로 가는 경로 라우팅, 의도적으로 설정한 타임아웃 어노테이션을 제공한다(기본값은 §6.6의 사고를 견디지 못한다). 이제 LB 하나가 네임스페이스 전체의 앞에 선다.}

\subsection*{API 게이트웨이: 제품이 사는 곳}

라우팅 위에는 \textbf{API 게이트웨이}가 있다. 기능상 L7 프록시이면서 \textit{제품}도 구현한다. API 키와 인증, 키별 속도 제한과 할당량(클라이언트가 존중하는 법을 배우게 될 \code{429 Too Many Requests}), 키·경로·토큰별 사용량 계측(청구서의 원재료), 요청 검증과 변환을 담당한다. 관리형 제품(AWS API Gateway, Apigee)과 자체 호스팅 제품(Kong, Envoy Gateway)은 운영 방식이 다를 뿐 개념은 같다. 이 책에서 게이트웨이는 제1장에서 시작한 고리를 닫는다. 계측되는 토큰, 계층별 할당량, 키를 제공하며 돈을 내고 사용하는 "모델 API 생태계"는 다른 조직의 플릿 앞에 놓인 \textit{바로 이 상자}다. 강의 프로젝트도 최소한 같은 형태를 요구한다. TLS 뒤에 키와 할당량을 갖춘 모델이 "curl에 응답하는 컨테이너"와 "누군가에게 비용을 청구할 수 있는 서비스"를 가른다. 중복에 관한 주의사항도 있다. 게이트웨이, Ingress, 메시(14주 차)는 모두 L7을 다루므로 성숙한 스택은 의도적으로 계층을 합친다. 경험 법칙은 약어마다 상자 하나를 두는 것이 아니라 \textit{관심사마다 소유자 하나}를 두는 것이다. TLS는 한 번만 종료하고, 속도 제한과 재시도는 각각 한곳에만 둔다.

\subsection*{타임아웃, 재시도, 연결 비우기: 추론을 망가뜨리는 설정}

출입구의 기본값은 수십 KB와 수백 ms인 웹페이지에 맞춰져 있으며 모델 서빙은 그 문장의 모든 가정을 어긴다. 대부분의 문제는 세 가지 설정에서 생긴다. \textbf{요청 타임아웃}: LB는 보통 30–60 s 안에 응답하지 않은 상위 서버를 포기하고 \code{504}를 반환한다. 제3장에서 추론 지연 시간에는 배치와 길이에 비례해 치솟는 \textit{형태}가 있다고 배웠다. 따라서 타임아웃은 최악 입력에서 측정한 정직한 p99(제1장의 SLO 규율)보다 길어야 한다. 그렇지 않으면 긴 요청은 출입구에서 죽고 GPU는 아무도 받지 않을 작업을 마친다. \textbf{유휴 타임아웃(idle timeout)}: 대부분의 LB는 N초 동안 \textit{바이트가 하나도 움직이지 않는} 연결을 별도로 정리한다(ALB 기본값: 60초). 생성 도중 토큰 사이에서 70 s 멈추는 스트리밍 LLM 응답(11주 차)은 이 제한에 걸린다. §6.6의 징후대로 연결이 \textit{정확히} N초에 끊긴다. 유휴 타임아웃을 늘리거나 연결 유지 프레임을 보내 해결한다. \textbf{재시도}: L7 출입구는 눈에 띄지 않게 실패를 재시도할 수 있다. 멱등 GET에는 선물이지만 추론에는 위험하다. 30초짜리 생성을 재시도하면 시스템이 느려진 바로 그때 GPU 부하가 \textit{두 배}가 된다. 게이트웨이 × 클라이언트 두 계층에서 타임아웃을 각각 재시도하면 3 × 3 = 사용자 동작 하나당 아홉 번 시도가 된다. 필요한 규율은 \textbf{재시도 예산(retry budget)}이다. 타임아웃은 절대 재시도하지 않고 연결 수준의 실패만 단일 계층에서 지터(jitter)를 넣은 지수 백오프(exponential backoff)와 함께 재시도한다. 9주 차 부하 테스트에서는 재시도 폭풍(retry storm)을 의도적으로 만들어 파형을 익힌다.

네 번째 설정은 모든 팀이 언젠가 만나는 증상인 \textbf{롤아웃 중 정확히 발생하는 502 급증}을 설명하므로 문단 하나를 온전히 쓸 가치가 있다. 연결을 비우는 두 시계가 일치해야 한다. LB는 자체 지연(\code{deregistration delay} / connection draining) 뒤 종료 중인 파드를 등록 해제하고, Kubernetes는 \code{terminationGracePeriodSeconds} 뒤 파드를 죽인다(제5장의 과정). LB가 여전히 유효한 대상으로 보는 동안 파드가 죽으면, 즉 유예 시간이 너무 짧거나 지연이 너무 길거나 서버가 SIGTERM을 보자마자 연결을 비우지 않고 종료하면(제3장의 핸들러가 다시 등장한다), 처리 중 요청이 시체에 부딪히고 출입구가 502를 보고한다. 해결책은 시계를 맞추는 것이다. 먼저 준비 상태 검사가 실패해 엔드포인트에서 빠지고 트래픽 유입이 멈춘다. 그런 다음 처리 중 작업을 비운 뒤 종료한다. 유예 기간은 가장 긴 정당한 요청보다 길어야 한다. 이 문장의 모든 계층은 이미 읽은 장에 있으며, 출입구는 불일치를 고객에게 보이게 할 뿐이다.

\begin{terminalbox}
\begin{Verbatim}[breaklines=true,breakanywhere=true,fontsize=\small,formatcom=\color{tbbtermfg},codes=\tbbcodeactive]
# the streaming incident, reproduced on purpose (lab, part C):
$ time curl -sN https://api.example.com/generate -d '{"n_tokens": 999}'
...tokens......tokens.......                       # then silence...
curl: (18) transfer closed with outstanding read data remaining
real  1m0.02s                       # <- suspiciously exactly 60s
# repeat it: 60.01s. 60.02s. the roundness IS the diagnosis:
# no bytes moved for idle-timeout seconds -> the LB reaped it.
$ kubectl annotate ingress serving-ingress \
    nginx.ingress.kubernetes.io/proxy-read-timeout="300" --overwrite
$ time curl -sN https://api.example.com/generate -d '{"n_tokens": 999}'
real  1m34.8s                       # the stream now outlives its pauses
\end{Verbatim}
\end{terminalbox}
\figcaption{터미널 기록 6-9  유휴 타임아웃에 의한 끊김. 정확히 60.0 s에 두 번 죽는 스트림은 서버 버그가 아니라 출입구의 기본값 때문이다. 11주 차 토큰 스트림의 생사는 이 설정에 달렸다.}

\begin{calloutbox}{tbbpurple}{tbbpurplefill}
{\normalsize\bfseries\color{tbbpurple}생각해 보기}\par\vspace{3pt}
\begin{tbbitemize}
  \item 상태 확인도 요청이다. LB 노드 2개 × 10초마다 × 대상 20개에 kubelet의 준비 상태 프로브까지 더해 백그라운드 QPS를 계산하자. 그런 다음 상태 확인을 실제 부하로 바꾸는 실수를 말하자. 순전파를 실행하는 비싼 /healthz를 생각하면 된다.
  \item p99가 8 s지만 가장 큰 배치와 길이를 사용한 최악의 정당한 요청은 45 s가 걸리는 서비스의 타임아웃 사다리를 설계하자. 클라이언트, 게이트웨이, LB 요청, uvicorn 타임아웃 가운데 무엇이 가장 커야 하며, 순서가 선택 사항이 아닌 이유는 무엇인가?
  \item 게이트웨이는 연결 실패를 한 번 재시도하고 클라이언트 라이브러리는 모든 5xx를 두 번 재시도하며 메시(14주 차)도 자체 재시도를 추가한다. 피크 부하에서 노드가 죽었다. 요청이 어떻게 늘어나고 플릿이 실제로 어떤 부하를 보는지 그린 뒤, 강제할 한 문장짜리 재시도 정책을 작성하자.
\end{tbbitemize}
\end{calloutbox}

\section{경계: VPC, 보안 그룹, 저렴한 경로}

지금까지는 패킷이 흐를 수 있다고 가정했다. 이 절에서는 \textit{어떤} 패킷이 흐를 수 있는지 정한다. 결정 단위는 클라우드 네트워크의 사설 영역인 \textbf{VPC}다. \code{10.0.0.0/\allowbreak 16} 같은 주소 공간을 고르고 가용 영역별 \textbf{서브넷(subnet)}으로 나눈 뒤 \textbf{라우팅 테이블}로 연결한다. 라우팅 테이블 항목 두 개가 두 가지 서브넷을 정의한다. \textbf{공개 서브넷}은 \code{0.0.0.0/\allowbreak 0}을 \textit{인터넷 게이트웨이}로 라우팅하므로 내부의 리소스가 공개 IP를 갖고 외부에서 접근될 수 있다. \textbf{사설 서브넷}에는 그런 경로가 없다. 인스턴스에 공개 IP가 없고, 방화벽 관리자의 성실함이 아니라 \textit{구조 자체로} 인터넷에서 접근할 수 없다. 서빙에 적합한 자세도 곧바로 나온다. 그림 6-6처럼 \textbf{로드 밸런서만 공개하고 GPU 플릿은 사설 영역에 둔다.} 시간당 \$30인 공개 IP 노드는 결제 계정이 붙은 공격 표면이지만 사설 노드는 공격자가 접근할 수 없다.

\begin{tbbfigure}
  \adjustbox{max width=0.94\linewidth,max totalheight=0.46\textheight}{\includegraphics{figures/fig6-6_vpc.png}}\\[3pt]
  \figcaption{그림 6-6  최소 서빙 VPC. 공개 출입구 하나(LB), 사설 하드웨어(플릿), 모델 스토어로 가는 두 경로인 무료 게이트웨이 엔드포인트와 사용량에 따라 과금되는 NAT 우회로로 구성된다. 의도적으로 선택해야 한다.}
\end{tbbfigure}

\subsection*{보안 그룹: 참조로 연결하는 상태 저장 허용 목록}

VPC 안에서 \textbf{보안 그룹}(SG)은 인스턴스별 방화벽이며, 세 가지 성질이 핵심이다. 첫째, \textbf{허용 전용}이다. SG는 연결을 허용할 포트, 프로토콜, 출발지를 나열하고 목록에 없는 것은 모두 거부하므로 "거부 규칙 순서" 퍼즐이 없다. 둘째, \textbf{상태를 저장}한다. 인바운드 연결을 허용하면 응답 패킷은 자동으로 흐른다. 클라우드는 제3장의 \code{-p 8080:8000} 경로에서 본 conntrack과 똑같이, 추상화 한 단계 위에서 연결을 추적한다. 셋째, 아키텍처를 읽기 쉽게 하는 성질로 \textbf{출발지에는 CIDR 범위뿐 아니라 다른 SG도 지정할 수 있다.} 서빙 연결 고리는 세 문장으로 정리된다. LB의 SG는 \code{0.0.0.0/\allowbreak 0}에서 오는 \code{:443}을 허용한다. 노드의 SG는 \textit{LB의 SG에서} 오는 서비스 포트를 허용한다. LB가 확장되거나 이동할 때 조용히 바뀌는 IP 범위가 아니라 SG를 참조한다. SSH는 어디에서 오든 허용하지 않으며, 사람이 꼭 들어가야 할 때는 감사 기록이 남는 SSM/IAP 세션 도구를 사용한다. 참조 연결은 인프라가 바뀌어도 의도를 보존한다. 내일 로드 밸런서의 IP가 무엇이든 "노드는 \textit{로드 밸런서에서} 오는 트래픽을 허용한다"는 사실이 유지된다. (SG 아래에는 상태가 없고 서브넷 수준에서 규칙 번호로 동작하는 네트워크 ACL도 있다. 관례상 열어 두고 SG로 의도를 표현하며, 서브넷 전체를 거칠게 차단할 때만 NACL을 사용한다.)

\begin{terminalbox}
\begin{Verbatim}[breaklines=true,breakanywhere=true,fontsize=\small,formatcom=\color{tbbtermfg},codes=\tbbcodeactive]
# perimeter, demonstrated: (1) the node, from the internet:
$ curl -m 5 http://203.0.113.40:30080/predict
curl: (28) Connection timed out       # no route, no public IP. good.
# (2) the door, from the internet:
$ curl -s https://api.example.com/predict -d '{"x":[1,2]}'
{"y": 0.87, "model": "v8"}            # the one sanctioned path
# (3) identity, from inside a pod - no keys anywhere:
$ kubectl exec -it model-server-...-2xk4p -- \
    aws sts get-caller-identity --query Arn --output text
arn:aws:sts::1234:assumed-role/model-server-irsa/...
$ kubectl exec -it model-server-...-2xk4p -- \
    aws s3 ls s3://ml-models/bert-classifier/v8/   # allowed prefix
2026-07-02  461301120 model.safetensors
$ kubectl exec -it model-server-...-2xk4p -- \
    aws s3 ls s3://corp-finance-exports/
An error occurred (AccessDenied)      # least privilege, working
\end{Verbatim}
\end{terminalbox}
\figcaption{터미널 기록 6-10  세 가지 프로브로 살펴본 경계. 외부에서 접근할 수 없는 노드, 응답하는 LB, 범위가 제한된 IAM 역할만 자격 증명으로 가진 파드가 보인다. 애초에 비밀 값이 아니었으므로 외부로 빼돌릴 수도 없다.}

\subsection*{키 없는 아이덴티티와 가중치가 지나야 할 경로}

터미널 기록의 세 번째 프로브는 자세히 풀어 볼 가치가 있다. 서빙 시스템이 팀을 가장 자주 곤란하게 하는 부분이 자격 증명 처리이기 때문이다. 안티패턴은 환경 변수에 액세스 키를 넣는 것이다. 수명이 길고 복사해서 붙여 넣을 수 있으며 모든 \code{kubectl describe}와 충돌 덤프에 노출되지만 한 번도 교체하지 않는다. 권장 패턴은 \textbf{워크로드 아이덴티티}(EKS의 IRSA, GKE의 Workload Identity)다. 제5장의 객체 목록에 나온 Kubernetes 아이덴티티인 파드의 \textit{서비스 계정}을 IAM 역할에 바인딩하면, 파드 안의 클라우드 SDK가 토큰 체인을 통해 수명이 짧은 자격 증명을 자동으로 얻는다. 누출할 것도 교체할 것도 없다. 운영상 특히 중요한 점은 \textbf{최소 권한을 워크로드별로 적용}할 수 있다는 것이다. 모델 서버 역할은 \code{ml-models/\allowbreak bert-classifier/\allowbreak *}만 읽고, 로더는 읽을 수 있지만 쓸 수 없으며, 다른 곳의 학습 파이프라인은 새 접두사를 쓸 수 있지만 덮어쓸 수 없다. §6.3의 변경 불가능성 규율을 팀 규범에서 강제 속성으로 바꾼다.

마지막으로 비용을 살펴보자. 사설 서브넷에서 인터넷으로 나가는 트래픽은 \textbf{NAT 게이트웨이(NAT gateway)}를 통과하고, NAT 게이트웨이는 시간당 비용에 더해 \textbf{처리한 GB마다} 비용을 받는다(대표값 GB당 약 \$0.045). 함정은 S3가 VPC \textit{밖에} 있다는 점이다. 그래서 사설 노드가 S3에서 13 GB 가중치를 가져오는 트래픽은 \textit{기본적으로} NAT 경로를 탄다. 해결책은 \textbf{게이트웨이 VPC 엔드포인트}다. S3/DynamoDB로 가는 트래픽을 제공자 백본을 통해 직접 보내는 라우팅 테이블 항목으로, \textbf{데이터 처리 비용이 0}이고 보통 처리량도 더 높다. ML 인프라에서 투자 대비 수익이 가장 큰 네트워크 설정 한 줄이라고 해도 과언이 아니다.

\begin{calloutbox}{tbbred}{tbbamberfill}
{\normalsize\bfseries\color{tbbamberdark}계산 예시 — 가중치에 붙는 NAT 세금}\par\vspace{3pt}
노드 20개 플릿이 매일 배포할 때마다 13 GB 모델을 다시 가져오고 오토스케일링에 따른 교체도 생긴다. 하루에 노드 가져오기 25회라 하면 S3 → 노드 트래픽은 \textbf{하루 325 GB}다.

\begin{tbbitemize}
  \item \textbf{NAT 게이트웨이 경유:} 325 GB × \$0.045 ≈ 하루 \$14.6 ≈ \textbf{월 \$440}다. 제공자 네트워크를 떠나지도 않은 바이트에 내는 비용이다. NAT 자체의 시간당 약 \$0.045, 즉 월 약 \$32의 기본료도 더해진다.
  \item \textbf{S3 게이트웨이 엔드포인트 경유:} 생성 시 라우팅 테이블 항목 하나만 추가하면 같은 바이트가 \textbf{\$0.00}이다.
\end{tbbitemize}

같은 아키텍처와 같은 트래픽에서 연간 약 \$5,000 차이가 난다. 더구나 실패는 조용하다. NAT를 통해서도 모든 것이 \textit{동작}하지만 비용이 청구될 뿐이다. 점검 규칙은 S3와 통신하는 사설 서브넷에 첫날부터 게이트웨이 엔드포인트를 두는 것이다. (리전 간 전송은 별도의 세금이다. 모델 스토어를 플릿과 같은 리전에 두지 않으면 GB당 약 \$0.02와 추가 지연 시간을 내야 한다. 12주 차 다중 리전 주제에서 다룬다.)
\end{calloutbox}

\vspace{3.0pt}

\begin{calloutbox}{tbbpurple}{tbbpurplefill}
{\normalsize\bfseries\color{tbbpurple}생각해 보기}\par\vspace{3pt}
\begin{tbbitemize}
  \item "노드 SG는 LB SG의 트래픽을 허용한다"가 "LB 서브넷의 CIDR에서 오는 트래픽을 허용한다"보다 확실히 나은 이유는 무엇인가? 시간이 지나면서 CIDR 규칙에 생기는 실패와 SG 참조에는 그 실패가 생길 수 없는 이유를 말하자.
  \item 공격자가 서빙 파드 안에서 코드 실행에 성공했다(악성 pickle로 7주 차 형식 경고가 현실이 됨). (a) 환경 변수에 s3:* 권한이 있는 액세스 키를 둔 경우와 (b) 접두사 하나의 읽기 전용으로 범위를 제한한 IRSA의 장애 영향 범위를 비교하자. 각각 어디까지 접근할 수 있으며 제3장의 공유 커널 주의사항은 무엇을 더하는가?
  \item 청구서에서 NAT 게이트웨이가 GPU 다음이자 스토리지보다 큰 세 번째 항목이다. 가능한 트래픽(가중치, 이미지 가져오기, 빌드 때의 pip 설치, 원격 측정)을 재구성하고, 그 가운데 게이트웨이 엔드포인트가 해결하는 것과 레지스트리 미러 또는 인터페이스 엔드포인트가 필요한 것을 구분하자.
\end{tbbitemize}
\end{calloutbox}

\section{출입구 읽기: 상태 코드 현장 안내서}

제5장의 현장 안내서는 클러스터 안에서 파드 상태를 읽는 법을 가르쳤다. 이 장의 실패는 \textit{바깥}에서 자신을 알린다. 사용자, 합성 검사, 또는 14주 차 경보가 상태 코드 하나만 건네준다. 제5장과 마찬가지로 좋은 소식이 있다. §6.4 출입구의 각 계층은 자기 어휘로만 실패하므로 명령을 하나도 실행하기 전에 \textbf{코드가 실패 위치를 좁혀 준다.} 포렌식을 시작하기 전 검색 공간을 절반으로 줄이는 동작도 하나 있다. \textbf{내부에서 요청을 다시 실행}한다(\code{kubectl port-forward}로 파드에 직접 연결). 내부는 정상이고 외부만 고장이라면 TLS, 라우팅, 타임아웃, 상태 확인, 경계 같은 출입구 문제다. 둘 다 고장이면 플릿 문제이며, 단서를 하나 얻은 채 제5장의 안내서로 돌아간다(그림 6-7).

\begin{tbbfigure}
  \adjustbox{max width=0.94\linewidth,max totalheight=0.46\textheight}{\includegraphics{figures/fig6-7_fieldguide.png}}\\[3pt]
  \figcaption{그림 6-7  해독표. 502(파드가 연결을 끊음), 503(응답할 대상 없음), 504(출입구 기준으로 너무 느림), 정확히 N초에 끊기는 스트림, 조용한 연결 타임아웃, 429(정상 작동하는 계량기)를 정리했다. 실습에서 모두 재현한다.}
\end{tbbfigure}

\subsection*{포렌식 따라 하기: 대시보드는 초록색인데 503}

\begin{terminalbox}
\begin{Verbatim}[breaklines=true,breakanywhere=true,fontsize=\small,formatcom=\color{tbbtermfg},codes=\tbbcodeactive]
$ curl -s -o /dev/null -w "%{http_code}\n" \
    https://api.example.com/predict -d '{"x":[1,2]}'
503                                  # "Service Unavailable" - from whom?
$ kubectl get pods -l app=model      # the fleet looks... fine?
NAME                            READY   STATUS    RESTARTS   AGE
model-server-59d8c7f9b6-2xk4p   1/1     Running   0          3h
model-server-59d8c7f9b6-m7qvw   1/1     Running   0          3h
$ kubectl port-forward pod/model-server-59d8c7f9b6-2xk4p 8000:8000 &
$ curl -s -o /dev/null -w "%{http_code}\n" \
    localhost:8000/predict -d '{"x":[1,2]}'
200                                  # inside OK + outside broken = DOOR
$ kubectl describe ingress serving-ingress | grep -A2 Backends
  /predict   model-svc:80   (<none>)          # <- the controller sees
$ kubectl get endpoints model-svc               #    no endpoints...
NAME        ENDPOINTS   AGE
model-svc   <none>      3h                      # Ch. 5's quiet failure -
# a readiness typo after this morning's change - now visible as a 503
# to every customer. fix the probe path; endpoints repopulate; 200.
\end{Verbatim}
\end{terminalbox}
\figcaption{터미널 기록 6-11  검색 공간을 나누는 동작. 외부는 503이고 포트 포워딩은 200이므로 출입구 문제다. 이어서 출입구 자체의 증거인 Ingress 백엔드 목록과 비어 있는 엔드포인트를 확인한다. 제5장의 조용한 실패를 거리에서 바라본 모습이다.}

터미널 기록처럼 연결 고리를 따라가면 각 홉이 질문 하나에 답한다. \textbf{상태 코드}는 실패 계층의 어휘를 알려 준다. \code{503}은 프록시가 "정상 상위 서버가 없다"고 말하는 것으로, DNS나 TLS가 아니라 엔드포인트, 준비 상태, 또는 LB 자체 상태 확인을 가리킨다. \textbf{포트 포워딩 프로브}는 파드가 서빙할 수 있음을 증명해 "API가 중단되었다"를 "출입구에 보낼 대상이 없다"로 바꾼다. 마지막으로 제5장에서 가장 가치 있는 습관인 \textbf{엔드포인트 확인}이 원인을 찾는다. 두 장이 하나의 진단 언어로 이어진다. 이 책의 인프라 절반을 \textit{완성}한다는 의미가 바로 이것이다.

\subsection*{제3–5장에 뿌리를 둔 나머지 해독표}

\begin{tbbitemize}
  \item \textbf{502 Bad Gateway — 호출 도중 파드가 연결을 끊었다.} 출입구는 연결했지만 백엔드가 닫거나 거부했다. 이미 아는 원인은 요청 중 충돌, 흔히 제3장의 OOM 강제 종료(종료 코드 137: 먼저 \code{Last State} 확인), 또는 §6.4의 롤아웃 시계 불일치다. LB가 여전히 대상으로 보는 동안 파드가 죽었다면 등록 해제 지연과 \code{terminationGracePeriodSeconds}를 비교하고 서버가 SIGTERM 뒤 연결을 비우는지 확인한다. \textit{배포 시점에 급증하는 502는 반증될 때까지 연결 비우기(connection draining) 버그다.}
  \item \textbf{504 Gateway Timeout — 출입구가 보기에는 너무 느리다.} 파드는 작업 중이지만 출입구가 먼저 포기했다. 제3장은 추론 지연 시간의 급증 형태(배치 × 시퀀스)를 보여 주었고 제1장은 SLO 언어를 제공했다. 해결책은 "LB/Ingress 요청 타임아웃을 정직하게 측정한 최악 p99보다 높이거나, 입력에 상한을 두거나, 바이트가 계속 흐르도록 스트리밍한다"다. 시스템이 가장 느릴 때 GPU 부하를 두 배로 만드는 "재시도 추가"는 답이 아니다(§6.4의 예산).
  \item \textbf{정확히 N초에 죽는 스트림} — §6.4의 유휴 타임아웃이며 딱 떨어지는 시간 자체가 진단의 전부다. 두 번 재현하고 숫자를 읽은 뒤 유휴 타임아웃을 늘리거나 연결 유지 프레임을 추가한다. 11주 차 토큰 스트림 때문에 이 책에서 가장 많이 재현하는 터미널 기록이 된다.
  \item \textbf{연결 타임아웃 / 거부 — HTTP에 도달하기 전.} 아무것도 응답하지 않았다. 잘못된 노출 단계(제5장의 사다리), 삭제된 LB를 가리키는 DNS, 또는 예상과 달리 자신을 상대로 \textit{설계대로 동작하는} §6.5의 경계가 원인이다. 사무실 IP 범위를 한 번도 허용하지 않은 SG나 공개라고 생각한 사설 엔드포인트가 예다. 누가 어디에서 어떤 포트로 무엇에 접근할 수 있는지 SG 연결 고리를 따라간다.
  \item \textbf{429 Too Many Requests — 제품 계층이 정상 작동한다.} 게이트웨이의 키별 할당량이 발동했다. 장애가 아니라 플릿과 청구서를 보호하는 계량기다. 클라이언트는 \code{Retry-After}를 지키고 지터를 넣은 백오프를 추가한다. 운영자는 현실에 맞는 키별 할당량을 설정하고 정상 사용량에서도 429가 계속되면 용량을 계획한다.
\end{tbbitemize}

\subsection*{두 계층 상태 확인의 함정}

한 가지 실패는 내부에서 보이지 않으므로 별도의 하위 절로 다룰 가치가 있다. 클라우드 LB는 Kubernetes 준비 상태와 독립적으로 대상에 \textbf{자체 상태 확인}을 실행한다. 두 상태 판단이 어긋날 수 있다. LB가 서버에 없는 \code{/\allowbreak health}를 검사해 404를 "비정상"으로 해석하거나, kube-proxy가 다른 곳으로 라우팅하는 노드의 NodePort를 검사하거나, 준비 상태보다 더 엄격한 타임아웃을 사용할 수 있다. 그 결과 \textbf{kubectl은 모두 준비 상태라고 표시하지만 LB의 정상 대상은 0이고 고객은 503을 본다.} 터미널 기록 6-11과 증상은 같지만 원인이 다르다. \code{kubectl get endpoints}에는 대상이 \textit{채워져} 나오므로 특히 혼란스럽다. 예방 규칙은 \textbf{상태 확인 엔드포인트 하나를 유일한 진실의 원천으로 삼는 것}이다. LB의 검사 경로, 준비 상태 프로브, 게이트웨이 검사가 모두 같은 의미의 \code{/\allowbreak ready}를 가리키게 한다. 실습에서 이 상황을 재현할 때 증거의 위치를 주목하자. Kubernetes 안이 아니라 LB의 대상 상태 콘솔에 있다. 이 책에서 처음으로 클러스터 밖의 제어판으로 이동하는 순간이며, 거부할 수 있는 모든 계층을 관측 가능성이 포괄해야 한다는 14주 차 주제를 미리 보여 준다.

\begin{tbbtablewrap}
  \begin{tabular}{@{}>{\raggedright\arraybackslash}m{\dimexpr 0.2181\tbltextwidth-2\tabcolsep\relax} >{\raggedright\arraybackslash}m{\dimexpr 0.3670\tbltextwidth-2\tabcolsep\relax} >{\raggedright\arraybackslash}m{\dimexpr 0.4149\tbltextwidth-2\tabcolsep\relax}@{}}
  \arrayrulecolor{tbbnavy}\specialrule{1pt}{0pt}{0pt}
  \rowcolor{tbbtablehead}{\centering \bfseries\color{tbbnavy}신호} & {\centering \bfseries\color{tbbnavy}출입구의 의미} & {\centering \bfseries\color{tbbnavy}우선 확인 항목} \\
  \arrayrulecolor{tbbnavy}\specialrule{0.7pt}{0pt}{2pt}
  {\textbf{502}} & {백엔드가 호출 도중 닫거나 거부} & {파드 재시작, Last State(OOM 137?), 롤아웃 연결 비우기 시계} \\
  \arrayrulecolor{tbbtableborder}\hline
  {\textbf{503}} & {정상 상위 서버가 전혀 없음} & {kubectl get endpoints, 이어서 LB 대상 상태(두 계층 함정)} \\
  \arrayrulecolor{tbbtableborder}\hline
  {\textbf{504}} & {상위 서버가 출입구 타임아웃보다 느림} & {LB/Ingress 타임아웃과 측정한 최악 p99 비교} \\
  \arrayrulecolor{tbbtableborder}\hline
  {\textbf{정확히 N s에 끊김}} & {N s 동안 바이트 없음 → 유휴 정리} & {time curl -N을 두 번, 유휴 타임아웃 상향 / 연결 유지} \\
  \arrayrulecolor{tbbtableborder}\hline
  {\textbf{타임아웃, HTTP 없음}} & {출입구에 도달하지 못함} & {DNS → SG 연결 고리 → 노출 단계(제5장 사다리)} \\
  \arrayrulecolor{tbbtableborder}\hline
  {\textbf{429}} & {설계대로 할당량 강제} & {키별 게이트웨이 사용량, Retry-After에 따른 클라이언트 백오프} \\
  \arrayrulecolor{tbbnavy}\specialrule{1pt}{2pt}{0pt}
  \end{tabular}\par
  \figcaption{표 6-3  한눈에 보는 출입구 오류 어휘}
\end{tbbtablewrap}

\begin{calloutbox}{tbbpurple}{tbbpurplefill}
{\normalsize\bfseries\color{tbbpurple}생각해 보기}\par\vspace{3pt}
\begin{tbbitemize}
  \item 사용자가 배포와 무관하게 요청의 약 1\%에서 간헐적인 502가 발생한다고 보고했다. 파드 재시작은 0이다. 이 절을 바탕으로 의심 목록을 만들고 구별할 측정값을 제시하자. 어떤 파드가 긴 요청을 처리하는지, LB 연결 재사용, uvicorn 워커 재활용을 생각해 보자.
  \item 합성 검사(14주 차)는 고객보다 \textit{먼저} 표 6-3의 각 행을 찾아야 한다. 503, 504, 스트림 끊김을 함께 포괄하는 프로브 세 개의 경로, 주기, 경보 조건을 설계하고, 각 프로브가 유효하려면 어느 계층 밖에서 실행되어야 하는지 밝히자.
  \item 두 계층 상태 확인 함정은 증거를 Kubernetes 밖에 두었다. 그림 6-5에서 kubectl 사용자가 볼 수 없는 상태를 유지하는 모든 계층을 나열하고, 각 계층에서 "여기는 초록색인데 저기는 빨간색"이 무엇을 뜻하는지 설명하자.
\end{tbbitemize}
\end{calloutbox}

\section{실습: S3에서 로딩하는 서빙 구조 완성}

이번 주 강의계획서의 문구는 "S3에서 모델을 로딩하는 서빙 구조를 완성한다"이며, 핵심 단어는 \textit{완성}이다. 실습을 마치면 스토어 → 로더 → 플릿 → Service → 출입구로 이어지는 이 책의 서빙 골격이 노트북에서 처음부터 끝까지 존재하고, 모든 연결부에 실제 측정값이 생긴다. 스토리지 절반은 minikube 클러스터 안에 배포한 S3 호환 저장소 \textbf{MinIO}를 사용한다. 같은 API와 SDK를 쓰고 \code{aws} 대신 \code{mc}를 사용하며 비용은 0이다. 클라우드 경계 절반은 §6.5의 체크리스트로 진행하고, 크레딧이 있는 팀을 위해 관리형 클러스터 선택 부록을 제공한다. 전체 과정은 \code{ml} 네임스페이스에서 작업한다.

\subsection*{파트 A — 스토어를 세우고 로더를 정직하게 만들기}

강의 매니페스트로 MinIO를 배포한다. Deployment, Service, PVC로 구성되어 있으며 모든 객체가 제5장의 명사이므로 읽어 보자. 버킷을 만들고 실습 모델을 \textbf{두 번}, 즉 매니페스트가 있는 \code{v7/\allowbreak }과 \code{v8/\allowbreak } 접두사로 업로드한다. \code{mc mb}, \code{mc cp}, \code{mc ls}를 사용해 첫 순간부터 변경 불가능성 규율을 연습한다. 이어서 이번 주의 기준이 될 측정을 수행한다. 제공된 순진한 로더(스트림 하나)를 4 GB 실습 모델에 실행해 시간을 기록한다. 이를 N개의 동시 범위 GET으로 확장한다. 유인물은 스레드 풀 골격을 제공하고 범위 계산과 재조립은 직접 작성한다. N ∈ \{1, 4, 8, 16, 32\}를 훑어 시간 대 N 그래프를 그리고 굴곡점을 찾는다. 마지막으로 승리를 선언하기 전에 \code{manifest.json} 체크섬으로 모든 파일을 검증한다. 제출 산출물은 탐색 표와 병목을 밝히는 한 문장이다. 클러스터 내부 MinIO에서는 실제 S3가 NIC를 포화시키기 전에 노드 디스크나 veth가 보통 병목이 되므로 \code{iostat} 증거와 함께 어느 쪽인지 말한다.

\subsection*{파트 B — 패턴을 연결하고 환경 변수 변경 하나로 모델 배포하기}

터미널 기록 6-6을 실제로 연결한다. initContainer(이제 이미지가 된 파트 A 로더), \code{emptyDir}, \code{MODEL\_URI} 환경 변수, 그리고 파트 A에서 측정한 다운로드+로드 시간으로 \code{failureThreshold × periodSeconds}를 \textit{계산한} \code{startupProbe}를 사용한다. YAML 주석에 계산을 적는다. 제5장의 구절에 이 장의 숫자를 넣는 작업이다. 배포한 뒤 자기 파드에서 아무것도 소유하지 않는 점검(터미널 기록 6-1)을 확인한다. 이어서 두 장면을 재현한다. \textbf{가중치가 있는 조정}: curl 루프 도중 파드를 삭제하고 대체 파드가 준비 상태가 될 때까지의 간격을 잰다. 직접 측정하는 첫 콜드 스타트다. 파드 이벤트와 로그를 바탕으로 스케줄링 / 초기 다운로드 / 로드 / 워밍업 항목을 그림 6-4 형식으로 나눈다. \textbf{환경 변수만 바꾸는 모델 배포}: curl 루프를 실행한 채 \code{kubectl set env MODEL\_URI=...v8/\allowbreak }을 적용하고, IMAGE 열이 그대로인 모습(터미널 기록 6-7), 실패 요청 0개, v8이 응답했음을 증명하는 서빙 응답을 캡처한다. 그런 다음 임시 접두사에서 \textit{의도적으로} 변경 불가능성을 깨뜨린다. 같은 URI 아래에 변경한 파일을 다시 업로드하고 플릿 절반을 재시작해 두 가지 진실이 생기는 증상을 문서화한다. 마지막으로 임시 접두사를 삭제하고 어긴 한 줄짜리 정책을 쓴다.

\subsection*{파트 C — 출입구와 실패 모드}

Ingress 애드온을 활성화하고(\code{minikube addons enable ingress}), 자체 서명 발급자를 사용해 cert-manager를 설치한 뒤, TLS를 포함한 터미널 기록 6-8의 Ingress를 \code{model-svc} 앞에 둔다. \code{curl -v}에 인증서가 보여야 한다. 이어서 현장 안내서의 세 행을 재현하고 각각 고친다. \textbf{(1) 대시보드는 초록색인데 503}: 준비 상태 경로 오타를 외부에서 내부 방향으로 진단해 터미널 기록 6-11의 연결 고리를 그대로 재현한다. \textbf{(2) 스트림 끊김}: 기본 60 s 프록시 타임아웃 뒤에서 실습 이미지의 \code{SLOW\_STREAM=90} 모드를 실행하고 정확히 60 s에서 두 번 끊기는 모습을 재현한다. 타임아웃 어노테이션으로 고친 뒤 스트림이 완료되는 모습을 보여 준다. \textbf{(3) 504}: 의도적으로 짧은 10 s 타임아웃을 통해 최악의 정당한 요청(최대 배치)을 보내 504를 확인한 뒤, 직접 측정한 최악의 경우에 여유를 더해 타임아웃을 설정한다. 경계는 \textbf{구축이 아닌 체크리스트}로 실습을 마무리한다. minikube 환경에서 실제로 배포할 §6.5 설계(공개·사설 분리, 참조로 연결한 SG 체인, 게이트웨이 엔드포인트, 버킷 접두사로 범위를 제한한 IRSA 역할)를 주석이 있는 한 페이지 그림으로 작성한다. 클라우드 크레딧이 있는 팀은 선택 부록에서 관리형 클러스터에 터미널 기록 6-10의 프로브 세 개를 재현한다.

\begin{calloutbox}{tbbpurple}{tbbpurplefill}
{\normalsize\bfseries\color{tbbpurple}실습 체크리스트}\par\vspace{3pt}
\begin{tbbitemize}
  \item 파트 A 산출물: 매니페스트와 v7/, v8/을 보여 주는 버킷 목록 · 로더 동시성 탐색 표 + 굴곡점 · 체크섬 검증 출력 · iostat 증거와 함께 병목을 밝힌 한 문장.
  \item 파트 B 산출물: 계산한 startupProbe 산술(YAML 내부 주석) · 자기 파드의 아무것도 소유하지 않는 점검 · 자기 숫자로 항목화한 콜드 스타트 막대 · 환경 변수만 바꾼 모델 배포 캡처(IMAGE 변화 없음, 실패 요청 0개) · 변경 불가능성 위반 보고서.
  \item 파트 C 산출물: curl -v TLS 핸드셰이크 · 503 포렌식 연결 고리(외부 → 포트 포워딩 → 엔드포인트) · 두 번의 60.0 s 스트림 끊김과 수정 결과 · 선택한 타임아웃과 근거를 포함한 504 수정 전후 · 경계 한 페이지 문서.
  \item 보고서: 한 페이지 + 산출물. 각 관찰에 관련 절을 명시한다(예: "정확히 60 s에서 끊겼다 — §6.4의 유휴 타임아웃, §6.6의 딱 떨어지는 시간 징후").
  \item 정리: kubectl delete namespace ml; minikube addons disable ingress; minikube stop. (MinIO 덕분에 이번 주 비용은 0이지만 정리 습관은 유지한다.)
\end{tbbitemize}
\end{calloutbox}

\vspace{3.0pt}

\begin{calloutbox}{tbbpurple}{tbbpurplefill}
{\normalsize\bfseries\color{tbbpurple}생각해 보기}\par\vspace{3pt}
\begin{tbbitemize}
  \item 환경 변수만 바꾼 배포는 실패 요청 0개로 v7 → v8을 롤아웃했다. 세 장에 걸쳐 협력해야 했던 모든 메커니즘을 열거한 뒤, 이를 깨뜨릴 가장 값싼 설정 실수 하나를 찾자. v7 크기에는 맞지만 v8에는 맞지 않는 startupProbe 예산을 생각해 보자.
  \item 실습 패턴을 4 GB에서 130 GB(11주 차 영역)로 확장하자. 배치, 매니페스트, 환경 변수 결합처럼 그대로 유지되는 부분, 프로브 예산, 동시성, 디스크처럼 다시 계산해야 하는 숫자, 노드 캐시와 사전 워밍업 풀처럼 필수가 되는 새 메커니즘을 구분하고 변경의 고통 순서대로 나열하자.
\end{tbbitemize}
\end{calloutbox}

\namedsection{요약}{열 가지로 정리한 이 장}

\qitem{1}{\textbf{모델 서버는 아무것도 소유하지 않는다.} 파드의 모든 바이트는 이미지에서 왔거나, 부팅할 때 저장소에서 가져왔거나, 버려도 되는 임시 데이터이며 절대 유일하지 않다. 이 원칙이 제5장의 모든 동작을 허가한다. 가치 있는 것이 파드와 함께 죽지 않기 때문에 파드를 가축처럼 다룰 수 있다.}

\qitem{2}{\textbf{스토리지는 세 제품이 아니라 세 계약이다.} 블록은 디스크(POSIX, 밀리초 미만, 소유자 하나, 프로비저닝 기준 과금)를 팔고, 객체는 HTTP 뒤의 내구성 높은 블롭(범위 GET, 제한 없는 읽기 주체, 일레븐 나인, 저장량 + 요청 + 외부 전송 과금)을 팔며, 파일은 블록의 약 4배 가격에 공유 POSIX를 판다. 접근 패턴, 공유, 가격 형태를 계산해 선택한다.}

\qitem{3}{\textbf{객체 스토리지의 형태는 "첫 바이트는 느리고 합계에는 상한이 없다"다.} 스트림 하나는 약 80 MB/s로 느리지만 16개의 범위 스트림은 NIC를 약 1 GB/s로 포화시킨다. 업계의 모든 빠른 로더는 이 관찰 하나를 적용한 결과다.}

\qitem{4}{\textbf{가중치는 소유한 바이트 가운데 가장 객체답다.} 매우 크고, 변경 불가능하며, 버전이 있고, 한 번 쓴 뒤 모든 레플리카가 순차적으로 읽는다. 하나의 버킷, 모델/버전/파일, sha256 매니페스트로 구성된 모델 스토어가 제3장 §3.4의 주의사항을 완전히 해결한다.}

\qitem{5}{\textbf{버전은 변경할 수 없다. 새 버전 = 새 접두사이며 절대 덮어쓰지 않는다.} 롤링 업데이트, 노드 캐시, 롤백이 모두 조용히 이 원칙에 의존하며 IAM으로 쓰기 주체는 생성만 하고 덮어쓰지 못하게 강제할 수 있다.}

\qitem{6}{\textbf{출시 주기를 분리하면 모델 배포가 소프트웨어 출시이기를 멈춘다.} 코드는 이미지에, 가중치는 스토어에 두고 환경 변수만 결합 지점으로 삼는다. kubectl set env MODEL\_URI=...v8은 IMAGE 열을 바꾸지 않은 채 여분 레플리카, 준비 상태 속도 조절, 롤백 같은 제5장의 전체 메커니즘으로 새 모델을 롤아웃한다. 13주 차에는 이 한 줄을 자동화한다.}

\qitem{7}{\textbf{콜드 스타트는 하나의 숫자가 아니라 합이다.} 스케줄링 + 이미지 가져오기 + 가중치 다운로드 + 로드 + 워밍업이다. 같은 하드웨어에서 순진한 방식은 약 260 s, 설계한 방식(캐시된 이미지, 16방향 로더, mmap)은 약 56 s다. 각 항목에는 책임자가 있고, 12주 차 제로 스케일링 논쟁은 이 막대를 두고 벌어진다.}

\qitem{8}{\textbf{출입구에는 계층이 있고 기본 설정은 서빙에 적대적이다.} L4는 연결을, L7은 요청을 로드 밸런싱해 제5장의 gRPC 주의사항을 해결한다. Ingress는 하나의 LB로 여러 Service를 제공하고 API 게이트웨이는 키, 할당량, 429, 사용량 계측 같은 제품 기능을 더한다. 요청 타임아웃은 정직하게 측정한 최악 p99보다 길게, 유휴 타임아웃은 스트리밍 토큰 간격보다 길게 잡고 재시도에는 예산을 둔다. LB 연결 비우기와 terminationGracePeriodSeconds를 맞추지 않으면 배포마다 502를 산다.}

\qitem{9}{\textbf{경계는 성실함이 아니라 구조로 만든다.} 사설 서브넷에는 인터넷에서 들어오는 경로가 없고 LB만 공개한다. 보안 그룹은 참조로 연결해 노드가 LB-SG에서 오는 트래픽을 허용하게 한다. 파드는 키가 아니라 워크로드 아이덴티티로 인증하고, 가중치는 GB당 \$0.045인 NAT 우회로가 아니라 게이트웨이 엔드포인트를 사용한다. 중간 규모 플릿 하나에서 연간 약 \$5,000 차이다.}

\qitem{10}{\textbf{상태 코드가 실패 위치를 좁힌다.} 502 = 파드가 연결을 끊음(OOM, 연결 비우기 불일치), 503 = 응답할 대상 없음(엔드포인트 또는 두 계층 상태 확인 함정), 504 = 출입구 기준으로 너무 느림, 정확히 N초에 끊김 = 유휴 타임아웃, 429 = 정상 작동하는 계량기다. 항상 첫 동작은 포트 포워딩으로 다시 요청하는 것이다. 내부 정상 + 외부 고장 = 출입구 문제다.}

\namedsection{연습문제}{}

\subsection*{A. 객관식}

\qitem{1}{서빙 파드 안에서 발견된 바이트 가운데 아무것도 소유하지 않는 원칙을 위반하는 것은?}

\begin{tbboptions}
  \item ① emptyDir 안의 JIT 컴파일 커널 캐시
  \item ② 부팅할 때 initContainer가 다운로드한 가중치
  \item ③ 나중에 가져오기 위해 컨테이너 파일시스템에 저장한 사용자 업로드 문서
  \item ④ 이미지에서 온 Python 인터프리터
\end{tbboptions}

\qitem{2}{PostgreSQL 데이터 디렉터리를 객체 스토리지가 아니라 블록 스토리지에 두는 주된 이유는?}

\begin{tbboptions}
  \item ① 객체 스토리지의 내구성이 더 낮다.
  \item ② 데이터베이스에는 밀리초 미만 지연 시간의 임의 제자리 쓰기가 필요하지만 객체 스토리지는 그 의미 체계를 판매하지 않는다.
  \item ③ 객체 스토리지는 그렇게 큰 파일을 담을 수 없다.
  \item ④ 블록 스토리지가 GB당 더 저렴하다.
\end{tbboptions}

\qitem{3}{13 GB 모델을 스트림 하나로 다운로드하면 162 s, 16개의 범위 스트림으로 다운로드하면 17 s가 걸린다. 활용한 객체 스토리지의 성질은?}

\begin{tbboptions}
  \item ① PUT 뒤의 강한 일관성
  \item ② 높은 첫 바이트 도착 시간
  \item ③ 동시 범위 GET에 따라 클라이언트 NIC 한계까지 확장되는 총 병렬 처리량
  \item ④ 비정기 접근 클래스로 옮기는 수명 주기 전환
\end{tbboptions}

\qitem{4}{모델 스토어 버전을 변경할 수 없게 해야 하는 이유는 무엇인가(새 버전 = 새 접두사)?}

\begin{tbboptions}
  \item ① S3는 객체를 덮어쓸 수 없다.
  \item ② 롤아웃 도중 "같은 버전"의 기존 레플리카와 새 레플리카가 달라질 수 있고, URI를 키로 사용하는 노드 캐시가 조용히 오염되며, 롤백 대상이 바뀔 수 있다.
  \item ③ 덮어쓰기가 새 접두사를 만드는 것보다 느리다.
  \item ④ 매니페스트는 새 접두사에서만 동작한다.
\end{tbboptions}

\qitem{5}{클라이언트가 스트리밍 엔드포인트가 매번 정확히 60초 뒤 끊긴다고 보고했다. 가장 가능성 높은 원인은?}

\begin{tbboptions}
  \item ① 파드가 60 s에 OOM으로 종료된다.
  \item ② GPU에서 타임아웃이 발생한다.
  \item ③ 긴 토큰 간격 동안 바이트가 움직이지 않은 연결을 로드 밸런서의 유휴 타임아웃이 정리한다.
  \item ④ DNS TTL 만료
\end{tbboptions}

\qitem{6}{사설 서브넷의 플릿이 매일 S3에서 가중치를 가져온다. 이 트래픽에 붙는 NAT 게이트웨이의 GB당 비용을 없애는 변경 하나는?}

\begin{tbboptions}
  \item ① 노드를 공개 서브넷으로 옮긴다.
  \item ② 가중치를 압축한다.
  \item ③ 사설 서브넷의 라우팅 테이블에 S3 게이트웨이 VPC 엔드포인트를 추가한다.
  \item ④ 버킷을 비정기 접근 스토리지 클래스로 바꾼다.
\end{tbboptions}

\subsection*{B. 단답형}

\qitem{1}{서빙 파드 바이트의 합법적인 출처 세 가지를 말하고, 이 장의 터미널 기록에서 각각 구체적인 예를 하나씩 들라.}

\qitem{2}{콜드 스타트 예산을 계산하라. 이미지 6 GB를 150 MB/s로 콜드 풀하고, 가중치 26 GB를 스트림 하나(80 MB/s)와 스트림 16개(800 MB/s)로 각각 가져오며, 로드 30 s, 워밍업 15 s, 스케줄링 2 s가 걸린다. 두 합계를 구하고 가장 큰 지렛대 하나를 말하라.}

\qitem{3}{L4와 L7을 두 문장으로 비교한 뒤, L7 홉만 고칠 수 있는 제5장의 주의사항과 이를 필수로 만드는 11주 차 워크로드를 말하라.}

\qitem{4}{롤아웃 중 약 20초 동안 502가 급증한다. 서로 어긋난 두 시계, 해결책이 작동하려면 서버가 구현해야 하는 제3장의 메커니즘, 연결 비우기를 깔끔하게 만드는 준비 상태 우선 순서를 말하라.}

\subsection*{C. 논술형}

\qitem{1}{"가중치를 이미지에 구워 넣기"와 "부팅할 때 모델 스토어에서 로딩하기"의 양쪽을 솔직하게 논하라. 구워 넣기가 이기는 엣지·망 분리 사례, 스토어가 이기는 출시 주기와 레지스트리 외부 전송 논거, 두 방식 모두가 마주하는 콜드 스타트 계산을 포함한다. 강의 프로젝트에 적용할 방식을 결정하고 가장 강한 반론에 맞서 옹호하라.}

\qitem{2}{새 공개 모델 API의 전체 경로를 설계하라. 스토어 배치와 IAM 쓰기·읽기 정책, 로더와 프로브 예산, TLS와 할당량을 갖춘 Ingress·게이트웨이 토폴로지, VPC·SG·엔드포인트 배치를 포함하고, 각 홉에 §6.6의 실패 모드와 이 장 계산 예시의 비용 항목을 주석으로 단다. 한 페이지 산문 또는 주석이 있는 그림으로 작성하라.}

\qitem{3}{"객체 스토리지가 머신러닝을 이겼다."를 평가하라. 계약의 변경 불가능성, 제한 없는 읽기 주체, 병렬 처리량, 가격 가운데 무엇이 ML 산출물에 잘 맞으며, 그 결과 생기는 첫 바이트 지연 시간, 콜드 스타트, 제자리 갱신 불가 비용을 나머지 서빙 아키텍처가 어떻게 흡수해야 하는가? 캐싱 단계, 콜드 스타트 청구서, 11주 차의 더 큰 모델을 근거로 사용하라.}

\namedsection{정답과 해설}{}

\subsection*{A. 객관식}

\qitem{1}{\textbf{③} 업로드는 다른 곳에 존재하지 않는다. 파드는 가축이기를 멈췄고 축출, 롤아웃, 재스케줄링 같은 제5장의 모든 메커니즘이 데이터 손실 경로가 되었다. ①, ②, ④는 임시 데이터, 가져온 데이터, 이미지라는 세 가지 합법적 출처다. (§6.1)}

\qitem{2}{\textbf{②} 내구성이나 크기에 관한 속설이 아니라 계약 조항 때문이다. B-tree는 디스크 지연 시간으로 페이지를 제자리에서 다시 쓰지만 객체 스토리지는 HTTP로 전체 객체 PUT/GET을 판매한다. 변경 불가능한 Parquet를 분석하는 반대 사례가 규칙을 입증한다. (§6.2)}

\qitem{3}{\textbf{③} 첫 바이트는 느리고 합계에는 상한이 없다는 객체 계약의 형태다. 범위 GET은 클라이언트 NIC 또는 디스크가 포화할 때까지 스트림을 늘린다. 10 Gbps ≈ 1.25 GB/s가 상한이다. (§6.2–6.3)}

\qitem{4}{\textbf{②} 세 실패는 모두 실제 사고다. 롤아웃 도중 두 가지 진실을 가진 플릿, 오염된 캐시, 변경된 롤백 대상이다. ①은 틀렸다. S3는 기꺼이 덮어쓰므로 규율과 IAM으로 금지해야 한다. (§6.3)}

\qitem{5}{\textbf{③} 딱 떨어지는 시간이 징후다. 유휴 타임아웃은 정확히 N초 동안 바이트가 흐르지 않은 연결을 정리한다. 최악의 토큰 간격보다 유휴 타임아웃을 높이거나 연결 유지 프레임을 보낸다. (§6.4, §6.6)}

\qitem{6}{\textbf{③} 게이트웨이 엔드포인트는 S3 트래픽을 \$0 데이터 처리 비용으로 제공자 백본에 전달하는 라우팅 테이블 항목이다. 연간 약 \$5,000 비용 항목을 고친다. ④는 전송이 아니라 저장 비용을 바꾼다. (§6.5)}

\subsection*{B. 단답형}

\qitem{1}{(a) 이미지에서 온 것 — /usr, Python 런타임(오버레이 마운트, 터미널 기록 6-1), (b) 부팅할 때 가져온 것 — initContainer를 통해 /models에 들어온 13 GB safetensors, (c) 버려도 되는 임시 데이터 — /dev/shm, emptyDir의 토크나이저·JIT 캐시다. 이 세 가지 밖의 것은 설계 버그다. (§6.1)}

\qitem{2}{이미지: 6 GB / 150 MB/s = 40 s다. 가중치: 26 GB / 80 MB/s = 325 s 대 26 GB / 800 MB/s ≈ 33 s다. 순진한 방식은 2 + 40 + 325 + 30 + 15 = \textbf{412 s}, 설계한 방식은 2 + 40 + 33 + 30 + 15 = \textbf{120 s}이며 노드에 캐시한 이미지라면 약 80 s다. 가장 큰 지렛대는 가중치 다운로드 병렬화(−292 s)다. (§6.3)}

\qitem{3}{L4는 TCP 연결을 전달하며 HTTP를 읽지 않는다. L7은 연결을 종료하고 개별 요청을 로드 밸런싱한다. Service는 연결을 로드 밸런싱하므로 다중화한 gRPC가 파드 하나에 고정된다는 주의사항이 있다. 요청 수준(L7) 밸런싱만 이를 고치며 11주 차 스트리밍 LLM 트래픽에서 필수가 된다. (§6.4)}

\qitem{4}{LB의 등록 해제 지연과 파드의 terminationGracePeriodSeconds가 서로 어긋났다. 서버는 SIGTERM에서 실제로 연결을 비워야 한다(제3장의 PID 1 핸들러, exec 형식 CMD). 깔끔한 순서는 준비 상태가 먼저 실패해 엔드포인트 제거로 새 트래픽을 막고 → 처리 중 요청을 비운 뒤 → 유예 기간 안에 종료하는 것이다. (§6.4, §6.6)}

\subsection*{C. 논술형(채점 기준)}

\qitem{1}{구워 넣는 쪽은 스토어에 접근할 수 없는 망 분리·엣지 환경, 단일 산출물의 감사 가능성, 스토어 가용성 SLA에 대한 부팅 시점 의존성 제거를 들어야 한다. 스토어 쪽은 출시 주기 분리(빌드 없는 모델 배포 — 터미널 기록 6-7), 레지스트리 외부 전송과 노드 디스크 복제, 이미지 태그의 버전 폭증, 제3장의 레이어 논거를 들어야 한다. 두 방식 모두 같은 콜드 스타트 합을 마주한다. 구워 넣기는 바이트를 이미지 가져오기 항목으로 옮길 뿐 없애지 않는다. 좋은 답은 이 대칭을 다루고 프로젝트의 입장을 정한다. 예를 들어 스토어를 선택하고 극단 사례를 예외 정책으로 명시한다. (§6.1, §6.3)}

\qitem{2}{다음 요소를 확인한다. 버전이 있는 접두사 + 매니페스트, 쓰기 주체는 생성만 하고 덮어쓰지 못하게 하는 IAM, 접두사별 범위가 정해진 읽기 주체, 측정한 다운로드+로드로 계산한 프로브 예산, 하나의 LB → Ingress → Service와 cert-manager의 TLS, 키·할당량용 게이트웨이, 사설 노드, 참조로 연결한 SG 체인, 첫날부터 둔 게이트웨이 엔드포인트다. 홉별 실패 주석은 503 엔드포인트, 504 타임아웃, N초 유휴 끊김, 502 연결 비우기이며 비용 항목은 몇 센트의 스토리지, 피한 NAT 세금, LB 시간당 비용, 거의 무시할 응답 외부 전송이다. 보너스는 LB·준비 상태·게이트웨이 전체에서 상태 확인의 진실 원천 하나를 쓰는 것이다. (§6.3–6.6)}

\qitem{3}{좋은 답은 계약과 산출물을 연결한다. ML 산출물은 크고, 변경 불가능하며, 버전이 있고, 여러 읽기 주체가 읽는다. 객체 스토리지가 월 GB당 약 \$0.023에 파는 정확한 약속이다. 부과되는 비용도 실제다. 수십 ms의 첫 바이트 지연과 제자리 갱신 불가는 로더(병렬 범위 GET), 캐시(단계), 프로브 예산, 12주 차 사전 워밍업으로 복잡성을 밀어낸다. 즉 스토어가 아키텍처에 맞추기보다 아키텍처가 스토어에 맞춰 휜다. 11주 차 130 GB 이상 확장은 양쪽을 더 선명하게 만든다. 스토어는 개의치 않고 확장하지만 하위 구성 요소의 예산, 캐시, NIC는 모두 다시 계산해야 한다. 이 휨의 비용을 계산하면 어느 결론도 점수를 받을 수 있다. (§6.2–6.3, §6.7)}

\namedsection{더 읽을거리}{}

\textbf{[이번 주]} 표시가 붙은 항목은 이번 주 필독 자료다. \textbf{[N주 차]} 표시는 해당 주에 도움이 될 자료를 뜻한다.

\begin{tbbitemize}
  \item \textbf{[이번 주] AWS S3 documentation: "Performance Guidelines" and "Optimizing Amazon S3 Performance"} — §6.3 로더의 1차 자료다. 접두사별 요청 병렬 처리, 범위 GET, 멀티파트(multipart) 임계값을 다룬다. 짧고 구체적이며 같은 아이디어가 GCS에도 적용된다.
  \item \textbf{[이번 주] safetensors documentation (Hugging Face)} — 형식의 헤더 배치, 설계에서 mmap 가능성이 자연스럽게 나오는 이유, 7주 차에 더 강하게 반복할 pickle 반대 보안 논거를 다룬다.
  \item \textbf{[이번 주] Kubernetes docs: "Ingress", "Init Containers", and "Volumes" (emptyDir)} — 이 장의 매니페스트를 구성하는 세 개념 페이지다. 타임아웃은 어노테이션에 있으므로 Ingress 문서를 컨트롤러의 어노테이션 레퍼런스와 함께 읽자.
  \item \textbf{AWS docs: "VPC endpoints for Amazon S3" and "NAT gateways" (pricing page included)} — 15분을 투자해 중간 규모 플릿에서 연간 약 \$5,000를 되찾는 §6.5 계산 예시의 공식 자료다.
  \item \textbf{Google SRE Book, ch. 19–20 ("Load Balancing at the Frontend" / "in the Datacenter")} — 대규모 L4/L7, 상태 확인, 연결 비우기의 정석이다. §6.4는 이 장들을 서빙 규모로 줄인 버전이다.
  \item \textbf{Kleppmann, "Designing Data-Intensive Applications", ch. 3} — 스토리지 엔진 장이다. 블록 기반 엔진이 제자리 쓰기로 실제로 무엇을 하는지, 즉 §6.2 계약표의 반대쪽을 보여 준다.
  \item \textbf{[7주 차] ONNX / TorchScript / safetensors format comparisons (framework docs)} — 다음 주 모델 형식 논의는 형식이 이 장의 스토어에 살고 변환된다는 사실을 전제로 한다.
  \item \textbf{[12주 차] Cluster Autoscaler FAQ + cloud provider docs on image streaming/pre-pulling} — 콜드 스타트 막대에 남은 설정인 사전 워밍업 풀, 지연 이미지 마운트, 제로 스케일링 경제학을 다룬다.
  \item \textbf{[13주 차] MLflow Model Registry (or equivalent) documentation} — 이 장의 스토어를 가리키는 포인터의 데이터베이스인 레지스트리를 다룬다. 단계, 승격, 계보가 추가되어도 메커니즘은 터미널 기록 6-7 그대로다.
\end{tbbitemize}

\vspace{5.0pt}

\begin{calloutbox}{tbbblue}{tbbbluefill}
{\normalsize\bfseries\color{tbbblue}다음 장 — 제7장: 서빙 시스템의 본격적인 시작}\par\vspace{3pt}
이 책의 인프라 절반인 커널, 컨테이너, 이미지, 오케스트레이션, 스토리지, 출입구가 완성되었다. 파트 2에서는 스택 위로 올라가 그곳에 머문다. 제7장은 서빙 문제 자체로 시작한다. \textbf{온라인 추론, 배치 추론, 스트리밍 추론}이 각각 제값을 하는 곳, 구호가 아니라 대기열로 설명하는 \textbf{지연 시간–처리량 절충 관계}, 이후 모든 벤치마크가 사용하는 \textbf{SLO와 백분위수(percentile) 어휘}(p95/p99), 이제 형식이 사는 위치(§6.3의 스토어)와 로딩 주체(이 장의 로더)를 알게 되었으므로 다룰 수 있는 \textbf{모델 형식} ONNX, TorchScript, safetensors와 서로 간의 변환 경로를 살펴본다. 팀 프로젝트 제안 주간이기도 하다. 7주 차에 약속하는 SLO는 15주 차 데모를 평가하는 숫자가 된다. 책상 위에 제1장의 삼각형을 펼쳐 놓고 선택하자.
\end{calloutbox}


\part{AI 서빙 시스템}
%% Chapter 7 — 서빙 시스템 입문
\setcounter{chapter}{6}
\chapter{서빙 시스템 입문}
\chaptersubtitle{파트 2의 시작 — 온라인, 배치, 스트리밍, 그리고 빠름의 언어: 꼬리, 약속, 대기열, 형식}

\begin{calloutbox}{tbbblue}{tbbbluefill}
{\normalsize\bfseries\color{tbbblue}학습 목표}\par\vspace{3pt}
\begin{tbbitemize}
  \item 팀 채팅의 논쟁 하나를 바탕으로, SLO가 없는 서비스에서는 모든 불만이 장애가 되고 조용한 날은 모두 우연이 되는 이유와 "평균은 182 ms"라는 말이 실제로는 아무도 설명하지 못하는 이유를 진단한다.
  \item "누가 무엇을 기다리는가?"를 물어 워크로드를 \textbf{온라인, 배치, 스트리밍} 추론으로 분류하고, 각 방식이 무엇을 최적화하고 무엇을 감수하며 무엇을 어려워하는지 설명한다.
  \item 지연 시간을 분포로 읽고 보고한다. 백분위수(p50/p95/p99), 팬아웃에서 꼬리가 증폭되는 이유(규모에서의 꼬리), 순진한 부하 테스트 도구가 거짓말하는 방식(협조적 누락)을 이해한다.
  \item SLI, 목표값, 평가 기간, 부하 조건, 비용 상한을 포함한 완전한 SLO를 작성하고 오류 예산을 도출하며, 이를 소망이 아니라 출시·동결 정책으로 사용한다.
  \item \textbf{Little's Law}(L = λ·W)와 사용률 하키 스틱으로 플릿 규모를 산정한다. 측정한 레플리카별 용량으로 레플리카 수를 계산하고, 의도적으로 여유분을 더하며, 결과의 비용을 정직하게 산정한다.
  \item 모델 산출물을 선택하고 변환한다. pickle의 로드 시 코드 실행 위험과 safetensors, 즉시 실행 방식과 TorchScript, ONNX, 컴파일 엔진을 비교하고, 변환이 시작 단계가 아니라 반드시 CI 단계여야 하는 이유를 설명한다.
\end{tbbitemize}
\end{calloutbox}

\vspace{3.5pt}

\begin{calloutbox}{tbbteal}{tbbtealfill}
{\normalsize\bfseries\color{tbbteal}핵심 용어}\par\vspace{3pt}
온라인 / 배치 / 스트리밍 추론 · 지연 시간 백분위수(p50 · p95 · p99) · 규모에서의 꼬리 · 협조적 누락 · SLI · SLO · SLA · 오류 예산 · 리틀의 법칙(L = λW) · 사용률 ρ · M/M/1 · 굴곡점 · 여유분 · 서빙 가능한 산출물 · safetensors · TorchScript / torch.export · ONNX · TensorRT 엔진 · GGUF
\end{calloutbox}

\vspace{3.5pt}

\begin{calloutbox}{tbblightgray}{tbbboxfill}
{\normalsize\bfseries\color{tbbink}장 구성}\par\vspace{3pt}
\textbf{7.1} 논쟁과 아무도 경험하지 않은 숫자    \textbf{7.2} 세 가지 추론 방식    \textbf{7.3} 지연 시간은 분포다 — 꼬리 읽기    \textbf{7.4} SLI, SLO, SLA, 오류 예산    \textbf{7.5} 대기열: 리틀의 법칙과 하키 스틱    \textbf{7.6} 모델 형식과 변환 경로    \textbf{7.7} 실습: 프로젝트 제안서 — 이어서 요약, 연습문제, 더 읽을거리.
\end{calloutbox}

\section{논쟁과 아무도 경험하지 않은 숫자}

파트 1은 현실적인 결과물로 끝났다. 컨테이너 안의 모델 서버를 Kubernetes가 스케줄링하고, 객체 스토어에서 가중치를 로딩하며, 잠긴 네트워크 안에서 로드 밸런서를 통해 접근할 수 있다. 시스템은 실행된다. 파트 2는 앞으로 6주를 관통할 질문으로 시작한다. \textbf{좋은 시스템인가?} 그런데 오늘 아침 현재 팀에는 "좋다"를 나타낼 숫자가 없다는 불편한 사실이 드러난다. 있는 것은 의견뿐이다. 프로덕션에서 의견이 어떤 모습인지 살펴보자. 아래 기록은 여러 사례를 합친 것이지만 이 장을 건너뛴 모든 엔지니어링 팀이 겪어 본 일이다.

\subsection*{사례 A: 논쟁}

\begin{terminalbox}
\begin{Verbatim}[breaklines=true,breakanywhere=true,fontsize=\small,formatcom=\color{tbbtermfg},codes=\tbbcodeactive]
[tue 10:12] pm     users are saying the demo "feels slow"
[tue 10:14] eng    avg latency is 182ms. that's fast. dashboard is green
[tue 10:31] pm     the CEO tried it in the all-hands. it spun for 4 seconds
[tue 10:33] eng    works fine for me?? just tried it. instant
[tue 11:02] cs     two pilot customers asked if "slow is normal for AI"
[tue 11:05] eng    define slow. compared to what? at what load?
[tue 11:06] pm     ...
[tue 11:06] eng    ...
[wed 09:00] cto    what did we actually PROMISE them?
[wed 09:04] pm     nothing. we never promised anything.
[wed 09:05] cto    then everything is a breach, and nothing is. fix that.
\end{Verbatim}
\end{terminalbox}
\figcaption{사례 7-A  논쟁. 모든 발언이 사실이라는 점이 바로 문제다. 합의된 숫자가 없으면 모든 불만은 장애가 되고 조용한 날은 모두 우연이 된다.}

다음으로 넘어가기 전에 부검해 보자. 실패가 미묘하기 때문이다. \textbf{기록 속 누구도 틀리지 않았다.} 엔지니어가 말한 평균 182 ms는 사실이고, CEO의 4초도 사실이며, "내 환경에서는 잘 된다"는 말도 사실이다. 모두 같은 서비스를 보면서 서로 다른 진실을 보고한다. 서비스의 속도는 하나의 숫자가 아니라 \textbf{분포}이고, 각자가 그중 다른 부분을 표본으로 뽑았기 때문이다. PM은 스스로 선택해 불만을 보낸 꼬리 사용자들을 표본으로 삼았다. 엔지니어는 곧 보겠지만 아무도 설명하지 못하는 요약인 평균을 표본으로 삼았다. CEO는 목격자들 앞의 실제 상황에서 한 번 표본을 뽑았고, 꼬리가 두꺼운 분포의 단일 표본은 복권과 같다. 수요일 CTO의 말은 §7.4 전체를 한 문장으로 요약한다. 목표를 명시하지 않으면 정상과 고장을 구별할 수 없다. \textit{"그러면 모든 것이 위반이면서 아무것도 위반이 아니다."}

\subsection*{사례 B: 아무도 경험하지 않은 숫자}

이제 같은 서비스가 의견의 대상에서 벗어나는 모습을 보자. 올바르게 실행한 부하 테스트 한 번 뒤에 팀이 다음 보고서를 얻었다. "올바르게"의 까다로운 조건은 §7.3에서 다룬다.

\begin{terminalbox}
\begin{Verbatim}[breaklines=true,breakanywhere=true,fontsize=\small,formatcom=\color{tbbtermfg},codes=\tbbcodeactive]
$ python report.py results.json     # 10,000 requests, tuesday 10:00-10:30
requests      10,000     errors  0.4%  (23x 504, 17x 502)
mean          182 ms     <- the dashboard number
p50            92 ms     <- half of all users: this or better
p90           220 ms
p95           340 ms
p99         2,140 ms     <- 100 users today: worse than this
p99.9       4,810 ms     <- 10 users: ~5 seconds. hi, CEO
max         9,400 ms
\end{Verbatim}
\end{terminalbox}
\figcaption{사례 7-B  사례 7-A와 같은 서비스를 분포로 본 모습. 평균 182 ms는 p50과 p90 사이에 있지만 실제 사용자 누구도 경험하지 않은 요약 통계다. 설명하지도 못하는 꼬리가 평균을 위로 끌어올렸다.}

논쟁을 번역한 표로 읽어 보자. \textit{"평균은 182 ms다"}는 사실이지만 쓸모없다. 오른쪽으로 치우친 분포의 평균은 사용자가 아무도 없는 곳에 놓인 허구다. 절반은 92 ms보다 빨랐고 운 나쁜 100명은 2.1 s보다 느렸다. \textit{"내 환경에서는 잘 된다"}는 당연하다. p50에서는 단일 표본 대부분이 잘 나온다. \textit{"CEO는 4초 동안 기다렸다"}는 p99.9 너머를 뽑은 것이다. 요청 10,000건에서는 실제 사람 10명이고 하루 요청 200,000건에서는 매일 200명이다. \textit{"고객 두 명이 느린 것이 AI에서는 정상인지 물었다"}는 꼬리가 극단 사례가 아니라 \textbf{매일 찾아오는 사용자 집단}임을 보여 준다. 그래서 서빙 시스템은 평균이 아니라 백분위수로 말한다. 이 강의의 나머지 모든 벤치마크, SLO, 대시보드가 p50/p95/p99를 보고하는 이유이며, 이 장에서 그 규율을 기초부터 세운다.

\subsection*{속도가 곧 매출이다: 고전적인 실험}

\begin{calloutbox}{tbbamber}{tbbamberfill}
{\normalsize\bfseries\color{tbbamberdark}모든 서빙 엔지니어가 기억해 둘 세 가지 실험}\par\vspace{3pt}
\begin{tbbitemize}
  \item \textbf{Amazon, 2006 (Greg Linden):} 페이지 로딩에 인위적으로 +100 ms를 더하자 매출이 측정 가능하게 감소했다. 널리 알려진 수치는 \textbf{100 ms당 매출 −1\%}다. 전자상거래 업계가 밀리초를 두려워하게 만든 실험이다.
  \item \textbf{Google, 2009 ("Speed Matters," Jake Brutlag):} 검색 결과를 +400 ms 늦추자 사용자당 검색이 \textbf{−0.59\%} 줄었고, 지연을 제거한 뒤에도 효과가 \textit{계속되었다.} 느린 경험이 사용자에게 검색을 덜 하도록 가르쳤다.
  \item \textbf{Google, 2006 (Marissa Mayer at Web 2.0):} 사용자는 페이지마다 결과 30개를 원했지만 제공하는 데 \textasciitilde{}+500 ms가 더 걸렸고 트래픽은 \textbf{\textasciitilde{}20\%} 감소했다. 사용자는 적더라도 빠른 것을 선호한다.
\end{tbbitemize}

한계도 솔직하게 밝히자. 2000년대 웹에서 나온 숫자라 상황에 따라 결과는 달라진다. 하지만 같은 방향의 결과가 20년간 반복되었고 핵심은 비대칭이다. 사용자는 지연이 없을 때보다 지연이 있을 때 훨씬 더 잘 알아챈다. 지연 시간은 엔지니어의 허영을 위한 지표가 아니라 비용 항목이다.
\end{calloutbox}

\vspace{3.0pt}

\subsection*{경첩: 소유한 것과 책임질 것}

그림 7-1은 이 장의 위치를 보여 준다. 선 아래는 파트 1인 플랫폼이다. 실제 본질을 조사할 수 있는 임대 컴퓨팅(제2장), 네임스페이스와 cgroup으로 만든 컨테이너(제3–4장), 조정하는 스케줄러(제5장), 상태·출입구·경계(제6장)가 있다. 이제 이 모든 것은 \textit{전제}다. 이번 주부터 "배포하라"는 이미 실행할 줄 아는 문장이다. 선 위는 파트 2인 서빙 시스템이다. 같은 배포를 이제 제품으로 평가한다. 요청 하나가 클라이언트에서 들어오며 유일하게 중요한 시계가 이때 시작된다. 6주 차에 만든 게이트웨이를 지나 대기열에서 기다리고, 배치에 들어가며, 모델 형식 하나를 실행하는 런타임을 통과한 뒤 돌아간다. 이 장의 모든 절은 그 경로에 표시된 조절점 하나를 다룬다.

\begin{tbbfigure}
  \adjustbox{max width=0.94\linewidth,max totalheight=0.46\textheight}{\includegraphics{figures/fig7-1_hinge.png}}\\[3pt]
  \figcaption{그림 7-1  강의의 경첩. 파트 1은 선 아래의 플랫폼을 만들었고 파트 2는 선 위의 서빙 시스템을 조정한다. 이 장의 각 절은 요청 경로에 방식(§7.2), 꼬리(§7.3), 약속(§7.4), 플릿(§7.5), 형식(§7.6), 제안서(§7.7)를 표시한다.}
\end{tbbfigure}

이번 주는 팀 프로젝트의 첫 제출물인 제안서 주간이기도 하다(§7.7). 행정적인 우연이 아니다. 제안서는 사례 7-A의 팀이 답하지 못한 질문에 숫자로 답하는 한 페이지 문서다. 무엇을 서빙하고(§7.6의 산출물), 어느 방식으로(§7.2), 무엇을 약속하며(§7.4의 SLO), 어느 부하에서 어떤 플릿(§7.5)을 사용하고, 월 비용 상한은 얼마인가(제1장의 비용 관점)? 여기서 §7.7까지의 모든 내용이 이 빈칸을 채울 재료다. 이 장 자체가 말 그대로 과제다.

\subsection*{이 장의 어휘}

\begin{tbbtablewrap}
  \begin{tabular}{@{}>{\raggedright\arraybackslash}m{\dimexpr 0.2660\tbltextwidth-2\tabcolsep\relax} >{\raggedright\arraybackslash}m{\dimexpr 0.7340\tbltextwidth-2\tabcolsep\relax}@{}}
  \arrayrulecolor{tbbnavy}\specialrule{1pt}{0pt}{0pt}
  \rowcolor{tbbtablehead}{\centering \bfseries\color{tbbnavy}용어} & {\centering \bfseries\color{tbbnavy}의미} \\
  \arrayrulecolor{tbbnavy}\specialrule{0.7pt}{0pt}{2pt}
  {\textbf{온라인 / 배치 / 스트리밍}} & {"누가 기다리는가"에 따라 나눈 세 가지 추론 방식. 요청마다 한 사람 / 아무도 없음 / 결과가 흐르는 모습을 지켜보는 사람으로 구분한다(§7.2).} \\
  \arrayrulecolor{tbbtableborder}\hline
  {\textbf{백분위수(pN)}} & {요청의 N\%가 이 값 이하로 빨랐음을 나타내는 지연 시간. p99는 가장 운 나쁜 1\%의 경험이며, 매일 나타나는 실제 사용자 집단이다.} \\
  \arrayrulecolor{tbbtableborder}\hline
  {\textbf{규모에서의 꼬리}} & {Dean과 Barroso의 관찰. 팬아웃에서는 드문 느림이 흔해진다. 1\% 느린 말단 100곳으로 팬아웃하면 요청의 63\%가 느려진다(§7.3).} \\
  \arrayrulecolor{tbbtableborder}\hline
  {\textbf{협조적 누락}} & {대표적인 부하 테스트 버그. 폐쇄 루프(closed loop) 시험기가 서버가 멈춘 동안 요청 전송도 멈춘 뒤, 멈춤을 제외한 백분위수를 보고한다(§7.3).} \\
  \arrayrulecolor{tbbtableborder}\hline
  {\textbf{SLI}} & {서비스 수준 지표(Service Level Indicator). "게이트웨이에서 500 ms 미만인 요청의 비율"처럼 신중하게 정의한 측정값으로, 현재 무슨 일이 일어나는지 나타낸다(§7.4).} \\
  \arrayrulecolor{tbbtableborder}\hline
  {\textbf{SLO}} & {서비스 수준 목표(Service Level Objective). 명시한 부하와 기간에 대해 스스로 정한 목표로, 우리 자신에게 하는 약속이다(§7.4).} \\
  \arrayrulecolor{tbbtableborder}\hline
  {\textbf{SLA}} & {서비스 수준 협약(Service Level Agreement). 비용이 걸린 외부 계약이며 항상 SLO보다 느슨하다(§7.4).} \\
  \arrayrulecolor{tbbtableborder}\hline
  {\textbf{오류 예산(error budget)}} & {1 − SLO 목표값으로 계산한, 쓸 수 있는 신뢰성 저하분. 30일 동안 99.9\%면 \textasciitilde{}43분이다. 성적표가 아니라 출시·동결 정책이다(§7.4).} \\
  \arrayrulecolor{tbbtableborder}\hline
  {\textbf{λ, W, L}} & {도착률(req/s), 시스템에 머무는 시간(s), 처리 중인 수. 안정적인 모든 시스템에서 Little's Law인 L = λ·W로 연결된다(§7.5).} \\
  \arrayrulecolor{tbbtableborder}\hline
  {\textbf{사용률 ρ · 여유분}} & {제공 부하 ÷ 용량. ρ → 1이면 지연 시간이 폭발하므로(하키 스틱) 온라인 플릿은 의도적으로 여유분을 구매한다(§7.5).} \\
  \arrayrulecolor{tbbtableborder}\hline
  {\textbf{서빙 가능한 산출물}} & {실제로 배포하는 가중치(safetensors), 구성 파일, 토크나이저다. 모델 스토어(제6장)에서 버전을 관리하며 런타임이 실행할 수 있는 형식이다(§7.6).} \\
  \arrayrulecolor{tbbnavy}\specialrule{1pt}{2pt}{0pt}
  \end{tabular}\par
  \figcaption{표 7-1  이 장에서 사용할 어휘}
\end{tbbtablewrap}

\begin{calloutbox}{tbbpurple}{tbbpurplefill}
{\normalsize\bfseries\color{tbbpurple}생각해 보기}\par\vspace{3pt}
\begin{tbbitemize}
  \item 사례 7-A를 좋아하는 실제 서비스의 규모로 재현해 보자. p99.9 = 5 s이고 하루 5,000만 건의 요청을 처리한다면 5초짜리 경험은 매일 몇 번 발생하는가? "p99.9"가 여전히 예외인가?
  \item 엔지니어의 대시보드는 평균을 그렸기 때문에 초록색이었다. 대신 그려야 할 두세 개 선과 경보 임계값을 어디서 가져올지 그려 보자. §7.4를 처음부터 설계한 뒤 읽고 답을 확인하자.
  \item 사례 7-B의 오류 행은 "0.4\% (23× 504, 17× 502)"다. 제6장의 상태 코드 현장 안내서에 따르면 두 코드 가운데 로드 밸런서가 자기 문제를 자백하는 것은 무엇이고, 서버를 탓하는 것은 무엇인가? §7.4의 SLO가 둘 다 세어야 하는 이유는 무엇인가?
\end{tbbitemize}
\end{calloutbox}

\section{세 가지 추론 방식: 누가 무엇을 기다리는가?}

제안서의 첫 번째 빈칸은 채우기는 가장 쉽지만 틀리면 가장 비싸다. 프레임워크나 GPU를 논하기 전에 한 가지 질문이 서빙 시스템의 형태를 결정한다. \textbf{누가 이 추론에서 무엇을 기다리는가?} 한 사람이 \textit{이} 답을 기다리는 온라인, 아무도 기다리지 않는 배치, 한 사람이 답이 \textit{흐르는} 모습을 지켜보는 스트리밍이라는 세 답은 모델 가중치가 바이트 단위로 같더라도 운영 설명서를 거의 공유하지 않을 만큼 다른 시스템을 만든다(그림 7-2).

\begin{tbbfigure}
  \adjustbox{max width=0.94\linewidth,max totalheight=0.46\textheight}{\includegraphics{figures/fig7-2_modes.png}}\\[3pt]
  \figcaption{그림 7-2  세 가지 추론 방식. 같은 모델로 만든 세 시스템에서 온라인은 꼬리를, 배치는 비용을, 스트리밍은 빠르다는 느낌을 관리한다.}
\end{tbbfigure}

\subsection*{온라인: 모두가 가정하는 기본값이자 이 강의가 가장 엄격하게 다룰 방식}

\textbf{온라인 추론}은 사람 또는 비용을 지불하는 상위 서비스가 기다리는 동안 실시간 요청에 응답한다. 대화의 한 차례, 검색 순위 결정, 결제 흐름의 사기 탐지가 예다. 이 방식의 물리는 다음과 같다. 작업 단위는 요청 하나다. 사용자가 실제로 이탈하는 곳이 꼬리이므로 중요한 시계는 지연 시간 분포의 \textbf{꼬리}다(§7.3). 부하는 사용자가 원하는 때 도착해 버스트를 보인다(제1장). 따라서 수요보다 \textit{먼저} 용량을 프로비저닝해야 한다. 사례 1-A의 대기열 붕괴부터 12주 차 오토스케일링까지 모두 이 방식에서 다루는 이유다. 이 책에서 수식어 없이 "서빙"이라 하면 온라인을 뜻한다. 꼬리와 절충하는 배치 처리, 콜드 스타트, SLO, 플릿 규모 산정 같은 파트 2의 어려운 주제가 모두 제값을 하는 방식이다.

\subsection*{배치: 쉬운 방식 — 언제나 먼저 검토하자}

\textbf{배치 추론}은 아무도 기다리지 않는 상태에서 데이터셋에 모델을 실행한다. 오늘 업로드된 모든 사진에 밤새 태그를 붙이고, 매주 일요일 문서 말뭉치를 다시 임베딩하며, 회의 전에 지난 분기 청구 건의 점수를 매긴다. 제1장에서 말한 "학습 형태의 서빙"이다. 배치 추론이 등장하면서 미뤄 둔 질문에 답할 수 있다. 제1장의 생각해 볼 문제는 \textit{무엇이 단순해지는가}를 물었다. 솔직하고 다소 민망한 목록은 다음과 같다. 꼬리 SLO가 \textbf{사라진다.} 야간 작업의 p99를 경험하는 사람은 없고 마감 시각만 중요하다. 오토스케일링도 \textbf{사라진다.} 직접 부하를 정하므로 뒤쫓아야 할 뜻밖의 상황이 없다. §7.5의 하키 스틱은 적이 아니라 친구가 된다. 기다리는 사람이 없으므로 온라인 원칙과 정반대로 GPU를 의도적으로 95\% 이상 사용률로 돌린다. 실패의 비용은 사용자가 아니라 재시도이므로 크게 할인되는 중단 가능 용량인 \textbf{스팟 인스턴스(spot instance)}(12주 차)가 배치에 완벽히 맞는다. 남는 것은 진행 상태 체크포인트, 멱등 재시도, \textbf{백만 항목당 비용}이라는 단일 지표를 중심으로 한 작업 형태의 엔지니어링이다.

여기서 얻는 실용적인 원칙이 있다. \textbf{세상에서 가장 저렴한 지연 시간 최적화는 실제로는 아무도 기다리지 않았음을 발견하는 것이다.} "이 답이 한 시간 뒤에 와도 괜찮은가?"라고 물으면 "실시간 추론이 필요하다"는 요구사항 가운데 놀랄 만큼 많은 것이 사라진다. 온라인에서 배치로 옮기는 모든 워크로드는 SLO 엔지니어링을 예약 작업으로 바꾼다. 제안서에는 트래픽 가운데 실제로 온라인인 부분을 명시하고 나머지는 쉬운 방식으로 보내야 한다.

\subsection*{스트리밍: 배관으로 빠르다는 느낌을 산다}

\textbf{스트리밍 추론}은 작업이 계속되는 동안 결과를 전달한다. 대화 화면을 살아 있는 것처럼 느끼게 하는 LLM 토큰 스트림이 가장 유명하지만 실시간 음성 전사와 이벤트 채점도 포함된다. 시계는 둘로 나뉜다. 첫 결과가 나타날 때까지의 시간인 \textbf{첫 단위 도착 시간}이 새로운 "지연 시간"이 되고, 이후에는 초당 단위 수를 나타내는 \textbf{전송 속도}가 중요하다. 기술만큼 심리적인 기법이다. 마지막 토큰이 8초에 도착해도 첫 토큰이 300 ms에 왔다면 답이 \textit{빠르게 느껴진다.} 그래서 11주 차에는 첫 토큰 도착 시간을 일급 SLI로 다룬다. 대가는 배관이다. 버퍼링 없이 제6장의 모든 프록시와 로드 밸런서를 통과해야 하는 SSE 또는 WebSocket 장기 연결, 소비자가 생산자보다 느릴 때의 \textbf{역압(backpressure)}, 답의 절반이 이미 전송된 상태에서 "재시도"의 의미를 정하는 부분 실패 의미 체계가 필요하다. 여기서는 방식의 이름만 소개하고 서빙 엔진 내부는 11주 차에 다룬다.

\begin{tbbtablewrap}
  \begin{tabular}{@{}>{\raggedright\arraybackslash}m{\dimexpr 0.2606\tbltextwidth-2\tabcolsep\relax} >{\raggedright\arraybackslash}m{\dimexpr 0.2447\tbltextwidth-2\tabcolsep\relax} >{\raggedright\arraybackslash}m{\dimexpr 0.2394\tbltextwidth-2\tabcolsep\relax} >{\raggedright\arraybackslash}m{\dimexpr 0.2553\tbltextwidth-2\tabcolsep\relax}@{}}
  \arrayrulecolor{tbbnavy}\specialrule{1pt}{0pt}{0pt}
  \rowcolor{tbbtablehead}{\centering \bfseries\color{tbbnavy}질문} & {\centering \bfseries\color{tbbnavy}온라인} & {\centering \bfseries\color{tbbnavy}배치} & {\centering \bfseries\color{tbbnavy}스트리밍} \\
  \arrayrulecolor{tbbnavy}\specialrule{0.7pt}{0pt}{2pt}
  {\textbf{누가 기다리는가?}} & {요청마다 사람 / 실시간 상위 서비스} & {아무도 없음, 기껏해야 마감 시각} & {결과가 흐르는 모습을 보는 사람} \\
  \arrayrulecolor{tbbtableborder}\hline
  {\textbf{시계}} & {꼬리 지연 시간(p95/p99)} & {작업 완료 · 백만 항목당 비용} & {첫 결과 도착 시간 + 이후 전송 속도} \\
  \arrayrulecolor{tbbtableborder}\hline
  {\textbf{부하 패턴}} & {사용자가 결정(버스트)} & {운영자가 결정(일정)} & {장기 연결, 세션 형태} \\
  \arrayrulecolor{tbbtableborder}\hline
  {\textbf{사용률 자세}} & {의도적인 여유분(§7.5)} & {의도적으로 뜨겁게 실행(ρ→0.95)} & {여유분 + 연결 예산} \\
  \arrayrulecolor{tbbtableborder}\hline
  {\textbf{실패 비용}} & {즉시 영향을 받는 사용자} & {보이지 않는 재시도} & {문장 중간에 끊긴 스트림} \\
  \arrayrulecolor{tbbnavy}\specialrule{1pt}{2pt}{0pt}
  \end{tabular}\par
  \figcaption{표 7-2  방식을 결정하는 다섯 가지 질문}
\end{tbbtablewrap}

마지막으로 현실을 인정하자. 실제 제품은 \textbf{혼합물}이다. 대화형 제품은 답을 스트리밍하고, 매일 밤 검색용 임베딩을 미리 계산하며(배치), 요청 경로 안에서 후보 문서의 순위를 정한다(온라인). 따라서 제안서 양식은 확인란 하나가 아니라 트래픽 \textit{분해}를 요구한다. 검토자는 이 분해를 보고 팀이 실제로 시스템을 고민했는지 처음 판단한다.

\begin{calloutbox}{tbbpurple}{tbbpurplefill}
{\normalsize\bfseries\color{tbbpurple}생각해 보기}\par\vspace{3pt}
\begin{tbbitemize}
  \item 팀이 "홈페이지에 필요하므로 추천 모델은 반드시 온라인이어야 한다"고 말한다. 홈페이지는 사용자마다 세션당 최대 한 번 렌더링되고 프로필은 하루에 한 번 바뀐다. 배치 재설계를 옹호해 보자. 무엇을 미리 계산하고, 남는 작은 온라인 구성 요소가 있다면 무엇이며, SLO는 어떻게 바뀌는가?
  \item LLM 대화형 제품에서 대화가 즉각적이라고 느껴질 때와 고장 났다고 느껴질 때를 가르는 두 스트리밍 시계, 즉 첫 토큰 도착 시간과 초당 토큰 수를 적어 보자. 사용하는 대화형 서비스에서 초시계로 숫자를 추정해 직관을 확인하자.
  \item 반드시 오전 6시까지 끝나야 하는 야간 작업은 어느 방식인가? 항목별 지연 시간 목표는 없지만 마감 시각은 있다. 그렇다면 "아무도 기다리지 않는다"는 말이 완전히 맞는가? 이 작업의 SLO 형태 문장은 무엇인가? §7.4의 어휘를 사용해 작성하자.
\end{tbbitemize}
\end{calloutbox}

\section{지연 시간은 분포다 — 꼬리 읽기}

사례 7-B는 이미 기습적으로 논지를 보여 주었다. 이 절에서는 이를 규율로 만든다. 백분위수가 붙지 않은 지연 시간 숫자는 정보가 아니라 사례 7-A가 일어나기를 기다리는 상태다. 그림 7-3을 정신 모형으로 기억하자. 모든 서비스의 지연 시간은 긴 꼬리가 있는 \textbf{오른쪽으로 치우친 분포}다. 분포를 읽는 법, 증폭되는 방식, 도구가 거짓말하는 방식을 다루는 아래 세 가지 사실은 서빙 시스템 운영자의 실무 지식이다.

\begin{tbbfigure}
  \adjustbox{max width=0.94\linewidth,max totalheight=0.46\textheight}{\includegraphics{figures/fig7-3_tail.png}}\\[3pt]
  \figcaption{그림 7-3  사례 7-B 뒤의 지연 시간 분포와 백분위수 표시. 오른쪽은 팬아웃에서 드문 느림을 흔한 경험으로 바꾸는 규모에서의 꼬리 계산이다.}
\end{tbbfigure}

\subsection*{백분위수 문법}

\textbf{pN은 요청의 N\%가 더 빠르게 끝난 지연 시간이다.} p50(중앙값)은 일반적인 경험이고 p95와 p99는 운 나쁜 5\%와 1\%이며 p99.9는 CEO가 전사 회의 데모에서 뽑는 집단이다. 네 가지 문법 규칙이 정직함을 지켜 준다. 첫째, \textbf{평균은 백분위수가 아니며} 치우친 분포에서는 아무도 설명하지 못한다. 꼬리에 끌려 올라가면서도 꼬리 크기에 관해서는 아무 말도 하지 않는다. 사례 7-B의 평균 182 ms는 p50과 p90 사이에 놓였지만 누구도 경험하지 않았다. 둘째, \textbf{백분위수는 기간과 모집단의 속성}이다. \textit{어디서, 어느 기간, 어느 부하에서 측정했는지} 없으면 "p99 = 2.1 s"는 아무 의미가 없다. 화요일 아침 p99와 블랙프라이데이 p99는 서로 다른 사실이다. 셋째, \textbf{백분위수는 더할 수 없다.} 단계 A와 B가 각각 p99 = 100 ms라도 파이프라인 p99는 \textit{200 ms가 아니다.} A와 B에서 운 나쁜 1\%는 대부분 서로 다른 요청이다. 파이프라인 p99는 그 사이 어딘가에 있으며 전체 분포만 위치를 알려 줄 수 있다. 넷째, 운영상 따라오는 결론으로 \textbf{백분위수를 평균할 수 없다.} 레플리카 10개의 p99 평균은 플릿 p99가 아니다. 레플리카나 시간 전체를 집계하려면 \textit{히스토그램(histogram)}을 합쳐야 한다. 그래서 진지한 지표 시스템(14주 차 Prometheus)은 미리 계산한 백분위수가 아니라 히스토그램 버킷으로 지연 시간을 전송한다.

\subsection*{규모에서의 꼬리: 증폭되는 드문 느림}

실무자가 p50보다 p99에 집착하는 이유는 무엇일까? "꼬리에서 사용자가 이탈한다"는 이유 외에 구조적인 증폭기가 있다. 모든 서빙 엔지니어가 결국 읽게 되는 Dean과 Barroso의 \textbf{"The Tail at Scale"}(CACM, 2013)에 이름이 붙어 있다. 현대 서빙은 팬아웃한다. 사용자 요청 하나가 검색 샤드, 임베딩 서비스, 순위 결정기, 안전 필터, 템플릿 렌더러를 거친다. 각 하위 호출이 1\%만 느리다고 하자(p99 오동작). 이런 호출 100개로 팬아웃해 모두 기다리면 사용자 요청이 느린 호출을 하나 이상 만날 확률은 1 − 0.99¹⁰⁰ ≈ \textbf{63\%}다. 각 구성 요소의 \textit{드문} 사건이 사용자의 \textit{흔한} 경험으로 바뀐다. 꼬리는 평균으로 사라지지 않고 증폭된다. 이 계산 하나가 서빙 아키텍처의 한 계열을 설명한다. 팬아웃을 얕게 유지하고, 호출별 타임아웃을 사용자 마감 시각보다 짧게 설정하며(제6장의 게이트웨이 규율), 논문에서 유명해진 \textbf{헤지 요청}을 사용한다. 일정 시간이 지나면 같은 질의를 두 번째 레플리카에 보내 먼저 온 답을 택함으로써 서버 하나의 꼬리를 둘의 \textit{최솟값}으로 바꾼다. 헤지는 중복 작업을 꼬리 지연 시간과 맞바꾼다. 다시 서빙 삼각형이며 비용을 써서 꼬리를 산다.

\subsection*{부하 테스트가 거짓말하는 방식: 협조적 누락}

백분위수는 측정 도구만큼만 정직하다. 업계에서 가장 흔한 측정 버그에는 알아 둘 만한 이름이 있다. Gil Tene가 이름 붙이고 HdrHistogram과 wrk2를 만들어 해결한 \textbf{협조적 누락(coordinated omission)}이다. 버그는 다음과 같다. \textit{폐쇄 루프} 부하 시험기는 처리 중인 요청 수를 고정한다. 각 가상 사용자가 요청을 보내고 응답을 기다린 뒤 다시 보낸다. 서버가 멈추면 시험기의 가상 사용자도 모두 기다리므로, 시험기는 \textit{서버가 최악인 바로 그때 정중하게 부하 생성을 멈춘다.} 그런 다음 우연히 보낸 요청만으로 백분위수를 계산해 보고한다. 데이터에는 멈춤이 거의 나타나지 않는다. 실제 사용자는 그렇게 정중하지 않다. 서버가 멈춘 동안에도 계속 도착하며 실제 대기 시간에는 멈춤도 포함된다. 터미널 기록 7-1은 의도적으로 조작한 서버에서 그 차이를 보여 준다.

\begin{terminalbox}
\begin{Verbatim}[breaklines=true,breakanywhere=true,fontsize=\small,formatcom=\color{tbbtermfg},codes=\tbbcodeactive]
# the server is rigged to stall 1.5 s every 60 s (a "GC pause")
$ wrk -t4 -c64 -d120s http://svc/predict      # closed loop, 64 in flight
  Latency   p50 11ms    p99 13ms    max 1.51s
# during each stall only the 64 in-flight requests wait; 64 samples out
# of ~60,000/window = 0.1% -> the p99 never sees the stall. "all green."
$ wrk2 -t4 -c64 -R1000 -d120s http://svc/predict   # open loop: 1,000/s,
  Latency   p50 12ms    p99 1,140ms p99.9 1.49s    # arrivals never pause
# 1,500 requests' worth of demand arrived DURING each stall and queued.
# wrk2 charges the server for the wait, as a user's clock would.
\end{Verbatim}
\end{terminalbox}
\figcaption{터미널 기록 7-1  협조적 누락 시연. 같은 서버와 같은 멈춤에서 폐쇄 루프 도구는 p99 = 13 ms를, 개방 루프 도구는 p99 = 1.14 s를 보고한다. 둘 중 하나만 사용자를 설명한다.}

교훈은 "wrk는 나쁘고 wrk2는 좋다"가 아니다. \textbf{부하 생성기는 도착 과정을 모델링}하므로 현실과 맞는 모델을 골라야 한다는 것이다. 동시성을 고정하는 폐쇄 루프는 "클라이언트가 차례를 기다릴 때 이 서버가 얼마나 감당할 수 있는가?"에 답한다. 용량 탐색에는 좋은 질문이다. wrk2의 -R이나 Locust의 고정 처리량 방식처럼 도착률을 고정하는 개방 루프(open loop)는 "초당 λ건의 요청이 올 때 사용자는 무엇을 경험하는가?"에 답한다. 실제 도착은 예의상 멈추지 않으므로 SLO에 맞는 질문이다. 9주 차 실습에서는 Locust로 프레임워크가 서빙하는 모델을 부하 테스트하며 개방 루프 답을 요구한다. 15주 차 데모도 이를 기준으로 평가한다.

\begin{calloutbox}{tbbpurple}{tbbpurplefill}
{\normalsize\bfseries\color{tbbx5B21B6}상자 7-1  부하 테스트 위생 — 숫자가 통과해야 할 체크리스트}\par\vspace{3pt}
\begin{tbbitemize}
  \item \textbf{먼저 워밍업하고 별도로 보고한다.} 첫 요청은 모델 로드, CUDA 컨텍스트, JIT, 캐시 같은 일회성 비용을 낸다. 제2장 실습에서 측정했다. 정상 상태 백분위수는 워밍업 뒤부터 시작하며 워밍업 자체는 콜드 스타트 데이터다(12주 차).
  \item \textbf{SLO 질문에는 개방 루프를 사용한다.} 목표 λ와 그 이상에서 도착률을 고정한다. 폐쇄 루프는 용량 탐색에만 사용하고 그렇게 표시한다.
  \item \textbf{꼬리를 채울 만큼 오래 실행한다.} 100 req/s에서 p99.9는 10초마다 요청 하나다. 30초 테스트에는 표본이 약 3개뿐이다. 초가 아니라 분 단위로 실행한다.
  \item \textbf{오류와 지연 시간을 반드시 함께 보고한다.} 요청의 5\%를 타임아웃하는 서버는 살아남은 요청만 보면 훌륭한 p99를 보인다. 사례 7-B가 같은 카드에 오류율을 인쇄한 이유다. 타임아웃은 트렌치코트를 입은 꼬리다.
  \item \textbf{사용자 위치에서 측정한다.} 파드 내부가 아니라 제6장의 게이트웨이를 통과한 클라이언트 쪽에서 측정한다. 놓치기 가장 쉬운 대기열은 자기 서버 앞에 있는 대기열이다.
  \item \textbf{최소 한 번 반복한다.} 실행마다 3배 움직이는 백분위수는 서버가 아니라 테스트에 관해 말하고 있다.
\end{tbbitemize}
\end{calloutbox}

\vspace{3.0pt}

\begin{calloutbox}{tbbpurple}{tbbpurplefill}
{\normalsize\bfseries\color{tbbpurple}생각해 보기}\par\vspace{3pt}
\begin{tbbitemize}
  \item 파이프라인이 게이트웨이 → 검색 → 모델 → 후처리로 이어지고 각각 p99 = 100 ms다. 팀원이 "전체 p99 ≈ 400 ms"로 예산을 잡았다. 이 문장의 서로 다른 오류 두 가지, 즉 백분위수를 더하는 문제와 파이프라인 꼬리를 지배하는 요소에 관한 문제를 설명하고, 정직한 예산을 위해 필요한 데이터를 말하자.
  \item 프로젝트에서 예상되는 팬아웃(검색 + 모델 + 안전 필터 = 호출 3개)의 규모에서의 꼬리를 계산하자. 각각 1\%가 느리다면 느린 호출을 만나는 사용자 요청 비율은 얼마인가? 팬아웃이 얼마일 때 10\%를 넘는가?
  \item 간단한 서버의 sleep()과 부하 도구 두 번 실행만으로 협조적 누락을 60초 동안 시연하는 방법을 설계하자. 각 실행이 어떤 숫자를 보고할 것으로 예상하는가? 그런 다음 실제로 실행하자. 이 장의 실습 선택지 가운데 하나다.
\end{tbbitemize}
\end{calloutbox}

\section{SLI, SLO, SLA와 오류 예산}

수요일 아침 CTO의 지시인 \textit{"고객에게 무엇을 약속했는가? 그 문제부터 해결하라"}에는 실전에서 검증된 표준 해결책이 있다. 이 강의는 이를 발명한 Google의 Site Reliability Engineering 규율에서 그대로 가져온다. SRE 책 제4장이 이번 주 필독 자료다. 해결책은 자주 혼동되지만 절대 그래서는 안 되는 세 단어의 작은 사다리인 \textbf{SLI, SLO, SLA}다(그림 7-4). 이 사다리를 한 번 익히면 사례 7-A는 구조적으로 불가능해진다. 논쟁에 참여한 사람마다 서로 다른 단이 빠져 있었다.

\begin{tbbfigure}
  \adjustbox{max width=0.94\linewidth,max totalheight=0.46\textheight}{\includegraphics{figures/fig7-4_slo.png}}\\[3pt]
  \figcaption{그림 7-4  사다리와 예산. SLI는 측정하고 SLO는 목표를 정하며 SLA는 계약한다. 오류 예산은 목표를 정책 기계로 바꾼다.}
\end{tbbfigure}

\subsection*{SLI: 사용자처럼 측정하고 회계사처럼 센다}

\textbf{서비스 수준 지표(Service Level Indicator)}는 사용자에게 보이는 동작을 신중하게 정의한 측정값이다. 기술은 정의에 있다. 좋은 SLI는 (a) \textbf{사용자와 가까워야} 한다. 파드 내부가 아니라 사용자가 있는 게이트웨이 또는 그 너머에서 측정해야 한다. 서버 앞 대기열은 내부에서 보이지 않는다는 §7.3의 위생 규칙이다. (b) \textbf{좋은 이벤트 수를 전체 이벤트 수로 나눈 비율}이어야 한다. 비율은 기간에 걸쳐 합성할 수 있고 트래픽 변화에도 견딘다. "평균 지연 시간"이 아니라 "500 ms 안에 성공적으로 응답한 요청의 비율"을 쓴다. (c) \textbf{논쟁 없이 셀 수 있어야} 한다. 엔지니어 두 명이 정의를 읽고 같은 숫자를 내야 한다. 사례 7-A의 대시보드는 세 조건을 동시에 어겼다. 비율이 아닌 평균을 그렸고, 사용자와 가까운 곳이 아닌 서버 쪽에서 계산했으며, 성공만 포함해 논쟁의 여지를 남겼다. 표 7-3은 추론 서비스에 사용할 수 있는 선택지다.

\begin{tbbtablewrap}
  \begin{tabular}{@{}>{\raggedright\arraybackslash}m{\dimexpr 0.2287\tbltextwidth-2\tabcolsep\relax} >{\raggedright\arraybackslash}m{\dimexpr 0.4362\tbltextwidth-2\tabcolsep\relax} >{\raggedright\arraybackslash}m{\dimexpr 0.3351\tbltextwidth-2\tabcolsep\relax}@{}}
  \arrayrulecolor{tbbnavy}\specialrule{1pt}{0pt}{0pt}
  \rowcolor{tbbtablehead}{\centering \bfseries\color{tbbnavy}SLI 계열} & {\centering \bfseries\color{tbbnavy}잘 정의한 예} & {\centering \bfseries\color{tbbnavy}피하는 함정} \\
  \arrayrulecolor{tbbnavy}\specialrule{0.7pt}{0pt}{2pt}
  {\textbf{가용성}} & {게이트웨이에서 LB가 만든 502/504를 포함해 5xx가 아닌 상태로 응답한 요청의 비율} & {코드가 반환한 오류만 세는 함정. LB의 504도 사용자에게 보이는 실패다(제6장).} \\
  \arrayrulecolor{tbbtableborder}\hline
  {\textbf{지연 시간}} & {게이트웨이에서 500 ms 안에 완료된 성공 요청의 비율(동일한 표현: p99 ≤ 500 ms)} & {평균 계산, 파드 내부 측정, 지연 시간 집합에서 오류를 조용히 제외하는 함정(§7.3)} \\
  \arrayrulecolor{tbbtableborder}\hline
  {\textbf{품질(대리 지표)}} & {스키마·안전성 검증을 통과한 응답의 비율. 스트리밍은 끝까지 완료된 스트림의 비율} & {모델 \textit{정확도}를 SLO로 약속하는 함정. 정확도는 인프라가 아니라 세상의 변화에 따라 드리프트한다(14주 차 드리프트 모니터링).} \\
  \arrayrulecolor{tbbtableborder}\hline
  {\textbf{최신성(배치)}} & {야간 출력이 06:00 전에 완전히 도착한 아침의 비율} & {배치 작업에는 SLO가 전혀 없다고 보는 함정. §7.2의 마감 시각도 약속이다.} \\
  \arrayrulecolor{tbbtableborder}\hline
  {\textbf{비용(예산 지표)}} & {매주 SLIs 옆에 보고하는 요청 1,000개당 \$} & {비용을 SLI로 취급하는 함정. 비용은 요청별 약속이 아니라 보고하는 \textit{제약 조건}이다(제1장의 삼각형).} \\
  \arrayrulecolor{tbbnavy}\specialrule{1pt}{2pt}{0pt}
  \end{tabular}\par
  \figcaption{표 7-3  추론 서비스의 SLI 선택지와 각각의 함정}
\end{tbbtablewrap}

\subsection*{SLO: 기간과 부하 조항이 있는 선택한 지점}

\textbf{서비스 수준 목표(Service Level Objective)}는 명시한 기간 동안의 SLI 목표값이며 팀이 \textit{자기 자신}에게 약속하는 숫자다. 다섯 가지 필수 요소가 있는데 학생 제안서는 역사적으로 마지막 두 가지를 빼먹었다. \textbf{(1)} 참조할 SLI, \textbf{(2)} 목표값("500 ms 안에 처리한 요청이 99.5\% 이상" / "p99 ≤ 500 ms"), \textbf{(3)} 평가 기간("이동식 30일마다" — "한 시간 동안 나빴다"는 말을 논쟁이 아닌 산술로 바꾼다), \textbf{(4)} \textbf{부하 조항}("초당 최대 200건의 지속 부하")이 필요하다. §7.5에서 지연 시간이 부하의 함수임을 볼 것이므로 λ에 상한을 두지 않은 p99 약속은 아무것도 약속하지 않는다. 부하 조항 없는 SLO는 소망이다. \textbf{(5)} 이 강의에서 추가한 \textbf{비용 상한}("월 \$900 이하")도 필요하다. 제1장의 삼각형에서 빠름·대용량·저비용 가운데 둘은 비용을 지불해 셋째를 살 수 있다고 배웠으므로, 셋째 꼭짓점을 무시한 약속은 아직 엔지니어링 문장이 아니다.

목표 \textit{숫자}는 어디서 나오는가? 야심에서 나오지 않는다. 규율은 사용자가 개선을 더는 알아채지 못하는 지점인 \textbf{제품이 방어할 수 있는 가장 느슨한 SLO}를 고르는 것이다. 불필요하게 9 하나를 더할 때마다 아래에서 볼 엔지니어링 비용이 배로 늘어난다. 대화형 제품의 p99 800 ms는 충분히 방어할 수 있다. 결제 흐름 안의 사기 탐지는 주변 페이지가 더 엄격한 예산을 강제한다. 야간 배치에는 그런 예산이 전혀 없다. 숫자로 하는 제품 대화이며, 바로 \textit{지금} 제안서에서 나누는 것이 15주 차 데모를 사례 7-A와 가른다. 데모는 이번 주에 작성한 문장을 기준으로 평가한다.

\subsection*{오류 예산: 써도 되는 신뢰성 저하분}

SLO의 운영상 천재성은 여집합에 있다. 30일 동안 99.9\%를 약속하면 0.1\%인 약 \textbf{43분}은 \textit{사용해도 되는} 몫이며, 이를 \textbf{오류 예산}이라 한다. 이 관점 전환은 심리적·정치적으로 중요하다. 신뢰성이 "감히 금요일에 배포하다니" 같은 미덕 경쟁에서 자원 장부로 바뀐다. 예산이 남으면 기능을 출시하고 실험하며 자유롭게 배포한다. 예산은 속도를 위해 쓰라고 존재한다. 예산을 다 쓰면 출시를 동결하고 기간이 회복될 때까지 팀 시간을 신뢰성에 투입한다. 부사장의 기분이 아니라 \textit{사전 합의}에 따른다. 예산이 계속 남으면 과도하게 설계한 것이다. SLO를 느슨하게 하거나 더 빠르게 출시해야 한다. 예산은 사례 7-A의 논쟁도 반대 방향으로 끝낸다. PM이 불만을 상향 보고하면 답은 더 이상 "내 환경에서는 잘 된다"가 아니라 "기간이 9일 남은 현재 예산의 61\%를 소진했다. 그래프는 다음과 같다"가 된다. 빠른 소진에는 긴급 호출, 느린 소진에는 티켓을 연결하는 소진율(burn rate) 경보는 14주 차 Prometheus에서 다룬다. 여기서는 개념만 이해하면 된다.

\subsection*{100\%를 약속하지 않는 이유와 감당 가능한 SLO의 필요성}

첫 초안은 언제나 지나치게 약속하므로 제안서를 쓰기 전에 100\%에 반대하는 SRE 논리를 체화할 가치가 있다. 사용자는 불안정한 Wi-Fi, 잠든 노트북, ISP 순간 장애를 통해 99.99\%와 100\%의 차이를 경험할 수 없다. 어느 지점을 넘으면 추가한 9는 사용자와 서비스 사이의 나머지 모든 잡음 바닥으로 사라진다. 반면 9 하나를 더할 때마다 엔지니어링 비용은 대략 10배 늘어난다. 99\%는 재시작하고 지켜보는 문화를 허용한다. 99.9\%는 자동화를 요구하며 제5장의 프로브와 자가 복구가 갑자기 핵심이 된다. 99.99\%는 다중 영역 중복, 카나리 규율, 당직 교대를 요구한다. \textit{학생 프로젝트}에는 터무니없다. 9의 개수는 가용성 축일 뿐이며 같은 논리가 지연 시간 축의 가격도 매긴다. §7.5가 곧 청구서를 내민다. 약속한 p99가 매달 매시간 레플리카 형태로 구매해야 하는 여유분을 결정한다. 감당할 수 있는 SLO를 작성하자. 증거를 바탕으로 나중에 강화하는 일은 평범한 화요일이지만, 데모한 약속을 완화하는 일은 후퇴다.

\subsection*{SLA: 사다리의 맨 위 단}

\textbf{서비스 수준 협약(Service Level Agreement)}은 건물 밖으로 나간 SLO다. 서비스 크레딧, 위약금, 해지 조항 같은 결과가 붙는 외부 계약이다. 이 강의에는 두 가지 사실이면 충분하다. 첫째, SLA는 내부 SLO보다 항상 \textbf{느슨하다}(그림 7-4의 피라미드). 계약을 위반하기 훨씬 전에 자체 경보를 위반해야 한다. 그렇지 않으면 계약 \textit{자체}가 경보다. 둘째, 이미 고객 입장에서 SLA를 만났다. 할당량과 429가 있는 제1장의 모델 API 등급은 계량기가 붙은 SLA다. 직접 서빙할 때 판매하는 약속은 이 절의 메커니즘으로 뒷받침해야 할 SLA다. 프로젝트의 "SLA"는 채점자와 맺는다. 제안서의 SLO 문장이 15주 차 데모의 평가 계약이고 결과가 따른다. 그에 맞게 선택하자.

\begin{calloutbox}{tbbblue}{tbbbluefill}
{\normalsize\bfseries\color{tbbblue}상자 7-2  제안서 SLO 워크시트 — 빈칸 다섯 개와 문장 하나}\par\vspace{3pt}
\textbf{문장:} "⟨산출물, §7.6⟩을 ⟨방식, §7.2⟩으로 서빙하며, ⟨평가 기간⟩ 동안 ⟨SLI⟩ ⟨목표값⟩을 약속한다. 최대 ⟨부하⟩에서 월 ⟨비용 상한⟩ 이하로 운영한다."

\textbf{계산 예제(이 책의 진행 프로젝트):} "양자화한 3B 대화 모델을 응답 스트리밍이 있는 온라인 방식으로 서빙한다. 이동식 30일 동안 첫 토큰 도착 시간 p99 ≤ 500 ms와 게이트웨이에서 가용성 ≥ 99.5\%를 약속하며, 최대 200 req/s의 지속 부하에서 월 \$900 이하로 운영한다."

\textbf{제출 전 타당성 검사:} 모든 명사를 측정할 수 있는가(§7.3)? 부하 조항이 있는가? 평가 기간을 명시했는가? 비용 상한을 소망이 아니라 제2장의 플릿 표로 계산했는가? 먼저 §7.5를 읽자. 첫 초안은 감당할 수 없을 가능성이 크며 이를 발견하는 것 \textit{자체가 과제}다.
\end{calloutbox}

\vspace{3.0pt}

\begin{calloutbox}{tbbpurple}{tbbpurplefill}
{\normalsize\bfseries\color{tbbpurple}생각해 보기}\par\vspace{3pt}
\begin{tbbitemize}
  \item 팀에 이 절의 메커니즘이 있었다면 사례 7-A가 어떻게 전개되었을지 다시 쓰자. PM의 10:12 메시지, 이제 SLI와 오류 예산 소진량을 인용하는 엔지니어의 10:14 답변, 수요일 CTO의 말을 숫자를 포함한 메시지 세 개로 작성한다.
  \item 팀원이 "평균 지연 시간 ≤ 200 ms, 가용성 99.99\%"라는 SLO를 제안한다. 사다리 각 단에서 이 절이 제공한 서로 다른 비판 세 가지를 제시하자. SLI 선택, 목표의 감당 가능성, 빠진 조항을 살펴본다.
  \item 배치 파이프라인은 매일 06:00까지 결과를 내야 한다. 표 7-3의 최신성 SLI를 사용해 평가 기간과 목표값, 하루 한 번 실행하는 작업에서 "오류 예산을 쓴다"는 의미를 포함한 전체 SLO 문장을 작성하자.
  \item 서비스가 호출하는 모델 API 제공자는 10\% 서비스 크레딧이 있는 월간 99.9\% SLA를 제공하고, 서비스는 사용자에게 99.5\%를 약속한다. 의존성을 따라가 보자. \textit{제공자가} SLA를 지킨 달에도 어떤 실패 패턴이 \textit{서비스의} 오류 예산을 날려 버리는가? 제1장의 429 생각해 볼 문제를 이제 어휘를 갖추고 다시 푼다.
\end{tbbitemize}
\end{calloutbox}

\section{대기열: 리틀의 법칙과 하키 스틱}

제1장은 기습으로 강의를 시작했다. 사례 1-A의 데모는 3.2 req/s 상한, 길이 50인 대기열, 30초에 포기하는 게이트웨이라는 산술 때문에 죽었다. 그 계산이 "7주 차에 정식으로" 돌아온다고 약속했고 바로 이 절이다. 두 결과가 전체를 떠받친다. 너무 일반적이라 반칙처럼 느껴지는 항등식인 \textbf{리틀의 법칙}과, 결과가 너무 중요해 전 세계 모든 용량 계획 팀의 벽에 붙어 있는 곡선인 \textbf{하키 스틱}이다. 둘을 합치면 "레플리카가 몇 개 필요한가?"가 감에서 계산으로 바뀐다. 제안서 용량 절의 수학적 척추다.

\subsection*{리틀의 법칙: 세 숫자와 하나의 항등식}

\textit{어떤} 안정적인 시스템이든 경계를 그려 보자. 서버 하나, 대기열이 있는 레플리카 하나, 전체 서빙 스택, 아래층 커피숍도 된다. 세 가지를 센다. \textbf{λ}는 경계 안으로 들어오는 초당 도착률, \textbf{W}는 요청 하나가 경계 안에서 보내는 평균 시간, \textbf{L}은 어느 순간이든 내부에 있는 평균 요청 수다. John Little은 1961년에 도착 패턴, 서비스 시간 분포, 대기열 규칙, 내부 동작에 \textit{아무 가정 없이} \textbf{L = λ · W}임을 증명했다(그림 7-6). 이 법칙의 힘은 바로 무관심에 있다. GPU 배치 대기열, Service 뒤의 Kubernetes 파드, 게이트웨이에서 응답까지 전체 경로에 모두 성립한다. 서빙의 보존 법칙이다.

\begin{tbbfigure}
  \adjustbox{max width=0.94\linewidth,max totalheight=0.46\textheight}{\includegraphics{figures/fig7-6_littles.png}}\\[3pt]
  \figcaption{그림 7-6  리틀의 법칙. 어떤 경계든 세 숫자와 하나의 항등식으로 설명한다. 아래에는 제안서의 용량 절에서 수행할 규모 산정 동작을 보여 준다.}
\end{tbbfigure}

주장이 너무 좋아 보이므로 터미널 기록 7-2에서 실제 서버로 확인한다. 워커 하나가 요청마다 고정 50 ms 작업을 수행하며 서로 다른 세 동시성으로 폐쇄 루프 부하를 준다.

\begin{terminalbox}
\begin{Verbatim}[breaklines=true,breakanywhere=true,fontsize=\small,formatcom=\color{tbbtermfg},codes=\tbbcodeactive]
# server: 1 worker, exactly 50 ms per request. three runs of hey:
$ hey -z 30s -c 1  http://svc/predict | grep -E "Requests/sec|Average"
  Requests/sec:  20.0      Average: 0.0500 s    # L = 20.0 x 0.050 =  1.0
$ hey -z 30s -c 10 http://svc/predict | grep -E "Requests/sec|Average"
  Requests/sec:  20.0      Average: 0.4999 s    # L = 20.0 x 0.500 = 10.0
$ hey -z 30s -c 40 http://svc/predict | grep -E "Requests/sec|Average"
  Requests/sec:  20.0      Average: 1.9998 s    # L = 20.0 x 2.000 = 40.0
# throughput is PINNED at 1/0.050 = 20/s no matter the concurrency;
# raising c only chose how much waiting existed. L = λW, every time.
\end{Verbatim}
\end{terminalbox}
\figcaption{터미널 기록 7-2  실제로 검증한 리틀의 법칙. 서버 하나, 요청당 50 ms라는 사실 하나, 세 가지 동시성에서 항등식이 소수 셋째 자리까지 성립한다. 처리량 상한(1 ÷ 서비스 시간)은 제1장 사례 A의 계산에 진짜 이름을 붙인 것이다.}

곧바로 세 가지 실용적인 용도가 나온다. 가장 중요한 아래의 \textbf{규모 산정}에서는 목표 λ와 정직한 W가 주어졌을 때 L = λW로 처리 중 요청 수를 구하고, 레플리카 하나가 처리 중으로 유지할 수 있는 수로 나눠 플릿을 정한다. \textbf{대시보드 타당성 검사}에서는 지표가 평균 400 ms에서 300 req/s라고 하면 처리 중 요청이 약 120개여야 한다. 처리 중 요청 게이지(gauge)가 15라면 세 숫자 가운데 하나가 거짓말하며, 보통 잘못된 위치에서 측정한 지연 시간이다(§7.3). \textbf{숨은 대기열 찾기}에서는 클라이언트에서 측정한 W가 서버 쪽 W보다 크면 그 차이에 리틀의 법칙을 적용해 LB, 커널 수락 대기열, GPU 배처(batcher) 가운데 \textit{보이지 않는 대기열의 위치를 좁힌다.} 제3장의 디버깅 감각이 이제 정량화되었다.

\subsection*{하키 스틱: 시스템이 90\%에서 죽는 이유}

리틀의 법칙은 \textit{현재 상태}를 세고 두 번째 결과는 운을 시험할수록 \textit{무슨 일이 일어날지} 예측한다. 레플리카 하나를 무작위 푸아송(Poisson) 도착과 지수 분포 서비스 시간을 가진 교과서적인 \textbf{M/M/1} 대기열로 모델링하고 사용률을 \textbf{ρ = λ ÷ 용량}으로 정의하자. 시스템 평균 시간은 \textbf{W = S ÷ (1 − ρ)}이며 S는 순수 서비스 시간이다. 분모를 읽어 보자. ρ = 0.5에서 W = 2S로 대기 시간이 이미 작업 시간과 같다. ρ = 0.8이면 5S, 0.9면 10S, 0.95면 20S다(그림 7-5). 곡선의 날이 SLO를 향해 치솟는 하키 스틱 모양이다. 물리적 원인은 고갈이 아니라 \textbf{변동성}이다. 무작위 도착은 \textit{뭉치며}, 용량에 가까워지면 이 뭉침을 흡수할 여유가 없다. 그래서 평균 수요가 평균 공급보다 낮아도 대기열이 생긴다. 실제 서빙 시스템에는 GPU 배치 처리, 여러 워커, 지수 분포가 아닌 서비스 시간이 있어 M/M/1보다 복잡하지만 형태는 모든 복잡성을 견디며 숫자만 달라진다.

\begin{tbbfigure}
  \adjustbox{max width=0.94\linewidth,max totalheight=0.46\textheight}{\includegraphics{figures/fig7-5_hockey.png}}\\[3pt]
  \figcaption{그림 7-5  하키 스틱 W = S ÷ (1 − ρ)와 세 가지 운영 교훈. 여유분은 SLO 입력이고, 배치는 의도적으로 뜨겁게 실행하며, 오토스케일링은 굴곡점을 따라가기 위해 존재한다.}
\end{tbbfigure}

\begin{terminalbox}
\begin{Verbatim}[breaklines=true,breakanywhere=true,fontsize=\small,formatcom=\color{tbbtermfg},codes=\tbbcodeactive]
# same 20/s-capacity server. open-loop sweeps (wrk2), 3 min per point:
  rate      rho     p50       p99
  10.0/s    0.50     69 ms    460 ms
  16.0/s    0.80    173 ms    1.2 s
  18.0/s    0.90    347 ms    2.3 s
  19.0/s    0.95    693 ms    4.6 s     # +5% load -> +100% latency
  19.8/s    0.99    3.5 s      23 s     # queue now barely drains at all
# nothing "broke," no error was logged. arithmetic happened.
# and note the tail: p99 hits YOUR SLO long before p50 looks sick.
\end{Verbatim}
\end{terminalbox}
\figcaption{터미널 기록 7-3  측정한 하키 스틱. 모든 행이 M/M/1의 예측(p50 ≈ 0.69·W, p99 ≈ 4.6·W)과 맞는다. 마지막 주석이 운영상 의미를 알려 준다. 중앙값이 아직 정상으로 보일 때 꼬리는 SLO를 위반한다.}

세 가지 결과를 기억하자(그림 7-5 오른쪽). \textbf{여유분은 SLO 입력}이다. "p99 ≤ 500 ms" 같은 약속에는 "ρ ≤ 약 0.7에서"라는 조건이 조용히 따라붙는다. 사용률은 공짜 돈이 아니라 쓰고 있는 지연 시간 예산이다. 온라인 경로에서 95\% 사용률을 좇는 플릿은 p99를 미리 써 버릴 뿐이다. \textbf{배치는 의도적으로 뜨겁게 실행}한다. 기다리는 사람이 없으면 ρ → 0.95가 \textit{옳다.} 같은 곡선의 서로 다른 두 운영 지점이며 이 구분이 §7.2의 전체 경제 논리다. \textbf{오토스케일링은 굴곡점을 따라간다.} 12주 차 컨트롤러는 λ가 움직일 때 ρ를 초록 구간에 유지하며 절벽 전에 레플리카를 추가하고 뒤에는 제거한다. 하키 스틱 때문에 HPA에 목표 사용률 필드가 있다.

\subsection*{정직한 플릿 규모 산정과 강의계획서를 가르치는 청구서}

이제 이미 가진 요소로 제안서 용량 절을 조립하자. \textbf{1단계 — 레플리카 하나를 정직하게 측정한다.} 게이트웨이를 통과하는 개방 루프 부하(§7.3)에서 SLO 지연 시간 목표가 깨질 때까지 도착률을 훑는다. L4 레플리카 하나의 양자화한 3B 모델이 \textbf{p99 = 460 ms에서 28 req/s}를 지속한다고 하자. 500 ms 약속 안에 넉넉히 들어오며 34 req/s로 밀자 꼬리가 무너졌다. 28/약 40의 포화 용량 ≈ ρ 0.7이므로 하키 스틱이 작동한 것이다. \textbf{2단계 — 리틀의 법칙으로 교차 확인한다.} 목표 200 req/s에서 정직한 평균 W ≈ 0.16 s라면 플릿 전체 L = 200 × 0.16 = \textbf{처리 중 요청 32개}다. 레플리카 8개에 각각 약 4개로 측정한 배처 깊이와 맞는다. 9주 차 실습에서 프레임워크의 처리 중 요청 게이지가 크게 다르면 보고 있지 않은 곳에 대기열이 있다. \textbf{3단계 — 나누고 올림한다.} 200 ÷ 28 = 7.1 → 피크에 \textbf{레플리카 8개}다. 포화가 아니라 SLO \textit{지점에서} 측정했으므로 여유분은 이미 28 안에 들어 있다.

\textbf{4단계 — 비용을 계산한다.} 제2장의 플릿 표를 사용하면 8 × L4(시간당 약 \$0.80) × 744 h = \textbf{월 \$4,762}다. 제안서 상한은 \$900다. 잠시 이 숫자를 바라보자. \textit{지금 종이 위에서 무료로} 발견하는 것이 이 장의 가장 가치 있는 제출물이다. 사례 1-B의 팀은 프로덕션에서 월말에 달러로 자기 숫자를 발견했다. 간격은 판결이 아니라 \textbf{강의계획서}다. 레플리카별 처리량을 두 배로 만드는 양자화(10주 차)는 플릿을 절반으로 줄인다(→ 약 \$2,381). 200 req/s 피크가 하루 약 4시간뿐이라면 일중 곡선(diurnal curve)을 따라가는 오토스케일링(12주 차)이 레플리카 시간을 다시 대략 절반으로 줄인다(→ 약 \$1,200). 더 똑똑한 배치 처리(10–11주 차)와 배치 구성 요소의 스팟 용량(12주 차)이 나머지를 메운다. \textbf{\$4,762에서 \$900로 가는 길은 거의 한 줄씩 이 강의의 후반부와 같다.} 제안서에 정직한 오늘 가격과 적용할 지렛대가 포함된 목표 가격을 모두 적자.

\begin{calloutbox}{tbbpurple}{tbbpurplefill}
{\normalsize\bfseries\color{tbbpurple}생각해 보기}\par\vspace{3pt}
\begin{tbbitemize}
  \item 이 절의 도구로 사례 1-A를 다시 계산하자. 서비스 시간 0.31 s, 레플리카 하나, 동시 사용자 100명일 때 리틀의 법칙에서 λ\_max, ρ, W는 얼마인가? 이제 게이트웨이의 30 s 타임아웃을 더하면 도착 요청 가운데 몇 \%가 504가 되는가? 제1장에서는 감으로 했지만 이제 두 줄 계산으로 보여 주자.
  \item 대시보드에 λ = 120 req/s, 평균 지연 시간 250 ms, "처리 중 요청" 게이지 90이 표시된다. 리틀의 법칙으로 세 숫자 가운데 하나가 틀렸음을 증명하고 의심 순위를 정하자. §7.3에는 이에 관한 견해가 있다.
  \item 팀원이 "유휴 GPUs는 낭비"라며 온라인 플릿을 ρ = 0.9로 실행하자고 제안한다(제1장). 터미널 기록 7-3의 행을 사용해 ρ = 0.9의 비용을 밝히는 두 문장 반론과, 그 직관이 실제로 맞는 워크로드에 적용할 §7.2의 대안을 작성하자.
  \item 4단계의 오토스케일링 추정치를 다시 도출하자. 피크 200 req/s가 하루 4시간이고 나머지는 약 60 req/s다. 레플리카마다 28 req/s라면 완벽한 오토스케일러는 하루에 레플리카 시간을 얼마나 쓰며, 고정 플릿 8개는 얼마나 쓰는가? 12주 차 콜드 스타트 문제는 "완벽"에 어떤 영향을 주는가?
\end{tbbitemize}
\end{calloutbox}

\section{모델 형식과 변환 경로}

제안서에는 빈칸 하나가 남았다. \textbf{런타임은 정확히 어떤 파일을 실행하는가?} 제6장은 객체 스토어에 버전과 다이제스트 주소가 있는 접두사를 만들고, 시작할 때 병렬 범위 GET으로 가져오는 선반을 세웠다. 선반에 \textit{무엇을 놓는지}는 이번 주에 다룬다고 약속했다. 질문은 둘로 깔끔하게 나뉜다. \textbf{가중치}를 파일에 담는 방식은 보안 교훈 하나가 붙은 이미 해결된 문제이고, \textbf{연산}을 배포하는 방식은 10주 차에 다시 볼 현재 진행형 절충 관계다. 그림 7-7이 지도다.

\begin{tbbfigure}
  \adjustbox{max width=0.94\linewidth,max totalheight=0.46\textheight}{\includegraphics{figures/fig7-7_formats.png}}\\[3pt]
  \figcaption{그림 7-7  서빙 가능한 산출물과 형식. 왼쪽은 실제 배포 요소다. 위는 pickle의 위험과 safetensors를 비교한 가중치 컨테이너다. 아래는 한 축에 놓인 연산 그래프 형식으로 오른쪽일수록 빠르고 적용 범위가 좁다.}
\end{tbbfigure}

\subsection*{산출물: 가중치와 읽는 방법, 대화하는 방법}

현대 모델 디렉터리를 벗겨 내면 세 역할이 남는다(그림 7-7 왼쪽). GB 단위의 \textbf{가중치}인 \code{model.safetensors}, 계층 수·차원·활성화 함수처럼 그 GB를 \textit{배열하는 방법}을 설명하는 KB 단위 \textbf{아키텍처 구성 파일}인 \code{config.json}, 사람의 텍스트와 텐서 인덱스 사이의 계약인 \textbf{토크나이저} \code{tokenizer.json}이다. 어느 하나를 잃어도 나머지를 쓸 수 없다. 구성 파일 없는 가중치는 꼬리표 없는 이진 데이터다. 구성 파일과 가중치에 \textit{정확한} 토크나이저가 없으면 유창한 헛소리를 만든다. 모델은 "동작"하지만 토큰 4821이 더 이상 학습 때와 같은 뜻이 아니어서 단어 샐러드로 답하는 대표적이고 답답한 실패다. 세 요소가 \textbf{서빙 가능한 산출물}이며 다이제스트로 고정한 접두사 하나 아래에서 단위로 버전을 관리한다. 제4장의 태그 대 다이제스트 규율을 모델 스토어에서 다시 적용한 것이다. 제안서에는 모델, 크기, 양자화 상태, 라이선스, 산출물 배치를 명시한다.

\subsection*{Pickle 문제: 현실이 된 경고}

제6장은 safetensors가 "실행할 수 없는 데이터"라는 사양 항목으로 이를 설명했고 이번 주에 \textit{강하게} 반복하겠다고 약속했다. 이제 그 의미를 느껴 보자. PyTorch의 고전 체크포인트 형식인 \code{.pt}/\code{.bin}은 직렬화한 Python \textit{객체 그래프}인 \textbf{pickle}이다. 역직렬화기는 파일이 지정한 함수를 호출해 기꺼이 객체를 재구성한다. pickle의 버그가 아니라 계약이며, \textbf{체크포인트를 로딩하면 파일 제작자가 고른 코드가 실행될 수 있음}을 뜻한다.

\begin{terminalbox}
\begin{Verbatim}[breaklines=true,breakanywhere=true,fontsize=\small,formatcom=\color{tbbtermfg},codes=\tbbcodeactive]
$ python - <<'EOF'
import pickle
class Innocent:                    # what a "checkpoint" can smuggle
    def __reduce__(self):          # pickle: "to rebuild me, call this"
        return (print, ("hello from INSIDE your checkpoint load",))
blob = pickle.dumps(Innocent())
pickle.loads(blob)                 # merely LOADING the file...
EOF
hello from INSIDE your checkpoint load
# loads() called a function the file chose. print() today;
# os.system() in a malicious "fine-tune" you downloaded tomorrow.
$ python -c "from safetensors.torch import load_file as l; l('m.safetensors')"
# JSON header parsed, tensors mmap'd. nothing executes, ever. (Ch. 6)
\end{Verbatim}
\end{terminalbox}
\figcaption{터미널 기록 7-4  무해하게 시연한 pickle 문제. pickle을 역직렬화하면 파일이 지정한 임의의 함수를 호출할 수 있다. 그래서 "체크포인트를 다운로드했다"는 말은 코드 실행 사건을 뜻하며 생태계가 safetensors로 옮겨 갔다.}

따라오는 운영 규칙은 짧다. 어디서나 \textbf{safetensors}를 우선한다. 제6장의 mmap과 무복사 로딩을 사용하고 역직렬화(deserialization) 비용이 없어 \textit{더 빠르기도} 하다. 기존 pickle을 피할 수 없으면 \code{torch.load(..., weights\_only=True)}로 로딩한다. 최신 PyTorch가 바로 이 터미널 기록 때문에 이를 기본값으로 사용한다. 예외가 생기면 우회할 불편이 아니라 위험 신호로 본다. 조직 차원에서는 모델 파일을 \textit{인터넷에서 온 입력}으로 취급한다. 무작위 Docker 이미지(제4장)에 주는 만큼만 신뢰한다. 즉 출처와 다이제스트가 없으면 사용하지 않는다. 제6장의 보안 연습문제인 "공격자가 악성 pickle로 서빙 파드 안에서 코드 실행에 성공했다"는 이 절을 전제로 했다. 이제 첫 줄을 직접 쓸 수 있다.

\subsection*{그래프 형식: 하나의 축, 네 정거장, 변환 시점 규칙}

가중치는 비활성 데이터이므로 무언가가 여전히 \textit{계산}해야 한다. 선택지는 하나의 축에 놓인다(그림 7-7 아래). 오른쪽으로 갈수록 이식성과 유연성을 속도와 특수성에 맞바꾼다. \textbf{즉시 실행(eager execution) PyTorch}는 가중치 옆에 Python 코드를 배포한다. 가장 유연하지만 환경과도 가장 깊이 얽힌다. 정확한 라이브러리 버전이 필요해 제4장의 "내 컴퓨터에서는 된다"는 질병이 생기며 이미 갖춘 이미지 규율만 이를 막는다. \textbf{TorchScript / `torch.export`}는 그래프 자체를 직렬화해 C++ 런타임이 Python 없이 실행하게 한다. \textbf{ONNX}는 프레임워크 중립 교환 형식이다. \code{torch.onnx.export}가 표준 그래프를 만들고 \textbf{ONNX Runtime}이 CPU, CUDA, TensorRT 등의 교체 가능한 \textit{실행 제공자}로 실행한다. 어디서나 학습하고 어디서나 서빙하는 대신 연산자 지원 범위 가장자리에서 비용을 낸다. 사용자 정의 연산과 특이한 동적 형상은 잘 이동하지 못한다. 대표적으로 TensorRT인 \textbf{컴파일 엔진}은 그래프를 GPU 모델 하나와 형상 범위 하나에 맞게 \textit{계획}한다. 커널을 융합하고 정밀도를 선택하며 메모리를 미리 배치한다. 가장 빠르지만 A10G용으로 만든 엔진이 다른 곳에서는 실행되지 않을 만큼 좁다. 이 모든 것과 나란히 llama.cpp의 양자화 단일 파일 형식인 \textbf{GGUF}가 있다. 이 강의의 GPU 없는 대체 경로로 실용적으로 중요하다. 축을 따라 측정한 속도 차이는 10주 차 벤치마크 실습에서 다루며 이번 주에는 \textit{선택하고 패키징}만 한다.

두 규칙이 선택을 운영 가능하게 만든다. \textbf{변환은 빌드 단계이며 절대 시작 단계가 아니다.} 내보내기와 검증은 CI에서 한 번 실행해 다이제스트로 고정한 새 산출물을 만든다. 파드 시작은 이미 콜드 스타트의 핵심 경로이고 제6장에서 비용을 계산했다. 확장 사건마다 변환을 N번 하면 미묘하게 서로 다른 모델 N개를 만들 기회도 N번 생긴다. 또한 \textbf{변환한 모델은 동등성을 증명해야 한다.} 산출물에 버전을 부여하기 \textit{전에} 고정 입력 모음을 원본과 변환한 모델에 실행해 허용 오차 안에서 출력을 비교한다. 그래프 변환은 부동소수점 계산 순서를 합법적으로 바꾸므로 비트 단위 동일성은 잘못된 검사다. 하지만 "1e-3 안에서 가깝고 1순위가 같음"을 확인하는 다섯 줄 스크립트는 실제로 조용한 수치 오류를 찾아냈다. 동등성 검사가 없는 변환은 실제로 한 번도 시험하지 않은 모델이다.

\begin{calloutbox}{tbbpurple}{tbbpurplefill}
{\normalsize\bfseries\color{tbbpurple}생각해 보기}\par\vspace{3pt}
\begin{tbbitemize}
  \item 팀원 제안서에 "유연성을 극대화하기 위해 파드 시작 때 TensorRT로 변환한다"고 적혀 있다. 이 절을 바탕으로 콜드 스타트, 산출물 정체성, 동등성 검사라는 세 부분으로 비판하고 CI 형태의 대안을 제시하자.
  \item 토크나이저를 서빙 이미지에 구워 넣지 않고 버전이 있는 산출물 \textit{안에} 두어야 하는 이유는 무엇인가? 모델 v7이 v6의 토크나이저와 배포될 때 생기는 실패를 구성하자. 서비스는 무엇을 반환하고 §7.4의 SLI 가운데 어떤 것이 이를 알아채는가? 알아채는 것이 있는지도 생각하자.
  \item 프로젝트에서 즉시 실행 방식 / TorchScript / ONNX / TensorRT / GGUF의 순위를 매기자. 제안서에 어느 형식을 배포하고 10주 차에는 무엇을 실험할 것인가? 선택을 결정한 제약 조건인 팀 역량, 하드웨어, 모델 아키텍처를 말하자.
\end{tbbitemize}
\end{calloutbox}

\section{실습: 프로젝트 제안서 — 한 페이지와 다섯 숫자}

이번 주 제출물은 팀 프로젝트 제안서다. 사례 7-A의 팀이 수요일 회의를 하지 않아도 되게 해 줄 \textbf{한 페이지}다. 이 장의 각 절이 빈칸 하나씩을 채웠으며, 실습에서는 용량 계산의 기준이 될 정직한 측정값 하나와 함께 조립한다. 채점 기준은 상자 7-2의 문장을 확장한 것이다.

\subsection*{제안서의 다섯 절}

\begin{tbbitemize}
  \item \textbf{1 · 산출물(§7.6).} 모델, 매개변수 수, 양자화 상태, 라이선스, 정확한 산출물 배치(가중치 형식, 토크나이저, 구성 파일), 스토리지 접두사를 적는다. 모델이 pickle로 왔다면 CI에서 동등성 검사를 거쳐 safetensors 산출물이 되는 방법을 설명한다.
  \item \textbf{2 · 방식과 트래픽 분해(§7.2).} 실제로 온라인인 요청, 배치로 실행할 작업, 응답의 스트리밍 여부를 밝히고 주 경로에 표 7-2의 다섯 질문으로 답한다.
  \item \textbf{3 · SLO 문장(§7.4).} SLI, 목표값, 평가 기간, 부하 조항, 비용 상한이라는 다섯 조항을 모두 넣는다. 15주 차 데모는 이 계약을 기준으로 평가하며 12주 차 점검에서 다시 읽는다.
  \item \textbf{4 · 용량과 비용 범위(§7.5).} SLO에서 측정한 단일 레플리카 지속 가능 도착률(개방 루프, 게이트웨이 경유 — 상자 7-1), 리틀의 법칙 교차 확인, 피크 레플리카 수, 제2장 플릿 표로 계산한 정직한 월 가격, 10–12주 차에 적용할 지렛대를 포함한 오늘 가격 → 목표 가격 간격을 적는다.
  \item \textbf{5 · 위험과 대체 경로.} 모든 팀이 밝혀야 할 대표 위험 두 가지는 GPU 할당량·가용성과 감당할 수 없는 SLO다. 전자의 대안은 더 작은 모델 또는 이 강의가 보장하는 경로인 CPU의 GGUF다. 후자의 대안은 14주 차 새벽 2시가 아니라 지금 차분하게 선택하고 방어할 §7.4의 완화안이다.
\end{tbbitemize}

\subsection*{제안서의 기준이 되는 측정}

제안서의 숫자 하나는 추정이 아니라 반드시 \textit{측정}해야 한다. 모델을 대신해 서빙할 수 있는 \textbf{어떤} 대상이든 단일 레플리카 벤치마크를 처음 실행한다. 실제 후보가 이미 실행되면 그것을 쓰고, 아니면 4주 차 FastAPI 실습 서버를 쓴다. 절차는 상자 7-1을 축소한 형태다. 게이트웨이 뒤에 배포하고, 워밍업 요청은 버리며, 3–5개 도착률에서 개방 루프 탐색을 수행하고, 각각의 오류율과 \textit{함께} p50/p95/p99를 기록한다. 초안 SLO의 지연 시간 목표를 지키는 가장 높은 도착률을 찾는다. SLO에서 레플리카당 초당 요청 수로 나타낸 이 값이 강의 나머지 기간에 계산할 모든 용량 숫자의 씨앗이다. 9주 차 프레임워크 실습에서 Locust로 올바르게 다시 측정해 도구가 숫자의 정직성을 높이는 모습을 직접 확인한다.

\begin{calloutbox}{tbbred}{tbbxFEF2F2}
{\normalsize\bfseries\color{tbbx991B1B}상자 7-3  채점자가 매년 보는 제안서 실수 여섯 가지}\par\vspace{3pt}
\begin{tbbitemize}
  \item \textbf{"평균 지연 시간 ≤ X"} — 이 장이 한 절을 들여 해체한 SLI다. 백분위수를 쓰거나 아무것도 쓰지 않는다(§7.3).
  \item \textbf{부하 조항 없음} — λ를 명시하지 않고 약속한 p99는 아무것도 약속하지 않는다. 첫 버스트에서 하키 스틱이 먹어 치운다(§7.5).
  \item \textbf{정확도를 SLO로 사용} — 모델 품질은 인프라가 아니라 세상과 함께 드리프트한다. 대신 유효성·완료율 SLI를 약속하고 14주 차에 드리프트를 모니터링한다(§7.4, 표 7-3).
  \item \textbf{비용 = GPU만} — LB, NAT 외부 전송(제6장의 세금), 스토리지, 종료를 잊은 스테이징 환경에서도 계량기가 돌아간다. 사례 1-B는 계속 살아 있다.
  \item \textbf{폐쇄 루프 벤치마크 숫자} — 정중하게 기다리는 클라이언트로 측정한 용량은 개방 루프 현실보다 과장된다. 루프를 표시하지 않으면 점수를 잃는다(§7.3, 터미널 기록 7-1).
  \item \textbf{시작 때 변환} — 콜드 스타트 경로는 신성하다. 변환은 CI에서 동등성 검사와 함께 한 번만 실행한다(§7.6).
\end{tbbitemize}
\end{calloutbox}

\vspace{3.0pt}

제출물 체크리스트는 한 페이지 문서(다섯 절, SLO 문장의 다섯 조항), 벤치마크 표(도착률 × p50/p95/p99/오류율), 그리고 사례 1-B를 기억하기 위한 이번 주 측정 비용을 보여 주는 결제 콘솔 화면 캡처다. 보통 답은 2달러 미만이다. 습관이 핵심이다.

\begin{calloutbox}{tbbpurple}{tbbpurplefill}
{\normalsize\bfseries\color{tbbpurple}생각해 보기}\par\vspace{3pt}
\begin{tbbitemize}
  \item 다른 팀과 제안서를 바꾸고 상자 7-2 체크리스트를 기준으로 SLO 문장만 평가하자. 어느 조항이 빠졌거나 측정할 수 없는가? 매년 가장 흔히 빠지는 것은 부하 조항이다.
  \item SLO에서 측정한 단일 레플리카 도착률은 9 req/s이고 목표는 200 req/s다. L4에서 레플리카 23개, 월 약 \$13,700다. 이 장과 강의계획서의 지렛대 가운데 처음 당길 세 가지는 무엇이며 각각 무엇을 현실적으로 얻는가? 올바른 첫 답이 있으며 "GPU를 더 산다"는 아니다.
  \item 새벽 2시 버전을 작성하자. 14주 차이고 SLO가 위반되었으며 데모는 6일 뒤다. 제안서의 어느 조항을 적용하는가? 이 장이 7주 차에 무엇을 미리 적게 했기에 새벽 2시 결정이 논쟁이 아니라 읽기 문제가 되는가?
\end{tbbitemize}
\end{calloutbox}

\namedsection{요약}{열두 가지로 정리한 이 장}

\qitem{1}{\textbf{합의한 숫자가 없는 서비스에서는 정상과 고장을 구별할 수 없다.} 사례 7-A의 논쟁에는 틀린 사람이 없었고 문장 하나만 빠져 있었다. 해결책은 SLI → SLO → SLA 사다리이며 그 문장을 작성하는 것이 이번 주 제출물이다.}

\qitem{2}{\textbf{지연 시간은 분포이고 평균은 아무도 설명하지 못한다.} 사례 7-B의 평균 182 ms는 p50(92 ms)과 p90 사이에 있었지만 누구도 경험하지 않았다. 반면 매일 사용자 100명이 p99 = 2.1 s보다 느린 경험을 했다. 서빙은 백분위수로 말한다. p50은 일반적 경험, p99는 매일 만나는 집단, p99.9는 CEO의 실시간 데모다.}

\qitem{3}{\textbf{속도가 곧 매출이다.} Amazon의 +100 ms → 매출 −1\%, Google의 +400 ms → 검색 −0.59\%와 지연 제거 뒤에도 지속된 효과, Mayer의 +500 ms → 트래픽 −20\% 같은 고전적 실험은 20년 뒤에도 방향이 반복된다. 지연 시간은 비용 항목이다.}

\qitem{4}{\textbf{"누가 무엇을 기다리는가?"라는 방식 질문이 전체 시스템의 형태를 결정한다.} 온라인은 꼬리를, 배치는 비용을(뜨겁게 실행, 스팟 사용, 기다리는 사람 없음), 스트리밍은 체감 속도를(첫 결과 도착 시간 + 이후 전송 속도) 관리한다. 가장 저렴한 지연 시간 최적화는 아무도 기다리지 않았음을 발견하는 것이다.}

\qitem{5}{\textbf{백분위수 문법:} 백분위수에는 평가 기간, 사용자와 가까운 측정 위치, 부하가 필요하다. 파이프라인 단계 전체에서 더할 수 없고 레플리카 전체에서 평균할 수도 없다. 백분위수가 아니라 히스토그램을 집계해야 한다. 14주 차 Prometheus가 버킷을 전송하는 이유다.}

\qitem{6}{\textbf{팬아웃에서는 꼬리가 증폭된다}(Dean \& Barroso). 호출마다 1\% 느림 × 100곳 팬아웃이면 사용자 요청의 63\%가 느린 호출을 만난다. 그래서 팬아웃을 얕게 하고 호출별 타임아웃과 헤지 요청을 사용한다. 비용을 써서 꼬리를 되산다.}

\qitem{7}{\textbf{부하 테스트는 협조적 누락으로 거짓말한다.} 폐쇄 루프 시험기는 서버가 멈춘 바로 그때 전송을 멈춘다(터미널 기록 7-1: p99 = 13 ms 대 개방 루프의 진실 1.14 s). SLO 질문에는 고정 도착률의 개방 루프 부하와 상자 7-1의 위생 규칙이 필요하다.}

\qitem{8}{\textbf{SLO에는 다섯 조항}인 SLI, 목표값, 평가 기간, 부하 조항, 비용 상한이 있고 여집합은 오류 예산(99.9\%/30일 ≈ 43분)이다. 성적표가 아니라 출시·동결 정책이다. 100\%는 잘못된 목표다. 9 하나마다 비용이 약 10배 늘지만 사용자의 잡음 바닥으로 사라진다.}

\qitem{9}{\textbf{리틀의 법칙(L = λ·W)은 안정적인 모든 경계에 성립}하는 서빙의 보존 법칙이다. 플릿 규모 산정, 대시보드 타당성 검사(λ·W = 처리 중 요청 수여야 함), 숨은 대기열 위치 찾기에 사용한다. 제1장 사례 A의 계산을 정식화한 것이다.}

\qitem{10}{\textbf{하키 스틱(W = S ÷ (1−ρ)) 때문에 시스템이 90\%에서 죽는다.} ρ = 0.5에서 이미 대기 시간과 작업 시간이 같고, ρ = 0.9에서 0.95로 가면 다시 두 배가 된다. 여유분은 SLO 입력이고 배치는 의도적으로 뜨겁게 실행하며 오토스케일링은 굴곡점을 따라가기 위해 존재한다(12주 차).}

\qitem{11}{\textbf{플릿 규모 산정은 측정 + 나눗셈이다.} SLO에서 단일 레플리카의 도착률(개방 루프, 게이트웨이 위치) → 레플리카 수 = ceil(목표 λ ÷ 레플리카당 도착률) → 플릿 표로 비용을 계산한다. 예제의 첫 청구서인 \$4,762 대 \$900 상한은 실패가 아니다. 간격이 10–12주 차 강의계획서다.}

\qitem{12}{\textbf{서빙 가능한 산출물 = 가중치 + 구성 파일 + 토크나이저이며 다이제스트로 고정한 단위 하나로 버전을 관리한다.} Pickle 체크포인트는 로딩 때 코드를 실행할 수 있으므로(터미널 기록 7-4) safetensors를 우선하고, 아니면 \code{weights\_only=True}를 사용한다. 그래프 형식은 이식성을 속도와 맞바꾼다(즉시 실행 → TorchScript → ONNX → TensorRT, CPU 대안은 GGUF). 변환은 동등성 검사를 포함한 CI 빌드 단계이며 절대 파드 시작 단계가 아니다.}

\namedsection{연습문제}{}

\subsection*{A. 객관식}

\qitem{1}{파이프라인이 매일 밤 전체 문서 말뭉치를 다시 임베딩하며 결과는 06:00까지 준비되어야 한다. 개별 임베딩을 기다리는 사용자는 없다. 알맞은 방식과 사용률 자세는?}

\begin{tbboptions}
  \item ① 온라인 — 꼬리를 위해 ρ ≤ 0.7 유지        ② 배치 — GPU를 뜨겁게 실행(ρ → 0.95), 스팟에 적합, 마감 시각 SLO
  \item ③ 스트리밍 — 첫 임베딩 도착 시간 최적화        ④ 온라인 — 06:00 마감 시각이 지연 시간 목표이기 때문
\end{tbboptions}

\qitem{2}{사례 7-B(요청 10,000건, p50 = 92 ms, 평균 = 182 ms, p99 = 2,140 ms)에서 옳은 설명은?}

\begin{tbboptions}
  \item ① 요청 10,000건 가운데 약 100건이 2,140 ms보다 오래 걸렸다.
  \item ② 사용자 대부분이 대략 182 ms를 경험했다.
  \item ③ p99가 평균의 10×를 넘으므로 반드시 잘못되었다.
  \item ④ 요청 절반이 182 ms보다 오래 걸렸다.
\end{tbboptions}

\qitem{3}{팬아웃한 하위 호출 10개가 각각 독립적으로 1\%만 느리고 요청은 모두를 기다린다. 느린 호출을 하나 이상 만나는 사용자 요청 비율은 대략 얼마인가?}

\begin{tbboptions}
  \item ① 0.1\%        ② 1\%        ③ \textasciitilde{}10\%        ④ 63\%
\end{tbboptions}

\qitem{4}{서버가 분마다 한 번씩 1.5 s 동안 완전히 멈춘다. 보고된 p99에 이 멈춤을 드러낼 부하 테스트 설정은?}

\begin{tbboptions}
  \item ① 폐쇄 루프, 동시 연결 64개, 2분 실행
  \item ② 더 많은 연결(256개)을 사용한 폐쇄 루프
  \item ③ 용량에 가까운 고정 도착률의 개방 루프, 여러 분 실행
  \item ④ 충분히 오래 실행하기만 하면 어떤 도구든 가능
\end{tbboptions}

\qitem{5}{§7.4의 구조에 따라 완전한 SLO는?}

\begin{tbboptions}
  \item ① "p99 ≤ 500 ms" 
  \item ② "게이트웨이에서 500 ms 안에 처리한 요청이 이동식 30일마다 99.5\% 이상이며, 최대 200 req/s 부하에서 월 \$900 이하"
  \item ③ "평균 지연 시간 ≤ 200 ms, 가용성 99.99\%를 영원히 유지"
  \item ④ "언제나 최대한 빠르게"
\end{tbboptions}

\qitem{6}{가용성 SLO가 30일 동안 99.9\%(오류 예산 약 43분)다. 9일째의 사고가 50분을 사용했다. 오류 예산 정책에 따라 어떤 일이 일어나는가?}

\begin{tbboptions}
  \item ① 아무 일도 없다. 예산은 권고 사항이다.
  \item ② 사전 합의에 따라 출시를 동결하고 이동식 평가 기간이 회복될 때까지 신뢰성 작업에 집중한다.
  \item ③ 사고에 맞게 SLO를 소급해 느슨하게 한다.
  \item ④ 사후 분석(postmortem) 뒤 오류 예산이 즉시 초기화된다.
\end{tbboptions}

\qitem{7}{서비스가 평균 시스템 시간 W = 0.3 s에서 λ = 80 req/s를 받는다. Little's Law에 따른 평균 처리 중 요청 수는?}

\begin{tbboptions}
  \item ① 24        ② 267        ③ 80        ④ 도착 분포를 모르면 결정할 수 없다.
\end{tbboptions}

\qitem{8}{온라인 플릿이 ρ = 0.9(평균 지연 시간 ≈ 서비스 시간의 10배, M/M/1)로 실행된다. 트래픽이 5\% 늘어 ρ ≈ 0.95가 되었다. 모델이 예측하는 평균 지연 시간은?}

\begin{tbboptions}
  \item ① 부하 증가와 같은 +5\%        ② 대략 두 배인 서비스 시간의 약 20배
  \item ③ ρ = 1.0까지 변화 없음        ④ +50\%
\end{tbboptions}

\subsection*{B. 단답형}

\qitem{1}{서버가 멈춘 바로 그동안 도구도 부하 생성을 멈추는 부하 테스트 병리의 이름을 말하고, 보고된 백분위수가 틀린 이유를 한 문장으로 설명하자.}

\qitem{2}{목표 부하는 120 req/s이고 레플리카 하나가 SLO 지연 시간 목표에서 25 req/s를 지속한다(개방 루프 측정). 피크에 레플리카가 몇 개 필요하며 별도의 여유 계수를 더하지 않아도 되는 이유는 무엇인가?}

\qitem{3}{기존 .pt 체크포인트를 로딩하면 공격자가 고른 코드를 실행할 수 있는 이유와 safetensors 파일은 그럴 수 없는 이유를 각각 한 문장으로 설명하자.}

\qitem{4}{가용성 SLO 99.5\%의 30일 오류 예산을 분 단위로 계산하자.}

\qitem{5}{이 강의에서 완전한 SLO 문장에 필요한 다섯 조항을 나열하자.}

\subsection*{C. 논술형}

\qitem{1}{§7.4의 메커니즘을 도입한 팀의 사고 서사로 사례 7-A를 다시 쓰자. PM의 보고, SLI·현재 소진량·§7.3의 히스토그램 증거가 있는 엔지니어 답변, 오류 예산 정책이 만드는 결정을 포함한다. 그런 다음 이 메커니즘이 제거한 조직적 논쟁과 의도적으로 보존한 엔지니어링 논쟁을 한 문단으로 밝히자.}

\qitem{2}{제안서의 용량 절을 설계하자. SLO에서 L4 레플리카마다 22 req/s를 측정했고 피크는 하루 \textasciitilde{}3시간 150 req/s, 나머지는 \textasciitilde{}40 req/s이며 L4 ≈ \$0.8/hr다. 고정 플릿 월 비용과 이상적인 오토스케일 비용을 계산하고 하나 또는 혼합안을 권고한다. 오토스케일링이 도입하는 위험(12주 차 예고)과 해결 전까지 보호해 주는 §7.4 조항을 말하자.}

\qitem{3}{문서 질의응답 제품은 매일 밤 말뭉치 임베딩을 수행하고, 요청 안에서 질의별 검색 + 재순위 결정을 하며, 답을 토큰별로 스트리밍한다. §7.2의 방식으로 분해하고 구성 요소마다 SLI와 목표값을 고른 뒤(표 7-3), §7.3의 팬아웃 논리에 따라 어느 구성 요소의 꼬리가 사용자 경험을 지배하는지 찾자.}

\qitem{4}{프로젝트 형식 계획을 옹호하자. 원본 체크포인트 형식과 위험(§7.6), 9주 차에 배포할 서빙 형식, 10주 차에 시도할 변환, 각 변환을 실행할 위치(CI 대 시작 단계와 근거), 승격을 제어할 동등성 검사를 포함한다. GGUF/CPU 대체 경로와 그에 따라 강제될 SLO 변경도 밝힌다.}

\namedsection{정답과 해설}{}

\subsection*{A. 객관식}

\qitem{1}{\textbf{②} 항목 하나를 기다리는 사람이 없으므로 배치다. 뜨겁게 실행하고 중단(스팟)을 감수하며 마감 시각을 최신성 SLO("아침의 99\% 이상에서 06:00까지 출력 완료")로 약속한다. ④는 작업 마감 시각과 요청별 지연 시간을 혼동했다. (§7.2)}

\qitem{2}{\textbf{①} p99 = 2,140 ms는 10,000의 1\%인 약 100건의 요청이 더 느렸다는 뜻이다. ②는 p50과 p90 사이에서 아무도 설명하지 못한 평균의 오류다. ④는 평균과 중앙값을 혼동했고, ③은 희망 사항이다. 꼬리가 두꺼우면 p99가 평균의 10–20배가 되는 일이 흔하다. (§7.3)}

\qitem{3}{\textbf{③} 1 − 0.99¹⁰ ≈ 9.6\%다. ④의 63\%는 본문의 100곳 팬아웃이다. 팬아웃이 10분의 1이어도 호출별 1\%가 10배 증폭된다. 꼬리는 증폭된다. (§7.3)}

\qitem{4}{\textbf{③} 개방 루프 부하만 멈춤 중에도 계속 도착하므로 실제 사용자가 만들 대기열을 측정한다(터미널 기록 7-1). 폐쇄 루프 연결을 늘리면(②) 멈춤 표본이 조금 늘지만 여전히 서버와 함께 멈춘다. 실행 시간을 늘려도(④) 누락을 고치지 못한다. 도구는 매번 멈춤 \textit{중}을 빼기 때문이다. (§7.3)}

\qitem{5}{\textbf{②} 다섯 조항이 모두 있다. SLI(게이트웨이에서 500 ms 미만인 비율), 목표값(99.5\%), 평가 기간(이동식 30일), 부하 조항(≤ 200 req/s), 비용 상한(\$900/월)이다. ①은 평가 기간·부하·비용이 없고, ③은 평균과 감당하기 어려운 9의 개수를 사용하며 끝나는 평가 기간이 없다. ④는 사례 7-A다. (§7.4)}

\qitem{6}{\textbf{②} 오류 예산(약 43분)을 넘겨 썼으므로 사전 합의한 정책이 발동한다. 출시를 동결하고 이동식 평가 기간이 회복될 때까지 신뢰성 작업을 한다. 결과를 사고 \textit{전에} 합의했다는 점이 이 메커니즘의 정치적 가치다. (§7.4)}

\qitem{7}{\textbf{①} L = λ·W = 80 × 0.3 = 24다. ④는 법칙의 일반성이 제거하는 함정이다. 리틀의 법칙에는 분포 가정이 전혀 필요하지 않으므로 대시보드에 안전하게 사용할 수 있다. (§7.5)}

\qitem{8}{\textbf{②} W = S ÷ (1−ρ)이므로 0.9에서는 10S, 0.95에서는 20S다. 절벽 가까이에서는 부하 5\% 증가가 지연 시간을 두 배로 만든다. 이 비대칭 때문에 여유분이 SLO 입력이고 오토스케일러는 지연 시간만이 아니라 사용률에 반응한다. 지연 시간이 움직일 때는 이미 절벽 위다. (§7.5)}

\subsection*{B. 단답형}

\qitem{1}{\textbf{협조적 누락.} 도구의 처리 중 요청이 서버와 함께 멈추므로 사용자가 대기열을 만들 바로 그 기간에는 표본이 거의 기록되지 않는다. 최악의 경험을 제외한 생존 모집단으로 백분위수를 계산한다. (§7.3)}

\qitem{2}{ceil(120 ÷ 25) = \textbf{레플리카 5개}다. 25 req/s는 포화가 아니라 SLO 지연 시간 목표 \textit{지점에서} 측정했다. 하키 스틱 절벽까지의 거리인 여유분이 이미 숫자 안에 있다. 계수를 하나 더 곱히면 이중 계산이다. (§7.5)}

\qitem{3}{.pt는 pickle이다. 역직렬화는 파일이 지정한 함수(\code{\_\_reduce\_\_})를 호출해 Python 객체 그래프를 재구성하므로 로딩 때 제작자가 고른 코드가 실행된다. safetensors는 JSON 헤더와 원시 텐서 바이트로 된 mmap 가능한 순수 데이터라 악용할 실행 단계가 없다. (§7.6)}

\qitem{4}{30일의 0.5\%는 43,200분 × 0.005 = \textbf{216분}(3.6시간)이다. 학생 SLO로 99.5\%는 방어할 수 있지만 99.99\%(약 4.3분)는 그렇지 않은 이유 가운데 하나다. (§7.4)}

\qitem{5}{\textbf{SLI · 목표값 · 평가 기간 · 부하 조항 · 비용 상한.} ("게이트웨이에서 500 ms 미만인 요청의 비율" · "≥ 99.5\%" · "이동식 30일" · "최대 200 req/s" · "월 \$900 이하") (§7.4)}

\subsection*{C. 논술형(채점 기준)}

\qitem{1}{엔지니어 답변에 평균이 아닌 \textit{정의된 SLI}, 현재 평가 기간의 소진량("43분 가운데 28분 사용" 또는 지연 시간 SLI 대응값), CEO 표본의 위치를 보여 주는 히스토그램 증거(p99.9)가 있는지 확인한다. 오류 예산에서 기계적으로 따라오는 동결 또는 출시 결정과 마지막 구분도 있어야 한다. 메커니즘은 \textit{"느린가?"}라는 논쟁을 산술로 바꿔 제거하면서, 제안 단계에서 제대로 다퉈야 할 제품 결정인 \textit{"무엇을 약속해야 하는가?"}는 보존했다. PM의 불만 접수 채널 자체가 치우친 꼬리 표본임을 알아채면 보너스다(§7.1). (§§7.1, 7.3, 7.4)}

\qitem{2}{고정 플릿은 ceil(150/22) = 레플리카 7개 × 월 \$595 ≈ \textbf{\$4,166}이다. 이상적인 오토스케일은 레플리카 7개 × 3시간 + ceil(40/22) = 레플리카 2개 × 21시간 ≈ 하루 63 레플리카 시간 대 168 → 약 37\%, 즉 \textbf{월 \$1,560}다. 좋은 답은 2개 아래로 내려가지 않는 하한 있는 오토스케일링을 권하고 콜드 스타트 위험을 말한다. 새 레플리카가 이미지를 가져오고 가중치를 로딩하는 데 몇 분이 걸려 버스트가 스케일러를 앞지른다(12주 차, 제6장에서 가져오기 비용 계산). 확장 문제를 해결하기 전에는 명시한 λ까지만 약속이 성립하도록 SLO의 부하 조항을 계약적 보호 장치로 적용한다. (§7.5)}

\qitem{3}{분해는 말뭉치 임베딩 = 배치(최신성 SLO, 뜨겁게 실행, 스팟에 적합), 검색 + 재순위 결정 = 온라인(게이트웨이의 꼬리 SLI, 팬아웃 구성원), 토큰 생성 = 스트리밍(첫 토큰 도착 시간 + 완료율 SLI)이다. 사용자의 체감 지연 시간은 스트림 앞의 \textit{온라인 팬아웃}이 지배한다. 검색 + 재순위 결정이 첫 토큰의 순차 사전 조건이므로 p99가 첫 토큰 도착 시간을 제어한다. 백분위수를 더할 수는 없지만 가장 느린 사전 조건이 지배한다. 그 뒤 스트림이 나머지를 숨긴다. 이 관찰을 스트리밍 전 팬아웃을 얕게 하고 타임아웃을 엄격하게 설정하는 것과 연결해야 한다. (§§7.2, 7.3)}

\qitem{4}{원본의 위험을 솔직히 밝힌다. pickle이면 터미널 기록 7-4, 허브의 safetensors면 다이제스트 고정을 언급한다. 9주 차 배포는 유연성을 위한 즉시 실행 방식 또는 TorchScript, CPU 중심이면 GGUF가 방어 가능하다. 10주 차 실험은 ONNX Runtime 또는 TensorRT이며 이식성 대 속도의 절충 관계를 밝힌다. CI에서 다이제스트로 고정한 출력으로 변환하고 콜드 스타트 논거에 따라 시작 경로에는 변환이 없음을 명시한다. 동등성 관문은 고정 입력 모음, 허용 오차 + 1순위 일치다. GGUF 대체 경로의 SLO 결과도 정성적으로 수량화한다. CPU 추론이라면 지연 시간 목표를 느슨하게 하거나 부하 조항을 낮추는 정직한 §7.4 재협상을 미리 선택한다. (§7.6)}

\namedsection{더 읽을거리}{}

\textbf{[이번 주]} 표시가 붙은 항목은 이번 주 필독 또는 강력 권장 자료다. \textbf{[N주 차]} 표시는 해당 주에 도움이 될 자료를 뜻한다.

\begin{tbbitemize}
  \item \textbf{[이번 주] Google SRE Book (2016), Chapter 4, "Service Level Objectives"} — 온라인에서 무료로 읽을 수 있으며 이 장이 그대로 받아들인 SLI/SLO/SLA 사다리와 오류 예산의 정석 자료다. 제안서 SLO 문장을 쓰기 전에 읽자. Appendix "Embracing Risk"(Ch. 3)와 잘 어울린다.
  \item \textbf{[이번 주] J. Dean \& L. A. Barroso, "The Tail at Scale," CACM 2013} — 팬아웃 계산, 꼬리를 평균으로 없앨 수 없는 이유, 헤지·연결 요청 완화책을 다룬 8페이지 논문이다. 해당 10년간 가장 많이 인용된 서빙 시스템 논문이자 §7.3의 63\%가 나온 곳이다.
  \item \textbf{[이번 주, 재미] Gil Tene, "How NOT to Measure Latency" (talk, various years; HdrHistogram docs)} — 쇼맨의 즐거움으로 진행하는 협조적 누락 강연이다. 보고 나면 폐쇄 루프 p99를 다시는 믿지 못한다. wrk2와 HdrHistogram이 가져갈 도구다.
  \item \textbf{속도가 곧 매출이라는 고전 자료} — G. Linden의 2006 talk notes("Make Data Useful", +100 ms Amazon experiment), J. Brutlag의 "Speed Matters for Google Web Search"(2009), M. Mayer의 Web 2.0 2006 talk(30-results 이야기)다. 짧고 오래되었지만 모두가 여전히 인용한다. §7.1의 한계에 따라 비판적으로 읽자.
  \item \textbf{J. D. C. Little, "A Proof for the Queuing Formula L = λW," Operations Research 9(3), 1961} — 놀랄 만큼 읽기 쉬운 원 논문이다. 교과서 깊이의 더 넓은 도구 모음은 L. Kleinrock의 \textit{Queueing Systems, Vol. 1}(1975)이 여전히 표준이다. 이 강의 필요를 훨씬 넘지만 M/M/1 장은 §7.5의 엄밀한 토대다.
  \item \textbf{[9주 차] Locust documentation (constant-throughput load shapes)} — 다음 주 실습 도구다. 개방 루프 도착을 표현하는 방식을 살펴보고 첫날부터 상자 7-1 체크리스트를 locustfile에 연결하자.
  \item \textbf{[10주 차] ONNX \& ONNX Runtime docs (execution providers) · TensorRT developer guide (builder, engines) · safetensors format spec} — §7.6의 상자를 구현 깊이로 다룬다. 지금 훑어보고 10주 차 실습에서 벤치마크한다.
  \item \textbf{PyTorch security note on `torch.load` and `weights\_only`} — 기존 체크포인트의 공식 마이그레이션 경로를 포함한 터미널 기록 7-4 교훈의 공식 설명이다.
\end{tbbitemize}

\vspace{5.0pt}

\begin{calloutbox}{tbbblue}{tbbbluefill}
{\normalsize\bfseries\color{tbbblue}다음 주 — 중간고사, 이어서 제9장: 모델 서빙 프레임워크}\par\vspace{3pt}
\textbf{8주 차에는 중간고사}를 치른다. 1–7주 차의 플랫폼(가상화, 컨테이너, Kubernetes, 스토리지·네트워킹)과 이 장의 서빙 어휘를 평가한다. 가장 좋은 준비는 이 책이 계속 설계해 온 방식이다. 터미널 기록을 다시 실행하고(터미널 기록 2-1부터 7-4까지 모든 어두운 상자를 재현할 수 있다), 기념품인 종료 코드 137, PID 41288, 1956 상자, 플릿 표, 삼각형, 이제 L = λW까지 따라가자. 각 기념품이 나온 장과 다시 등장한 지점을 설명할 수 있으면 준비가 된 것이다.

\textbf{제9장}에서는 모델 아래에 전문 서빙 프레임을 놓는다. Triton Inference Server, TorchServe, BentoML, Ray Serve를 살펴보고, "서빙 프레임워크"가 4주 차 FastAPI 서버에 실제로 추가하는 동적 배치 처리(dynamic batching), 모델 관리, 지표 엔드포인트, 다중 모델 라우팅, 선택 방법, 제안서 용량 숫자를 정직하게 만드는 Locust 부하 테스트를 다룬다. 이 장의 백분위수 규율을 가지고 가자. 곧 도구를 만난다.
\end{calloutbox}

%% Chapter 9 — 모델 서빙 프레임워크
\setcounter{chapter}{8}
\chapter{모델 서빙 프레임워크}
\chaptersubtitle{모델 서버의 구조, 배치 처리의 절충 관계, 대표 프레임워크 네 가지 — 그리고 모두를 정직하게 만드는 Locust}

\begin{calloutbox}{tbbblue}{tbbbluefill}
{\normalsize\bfseries\color{tbbblue}학습 목표}\par\vspace{3pt}
\begin{tbbitemize}
  \item Locust 표 하나와 \code{nvidia-smi} 줄 하나를 보고 직접 만든 모델 서버가 GPU 성능의 일부에서 정체되는 이유를 진단하고, 간격을 메우는 세 가지 엔지니어링인 배치 처리, 파이프라이닝, 동시성을 말한다.
  \item 프로덕션 모델 서버의 \textbf{여섯 기관}인 프런트엔드, 대기열+배처, 인스턴스, 모델 관리, 관측 가능성, 상태 확인/수명 주기를 그리고 각각을 필요하게 만든 실패와 연결한다.
  \item \textbf{동적 배치 처리}의 두 설정값, 지연 시간–처리량 절충 관계, 상한을 정하는 VRAM 경계를 구성하고 논리적으로 분석하며, 벤치마크 차트가 아니라 \textit{SLO를 기준으로} 배치 상한을 고른다.
  \item 같은 여섯 기관을 바라보는 네 가지 철학으로 \textbf{Triton, TorchServe, BentoML, Ray Serve}를 비교하고 네 질문 체크리스트로 선택하며, vLLM과 KServe가 경쟁자가 아니라 이웃인 이유를 이해한다.
  \item 개방 루프 페이싱, 유효 처리량, 리틀의 법칙 감사를 포함한 정직한 Locust 스윕을 실행하고 제안서에 필요한 단 하나의 숫자인 SLO에서 레플리카당 req/s를 추출한다.
  \item 의도적으로 \textbf{재시도 폭풍}을 재현하고 증폭된 도착 요청, 붕괴한 유효 처리량, 복구 불가라는 파형을 알아보며, 이를 막는 재시도 예산 규칙을 설명한다.
\end{tbbitemize}
\end{calloutbox}

\vspace{3.5pt}

\begin{calloutbox}{tbbteal}{tbbtealfill}
{\normalsize\bfseries\color{tbbteal}핵심 용어}\par\vspace{3pt}
모델 서버 · 서빙 프레임워크 · 동적 배치 처리 · 최대 배치 크기 · 최대 대기열 지연 · 인스턴스 그룹 · 모델 저장소 · 핸들러 · 앙상블 · 유효 처리량 · 개방 루프 페이싱(고정 처리량) · 재시도 폭풍 · 부하 차단 · Triton · TorchServe · BentoML · Ray Serve · 이웃: vLLM · KServe
\end{calloutbox}

\vspace{3.5pt}

\begin{calloutbox}{tbblightgray}{tbbboxfill}
{\normalsize\bfseries\color{tbbink}장 구성}\par\vspace{3pt}
\textbf{9.1} 정중한 서버와 6개월 로드맵    \textbf{9.2} 모델 서버의 구조: 여섯 기관    \textbf{9.3} 동적 배치 처리: 두 설정값과 하나의 절충 관계    \textbf{9.4} 대표 프레임워크 네 가지와 선택 방법    \textbf{9.5} Locust: 중요한 숫자를 제대로 측정하기    \textbf{9.6} 실습: 서빙하고, 스윕하고, 의도적으로 폭풍 만들기 — 이어서 요약, 연습문제, 더 읽을거리.
\end{calloutbox}

\section{정중한 서버와 6개월 로드맵}

4주 차 FastAPI 서버는 지금까지 제 몫을 다했다. 컨테이너를 가르치고(제4장), 스케줄링되고(제5장), 스토어에서 가중치를 가져오는 법을 배우고(제6장), 제안서의 첫 벤치마크 기준이 되었다(제7장). 이번 주에는 프로덕션 트래픽의 첫 질문인 \textit{얼마나 많이 처리할 수 있는가?}를 마주한다. 정직하게 답하려면 사례가 필요하다. 실습에서 사용한 이미지 분류기를 강의 T4의 레플리카 하나에서 실행하고, Locust에게 사용자 60명이 낼 수 있는 만큼 강하게 밀도록 한다.

\subsection*{사례 A: 정체 구간}

\begin{terminalbox}
\begin{Verbatim}[breaklines=true,breakanywhere=true,fontsize=\small,formatcom=\color{tbbtermfg},codes=\tbbcodeactive]
$ locust -f locustfile.py --headless -u 60 -r 10 --run-time 5m \
         --host http://clf.demo.internal
 Type  Name         # reqs  # fails |  Avg    p50    p99  | req/s
 POST  /classify    3,529   0(0.0%) | 4,913  4,800  8,102 | 11.8
# sixty users pushing as hard as they can: 11.8 req/s. the ceiling.
$ kubectl exec deploy/clf -- nvidia-smi \
         --query-gpu=utilization.gpu --format=csv,noheader
7 %                        # the "AI server" is 93% idle...
$ kubectl top pod -l app=clf
NAME         CPU(cores)   MEMORY(bytes)
clf-6d9f4    996m         2913Mi
# ...while ONE cpu core runs flat out. there is your server.
\end{Verbatim}
\end{terminalbox}
\figcaption{사례 9-A  정체 구간. 기꺼이 요청을 보내는 사용자 60명, 11.8 req/s, 5초 대기인데 GPU는 7\%에서 놀고 CPU 코어 하나만 포화한다. 고장 난 것은 없고 모두가 지나치게 정중할 뿐이다.}

제7장의 도구로 부검해 보자. 사례의 모든 숫자는 한 가지 설계 사실에 의해 산술적으로 강제된다. \textbf{서버가 Python 스레드 하나에서 한 번에 요청 하나를 처음부터 끝까지 처리한다.} 요청마다 CPU에서 이미지 디코딩과 크기 조정에 \textasciitilde{}60 ms, 실제 GPU 추론에 \textasciitilde{}8 ms, 답 직렬화에 \textasciitilde{}15 ms를 쓴다. 직렬 작업이 \textasciitilde{}83 ms이므로 상한은 1 ÷ 0.083 ≈ \textbf{12 req/s}다. 사례 1-A의 계산이 세 번째 등장했다. 리틀의 법칙도 자백에 서명한다. 11.8 req/s로 동작하는 폐쇄 루프 사용자 60명은 각각 W = 60 ÷ 11.8 ≈ \textbf{5.1 s}를 기다려야 한다. 평균 열을 소수점까지 예측했다. GPU는 매초 11.8 × 8 ms ≈ \textbf{9\%}만 바쁘며 \code{nvidia-smi}가 이를 보여 준다. 고정된 코어 하나는 초당 11.8번 도착하는 60 ms 전처리를 담당한다. 제1장은 GPU를 수천 명의 작업자 부대라고 불렀다. 이 서버는 부대에 티스푼으로 일을 먹이면서 부대 가격을 낸다(그림 9-1).

\begin{tbbfigure}
  \adjustbox{max width=0.94\linewidth,max totalheight=0.46\textheight}{\includegraphics{figures/fig9-1_teaspoon.png}}\\[3pt]
  \figcaption{그림 9-1  숫자로 본 티스푼. 같은 모델과 같은 T4에서 정중한 서버는 12 req/s와 GPU 7\%에서 정체하고, 서빙 프레임워크의 배치 처리와 파이프라이닝은 SLO 안에서 240 req/s와 71\%에 도달한다.}
\end{tbbfigure}

세 가지 엔지니어링으로 간격을 메울 수 있으며 더 똑똑한 모델은 하나도 없다. \textbf{배치 처리}는 요청을 잠깐 모아 한 번의 넓은 GPU 실행으로 처리한다. 이 장의 주인공인 §9.3에서 다룬다. \textbf{파이프라이닝}은 번갈아 처리하지 않고 GPU가 배치 N을 계산하는 \textit{동안} CPU가 요청 N+1을 전처리한다. \textbf{동시성}은 여러 실행 스트림을 사용해 머신이 요청 하나의 왕복을 정중하게 기다리지 않게 한다. 세 가지 모두 직접 구현할 수 있다. 그렇다면 이 장의 실제 질문은 \textit{직접 해야 하는가?}다.

\subsection*{사례 B: 6개월 로드맵 또는 수렴 진화}

"그렇다"고 답한 모든 팀에 일어나는 일을 보자. 아래 git 로그는 여러 사례를 합쳤지만 ML 인프라 엔지니어에게 검토해 달라고 하면 알아보고 몸서리칠 것이다. 대부분의 커밋을 직접 작성해 보았기 때문이다.

\begin{terminalbox}
\begin{Verbatim}[breaklines=true,breakanywhere=true,fontsize=\small,formatcom=\color{tbbtermfg},codes=\tbbcodeactive]
$ git log --oneline --reverse -- serving/   # one team, eight months
a3f1e02  feat: flask server for resnet, MVP           (day 1)
9c02b11  fix: load model ONCE at startup not per req   (day 3. oops.)
5e77a4d  feat: /healthz + /ready (k8s kept killing us) (week 2)
b8d0c9e  feat: worker pool + request queue             (week 5)
1f6602a  feat: micro-batching, 10ms window (GPU util!) (week 7)
77ab3f0  feat: /metrics for prometheus                 (week 9)
c90d1f4  feat: serve model v2 alongside v1, split      (week 14)
e01b772  fix: drain in-flight reqs on SIGTERM          (week 15. outage.)
4d9a8c3  feat: config-driven model loading from s3     (week 19)
f33c6b8  docs: "how our serving works" (34 pages)      (week 23)
# congratulations: you have rebuilt 2016's TF-Serving.
# minus the tests, the docs, and the other users.
\end{Verbatim}
\end{terminalbox}
\figcaption{사례 9-B  수렴 진화. 직접 만든 서버를 유지하는 모든 팀은 정확히 이 커밋 이력을 걷는다. 각 커밋은 프로덕션이 어김없이 제공하는 실패에 답하기 때문이다. 도착점에는 서빙 프레임워크라는 이름이 있다.}

로그를 생물학처럼 읽어 보자. 추측으로 만든 것은 아무것도 없다. 각 커밋은 특정 실패가 남긴 흉터 조직이다. 3일째의 요청별 재로딩, 2주째의 활성 상태 검사에 따른 종료, 15주째 배포 중 누락된 요청이 있으며 제5장과 제3장이 뒤의 두 가지를 예측했다. 프로덕션은 어디서나 같은 선택 압력을 가하므로 팀마다 \textit{순서}도 거의 다르지 않다. 서로 다른 조직이 내부 서빙 스택을 오픈 소스로 공개했을 때 결과물에 모두 같은 기관이 있었던 이유다. 이 실패를 다른 조직보다 몇 년 먼저 겪은 Google의 \textbf{TensorFlow Serving}(2016)이 먼저 나왔고, 이어 NVIDIA의 TensorRT Inference Server(2018, 2020년 \textbf{Triton}으로 이름 변경), \textbf{TorchServe}(2020), Python 우선 물결인 \textbf{BentoML}과 \textbf{Ray Serve}가 등장했다. 서빙 프레임워크는 제품 범주라기보다 화석 기록이다. 모두가 보낸 6개월을 추출해 단단하게 만들고 다른 누군가가 유지한다. 따라서 이 장이 답할 수 있게 하는 엔지니어링 질문은 "만들 수 있는가?"가 아니다. 만들 수 있고 사례 9-B가 일정표다. 질문은 \textbf{"어느 실패부터 다시 발견하기를 멈출 것인가?"}다.

\subsection*{이 장의 어휘}

\begin{tbbtablewrap}
  \begin{tabular}{@{}>{\raggedright\arraybackslash}m{\dimexpr 0.2660\tbltextwidth-2\tabcolsep\relax} >{\raggedright\arraybackslash}m{\dimexpr 0.7340\tbltextwidth-2\tabcolsep\relax}@{}}
  \arrayrulecolor{tbbnavy}\specialrule{1pt}{0pt}{0pt}
  \rowcolor{tbbtablehead}{\centering \bfseries\color{tbbnavy}용어} & {\centering \bfseries\color{tbbnavy}의미} \\
  \arrayrulecolor{tbbnavy}\specialrule{0.7pt}{0pt}{2pt}
  {\textbf{모델 서버}} & {모델 아티팩트(§7.6)를 로딩하고 추론 요청에 응답하는 프로세스. 사례 9-A의 피고이자 §9.2의 환자다.} \\
  \arrayrulecolor{tbbtableborder}\hline
  {\textbf{서빙 프레임워크}} & {모델 서버의 여섯 기관(§9.2)을 유지 관리되는 형태로 구현한 것. Triton, TorchServe, BentoML, Ray Serve가 있다.} \\
  \arrayrulecolor{tbbtableborder}\hline
  {\textbf{동적 배치 처리}} & {동시 요청을 서버에서 한 번의 GPU 실행으로 모으는 방식. 최대 배치 크기와 최대 대기열 지연으로 제한한다(§9.3).} \\
  \arrayrulecolor{tbbtableborder}\hline
  {\textbf{인스턴스 그룹}} & {GPU 하나 또는 여러 GPU에서 실행되는 모델 복사본 N개. 계산과 전처리/후처리를 겹치게 한다. Triton의 명칭이며 모든 프레임워크에 대응 요소가 있다.} \\
  \arrayrulecolor{tbbtableborder}\hline
  {\textbf{모델 저장소}} & {프레임워크가 지켜보는 디렉터리 트리. 로컬 또는 제6장의 스토어를 직접 쓰는 s3://이며 모델 × 버전 × 설정을 런타임에 로딩할 수 있다.} \\
  \arrayrulecolor{tbbtableborder}\hline
  {\textbf{핸들러}} & {프레임워크가 모델 앞뒤에서 실행하는 사용자 코드(전처리/후처리). TorchServe의 용어이며 BentoML과 Ray는 일반 Python으로 표현한다.} \\
  \arrayrulecolor{tbbtableborder}\hline
  {\textbf{앙상블/그래프}} & {클라이언트 왕복 없이 실행하는 모델과 단계의 서버 측 파이프라인(전처리 → 모델 → 후처리).} \\
  \arrayrulecolor{tbbtableborder}\hline
  {\textbf{유효 처리량}} & {SLO 안에서 성공한 요청만 세는 처리량. 시간 초과와 오류가 표를 채울 때 중요한 숫자다(§9.5).} \\
  \arrayrulecolor{tbbtableborder}\hline
  {\textbf{개방 루프 페이싱}} & {응답과 무관하게 고정 도착률로 부하를 생성하는 방식(Locust의 constant\_throughput). §7.3의 정직성 요구를 구현한다.} \\
  \arrayrulecolor{tbbtableborder}\hline
  {\textbf{재시도 폭풍}} & {이미 포화한 서버에서 클라이언트 재시도가 도착 요청을 늘려 스스로 유지되는 과부하다. 유효 처리량이 붕괴하고 복구되지 않는다(§9.5).} \\
  \arrayrulecolor{tbbnavy}\specialrule{1pt}{2pt}{0pt}
  \end{tabular}\par
  \figcaption{표 9-1  이 장에서 사용할 어휘}
\end{tbbtablewrap}

\begin{calloutbox}{tbbpurple}{tbbpurplefill}
{\normalsize\bfseries\color{tbbpurple}생각해 보기}\par\vspace{3pt}
\begin{tbbitemize}
  \item 사례 9-A의 내부 일관성을 직접 검증하자. 서비스 시간 83 ms와 폐쇄 루프 사용자 60명으로 상한, 평균 대기 시간, GPU 사용률을 도출한다. 각각 제7장의 두 도구 가운데 무엇을 사용했는가? 이 장에서 시험 문제를 낸다면 바로 이것이다.
  \item 사례 9-B의 각 커밋을 자라나게 한 §9.2의 기관과 연결하자. 일부는 두 기관에 해당한다. 어떤 커밋 두 개가 기능이 아니라 \textit{장애}였으며 이 책의 어느 장이 예측했는가?
  \item 사례 9-B의 팀은 \textasciitilde{}8개월을 썼다. 엔지니어 두 명의 비용을 "2주 동안 Triton 학습"과 비교한 뒤, 그래도 직접 만드는 것이 옳은 상황을 말하자. 극단적인 사용자 정의 프로토콜, 특이한 하드웨어처럼 실제 사례가 있지만 경력 과시용 개발은 포함되지 않는다.
\end{tbbitemize}
\end{calloutbox}

\section{모델 서버의 구조: 여섯 기관}

제품을 비교하기 전에 기본 구조부터 익히자. 8개월 동안 직접 만든 서버든 바로 설치한 프레임워크든 진지한 모델 서버는 모두 같은 여섯 기관으로 수렴한다(그림 9-2). 각 기관이 프로덕션이 어김없이 제공하는 실패 하나에 답하기 때문이다. 프레임워크의 철학은 다르지만(§9.4) 기관은 같다. 각 기관이 존재하는 이유를 처음 만난 장과 연결하며 한 번 살펴보자.

\begin{tbbfigure}
  \adjustbox{max width=0.94\linewidth,max totalheight=0.46\textheight}{\includegraphics{figures/fig9-2_anatomy.png}}\\[3pt]
  \figcaption{그림 9-2  모델 서버의 여섯 기관. 각 상자는 특정 프로덕션 실패에 답하고, §9.4의 프레임워크 네 가지는 같은 구조를 바라보는 네 가지 견해다.}
\end{tbbfigure}

\subsection*{① 프런트엔드: 프로토콜, 검증, 출입구 정책}

프런트엔드는 제6장 게이트웨이의 공용어인 HTTP를 사용하고, 진지한 배포에서는 초당 요청량이 많은 내부 호출자를 위해 \textbf{gRPC}도 사용한다. gRPC의 이점은 구체적이다. 이미지나 텐서는 base64로 부풀린 JSON(+33\% 바이트와 두 번의 디코딩 단계)이 아니라 압축된 이진 프레임의 원시 바이트로 이동하고, 하나의 다중화 연결이 계속 반복되는 HTTP 핸드셰이크를 대신한다. 작은 텐서를 1,000 req/s로 보낼 때 직렬화 비용은 지연 시간 예산에서 눈에 보이는 몫이지만 큰 이미지를 10 req/s로 보낼 때는 차이를 알아채기 어렵다. 프런트엔드는 출입구 정책도 강제한다. \textbf{GPU 전에 검증}해야 한다. 잘못된 이미지는 배치 자리를 소비한 뒤 CUDA 커널 깊은 곳에서 500 오류로 죽는 대신 몇 마이크로초 안에 출입구에서 400 오류로 거부되어야 한다. 방식이 요구하면 스트리밍 응답도 구현한다(§7.2). 그래서 "프록시 체인이 서버 전송 이벤트를 버퍼링 없이 전달할 수 있는가?"가 제6장의 질문이었다.

\subsection*{② ③ 대기열, 배처, 인스턴스: 성능 기관}

\textbf{대기열과 배처}에는 이 장의 경제학이 담겨 있다. §9.3 전체에서 다루므로 여기서는 위치만 중요하다. 프런트엔드 뒤의 서버 \textit{내부}에 있어 요청을 검사하고 모델과 형태별로 묶으며 의도적으로 스케줄링할 수 있다. 4주 차 서버에 가장 부족했던 부분이다. 그 서버의 "대기열"은 우연히 TCP 연결 수락 대기 목록에 들어간 것이 전부였고 보이지 않으며 측정하거나 관리할 수 없었다. §7.5의 리틀 감사가 계속 찾아내던 숨은 대기열이다. \textbf{인스턴스}는 세 번째 성능 기관이다. 프레임워크는 GPU 하나에서 모델의 \textit{실행 가능한 복사본 M개}를 돌린다. 복사본 하나가 배치를 계산하는 동안 두 번째 복사본의 입력을 디코딩해 장치로 복사한다. 사례 9-A의 번갈아 처리하는 낭비를 고치는 파이프라이닝이다. M은 실제 설정값이다. 2는 계산과 데이터 전송을 겹쳐 흔히 공짜 성능을 주지만, 8은 복사본이 하나의 SM 집합과 PCIe 버스를 공유하므로 대부분 경합만 산다. 제5장의 나눌 수 없는 정수 GPU를 내부에서 체감하는 셈이다. §9.5 스윕에 실행 하나를 더하면 경험적으로 조정할 수 있다.

\subsection*{④ 모델 관리: 지켜보는 선반}

직접 만든 서버가 가장 늦게 키우고 프레임워크는 가장 먼저 제공하는 기관이 모델 관리다. 서버가 \textit{지켜보는} 모델 × 버전 디렉터리 트리인 \textbf{모델 저장소}를 두고 로딩과 언로딩을 배포가 아니라 API 호출로 수행한다. 여기서 강의의 두 흐름이 맞물린다. 첫째, 저장소는 객체 스토어의 접두사인 \code{s3:/\allowbreak /\allowbreak models/\allowbreak prod/\allowbreak } 자체일 수 있다. 제6장의 선반을 그대로 사용한다. Triton은 이를 폴링하고 예상한 병렬 범위 GET으로 새 버전을 가져오며 코드에 넣은 키 대신 IRSA로 인증한다. 둘째, 버전을 나란히 두면 13주 차 배포가 영웅적인 일이 아니라 기계적인 일이 된다. v7과 v6가 모두 로딩되어 예열된 상태이고, 트래픽 분할은 상류 시스템이 결정하며, 롤백은 레이블을 다시 가리키는 일이다. 재배포가 아니라 몇 초면 된다. 하나의 서버 프로세스가 모델별 설정과 메트릭을 갖춘 \textit{여러} 모델을 호스팅하는 방식은 롱테일 플릿의 경제적 답이기도 하다. 각각 사용률 3\%인 작은 모델 40개에 GPU 40개를 줄 이유는 없다. 제1장의 유휴 부대 계산을 플릿 규모로 적용한 것이다.

\subsection*{⑤ 관측 가능성: 제7장이 요구한 게이지}

관측 가능성 기관은 Prometheus 형식의 \code{/\allowbreak metrics} 엔드포인트다. 이를 읽는 일은 제7장 시험을 축소해 놓은 것과 같다.

\begin{terminalbox}
\begin{Verbatim}[breaklines=true,breakanywhere=true,fontsize=\small,formatcom=\color{tbbtermfg},codes=\tbbcodeactive]
$ curl -s localhost:8002/metrics | grep -E "queue|batch|count" | head -12
# latency as HISTOGRAM BUCKETS - rule 4: aggregate buckets, never p99s
nv_inference_request_duration_us_bucket{model="resnet50",le="50000"} 9871
nv_inference_request_duration_us_bucket{model="resnet50",le="100000"} 11209
# where time is spent: queue vs compute, separately measured
nv_inference_queue_duration_us{model="resnet50"}    8.31e+08
nv_inference_compute_infer_duration_us{...}         2.94e+09
# the in-flight facts for the Little audit (§7.5):
nv_inference_pending_request_count{model="resnet50"} 19
# and how well batching is working - the histogram that explains
# "throughput fell when traffic went trickle":
nv_inference_batch_size_bucket{le="8"} 4102   {le="16"} 9877  ...
\end{Verbatim}
\end{terminalbox}
\figcaption{터미널 기록 9-1  주석을 단 모델 서버의 /metrics. 지연 시간은 히스토그램 버킷으로 전송되고(백분위수를 평균할 수 없다는 제7장 규칙 ④), 대기 시간과 계산 시간은 분리되며, 처리 중 요청 게이지는 리틀 감사에 쓰이고, 배치 크기 히스토그램은 §9.3의 절충 거래가 실제로 이루어지는지 알려 준다.}

두 가지 세부사항을 강조할 가치가 있다. 대기 시간과 계산 시간은 \textit{별도} 카운터로 온다. 그래서 "지연 시간이 늘었다"를 쿼리 하나로 "대기가 늘었다"(용량/변곡점 문제, §7.5)와 "계산이 늘었다"(모델/하드웨어 문제)로 나눌 수 있다. 14주 차 대시보드 설계의 절반을 대신 해 준 셈이다. 배치 크기 분포는 배처의 정직성 계기다. 트래픽 구성이 바뀌면 배치 처리 효과가 조용히 달라지고 이 히스토그램만 이를 목격한다.

\subsection*{⑥ 상태 확인과 수명 주기: 과거 장애의 재발을 막는 기관}

상태 확인과 수명 주기는 제5장의 프로브 규율과 제3장의 종료 과정을 다시 겪지 않도록 구현한 기관이다. \code{/\allowbreak live}("프로세스가 멈추지 않음")와 \code{/\allowbreak ready}("\textbf{모델 로딩과 예열 완료}, 트래픽 전송 가능")를 나눠 롤링 업데이트가 스토어에서 가중치를 아직 스트리밍하는 레플리카에 요청을 보내지 않게 한다. SIGTERM에서는 새 요청 수락을 멈추고 처리 중인 배치를 끝낸 뒤 종료하는 정상 드레이닝을 수행해 배포가 처리 중이던 요청을 먹어 버리지 않게 한다. 사례 9-B의 15주째 장애를 영구히 고친 것이다. 실습 매니페스트에서 하나만 올바르게 연결한다면 준비 상태 프로브를 연결하자. "14:00에 배포했다"와 "14:00에 40초 동안 오류가 급증했는데 이유를 모른다"를 가르는 요소다.

\begin{calloutbox}{tbbpurple}{tbbpurplefill}
{\normalsize\bfseries\color{tbbpurple}생각해 보기}\par\vspace{3pt}
\begin{tbbitemize}
  \item 직접 만든 서버가 가중치 로딩이 끝나기 40 s 전인 HTTP 포트가 열릴 때 "준비됨"이라고 표시한다. 롤링 업데이트 중 제5장의 메커니즘을 포함해 사용자에게 정확히 어떤 실패가 보이는지 추적하고 이를 고치는 기관과 설정을 말하자.
  \item 배치 크기 분포가 메트릭 엔드포인트에 있어야 하는 이유는 무엇인가? 이 히스토그램만 설명할 수 있는 "트래픽 구성이 바뀐 뒤 처리량이 30\% 줄었다" 같은 프로덕션 미스터리를 설명하자. 요청이 조금씩 들어오면 동적 배치 처리에 무슨 일이 생기는지 생각한다.
  \item 게이트웨이와 모델 서버 사이 경로에서 gRPC와 HTTP를 비교하자. 어떤 비용이 사라지며, 큰 이미지 10 req/s보다 작은 텐서 1,000 req/s에서 차이가 더 중요한 이유는 무엇인가?
\end{tbbitemize}
\end{calloutbox}

\section{동적 배치 처리: 두 설정값과 하나의 절충 관계}

제1장은 배치 처리가 첫 요청의 대기 시간으로 처리량을 산다는 긴장을 설명했고, 제7장은 일반적인 대기 비용을 계산했으며, 사례 9-A는 이 거래를 거부할 때의 비용을 보여 주었다. 이제 메커니즘 자체를 살펴본다. 이 절부터 강의 끝까지 어떤 형태의 배치 처리든 "GPU 비용을 어떻게 감당하는가?"라는 질문에 반복해서 나오는 \textit{그} 답이기 때문이다.

GPU가 이를 중요하게 여기는 이유를 떠올리자. 제1장의 부대와 모습만 바꾼 제2장의 교훈이다. GPU는 수천 개의 스레드를 병렬로 실행한다. 배치 크기 1의 추론은 다이의 작은 조각만 밝히면서 커널 실행과, 큰 모델의 지배적 비용인 \textit{전체 가중치 텐서}를 연산 유닛으로 스트리밍하는 비용을 전부 낸다. 한 번의 실행에서 입력 16개를 처리하면 가중치를 \textbf{16개 전체에 한 번만} 읽는다. 제2장의 virtio 일괄 링 버퍼 디스크립터가 VM 탈출 비용을 분산한 것처럼 항목별 비용이 급락한다. 한 계층 위의 같은 물리다. \textbf{비싼 경계 통과 비용을 여러 작업 단위에 분산한다.}

\begin{tbbfigure}
  \adjustbox{max width=0.94\linewidth,max totalheight=0.46\textheight}{\includegraphics{figures/fig9-3_batching.png}}\\[3pt]
  \figcaption{그림 9-3  시간축으로 본 동적 배치 처리. 최대 배치 크기와 최대 대기열 지연이라는 두 설정값과 절충 관계의 가격을 보여 준다. 처리량은 선형보다 느리게 늘고 꼬리는 더 빠르게 늘며 SLO가 운영 지점을 고른다.}
\end{tbbfigure}

\subsection*{두 설정값과 결정 주체}

동적 배처는 정확히 두 가지 약속을 지키며 두 상한으로 표현한다. \textbf{max\_batch\_size}는 실행 \textit{너비}의 상한이다. 취향이 아니라 \textbf{VRAM}이 결정한다. 활성화 메모리는 배치와 함께 늘어나므로(제1장의 급증 교훈) 실제 상한은 "가중치 옆에 들어가는 가장 큰 배치 × 가장 긴 입력"이다. 제5장에서 \code{memory.max} 기반 상한을 설정하기 전에 측정하라고 한 최악의 경우와 정확히 같다. §9.5의 Locust 실행으로 의도적으로 생성한다. \textbf{max\_queue\_delay}는 첫 요청이 동료를 얼마나 \textit{오래} 기다릴지 제한한다. 이 설정값은 \textbf{SLO가 기계적으로} 정한다. p99 ≤ 500 ms이고 배치 계산 비용이 \textasciitilde{}35 ms라면 10 ms 수집 구간은 예산의 2\%를 써 \textasciitilde{}20× 처리량을 얻는다(그림 9-1의 계산). 100 ms 구간은 작업이 시작되기도 전에 예산의 5분의 1을 쓴다. 부하가 크면 배치가 즉시 차므로 지연 설정은 거의 영향을 주지 않는다. \textit{가벼운 부하}의 최악을 제한하기 위해 존재한다. 새벽 3시의 외로운 요청이 오지도 않을 동료를 구간 끝까지 기다려서는 안 된다.

실습에서 사용할 Triton의 설정 언어로 이 절 전체를 표현하면 여섯 줄이다. 바로 그것이 핵심이다.

\begin{terminalbox}
\begin{Verbatim}[breaklines=true,breakanywhere=true,fontsize=\small,formatcom=\color{tbbtermfg},codes=\tbbcodeactive]
$ cat model_repository/resnet50/config.pbtxt
platform: "onnxruntime_onnx"        # §7.6's format, meeting its runtime
max_batch_size: 32
dynamic_batching {
  preferred_batch_size: [ 8, 16 ]
  max_queue_delay_microseconds: 10000     # 10 ms: 2% of the SLO budget
}
instance_group [ { count: 2, kind: KIND_GPU } ]   # organ ③: pipelining
# eight months of Exhibit 9-B, expressed as configuration.
\end{Verbatim}
\end{terminalbox}
\figcaption{터미널 기록 9-2  설정으로 표현한 절충 관계. max\_batch\_size는 VRAM이 제한하므로 최악의 경우를 측정하고, max\_queue\_delay는 SLO가 정하므로 지연 시간 예산의 작고 알려진 몫을 쓰며, instance\_group은 파이프라이닝을 산다.}

\subsection*{엔지니어처럼 절충 관계 읽기}

그림 9-3의 가격표는 제7장의 관점으로 읽을 때 의미가 커진다. 배치 처리는 커널 실행과 가중치 읽기 비용을 분산하지만 계산 자체는 줄이지 못하므로 처리량은 1 → 64에서 95 → 460 req/s로 \textbf{선형보다 느리게} 늘어난다. 표 중간을 지나면 꼬리가 \textbf{처리량보다 빠르게} 늘어난다. 큰 배치에서는 함께 묶인 모든 요청이 가장 넓고 느린 실행을 기다리기 때문이다. 배치는 작은 팬아웃이고 §7.3에서 모두 기다릴 때 꼬리에 일어나는 일을 배웠다. 따라서 운영 지점은 "들어가는 가장 큰 배치"가 아니라 \textbf{"목표 부하에서 p99가 여전히 SLO를 지키는 가장 큰 배치"}다. 표에서는 500 ms 약속에 16–32가 해당하며, 꼬리를 두 배로 늘려 처리량 7\%를 더 얻는 64는 분명 아니다. 솔직한 사실을 하나 더 보자. 동적 배치 처리는 \textit{동시} 트래픽에 도움이 된다. 사용자 하나의 조금씩 들어오는 요청은 배치를 만들지 못하고 지연 구간 비용만 낸다. 그래서 배치 크기 히스토그램이 메트릭 엔드포인트에 있어야 한다(§9.2 생각해 볼 문제). 트래픽 구성이 바뀌면 배처 효과도 조용히 달라진다.

\subsection*{두 매개변수로 설명하는 전체 표 — 기억할 만한 모델}

그림 9-3의 다섯 행은 경험적 결과처럼 보이지만 기억할 만큼 단순한 법칙이 숨어 있다. 배치 실행 하나의 시간을 \textbf{t(B) = a + b·B}로 모델링하자. \textit{a}는 실행마다 한 번 내는 고정 비용(커널 실행, 스케줄링, 가중치 읽기)이고 \textit{b}는 배치 항목마다 내는 한계 비용(입력마다 실제로 해야 하는 계산)이다. 표에 맞추면 a ≈ 8.4 ms, b ≈ 2.05 ms가 나오고 모든 숫자가 재현된다. t(1) ≈ 10.5 ms → 95 req/s, t(16) ≈ 41 ms → 380 req/s, t(64) ≈ 140 ms → 460 req/s다. 대수에서 세 결과가 나온다. \textbf{상한은 1/b}로 여기서는 ≈ 488 req/s이며 어떤 배치 크기도 도달하지 못한다. 배치 처리는 \textit{a}를 분산하지만 \textit{b}는 건드리지 못해 곡선이 포화한다. \textbf{이익은 a/b}다. 고정 비용이 한계 비용보다 훨씬 크면(작은 모델, 큰 가중치, 빠른 GPU) 배치 처리가 놀라운 효과를 내고, \textit{b}가 지배하면(항목별 계산이 큼) 이익이 적어 다른 곳에 투자해야 한다. \textbf{10주 차에는 두 상수를 직접 공격한다.} 컴파일(TensorRT)은 커널을 융합해 \textit{a}를 줄이고, 양자화는 계산을 저렴하게 해 \textit{b}를 줄인다. 그래서 다음 주 벤치마크 실습에서 정확히 이 스윕을 다시 실행하고 두 숫자를 다시 적합한다. 실습에서 자기 모델의 a와 b를 측정하면(배치 1과 배치 32 두 번이면 충분) Locust가 확인하기 전에 전체 처리량 곡선을 예측할 수 있다.

\begin{calloutbox}{tbbteal}{tbbtealfill}
{\normalsize\bfseries\color{tbbx115E59}상자 9-1  배치 처리 이야기가 이어지는 곳}\par\vspace{3pt}
\textbf{이 장:} 이미지 분류기, 임베딩 모델, 랭커처럼 독립적이고 형태가 비슷한 요청의 동적 배치 처리를 다룬다. 요청 하나가 단위 하나이며 배치 처리는 투명하다.

\textbf{11주 차:} LLM 생성은 "형태가 비슷하다"는 가정을 치명적으로 깨뜨린다. 시퀀스마다 길이와 종료 시점이 다르므로 고전 동적 배치 처리는 가장 긴 시퀀스를 기다리며 GPU 절반을 비운다. 해결책인 토큰마다 배치에 들어오고 나가는 \textbf{연속 배치 처리}와 이를 가능하게 하는 KV 캐시가 vLLM의 존재 이유다. 그림 9-4에서 일반 프레임워크 지도 밖에 둔 이유이기도 하다.
\end{calloutbox}

\vspace{3.0pt}

\begin{calloutbox}{tbbpurple}{tbbpurplefill}
{\normalsize\bfseries\color{tbbpurple}생각해 보기}\par\vspace{3pt}
\begin{tbbitemize}
  \item 모델의 배치 계산은 배치 16에서 35 ms가 들고 SLO는 게이트웨이에서 측정한 p99 ≤ 200 ms로 진행 예보다 엄격하다. 네트워크 + 게이트웨이 + 전처리(\textasciitilde{}40 ms, 제6장)를 고려해 대기열 지연 설정 예산을 잡자. 방어할 수 있는 가장 큰 max\_queue\_delay는 얼마인가?
  \item 팀원이 배치 64를 "최고 처리량"으로 벤치마크하고 프로덕션에 max\_batch\_size = 64를 설정한다. 밤 트래픽의 40\%는 단일 사용자의 간헐적 요청이다. 이 설정이 실망을 주는 서로 다른 두 방식인 가격표의 꼬리 열과 간헐적 요청 문제를 설명하고 대신 설정할 값을 제안하자.
  \item 인스턴스 그룹(count: 2)은 배치 처리 \textit{없이도} 사례 9-A 형태 서버의 처리량을 거의 두 배로 늘렸다. §9.2의 기관 ③을 사용해 메커니즘을 한 문장으로 설명하고 count: 8이 8×를 주지 못하는 이유를 말하자. 복사본이 공유하는 리소스는 무엇인가?
\end{tbbitemize}
\end{calloutbox}

\section{대표 프레임워크 네 가지 — 같은 기관, 네 철학}

이제 제품을 보자. 강의계획서에는 \textbf{Triton Inference Server, TorchServe, BentoML, Ray Serve} 네 가지가 나온다. 혼동을 멈추는 가장 빠른 방법은 같은 여섯 기관을 바라보는 네 가지 \textit{견해}이며 두 축에서 다르다고 보는 것이다(그림 9-4). 무게 중심이 아티팩트 하나를 최대한 강하게 서빙하는 \textbf{모델 엔진}에 있는가, 모델과 로직을 조합하는 \textbf{애플리케이션 프레임워크}에 있는가가 첫 축이다. \textbf{설정}으로 구동하는 운영 우선 방식인가, \textbf{코드}로 구동하는 Python 우선 방식인가가 두 번째 축이다. 지도에 프레임워크를 놓으면 문서 대부분을 예측할 수 있다.

\begin{tbbfigure}
  \adjustbox{max width=0.94\linewidth,max totalheight=0.46\textheight}{\includegraphics{figures/fig9-4_map.png}}\\[3pt]
  \figcaption{그림 9-4  두 축에 놓인 대표 프레임워크 네 가지. LLM 전문 엔진인 vLLM(11주 차)과 어느 프레임워크든 Kubernetes에서 실행하는 플랫폼인 KServe(12주 차)는 의도적으로 지도 밖에 두었다.}
\end{tbbfigure}

\subsection*{Triton: GPU를 끝까지 활용하는 설정 우선 엔진}

\textbf{Triton}(NVIDIA, 2018년 TensorRT Inference Server로 탄생해 2020년 이름 변경)은 엔진 영역을 진지하게 밀어붙인 모습이다. C++ 서버의 \textbf{백엔드}가 §7.6에서 만든 TensorRT 엔진, 실습에서 사용할 ONNX, TorchScript, 임의의 Python까지 통일된 API 뒤에서 모델 하나 또는 수백 개를 나란히 실행한다. 모든 것은 모델 저장소와 \code{config.pbtxt}(터미널 기록 9-2)로 표현한다. 동적 배치 처리, 인스턴스 그룹, 버전 정책을 코드가 아니라 선언으로 지정한다. 강력한 기능도 GPU에 집착하는 엔진다운 것들이다. 클라이언트 왕복 없이 전처리 → 모델 → 후처리를 실행하는 서버 측 DAG인 \textbf{앙상블}, 분리형/스트리밍 응답, 배치 크기와 동시성을 대신 스윕하는 내장 부하 생성기 \code{perf\_analyzer}를 제공한다. 명확한 철학을 지닌 Locust 형태 도구로 실습에서는 둘 다 쓴다. 이 모든 것의 비용은 철학에 있다. 전처리/후처리는 앙상블 안의 모델(Python 백엔드의 "모델")이 되거나 클라이언트에 있거나 존재하지 않아야 한다. Triton은 애플리케이션 코드를 서버 밖에 두기를 선호한다. 저장소 배치는 다음과 같다.

\begin{terminalbox}
\begin{Verbatim}[breaklines=true,breakanywhere=true,fontsize=\small,formatcom=\color{tbbtermfg},codes=\tbbcodeactive]
$ tree model_repository/          # local dir or s3://models/prod (Ch. 6)
model_repository/
└── resnet50/
    ├── config.pbtxt
    ├── 1/model.onnx              # version 1, kept for rollback
    └── 2/model.onnx              # version 2, live (policy: latest)
$ curl -s localhost:8000/v2/models/resnet50/versions/2/ready
{"ready": true}                    # organ ⑥, speaking K8s's language
# metrics on :8002/metrics, gRPC on :8001 - organs ⑤ and ①, batteries included.
\end{Verbatim}
\end{terminalbox}
\figcaption{터미널 기록 9-3  Triton의 모델 저장소. 하나의 접두사 아래 버전을 나란히 두어 제6장의 선반을 그대로 사용하고 제13장의 롤백 이야기를 미리 연결해 놓았다.}

\subsection*{TorchServe: PyTorch 네이티브 베테랑과 생명력에 관한 교훈}

\textbf{TorchServe}(AWS + Meta, 2020)는 PyTorch 색채가 있는 엔진 영역이다. \code{torch-model-archiver}로 가중치와 서버 안에서 실행할 전처리/후처리 Python \textbf{핸들러}를 \code{.mar} 아카이브로 패키징하면 서버가 배치 처리, 버전이 있는 관리 API, 메트릭 같은 기관을 제공한다. 오랫동안 "PyTorch 모델이 있다"에 대한 기본 답이었고 바로 그 이유로 현장에서 만날 것이다. 이 절의 경고 사례도 담고 있다. 2024년 프로젝트가 \textbf{유지 관리 모드}에 들어가 기능은 유지하지만 적극적으로 개발하지 않게 되었다. 교훈은 "TorchServe가 나쁘다"가 아니다. \textbf{프레임워크에는 맥박이 있으므로 결혼하기 전에 확인해야 한다.} 릴리스 주기, 이슈 분류, 유지 관리자의 고용 주체를 살펴본다. 탈출 비용도 낮게 유지해야 한다. §7.6은 \textit{아티팩트}를 개방형 형식으로 보관하라고 했고 §9.2는 기관이 보편적이라고 했다. 따라서 TorchServe 플릿을 Triton 또는 후속 제품으로 옮기는 일은 작업이지만 재앙은 아니다. 비즈니스 로직을 프레임워크 내부와 결합한 팀은 그 차이를 비싸게 배웠다.

\subsection*{BentoML: 코드 우선 패키징 도구}

\textbf{BentoML}은 다른 개발자에게 답한다. \textit{"서비스가 Python이고 모델, 약간의 전처리/후처리, 어쩌면 모델 두 개와 규칙 하나가 있다. 이를 배포하고 싶다."} 형식이 지정된 엔드포인트가 있는 일반 Python 클래스를 작성하면 BentoML이 주위에 서버, 적응형 배치 처리, 메트릭 같은 기관을 제공한다. 대표 동작은 \textbf{전체를 패키징하는 것}이다. \code{bentoml build}가 코드, 의존성, 모델 참조를 버전이 있는 "bento"로 만들고 \code{bentoml containerize}가 Docker 이미지를 내보낸다. 마지막 단계는 익숙할 것이다. 사례 4-A에서 극적으로 보여 준 제4장의 규율인 의존성 고정, 계층형 이미지, "내 컴퓨터에서는 됨" 제거를 자동화한다. 도구로 만든 안전장치다. 절충 관계는 Triton의 거울상이다. Triton이 최대한 선언적인 곳에서 BentoML은 최대한 표현력이 높지만 GPU를 끝까지 활용하는 상한과 운영팀의 가독성에서 비용을 낸다. 철학을 여덟 줄로 보면 다음과 같다.

\begin{terminalbox}
\begin{Verbatim}[breaklines=true,breakanywhere=true,fontsize=\small,formatcom=\color{tbbtermfg},codes=\tbbcodeactive]
# service.py - the service IS this class; the organs wrap it
import bentoml

@bentoml.service(resources={"gpu": 1})
class Classifier:
    def __init__(self):                    # organ ⑥: load = ready
        self.model = load_onnx("clf:v7")   # from the bento's models

    @bentoml.api(batchable=True, max_batch_size=16,
                 max_latency_ms=10)         # §9.3's two knobs, as kwargs
    def classify(self, images: list) -> list:
        return self.model.run(preprocess(images))   # YOUR code, in-server
$ bentoml build && bentoml containerize clf:latest   # Ch. 4, automated
\end{Verbatim}
\end{terminalbox}
\figcaption{터미널 기록 9-4a  BentoML의 철학. 서비스는 일반 Python이고 배치 처리 설정값은 데코레이터 인수이며 고정된 컨테이너 이미지로 패키징하는 일은 하위 명령이다.}

\subsection*{Ray Serve: 분산 프로그램으로서의 서빙}

\textbf{Ray Serve}는 코드 우선 철학을 바깥으로 확장한다. \textbf{배포 단위}는 Ray 클러스터 전체에 배치하는 Python 클래스이며 각각 자체 레플리카 수, 리소스 요청, 오토스케일링 정책이 있다. \textbf{분할 GPU}도 지원해 \code{num\_gpus=0.25}로 작은 모델 네 개가 카드 하나를 공유할 수 있다. 제2장의 공유 이야기를 프레임워크 계층에서 구현했으며 이웃을 하드웨어로 격리하지 못한다는 같은 주의사항이 따른다. 일반 메서드 호출로 배포 단위를 그래프로 조합한다. 서빙 \textit{자체가} 검색 → 순위화 → 필터링 → 생성처럼 각 단계가 독립적으로 확장하는 애플리케이션이거나 팀이 이미 데이터와 학습에 Ray를 쓸 때 자연스러운 답이다. 비용은 이제 Ray 클러스터를 운영한다는 것이다. Kubernetes 위에 사는 두 번째 스케줄러로 제5장의 "두 두뇌" 패턴이다. 강력하지만 호출기를 울릴 대상이 하나 더 생긴다.

\begin{terminalbox}
\begin{Verbatim}[breaklines=true,breakanywhere=true,fontsize=\small,formatcom=\color{tbbtermfg},codes=\tbbcodeactive]
# app.py - an application of models, composed and scaled per stage
from ray import serve

@serve.deployment(num_replicas=2, ray_actor_options={"num_gpus": 0.25})
class Reranker: ...                        # a quarter of a GPU each

@serve.deployment(autoscaling_config={"min_replicas": 1,
                                      "max_replicas": 8})
class Retriever: ...                       # scales on ITS own load

@serve.deployment
class QueryPath:                           # the graph: plain method calls
    def __init__(self, retriever, reranker):
        self.r, self.rr = retriever, reranker
    async def __call__(self, q):
        return await self.rr.rank.remote(await self.r.fetch.remote(q))
\end{Verbatim}
\end{terminalbox}
\figcaption{터미널 기록 9-4b  Ray Serve의 철학. 각 단계는 자체 레플리카, GPU 할당 비율, 오토스케일링을 가진 배포 단위다. 설정한 서버가 아니라 분산 Python 프로그램으로 서빙한다.}

\begin{tbbtablewrap}
  \begin{tabular}{@{}>{\raggedright\arraybackslash}m{\dimexpr 0.1915\tbltextwidth-2\tabcolsep\relax} >{\raggedright\arraybackslash}m{\dimexpr 0.2074\tbltextwidth-2\tabcolsep\relax} >{\raggedright\arraybackslash}m{\dimexpr 0.1968\tbltextwidth-2\tabcolsep\relax} >{\raggedright\arraybackslash}m{\dimexpr 0.1968\tbltextwidth-2\tabcolsep\relax} >{\raggedright\arraybackslash}m{\dimexpr 0.2074\tbltextwidth-2\tabcolsep\relax}@{}}
  \arrayrulecolor{tbbnavy}\specialrule{1pt}{0pt}{0pt}
  \rowcolor{tbbtablehead}{\centering \bfseries\color{tbbnavy} } & {\centering \bfseries\color{tbbnavy}Triton} & {\centering \bfseries\color{tbbnavy}TorchServe} & {\centering \bfseries\color{tbbnavy}BentoML} & {\centering \bfseries\color{tbbnavy}Ray Serve} \\
  \arrayrulecolor{tbbnavy}\specialrule{0.7pt}{0pt}{2pt}
  {\textbf{배포 대상}} & {모델 저장소(설정 기반)} & {.mar 아카이브 + 핸들러} & {"bento"(코드 + 의존성 + 모델)} & {Ray 클러스터의 Python 배포 단위} \\
  \arrayrulecolor{tbbtableborder}\hline
  {\textbf{전처리/후처리 위치}} & {클라이언트 또는 앙상블/Python 백엔드} & {서버 내부 핸들러} & {서비스 코드} & {배포 단위 코드} \\
  \arrayrulecolor{tbbtableborder}\hline
  {\textbf{배치 처리}} & {동적, 선언형, 모델별} & {내장형, 모델별} & {적응형, 엔드포인트별} & {배포 단위별(@serve.batch)} \\
  \arrayrulecolor{tbbtableborder}\hline
  {\textbf{형식}} & {TensorRT · ONNX · TorchScript · Python… (§7.6 선택지)} & {PyTorch 계열} & {Python이 로딩할 수 있는 모든 것} & {Python이 로딩할 수 있는 모든 것} \\
  \arrayrulecolor{tbbtableborder}\hline
  {\textbf{여러 모델}} & {서버마다 수백 개, 앙상블} & {여러 개, 관리 API} & {서비스별 조합} & {배포 단위 그래프} \\
  \arrayrulecolor{tbbtableborder}\hline
  {\textbf{주의점}} & {애플리케이션 로직을 서버 밖에 두기} & {유지 관리 모드의 맥박(2024)} & {GPU 활용 상한} & {Ray 자체 운영} \\
  \arrayrulecolor{tbbtableborder}\hline
  {\textbf{적합한 곳}} & {GPU 극대화, 혼합 형식 플릿, 운영팀이 관리하는 플랫폼} & {기존 PyTorch 조직, 레거시 플릿} & {빠르게 배포하는 Python 팀} & {다중 모델 파이프라인, Ray 조직} \\
  \arrayrulecolor{tbbnavy}\specialrule{1pt}{2pt}{0pt}
  \end{tabular}\par
  \figcaption{표 9-2  대표 프레임워크 네 가지 한눈에 보기}
\end{tbbtablewrap}

\subsection*{정직하게 선택하기}

그림 9-7은 결정을 네 질문으로 압축한다. 워크로드 형태가 인기 모델 하나인지 애플리케이션인지, 전처리/후처리가 어디에 살아야 하는지, 누가 운영하는지, 프로젝트가 살아 있고 탈출 비용이 낮은지를 묻는다. 이어 중요한 벤치마크 규칙은 하나뿐이다. \textbf{자기 모델, 자기 SLO, 자기 하드웨어에서 후보를 벤치마크한다.} 공개 막대 차트는 다른 사람의 모델, 배치 크기, "지연 시간" 정의로 만들었다. 배치 처리가 도움을 주고 엔진이 원시 처리량에서 프레임워크를 이긴다는 \textit{형태}는 옮겨 오지만 \textit{숫자}는 절대 그렇지 않다. §9.5 덕분에 자기 벤치마크를 오후 하나에 저렴하게 실행할 수 있다. 이 강의는 9주 차 실습에 \textbf{Triton}을 고른다. 동적 배치 처리를 눈에 보이게 하고 §7.6의 ONNX를 바로 넣을 수 있다. 코드 우선 팀을 위한 BentoML 경로도 있다. 11주 차에 연속 배치 처리가 양보할 수 없는 조건이 되면 \textbf{vLLM}이 이어받는다. 12주 차에는 \textbf{KServe}가 등장해 선택한 프레임워크를 클러스터에서 실행한다. 프레임워크는 \textit{레플리카 하나가 잘 서빙하는 방법}에, 플랫폼은 \textit{몇 개 레플리카가 살아야 하는지}에 답한다.

\begin{tbbfigure}
  \adjustbox{max width=0.94\linewidth,max totalheight=0.46\textheight}{\includegraphics{figures/fig9-7_choose.png}}\\[3pt]
  \figcaption{그림 9-7  선택 체크리스트. 네 질문 뒤에 자기 워크로드에서 후보를 벤치마크한다. 아래에는 실습을 진행하기 위한 이 강의의 선택을 표시했다.}
\end{tbbfigure}

\begin{calloutbox}{tbbpurple}{tbbpurplefill}
{\normalsize\bfseries\color{tbbpurple}생각해 보기}\par\vspace{3pt}
\begin{tbbitemize}
  \item 세 팀을 그림 9-4에 놓고 각각의 선택을 옹호하자. (a) ML 엔지니어 두 명, 미세 조정한 비전 모델 하나, 월 30\% 성장하는 트래픽, (b) 공유 GPU 풀에서 여섯 팀의 모델 40개를 실행하는 플랫폼 팀, (c) 쿼리 경로가 검색 → 재순위화 → 필터링이고 색인에 이미 Ray를 쓰는 검색 팀이다.
  \item 회사가 2022년에 TorchServe를 표준화했고 지금은 2026년이다. 무엇이 바뀌었는지(§9.4의 맥박 교훈), 무엇이 보호하는지(§7.6 + §9.2), 처음 벤치마크할 대상을 포함해 어떤 마이그레이션을 제안할지 세 줄 메모로 작성하자.
  \item Triton은 전처리를 서버 밖으로 밀고 TorchServe와 BentoML은 안으로 끌어들인다. 실습의 이미지 분류기에서 (a) GPU 사용률, (b) 클라이언트 단순성, (c) 꼬리 지연 시간 측면에서 어느 배치가 이기는가? 사례 9-A의 60 ms가 실제로 어디에 쓰였는지 떠올리자.
\end{tbbitemize}
\end{calloutbox}

\section{Locust: 중요한 숫자를 제대로 측정하기}

제7장은 정직한 측정을 법으로 정했다(상자 7-1). 이 절에서는 실습 부하 도구인 \textbf{Locust}로 구현한다. Python 네이티브이고 스크립트로 제어할 수 있으며 순진하게 사용하기 전까지는 쾌적하다. 함정은 §7.3에서 이름 붙인 것이다. Locust 사용자는 기본적으로 \textbf{폐쇄 루프}라 각각 응답을 기다린 뒤 다시 보낸다. 별도로 설정하지 않은 부하 집단은 서버가 힘들어하는 바로 그때 정중하게 느려지고 조정 누락이 조용히 백분위수를 편집한다. 해결책은 한 줄이며 파일에서 가장 중요한 줄이다.

\begin{terminalbox}
\begin{Verbatim}[breaklines=true,breakanywhere=true,fontsize=\small,formatcom=\color{tbbtermfg},codes=\tbbcodeactive]
$ cat locustfile.py
from locust import HttpUser, task, constant_throughput

class Client(HttpUser):
    # THE line: each user fires 2 req/s on a fixed clock,
    # whether or not replies come back -> arrivals never pause.
    wait_time = constant_throughput(2)

    @task
    def classify(self):
        with self.client.post("/v2/models/resnet50/infer",
                data=PAYLOAD, catch_response=True, timeout=2.0) as r:
            if r.elapsed.total_seconds() > 0.5:   # the SLO, enforced
                r.failure("over SLO")             # -> goodput, not throughput

# offered load = users x 2 req/s: sweep by stepping -u 60,90,120,135...
$ locust -f locustfile.py --headless -u 120 -r 20 --run-time 4m ...
\end{Verbatim}
\end{terminalbox}
\figcaption{터미널 기록 9-5  상자 7-1을 따르는 locustfile. constant\_throughput 페이싱(개방 루프 도착), 코드 안에서 강제하는 SLO(SLO 초과 응답을 실패로 세어 구조적으로 유효 처리량 계산), 사용자 수의 단계별 조정으로 부하 스윕을 구현한다.}

세 가지 습관이 방법을 완성한다. \textbf{사용자를 충분히 프로비저닝한다.} 제공 부하는 사용자 수 × 사용자당 요청률이며 각 사용자는 평균적으로 rate × timeout ≤ 1 request outstanding이어야 한다. 사용자가 너무 적으면 변장한 폐쇄 루프가 된다. 경험 법칙은 users ≥ 2 × target λ × worst-case latency다. \textbf{한 지점만 확인하지 말고 스윕한다.} 요청률 하나로는 변곡점을 알 수 없다. 의심하는 용량의 50/75/90/105/115\%를 단계별로 3분 이상 실행하고 첫 1분은 버린다(그림 9-5). \textbf{유효 처리량을 읽는다.} locustfile이 SLO 초과 응답을 실패 처리하므로 Locust 자체 실패 열이 \textit{곧} SLO 준수율이다. 사후 필터링도 없고 별도 표에 숨은 시간 초과도 없다.

\begin{tbbfigure}
  \adjustbox{max width=0.94\linewidth,max totalheight=0.46\textheight}{\includegraphics{figures/fig9-5_locust.png}}\\[3pt]
  \figcaption{그림 9-5  스윕, 변곡점, 그리고 유지할 숫자인 SLO에서 레플리카당 req/s. 프로토콜 패널은 상자 7-1을 적용한 것이며 숫자는 제안서 추정치를 대체한다.}
\end{tbbfigure}

\begin{terminalbox}
\begin{Verbatim}[breaklines=true,breakanywhere=true,fontsize=\small,formatcom=\color{tbbtermfg},codes=\tbbcodeactive]
# the sweep, tabulated (Triton, batching on, 1 replica, T4):
  offered   p50    p95    p99    err%   goodput   verdict
  120/s     58ms   96ms  145ms   0.0%   120/s     cruising
  180/s     66ms  118ms  190ms   0.0%   180/s     healthy
  240/s     84ms  168ms  310ms   0.2%   239/s     approaching knee
  260/s    118ms  290ms  520ms   1.9%   249/s     p99 breaches SLO
  290/s    460ms  1.9s   4.1s   11.4%   182/s     past the cliff:
                                                   goodput FALLING
# the number you keep: ~238 req/s at p99<=500ms. write it down.
# Little audit: 240/s x 0.084s = ~20 in flight; Triton's queue+exec
# gauges sum to 19-22 across the run. no hidden queue. (§7.5 ✓)
\end{Verbatim}
\end{terminalbox}
\figcaption{터미널 기록 9-6  스윕. 유효 처리량은 변곡점까지 제공 부하를 따라가다가 원시 처리량이 여전히 좋아 보이는 동안 감소한다. 리틀의 법칙 감사로 대시보드가 거짓말하지 않음을 확인한다. 238 req/s @ SLO를 제안서에 넣는다.}

\subsection*{의도적으로 만드는 재시도 폭풍 하나}

제6장은 재시도 논의를 약속으로 끝냈다. \textit{"9주 차 부하 테스트에서 재시도 폭풍을 의도적으로 보여 주어 파형을 알아보게 한다."} 이제 약속을 지킬 때다. 설정은 무해해 보인다. 수많은 HTTP 라이브러리의 기본값처럼 세 번 시도하고 백오프하지 않는 클라이언트 측 재시도를 locustfile에 추가하고 제공 부하를 변곡점보다 5\% 높여 2분 동안 밀어 넣는다. 결과인 그림 9-6은 오래 바라볼 가치가 있다. \textbf{3분째 사용자 수요는 정상으로 돌아오지만 시스템은 복구되지 않는다.} 시간 초과마다 재시도가 생기고 이미 용량을 넘은 서버의 도착 요청이 세 배가 된다. 추가 도착 요청이 더 많은 시간 초과를 만들며 루프가 닫힌다. 더 나쁘게도 서버의 처리 자리는 원래 클라이언트가 오래전에 연결을 끊은 작업으로 가득하다. 처리량은 \textit{바빠} 보이지만 유효 처리량은 0에 가깝다. 수요가 정상화된 뒤에도 도착 요청이 높게 고정되고 유효 처리량은 바닥에 붙는 징후가 재시도 폭풍이다. 이를 알아보지 못한 팀은 레플리카를 더해 "고친다." 쓰레기를 더 빨리 서빙하기 위해 3× 비용을 낸다.

\begin{tbbfigure}
  \adjustbox{max width=0.94\linewidth,max totalheight=0.46\textheight}{\includegraphics{figures/fig9-6_retrystorm.png}}\\[3pt]
  \figcaption{그림 9-6  재현한 재시도 폭풍. 변곡점 너머로 2분 밀어 넣기 + 순진한 클라이언트 재시도 = 지속되는 3× 도착 요청과 버스트보다 오래가는 유효 처리량 붕괴다. 해결책은 제6장의 재시도 예산 그대로다.}
\end{tbbfigure}

해결책은 이제 파형이 붙은 제6장의 재시도 예산이다. \textbf{연결 실패만 재시도하고 시간 초과는 절대 재시도하지 않는다.} 시간 초과는 서버가 요청을 가지고 있다는 뜻이므로 다시 보내면 순수한 증폭이다. \textbf{단일 계층만 재시도}한다. 클라이언트 \textit{또는} 게이트웨이 중 하나이며 둘이 곱하지 않는다. \textbf{지터를 적용한 지수 백오프}로 재시도가 동기화되지 않고 퍼지게 한다. 서버 자체의 방어는 \textbf{일찍 부하를 버리는 것}이다. 깊이 임계값을 넘으면 헛되이 대기열에 넣지 말고 즉시 429/503으로 거부한다. 빠른 \textit{거절}은 클라이언트가 백오프하게 하지만 느린 \textit{어쩌면}은 루프를 키운다. 429가 익숙해야 한다. 제1장의 울타리를 카운터 뒤에서 본 것이며 이제 양쪽에 서 보았다. 실습에서 폭풍을 한 번 실행하고 해결책을 적용해 다시 실행하면 같은 버스트를 가볍게 떨쳐 내는 모습을 볼 수 있다. Grafana 화면 캡처 두 장은 14주 차 새벽 2시에 파형을 다시 볼 때 산문 한 장보다 가치 있을 것이다.

\begin{calloutbox}{tbbpurple}{tbbpurplefill}
{\normalsize\bfseries\color{tbbpurple}생각해 보기}\par\vspace{3pt}
\begin{tbbitemize}
  \item 스윕의 최고 단계인 목표 290 req/s, 시간 초과 2.0 s에 사용자 프로비저닝 규칙을 적용하자. constant\_throughput(2)에 사용자가 몇 명 필요하며 사용자 150명만으로 실행하면 도착 과정에 어떤 일이 조용히 생기는가?
  \item 터미널 기록 9-6에서 290/s의 제공 부하에 대한 원시 처리량은 \textasciitilde{}205/s였다. 순진한 보고서는 자랑스럽게 인쇄할 숫자다. 유효 처리량이 182/s인 이유를 재구성하고 각 숫자가 §7.4의 어느 SLI에 해당하는지 말하자.
  \item 게이트웨이(제6장)는 5xx를 두 번 재시도하고 클라이언트 라이브러리는 시간 초과를 3번 재시도한다. 60초 성능 저하 동안 사용자 요청 하나가 서버 측 실행 몇 개가 될 수 있는가? 곱셈을 그린 뒤 "한 계층만 재시도" 규칙을 적용해 다시 계산하자.
\end{tbbitemize}
\end{calloutbox}

\section{실습: 서빙하고, 스윕하고, 의도적으로 폭풍 만들기}

실습에서는 강의 T4에서 이 장 전체를 조립한다. 또는 SLO를 낮춘 minikube + CPU를 사용해도 그림 9-1 각주처럼 형태는 유지된다. Triton으로 이미지 분류기를 서빙하고, 제7장의 방식으로 측정하며, 배치 처리 스위치를 켜 절충 거래가 일어나는 모습을 본 뒤, 통제된 재난 하나를 재현한다. 제출물은 스윕 표, GPU 사용률 화면 캡처 두 장, 폭풍과 수정의 전후 쌍, \textit{측정한} 숫자로 갱신한 제안서 용량 절이다.

\subsection*{여덟 단계}

\begin{tbbitemize}
  \item \textbf{① 패키징(§7.6 적용).} 스크립트에서 분류기를 ONNX로 내보내고 동등성 검사 묶음(eager와 허용 오차 + top-1 일치)을 실행한 뒤 모델 저장소를 구성한다. 자동화할 대상을 이해하도록 CI 변환 단계를 한 번 직접 수행한다.
  \item \textbf{② 선반에 올리기(제6장).} 저장소를 \code{s3:/\allowbreak /\allowbreak <team>/\allowbreak models/\allowbreak } 또는 로컬 MinIO에 푸시한다. Triton이 직접 읽는다. 자격 증명은 클러스터 내부 IRSA로 제공하고 이미지에 키를 넣지 않는다.
  \item \textbf{③ 배포(제5장).} 강의 Triton 매니페스트에서 \code{/\allowbreak v2/\allowbreak health/\allowbreak ready}에 준비 상태 프로브를 둔다. 모델이 예열될 때까지 트래픽을 막는 기관 ⑥이다. 제5장 방법에 따라 요청량/상한을 정한다. 먼저 최악의 경우(최대 배치 × 가장 큰 이미지)를 실행하고 \code{memory.max}를 설정하기 전에 \code{memory.current}를 읽는다. 종료 코드 137은 여전히 존재한다.
  \item \textbf{④ 기준선 스윕, 배치 처리 끄기.} \code{dynamic\_batching}이 없는 \code{config.pbtxt}, 터미널 기록 9-5의 locustfile로 다섯 단계를 실행해 표로 만든다. 변곡점에서 GPU 사용률을 기록한다. 사례 9-A 근처를 예상한다.
  \item \textbf{⑤ 배치 처리 켜기.} 터미널 기록 9-2의 여섯 줄을 추가하고 다시 스윕한다. 같은 실리콘과 새 표로 자기만의 그림 9-1을 만든다. 두 실행 모두 \code{nvidia-smi dmon}을 화면 캡처한다.
  \item \textbf{⑥ 설정값 하나 실험.} \code{instance\_group count: 1 → 2}로 바꾸고 한 번 다시 스윕한다. 배치 처리는 GPU 비용을 분산하고 인스턴스는 전처리/후처리를 계산과 겹친다(§9.2 기관 ③). \textit{이} 모델에서는 무엇이 더 중요했는지 설명한다.
  \item \textbf{⑦ 재시도 폭풍(§9.5).} 세 번 재시도/백오프 없음으로 변곡점의 1.1×에서 120초 실행하고 도착 요청 대 유효 처리량을 화면 캡처한다. 이어 재시도 예산을 적용한다. 시간 초과 재시도 제거, 지터를 적용한 백오프, Triton의 \code{max\_queue\_size} + \code{default\_timeout}을 통한 대기열 깊이 기반 부하 차단을 적용해 버스트를 반복하고 복구를 화면 캡처한다.
  \item \textbf{⑧ 제안서 갱신.} §7.7에서 추정한 레플리카별 요청률을 ⑤에서 측정한 숫자로 바꾼다. 플릿 규모와 월 청구서를 다시 계산한다(제2장 표). 숫자가 크게 움직였다면 보통 그렇듯 이유를 한 문장으로 밝힌다.
\end{tbbitemize}

\begin{calloutbox}{tbbpurple}{tbbpurplefill}
{\normalsize\bfseries\color{tbbpurple}실습 체크리스트(크레딧과 점수를 모두 지키는 버전)}\par\vspace{3pt}
\begin{tbbitemize}
  \item 동등성 검사 묶음의 출력 저장(①), 버킷에서 저장소 확인(②), 준비 상태 게이팅 시연. 모델 로딩 뒤에야 파드가 준비 상태가 되는 모습을 화면 캡처한다(③).
  \item 행마다 p50/p95/p99 + 오류율 + 유효 처리량을 담은 스윕 표 두 개(④⑤). 첫 1분은 버리고 루프 유형을 \textbf{개방형}(constant\_throughput)으로 표시한다. 폐쇄 루프 숫자는 수정 요청으로 돌려보낸다.
  \item 두 구성의 GPU 사용률 증거, 폭풍 쌍(⑦): 깨진 파형 + 수정한 파형.
  \item 측정한 숫자와 한 줄 변화 설명으로 제안서 갱신(⑧).
  \item 인스턴스 종료, 볼륨 확인, 청구 화면 캡처. 사례 1-B의 의식이 이제 전통이 되었다. 일반적인 실습 비용은 T4 시간 \textasciitilde{}\$2–3다.
\end{tbbitemize}
\end{calloutbox}

\vspace{3.0pt}

\begin{calloutbox}{tbbpurple}{tbbpurplefill}
{\normalsize\bfseries\color{tbbpurple}생각해 보기}\par\vspace{3pt}
\begin{tbbitemize}
  \item ⑤ 전에 예측하자. 사례 9-A의 시간 분해(60 ms CPU / 8 ms GPU / 15 ms 직렬화)를 보면 배치 처리 켜기가 이 워크로드의 처리량을 대략 몇 배로 늘릴까? 배치 처리가 빛을 발하려면 전처리를 어디에 두어야 하며, 답은 §9.4의 어느 프레임워크 철학을 선호하는가?
  \item ③에서 최악의 메모리가 6.1 GiB였고 memory.max = 8 GiB로 설정했다. 다음 달 ⑤에서 max\_batch\_size를 두 배로 늘린다. 아무도 ③을 다시 실행하지 않으면 이어질 실패 연결 고리를 재구성하자. 제1장의 급증, 제3장의 137, 제5장의 QoS가 모두 등장한다.
  \item ⑧에서 측정한 숫자가 §7.7 추정치보다 40\% 낮다. 가장 가능성 높은 원인 세 가지를 확인 비용 순서로 나열하자. 하나는 locustfile, 하나는 config.pbtxt, 하나는 매니페스트의 CPU 상한에 있다. 제5장의 스로틀링 이야기를 떠올리자.
\end{tbbitemize}
\end{calloutbox}

\namedsection{요약}{열두 가지로 정리한 이 장}

\qitem{1}{\textbf{정중한 서버는 우연이 아니라 산술 때문에 정체한다.} 한 번에 하나씩 처리하자 사례 9-A의 상한이 1 ÷ 83 ms ≈ 12 req/s가 되고 GPU는 93\% 놀며 CPU 코어 하나가 고정되었다. 모든 숫자는 제7장의 두 도구로 도출할 수 있다. 240 req/s까지의 간격은 배치 처리, 파이프라이닝, 동시성이라는 엔지니어링이 메운다.}

\qitem{2}{\textbf{프레임워크는 추출한 흉터 조직이다.} 사례 9-B의 git 로그는 수렴 진화다. 모든 팀이 8개월 동안 같은 순서로 같은 기관을 키우므로 TF Serving(2016), Triton(2018/2020), TorchServe(2020), BentoML, Ray Serve의 구조가 같다. 자체 구축 대 외부 도입 질문은 실제로 "어느 실패부터 다시 발견하기를 멈출 것인가?"다.}

\qitem{3}{\textbf{여섯 기관:} 프런트엔드(HTTP/gRPC, GPU 전 검증), 대기열+배처, 인스턴스 그룹(파이프라이닝), 모델 관리(지켜보는 저장소 — 제6장의 S3 선반을 그대로 사용하고 버전을 나란히 두어 13주 차 롤백 연결), 관측 가능성(히스토그램, 대기열 깊이, 배치 크기가 있는 /metrics), 상태 확인/수명 주기(/ready = 모델 예열 완료, SIGTERM에서 드레이닝)다.}

\qitem{4}{\textbf{동적 배치 처리는 실리콘에서 재현한 제2장의 virtio 교훈이다.} 커널 실행 + 가중치 스트리밍이라는 비싼 경계 통과 비용을 여러 단위에 분산한다. 두 설정값은 최악 VRAM이 제한하는 max\_batch\_size(측정해야 하며 제5장의 memory.max 입력이기도 함)와 SLO에서 지연 시간 예산의 작고 알려진 몫으로 정하는 max\_queue\_delay다.}

\qitem{5}{\textbf{절충 관계의 작은 조항:} 처리량은 선형보다 느리게 늘고 꼬리는 더 빠르게 늘어난다. 배치는 작은 팬아웃이다(§7.3). 운영 지점은 벤치마크에서 이긴 배치가 아니라 목표 부하에서 p99가 SLO를 지키는 가장 큰 배치다. 배치 처리에는 동시성이 필요해 간헐적 트래픽은 배치를 만들지 못한다. 배치 크기 히스토그램을 지켜보자.}

\qitem{6}{\textbf{지도에는 두 축이 있다.} 모델 엔진 ↔ 애플리케이션 프레임워크, 설정 기반 ↔ 코드 기반이다. Triton은 모든 §7.6 형식의 백엔드와 앙상블로 GPU를 끝까지 활용하는 설정 우선 엔진이다. TorchServe는 PyTorch 네이티브 .mar + 핸들러이며 2024년 유지 관리 모드 교훈에 따라 프로젝트 맥박을 확인하고 탈출 비용을 낮게 유지해야 한다. BentoML은 코드 우선 패키징 도구로 bento → 컨테이너 이미지를 만들어 제4장 규율을 자동화한다. Ray Serve는 배포 단위를 분할 GPU가 있는 분산 Python 프로그램으로 조합하지만 Ray 운영 비용을 낸다.}

\qitem{7}{\textbf{네 질문으로 선택}한다. 워크로드 형태, 전처리/후처리 위치, 운영 주체, 프로젝트의 생명력 + 탈출 비용을 묻고 \textbf{자기 모델, 자기 SLO, 자기 하드웨어에서 후보를 벤치마크한다.} 공개 막대 차트는 형태만 옮겨 오며 숫자는 절대 그렇지 않다.}

\qitem{8}{\textbf{vLLM과 KServe는 경쟁자가 아니라 이웃이다.} vLLM은 생성 작업이 비슷한 형태라는 가정을 깨뜨릴 때 연속 배치 처리를 제공하는 LLM 전문 엔진(11주 차)이다. KServe는 어느 프레임워크든 Kubernetes에서 실행하는 플랫폼(12주 차)이다. 프레임워크는 레플리카 하나가 잘 서빙하는 방법, 플랫폼은 몇 개 레플리카가 살아야 하는지에 답한다.}

\qitem{9}{\textbf{정직한 Locust = 한 줄 + 세 습관이다.} 개방 루프 도착에는 wait\_time = constant\_throughput(r), 충분한 사용자 수(≥ 2 × λ × worst latency), 예열을 버리는 단계별 스윕, 실패 열이 곧 유효 처리량이 되도록 locustfile 안에서 SLO 강제를 사용한다.}

\qitem{10}{\textbf{스윕은 중요한 숫자 하나}인 SLO에서 레플리카당 req/s(계산 예제 238)를 만든다. 서버 자체의 처리 중 요청 게이지를 상대로 리틀의 법칙으로 감사하고 제안서 용량 절의 추정치를 대체한다.}

\qitem{11}{\textbf{재시도 폭풍의 파형:} 수요가 정상화된 뒤 도착 요청이 \textasciitilde{}3×에 고정되고 유효 처리량은 0 근처 바닥에 붙는다. 용량 문제가 아니라 스스로 유지되는 루프이며 레플리카는 3× 비용으로 증상만 처리한다. 해결책은 제6장의 재시도 예산이다. 시간 초과를 재시도하지 않고, 단일 계층만 재시도하며, 지터를 적용한 백오프를 사용하고, 부하를 일찍 버린다. 느린 어쩌면보다 빠른 거절이 낫다. 카운터 뒤에서 본 429다.}

\qitem{12}{\textbf{실습의 흐름이 곧 이 장의 흐름이다.} 동등성 검사와 함께 패키징 → S3 선반에 보관 → 준비 상태로 게이트한 배포 → 끄기 대 켜기 스윕 → 설정값 하나 → 깨뜨린 뒤 고친 폭풍 하나 → 측정한 숫자와 한 줄 변화 설명으로 제안서 갱신이다.}

\namedsection{연습문제}{}

\subsection*{A. 객관식}

\qitem{1}{사례 9-A(요청마다 직렬 83 ms: 60 CPU + 8 GPU + 15 직렬화)에서 단독으로 처리량 상한을 가장 많이 높이는 변경은?}

\begin{tbboptions}
  \item ① 두 배 빠른 GPU(8 ms → 4 ms)
  \item ② 요청이 겹치도록 전처리를 요청 스레드 밖으로 이동(파이프라이닝/동시성)
  \item ③ RAM 두 배        ④ HTTP 대신 gRPC
\end{tbboptions}

\qitem{2}{모델 서버의 롤링 업데이트를 안전하게 만들어 콜드 레플리카에 트래픽을 보내지 않고 처리 중 작업도 누락하지 않게 하는 기관 쌍은?}

\begin{tbboptions}
  \item ① 프런트엔드 + 배처        ② 메트릭 + 저장소
  \item ③ 준비 상태 게이팅(/ready = 모델 예열 완료) + SIGTERM의 정상 드레이닝
  \item ④ 인스턴스 그룹 + 앙상블
\end{tbboptions}

\qitem{3}{SLO는 p99 ≤ 300 ms이고 배치 계산은 40 ms, 게이트웨이+네트워크 ≈ 50 ms다. 팀원이 "더 큰 배치를 채우기 위해" max\_queue\_delay = 150 ms로 설정했다. 올바른 반론은?}

\begin{tbboptions}
  \item ① 지연 구간은 부하가 클 때 아무 영향이 없으므로 무해하다.
  \item ② 작업 전에 동료를 기다리는 데 지연 시간 예산 절반을 써 가벼운 부하의 최악의 경우만으로도 SLO를 위반할 수 있다.
  \item ③ 150 ms는 Triton 최대값을 넘는다.        ④ 더 큰 배치는 언제나 꼬리에 도움이 된다.
\end{tbboptions}

\qitem{4}{배치 서버의 트래픽이 동시 300 req/s에서 밤새 5 req/s의 간헐적 요청으로 바뀐다. 무슨 일이 생기며 어느 메트릭이 보여 주는가?}

\begin{tbboptions}
  \item ① 아무 일도 없다. 배치 처리는 부하와 무관하다.
  \item ② 배치가 형성되지 않아 요청마다 지연 구간 비용을 내고 배치 크기 1에 가깝게 실행되며 배치 크기 히스토그램이 1 쪽으로 무너진다.
  \item ③ GPU 사용률이 오른다.        ④ 대기열이 넘친다.
\end{tbboptions}

\qitem{5}{공유 GPU에서 여섯 팀의 모델 40개(ONNX, TensorRT, 일부 Python)를 서빙하고 코드보다 설정을 선호하는 SRE가 운영하는 플랫폼 팀에 가장 알맞은 프레임워크는?}

\begin{tbboptions}
  \item ① Ray Serve        ② BentoML        ③ Triton        ④ 팀마다 직접 만든 FastAPI 서버
\end{tbboptions}

\qitem{6}{주기적으로 2 s 멈추는 서버에 Locust 기본 대기 시간(폐쇄 루프)을 사용한다. constant\_throughput 페이싱과 비교해 보고된 p99는?}

\begin{tbboptions}
  \item ① 더 높다. 폐쇄 루프가 멈춤을 과장한다.        ② 충분히 오래 실행하면 같다.
  \item ③ 오해를 부를 만큼 낮다. 사용자가 서버와 함께 멈춰 정지 구간이 거의 표본에 잡히지 않는다(조정 누락, §7.3).
  \item ④ 예측할 수 없다.
\end{tbboptions}

\qitem{7}{사고 중 클라이언트 수요는 10분 전에 정상으로 돌아왔지만 서버 도착 요청은 평소의 3×이고 유효 처리량은 0에 가까우며 CPU와 GPU는 "바빠" 보인다. 진단과 첫 해결책은?}

\begin{tbboptions}
  \item ① 용량 부족 — 즉시 레플리카 추가
  \item ② 재시도 폭풍 — 시간 초과 재시도 중지 / 부하를 일찍 버리기. 레플리카는 사용자가 아니라 증폭을 서빙한다.
  \item ③ 메모리 누수 — 파드 재시작        ④ 네트워크 분할 — 장애 조치
\end{tbboptions}

\qitem{8}{실습 ⑧의 숫자(SLO에서 측정한 레플리카당 req/s)의 주된 목적은?}

\begin{tbboptions}
  \item ① 다른 팀과의 벤치마크 논쟁에서 이기기
  \item ② 제안서의 추정 용량을 측정값으로 교체하기. 플릿 규모, 오토스케일러 목표(12주 차), 15주 차 데모 모두 이 값으로 나눈다.
  \item ③ memory.max 설정        ④ HTTP와 gRPC 가운데 선택
\end{tbboptions}

\subsection*{B. 단답형}

\qitem{1}{동적 배치 처리의 두 설정값과 각각을 정하는 제약 조건을 말하자. 하나는 물리적이고 하나는 계약적이다.}

\qitem{2}{리틀의 법칙 감사에서 Locust는 평균 84 ms에 240 req/s를 보고하지만 Triton의 대기열+실행 게이지 합은 \textasciitilde{}45다. 시스템의 이야기가 일관적인가? 차이가 있을 가능성이 가장 높은 위치는 어디인가?}

\qitem{3}{constant\_throughput(r) 부하 생성의 사용자 프로비저닝 규칙을 제시하고 λ = 200 req/s, r = 2, 시간 초과 2 s일 때 최소 사용자 수를 계산하자.}

\qitem{4}{재시도 폭풍을 막는 재시도 예산 규칙 네 가지를 설명하자.}

\qitem{5}{이 책이 vLLM과 KServe를 "대표 네 가지" 밖에 두는 이유는 무엇인가? 각각 네 프레임워크가 답하지 않는 어떤 질문에 답하는가?}

\subsection*{C. 논술형}

\qitem{1}{사례 9-A를 처음부터 끝까지 재구성하자. 83 ms 분해에서 상한, 사용자 60명의 평균 대기, GPU 사용률, 고정된 CPU 코어를 도출한다. 이어 파이프라이닝, 배치 처리, 인스턴스라는 세 단계 해결책을 설계하고 각 단계가 처리량과 p99를 대략 어디까지 옮길지 계산으로 예측한다. 모든 가정을 밝힌다.}

\qitem{2}{사례 9-B 로드맵의 5주째에 작업자 풀을 막 병합한 엔지니어 4명 팀을 위한 자체 구축 대 외부 도입 메모를 작성하자. 로드맵을 끝내는 비용, 지금 Triton 또는 BentoML로 마이그레이션하는 비용(하나를 골라 그림 9-7 질문으로 정당화), 각각의 위험, 소문보다 마이그레이션을 저렴하게 만드는 아티팩트/기관 사실(§7.6, §9.2)을 포함한다.}

\qitem{3}{자기 프로젝트 모델의 전체 측정 계획을 설계하자. locustfile 개요(페이싱, SLO 강제, 페이로드 현실성), 스윕 일정, 리틀 감사, 수정 검증을 포함한 폭풍 훈련, ⑧이 대체할 제안서의 정확한 문장을 포함한다. CPU 전용 대체 경로(GGUF 또는 eager)가 계획을 바꾸는 지점도 밝힌다.}

\qitem{4}{유지 관리 모드 문제를 다루기 위해 중간 규모 회사의 서빙 프레임워크 도입 정책을 제안하자. 도입 전 기준, 도입 중 결합 규칙, 탈출 조건과 플레이북을 포함하고 이 장의 TorchServe 교훈과 §7.6의 개방형 아티팩트 규율을 근거로 삼는다.}

\namedsection{정답과 해설}{}

\subsection*{A. 객관식}

\qitem{1}{\textbf{②} 상한은 83 ms \textit{직렬} 경로가 결정한다. GPU의 8 ms를 절반으로 줄여도 4 ms만 절약해 → \textasciitilde{}12.7 req/s가 된다. 반면 60 ms 전처리를 계산과 겹치면 임계 경로에서 지배적 항목을 완전히 제거한다. 서빙판 Amdahl의 법칙이다. (§9.1)}

\qitem{2}{\textbf{③} 준비 상태 게이팅은 가중치가 예열될 때까지 트래픽을 막고(제5장의 프로브 적용), 정상 드레이닝은 SIGTERM에서 처리 중인 배치를 완료한다(제3장의 과정). 기관 ⑥이자 사례 9-B의 "장애" 커밋 두 개다. (§9.2)}

\qitem{3}{\textbf{②} 예산은 300 ms이고 서버 전에 50 ms, 계산에 40 ms를 써 \textasciitilde{}210 ms가 남는다. 수집 구간에 150 ms를 쓰면 대기열 지터에 남는 것이 없고 가벼운 부하에서는 모든 요청이 구간 끝까지 기다린다. 지연 구간은 처리량 소망이 아니라 SLO의 작고 알려진 몫으로 설정한다(§9.3).}

\qitem{4}{\textbf{②} 배치 처리에는 동시성이 필요하다. 간헐적 요청은 배치를 만들지 못하므로 요청마다 지연 비용을 내고 배치 크기 1의 효율에 가깝게 실행한다. 배치 크기 히스토그램(기관 ⑤)이 1 쪽으로 무너지는 것이 징후이자 그 메트릭의 존재 이유다. (§9.3)}

\qitem{5}{\textbf{③} 혼합 형식 + 다수 모델 + 운영팀 관리 + 설정 우선은 Triton의 정확한 사분면이다(그림 9-4). 앙상블과 모델 저장소가 여러 팀의 플릿을 처리한다. Ray Serve는 아무도 운영을 요청하지 않은 클러스터를 강제하고 BentoML은 이 팀이 설정을 원하는 곳에 코드를 둔다. (§9.4)}

\qitem{6}{\textbf{③} 폐쇄 루프 사용자가 서버와 함께 멈추는 조정 누락(§7.3, 터미널 기록 7-1) 때문에 정지 구간은 표본에 거의 들어가지 않아 p99가 지나치게 좋아 보인다. constant\_throughput은 도착 요청을 정직하게 유지한다. (§9.5)}

\qitem{7}{\textbf{②} 수요가 정상화된 뒤 도착 요청 ≫ 수요이면서 유효 처리량이 붕괴한 것은 폭풍 징후다(그림 9-6). 레플리카를 추가하면 3× 가격으로 증폭을 먹인다. 시간 초과 재시도 금지와 부하 차단을 비롯한 예산 규칙이 루프를 끊는다. (§9.5)}

\qitem{8}{\textbf{②} 측정값은 제안서 플릿 규모, 12주 차 오토스케일러 목표, 15주 차 채점을 포함한 모든 용량 계산의 씨앗이며 §7.7 추정치를 대체한다. (§9.5–9.6)}

\subsection*{B. 단답형}

\qitem{1}{\textbf{max\_batch\_size}는 최대 배치 × 가장 긴 입력에서 가중치 + 활성화가 차지하는 최악 VRAM이 물리적으로 제한한다. 제5장의 memory.max 입력이기도 하다. \textbf{max\_queue\_delay}는 동료를 모으는 데 의도적으로 쓰는 지연 시간 예산의 작고 알려진 몫으로 SLO가 계약적으로 제한한다. (§9.3)}

\qitem{2}{리틀의 법칙에 따르면 L = 240 × 0.084 ≈ \textbf{처리 중 20개}인데 서버는 \textasciitilde{}45라고 한다. 일관적이지 않아 측정이 잘못되었거나 무언가 숨었다. 가장 가능성 높은 것은 클라이언트에서 \textit{평균}을 측정했지만 서버가 세는 대기열에서 추가 요청 \textasciitilde{}25개가 기다리고 해당 지연 시간에는 포함되지 않은 경우다. 또는 게이지가 대기 중 + 실행 중의 중첩을 이중 계산했을 수 있다. 어느 숫자든 믿기 전에 숨은 대기열을 찾는다. (§7.5, §9.5)}

\qitem{3}{Users ≥ 2 × λ × worst-case latency ÷ ... 실용적으로 사용자마다 요청률 r에서 \textasciitilde{}r × 시간 초과만큼 미완료 요청을 가질 수 있으므로 시간 초과가 쌓일 때 여유를 두고 users ≥ λ/r로 잡는다. λ/r = 최소 100, ×2 여유 → r = 2, 2 s 시간 초과에서 200 req/s에는 \textbf{사용자 ≥ 200명}이다. 사용자를 적게 잡으면 실행이 조용히 폐쇄 루프로 바뀐다. (§9.5)}

\qitem{4}{① 연결 실패만 재시도하고 시간 초과는 절대 재시도하지 않는다. ② 정확히 한 계층만 재시도한다. ③ 지터를 적용한 지수 백오프를 사용한다. ④ 서버가 대기열 깊이 임계값 너머에서 빠른 429/503으로 부하를 일찍 버린다. (§9.5, 제6장)}

\qitem{5}{vLLM은 "길이와 종료 시간이 다른 시퀀스를 가진 \textit{생성} 워크로드를 어떻게 배치하는가?"에 연속 배치 처리 + KV 캐시 관리로 답하는 한 모델 계열 전문 엔진이며 11주 차에 다룬다. KServe는 "이 서버들 가운데 어느 것이든 여러 레플리카가 Kubernetes에서 어떻게 사는가?"에 오토스케일링, 카나리, 플랫폼 배관으로 답하며 12주 차에 다룬다. 세 번째와 네 번째 경쟁자가 아니라 엔진 전문화와 플랫폼 계층이다. (§9.4)}

\subsection*{C. 논술형 채점 기준}

\qitem{1}{상한은 1/0.083 ≈ 12 req/s, W = 60/11.8 ≈ 5.1 s(리틀), GPU 가동 비율은 ≈ 11.8 × 8 ms ≈ 9\%, CPU는 ≈ 11.8 × 60 ms ≈ 0.7 코어에 직렬화 비용을 더한 값이다. 타당한 산술을 곁들인 개선안은 다음과 같다. (a) 전처리를 요청 스레드 또는 서버 밖으로 옮기면 임계 경로가 \textasciitilde{}23 ms로 줄어 \textasciitilde{}40+ req/s, (b) 동적 배치 처리 16을 적용하면 GPU 실행이 \textasciitilde{}35 ms로 분산되어 200–300 req/s, p99는 대략 수집 구간 + 계산 + 지터, (c) 인스턴스를 2로 늘리면 남은 전처리/후처리가 겹쳐 +20–60\%가 된다. 만점을 받으려면 배치 확장 계수와 전처리 목적지 같은 가정을 밝히고, 처리량뿐 아니라 p99의 변화도 설명해야 한다. (§§9.1–9.3)}

\qitem{2}{좋은 답은 로드맵의 비용(5주째→23주째 ≈ 엔지니어 4–5개월분 + 영구적인 운영 책임)을 솔직하게 계산하고, 막연한 선호가 아니라 네 질문을 통해 Triton(향후 혼합 형식, 운영 우선) 또는 BentoML(모두 Python, 개발 경험 우선)을 고른다. 또한 탈출 비용이 낮은 근거를 든다. 아티팩트가 이미 ONNX/safetensors 형식이고(§7.6) 기관의 구조도 같으므로(§9.2), 마이그레이션은 전면 재작성 대신 저장소 배치 + 설정 + 부하 테스트에 가깝다. 양쪽 위험, 즉 프레임워크의 활력(TorchServe 교훈)과 자체 구축의 핵심 인력 의존성 및 아직 발견하지 못한 기관도 비교해야 한다. 2주짜리 사전 검증을 의사 결정 절차로 제안하면 가점한다. (§§9.1, 9.4)}

\qitem{3}{다음을 확인한다. 실제 이미지와 실제 프롬프트 길이를 반영한 페이로드 현실성(배치 형태가 여기에 달려 있다), 프로비저닝 규칙을 계산한 constant\_throughput, 파일 안에서 SLO를 강제해 유효 처리량을 직접 산출하는 방식, 예상 용량의 양쪽을 포괄하는 스윕 단계, 예열 제외, 이름을 밝힌 서버 게이지와의 리틀 감사, 버스트 × 무분별한 재시도로 폭풍을 만든 뒤 예산 적용 전후를 입증하는 훈련, 그리고 제안서에서 대체할 정확한 문장이다. CPU 대체 경로는 요청률만 낮아질 뿐 방법은 같다. 이 절차는 하드웨어에 독립적이다. (§§9.5–9.6)}

\qitem{4}{설득력 있는 정책은 릴리스 주기, 여러 회사에 걸친 유지 관리자, 탈출 가능한 형식을 갖춘 프로젝트만 채택한다. 도입 후에는 비즈니스 로직을 프레임워크 내부에 넣지 않고 아티팩트를 개방형 형식으로 유지하며, 매년 하루 안에 서비스 하나를 대안에 재배포하는 "탈출 훈련"을 한다. 탈출 조건은 유지 관리 모드 징후, 늦어지는 보안 대응, 특정 주체에 사로잡힌 로드맵이다. 각 항목은 일반적인 공급자 위험 문구가 아니라 TorchServe-2024 또는 §7.6/§9.2의 구체적인 작동 원리를 근거로 삼아야 한다. (§9.4)}

\namedsection{더 읽을거리}{}

\textbf{[이번 주]} 표시가 붙은 자료는 실습의 기억이 생생할 때 읽고, \textbf{[N주 차]} 자료는 뒤에서 배울 내용과 연결될 때 다시 펼쳐 보자.

\begin{tbbitemize}
  \item \textbf{[이번 주] Triton Inference Server 문서 — "Model Configuration"과 "Model Analyzer / perf\_analyzer"} — 터미널 기록 9-2/9-3의 공식 설명과 배치/인스턴스 스윕을 자동화하는 도구를 제공한다. 그림 9-3을 곁에 두고 동적 배치 처리 페이지를 읽자.
  \item \textbf{[이번 주] Locust 문서 — wait\_time 전략과 사용자 정의 LoadTestShape} — constant\_throughput의 정확한 의미와 단계별 스윕을 스크립트로 만드는 법을 설명한다. 첫 실행부터 상자 7-1의 규칙을 반영하자.
  \item \textbf{[이번 주] C. Olston 외, "TensorFlow-Serving: Flexible, High-Performance ML Serving"(2017)} — 사례 9-B의 결론을 뒷받침하는 논문이다. Google의 기관을 처음 키운 사람들이 직접 그 구조를 설명한다. 짧고 읽기 쉬우니 여섯 기관을 모두 찾아보자.
  \item \textbf{BentoML과 Ray Serve 문서(아키텍처 + 배치 처리 페이지)} — 그림 9-4에서 코드 우선 쪽 절반을 원자료로 확인할 수 있다. Ray Serve의 "resource allocation" 페이지는 분할 GPU를 §2.6에서 요구한 수준으로 솔직하게 다룬다.
  \item *\textit{M. Nygard, }Release It!\textit{(2판, 2018)의 안정성 패턴 장}* — 재시도 폭풍, 격벽, 부하 차단을 시대를 타지 않는 패턴으로 설명한다. §9.5의 폭풍은 그가 말한 "self-denial attack"이 ML 옷을 입은 모습이다.
  \item \textbf{[10주 차] ONNX Runtime + TensorRT 실행 공급자 문서} — 다음 주에는 같은 실습 서버의 백엔드를 바꾸고 스윕을 다시 실행한다. 재현 가능한 스윕을 만든 이유가 바로 여기에 있다.
  \item \textbf{[11주 차] Yu 외, "Orca: A Distributed Serving System for Transformer-Based Generative Models"(OSDI 2022)} — 상자 9-1의 가정이 깨지는 이유와 vLLM이 필요한 이유를 설명한 연속 배치 처리 논문이다. KV 캐시 강의 뒤에 읽자.
  \item \textbf{[12주 차] KServe 문서(InferenceService)} — 이번 주 Triton 파드가 오토스케일링과 카나리를 갖춘 플랫폼 구성원이 되는 과정을 보여 준다. 지금은 CRD를 훑고 12주 차에 직접 배포한다.
\end{tbbitemize}

\vspace{5.0pt}

\begin{calloutbox}{tbbblue}{tbbbluefill}
{\normalsize\bfseries\color{tbbblue}다음 장 — 제10장: 추론 최적화}\par\vspace{3pt}
이번 주에는 레플리카 하나가 정직하게 서빙하도록 만들었고, 다음 주에는 그것이 \textbf{저렴하게} 서빙하도록 만든다. 제10장은 제7장의 청구서를 \$4,762에서 \$2,381로 낮춘 지렛대를 실제로 움직인다. \textbf{양자화}(INT8/FP8에서 숫자가 무엇을 잃고도 추론이 대개 문제없이 작동하는 이유), \textbf{그래프 컴파일}(ONNX Runtime과 TensorRT로 §7.6의 변환 축을 마침내 벤치마크), 그리고 방금 설정한 배치 처리의 더 깊은 작동 원리를 다룬다. 실습 결과물은 자기 모델의 전후 벤치마크 보고서이며, 당연히 이번 주에 만든 것과 똑같은 Locust 방법을 쓴다. 스윕 표를 가져오자. 곧 보고서의 "이전" 열이 된다.
\end{calloutbox}

%% Chapter 10 — 추론 최적화
\setcounter{chapter}{9}
\chapter{추론 최적화}
\chaptersubtitle{메모리 벽, 양자화, 그래프 컴파일 — 같은 GPU를 절반의 비용으로}

\begin{calloutbox}{tbbblue}{tbbbluefill}
{\normalsize\bfseries\color{tbbblue}학습 목표}\par\vspace{3pt}
\begin{tbbitemize}
  \item 한 장의 4열 비용표를 바탕으로, 서빙에서 "최적화"란 \textbf{같은 답을 요청당 더 낮은 비용으로 제공하는 것}임을 설명한다. 가장 큰 개선에 새 하드웨어도 새 모델도 필요하지 않았던 이유도 이해한다.
  \item 두 줄짜리 산술 집약도 추정으로 워크로드가 \textbf{메모리 벽}의 어느 쪽에 있는지 진단하고, 어떤 최적화가 효과를 내고 어떤 최적화가 기대에 못 미칠지 예측한다.
  \item 스케일, 영점, 채널별 세분성, 보정이라는 양자화의 작동 원리를 설명하고, 학습과 달리 추론이 거친 눈금을 견디는 이유를 이미 익숙한 JPEG 사례와 연결해 이해한다.
  \item \textbf{PTQ 워크플로}를 처음부터 끝까지 실행하고 엔지니어답게 게이트를 적용한다. 작업 지표와 명시된 허용 오차를 비교하고, 슬라이스를 점검하며, 기존 산출물을 제자리에서 바꾸지 않고 다이제스트에 고정된 새 산출물을 만든다.
  \item ONNX Runtime, TensorRT, torch.compile에서 \textbf{그래프 컴파일러}가 제공하는 이점(융합, 폴딩, 레이아웃 계획, 오토튜닝)과 그 대가(적용 범위의 제약, 빌드 시간, 워밍업)를 설명한다.
  \item 컴파일, 양자화, 배치라는 수단을 올바른 순서로 조합하고, 제9장의 t(B) = a + b·B를 다시 적합해 개선의 출처를 입증한 뒤 프로젝트에 필요한 전후 벤치마크 보고서를 완성한다.
\end{tbbitemize}
\end{calloutbox}

\vspace{3.5pt}

\begin{calloutbox}{tbbteal}{tbbtealfill}
{\normalsize\bfseries\color{tbbteal}핵심 용어}\par\vspace{3pt}
요청당 비용 · 메모리 벽 · 산술 집약도 · 루프라인 · 고정 메모리와 PCIe · 양자화 · FP16 / BF16 / FP8 / INT8 / INT4 · 스케일과 영점 · 채널별 / 그룹별 · 보정(최솟값–최댓값, 백분위수, KL) · PTQ 대 QAT · 활성값 이상치 · 게이트(정확도 허용 오차) · 커널 융합 · 상수 폴딩 · 오토튜닝 · ONNX Runtime EP · TensorRT 엔진 · torch.compile · 형상 프로파일 · 패딩과 버킷팅
\end{calloutbox}

\vspace{3.5pt}

\begin{calloutbox}{tbblightgray}{tbbboxfill}
{\normalsize\bfseries\color{tbbink}장 구성}\par\vspace{3pt}
\textbf{10.1} 하나의 모델, 네 가지 비용 — 이미 살아가는 손실 압축의 세계    \textbf{10.2} 메모리 벽: 최적화보다 진단이 먼저다    \textbf{10.3} 양자화: 더 거친 눈금, 정직한 게이트    \textbf{10.4} 그래프 컴파일: 더 적은 왕복과 실행    \textbf{10.5} 수단 조합하기 — a와 b 다시 적합하기    \textbf{10.6} 실습: 전후 벤치마크 보고서 — 이어서 요약, 연습문제, 더 읽을거리
\end{calloutbox}

\section{하나의 모델, 네 가지 비용 — 이미 살아가는 손실 압축의 세계}

제7장은 \$900라는 상한 앞에 월 \$4,762라는 청구서를 내놓고, 그 격차 자체가 앞으로 배울 내용이라고 예고했다. 제9장에서는 하나의 레플리카가 정직하게 동작하도록 만들었다. 이번 장은 그 약속을 처음으로 크게 갚는다. 이번 주를 규정하는 다음 사례부터 살펴보자. 프로젝트의 같은 모델을, 플릿이 임대한 같은 L4와 같은 SLO 및 Locust 스윕 아래 네 가지 방식으로 서빙한 결과다.

\begin{terminalbox}
\begin{Verbatim}[breaklines=true,breakanywhere=true,fontsize=\small,formatcom=\color{tbbtermfg},codes=\tbbcodeactive]
# the project's 3B chat model. same L4, same SLO (p99<=500ms),
# same open-loop sweep (§9.5). only the execution stack changes:
  stack                p99@peak   req/s@SLO   $/M req    vs
  A eager FP32         -- 12 GB weights + KV cache + activations:
                          OOM at peak load on the 24 GB card --
  B eager FP16          492ms       28/s       $7.94     1.0x
  C + ONNX Runtime      455ms       39/s       $5.70     1.4x
  D + TensorRT FP16     448ms       46/s       $4.83     1.6x
  E + TensorRT INT8     461ms       57/s       $3.90     2.0x
# same weights (within a gated tolerance - §10.3), same GPU.
# the fleet: 200 req/s / 57 = 4 replicas = $2,381/mo. as promised.
\end{Verbatim}
\end{terminalbox}
\figcaption{사례 10-A  하나의 모델, 네 가지 비용. B행에서 E행으로 가는 동안 모델의 지식은 전혀 바뀌지 않았지만 비용은 절반이 되었다. A행도 별도의 교훈을 준다. 정밀도는 속도의 문제이기 전에 모델이 메모리에 들어가기나 하는가의 문제다.}

작동 원리를 보기 전에 이 사례가 주는 세 가지 교훈을 읽어 보자. 첫째, 여기서 \textbf{"최적화"란 같은 SLO에서 요청당 비용을 낮추는 것}이다. 실제로는 아무도 서빙하지 않는 배치 512의 과시용 벤치마크가 아니라, §9.5의 정직한 수치를 각 행마다 다시 측정했다. 제1장에서 기법이라면 그래야 한다고 했듯 서빙 삼각형 자체가 \textit{이동}하고 있다. 둘째, \textbf{개선 효과는 합성된다}. C는 컴파일러 패스, D는 더 깊은 컴파일러, E는 컴파일러 \textit{내부}에 양자화를 추가한다. 각 행은 이전 행의 개선을 그대로 이어받는다(§10.5에서 이 합성을 정식화한다). 셋째, \textbf{E행의 별표가 이 규율의 전부다}. INT8의 답은 FP16의 답과 비트 단위로 같지 않다. 하지만 \textit{벤치마크 전에 누군가 명시한 허용 오차 안에서는 동등}하다. 이는 이 장에서 SLO 문장에 대응하는 원칙이자, 엔지니어링과 막연한 기대를 가르는 기준이다(§10.3의 게이트).

\begin{tbbfigure}
  \adjustbox{max width=0.94\linewidth,max totalheight=0.46\textheight}{\includegraphics{figures/fig10-1_fourprices.png}}\\[3pt]
  \figcaption{그림 10-1  사례 10-A를 도식화한 모습. 같은 실리콘에서 달러당 유효 처리량이 2×가 되었다. 제7장이 예고한 플릿의 결과도 그대로 나타난다. 레플리카 8개가 4개로, \$4,762가 \$2,381로 줄었다.}
\end{tbbfigure}

\subsection*{우리는 평생 손실 압축을 믿고 살아왔다}

정밀도를 버리고도 결과물을 유지할 수 있다는 생각은 무모하게 들린다. 그러나 \textbf{우리가 감각으로 소비하는 세계 전체가 이미 이 원리 위에 세워져 있다}. 전화망이 음성을 디지털화했을 때(Bell System의 T1 반송파, 1962), 원음 그대로를 전송하지 않았다. 사람의 귀는 절대적인 간격보다 비율에 민감하므로 로그 눈금(μ-law 컴팬딩)의 \textbf{8-bit} 표본을 보냈다. 연속적인 아날로그 신호를 256단계로 줄였지만, 적절히 배치한 8비트만으로도 음성처럼 들렸다. 통화 속 어머니의 목소리가 잘못 들리는 일은 없었다. 30년 뒤에는 \textbf{JPEG (1992)}가 이미지에 같은 원리를 적용했다. 눈이 거의 감지하지 못하는 고주파 세부를 버려 바이트를 10분의 1로 줄여도 휴가 사진은 멀쩡했다. \textbf{MP3 (1993)}도 귀가 다른 소리에 가려 듣지 못하는 부분을 모델링해 음악에 이를 적용했다. 패턴은 언제나 같다. \textit{신호의 소비자가 실제로 필요한 것을 찾아내고, 나머지에는 비용을 쓰지 않는다.}

양자화는 이 패턴을 신경망에 적용한다. 단, 한 가지 멋진 치환이 있다. 여기서 "소비자"는 눈이나 귀가 아니라 \textbf{모델 자체의 결정}이다. 학습에는 실제로 높은 정밀도가 필요하다. 경사 하강법은 가중치를 10⁻⁵ 단위로 조금씩 움직이고 수백만 단계에 걸쳐 이를 누적한다. 그 작은 변화를 반올림해 없애면 학습이 멈춘다. 반면 추론에서 살아남아야 하는 것은 \textit{출력}뿐이다. 최댓값의 위치, 순위, 생성된 토큰이 유지되면 된다. 학습된 네트워크는 놀라울 만큼 관대한 소비자다. 좋은 최솟값 주변의 곡면은 평평해 모든 가중치를 몇 분의 1퍼센트 움직여도 손실이 거의 변하지 않고, 활성값에는 어차피 잡음이 있으며, 결정 경계에는 여유가 있다. 가중치 0.087231을 0.0872로 반올림하는 변화는 네트워크가 이미 무시하는 잡음보다 훨씬 작다. \textbf{대부분의 추론에서 FP32를 쓰는 것은 전화 통화에 스튜디오 마스터 음원을 보내는 것과 같다.} §10.3에서는 Bell Labs가 했던 만큼 세심하게 μ-law 변환을 수행한다. 신호를 측정하고 눈금을 잘 배치한 다음, \textit{듣는 사람이 차이를 알아채지 못하는지 검증한다}.

작동 원리로 들어가기 전에 이 장 전체를 조직하는 틀을 하나 더 세우자. 제9장은 두 매개변수 법칙 \textbf{t(B) = a + b·B}를 남겼다. a는 패스마다 드는 고정 비용(실행, 스케줄링, 가중치 스트리밍)이고, b는 항목마다 추가되는 한계 비용(연산)이다. 이번 장은 두 상수를 정조준한다. \textbf{컴파일은 a를 줄인다}(§10.4: 커널을 융합하고 상수를 폴딩하며, 활성값이 메모리를 불필요하게 왕복하지 않게 한다). \textbf{양자화는 b를 줄이고 a의 일부도 줄인다}(§10.3: 가중치당 바이트가 절반이면 각 패스의 가중치 스트리밍 부분도 절반이 된다). §10.2는 두 문장에 \textit{바이트}가 계속 등장하는 이유를 설명한다. 대부분의 추론에서 병목은 GPU의 연산 성능이 아니기 때문이다.

\begin{calloutbox}{tbbpurple}{tbbpurplefill}
{\normalsize\bfseries\color{tbbpurple}생각해 보기}\par\vspace{3pt}
\begin{tbbitemize}
  \item 사례 10-A의 A행에서 OOM 계산을 재구성하라. 3B params × 4 bytes = 12 GB의 가중치가 24 GB L4에 올라간다. 최대 부하에서 VRAM을 요구하는 다른 요소(배치 활성값, 채팅 모델의 시퀀스별 KV 캐시 — 제1장의 메모리 이야기와 제9장의 max\_batch\_size 상한)를 차례로 살펴보고, FP16 (6 GB)이 같은 서버에 여유를 준 이유를 설명하라.
  \item μ-law 비유에는 한계가 있다. 전화 통화의 품질은 조금 더 흐릿한 목소리처럼 점진적으로 나빠지지만, 분류기는 어느 순간 결정이 불연속적으로 뒤집힐 수 있다. 바로 이 차이 때문에 존재하는 §10.3의 안전장치는 무엇인가? 방송 엔지니어링에도 이에 해당하는 장치가 있는가?
  \item C행(ONNX Runtime)은 정밀도를 바꾸지도 양자화하지도 않고 1.4×의 개선을 얻었다. 제9장에서 배운 a와 b를 바탕으로 볼 때 이 개선은 어디에서 나왔어야 하는가? §10.4의 여러 메커니즘 중 가장 유력한 것은 무엇인가?
\end{tbbitemize}
\end{calloutbox}

\section{메모리 벽: 최적화보다 진단이 먼저다}

이 장의 모든 최적화 이야기는 현대 가속기의 화려하지 않은 사실 하나에 기대고 있다. \textbf{연산 속도가 빨라진 속도를 바이트가 가까워지는 속도가 따라가지 못했다.} T4는 초당 65조 회의 FP16 연산을 수행할 수 있지만, 같은 시간에 메모리가 공급할 수 있는 데이터는 320기가바이트뿐이다. 둘을 나누면 이 장치의 성격이 드러난다. 연산 장치가 굶지 않으려면 가져오는 \textbf{바이트당 약 200회의 부동소수점 연산}이 필요하다. 이 비율이 \textit{메모리 벽}이다. 이 이름은 1995년 Wulf와 McKee가 CPU 속도는 해마다 \textasciitilde{}50\%씩 늘지만 메모리 접근 성능은 \textasciitilde{}7\%만 개선된다는 사실을 발견하고, 업계가 전속력으로 벽을 향해 달리고 있다고 유명한 짧은 논문에서 경고하며 붙였다. GPU는 막대한 대역폭으로 충돌을 늦췄지만(HBM은 메모리 벽에 바치는 몸값이다), 그 비율은 계속 악화되었다. 워크로드가 벽의 어느 쪽에 있는지를 알면 벤치마크를 하기도 전에 어떤 최적화가 효과를 낼지 판단할 수 있다.

진단에 필요한 것은 하나의 수치다. \textbf{산술 집약도(AI), 즉 이동한 바이트당 수행한 FLOP 수}다. AI가 장치의 비율보다 큰 워크로드는 \textbf{연산 병목}이다. 다이가 실제로 바쁘므로 연산 자체를 더 저렴하게 만드는 것만 도움이 된다. 반대로 그 비율보다 낮으면 \textbf{메모리 병목}이다. 텐서가 흘러들어오는 동안 연산 장치는 쉬고 있으므로 "더 빠른 연산"은 아무것도 바꾸지 못하며, 오직 \textit{더 적은 바이트}만 효과가 있다. 그림 10-2는 이를 고전적인 \textbf{루프라인(roofline)}으로 나타낸다. 기울어진 대역폭 지붕과 평평한 연산 지붕이 있고, 둘이 만나는 능선이 있다. 서빙에는 불편한 소식이 있다. \textbf{서빙에 적합한 배치 크기에서 대부분의 추론은 능선 왼쪽}, 즉 메모리 병목 영역에 놓인다. 배치 1의 LLM 토큰 생성은 극단적인 사례다. 토큰 하나마다 모든 가중치를 다시 스트리밍해 각각 몇 번의 FLOP만 수행하므로 AI ≈ 1–10에 불과하고, 대역폭 병목에서 벗어날 가망이 없다. 이것이 제11장 전체의 난제다. 큰 배치의 CNN과 규모가 큰 GEMM만 연산 지붕에 도달한다.

\begin{tbbfigure}
  \adjustbox{max width=0.94\linewidth,max totalheight=0.46\textheight}{\includegraphics{figures/fig10-2_memorywall.png}}\\[3pt]
  \figcaption{그림 10-2  루프라인으로 표현한 메모리 벽. 배치는 오른쪽으로 이동시킨다(같은 가중치로 바이트당 더 많은 작업을 수행한다). 양자화는 기울어진 지붕을 들어 올린다(가중치당 바이트가 줄어든다). 컴파일은 불필요한 왕복을 없앤다. 먼저 진단하라. 무엇이 효과를 낼지는 능선이 결정한다.}
\end{tbbfigure}

\subsection*{능선 적용하기: 이미 다뤄 본 세 가지 워크로드}

2주 차부터 실행해 온 세 연산의 AI를 직접 계산하면 이 정의는 곧 직관이 된다. 숫자를 모두 펼쳐서 계산해 보자. \textbf{첫째, 토큰을 생성하는 행렬–벡터 곱이다.} 어떤 트랜스포머 계층의 FP16 가중치 행렬이 4,096 × 4,096이라면 크기는 33.6 MB다. 입력 벡터 하나를 곱할 때 2 × 4,096² ≈ 33.6 MFLOP을 수행한다. 둘을 나누면 \textbf{AI ≈ 바이트당 1 FLOP}이다. 가중치를 한 번 읽어 한 번 쓴다. 능선이 200인 장치에서 이는 워크로드라기보다 약간의 연산을 곁들인 바이트 배달 서비스다. 이제 B개의 요청을 배치로 묶자. 같은 33.6 MB의 가중치가 B개의 벡터를 처리하므로 FLOP은 B배가 되지만 바이트는 거의 변하지 않는다. 따라서 \textbf{AI ≈ B}다. 이 한 줄이 배치 처리의 이론 전체이며, 뒤에 나올 표 10-3의 핵심이기도 하다. T4의 능선에 닿으려면 B ≈ 200이어야 한다. 하지만 SLO 때문에 이 모델의 배치는 약 32가 상한이므로 생성은 늘 메모리 병목 영역에 \textit{머문다}. 제11장의 첫 문제가 여기서는 산술로 드러난다.

\textbf{둘째, ReLU}다. max(0, x)는 원소마다 한 번의 FLOP을 수행한다. FP16 원소 하나를 읽고(2 bytes) 써야 하므로(2 bytes) \textbf{AI = 1 ÷ 4 = 0.25}이며, 능선보다 800배 낮다. 지금껏 만들어진 어떤 하드웨어도 원소별 연산을 메모리 속도보다 빠르게 실행하지 못한다. 그래서 §10.4의 컴파일러는 ReLU에 별도 커널을 주지 않는다. 값을 만드는 연산과 융합하면 그 바이트 자체가 사라진다. \textbf{셋째, 실제 합성곱}을 보자. 56×56 특징 맵에서 3×3, 64→64 채널인 연산이다. FLOP은 2 × 9 × 64 × 64 × 56² ≈ 231 M이다. 바이트는 입력 0.40 MB + 가중치 0.07 MB + 출력 0.40 MB ≈ 0.88 MB다. 따라서 \textbf{AI ≈ 264}로 T4의 능선 \textit{바로 위}에 놓인다. 합성곱은 3,136개 공간 위치 전체에서 각 가중치를 재사용한다. 데이터 재사용이 연산 자체에 내장되어 있다. 이 때문에 CNN은 루프라인의 모범 시민이고, 9주 차 분류기의 개선은 가중치 스트리밍보다 연산 속도와 실행 오버헤드에서 나온다. 같은 GPU가 한 모델에서는 굶고 다른 모델에서는 땀을 흘리는 이유도 여기에 있다.

세 수치 \textbf{≈ B, 0.25, 264}를 머릿속에 두자. 앞으로 서빙할 거의 모든 계층은 이 셋 중 하나가 다른 옷을 입은 모습이다. 첫 진단에는 프로파일러도 필요 없다. 두 줄의 계산이면 충분하며, 실습 보고서도 바로 이 두 줄로 시작해야 한다.

\begin{terminalbox}
\begin{Verbatim}[breaklines=true,breakanywhere=true,fontsize=\small,formatcom=\color{tbbtermfg},codes=\tbbcodeactive]
# is my model memory-bound at this batch size? no profiler needed.
# the project's 3B chat model, on the fleet's L4 (24 GB, 300 GB/s):
# floor set by BYTES: weights must stream once per forward pass:
$ python -c "print(6.0e9 / 300e9 * 1e3, 'ms floor per pass')"
20.0 ms floor per pass      # 3B params x 2B (FP16) / 300 GB/s
# floor set by MATH: ~2 FLOPs/param/item -> batch 1: 6 GFLOP:
$ python -c "print(6.0e9 * 2 / 121e12 * 1e3, 'ms of compute')"
0.1 ms of compute           # vs 20.0 ms of streaming: 200x apart
# verdict: memory-bound by two orders of magnitude at batch 1.
# bytes are the bill. (batch 32 closes the gap ~32x - still bound.)
\end{Verbatim}
\end{terminalbox}
\figcaption{터미널 기록 10-1  두 줄짜리 진단. 배치 1에서 이 모델은 가중치로 연산하는 시간보다 가중치를 스트리밍하는 데 200× 더 오래 걸린다. 다이는 정체 구간에서 공회전하는 페라리와 같다. 이후의 모든 최적화 결정은 이 비율에서 출발한다.}

이제 메모리 벽을 염두에 두고 제9장의 두 상수를 다시 읽어 보자. 이 작업은 이 책이 계속 사용하는 두 모델 \textit{모두}에 해야 한다. 둘은 벽의 반대편에 있으며, 그 대비를 읽는 것이 진단의 핵심이다. \textbf{3B 채팅 모델}의 고정 비용 \textbf{a}는 가중치 스트리밍이 지배한다. FP16에서는 패스마다 피할 수 없는 바이트에 20 ms가 든다. 어떤 컴파일러도 이 하한을 낮출 수 없지만 양자화는 절반으로 줄일 수 있다(INT8 → \textasciitilde{}10 ms). 반면 \textbf{9주 차 분류기}의 가중치는 \textasciitilde{}50 MB로 스트리밍에 0.16 ms밖에 들지 않아 반올림 오차에 가깝다. 따라서 a = 8.4 ms는 거의 전부 \textit{실행, 스케줄링, 연산별 오버헤드}다. 컴파일은 바로 이런 낭비를 없애기 위해 존재한다(§10.4). 양자화의 이점은 주로 \textbf{b}에 나타난다(INT8 연산 속도와 더 가벼운 활성값 트래픽). 배치는 스트리밍한 가중치 하나를 B개 항목에 재사용해 AI를 높였고, 루프라인에서 오른쪽으로 이동시켰다. 이것이 제9장의 거래였다. 이제 이 장의 두 수단이 맡을 역할은 이미 보인다. \textbf{양자화는 바이트를 공격하고} 기울어진 지붕을 들어 올린다. \textbf{컴파일은 낭비를 공격해} 실행을 줄이고 활성값이 VRAM을 왕복하지 않고 레지스터에 머물게 한다. 최적화보다 진단이 먼저라는 말은 상투적인 구호가 아니다. 메모리 병목 모델에 "2× 빠른 연산"을 배포한 팀이 3\% 개선만 측정하고 오후 내내 혼란에 빠지는 이유를 설명하는 원칙이다.

\begin{tbbtablewrap}
  \begin{tabular}{@{}>{\raggedright\arraybackslash}m{\dimexpr 0.2021\tbltextwidth-2\tabcolsep\relax} >{\raggedright\arraybackslash}m{\dimexpr 0.1862\tbltextwidth-2\tabcolsep\relax} >{\raggedright\arraybackslash}m{\dimexpr 0.1809\tbltextwidth-2\tabcolsep\relax} >{\raggedright\arraybackslash}m{\dimexpr 0.1915\tbltextwidth-2\tabcolsep\relax} >{\raggedright\arraybackslash}m{\dimexpr 0.2394\tbltextwidth-2\tabcolsep\relax}@{}}
  \arrayrulecolor{tbbnavy}\specialrule{1pt}{0pt}{0pt}
  \rowcolor{tbbtablehead}{\centering \bfseries\color{tbbnavy}GPU (VRAM)} & {\centering \bfseries\color{tbbnavy}FP16 텐서 연산} & {\centering \bfseries\color{tbbnavy}메모리 대역폭} & {\centering \bfseries\color{tbbnavy}능선(FLOPs/byte)} & {\centering \bfseries\color{tbbnavy}이 책에서의 역할} \\
  \arrayrulecolor{tbbnavy}\specialrule{0.7pt}{0pt}{2pt}
  {\textbf{T4 (16 GB)}} & {65 TFLOPS} & {320 GB/s} & {≈ 200} & {실습용 카드 — 스윕, 2× 속도의 INT8 (130 TOPS)} \\
  \arrayrulecolor{tbbtableborder}\hline
  {\textbf{L4 (24 GB)}} & {121 TFLOPS} & {300 GB/s} & {≈ 400} & {플릿 카드 — 사례 10-A, \$0.80/hr의 비용} \\
  \arrayrulecolor{tbbtableborder}\hline
  {\textbf{A100 (80 GB)}} & {312 TFLOPS} & {2,039 GB/s} & {≈ 153} & {생각해 볼 문제의 카드 — 프리미엄의 원천은 대역폭} \\
  \arrayrulecolor{tbbtableborder}\hline
  {\textbf{H100 (80 GB)}} & {\textasciitilde{}990 TFLOPS} & {3,350 GB/s} & {≈ 295} & {최첨단 기준 — 11주 차의 FP8이 동작하는 곳} \\
  \arrayrulecolor{tbbnavy}\specialrule{1pt}{2pt}{0pt}
  \end{tabular}\par
  \figcaption{표 10-1  비율로 읽는 이 책의 장치들(공개된 사양의 근삿값. 기억해야 할 수치는 능선 = TFLOPS ÷ 대역폭이다)}
\end{tbbtablewrap}

표를 비율로 읽으면 세 가지 교훈이 나온다. 첫째, \textbf{능선은 자연 상수가 아니다.} 이는 설계 선택이며, 당장 빌릴 수 있는 카드 사이에서도 2.5×나 차이 난다. L4가 가장 극명하다. T4보다 연산 성능을 두 배로 높이면서도 대역폭은 \textit{낮춰} 능선을 약 400에 놓았다. 건강한 배치 크기에서 INT8을 공급하도록 만든 카드이며, 배치 1 FP16 생성에서는 예상대로 고전한다. 사례 10-A의 B행을 실리콘 관점에서 설명한 셈이다. 둘째, \textbf{A100의 프리미엄은 대부분 대역폭에서 나온다.} T4보다 연산은 \textasciitilde{}5×지만 바이트는 \textasciitilde{}6.4×다. 터미널 기록 10-1을 A100에서 다시 계산하면 연산 하한은 여전히 별 의미가 없는 반면 스트리밍 하한은 20 ms에서 2.9 ms로 떨어진다. 메모리 병목 추론에서 비싼 GPU란 연산을 곁들인 비싼 \textit{메모리}다. 셋째, 세대가 바뀔수록 능선은 \textit{위로} 이동하고 벽은 계속 높아진다. 따라서 \textbf{양자화의 가치는 하드웨어 세대가 바뀔수록 줄어드는 것이 아니라 더 커진다.} 바이트는 가치가 오르는 통화다. Hopper급 카드가 FP8을 일급 형식으로 추가한 이유도 여기에 있다. 하드웨어 자체가 더 적은 바이트를 보내 달라고 요구하며, §10.3은 그 요구에 응한다.

\subsection*{또 하나의 벽: PCIe와 바이트가 먼저 거치는 여정}

루프라인의 300 GB/s는 VRAM에서 연산 장치로 이어지는 \textit{안쪽} 벽이다. 바깥쪽에도 벽이 하나 있다. 제6장은 바로 그 앞에 독자를 세워 두었다. 가중치는 객체 스토어에서 \textbf{호스트 페이지 캐시}로 도착했다(mmap한 safetensors를 처음 건드릴 때 커널이 텐서를 페이지 인한다). 이번 장에서 그 여정을 마치겠다고 약속했다. 마지막 구간은 \textbf{PCIe를 통한 호스트-장치(H2D) 복사}다. PCIe는 장치 안에서 가장 느린 도로다. gen3 ×16 슬롯의 이론적 속도는 \textasciitilde{}16 GB/s지만, 요청 방식에 따라 실제 처리량은 그 절반이나 3분의 1에 불과하다. 제3장에서 본 메모리 메커니즘 때문에 요청 방식이 중요하다. GPU의 DMA 엔진은 \textbf{고정된(page-locked) 메모리}에서만 데이터를 가져올 수 있다. 이는 커널이 스왑하거나 이동하지 않겠다고 보장한 페이지다. 일반적인 페이지 가능 메모리(pageable memory)에서 복사하면 CUDA는 모든 데이터를 내부의 고정 바운스 버퍼(bounce buffer)에 조용히 옮긴다. memcpy가 한 번 더 생기고 처리량은 대략 절반이 된다. 버퍼를 직접 고정하면 DMA 엔진이 바로 읽을 수 있다.

\begin{terminalbox}
\begin{Verbatim}[breaklines=true,breakanywhere=true,fontsize=\small,formatcom=\color{tbbtermfg},codes=\tbbcodeactive]
# how fast can bytes actually get INTO the GPU? (lab T4, PCIe gen3 x16)
$ python h2d_bench.py --mb 256
  pageable host -> device    6.3 GB/s   # malloc'd memory: CUDA stages
                                        # it through a pinned bounce buffer
  pinned   host -> device   12.6 GB/s   # page-locked: DMA reads directly
  VRAM     device internal 298.1 GB/s   # the inner wall, for scale: ~25x
# lesson 1: PCIe is ~25x slower than VRAM. weights cross ONCE, at startup,
#           and live in VRAM. nothing model-sized crosses per request.
# lesson 2: pin what does cross (DataLoader/serving: pin_memory=True),
#           and overlap copy with compute (CUDA streams - free otherwise).
\end{Verbatim}
\end{terminalbox}
\figcaption{터미널 기록 10-2  측정한 바깥쪽 벽. 페이지 가능 메모리 복사는 숨겨진 스테이징 memcpy 비용을 내고, 고정 메모리는 H2D 처리량을 두 배로 높인다. 하지만 둘 다 VRAM 대역폭과 비교하면 미미하다. 그래서 전송 속도가 아니라 상주가 설계 원칙이 된다.}

이제 제2장 실습에서 관찰한 현상을 아래 계층에서 설명할 수 있다. 배치 1 GPU 벤치마크가 기대에 못 미친 이유는 작은 작업이 PCIe 이동과 실행 오버헤드를 어느 쪽도 상각하지 못한 채 모두 부담했기 때문이다. 또한 모델의 여정은 세 가지 속도로 이루어진 공급망으로 완성된다. \textbf{객체 스토어 → 호스트 페이지 캐시(GB/s, 파드마다 한 번, §6.4) → VRAM(고정 메모리에서 12 GB/s, 프로세스마다 한 번) → 연산 장치(300 GB/s, 모든 패스마다)}다. 각 단계는 이전 단계보다 대략 한 자릿수만큼 빠르고, 빠른 단계일수록 건너는 횟수는 비례해 줄어든다. 설계 원칙은 자연스럽게 도출된다. 가중치는 시작 시점, 즉 제5장에서 예산을 잡은 콜드 스타트 경로에서 PCIe를 정확히 한 번 건넌다. 요청별 텐서는 작지만(이미지 \textasciitilde{}150 KB, 프롬프트 몇 KB로 12 GB/s에서는 수십 마이크로초) 반드시 고정하고 비동기로 처리해야 직렬화를 피할 수 있다. 결과는 반대 방향으로 돌아온다. 서빙 프레임워크(제9장의 기관 ⑤)는 고정과 스트림 중첩을 대신 처리한다. 그 작동 원리를 알아야 하는 까닭은 제9장 §9.4의 교훈 때문이다. 프레임워크의 기본값이 잘못되었을 때 두 줄짜리 진단에는 이제 세 번째 줄이 추가되며, \code{nvidia-smi dmon}의 PCIe 열을 확인해야 한다.

\begin{calloutbox}{tbbpurple}{tbbpurplefill}
{\normalsize\bfseries\color{tbbpurple}생각해 보기}\par\vspace{3pt}
\begin{tbbitemize}
  \item 같은 3B 모델을 서빙하는 A100-80GB (312 TFLOPS FP16, \textasciitilde{}2,000 GB/s)에 대해 터미널 기록 10-1의 두 줄을 계산하라. 어느 하한이 더 크게 움직이는가? 추론에서 고가 GPU의 프리미엄이 어디에서 나오는지를 설명하라.
  \item 팀원이 배치 4 LLM 엔드포인트의 지연 시간을 해결하려고 플릿을 "TFLOPS가 4×"이지만 대역폭은 비슷한 GPU로 업그레이드하자고 제안한다. 메모리 벽을 이용해 결과를 한 문장으로 예측하고, 실제로 수치를 움직일 두 가지 구매 항목을 제시하라.
  \item 배치 처리는 왜 "지붕을 들어 올리는" 대신 "오른쪽으로 이동"시키는가? B가 1에서 16으로 바뀔 때 무엇이 바이트와 FLOP 관점에서 변하며, 무엇은 변하지 않는지 서술하라. 이는 제9장의 a + b·B를 다른 방식으로 표현한 것이다.
  \item 비디오 추론 서비스가 요청마다 90 MB의 디코딩된 프레임을 보낸다. 터미널 기록 10-2의 수치로 페이지 가능 메모리와 고정 메모리(pinned memory)의 H2D 시간을 각각 계산하고, 모델의 11 ms 순방향 패스와 비교하라. 모델을 건드리기 전에 적용할 §10.2의 두 가지 해결책도 제시하라.
\end{tbbitemize}
\end{calloutbox}

\section{양자화: 더 거친 눈금, 정직한 게이트}

양자화의 작동 원리는 소박하다. \textbf{텐서를 눈금이 더 적은 자로 다시 표현하는 것}이다. FP32 가중치는 넓은 동적 범위에 퍼진 \textasciitilde{}40억 개의 표현 가능한 값 중 하나다. 하지만 특정 계층의 가중치는 작은 구간에 모여 있다. 흔히 ±0.4 정도이며, 그림 10-3의 종 모양 분포가 이를 보여 준다. INT8 양자화는 바로 그 구간 위에 256개 눈금의 자를 놓는다. 식은 \textbf{q = round(w / scale) + zero\_point}이고, 가중치마다 1바이트와 소수의 스케일 매개변수를 저장한다. 연산할 때 다시 곱해 역양자화할 수도 있다. 더 좋은 방법은 하드웨어가 \textbf{INT8로 직접 연산하는 것}이다. 이 때문에 단순히 저장 공간만 줄지 않고 이 장의 수치가 두 배로 좋아진다. T4의 텐서 코어는 INT8을 FP16보다 두 배 빠른 130 TOPS로 실행한다. 따라서 터미널 기록 10-1의 두 하한이 동시에 낮아진다. 스트리밍할 바이트는 절반이 되고 연산 처리량은 두 배가 된다.

\begin{tbbfigure}
  \adjustbox{max width=0.94\linewidth,max totalheight=0.46\textheight}{\includegraphics{figures/fig10-3_quantmap.png}}\\[3pt]
  \figcaption{그림 10-3  양자화는 텐서가 실제로 사용하는 범위 위에 더 거친 자를 놓는 작업이다. 기술의 핵심은 눈금을 배치하는 방식(스케일, 채널별)에 있다. 주의할 점은 채널마다 범위가 다르고 트랜스포머 활성값에 이상치가 숨어 있다는 것이다.}
\end{tbbfigure}

\subsection*{자 자체의 차이: 부동소수점은 관대하고 정수는 그렇지 않은 이유}

눈금을 놓기 전에 눈금의 재료부터 살펴보자. 이 절의 어려움은 하나의 구조적 차이에서 비롯된다. \textbf{부동소수점 수는 자기만의 자를 지닌다.} 지수 비트가 각 값의 스케일을 개별적으로 정하므로 0.001 근처의 수와 1,000 근처의 수가 모두 가수부의 전체 해상도를 얻는다. 지난 5주간 실습하면서 범위를 한 번도 고민하지 않아도 되었던 이유다. \textbf{정수에는 지수가 없다.} 하나의 고정된 자 위에 놓인 위치일 뿐이므로, \textit{다른 누군가}가 텐서 전체에 적용할 자의 범위를 한 번 정해야 한다. 그 누군가가 바로 보정이다. 아래에서 볼 모든 함정, 즉 서로 다른 채널과 범위를 늘리는 이상치(outlier)는 많은 값이 하나의 스케일을 함께 써야 한다는 데서 생긴다. 부동소수점은 관대하지만, 정수는 그 관대함을 사용자에게 위임한다.

부동소수점 계열은 30초쯤 들여 살펴볼 가치가 있다. 이제 모든 GPU 사양표에 등장하기 때문이다. \textbf{FP32} (부호 1비트 · 지수 8비트 · 가수 23비트)는 기준이 되는 자다. \textbf{FP16} (1·5·10)은 바이트를 절반으로 줄이는 대신 \textit{범위}도 줄여 최댓값이 65,504다. FP16 학습에서 작은 경사가 0으로 사라지지 않게 손실 스케일링 기법이 필요한 이유다. \textbf{BF16} (1·8·7)은 반대의 절충을 택한다. FP32의 지수 범위를 그대로 유지하고 거친 가수부(십진수 약 3자리)를 대가로 치른다. 경사가 오버플로하지 않아 학습에서 각광받지만, \textit{서빙}에서는 FP16보다 이득인 경우가 드물다. 둘 다 2바이트이고 대부분의 GPU에서 속도도 같다. \textbf{FP8}은 Hopper/Ada 하드웨어에서 지원한다. H100과 플릿의 L4에도 관련 회로가 있다. E4M3 (±448, 가중치와 활성값용)와 E5M2 (±57,344, 경사용)의 두 형식이 있으며, \textit{정수 수준의 저장 비용에 약간의 자체 스케일링을 남긴 형식}으로 이해하면 좋다. 실제로는 여전히 보정하지만 INT8보다 불완전한 스케일을 훨씬 잘 견딘다. \textbf{INT8/INT4}는 이 장에서 다루는 맨몸의 자다. 가장 저렴하고 빠르지만 아래에서 설명할 눈금 배치 기술에 전적으로 의존한다.

반복되는 혼란을 막아 줄 구현상의 각주가 하나 있다. INT8 계층 안에서 INT8 × INT8의 곱은 \textbf{INT32에 누적}된다. 수천 개 항을 오버플로 없이 더한 뒤, 결과가 나갈 때 출력 스케일로 \textit{재양자화}한다. 따라서 양자화 모델의 정확도 문제는 사실상 "누산기가 오버플로했다"는 데서 생기지 않는다. 언제나 "이 텐서에 맞지 않게 눈금을 배치했다"는 문제다. 이에 맞게 진단해야 한다. 해결책은 연산 폭이 아니라 스케일과 보정(calibration)에 있다.

\subsection*{눈금 배치하기: 스케일, 세분성, 보정}

실제 수치로 작동 원리를 한 번 살펴보자. 네 줄이면 충분하며 이후의 모든 약어에서 신비를 걷어 낼 수 있다. 가중치 범위가 [−0.42, +0.38]인 계층을 생각하자. 대칭 INT8은 −127…+127의 눈금을 사용하므로 \textbf{scale = max|w| ÷ 127 = 0.42 ÷ 127 ≈ 0.00331}이고 zero\_point = 0이다. 가중치 0.087231은 q = round(0.087231 ÷ 0.00331) = \textbf{26}이 되어 1바이트에 저장된다. 연산할 때 역양자화하면 26 × 0.00331 = 0.08606이다. 모두가 논하는 "손실"이란 약 0.0012, 즉 눈금 하나의 3분의 1에 해당하는 이 반올림 오차가 전부다. 이는 네트워크가 이미 견디는 잡음 하한 \textit{아래}에 있다. §10.1에서 본 평평한 최솟값 논증을 수치로 확인한 셈이다. 25 MB짜리 FP32 계층 전체는 \textasciitilde{}6.3 MB와 몇 개의 부동소수점 스케일로 바뀐다. 양자화란 \textit{그게 전부}다. 이 절의 나머지 내용인 세분성, 보정, 이상치, 게이트는 이 네 줄짜리 이야기가 어긋나는 경우를 다룬다.

앞의 예에서는 가중치가 0을 중심으로 양쪽에 걸쳐 있어 \textit{대칭} 자(zero\_point = 0)를 사용했다. 활성값은 그렇지 않은 경우가 많다. ReLU를 지난 계층의 출력 범위가 \textbf{[0, 6.2]}라면 대칭 자는 256개 눈금의 절반을 절대 나타나지 않는 음수에 낭비한다. 그 결과 아무 이유 없이 간격이 두 배가 된다(6.2 ÷ 127 = 0.049). \textbf{zero\_point}는 바로 이 문제를 해결한다. \textit{비대칭} 자는 scale = 6.2 ÷ 255 = 0.024, zero\_point = −128로 설정해 실수 0.0을 코드 −128에 대응시킨다. 이제 256개 눈금 모두가 실제 값의 범위를 덮고, 간격은 절반이 된다. q = round(w/scale) + zero\_point의 알쏭달쏭한 두 번째 매개변수는 자를 옮겨 어느 부분도 낭비하지 않게 하는 역할을 한다.

양자화의 모든 기술은 눈금 배치로 귀결된다. \textbf{텐서마다 하나의 스케일}을 쓰는 방식이 가장 둔감하며, 계산해 보면 그 정도가 드러난다. 채널 7의 가중치 범위는 ±2.0이고 채널 300은 ±0.05라고 하자. 공유 스케일은 채널 7을 수용해야 하므로 scale = 2.0 ÷ 127 ≈ 0.0157이다. 그러면 채널 300의 \textit{전체 범위}가 약 세 눈금 안에 갇힌다. 전형적인 가중치 0.02는 q = round(0.02 ÷ 0.0157) = 1로 양자화되고 0.0157로 역양자화된다. 모든 가중치에서 체계적으로 \textbf{21\% 오차}가 난다. 채널 300에 자체 자(scale = 0.05 ÷ 127 ≈ 0.00039)를 주면 같은 가중치가 q = 51 → 0.0201이 되어 \textbf{0.4\% 오차}에 그친다. 채널마다 부동소수점 하나만 더 저장하면 오차를 50분의 1로 줄일 수 있다. 그래서 \textbf{채널별 스케일}, 즉 출력 채널마다 하나의 자를 두는 방식이 가중치의 보편적인 표준이다. 4비트 LLM 양자화는 \textbf{그룹별 스케일}로 더 세분화한다. 가중치 128개마다 하나를 두는 방식이며, 11주 차의 GPTQ/AWQ 용어를 여기서 미리 만난다. 범위는 어디서 구할까? 가중치는 알려진 상수이므로 파일에서 읽으면 된다. \textit{활성값}은 누군가 관찰해야 한다. \textbf{보정}은 대표 입력 수백 개를 모델에 통과시키고 계층마다 활성값 범위를 기록해 그에 맞는 스케일을 고른다. 이것이 \textbf{학습 후 양자화(PTQ)}다. 재학습이 필요 없고 CI에서 몇 분이면 끝난다. 품질은 보정 세트의 품질과 정확히 같다. 이미지 모델을 낮 사진만으로 보정하면 밤 입력은 클리핑된다. 대안인 \textbf{양자화 인식 학습(QAT)}은 미세 조정 중 반올림을 모사해 네트워크가 그 영향을 피해 학습하게 한다. PTQ가 잃은 정확도를 되찾지만 학습 파이프라인만큼 비용이 든다. 이 책은 PTQ를 사용하지만, 언제 그 너머의 방법이 필요한지는 알아야 한다. 그림 10-4를 참고하자.

보정은 관찰한 활성값 더미를 어떻게 하나의 스케일로 바꿀까? 정교함이 높아지는 순서로 세 가지 방법이 있으며, 그 차이는 결코 이론에 그치지 않는다. \textbf{최솟값–최댓값}은 실제 극값을 택한다. 정직하지만 취약하다. 활성값 999개가 6.0 아래에 있는데 하나의 특이한 값만 41.0에 이르면 스케일은 41 ÷ 127 ≈ 0.32가 된다. 실제 값 전체가 사용 가능한 127개 중 약 \textbf{열아홉 눈금}(6 ÷ 0.32)에 몰린다. 이상한 입력 하나가 영구적으로 해상도 3비트를 앗아 간 셈이다. \textbf{백분위수} 방식은 예컨대 99.99번째 백분위수를 기준으로 이상치를 클리핑한다. scale = 6.3 ÷ 127 ≈ 0.05가 되어 일반 값은 자 전체를 되찾고, 드문 41.0은 맨 위 눈금에서 포화된다. 드문 값에 유한한 오차를 허용하는 대신 나머지 모든 값의 해상도를 얻는다. 거의 언제나 올바른 절충이다. \textbf{엔트로피/KL 보정}은 TensorRT의 기본값이며 이 절충을 자동화한다. 후보 클리핑 임곗값을 훑어 FP32 분포와 그 INT8 그림자 사이의 정보 손실이 가장 작은 값을 고른다. 어떤 방식을 택하든 계층 관계를 기억해야 한다. \textit{방법}은 보정 세트가 만든 세계 안에서 자를 조절하지만, 그 세계 자체를 정의하는 것은 \textit{세트}다. 앞의 낮/밤 실패가 어떤 방법 선택보다 중요한 이유다.

두 번째 함정은 현대 LLM 스택의 형태를 좌우했으므로 따로 살펴볼 가치가 있다. 가중치는 얌전한 종 모양 분포를 보이지만, \textbf{대형 트랜스포머의 활성값은 그렇지 않다.} 일부 채널에서 전형적인 크기의 \textasciitilde{}100×에 이르는 드물고 거대한 스파이크가 나타난다. 2022년 LLM.int8() 연구로 널리 알려진 현상이다. 이상치 하나가 자를 늘려 정상 값을 모두 뭉개진 값으로 반올림하게 만든다. 활성값을 단순하게 양자화하면 13B 모델의 퍼플렉시티가 폭발한다. 실용적인 대응책은 앞으로 만날 순서대로 다음과 같다. \textbf{가중치 전용 양자화(weight-only quantization)}는 가중치를 INT8/INT4로 저장하되 FP16으로 연산한다. §10.2에서 바이트를 지배하는 것이 가중치임을 확인했으므로 대부분의 이점이 남으며, LLM의 기본 방식이다. \textbf{평활화}는 SmoothQuant 방식의 재스케일링으로 활성값의 범위 부담을 가중치 쪽으로 옮긴 뒤 둘 다 양자화한다. \textbf{혼합 정밀도(mixed precision)}는 이상치 채널만 FP16으로 유지한다. 이번 주 CNN급 실습 모델의 활성값은 온순하므로 전체 W8A8 INT8이 잘 동작한다. 이상치 이야기는 11주 차에 다시 꺼낼 수 있게 기억해 두자.

\begin{tbbfigure}
  \adjustbox{max width=0.94\linewidth,max totalheight=0.46\textheight}{\includegraphics{figures/fig10-4_quantmenu.png}}\\[3pt]
  \figcaption{그림 10-4  선택지와 게이트. 정밀도의 단계(제1장의 70B 모델은 INT4에서 140 GB → \textasciitilde{}35 GB), PTQ 대 QAT, 그리고 모든 손실 산출물이 버전을 얻기 전에 통과해야 할 게이트를 보여 준다.}
\end{tbbfigure}

\subsection*{게이트: 손실 산출물이 배포 자격을 얻는 법}

사례 10-A의 E행을 엔지니어링이라 부를 수 있는 것은 처리량을 보기 전에 \textbf{게이트}를 거쳤기 때문이다. 양자화 모델은 \textit{변경된} 모델이다. 본질적으로 무손실인 변환을 위해 만든 §7.6의 동등성 검사를 손실 변환에 맞게 강화해야 한다. 게이트는 세 부분으로 구성된다. 첫째, \textbf{감이 아니라 작업 지표}를 사용한다. 별도 보관 세트에서 정확도, F1 또는 퍼플렉시티를 기준 모델과 양자화 모델 사이에 비교한다. 들뜬 사람이 임의로 실행하는 것이 아니라 CI가 실행한다. 둘째, \textbf{허용 오차를 문서로 명시}한다. "≤ 0.5 points top-1 loss"는 벤치마크 \textit{전에} 내리는 엔지니어링 결정이다. 정확도에 대응하는 SLO이며, SLO와 마찬가지로 논쟁을 계산으로 바꾼다. 셋째, \textbf{슬라이스를 점검}한다. 집계값은 특정 구간에 집중된 손상을 감춘다. 전체적으로 0.2포인트 하락했지만 제품의 존재 이유인 클래스 하나에서는 4포인트가 무너질 수 있다. 따라서 게이트는 이해관계자가 중요하게 여기는 슬라이스를 검사한다. 터미널 기록 10-3은 통과한 실행을 보여 준다.

\begin{terminalbox}
\begin{Verbatim}[breaklines=true,breakanywhere=true,fontsize=\small,formatcom=\color{tbbtermfg},codes=\tbbcodeactive]
$ python gate.py --baseline clf_fp16 --candidate clf_int8 --tol 0.5
held-out set: 10,000 images (never seen by calibration)
  baseline   top-1  76.14%
  candidate  top-1  75.92%     delta -0.22 pt   (tol 0.5)  PASS
  worst slice: "streetcar" -1.8 pt  (slice fail-line 3.0)  flag
  logit drift |p99|: 0.041
verdict: PROMOTE -> s3://models/clf/v8-int8/  (new digest)
# fail path exists and is honorable: ship FP16, file the finding.
# an ungated quant is how "the model got weird" tickets are born.
\end{Verbatim}
\end{terminalbox}
\figcaption{터미널 기록 10-3  게이트 통과. 허용 오차는 실행 전에 명시했고, 슬라이스 검사는 집계값이 손상을 감추는 곳을 살폈다. 출력은 자체 다이제스트를 지닌 새 산출물이다. 기존 산출물을 제자리에서 변경하지 않는다.}

두 가지 운영상의 귀결로 이 절을 마무리하자. \textbf{양자화 모델은 새로운 산출물}이다. 새 다이제스트를 갖고, 자체 게이트 보고서를 옆에 저장하며, 모델 스토어(제6장)에 별도 버전으로 둔다. 저장소가 여러 버전을 나란히 두도록 설계된 이유(제9장의 기관 ④)가 바로 이런 롤백에 있다. 또한 \textbf{양자화는 하드웨어의 가격표를 다시 쓴다.} A100의 VRAM이 필요하던 모델이 이제 T4 플릿에 들어갈 수 있다. 제1장의 대표 사례였던 70B 모델은 INT4에서 140 GB에서 \textasciitilde{}35 GB로 줄어 저렴한 GPU 네 장이나 큰 GPU 한 장이면 된다. 정밀도를 바꿀 때마다 제2장의 플릿 표를 다시 읽어야 한다. 제5장의 \code{nodeSelector} 질문인 "어느 GPU 레이블을 겨냥할 것인가?"에는 새로운 정답이 생기며, 구속 차원이 달라졌을 수도 있다.

\begin{calloutbox}{tbbpurple}{tbbpurplefill}
{\normalsize\bfseries\color{tbbpurple}생각해 보기}\par\vspace{3pt}
\begin{tbbitemize}
  \item 보정 세트는 지난달 트래픽에서 표본으로 뽑은 이미지 500장이다. 다음 주에 새로운 국가에서 제품을 출시한다. PTQ가 겪도록 설정된 실패 유형과 이에 해당하는 §10.3의 조절점(눈금 배치)을 설명하고, 저렴하게 완화할 방법을 제시하라.
  \item 자신의 프로젝트 모델에 적용할 게이트 문장을 작성하라. 지표, 별도 보관 세트, 허용 오차, 반드시 점검할 슬라이스 하나를 포함해야 한다. 이어서 회의적인 팀원에게 그 허용 오차 수치의 출처를 설명하고 정당화하라.
  \item 가중치 전용 INT4는 활성값을 FP16으로 유지하므로 연산 자체는 변하지 않는다. 그런데도 배치 1 LLM의 지연 시간은 거의 4× 개선된다. 이를 §10.2와 모순 없이 두 문장으로 설명하라. 어느 하한이 움직였는가?
  \item 어떤 공급업체가 "정확도 손실 0을 보장하는 INT8이며 보정도 필요 없다"고 주장한다. 이 절을 바탕으로 그 주장의 실제 의미를 드러낼 세 가지 질문을 제시하라.
  \item BF16은 FP32의 지수를 유지하면서 가수부 16비트를 버리고, FP16은 가수부를 더 많이 유지하면서 범위를 버린다. "자 자체의 차이" 소절을 바탕으로 BF16이 학습에서 각광받은 이유와 서빙에서는 둘 중 어느 것도 상대보다 우위가 아닌 이유를 설명하라. FP8-E4M3이 INT8에는 없는 무엇을 제공하는지도 덧붙여라.
\end{tbbitemize}
\end{calloutbox}

\section{그래프 컴파일: 더 적은 왕복과 실행}

양자화가 바이트를 줄였다면 컴파일은 \textit{낭비}를 공격한다. 그 낭비를 보려면 즉시 실행이 정직하고 순진하게 일하는 모습을 관찰하면 된다. 즉시 실행 PyTorch는 Python이 모델을 만나는 순서대로 실행한다. 각 연산이 자체 GPU 커널을 실행하고, VRAM에서 입력을 읽은 뒤 출력을 다시 기록한다. conv → batch-norm → ReLU → add 연쇄에서는 커널을 네 번 실행한다. 각 실행은 고정 비용으로 제9장의 \textbf{a} 일부를 차지한다. 메모리 병목 환경에서는 더 나쁜 일도 생긴다. \textbf{중간 활성값이 VRAM을 세 번 왕복}한다. §10.2에서 비용으로 보아야 한다고 배운 바로 그 바이트를 몇 마이크로초만 존재하는 값에 쓴다. 즉시 실행의 유연성은 어떤 Python과 제어 흐름도 허용하므로 프로토타입에 적합하다. 이 절은 더는 쓰지 않는 유연성에 비용을 내지 않을 때 어떤 일이 생기는지를 다룬다.

\begin{tbbfigure}
  \adjustbox{max width=0.94\linewidth,max totalheight=0.46\textheight}{\includegraphics{figures/fig10-5_compiler.png}}\\[3pt]
  \figcaption{그림 10-5  컴파일러의 역할. 즉시 실행에서는 연산마다 커널 하나가 실행되고 활성값이 메모리를 왕복한다. 컴파일하면 커널이 융합되고 값이 레지스터에 머문다. a가 줄고, 메모리 벽을 세 번이 아니라 한 번만 건넌다.}
\end{tbbfigure}

\subsection*{도구 상자: "그래프 최적화"의 실제 의미}

컴파일러에 전체 그래프, 즉 §7.6의 내보낸 산출물을 건네면 네 계열의 재작성이 가능해진다. 내보내기 형식이 존재하는 이유다. \textbf{융합}은 원소별 연산과 정규화 연산의 연쇄를 값을 생성하는 커널에 합친다. conv+bn+ReLU+add가 \textit{한 번}의 실행이 되고 중간값은 레지스터를 떠나지 않는다. 네 번 실행하고 세 번 왕복하는 대신 벽을 한 번만 건넌다. \textbf{상수 폴딩(constant folding)}은 빌드 시점에 계산할 수 있는 것을 미리 계산한다. 학습이 끝나면 상수일 뿐인 batch-norm의 스케일과 이동값을 conv의 가중치에 합치고, reshape를 접으며, 실행되지 않는 분기를 없앤다. \textbf{레이아웃과 정밀도 계획}은 텐서 코어가 요구하는 특정 메모리 레이아웃(NHWC 대 NCHW)과 위치별 정밀도를 맞춘다. 연산마다 전치하지 않고 컴파일러가 전체 텐서를 한 번에 다시 배치한다. \textbf{커널 오토튜닝(autotuning)}은 각 융합 연산에 존재하는 수십 가지 후보 구현(타일링, 루프 펼치기, 텐서 코어 경로)을 평가한다. 빌드 시점에 \textit{실제 GPU와 실제 형상에서 벤치마크}하고 가장 좋은 구현을 결과물에 넣는다. 마지막 항목은 단일 기법으로 가장 큰 개선을 내는 동시에 §7.6에서 이미 본 절충의 원인이 된다. \textbf{출력은 하나의 GPU 모델과 하나의 형상 범위에 결합된다.} 가장 빠르지만 가장 좁다.

\subsection*{직접 계산하기: 8.4 → 3.1 ms는 어디에서 사라졌는가}

"컴파일러가 a를 줄였다"는 말을 냅킨 한 장에서 계산할 수 있다면 구호로 남겨 둘 이유가 없다. ResNet급 분류기는 즉시 실행 순방향 패스마다 약 \textbf{175개의 CUDA 연산}을 수행한다. conv 약 53개와 각각의 batch-norm, \textasciitilde{}49개 ReLU, 16개 잔차 덧셈, 풀링과 헤드가 포함된다. 즉시 실행에서는 각 연산이 Python 디스패치와 커널 실행 비용을 낸다. 이를 \textbf{연산당 \textasciitilde{}25 μs의 오버헤드}라고 하면 한 패스에 175 × 25 μs ≈ \textbf{4.4 ms의 순수 관리 비용}이 붙는다. GPU가 실제 작업으로는 전혀 보지 못하는 시간이다. TensorRT는 이 175개 연산을 \textbf{\textasciitilde{}60개 커널}로 융합하고 연산마다 Python을 거치지 않은 채 C++에서 실행한다. 오버헤드는 \textasciitilde{}0.3 ms로 줄어든다. 메모리 벽의 몫도 더해 보자. 배치 1의 중간 활성값은 총 \textasciitilde{}45 MB이고, 즉시 실행의 개별 bn/ReLU/add 커널이 이를 다시 읽고 쓰면서 320 GB/s에서 \textbf{피할 수 있는 트래픽 90 MB ≈ 0.3 ms}를 만든다. 관리 비용(≈4.1 ms 절감) + 왕복(≈0.3 ms) + 상수 폴딩으로 완전히 삭제한 수십 개 연산을 합치면 \textbf{a에서 사라진 5.3 ms}가 된다. 이 계산은 터미널 기록 10-6에서 확인할 수 있는 미묘한 현상도 예측한다. 배치 32에서는 활성값이 32× 커져 피할 수 있는 트래픽이 \textasciitilde{}3 GB, 즉 패스당 ≈ 9 ms이고 \textit{항목당} \textasciitilde{}0.3 ms가 된다. 따라서 융합은 b도 조금씩 줄인다. C에서 D로 갈 때 적합 결과의 b가 1.74 → 1.62 ms로 소폭 내려가는 이유다.

상수 폴딩의 묘기도 온전히 살펴보자. 세 줄의 대수만으로 네트워크에서 연산자 계열 하나를 통째로 삭제한다. 추론 시 batch-norm은 conv 출력 z = Wx + b에 y = γ·(z − μ)/σ + β를 적용한다. 학습이 끝나면 네 매개변수는 모두 \textbf{고정된 상수}다. 대입하고 항을 다시 묶으면 \textbf{y = (γ/σ)·W·x + (γ(b − μ)/σ + β)}가 된다. 이는 가중치 W′ = (γ/σ)W, 편향 b′ = γ(b − μ)/σ + β인 또 다른 합성곱일 뿐이며 \textit{빌드 시점}에 계산할 수 있다. 융합을 시작하기도 전에 53개의 batch-norm이 0개가 된다. 이는 §7.6의 동등성 검사에도 조용한 경고를 보낸다. 폴딩한 네트워크는 대수적으로 같지만 부동소수점 연산의 \textit{결합 순서가 달라진다}. 컴파일러가 손대는 순간 비트 단위 동일성은 사라진다. 그래서 검사는 비트가 아니라 허용 오차를 확인한다.

\subsection*{세 가지 컴파일러, 하나의 결정 패턴}

\textbf{ONNX Runtime}은 이식성 있는 중간 지점이다. 그래프 패스(융합, 폴딩)를 적용한 뒤 CUDA, TensorRT, CPU 등의 플러그형 \textbf{실행 공급자}에 디스패치한다. 하나의 산출물을 어디서나 최고는 아니지만 준수한 속도로 서빙할 수 있다. 사례 10-A의 C행에서 별다른 대가 없이 1.4×를 얻었으며, 이기종 플릿에서 합리적인 기본값이다. \textbf{TensorRT}는 NVIDIA에 모든 것을 건다. \textit{빌더}가 그래프와 \textbf{최적화 프로파일}, 즉 전송을 약속한 형상 범위인 min/opt/max 배치와 차원을 받아 정확한 GPU에 맞춰 계층별 커널을 오토튜닝한다. 선택적으로 INT8도 보정하며, §10.3의 스케일을 같은 빌더가 계산한다. 그 결과 빠르게 로드되고 가장 빠르게 실행되는 직렬화된 계획인 \textbf{엔진}이 나온다(D–E행). 엔진의 좁은 범위는 이론이 아니라 운영 현실이다. 프로파일 밖의 형상을 거부하고 다른 GPU 세대에서는 실행되지 않는다. 실습에서 산출물 이름을 \code{t4/\allowbreak clf\_int8.plan}으로 지정하는 이유다. 제4장의 다이제스트 규율에 하드웨어 좌표가 추가되었다. \textbf{torch.compile} (PyTorch 2.x)은 즉시 실행에 자연스러운 경로다. 데코레이터 하나가 일반 Python에서 그래프를 캡처해 JIT 컴파일한다(Inductor). 워크플로를 거의 바꾸지 않고 그래프 중심 모델에서 큰 개선을 얻는다. 다만 비용이 다른 위치에 발생한다. 컴파일이 \textit{런타임}에 일어나므로 첫 요청이 워밍업 비용을 부담하고 준비 상태 프로브와 콜드 스타트 예산이 이를 흡수해야 한다. 제5장과 제12장에서 이 문제를 다룬다. 동적 제어 흐름은 그래프를 여러 작은 컴파일로 쪼갤 수 있다.

\begin{terminalbox}
\begin{Verbatim}[breaklines=true,breakanywhere=true,fontsize=\small,formatcom=\color{tbbtermfg},codes=\tbbcodeactive]
# the lab's two builds, in full (CI steps - never at pod start):
$ trtexec --onnx=clf.onnx --saveEngine=t4/clf_fp16.plan --fp16 \
          --minShapes=input:1x3x224x224 \
          --optShapes=input:16x3x224x224 \
          --maxShapes=input:32x3x224x224    # the shape PROMISE
$ trtexec --onnx=clf.onnx --saveEngine=t4/clf_int8.plan --int8 \
          --calib=calib.cache               # §10.3, inside the builder
# engines are per-GPU, per-profile: the t4/ prefix is load-bearing.
# send a 64-batch to a max=32 engine -> refused. the narrowness is real.
\end{Verbatim}
\end{terminalbox}
\figcaption{터미널 기록 10-4  TensorRT 엔진 빌드. 최적화 프로파일은 형상에 대한 약속이고, INT8 보정은 빌더 안에서 실행된다. 엔진은 GPU 좌표 없이는 살아남지 못하므로 산출물 경로에도 이를 포함한다.}

\begin{tbbtablewrap}
  \begin{tabular}{@{}>{\raggedright\arraybackslash}m{\dimexpr 0.1862\tbltextwidth-2\tabcolsep\relax} >{\raggedright\arraybackslash}m{\dimexpr 0.2128\tbltextwidth-2\tabcolsep\relax} >{\raggedright\arraybackslash}m{\dimexpr 0.2021\tbltextwidth-2\tabcolsep\relax} >{\raggedright\arraybackslash}m{\dimexpr 0.1862\tbltextwidth-2\tabcolsep\relax} >{\raggedright\arraybackslash}m{\dimexpr 0.2128\tbltextwidth-2\tabcolsep\relax}@{}}
  \arrayrulecolor{tbbnavy}\specialrule{1pt}{0pt}{0pt}
  \rowcolor{tbbtablehead}{\centering \bfseries\color{tbbnavy}도구} & {\centering \bfseries\color{tbbnavy}컴파일 시점} & {\centering \bfseries\color{tbbnavy}산출물과 범위} & {\centering \bfseries\color{tbbnavy}사례 10-A의 개선} & {\centering \bfseries\color{tbbnavy}치러야 할 대가} \\
  \arrayrulecolor{tbbnavy}\specialrule{0.7pt}{0pt}{2pt}
  {\textbf{ONNX Runtime}} & {로드 시점(그래프 패스) + EP 디스패치} & {EP가 있으면 어디서나 실행되는 하나의 .onnx} & {1.4× (C행)} & {준수하지만 최고는 아닌 속도, 연산 지원 범위의 경계} \\
  \arrayrulecolor{tbbtableborder}\hline
  {\textbf{TensorRT}} & {CI의 빌드 시점(GPU-프로파일마다 수분)} & {하나의 GPU + 형상 프로파일(shape profile)에 결합된 엔진} & {1.6–2.0× (D–E행)} & {좁은 범위: GPU별 산출물, 형상 거부} \\
  \arrayrulecolor{tbbtableborder}\hline
  {\textbf{torch.compile}} & {첫 호출의 런타임(JIT)} & {산출물 없음 — 데코레이터 하나, 캐시는 최선 노력} & {(벤치마크하지 않음)} & {콜드 스타트 경로의 워밍업, 새 형상마다 재컴파일} \\
  \arrayrulecolor{tbbnavy}\specialrule{1pt}{2pt}{0pt}
  \end{tabular}\par
  \figcaption{표 10-2  세 가지 컴파일러 한눈에 보기. 어느 행을 고를지는 속도의 문제이기 전에 워크플로에 관한 결정이다.}
\end{tbbtablewrap}

\subsection*{torch.compile 시간 측정 — JIT 비용은 실제로 어디에 붙는가}

표 10-2의 셋째 행은 별도의 기록으로 살펴볼 가치가 있다. 공간이 아니라 \textit{시간}의 절충이기 때문이다. 비용을 \textit{어디서}가 아니라 \textit{언제} 치르는지가 달라진다. 게다가 한 줄이면 적용할 수 있어 팀원이 가장 먼저 손댈 도구다. 분류기를 관찰해 보자.

\begin{terminalbox}
\begin{Verbatim}[breaklines=true,breakanywhere=true,fontsize=\small,formatcom=\color{tbbtermfg},codes=\tbbcodeactive]
>>> model_c = torch.compile(model)        # one line. too easy?
>>> t(model_c, batch=1)   # first call: Dynamo traces, Inductor
call 1:   38,400 ms       # autotunes kernels... on the SERVING pod
>>> t(model_c, batch=1)
steady:        6.1 ms     # vs eager 10.5 -> a real 1.7x. BUT:
>>> t(model_c, batch=7)   # a shape it has not seen ->
call 1:    9,200 ms       # recompile. every novel shape, again.
# the win is real and the bill is real: ~40 s on the cold-start path
# (Ch. 5's probes + Ch. 12's scale-out both notice), and dynamic
# shapes -> recompiles (mark_dynamic / bucketing, §10.5, tames it).
\end{Verbatim}
\end{terminalbox}
\figcaption{터미널 기록 10-5  torch.compile은 런타임에 비용을 낸다. 정상 상태에서는 컴파일러다운 개선을 주지만 첫 호출과 처음 보는 모든 형상도 컴파일러만큼 비용을 청구한다. TensorRT처럼 CI가 아니라 파드의 준비 상태 구간에 정확히 38 s가 붙는다.}

운영 관점에서는 이렇게 읽을 수 있다. \textbf{TensorRT는 컴파일을 CI로 옮겼고, torch.compile은 파드 위로 옮겼다.} 어느 쪽도 틀리지 않다. 후자는 반복 개발 속도에서 압도적이며 Python 때문에 내보낼 수 없는 모델에도 유용하다. 다만 38초는 반드시 \textit{어딘가}에 자리 잡아야 한다. 합당한 장소는 세 곳이다. \textbf{준비 상태 프로브 워밍업}은 파드가 Ready를 보고하기 전에 대표 형상을 실행하며, 제5장의 프로브가 설계 목적을 정확히 수행한다. \textbf{영속 컴파일 캐시}는 이미지나 볼륨에 담아 배포한다. 캐시 키에 GPU, 버전, 형상이 포함되므로 최선 노력 방식이다. \textbf{사전 내보내기}는 \code{torch.export} + AOTInductor를 사용한다. 즉시 실행 모델에 CI에서 빌드한 산출물을 제공하려는 생태계의 현재 시도로 지켜볼 만하다. 존중받기 어려운 방식은 12주 차 오토스케일링 실습에서야 비용을 발견하는 것이다. 스케일 아웃한 파드가 처음 40초를 컴파일에 쓰는 동안 흡수하려던 대기열은 계속 늘어난다. 제7장에서 본 재시도 폭풍 파형을 스스로 만드는 셈이다.

\subsection*{컴파일이 기대에 못 미칠 때 — 이를 예측하는 점검표}

컴파일이 어디서나 공짜 속도를 주는 것은 아니며, 실패 유형은 이 장의 두 아이디어로 예측할 수 있다. 모델이 배치 1 LLM 디코드(decode)처럼 \textbf{가중치 때문에 완전히 메모리 병목}이라면 활성값 트래픽을 융합해도 전체는 거의 움직이지 않는다. 가중치가 지배하며 이 하한은 컴파일이 아니라 양자화의 영역이다. 다시 한번 §10.2의 진단이 먼저다. 형상이 제각각인 텍스트 길이나 가변 이미지처럼 \textbf{매우 동적}이면 엔진이 여러 프로파일로 조각나고 torch.compile은 재컴파일한다. 컴파일러보다 \textit{먼저} 적용할 답은 §10.5의 버킷팅이다. Python이 \textbf{데이터에 의존}하면 그래프가 끊긴다. \code{if x.mean().item() > 0:} 한 줄이 GPU→CPU 동기화를 강제하고 그 지점에서 모델을 두 그래프로 나눈다. 토큰을 도는 Python 루프는 반복마다 컴파일 하나가 된다. \code{torch.compile}은 이를 보고하며, \code{fullgraph=True}를 쓰면 읽을 가치가 있는 오류로 바꾼다. 그래프가 \textbf{사용자 정의 연산으로 가득}하면 컴파일러가 연산별 실행으로 폴백해 아무것도 융합하지 못한다. §7.6의 내보내기 지원 범위 문제가 이제 성능 비용으로 나타난다. 모든 컴파일 산출물은 파이프라인에 \textbf{빌드 시간과 워밍업}도 추가한다. 엔진은 GPU-프로파일 쌍마다 빌드에 몇 분이 걸리는 CI 비용이 있고, torch.compile은 첫 요청에서 콜드 스타트 비용을 거둔다. 네 경우의 공통 원칙은 벽의 어느 쪽에 있는지 알고 형상을 파악한 뒤, 홍보 자료가 아니라 \textit{게이트 + 스윕}으로 결정하는 것이다. 수치적으로 컴파일러는 합법적으로 부동소수점 연산 순서를 바꾼다. §7.6의 동등성 검사가 이미 이를 다루므로 양자화 산출물과 마찬가지로 컴파일 산출물마다 실행해야 한다. 융합하고 순서를 바꾼 결과 역시 변경된 결과다.

\begin{calloutbox}{tbbpurple}{tbbpurplefill}
{\normalsize\bfseries\color{tbbpurple}생각해 보기}\par\vspace{3pt}
\begin{tbbitemize}
  \item 사례 10-A의 C→D행은 대부분 a에서 개선되었고(그림 10-6의 적합에서 8.4→3.1 ms), b는 거의 변하지 않았다. 융합과 오토튜닝은 a에, INT8은 b에 주로 영향을 주는 이유를 메커니즘별로 설명하라. 각 기법에서 가장 큰 이점을 얻는 모델도 예측하라.
  \item 팀원이 채팅 모델에 torch.compile을 적용해 정상 상태에서 1.7× 개선되었다며 기뻐한다. 그런데 12주 차 오토스케일러가 플래핑하기 시작하고 스케일 아웃할 때마다 p99가 치솟는다. 두 사실을 연결하고, 제5장과 제12장의 도구 상자에서 완화책을 하나씩 제시하라. 확인을 위해 측정할 항목도 말하라.
  \item FP16과 INT8 변형을 갖춘 두 GPU 플릿(T4 + A100)의 CI 행렬을 설계하라. 엔진 빌드와 게이트 실행은 각각 몇 번 필요하며, 모델 스토어의 레이아웃은 어떻게 구성되는가(제6장의 접두사)? 새 GPU 세대가 도입될 때 어느 부분에서 비용이 발생하는가?
\end{tbbitemize}
\end{calloutbox}

\section{수단 조합하기 — a와 b 다시 적합하기}

제9장은 이번 주에 "바로 그 스윕을 그대로 다시 실행하고, 바로 그 두 수치를 다시 적합하겠다"고 약속했다. 이제 그 약속을 지킬 차례다. 9주 차 분류기에서 a = 8.4 ms, b = 2.05 ms로 적합한 법칙 \textbf{t(B) = a + b·B}를 떠올리자. 각 변형을 두 번, 배치 1과 배치 32에서 실행하면 최적화할 때마다 다시 적합할 수 있다. 이 적합은 진단이다. 변형이 더 빠르다는 사실뿐 아니라 \textbf{속도가 어디에서 왔는지} 알려 준다. 바로 이 점이 벤치마크 보고서와 막대그래프를 구분한다.

\begin{terminalbox}
\begin{Verbatim}[breaklines=true,breakanywhere=true,fontsize=\small,formatcom=\color{tbbtermfg},codes=\tbbcodeactive]
# refit a,b from two timed points per variant (t1, t32):
$ python fit.py --t1 10.5 --t32 74.0     # B: eager FP16 (Wk 9)
  a = 8.4 ms    b = 2.05 ms    ceiling 1/b = 488/s
$ python fit.py --t1  7.6 --t32 61.6     # C: ONNX Runtime
  a = 5.9 ms    b = 1.74 ms    ceiling 575/s
$ python fit.py --t1  4.7 --t32 55.0     # D: TensorRT FP16
  a = 3.1 ms    b = 1.62 ms    ceiling 617/s
$ python fit.py --t1  3.8 --t32 32.3     # E: TensorRT INT8
  a = 2.9 ms    b = 0.92 ms    ceiling 1,087/s
# compilation attacked a (8.4->3.1): launches, round-trips, folding.
# quantization attacked b (2.05->0.92): the bytes per item.
# exactly as Chapter 9 predicted. write BOTH stories in the report.
\end{Verbatim}
\end{terminalbox}
\figcaption{터미널 기록 10-6  다시 적합한 결과. 각 최적화는 자신의 흔적을 남긴다. 컴파일러는 고정 비용 a를 줄이고 양자화는 한계 비용 b를 줄이며, 상한 1/b는 두 배 이상 높아진다. 변형마다 두 번만 실행하면 적합할 수 있다.}

\begin{tbbfigure}
  \adjustbox{max width=0.94\linewidth,max totalheight=0.46\textheight}{\includegraphics{figures/fig10-6_levers.png}}\\[3pt]
  \figcaption{그림 10-6  여러 수단은 합성된다. 전체 a/b 표, SLO에서 측정한 수치(238 → 507 req/s), 조합 규칙을 보여 준다. 순서가 중요하며 가변 형상은 모든 단계에 비용을 더한다.}
\end{tbbfigure}

\subsection*{배치 처리가 효과를 내는 이유 — 대수로 한 번에 정리한 원리}

강의계획서에서는 이번 주의 셋째 주제를 "동적 배치의 원리"라고 부른다. 두 장에 걸쳐 이를 사용했으니 이제 한 줄로 정리할 수 있다. t(B) = a + b·B를 B로 나누면 \textbf{항목당 비용 c(B) = b + a/B}가 된다. b를 향해 내려가는 쌍곡선이다. 원리는 이것이 전부다. \textbf{배치는 오직 a를 상각하는 기법}이다. 실행, 스케줄링, 가중치 스트림이라는 고정 비용을 B명의 동승자에게 나누지만 b는 건드릴 수 없다. 수익이 줄어드는 까닭도 속설이 아니라 \textit{산술}에 있다. B에서 2B로 늘리면 항목마다 a/2B를 회수하므로 배치를 두 배로 키울 때마다 얻는 이점은 직전의 절반이다. 표 10-3은 최적화된 분류기(a = 2.9, b = 0.92)의 수치를 계산한다.

\begin{tbbtablewrap}
  \begin{tabular}{@{}>{\raggedright\arraybackslash}m{\dimexpr 0.1383\tbltextwidth-2\tabcolsep\relax} >{\raggedright\arraybackslash}m{\dimexpr 0.2128\tbltextwidth-2\tabcolsep\relax} >{\raggedright\arraybackslash}m{\dimexpr 0.2553\tbltextwidth-2\tabcolsep\relax} >{\raggedright\arraybackslash}m{\dimexpr 0.2021\tbltextwidth-2\tabcolsep\relax} >{\raggedright\arraybackslash}m{\dimexpr 0.1915\tbltextwidth-2\tabcolsep\relax}@{}}
  \arrayrulecolor{tbbnavy}\specialrule{1pt}{0pt}{0pt}
  \rowcolor{tbbtablehead}{\centering \bfseries\color{tbbnavy}배치 B} & {\centering \bfseries\color{tbbnavy}t(B) = a + b·B} & {\centering \bfseries\color{tbbnavy}항목당 c(B) = b + a/B} & {\centering \bfseries\color{tbbnavy}이상적 req/s = B/t(B)} & {\centering \bfseries\color{tbbnavy}상한(1/b) 대비} \\
  \arrayrulecolor{tbbnavy}\specialrule{0.7pt}{0pt}{2pt}
  {1} & {3.8 ms} & {3.82 ms} & {262/s} & {24\%} \\
  \arrayrulecolor{tbbtableborder}\hline
  {4} & {6.6 ms} & {1.65 ms} & {608/s} & {56\%} \\
  \arrayrulecolor{tbbtableborder}\hline
  {16} & {17.6 ms} & {1.10 ms} & {908/s} & {84\%} \\
  \arrayrulecolor{tbbtableborder}\hline
  {32} & {32.3 ms} & {1.01 ms} & {990/s} & {91\%} \\
  \arrayrulecolor{tbbtableborder}\hline
  {\textbf{∞}} & {—} & {\textbf{→ b = 0.92 ms}} & {\textbf{→ 1,087/s}} & {100\%} \\
  \arrayrulecolor{tbbnavy}\specialrule{1pt}{2pt}{0pt}
  \end{tabular}\par
  \figcaption{표 10-3  상각 계산(INT8 변형: a = 2.9 ms, b = 0.92 ms). 셋째 열에서 굴곡점이 보인다. "이상적 req/s"는 대기열 처리를 무시하므로 SLO에서 측정한 실제 수치는 이보다 낮다(제7장의 여유 용량).}
\end{tbbtablewrap}

표를 두 가지로 읽으면 앞의 대수가 제값을 한다. 첫째, \textbf{어려운 부분은 "동적"이라는 말이며, 이는 GPU가 아니라 대기열에 관한 용어다.} 트래픽은 한 번에 하나의 요청으로 도착한다. 따라서 §9.3의 배처는 \textit{기다림으로} 상각을 산다. 함께 탈 요청을 기대하며 첫 요청을 max\_queue\_delay만큼 붙잡되, 기다림 + t(B)가 SLO 예산 안에 들어야 한다. 제9장에서 관찰한 자기 조절에는 이제 작동 원리가 생겼다. 부하가 낮으면 거의 비어 있는 상태에서 창이 만료되어 작은 배치가 된다. 유휴 시스템에는 효율이 필요 없으므로 괜찮다. 부하가 높으면 즉시 차서 큰 배치가 되고, 효율이 필요한 바로 그때 나타난다. 둘째, \textbf{최적화는 이제 직관적으로 예상할 수 있는 방향으로 굴곡점을 옮겼다.} a가 8.4에서 2.9 ms로 줄면 \textit{상각할 고정 비용 자체가 적다.} 배치 4에서 이미 상한의 56\%를 얻으며, 최적화 전에는 \textasciitilde{}49\%였다. 배치를 더 깊게 해도 얻는 몫은 얇아진다. 따라서 지연 창을 \textit{줄여} 지연 시간 개선을 취하는 것이 흔히 옳다. 아래 규칙 ②는 주장한 것이 아니라 도출한 결론이다. "최적화 뒤에 조절점을 다시 튜닝하라"는 말이 선택적인 위생 관리가 아닌 이유도 여기에 있다. 같은 두 상수가 움직였기 때문이다.

\subsection*{조합 규칙}

수단을 쌓는 방식은 세 가지 규칙이 지배한다. 이를 모르면 동료 중 누군가는 반드시 같은 실수를 한다. \textbf{규칙 ① — 먼저 컴파일하고 컴파일러 안에서 양자화한다.} TensorRT의 INT8 모드는 정밀도를 \textit{중심으로} 커널을 계획한다. 어느 계층과 융합 그룹을 쓸지, 역양자화(dequantization) 경계를 어디에 둘지 결정한다. 컴파일러가 보지 못한 그래프를 양자화하면 두 도구 모두 눈을 가린 채 일한다. 그래서 사례 10-A의 E행은 두 부품을 따로 붙인 것이 아니라 한 번에 빌드한 "TensorRT INT8"이다. \textbf{규칙 ② — 배치 처리는 남은 것을 곱한다.} 배처는 줄어든 a를 B개 항목에 상각하므로 최적화 후 §9.3의 조절점을 다시 튜닝해야 한다. a = 8.4가 아니라 2.9 ms이면 더 작은 배치가 더 일찍 효율에 도달하고 굴곡점은 \textasciitilde{}2× 더 멀어진다. 처리량이 오르는 동안에도 max\_queue\_delay를 \textit{줄일} 수 있는 경우가 많다. 가정하지 말고 다시 스윕하라. \textbf{규칙 ③ — 가변 형상은 모든 것에 비용을 부과한다.} 배치는 모든 항목을 가장 넓은 형상에 맞춰 패딩한다. 6토큰 쿼리를 500토큰 문서와 같은 배치로 묶으면 500토큰 분량의 작업 비용을 낸다. 해결책은 \textbf{버킷팅(bucketing)}이다. 이미지 해상도나 시퀀스 길이 구간 같은 크기 등급으로 요청을 묶어 패딩(padding) 낭비를 제한한다. 버킷은 컴파일러의 형상 프로파일(§10.4)과 \textit{함께} 설계해야 한다. 버킷과 프로파일은 서로 다른 두 도구에서 본 같은 결정이다. 11주 차에는 이를 논리적 끝까지 밀고 간다. LLM 생성에서는 매 단계마다 모든 시퀀스의 길이가 달라지므로 연속 배치가 필요하다.

보고서에는 정직한 각주가 하나 필요하다. 이 개선은 \textbf{레플리카별 경제성이지 마법이 아니다.} 507 req/s 변형도 여전히 하키 스틱을 따른다. 굴곡점이 옮겨졌을 뿐이다. SLO를 지키려면 여유 용량이 필요하고, 보정을 소홀히 하면 게이트에 실패한다. 여러 수단을 합성한 개선이 플릿의 가격표를 다시 쓴다. 프로젝트 모델이 SLO에서 28 → 57 req/s로 개선되면 제7장이 예고한 그대로 레플리카가 8개 → 4개, 비용이 \$4,762 → \$2,381로 바뀐다. 12주 차 오토스케일링은 이 비용을 다시 절반으로 줄일 차례를 기다리고 있다.

\begin{calloutbox}{tbbpurple}{tbbxF5F3FF}
{\normalsize\bfseries\color{tbbx5B21B6}상자 10-2  나머지 선택지 — 이번 주에 의도적으로 사용하지 않은 수단}\par\vspace{3pt}
이 장의 세 수단에는 공통점이 있다. \textbf{학습 루프가 필요 없다.} 이미 가진 산출물을 대상으로 CI에서 몇 분 안에 실행한다. 더 넓은 선택지는 이 장점을 포기하는 대신 더 크거나 다른 개선을 얻는다. 어떤 방법이 있고 대가가 무엇인지 알아 두자.

\begin{tbbitemize}
  \item \textbf{증류(distillation)}(Hinton et al., 2015): 작은 \textit{학생} 모델이 큰 \textit{교사} 모델의 출력을 모방하도록 학습한다. DistilBERT는 BERT보다 매개변수를 40\% 줄이면서도 성능의 \textasciitilde{}97\%를 유지한 것으로 유명하다. 구조적이고 영구적인 개선이지만, 전체 학습 파이프라인과 데이터 및 수 주의 시간이 필요하다. CI 단계가 아니라 \textit{모델링} 프로젝트다.
  \item \textbf{가지치기(pruning)와 구조적 희소성(structural sparsity)}: 가중치를 완전히 삭제한다. 비구조적 가지치기는 GPU에서 효과를 내는 경우가 드물다. 불규칙한 0은 작업을 건너뛰게 하지 못하기 때문이다. \textbf{2:4 구조적 희소성}은 가중치 네 개마다 두 개를 0으로 만들며 Ampere+ 텐서 코어에서 실제 연산 속도를 두 배로 높인다. 다만 희소화와 미세 조정 주기를 거쳐야 하므로 역시 학습이다.
  \item \textbf{추측 디코딩(speculative decoding)}(11주 차): 작은 모델이 토큰 초안을 만들고 큰 모델이 한 번의 배치 패스로 검증한다. LLM의 토큰별 가중치 스트림을 직접 공격한다. 재학습은 필요 없지만 생성에 특화된 메커니즘이므로 생성 자체가 온전한 주제가 되는 장에서 다룬다.
  \item \textbf{캐스케이드와 조기 종료}: 쉬운 입력은 작은 모델로, 어려운 입력은 큰 모델로 보낸다. \textit{시스템} 차원의 수단이며 제12장의 라우터 영역에 속한다. 게이트도 정확도 델타가 아니라 라우팅 지표다.
\end{tbbitemize}

선택지 어디에서나 규율은 같다. 먼저 진단하고(§10.2), 정직하게 게이트를 적용하며(§10.3), 개선의 원인을 밝힌다(§10.5). 통화만 다를 뿐이다.
\end{calloutbox}

\vspace{3.0pt}

\begin{calloutbox}{tbbpurple}{tbbpurplefill}
{\normalsize\bfseries\color{tbbpurple}생각해 보기}\par\vspace{3pt}
\begin{tbbitemize}
  \item 팀원이 "배치 32에서 INT8이 1.26×밖에 개선하지 못했다. 홍보 자료에는 2×라고 했는데 실망스럽다"고 보고한다. 터미널 기록 10-6의 수치로 배치 32의 속도 향상이 SLO에서의 개선을 과소평가하는 이유를 보여라. t(32)의 비율과 상한의 비율을 계산하면 된다. 대신 보고해야 할 단일 수치는 무엇인가?
  \item 최적화 후 기존 max\_queue\_delay=10 ms 창이 SLO 굴곡점을 지나 포화되는 배치를 모은다. a = 2.9, b = 0.92, SLO 500 ms에서 §9.3의 두 조절점을 다시 튜닝하는 과정을 설명하라. 무엇이 변하고 무엇은 변하지 않는가? 제5장의 memory.max도 다시 측정해야 하는 이유는 무엇인가?
  \item 텍스트 서비스의 길이가 5–800 tokens 범위에서 균일하게 분포한다. 버킷의 수와 경계를 설계하고 이에 맞는 TensorRT 프로파일을 구성하라. 버킷이 있을 때와 없을 때 최악의 패딩 낭비도 계산하라.
\end{tbbitemize}
\end{calloutbox}

\section{실습: 전후 벤치마크 보고서}

이번 실습은 이 장의 내용을 자신의 모델에서 직접 실행한다. 결과물은 프로젝트 평가 기준이 1주 차부터 요구해 온 \textbf{전후 벤치마크 보고서}, 즉 "최적화 전후 비교"다. 그림 10-7의 구조를 보자. 9주 차 스윕은 이미 \textit{이전} 결과다. 이번 주에는 게이트를 적용한 변형을 빌드하고 같은 스윕을 그대로 다시 실행한 뒤, 표를 보고서로 바꾸는 두 문단을 작성한다.

\begin{tbbfigure}
  \adjustbox{max width=0.94\linewidth,max totalheight=0.46\textheight}{\includegraphics{figures/fig10-7_lab.png}}\\[3pt]
  \figcaption{그림 10-7  실습 파이프라인과 보고서 구조. CI에서 빌드한 각 변형은 게이트와 동일한 스윕을 모두 거친다. 결과물은 표 하나와 정직한 두 문단이다.}
\end{tbbfigure}

\subsection*{일곱 단계}

\begin{tbbitemize}
  \item \textbf{① 진단한다(§10.2).} 자신의 GPU에서 자신의 모델을 대상으로 터미널 기록 10-1의 두 줄을 계산한다. 메모리 벽의 어느 쪽인가? 어떤 수단이 가장 큰 효과를 낼지 어떤 빌드도 하기 \textit{전에} 예측해 적는다. 예측이 틀렸다는 사실을 명확히 드러내도 점수를 얻을 수 있다. 보고서를 통해 무언가 배웠다는 뜻이기 때문이다.
  \item \textbf{② 변형 C를 빌드한다(§10.4).} 내보내기 → CUDA EP를 사용하는 ONNX Runtime 순서로 진행한다. §7.6의 동등성 검사를 실행한다. 이 변형은 거의 정확히 통과해야 한다.
  \item \textbf{③ 변형 D와 E를 빌드한다(§10.4 + §10.3).} 명시적인 형상 프로파일로 TensorRT FP16을 빌드한 뒤, 엄선한 보정 세트로 INT8을 빌드한다. 최근 트래픽에서 500–1,000개의 입력을 뽑되 클래스 \textit{및 게이트가 검사하는 슬라이스}에 걸쳐 층화한다. 세트의 매니페스트와 다이제스트를 calib 캐시 옆에 저장한다. 게이트의 신뢰성은 보정 세트의 문서화 수준과 정확히 같다.
  \item \textbf{④ 각 변형에 게이트를 적용한다(§10.3).} 작업 지표를 미리 명시한 허용 오차와 비교하고 슬라이스를 검사한다. 각 게이트 보고서를 산출물 옆에 저장한다. 게이트에 실패하면 FP16을 배포하고 발견 사항을 기록한다. 이 역시 \textit{완전하며 통과한} 실습 결과다.
  \item \textbf{⑤ 살아남은 각 변형을 스윕한다(§9.5, 변경 없음).} 같은 locustfile, 같은 속도, 같은 워밍업 폐기 구간을 사용한다. TensorRT 엔진 워밍업은 정상 상태 평균에 섞지 않고 콜드 스타트 데이터로 보고한다.
  \item \textbf{⑥ a와 b를 다시 적합한다(§10.5).} 변형마다 시간을 측정한 지점 두 개를 추가하고 모든 개선을 해당 메커니즘에 귀속한다. 귀속 결과가 예상과 다르면 그대로 적는다. 그 문장이 보고서에서 가장 가치 있을 수 있다.
  \item \textbf{⑦ 가격을 다시 계산한다(제7장의 청구서).} SLO에서 측정한 레플리카별 새 수치 → 새 플릿 크기 → 새 월 비용을 제안서에 반영한다. 프로젝트의 비용 이야기는 이제 추정 → 측정(9주 차) → 최적화(10주 차)로 이어지고, 12주 차 오토스케일링을 다음 수단으로 명시해야 한다.
\end{tbbitemize}

"정직한 두 문단"은 실제로 어떤 모습일까? 이미 본 수치를 모두 활용해 분류기용 예시를 작성했다. 문장이 아니라 \textit{구조}를 참고하라.

\begin{calloutbox}{tbbteal}{tbbxF0FDFA}
{\normalsize\bfseries\color{tbbx115E59}사례 10-B  완성된 보고서의 두 문단}\par\vspace{3pt}
\textbf{¶1 — 개선은 어디에서 왔는가.} "최적화로 분류기의 성능은 p99 ≤ 500 ms SLO에서 238 req/s에서 507 req/s로 이동했다(2.1×). T4의 요청 백만 건당 비용으로는 \$0.62 → \$0.29다. 다시 적합한 결과는 개선의 출처를 보여 준다. 컴파일은 디스패치와 실행 오버헤드를 줄여 a를 8.4에서 3.1 ms로 낮췄다. 즉시 실행 연산 \textasciitilde{}175개를 커널 \textasciitilde{}60개로 융합한 §10.4의 계산과 일치한다. 활성값 왕복을 융합해 b도 2.05에서 1.62 ms로 줄였다. 이어 INT8은 텐서 코어 속도가 두 배가 된 덕분에 b를 0.92 ms로 줄였다. ①단계의 예측은 절반만 맞았다. INT8이 바이트를 절반으로 줄이는 효과가 지배할 것으로 예상했지만, 적합 결과 이 모델은 배치 1에서 \textit{실행 병목}이고 배치 16에서는 연산 병목에 가까웠다. INT8의 개선은 바이트가 아니라 연산 속도에서 나왔다. 메커니즘에 대한 예상은 틀렸지만 크기는 맞았다. 다음 모델은 60× 더 크므로 이 차이가 중요하다."

\textbf{¶2 — 믿기 전에 무엇을 확인했는가.} "모든 변형이 §7.6의 동등성 검사를 통과했다. INT8은 벤치마크 \textit{전에} 제안서에 기록한 0.5 pt 허용 오차에 대해 −0.22 pt로 게이트를 통과했다. 보고서는 t4/clf\_int8/gate.json에 저장했다. 슬라이스 검사에서는 \textit{streetcar}가 −1.8 pt로 표시되었다. 보정 매니페스트를 보니 streetcar 이미지가 3장뿐이어서 40장으로 다시 보정했고, 전체 델타는 그대로인 채 해당 슬라이스가 −0.6 pt로 회복되었다. 모든 변형은 동일한 locustfile(다이제스트 9f3c…), 동일한 속도를 사용했고 첫 60 s를 폐기했다. TensorRT 엔진 로드는 정상 상태 평균에 넣지 않고 콜드 스타트로 보고했다. 엔진에는 t4/ 접두사가 있으며 팀의 A10G 한 대는 설계대로 이를 거부한다."

\textit{이 글이 좋은 평가를 받는 이유는 ¶1의 모든 수치를 메커니즘에 귀속하고 ¶2의 모든 주장을 저장된 산출물로 확인할 수 있기 때문이다. 틀린 예측 하나도 숨기지 않고 명시했다. 채점자와 미래의 팀원, 14주 차 사고 검토자가 바로 이런 내용을 읽는다.}
\end{calloutbox}

\vspace{3.0pt}

\begin{calloutbox}{tbbred}{tbbxFEF2F2}
{\normalsize\bfseries\color{tbbx991B1B}상자 10-1  보고서의 위험 신호(공개된 채점자 점검표)}\par\vspace{3pt}
\begin{tbbitemize}
  \item \textbf{게이트 결과가 나란히 없는 속도 향상} — 빠르지만 틀린 모델은 최적화가 아니다.
  \item \textbf{변형마다 서로 다른 스윕}(속도, 페이로드, 워밍업) — 비교 자체가 결과물이므로 지켜야 한다.
  \item \textbf{배치 N 지연 시간 비율을 "속도 향상"으로 제시} — 중요한 수치는 SLO에서의 req/s(유효 처리량)와 a/b 귀속이다.
  \item 경로나 메타데이터에 \textbf{GPU 좌표가 없는 엔진 산출물} — 언젠가는 반드시 잘못된 플릿에 배포된다. 제4장의 다이제스트 교훈을 하드웨어에 적용한 사례다.
  \item \textbf{"시작할 때 파드 안에서 양자화했다"} — 변환은 CI에서 한 번 수행하고 게이트를 적용한다(§7.6, §10.4). 콜드 스타트 경로는 계속 신성하게 지킨다.
\end{tbbitemize}
\end{calloutbox}

\vspace{3.0pt}

제출물은 표 하나와 두 문단(개선이 어디에서 왔는가, 믿기 전에 무엇을 확인했는가)으로 구성한 보고서, 변형별 게이트 출력, 다시 적합한 상수, 갱신한 제안서 비용 항목이다. 결제 내역 스크린샷도 계속 제출한다. 엔진 빌드를 포함한 일반적인 실습 비용은 T4 시간 기준 \textasciitilde{}\$3–4다.

\begin{calloutbox}{tbbpurple}{tbbpurplefill}
{\normalsize\bfseries\color{tbbpurple}생각해 보기}\par\vspace{3pt}
\begin{tbbitemize}
  \item 보고서의 첫 문단을 "§10.2의 내용에 따라 컴파일이 지배할 것으로 예상한다..."와 같은 가설 형태로 미리 작성한 뒤 ⑥단계가 끝나면 채점하라. 자신의 모델은 실제로 메모리 벽의 어느 쪽에 있었는가?
  \item INT8 변형이 0.5 pt 허용 오차에 대해 0.9 pt로 게이트에 실패한다. 네 가지 선택지를 비용이 늘어나는 순서대로 나열하라. 더 나은 보정 → 채널별/그룹별 세분성 → 가중치 전용/혼합 정밀도 → QAT 순이다. 각 단계에서 다음 단계로 넘어갈 근거도 제시하라.
  \item 최적화 모델은 제2장의 플릿 답을 바꾼다. T4 한 대당 507 req/s일 때 가격이 ≈5×이고 이 모델 처리량이 ≈2.5×인 A100이 이 분류기 서빙에서 유리한 부하가 하나라도 있는가? 계산으로 보여라. A100이 유리해지는 워크로드도 제시하라. 힌트는 §10.2의 메모리 벽이다.
\end{tbbitemize}
\end{calloutbox}

\namedsection{요약}{열다섯 가지 핵심 내용}

\qitem{1}{\textbf{최적화란 같은 SLO에서 요청당 비용을 낮추는 것이다.} 사례 10-A에서는 하나의 모델과 하나의 GPU를 네 스택에서 실행했다. 같은 p99에서 28 → 57 req/s, 백만 건당 \$7.94 → \$3.90이 되었고, 제7장이 예고한 대로 플릿은 레플리카 8개 → 4개, \$4,762 → \$2,381/month로 줄었다. 모델의 지식은 전혀 바뀌지 않았다.}

\qitem{2}{\textbf{우리는 평생 손실 압축(lossy compression)을 믿어 왔다.} 8-bit μ-law 전화(1962), JPEG (1992), MP3 (1993)는 신호 소비자에게 필요한 것을 찾아 나머지에 비용을 내지 않는 방식이다. 양자화는 모델 자체의 결정을 소비자로 삼아 이 패턴을 적용한다. 학습에는 작은 경사 단계를 위한 세밀한 눈금이 필요하지만 추론에서는 반올림 뒤에도 최댓값의 위치나 토큰만 살아남으면 된다. 평평한 최솟값이 이를 가능하게 한다.}

\qitem{3}{\textbf{메모리 벽은 가장 중요한 진단 기준이다.} T4는 65 TFLOPS를 수행하지만 320 GB/s만 이동하므로 바이트당 \textasciitilde{}200 FLOP이 필요하다. 이 능선과 비교한 산술 집약도(FLOPs/byte)가 모든 것을 결정한다. 서빙 배치에서 대부분의 추론은 메모리 병목이므로 비용은 연산이 아니라 바이트다. 터미널 기록 10-1처럼 프로파일러 없이 두 줄로 진단할 수 있다.}

\qitem{4}{\textbf{루프라인에서 각 수단은 기하학적 의미를 지닌다.} 배치는 가중치를 B개 항목에 재사용해 오른쪽으로 이동시키고, 양자화는 가중치당 바이트를 줄여 대역폭 지붕을 높이며, 컴파일은 융합 커널로 활성값이 VRAM에 가지 않게 해 벽의 비용을 중복해서 내지 않는다. 능선은 카드마다 다르다(표 10-1). L4는 약 400이고 A100 프리미엄의 대부분은 대역폭이다. 세대마다 벽이 높아지므로 바이트는 가치가 오르는 통화다.}

\qitem{5}{\textbf{바이트는 세 가지 속도의 여정을 거친다(제6장의 약속을 지켰다).} 객체 스토어 → 호스트 페이지 캐시(파드마다 한 번) → PCIe의 고정 H2D(\textasciitilde{}12 GB/s, 프로세스마다 한 번) → VRAM에서 연산 장치(300 GB/s, 모든 패스)로 이동한다. PCIe는 VRAM보다 \textasciitilde{}25× 느리고 페이지 가능 복사는 이를 다시 절반으로 줄인다. 따라서 가중치는 시작 시 한 번만 건너고 요청별 텐서는 고정 버퍼와 비동기 복사를 사용한다. nvidia-smi dmon의 PCIe 열은 진단의 셋째 줄이다.}

\qitem{6}{\textbf{부동소수점은 관대하고 정수는 위임한다.} 부동소수점은 지수로 값마다 자체 스케일을 조정하지만 정수는 텐서 전체가 보정된 스케일 하나를 공유한다. 이것이 §10.3의 모든 함정의 근원이다. FP16은 범위를 희생해 학습 시 손실 스케일링이 필요하고, BF16은 FP32의 범위를 유지해 학습에서 각광받았다. FP8 (Hopper/Ada)은 정수 수준의 저장 비용에 자체 스케일링을 남겼다. INT8×INT8의 합은 INT32에 누적되므로 정확도 저하는 언제나 눈금 배치 문제이지 오버플로 문제가 아니다.}

\qitem{7}{\textbf{양자화는 거친 자를 잘 배치하는 일이다.} q = round(w/scale) + zero\_point에서 비대칭 zero\_point는 자를 옮겨 ReLU 이후 범위에서 눈금을 낭비하지 않게 한다. 가중치에는 채널별 스케일을 쓴다. 하나의 공유 스케일은 ±0.05 채널에 21\% 오차를 내지만 자체 자는 0.4\%만 낸다. INT4 LLM은 그룹별 스케일을 쓴다. 보정은 최솟값–최댓값, 백분위수 또는 KL로 활성값 범위를 정하며, 보정의 품질이 곧 양자화 품질이다. PTQ는 CI에서 몇 분이면 되고 PTQ가 게이트에 실패하면 QAT로 확대한다.}

\qitem{8}{\textbf{두 가지 함정이 있다.} 채널마다 범위가 다르므로 채널별 스케일이 필요하고, 트랜스포머 활성값에는 이상치가 숨어 있다(LLM.int8(), 2022). 그래서 LLM은 기본적으로 가중치 전용 양자화를 쓰며, 평활화나 혼합 정밀도를 적용하기도 한다. CNN급 실습 모델은 전체 W8A8을 견딘다.}

\qitem{9}{\textbf{게이트는 손실 변환에 정당성을 부여한다.} 작업 지표를 벤치마크 전에 명시한 허용 오차와 비교하고 슬라이스를 검사한다. 출력은 게이트 보고서를 나란히 둔 새 다이제스트 고정 산출물이다. 게이트에 실패해 FP16을 배포하는 것은 통과한 결과지만, 게이트 없이 배포하면 "모델이 이상해졌다"는 티켓이 생긴다.}

\qitem{10}{\textbf{컴파일러는 a를 줄인다. 융합, 폴딩, 레이아웃 계획, 오토튜닝이 핵심이다.} 네 번 대신 한 번만 실행하고 활성값 왕복을 없애며, 빌드 시점에 실제 GPU와 형상에서 커널을 벤치마크한다. 대가는 좁은 범위다. 엔진은 GPU별, 형상 프로파일별 산출물이며 t4/ 접두사가 꼭 필요하다. torch.compile의 워밍업은 콜드 스타트 경로에 붙는다.}

\qitem{11}{\textbf{세 컴파일러는 하나의 패턴을 따른다.} ONNX Runtime은 플러그형 EP를 갖춘 이식성 있는 중간 지점이고, TensorRT는 가장 빠르지만 가장 좁으며 빌더 안에서 INT8을 보정한다. torch.compile은 즉시 실행에 자연스러운 런타임 JIT다. 가중치 병목 디코드, 제각각인 형상, 사용자 정의 연산에서는 컴파일이 예상대로 기대에 못 미치므로 먼저 진단해야 한다.}

\qitem{12}{\textbf{배치는 오직 a를 상각한다.} c(B) = t(B)/B = b + a/B는 b를 향해 내려가는 쌍곡선이다. B를 두 배로 늘릴 때마다 얻는 이점은 직전의 절반이므로 굴곡점은 속설이 아니라 산술에서 나온다. 표 10-3에서 배치 4는 최적화 상한의 56\%를 이미 얻는다. \textit{동적}이라는 말은 대기열 측면을 뜻한다. SLO 예산 안에서 동승자를 기다리는 제한된 시간을 효율과 맞바꾸고, 부하에 따라 스스로 조절한다.}

\qitem{13}{\textbf{조합 규칙을 지켜야 한다.} 먼저 컴파일하고 컴파일러 안에서 양자화한다. 배치 처리는 남은 것을 곱하므로 §9.3의 조절점을 다시 튜닝한다. 굴곡점이 움직였고 상각할 a가 줄어 지연 창을 흔히 \textit{줄일} 수 있다. 가변 형상은 모든 것에 비용을 부과하므로 형상 프로파일과 함께 크기별 버킷을 설계한다.}

\qitem{14}{\textbf{다시 적합한 결과가 보고서다.} t(B) = a + b·B에서 변형마다 두 번 실행한다. 컴파일의 흔적은 a (8.4 → 3.1 ms), 양자화의 흔적은 b (2.05 → 0.92 ms)에 나타난다. 상한은 488 → 1,087/s, 측정값은 SLO에서 238 → 507 req/s가 된다. 원인 귀속이 막대그래프를 엔지니어링으로 바꾼다.}

\qitem{15}{\textbf{최적화는 하드웨어의 가격표를 다시 쓴다.} 대표 70B 모델은 INT4에서 140 → \textasciitilde{}35 GB로 줄고, A100이 필요하던 모델이 T4 플릿에 들어간다. 제2장의 표와 제5장의 nodeSelector에는 새로운 정답이 생긴다. 정밀도를 바꿀 때마다 플릿 표를 다시 읽어야 한다.}

\namedsection{연습문제}{}

\subsection*{A. 객관식}

\qitem{1}{FP16의 7B 모델(14 GB의 가중치)이 대역폭 700 GB/s인 GPU에서 배치 1로 토큰을 생성한다. 토큰당 지연 시간 하한과 그 원인을 가장 잘 추정한 것은 무엇인가?}

\begin{tbboptions}
  \item ① \textasciitilde{}2 ms — 커널 실행 오버헤드        ② \textasciitilde{}20 ms — 토큰마다 가중치를 한 번 스트리밍해야 한다(14 GB ÷ 700 GB/s)
  \item ③ \textasciitilde{}0.1 ms — 연산이 단순하다        ④ 프로파일러 없이는 추정할 수 없다
\end{tbboptions}

\qitem{2}{가중치를 FP16에서 INT8로 양자화할 때 직접 생기지 않는 변화는 무엇인가?}

\begin{tbboptions}
  \item ① 패스마다 스트리밍하는 바이트가 절반이 된다        ② 가중치용 VRAM이 절반이 된다
  \item ③ INT8 지원 GPU에서 텐서 코어 연산 속도가 두 배가 된다        ④ 모델 정확도가 높아진다
\end{tbboptions}

\qitem{3}{트랜스포머 활성값 이상치가 단순한 W8A8 양자화를 깨뜨리는 이유는 무엇인가?}

\begin{tbboptions}
  \item ① INT8은 음수를 표현할 수 없다
  \item ② 소수의 큰 값이 스케일을 늘려 일반 값을 거의 0에 가깝게 반올림하게 한다
  \item ③ 활성값은 요청마다 변하므로 어떤 스케일도 존재할 수 없다
  \item ④ GPU는 활성값을 FP64로 처리한다
\end{tbboptions}

\qitem{4}{가장 큰 이점을 얻기 위한 올바른 작업 순서는 무엇인가?}

\begin{tbboptions}
  \item ① ONNX 파일을 양자화한 뒤 TensorRT에 건넨다        ② 컴파일러가 정밀도에 맞춰 커널을 계획하도록 INT8 모드 + 보정으로 TensorRT 엔진을 빌드한다
  \item ③ 개선은 독립적이므로 순서는 중요하지 않다        ④ 먼저 배치하고 컴파일한 뒤 파드 시작 시 양자화한다
\end{tbboptions}

\qitem{5}{다시 적합한 결과 변형 X에서 b는 변하지 않고 a만 8.4 → 3.2 ms로 줄었다. 변형 X일 가능성이 가장 큰 것은 무엇인가?}

\begin{tbboptions}
  \item ① INT8 양자화        ② 그래프 컴파일(융합 + 오토튜닝)
  \item ③ 더 큰 배치 크기        ④ 더 빠른 토크나이저
\end{tbboptions}

\qitem{6}{T4에서 maxShapes batch=32로 빌드한 TensorRT 엔진이 A10G에서 배치 48을 받는다. 어떤 일이 일어나는가?}

\begin{tbboptions}
  \item ① 조금 느리지만 실행된다        ② 배치를 조용히 잘라 낸다
  \item ③ 엔진용 GPU가 잘못되었고 형상도 프로파일 밖이므로 두 조건 모두에서 실패한다. 엔진은 GPU별, 프로파일별 산출물이다
  \item ④ TensorRT가 즉석에서 엔진을 다시 빌드한다
\end{tbboptions}

\qitem{7}{게이트의 허용 오차("≤ 0.5 pt top-1 loss")를 벤치마크 전에 명시해야 하는 이유는 무엇인가?}

\begin{tbboptions}
  \item ① 규정이 요구한다        ② 그렇지 않으면 관찰된 손실을 보고 나서야 "허용 가능"의 뜻을 정하게 된다. 측정 전에 SLO를 약속하는 것에 대응하는 정확도 규율이다(§7.4)
  \item ③ TensorRT가 빌드 플래그로 사용한다        ④ 보고서가 짧아진다
\end{tbboptions}

\qitem{8}{최적화 후(a: 8.4→2.9, b: 2.05→0.92) §9.3의 배치 조절점에 관한 설명으로 옳은 것은 무엇인가?}

\begin{tbboptions}
  \item ① 영향을 받지 않는다. 배치 처리는 직교한다
  \item ② 굴곡점이 바깥으로 \textasciitilde{}2× 이동하고 더 작은 배치가 더 빨리 효율에 도달하므로 두 조절점과 제5장의 memory.max를 모두 다시 스윕해야 한다
  \item ③ max\_queue\_delay를 두 배로 늘려야 한다        ④ max\_batch\_size에는 더 이상 VRAM 상한이 없다
\end{tbboptions}

\subsection*{B. 단답형}

\qitem{1}{320 GB/s GPU에서 배치 1로 실행하는 1.5 GB INT8 모델의 두 줄짜리 진단(바이트 하한과 FLOP 하한)을 제시하고 결론을 내려라.}

\qitem{2}{게이트의 세 부분과 이를 통과한 모든 변형에 적용되는 산출물 규칙을 서술하라.}

\qitem{3}{LLM에서는 가중치 전용 INT4가 기본인 반면 실습의 CNN에서는 전체 W8A8이 동작하는 이유는 무엇인가? 모델 계열마다 한 문장으로 답하라.}

\qitem{4}{TensorRT의 "최적화 프로파일"이 무엇이며 어떤 운영 약속을 표현하는지 설명하라.}

\qitem{5}{PTQ와 QAT를 각각 한 문장으로 설명하고, 첫 방법에서 둘째 방법으로 확대하는 조건을 제시하라.}

\qitem{6}{FP16에서 ReLU의 산술 집약도(arithmetic intensity)를 계산하고, 이 수치로 그래프 컴파일러가 ReLU에 별도 커널을 주지 않는 이유를 설명하라.}

\subsection*{C. 서술형}

\qitem{1}{가상의 서비스에 대해 "네 가지 비용" 분석을 작성하라. A100 (weights 26 GB FP16, 2,000 GB/s, 312 TFLOPS)에서 배치 1로 LLM을 디코드한다. 메모리와 연산 하한을 계산하고, 이 장의 수단 중 무엇이 가장 큰 효과와 가장 작은 효과를 낼지 이유와 함께 예측하라. 11주 차 연속 배치가 이들 수단으로는 얻을 수 없는 무엇을 추가하는지도 설명하라.}

\qitem{2}{제도로서 게이트를 옹호하라. PM은 "정확도 0.2 points를 대가로 \$2,400/month를 아낄 수 있는데 당연히 배포해야지, 왜 이런 절차가 필요한가?"라고 묻는다. 허용 오차가 수치보다 먼저 정해지는 이유, 집계값이 보호할 수 없는 무엇을 슬라이스 검사가 보호하는지, 14주 차 사고에서 다이제스트 고정 산출물 + 저장된 게이트 보고서가 주는 이점, PM이 실제로 옳은 지점을 포함해 답하라.}

\qitem{3}{자신의 프로젝트 모델을 위한 전체 최적화 파이프라인을 CI로 설계하라. 단계(내보내기, 빌드, 보정, 게이트, 스윕, 다시 적합, 승격), 산출물과 저장소 레이아웃(제6장), 커밋마다 실행할 작업과 릴리스마다 실행할 작업, 두 실패 경로(게이트 실패, 스윕 회귀)와 각각의 결과를 포함하라. 둘째 GPU 유형이 플릿에 들어오면 어느 단계가 달라지는지도 적어라.}

\qitem{4}{경제성에 관한 글을 작성하라. 최적화 후 분류기는 T4 한 대(≈\$0.53/hr, 표 2-3)에서 507 req/s를 처리한다. 요청 백만 건당 새 비용을 계산한 뒤 같은 모델에서 처리량이 2.5×라고 가정한 A100의 손익분기 가격을 구하라. 메모리 벽을 이용해 메모리 병목 LLM에서는 A100의 이점이 더 커지는 이유를 설명하고, 플릿 권장안과 제5장 nodeSelector에 미치는 영향으로 결론지어라.}

\namedsection{정답과 해설}{}

\subsection*{A. 객관식}

\qitem{1}{\textbf{②} 생성되는 토큰마다 14 GB의 가중치 전체가 연산 장치를 통과해야 한다. 14 ÷ 700 GB/s = 20 ms다. 더 빠른 연산으로도 넘을 수 없는 하한이며, 배치 1 디코드가 루프라인 맨 왼쪽에 놓이는 이유다. (§10.2)}

\qitem{2}{\textbf{④} 양자화는 손실 변환이다. 게이트는 개선을 기대하기 위해서가 아니라 손실의 범위를 제한하기 위해 존재한다. 드물게 우연한 정규화 효과가 있었다는 일화는 엔지니어링 계획이 아니다. ①–③은 정확히 양자화 메커니즘이다. 바이트와 VRAM이 줄고 정수 연산이 빨라진다. (§10.3)}

\qitem{3}{\textbf{②} 드문 스파이크 하나가 자를 늘린다. 256개 눈금이 ±100에 퍼지면 일반적인 ±1 값에는 눈금 몇 개만 남아 뭉개진 값으로 반올림된다. 그래서 트랜스포머에는 가중치 전용 양자화, 평활화 또는 혼합 정밀도를 쓴다. (§10.3)}

\qitem{4}{\textbf{②} 규칙 ①에 따라 컴파일러가 정밀도를 알아야 융합 INT8 커널과 역양자화 경계를 계획할 수 있다. 컴파일러가 보지 못하게 미리 양자화하거나 시작 시점에 양자화하면(④, 콜드 스타트에도 해롭다) 두 도구를 모두 낭비한다. (§10.5)}

\qitem{5}{\textbf{②} a는 패스마다 드는 고정 비용인 실행, 왕복, 폴딩 기회이며 컴파일러가 주로 다루는 영역이다. 양자화는 주로 항목당 바이트와 연산을 뜻하는 b에 나타난다. 다시 적합하는 이유가 바로 이처럼 원인을 귀속하기 위해서다. (§10.5)}

\qitem{6}{\textbf{③} 엔진은 한 GPU 세대에 오토튜닝한 커널과 약속된 형상 범위를 인코딩한다. 둘 중 하나라도 어기면 실행을 거부하거나 아예 로드되지 않는다. 산출물 경로에 하드웨어 좌표를 넣는 이유다. (§10.4)}

\qitem{7}{\textbf{②} 결과를 본 뒤 정의한 허용 오차는 허용 오차가 아니라 합리화다. 약속이 측정보다 앞서야 한다는 점에서 SLO와 같은 규율이다. (§10.3, §7.4)}

\qitem{8}{\textbf{②} a가 작아지면 배치가 상각할 오버헤드가 줄어 효율에 더 일찍 도달하고 전체 곡선이 압축된다. 굴곡점은 req/s 기준으로 더 바깥에 놓인다. 두 조절점과 memory.max의 근거가 되는 최악의 메모리 측정을 모두 다시 살펴봐야 한다. (§10.5, 제5장)}

\subsection*{B. 단답형}

\qitem{1}{바이트: 1.5 GB ÷ 320 GB/s ≈ 패스당 \textbf{4.7 ms}의 하한이다. FLOP: \textasciitilde{}1.5B params × 2 ≈ 3 GFLOP ÷ 65 TFLOPS ≈ \textbf{0.05 ms}다. 결론은 \textasciitilde{}100× 메모리 병목이므로 연산이 아니라 바이트를 줄이는 데 투자해야 한다. 추가 양자화와 배치 처리가 해당한다. (§10.2)}

\qitem{2}{① 별도 보관 세트에서 기준 모델과 비교한 작업 지표, ② 실행 전에 명시한 허용 오차, ③ 집계값이 손상을 감추는 곳의 슬라이스 검사다. 통과한 출력은 게이트 보고서를 옆에 저장한 \textbf{새 다이제스트 고정 산출물}이어야 하며 제자리에서 변경하지 않는다. (§10.3)}

\qitem{3}{LLM: 활성값에 자를 늘리는 스파이크인 이상치가 있으므로 활성값은 FP16으로 두고, 어차피 바이트를 지배하는 가중치만 INT4로 낮춘다. CNN: 활성값이 온순한 종 모양 분포이므로 전체 W8A8이 게이트를 통과하고 정수 연산 속도도 얻는다. (§10.3)}

\qitem{4}{엔진을 빌드할 때 사용한 min/opt/max 형상 범위이며 무엇을 전송할지에 관한 약속이다. 범위 안에서는 오토튜닝한 커널을 쓰고 밖에서는 거부한다. 운영상 §10.5의 버킷팅과 연결되므로 버킷과 프로파일은 하나의 결정이다. (§10.4)}

\qitem{5}{PTQ는 학습 후 보정 데이터에서 스케일을 고르며 학습 루프 없이 몇 분이면 된다. QAT는 미세 조정 중 반올림을 모사해 네트워크가 적응하게 하므로 학습 비용이 든다. PTQ가 명시된 게이트에 실패하고 정밀도 개선이 학습 파이프라인 보유 비용을 정당화할 때 QAT로 확대한다. (§10.3)}

\qitem{6}{원소당 1 FLOP ÷ (읽기 2 B + 쓰기 2 B) = \textbf{0.25 FLOP/byte}다. T4의 능선보다 \textasciitilde{}800× 낮아 ReLU는 어떤 하드웨어에서도 메모리 속도로 실행된다. 독립 ReLU 커널은 순수한 왕복 비용이므로 컴파일러가 값을 생성하는 연산의 에필로그에 융합해 바이트를 완전히 없앤다. (§10.2, §10.4)}

\subsection*{C. 서술형(채점 기준)}

\qitem{1}{하한: 메모리는 26 ÷ 2,000 GB/s = \textbf{13 ms/token}, 연산은 \textasciitilde{}13 GFLOP ÷ 312 TFLOPS ≈ 0.04 ms이므로 \textasciitilde{}300× 메모리 병목이다. 예측: 가중치 전용 양자화가 가장 큰 효과를 내며 INT4는 하한을 \textasciitilde{}3.3 ms로 줄인다. 융합은 가중치 스트림을 건드릴 수 없어 컴파일의 효과가 가장 작다. 동시 시퀀스가 있을 때만 배치가 도움이 된다. 연속 배치는 길이와 단계가 다른 시퀀스 사이에 그 동시성을 만들어 스트리밍한 가중치를 매 단계 상각한다. §9.3의 가정을 복구하며 vLLM이 존재하는 이유다. (§§10.2–10.5, 11주 차)}

\qitem{2}{좋은 답은 다음 내용을 포함한다. 허용 오차가 수치보다 먼저 정해져야 바람에 이끌린 결정을 피할 수 있다(§7.4의 규율). 슬라이스는 제품을 지탱하는 클래스가 평균에 묻히지 않게 보호한다. 사고 중에는 다이제스트 + 저장된 게이트 보고서가 "모델이 바뀌었고 검사를 거쳤는가?"에 몇 분 안에 답하며 제13장의 롤백에 필요하다. 최종 결정이 제품 차원의 절충이라는 점에서는 PM이 옳다. 게이트는 이를 \textit{근거 있는} 절충으로 만들며, 허용 오차는 배포 전에 공개적으로 다시 협의할 수 있다. (§10.3)}

\qitem{3}{포함할 내용: 커밋마다 내보내기 + 동등성 검사, 릴리스마다 GPU 유형별 엔진 빌드 + 보정 + 게이트를 실행한다. 둘째 플릿 구성원이 들어오면 행렬이 두 배가 된다. 스윕 + 다시 적합은 일정에 따라 또는 릴리스마다 실행한다. 승격은 가중치 옆에 게이트 산출물을 둔 저장소 이동이며 t4/, a100/ 접두사를 쓴다. 게이트 실패 → 이전 정밀도 배포 + 발견 사항 기록, 스윕 회귀 → 승격 차단 후 a/b 귀속으로 변형을 이분 탐색한다. 새 GPU가 들어오면 엔진을 다시 빌드하고, 수치가 달라질 수 있으므로 게이트와 스윕을 다시 실행한다. ONNX를 이식 가능한 단일 진실 공급원으로 유지해야 하는 이유다. (§§10.3–10.6, 제4/6장)}

\qitem{4}{분류기: \$0.53 ÷ (507 × 3,600) × 10⁶ ≈ \textbf{\$0.29/M}이다. 처리량이 2.5× (1,268/s)일 때 A100의 손익분기 가격은 0.53 × 2.5 ≈ \textbf{\$1.33/hr} 아래다. 온디맨드(on-demand) 실제 가격 \$3–5의 3분의 1 수준이므로(표 2-3) 연산 비중이 충분한 이 CNN에는 T4가 유리하다. 메모리 병목 LLM에서는 FLOP이 아니라 A100의 6× 대역폭이 배수가 되고 처리량이 대역폭에 거의 비례하므로 같은 계산이 대형 카드 쪽으로 뒤집힌다. 플릿 권장안은 분류기를 T4 레이블에, LLM을 A100 레이블에 두는 것이며 nodeSelector 분할에 비용 근거가 생긴다. (§10.2, §10.6, 제2/5장)}

\namedsection{더 읽을거리}{}

\textbf{[이번 주]}로 표시한 자료는 실습을 직접 뒷받침하고, \textbf{[N주 차]}는 나중에 도움이 될 자료를 뜻한다.

\begin{tbbitemize}
  \item \textbf{[이번 주] TensorRT Developer Guide — "Working with INT8" and "Optimization Profiles"} — 터미널 기록 10-4의 권위 있는 설명이다. 보정, 계층별 정밀도 결정, 형상 약속 메커니즘을 다룬다. 그림 10-5와 함께 읽어 보자.
  \item \textbf{[이번 주] ONNX Runtime 문서 — Graph Optimizations and Execution Providers} — "공짜" 1.4×가 실제로 수행한 작업을 설명한다. 최적화 수준의 단계와 EP가 TensorRT 등을 포함한 하위 그래프를 위임하는 방식을 다룬다.
  \item \textbf{S. Williams, A. Waterman, D. Patterson, "Roofline: An Insightful Visual Performance Model," CACM 2009} — 원래의 루프라인 논문이다. §10.2는 이를 서빙에 맞게 요약했다. 짧고 시간이 지나도 유효하다. GPU가 지금 겪는 문제에 이름을 붙인 CPU 시대의 네 쪽짜리 논문 \textbf{Wulf \& McKee, "Hitting the Memory Wall" (1995)}도 함께 읽어 보자.
  \item \textbf{J. Dettmers et al., "LLM.int8(): 8-bit Matrix Multiplication for Transformers at Scale" (2022)} — 활성값 이상치의 발견을 보기 드물게 명료하게 설명한다. §10.3에서 LLM의 기본값이 가중치 전용인 이유다. 재스케일링 해결책은 \textbf{Xiao et al., "SmoothQuant" (2022)}에서 이어서 볼 수 있다.
  \item \textbf{[11주 차] Frantar et al., "GPTQ" (2022) · Lin et al., "AWQ" (2023)} — 다음 주 vLLM으로 서빙할 모델의 기반이 되는 4비트 가중치 전용 기법이다. 이 장에서 배운 그룹별 스케일 용어를 알면 읽기 쉽다.
  \item \textbf{PyTorch 문서 — torch.compile(Inductor) 튜토리얼과 문제 해결} — 그래프 단절, 재컴파일 조건, 콜드 스타트 예산이 부담할 워밍업 비용 등 즉시 실행에 자연스러운 경로의 실제 동작을 설명한다.
  \item \textbf{[이번 주, 재미] 어떤 이미지 편집기의 JPEG 품질 슬라이더} — 사진에서 100부터 10까지 내려 보며 직접 차이를 알아채는 지점을 찾아보자. 자신의 시각 피질에 보정 + 게이트를 수행하는 셈이며, 실습은 여기에 허용 오차를 명시한 같은 실험이다.
\end{tbbitemize}

\vspace{5.0pt}

\begin{calloutbox}{tbbblue}{tbbbluefill}
{\normalsize\bfseries\color{tbbblue}다음 장 — 제11장: LLM 서빙}\par\vspace{3pt}
지금까지 최적화한 모든 시스템은 한 덩어리로 답했다. 요청이 도착하고 순방향 패스가 실행되면 응답이 나갔다. 다음 주부터 모델은 \textbf{한 번에 한 토큰씩 답하며} 규칙이 달라진다. 토큰마다 가중치를 다시 스트리밍하고 메모리 벽은 가장 가혹한 모습으로 나타난다. 모든 시퀀스의 길이가 달라 제9장의 배치 가정이 깨지고, 계속 커지는 \textbf{KV 캐시} 때문에 메모리 관리가 게임의 전부가 된다. 제11장에서는 해결책인 연속 배치, PagedAttention, 이 기법을 유명하게 만든 vLLM 엔진을 다룬다. §7.2에서 스트리밍을 위해 약속한 두 시계인 TTFT와 tokens/second에도 마침내 SLO를 부여한다. 실습에서는 작은 오픈 LLM을 vLLM으로 서빙하고 스트리밍하며 정직하게 측정한다.
\end{calloutbox}

%% Chapter 11 — LLM 서빙
\setcounter{chapter}{10}
\chapter{LLM 서빙}
\chaptersubtitle{자기회귀, KV 캐시, 연속 배치 — 엔진 안으로 들어간 운영체제}

\begin{calloutbox}{tbbblue}{tbbbluefill}
{\normalsize\bfseries\color{tbbblue}학습 목표}\par\vspace{3pt}
\begin{tbbitemize}
  \item 생성 트래픽에서 요청 수준 지연 시간 SLO가 의미를 잃는 이유를 설명하고, 두 가지 인간 상수(Miller 1968, 읽기 속도)를 예산으로 삼는 \textbf{두 시계} TTFT와 TPOT로 다시 작성한다.
  \item 한 번의 생성을 \textbf{프리필과 디코드}로 분해하고 각 단계의 산술 집약도를 계산한 뒤, 제10장의 루프라인으로 각 단계가 어떤 자원을 소진하는지 예측한다.
  \item 어텐션이 다시 읽는 값에서 \textbf{KV 캐시}를 도출하고 모델 카드로 크기를 계산한다. VRAM을 가중치, 런타임, 캐시 풀로 나누어 예산을 잡고 동시성을 메모리 수치로 표현한다.
  \item 프리필–디코드 간섭과 청크 프리필을 포함해 \textbf{연속 배치}(반복 수준 스케줄링)가 제9장의 깨진 가정을 어떻게 복구하는지 설명한다.
  \item \textbf{PagedAttention}을 어텐션에 이식한 가상 메모리로 이해하고, 블록 테이블, 공유, 쓰기 시 복사를 설명한다. 낭비를 회수한 결과 측정 처리량이 2–4×가 된 이유도 이해한다.
  \item \textbf{vLLM}을 운영한다. 제9장의 여섯 기관에 대응시키고, 메트릭을 읽고, 양자화한 AWQ 모델을 서빙하며, 동시 부하에서 TTFT/TPOT를 정직하게 측정한다. 이것이 실습이자 프로젝트의 새 SLO 문장이다.
\end{tbbitemize}
\end{calloutbox}

\vspace{2.5pt}

\begin{calloutbox}{tbbteal}{tbbtealfill}
{\normalsize\bfseries\color{tbbteal}핵심 용어}\par\vspace{3pt}
자기회귀 생성 · 프리필 대 디코드 · TTFT · TPOT / 초당 토큰 수(tokens-per-second) · 스트리밍(SSE) · 읽기 속도 하한 · KV 캐시 · GQA / MQA · KV 예산 · 입장 제어와 선점 · 연속 배치 · 반복 수준 스케줄링(Orca) · 프리필–디코드 간섭 · 청크 프리필 · PagedAttention · 블록 테이블 · 접두사 캐싱 · 쓰기 시 복사 · vLLM · 추측 디코딩 · 토큰 인식 부하 테스트
\end{calloutbox}

\vspace{2.5pt}

\begin{calloutbox}{tbblightgray}{tbbboxfill}
{\normalsize\bfseries\color{tbbink}장 구성}\par\vspace{3pt}
\textbf{11.1} 모델이 말을 걸다 — 대시보드는 거짓말한다    \textbf{11.2} 생성의 구조: 두 단계, 두 시계    \textbf{11.3} KV 캐시: 메모리가 새로운 동시성이다    \textbf{11.4} 연속 배치: 시내버스    \textbf{11.5} PagedAttention과 vLLM: 운영체제가 안으로 들어가다    \textbf{11.6} 실습: 실제 LLM 서빙, 스트리밍, 측정 — 이어서 요약, 연습문제, 더 읽을거리
\end{calloutbox}

\section{모델이 말을 걸다 — 대시보드는 거짓말한다}

지금까지 만든 모든 서빙 시스템에는 말하지 않은 가정 하나가 있었다. \textbf{요청 하나가 작업 단위다.} 분류 호출은 순방향 패스 한 번의 비용을 내고, 패스는 몇 밀리초가 걸리며, 그 수치의 백분위수를 약속하고 이에 맞게 프로비저닝한다. 제7장부터 제10장까지는 이 세계를 측정하고 최적화하며 가격을 매겼다. 지난주는 승리로 끝났다. TensorRT-INT8로 최적화한 3B 채팅 모델은 L4 한 대당 57 req/s에서 p99를 500 ms 아래로 유지했고, 약속대로 플릿 비용이 절반이 되었다. 그런데 월요일에 베타 사용자가 도착했다. 이들은 스윕에 넣어 둔 정형 프롬프트를 보내지 않았다. 자기소개서를 써 달라고 요청했다.

\begin{terminalbox}
\begin{Verbatim}[breaklines=true,breakanywhere=true,fontsize=\small,formatcom=\color{tbbtermfg},codes=\tbbcodeactive]
# monday 09:12. the Week-10 stack (INT8, 4 x L4), meeting REAL
# chat traffic for the first time. locust replays the beta logs:
$ locust -f chat_replay.py --headless -u 32 -r 4 --run-time 10m
 Type  Name     # reqs   # fails |  Avg     p50      p99   | req/s
 POST  /chat     4,102   0(0.0%) | 3,846   2,910   11,400  |  6.8
# p99 = 11.4 SECONDS against a 500 ms SLO. dashboard: solid red.
$ kubectl top pod -l app=chat      # surely something is dying?
NAME         CPU(cores)   MEMORY
chat-7f9b4   3210m        11.2Gi   # ...healthy.
$ nvidia-smi --query-gpu=utilization.gpu --format=csv,noheader
94 %                                # busy AND healthy.
# nothing is broken. the ANSWERS got long:
#   "reset my password?"        ->  40 tokens,   0.9 s
#   "write my cover letter..."  -> 700 tokens,  14.8 s
# the request stopped being a unit of work. the SLO measures noise.
\end{Verbatim}
\end{terminalbox}
\figcaption{사례 11-A  p99 폭발. 모든 구성 요소가 상태 검사를 통과하지만 SLO는 20배나 위반한다. 원인은 버그가 아니라 분포다. 응답 길이가 이제 100× 차이 나고 요청 지연 시간이 그 차이를 모두 물려받는다. 죽은 것은 시스템이 아니라 메트릭이다.}

이 사례의 낯섦을 잠시 곱씹어 보자. 이 책에서 처음 만나는 \textit{종류}의 실패다. 제7장의 용어로 보면 어떤 것도 회귀하지 않았다. 장치는 금요일과 똑같이 토큰마다 같은 작업을 수행한다. 깨진 것은 \textbf{계약의 단위}다. 이제 "요청" 하나에 필요한 작업은 40 tokens에서 4,000까지 세 자릿수 범위에 걸친다. 용량 계획이 아니라 \textit{사용자가 무엇을 묻고 모델이 어떻게 답하기로 하는지}가 양을 결정한다. 이 분포 위에서 요청 수준 p99를 약속하는 것은 날씨를 약속하는 것과 같다. 제10장의 최적화를 아무리 적용해도 해결되지 않는다. 토큰당 비용을 절반으로 줄여도 자기소개서는 7초가 걸린다. 평균은 아무도 설명하지 못했고, 이제 p99도 마찬가지다. 시스템 상태가 아니라 답변 길이를 추적하기 때문이다.

\subsection*{컴퓨팅에서 가장 오래된 인터페이스의 재발견}

해결책은 55년 간격을 두고 두 번 발견되었다. 이 책의 2주 차는 CP-67이 하나의 메인프레임을 여러 개로 나누는 시분할에서 시작했다. 지금 중요한 세부는 \textit{터미널이 무엇을 했는가}다. 터미널은 \textbf{답을 타이핑해 주었다.} 초당 10 characters 속도의 텔레타이프가 컴퓨터의 답을 스트림으로 전달하자 사용자들은 이를 좋아했다. 장치가 살아 있고 자신에게 집중하며 \textit{나를 위해 일하는 것}처럼 느껴졌다. 1968년 IBM의 Robert Miller는 이 감각 아래의 심리를 측정해 이후 모든 인터페이스 혁명에서도 살아남은 세 수치를 제시했다. \textbf{0.1 s} 아래에서는 즉각적으로 느껴지고, \textbf{1 s}까지는 사고의 흐름이 끊기지 않으며, \textbf{10 s}가 지나면 주의가 무너져 다른 작업으로 전환한다. 사람은 분당 약 250 words, 즉 초당 5–7 tokens를 읽는다는 인간 상수 하나를 더하면 대화형 AI의 전체 설계 공간이 이미 1968년의 종이 위에 그려진다.

이제 그 상수를 기준으로 사례 11-A를 다시 읽어 보자. 스트리밍하지 않은 자기소개서는 14.8 s 뒤에 도착해 Miller의 포기 구간 깊숙이 들어간다. 하지만 모델은 수백 밀리초 안에 \textit{첫 단어를 생성}했고, 나머지는 사람이 읽는 속도의 일곱 배인 \textasciitilde{}50 tokens/second로 만들었다. 같은 연산을 스트림으로 전달하면 1초의 흐름 구간 \textit{안쪽에 여유롭게} 들어오고 이후 영원히 독자보다 앞서간다. \textbf{스트리밍은 답을 더 빨리 만들지 않는다. 사용자가 경험하는 시계를 바꾼다.} 이것이 ChatGPT가 1억 명에게 선보인 진짜 혁신이다. 지연 시간 최적화가 아니라 계약의 재협상이었고, 이후 모든 채팅 제품이 1968년 텔레타이프처럼 답을 타이핑하는 이유다. §7.2는 이 모드에 이름을 붙이고 "11주 차"에 시계에 SLO를 부여하겠다고 약속했다. 이제 11주 차다.

\begin{tbbfigure}
  \adjustbox{max width=0.94\linewidth,max totalheight=0.46\textheight}{\includegraphics{figures/fig11-1_twoclocks.png}}\\[3pt]
  \figcaption{그림 11-1  하나의 스트리밍 답변과 두 시계. TTFT(대기열 + 프리필)는 "반응성이 있는가?"에 답하고, TPOT는 "읽는 속도보다 빠른가?"에 답한다. 기존의 요청부터 마지막 바이트까지 p99는 아무도 묻지 않은 질문에 답한다. Miller가 1968년에 제시한 임곗값과 읽기 속도가 예산이다.}
\end{tbbfigure}

따라서 프로젝트의 SLO 문장은 폐기하는 대신 다시 쓴다. 새 문장은 생각보다 관대하다. 기존 약속인 "p99 ≤ 500 ms"는 자신도 모르게 \textit{반응성}에 관한 약속이었다. 500 ms는 Miller의 1초 선 아래에 여유 있게 들어가므로 그대로 \textbf{TTFT 예산}이 된다. 여기에 \textbf{속도 조항}을 추가한다. 스트림당 tokens/second의 하한을 읽기 속도보다 넉넉히 높게 둔다. \textbf{용량 표현}도 바꾼다. 길이가 100× 차이 나는 환경에서 무의미한 초당 요청 수가 아니라 \textit{동시 스트림}을 쓴다. 실습에서 측정할 완전한 문장은 다음과 같다. \textbf{"TTFT p99 ≤ 500 ms, 스트림당 ≥ 15 tokens/s p50, 4-레플리카 L4 플릿에서 동시 스트림 100개, 월 \$900 이내."} 제7장 규율의 모든 조항인 백분위수, 임곗값, 부하, 비용은 살아남는다. 단위의 종만 바뀌었다.

\subsection*{두 번째 사례: 무언가가 기억하고 있다}

같은 날 오후, 이 장의 또 다른 수수께끼가 나타난다. 팀원이 수업을 위해 다중 턴 데모를 실행하면서 제2장에서 들인 습관대로 \code{nvidia-smi}를 계속 지켜본다. 모델은 6 GB다. 결과를 살펴보자.

\begin{terminalbox}
\begin{Verbatim}[breaklines=true,breakanywhere=true,fontsize=\small,formatcom=\color{tbbtermfg},codes=\tbbcodeactive]
# 14:00, demo day. multi-turn chat, ~30 slow-growing conversations.
$ watch -n 60 nvidia-smi --query-gpu=memory.used --format=csv,noheader
14:00    9,812 MiB    # weights (6 GB) + runtime + a few young chats
14:07   14,203 MiB    # ...the demo continues, nothing new deployed...
14:15   19,540 MiB    # no bigger model. no leak. no new users.
14:21   22,911 MiB    # the server is REMEMBERING something.
14:22   torch.cuda.OutOfMemoryError: CUDA out of memory. Tried to
        allocate 448.00 MiB (GPU 0; 23.99 GiB total capacity)
# exit 137's GPU cousin (Chs. 3/5) - and note the amount it wanted:
# 448 MiB. remember that number. §11.3 will compute it from scratch.
\end{Verbatim}
\end{terminalbox}
\figcaption{사례 11-B  서서히 증가하는 메모리. 가중치는 그대로인데 모든 대화의 토큰마다 메모리가 늘어난다. 결국 할당기가 정확히 448 MiB를 더 요구하다 실패한다. 누수는 없다. 모델은 열려 있는 모든 대화에 대해 계속 커지는 공책을 유지하지만, 아무도 그 공책의 예산을 잡지 않았다.}

두 사례와 두 수수께끼를 한 장에서 다룬다. 사례 A에서는 \textit{시간} 축이 깨졌다. 지연 시간을 토큰 단위로 다시 계약해야 한다. 사례 B에서는 \textit{메모리} 축이 깨졌다. 토큰마다 무언가가 자라며, 연산이 아니라 그것이 GPU에 몇 명의 사용자를 넣을 수 있는지 결정한다. 둘은 한 문장으로 정리할 수 있는 기계적 사실에서 비롯된다. \textbf{언어 모델은 출력 토큰마다 전체 순방향 패스를 한 번씩 실행하고 지금까지 생성한 모든 것을 참고해 답한다.} 그 루프가 §11.2다. 참고하는 공책은 §11.3이며, 448 MiB가 4,096토큰 대화 하나의 공책 크기와 메가바이트 단위까지 일치한다는 사실을 확인한다. GPU를 이런 루프로 가득 채우는 스케줄러는 §11.4, 공책의 선반 공간 낭비를 막는 메모리 관리자는 §11.5에서 다룬다. §11.6 실습에서는 vLLM, 초시계, 10달러어치 T4 시간을 사용해 이 모든 것을 직접 다룬다.

이 장은 책이 미뤄 둔 약속을 어느 장보다 많이 갚는다. 제1장은 KV 캐시를 "11주 차의 주인공"이라 부르고 Kwon 등의 논문을 "필독"으로 지정했다. 제6장에는 고정 세션의 결론과 정확히 60.0초에 끊기는 스트림이라는 빚이 있다. 제7장은 스트리밍에 SLO를 부여해야 한다. 제9장 상자 9-1의 작은 글씨, 즉 \textit{요청 형상이 비슷해야 한다}는 가정은 여기서 폭발하며 더 읽을거리에서는 "KV 캐시 강의 뒤"에 Orca를 보라고 안내했다. 제10장은 그룹별 INT4 용어의 실제 활용, 추측 디코딩 수단, AI ≈ 1인 matvec에 대한 답을 남겼다. 아래에서 모두 갚는다.

\begin{calloutbox}{tbbpurple}{tbbpurplefill}
{\normalsize\bfseries\color{tbbpurple}생각해 보기}\par\vspace{3pt}
\begin{tbbitemize}
  \item 사례 11-A의 서비스는 "모든 대시보드에서 정상"이지만 SLO를 20× 위반한다. 제5–9장에서 모두 녹색으로 나타날 대시보드 세 개와, 아직 어느 대시보드에도 없지만 실제로 사용자 불만을 예측하는 수치 하나를 제시하라. 그 수치는 기계적으로 어디에서 나오는가?
  \item Miller의 임곗값은 1968년 시분할 터미널에서 측정되었고 읽기 속도는 그보다도 오래된 상수다. "대화형 AI 제품 설계는 반세기 전에 생긴 상수로 과도하게 결정된다"는 주장에 찬성하거나 반박하라. LLM이 실제로 새로운 상수를 추가한다면 무엇인가?
  \item PM이 max\_tokens를 100으로 제한해 사례 11-A를 해결하자고 제안한다. "자기소개서는 더 이상 없다"는 뜻이다. 서빙 삼각형과 두 시계를 이용해 무엇을 포기하게 되고 어느 사용자가 이를 느끼는지 나열하라. 같은 자원을 나누는 방법 중 스트리밍이 엄격히 더 관대한 이유도 설명하라.
  \item 사례 11-B의 할당기는 448 MiB를 요구하다 실패했다. 다음 절을 읽지 않고 모델 카드에서 이 수치를 예측하려면 제1–10장에서 배운 어떤 양이 필요한지 모두 나열하고, 무엇이 커지는지 추측하라. §11.3에서 그 추측을 채점한다.
\end{tbbitemize}
\end{calloutbox}

\section{생성의 구조: 두 단계, 두 시계}

먼저 신비를 걷어 내자. 트랜스포머 언어 모델은 토큰 시퀀스를 읽고 \textbf{다음 토큰 하나의 확률 분포}를 출력하는 함수다. 할 수 있는 일은 이것뿐이다. 따라서 200토큰 답변을 만들려면 서빙 스택이 모델을 200번 실행해야 한다. 프롬프트를 넣고 토큰 \#1을 표본 추출한 뒤 이를 붙여 다시 실행하고 토큰 \#2를 뽑는다. 붙이고, 실행하고, 표본 추출하는 루프가 이어지며, 이를 \textbf{자기회귀(autoregressive) 생성}이라고 한다. 모델이 종료 토큰을 내거나 한도에 도달해야 끝난다. 두 시계, 증가하는 메모리, 깨진 배처, 이색적인 엔진 등 LLM 서빙의 낯선 모든 현상은 이 루프 하나와 이미 배운 루프라인에서 나온다. 제10장의 도구로 요청 하나를 분해해 보자.

\begin{tbbfigure}
  \adjustbox{max width=0.94\linewidth,max totalheight=0.46\textheight}{\includegraphics{figures/fig11-2_autoregress.png}}\\[3pt]
  \figcaption{그림 11-2  요청 하나와 두 장치. 프리필은 한 번의 연산 병목 패스로 프롬프트 전체를 처리해 캐시를 채운다. 이후 디코드는 토큰마다 전체 순방향 패스를 실행하며 모든 가중치를 다시 스트리밍하므로 영원히 메모리 병목인 루프를 돈다. 엔진의 역할은 두 단계가 서로 방해하지 않게 하는 것이다.}
\end{tbbfigure}

\subsection*{프리필: 이미 알고 있는 큰 GEMM}

첫 패스는 특별하다. 프롬프트의 모든 토큰은 이미 \textit{알려져} 있어 그 사이에서 표본 추출할 필요가 없다. 따라서 모델은 500개 토큰을 \textbf{한 번의 패스에서 함께} 처리한다. 서로에게 어텐션하는 500개 항목을 9주 차 배치로 묶은 것과 같다. 제10장의 세 수치로 보면 이는 큰 GEMM이다. 산술 집약도 ≈ 토큰 수다. 스트리밍한 가중치마다 500토큰 분량의 작업을 수행하므로 500토큰 패스는 L4의 \textasciitilde{}400 능선 \textit{위}에 놓여 실제로 \textbf{연산 병목}이다. 루프라인에서 가장 좋은 영역이다. 비용은 정직한 작업에서 나온다. \textasciitilde{}2 FLOPs per parameter per token × 3B × 500 ≈ \textbf{3 TFLOP}이며, 현실적인 효율에서 L4가 50–60 ms 동안 처리한다. 프리필(prefill)은 두 가지를 만든다. 첫째는 \textbf{첫 출력 토큰}이므로 프리필 시간이 TTFT의 거의 전부다. 둘째는 앞의 두 사례에서 이름을 기다려 온 부산물이다. 모든 프롬프트 토큰의 K와 V 텐서를 만들어 이후 루프를 위해 캐시한다(§11.3).

\subsection*{디코드: 끝나지 않는 matvec}

이어서 루프가 시작되면 장치의 종이 바뀐다. 각 디코드 단계는 방금 표본 추출한 \textbf{토큰 하나}를 전체 네트워크에 넣어 다음 분포를 얻는다. 토큰 하나는 처음부터 끝까지 행렬–\textit{벡터} 연산을 뜻한다. 제10장 첫 워크로드였던 AI ≈ 1, 즉 연산을 곁들인 바이트 배달 서비스다. 단계마다 \textasciitilde{}6 GFLOP을 수행하기 위해 6 GB의 가중치 전체를 다시 스트리밍하므로 §10.2의 하한인 \textbf{6 GB ÷ 300 GB/s = 20 ms}를 매번 부담한다. 여기에 실행과 표본 추출에 1–2밀리초가 더해진다. 한 \textit{시퀀스 안에서는} 단계가 중첩될 수 없다. token t+1의 입력이 바로 token t이므로 루프는 본질적으로 직렬이다. 배치도 컴파일러도 어떤 묘수도 이 의존성을 없애지 못한다. §11.5의 추측 디코딩만 조금 휘게 한다. 디코드가 TPOT를 결정하고 TPOT는 메모리 벽에 고정된다. 제10장이 배치 1 생성을 거듭 "가장 가혹한 모습의 벽"이라 부른 이유다.

\begin{terminalbox}
\begin{Verbatim}[breaklines=true,breakanywhere=true,fontsize=\small,formatcom=\color{tbbtermfg},codes=\tbbcodeactive]
# one instrumented request against a single unloaded L4 replica:
$ python two_clocks.py --url http://chat-l4:8000 --prompt 500tok.txt
tokenize + queue + prefill (500 tok):     73.8 ms   <- TTFT
decode: 200 tokens in 4.28 s:             21.5 ms/token  (46.5 tok/s)

# sanity, against §10.2 (no profiler, as always):
#   decode floor: 6 GB / 300 GB/s        = 20.0 ms  -> observed 21.5
#   overhead: launches+sampling+attention =  1.5 ms  -> a small a
#   prefill: 3 TFLOP / (~55 TFLOPS eff.) = ~55 ms   -> observed ~65
# prefill did 500 tokens of work in ~3.4 decode-tokens' time.
# two machines, one API. every engine decision follows from this.
\end{Verbatim}
\end{terminalbox}
\figcaption{터미널 기록 11-1  두 시계를 측정하고 제10장의 하한과 맞춘 결과. 프리필은 연산이며 디코드 약 3.4단계의 비용으로 500 tokens를 처리한다. 디코드는 대역폭이며 20 ms 하한에 소량이 더해진다. TTFT ≈ 대기열 + 프리필, TPOT ≈ 하한이다.}

\subsection*{두 시계에 SLO를 부여하다 — 스트림에 돌아온 Little의 법칙}

이제 §11.1의 계약을 선언하는 데 그치지 않고 설계할 수 있다. \textbf{TTFT}는 대기열 + 토큰화 + 프리필이다. 하한은 프리필 연산이며 500토큰 프롬프트에서 \textasciitilde{}65 ms다. 따라서 500 ms 예산 중 \textasciitilde{}430 ms를 대기열 처리에 쓸 수 있다. 배치 입장 제어(§11.4)에 실제로 사용할 수 있는 예산이며, 부하가 TPOT보다 TTFT를 먼저 무너뜨리는 이유다. \textbf{TPOT}는 디코드 하한에 스케줄러가 더하는 비용을 합친 값이다. 부하가 없을 때 21.5 ms (46 tok/s)이고 동시 스트림이 대역폭을 나눌수록 점진적으로 나빠진다. §11.3에서 정확히 계산한다. \textbf{종단 간 시간}도 여전히 존재한다. 200토큰 답변은 73.8 ms + 199 × 21.5 ms ≈ \textbf{4.4 s} 뒤에 도착한다. 하지만 이제 답변 길이가 지배하는 \textit{파생} 수치로서 보고는 하되 약속하지 않는다. 스트림이 주는 이점도 보자. 독자가 200토큰을 읽는 데 \textasciitilde{}30초가 걸리지만 시스템은 4.4초 만에 전달한다. 스트림은 이를 읽는 눈보다 일곱 배 앞서간다.

용량 계획도 익숙한 법칙을 통해 살아남는다. 제7장의 \textbf{L = λW}는 "고객"이 무엇인지 개의치 않으므로 다시 정의하면 된다. \textit{스트림} 하나가 W ≈ TTFT + N·TPOT초 동안 좌석을 점유한다. 200토큰 답변에서 W ≈ 4.4 s다. 동시 스트림 32개를 수용하는 레플리카가 정상 상태에서 받아들일 수 있는 도착률은 λ = 32 ÷ 4.4 ≈ \textbf{7.3 requests/s}다. 32가 어디에서 나오는지는 §11.3에서 설명한다. 이를 사례 11-A와 비교해 보자. 실패한 플릿은 사용률이 고정된 상태에서 6.8 req/s로 측정되었다. 모든 것이 "깨진" 날에도 법칙은 잡음 범위 안에서 들어맞았다. 제2장의 플릿 크기를 정한 공식이 여기서도 통한다. W의 의미만 달라졌다.

\begin{calloutbox}{tbbpurple}{tbbpurplefill}
{\normalsize\bfseries\color{tbbpurple}생각해 보기}\par\vspace{3pt}
\begin{tbbitemize}
  \item 터미널 기록 11-1의 프리필은 디코드 3.4토큰 시간에 500토큰 분량의 작업을 했다. L4의 121 TFLOPS를 기준으로 각 단계의 유효 FLOP 사용률을 계산하라. 루프라인 용어로 디코드의 나머지 \textasciitilde{}97\% 연산 용량이 어디로 갔는지 설명하라.
  \item 한 팀원은 "프롬프트를 줄여 TTFT를 낮추자"고 하고 다른 팀원은 "디코드보다 프리필을 먼저 받아 TTFT를 낮추자"고 한다. 각 제안은 누구의 어느 시계에 비용을 부과하는가? 두 번째 제안을 정식화하는 §11.4의 메커니즘은 무엇인가?
  \item 제10장에서 배포한 INT8 변형(3 GB 가중치, 같은 GPU)을 대상으로 터미널 기록 11-1을 다시 계산하라. 어느 시계가 얼마나 개선되며, 다른 시계는 왜 거의 움직이지 않는가? 힌트: 프리필은 연산 병목이다.
  \item L = λW로 전체 플릿을 계산하라. 동시 스트림 100개가 필요하고 W ≈ 4.4 s다. 제품이 감당할 수 있는 도착률은 얼마인가? 이제 평균 답변 길이를 두 배로 늘려라. 세 SLO 조항 중 무엇이 먼저 깨지며, 이를 해결하려면 무엇을 구매해야 하는가?
\end{tbbitemize}
\end{calloutbox}

\section{KV 캐시: 메모리가 새로운 동시성이다}

사례 11-B의 서버는 "무언가를 기억"하고 있었다. 이제 그 정체를 알아보자. 모든 트랜스포머 계층 안에서 어텐션(attention)은 현재 토큰이 이전의 모든 토큰을 참고하게 한다. 기계적으로 현재 토큰의 \textit{쿼리} 벡터를 각 과거 토큰의 \textbf{키} 벡터와 비교하고 그 \textbf{값} 벡터를 섞는다. 각 토큰이 위치에 고정되고 나면 K와 V 벡터는 그 토큰의 결정론적 함수가 되므로 두 선택지가 생긴다. 매 단계 \textbf{다시 계산}하면 t단계에서 이미 알던 값을 다시 얻기 위해 이전 t−1개 토큰 전체를 모든 계층에 다시 통과시켜야 한다. 대화 중간쯤에는 새 토큰마다 프리필 한 번의 비용이 들고, 2,000토큰 대화는 약 \textit{200만 번}의 불필요한 토큰-계층 연산을 수행한다. 아니면 \textbf{캐시}할 수 있다. 각 토큰의 K와 V를 정확히 한 번 계산해 저장하고 이후 모든 단계가 VRAM에서 읽게 한다. 세상의 모든 서빙 스택은 캐시를 택한다. 루프의 토큰당 연산은 일정해지는 대신 대화가 살아 있는 동안 토큰마다, 계층마다 K,V 행 하나씩 메모리가 늘어난다. 사례에서 본 점진적 증가가 이것이다.

\subsection*{모델 카드에서 크기 계산하기 — 속설은 필요 없다}

토큰당 비용은 아키텍처에서 바로 읽을 수 있다. 책의 3B 모델은 28개 계층을 갖고, 거의 모든 현대 LLM처럼 \textbf{그룹 쿼리 어텐션(GQA)}을 사용한다. 쿼리 헤드 24개가 각각 128차원인 KV 헤드 8개를 \textit{공유}한다. 연산이 아니라 캐시를 3× 줄이기 위한 구조다. 계층마다 토큰 하나에 8 × 128개 값의 K 벡터 하나와 V 벡터 하나를 저장한다. FP16에서는 2 × 28 × 8 × 128 × 2 bytes = \textbf{114,688 bytes ≈ 112 KB per token}이다. 따라서 4,096토큰 대화는 4,096 × 114,688 bytes = \textbf{469,762,048 bytes = 448.00 MiB}를 차지한다. 사례 11-B에서 요청하다 실패한 할당량과 바이트 단위까지 일치한다. 수수께끼의 원인은 누수가 아니었다. 대화 30개가 아무도 예산을 잡지 않은 정당한 메모리를 각각 최대 0.5기가바이트씩 들고, 한 번에 112 KB 행 하나만큼 자라고 있었다.

\begin{terminalbox}
\begin{Verbatim}[breaklines=true,breakanywhere=true,fontsize=\small,formatcom=\color{tbbtermfg},codes=\tbbcodeactive]
$ python kv_math.py --layers 28 --kv-heads 8 --head-dim 128 --dtype fp16
per token:            2 x 28 x 8 x 128 x 2 B = 114,688 B  (112 KB)
per 4k conversation:  4096 x 112 KB           = 448.00 MiB   # !!

# now budget the L4 (24 GB), vLLM-style at 90% utilization:
  usable            24.0 x 0.90 = 21.6 GB
  weights (FP16)                -  6.0 GB   # resident forever
  activations + CUDA + runtime  -  2.5 GB
  KV pool                       = 13.1 GB  -> ~115,000 tokens
# = 28 conversations at 4k. or 115 at 1k. THAT is your concurrency.
# note what is NOT in this arithmetic: FLOPS. compute stopped being
# the binding resource the moment the model started talking back.
\end{Verbatim}
\end{terminalbox}
\figcaption{터미널 기록 11-2  모델 카드에서 계산한 캐시. 사례 11-B의 448 MiB를 관찰이 아니라 도출했다. L4의 좌석 수는 13 GB 풀을 토큰당 112 KB로 나누는 문제다. 동시성이 메모리 수치가 되었다.}

\begin{tbbfigure}
  \adjustbox{max width=0.94\linewidth,max totalheight=0.46\textheight}{\includegraphics{figures/fig11-3_kvcache.png}}\\[3pt]
  \figcaption{그림 11-3  캐시가 존재하는 이유와 토큰당 비용, L4의 24 GB를 가중치, 런타임, KV 풀로 나눈 예산. GQA가 이미 비용을 3× 줄였고 다른 경량화 기법이 더 줄인다. §11.5는 남은 공간의 낭비를 막는다.}
\end{tbbfigure}

\subsection*{캐시가 디코드 단계에 미치는 영향 — 증가하는 b}

캐시는 점유 공간일 뿐 아니라 \textbf{트래픽}이다. 어텐션이 과거 전체를 참고하므로 각 디코드 단계는 시퀀스가 보유한 캐시의 모든 K,V를 \textit{읽어야} 한다. 한 단계의 바이트는 가중치 + 배치의 전체 캐시다. 제9장의 법칙을 조금 바꾸어 \textbf{t(step) ≈ (W\_bytes + Σ KV\_bytes(context)) ÷ bandwidth}라고 쓰자. 평균 컨텍스트가 4k인 배치 32에서 책의 모델을 실행하면 bytes = 6 GB + 32 × 0.47 GB ≈ \textbf{21 GB per step → \textasciitilde{}70 ms → \textasciitilde{}14 tokens/s per stream, \textasciitilde{}460 tokens/s aggregate}다. 배치 1의 6.47 GB → 21.5 ms → 46 tok/s와 비교해 보자. 제9장의 거래를 스트림에 맞게 다시 표현하면 \textbf{스트림당 지연 시간을 3× 늘리는 대신 집계 처리량을 10× 얻는 것}이다. 모든 스트림은 여전히 읽기 속도보다 빠르다. 새 항의 동작에 주목하자. 가중치는 같은 6 GB로 모두를 처리해 배치 전체에 상각된다. 이는 9주 차의 a다. 반면 \textbf{KV 바이트는 배치에 따라 증가}한다. 각 스트림이 자체 과거를 끌고 오므로 컨텍스트와 함께 커지는 b다. n × 0.47 = 6으로 놓으면 교차점을 찾을 수 있다. \textbf{동시 4k 대화가 \textasciitilde{}13개를 넘으면 캐시가 모델보다 무겁다.} 규모가 커지면 LLM 서버는 답을 만드는 네트워크보다 대화를 읽는 데 더 많은 대역폭을 쓴다.

이 계산 하나가 현대 LLM 서빙 연구 과제의 절반을 조용히 설명하므로 그 귀결을 더 살펴보자. \textit{긴 컨텍스트는 모두에게 부과되는 세금}이다. 32k-token 문서 대화 하나는 \textbf{모든 단계에서} 3.8 GB의 캐시를 벽 너머로 끌고 가 함께 배치된 모든 스트림을 늦춘다. 길이 제한과 컨텍스트 가격 책정은 인색함이 아니라 물리학이다. \textit{캐시 때문에 GQA/MQA가 존재한다.} Shazeer는 2019년 MHA의 많은 KV 헤드가 추론에서 무엇을 사는지 묻고 "대부분 바이트"라고 답했다. \textit{KV 캐시 양자화}가 표준 수단인 이유도 같다. FP8은 112 KB를 절반으로 줄이며, §10.3의 게이트는 가중치와 똑같이 캐시에 적용된다. \textit{슬라이딩 윈도 모델}이 과거 범위를 아예 제한하는 이유도 여기에 있다. 입장 문제도 더 까다로워진다. 새 스트림은 \textbf{예측할 수 없이 커질} 풀 공간을 필요로 한다. 풀이 가득 차면 엔진은 입장을 대기열에 두거나 실행 중 스트림을 \textbf{선점(preemption)}해야 한다. 블록을 퇴거시키고 나중에 다시 프리필하여 메모리를 되찾는 대신 연산 비용을 낸다. 대기, 선점, 거부 중 \textit{무엇을} 고를지가 스케줄링이며 §11.4의 주제다.

\subsection*{다중 턴 채팅: 캐시와 상태 비저장 웹의 만남}

배처로 넘어가기 전에 제6장의 빚을 갚는 귀결 하나를 더 살펴보자. HTTP API는 상태를 저장하지 않으므로 대화의 각 턴이 \textbf{전체 기록}을 프롬프트로 다시 보낸다. 기계적으로 서버가 매 턴 전체 대화를 다시 \textit{프리필}한다는 뜻이다. 3,000토큰 대화의 5번째 턴은 4턴 전에 똑같이 계산한 K,V 행을 다시 만들기 위해 3,000토큰 프리필 비용, 즉 순수 연산 \textasciitilde{}350 ms를 낸다. 누군가 보관해 두지 않았다면 그렇다. 이를 보관하는 방식을 \textbf{접두사 캐싱(prefix caching)}이라 하며 §11.5의 블록 메커니즘은 거의 공짜로 만든다. 하지만 제6장의 고정 세션 긴장을 정확히 되살린다. 보관한 행은 \textbf{한 레플리카의 VRAM}에 있으므로 5번째 턴을 다른 레플리카로 라우팅하면 조용히 전체 프리필 비용을 다시 낸다. 이 긴장을 기억해 두자. §11.5에서는 어피니티를 정확성 요구사항이 아닌 \textit{성능 힌트}로만 취급하는 성숙한 방식으로 해결한다. 아래 생각해 볼 문제에서 먼저 양쪽 비용을 계산한다.

\begin{calloutbox}{tbbpurple}{tbbpurplefill}
{\normalsize\bfseries\color{tbbpurple}생각해 보기}\par\vspace{3pt}
\begin{tbbitemize}
  \item 사례 11-B의 시간선을 정량적으로 재구성하라. 24 GB 카드에서 9.8 GB를 바닥으로 삼고 대화 \textasciitilde{}30개가 각각 \textasciitilde{}2 tokens/s씩 자란다. 실패할 시각을 예측한 뒤 실제 곡선이 더 울퉁불퉁한 이유를 설명하라. 힌트: 답변은 끝나고 사용자는 멈추며 새 입장이 도착한다.
  \item PM은 "모델이 지원하니까" 128k 컨텍스트 요금제를 원한다. 책의 3B 모델에서 128k 대화 하나가 차지하는 KV 바이트, 배치 8의 단계별 트래픽, 퇴거되는 좌석 수로 비용을 계산하라. 이어 요금제의 가격 구조와 입장 규칙을 제안하라.
  \item 터미널 기록 11-2의 좌석 수를 (a) 제10장에서 배포한 INT8 가중치, (b) 여기에 FP8 KV 캐시 추가, (c) 둘에 GQA→MQA(KV 헤드 1개)까지 적용한 경우에 대해 다시 계산하라. 어느 단일 변경이 가장 많은 좌석을 사며, 무엇에 §10.3의 게이트가 가장 시급한가?
  \item 단계 비용 모델에 따르면 32k 문서 대화 하나를 짧은 대화 31개와 함께 배치하면 모두 느려진다. 큰 대화가 있을 때와 없을 때 각 스트림의 TPOT를 계산한 뒤 실제로 배포할 공정성 정책을 설계하라. 좌석별 KV 할당량, 별도 풀, 가격 책정을 고려할 수 있다.
  \item 제3장은 페이지 캐시를 "회수 가능한 짐", 익명 메모리를 "필수 메모리"라고 설명했다. KV 캐시는 어느 쪽인가? 답을 바탕으로 메모리 압박에서 엔진이 웹 서버에는 불가능했던 어떤 작업을 할 수 있는지 설명하라. §11.4의 선점이 그 결론이다.
\end{tbbitemize}
\end{calloutbox}

\section{연속 배치: 시내버스}

배치는 여전히 경제성의 전부다. §11.3의 계산에서 스트림 하나는 가중치 스트림 가치의 \textasciitilde{}97\%를 쓰지 못하지만 동승자 30명이 이를 회수했다. 그렇다면 9주 차의 동적 배처로 생성 요청을 모으면 될 것 같지만, 작은 글씨의 조항 하나 때문에 실패한다. 상자 9-1은 동적 배치가 \textbf{형상이 비슷한 요청}을 가정한다고 경고했다. 제10장의 규칙 ③은 \textit{입력}의 제각각인 길이에 패딩 비용을 부과하고 버킷팅을 제시했다. 생성은 버킷팅으로 해결할 수 없는 문제를 만든다. \textbf{출력 길이가 제각각이고 미리 알 수도 없다.} 요청 네 개가 함께 들어와 배치가 공유 디코드 루프를 돈다. 하나는 token 32에, 다른 하나는 210에 끝난다. \textit{정적} 배치에서는 일찍 끝난 요청의 좌석이 가장 긴 승객이 끝날 때까지 \textbf{죽은 채} 따라다닌다. 빈 좌석으로 가득한 버스 밖에는 새 요청이 줄을 선다. 수치는 명확하다. 답변 길이가 10–100× 차이 나면 좌석-단계 사용률은 약 50\%에 그친다. 더 가혹하게도 새 요청의 TTFT에는 먼저 탄 사람이 쓰고 있는 \textit{남은 글 전체}가 포함된다. 사례 11-A의 11.4초 p99에는 또 하나의 원인이 있었다.

\begin{tbbfigure}
  \adjustbox{max width=0.94\linewidth,max totalheight=0.46\textheight}{\includegraphics{figures/fig11-4_contbatch.png}}\\[3pt]
  \figcaption{그림 11-4  전세버스 대 시내버스. 정적 배치는 모두 함께 끝나거나 아무도 끝나지 않는다. 죽은 좌석과 대기 중 도착이 생기고 TTFT는 낯선 사람의 긴 글에 붙잡힌다. 반복 수준 스케줄링은 \textasciitilde{}20 ms 단계마다 승객 명단을 다시 정한다. 끝난 시퀀스가 내리고 기다리던 시퀀스가 타므로 대기열이 있는 동안 어느 좌석도 쉬지 않는다.}
\end{tbbfigure}

\subsection*{반복 수준 스케줄링 — Orca의 아름다운 아이디어 하나}

해결책은 Yu 등의 \textbf{Orca}(OSDI 2022)가 제시했다. 제9장의 더 읽을거리에서 "KV 캐시 강의 뒤"까지 미뤄 둔 논문인데, 그 강의가 20분 전에 끝났다. 해결책은 한 문장이다. \textbf{요청 스케줄링을 멈추고 반복을 스케줄링하라.} 디코드 루프는 이미 단계별로 실행되므로 \textit{어떤} 두 단계 사이에서도 배치 구성원을 바꿀 수 있다. 스케줄러는 \textasciitilde{}20 ms마다 세 박자의 주기를 실행한다. 방금 종료 토큰을 낸 시퀀스를 \textbf{퇴거}해 좌석과 KV 블록을 즉시 비운다. KV 풀이 수용할 수 있으면 기다리던 시퀀스를 \textbf{입장}시킨다. 이제 §11.3의 입장 규칙이 시스템을 지탱한다. 그리고 탑승한 모든 시퀀스를 순방향 패스 한 번으로 \textbf{진행}한다. 이것이 연속 배치다. "배치"는 시작과 끝이 있는 사건이 아니라 매 정류장에서 인구 조사가 달라지는 \textbf{집단}이 된다. 전세버스가 시내버스로 바뀌며 그림 11-4 위쪽의 두 실패 유형은 아예 일어날 수 없게 된다. 누군가 기다리는 동안 좌석이 쉬지 않고, 새 승객의 대기 시간은 낯선 사람의 글 길이가 아니라 \textit{단계} 단위로 측정된다.

왜 여기서는 \textit{가능}하지만 9주 차에는 불가능했는지 짚고 넘어갈 가치가 있다. CNN 순방향 패스는 \textbf{나눌 수 없는 커널 하나}의 실행이므로 Triton 동적 배처는 실행 도중에 퇴거하거나 입장시킬 수 없었다. 개입할 "단계 사이"가 없다. 반면 토큰마다 한 단계씩 도는 Python에서 보이는 바깥 루프라는 생성의 비효율은 스케줄링의 기회가 된다. 루프의 이음매가 스케줄러의 자리다. Orca의 두 번째 기여인 \textit{선택적 배치}는 길이가 다른 공존 시퀀스의 텐서 처리를 맡는다. 어텐션은 시퀀스별 캐시를 대상으로 시퀀스마다 실행하고, 큰 선형 계층은 모두를 함께 배치한다. vLLM 커널도 이 패턴을 이어받았다. 9주 차의 유산 하나는 변환 과정에서 사라진다. \textbf{max\_queue\_delay가 없어졌다.} 다음 출발까지 최대 한 단계뿐이므로 기다릴 이유가 없다. 후속 조절점은 \textit{한 단계에 얼마나 많은 작업을 태울 것인가}이며, 여기서 두 시계가 싸우기 시작한다.

\subsection*{프리필–디코드 간섭: 두 시계와 하나의 엔진}

새 스트림을 입장시키는 일은 공짜 좌석 배정이 아니다. 새 승객에는 50–350 ms의 연산 병목 GEMM인 \textbf{프리필}이 필요하다. 디코드 스트림 30개가 20 ms씩 나눠 쓰는 바로 그 GPU가 실행한다. 프리필을 나눌 수 없는 한 단계로 실행하면 모든 기존 스트림의 다음 토큰이 뒤에서 기다린다. 새 승객 한 명의 TTFT를 최소화하는 동안 30개 스트림이 \textbf{딸꾹질}하고 TPOT가 21에서 \textasciitilde{}250 ms로 치솟는다. 독자의 눈이 실제로 알아채는 끊김이다. 반대로 프리필을 미루면 디코드는 매끄럽지만 TTFT가 부푼다. 새로운 꼭짓점 이름으로 서빙 삼각형이 다시 태어난다. \textbf{TTFT 대 TPOT 대 비용}이다. 모서리를 따라 움직이는 조절점은 \textbf{청크 프리필(chunked prefill)}이다. 예를 들어 500토큰 프리필을 네 청크로 나누고 디코드 단계마다 한 청크를 함께 스케줄링한다. 새 승객의 프리필은 네 단계에 걸쳐 끝나 TTFT가 +\textasciitilde{}60 ms 늘지만 기존 승객의 TPOT는 단계마다 \textasciitilde{}2 ms만 증가한다. 엔진은 이 절충을 예산으로 노출한다. vLLM의 \code{max\_num\_batched\_tokens}는 \textit{반복마다 총 토큰 수(프리필 청크 + 디코드 단계)}를 제한한다. 이를 튜닝하는 일은 §9.3의 조절점 스윕이 다시 태어난 모습이며, 보고서의 지연 시간 시계가 하나가 아니라 둘이다.

\begin{terminalbox}
\begin{Verbatim}[breaklines=true,breakanywhere=true,fontsize=\small,formatcom=\color{tbbtermfg},codes=\tbbcodeactive]
# vLLM /metrics, sampled every 30 s while load ramps (c = 4 -> 64):
$ watch -n 30 'curl -s :8000/metrics | grep -E "running|waiting|cache"'

#            running   waiting   gpu_cache_usage    <- read as:
t+0:00          4          0          9 %           # roomy. TTFT=prefill
t+1:30         19          0         41 %           # filling. all admitted
t+3:00         31          0         88 %           # full-ish. TPOT sagging
t+4:30         32          9         97 %           # pool full: queue forms
t+6:00         32         26         98 %           # waiting grows: TTFT now
                                                    # queue-dominated (L=λW!)
# the three regimes every LLM service lives in: ROOMY (TTFT=prefill),
# FULL (TPOT shares bandwidth), QUEUED (TTFT explodes first).
# your SLO breaks in regime 3 at the TTFT clause - almost never TPOT.
\end{Verbatim}
\end{terminalbox}
\figcaption{터미널 기록 11-3  자체 메트릭으로 관찰한 연속 배치. running은 KV 풀의 좌석 수까지 올라간 뒤 멈추고, 풀이 차는 순간 waiting이 나타난다. 부하에서 실패하는 순서는 언제나 같다. 제7장의 하키 스틱 예측대로 대기열 때문에 TTFT가 먼저 무너진다.}

\subsection*{조절점은 무엇이 되었는가 — 9주 차 스윕의 변환}

\begin{tbbtablewrap}
  \begin{tabular}{@{}>{\raggedright\arraybackslash}m{\dimexpr 0.2447\tbltextwidth-2\tabcolsep\relax} >{\raggedright\arraybackslash}m{\dimexpr 0.2660\tbltextwidth-2\tabcolsep\relax} >{\raggedright\arraybackslash}m{\dimexpr 0.4894\tbltextwidth-2\tabcolsep\relax}@{}}
  \arrayrulecolor{tbbnavy}\specialrule{1pt}{0pt}{0pt}
  \rowcolor{tbbtablehead}{\centering \bfseries\color{tbbnavy}9주 차(동적 배치)} & {\centering \bfseries\color{tbbnavy}11주 차(연속 배치)} & {\centering \bfseries\color{tbbnavy}달라진 점과 이유} \\
  \arrayrulecolor{tbbnavy}\specialrule{0.7pt}{0pt}{2pt}
  {max\_batch\_size} & {\textbf{max\_num\_seqs}} & {좌석 수 — 다만 보통 KV 풀이 먼저 구속한다(§11.3). 주로 최악의 간섭을 제한한다} \\
  \arrayrulecolor{tbbtableborder}\hline
  {max\_queue\_delay} & {\textbf{(사라짐)}} & {기다릴 이유가 없다. 다음 출발까지 한 단계(\textasciitilde{}20 ms)뿐이다. 성공해서 사라진 조절점이다} \\
  \arrayrulecolor{tbbtableborder}\hline
  {(대응 항목 없음)} & {\textbf{max\_num\_batched\_tokens}} & {단계 예산: 반복마다 프리필 청크 + 디코드 토큰. TTFT ↔ TPOT 다이얼이다} \\
  \arrayrulecolor{tbbtableborder}\hline
  {memory.max (제5장, 외부)} & {\textbf{gpu\_memory\_utilization}} & {제한이 엔진 안으로 이동하며 반전되었다. 예산을 부여하면 엔진이 자체 좌석 수를 계산한다} \\
  \arrayrulecolor{tbbtableborder}\hline
  {(오버플로 시 충돌)} & {\textbf{선점 정책}} & {초과 인출 계획: 시퀀스 블록을 퇴거하고 나중에 다시 계산한다. 연산을 써서 메모리를 회수한다} \\
  \arrayrulecolor{tbbnavy}\specialrule{1pt}{2pt}{0pt}
  \end{tabular}\par
  \figcaption{표 11-1  변환된 배처 패널. 9주 차의 모든 직관은 11주 차에도 자리를 찾았지만 단위는 요청에서 토큰으로, 구속 자원은 연산에서 KV 풀로 이동했다.}
\end{tbbtablewrap}

§9.3의 배처에는 조절점이 두 개 있었고 연속 배처의 패널이 이를 변환한다. 좌석 수인 \textbf{max\_num\_seqs}가 max\_batch\_size를 잇는다. 하지만 §11.3에서 실제 상한은 대개 KV 풀임을 보았으므로 이 조절점은 주로 최악의 간섭을 제한한다. 단계 예산인 \textbf{max\_num\_batched\_tokens}는 TTFT↔TPOT 다이얼로서 max\_queue\_delay를 잇는다. \textbf{gpu\_memory\_utilization}은 풀 자체의 크기를 정한다. 제5장의 memory.max 규율이 이제 엔진 \textit{안쪽}으로 들어가 기분 좋게 반전되었다. 최악의 사용량을 측정해 제한을 정하는 대신 예산을 \textit{부여}하면 엔진이 자체 좌석 수를 계산한다. \textbf{선점 정책}은 풀의 초과 인출을 처리한다. vLLM은 시퀀스 전체의 블록을 퇴거하고 공간이 생기면 나중에 프리필을 \textbf{다시 계산}한다. 연산 비용을 내고 메모리를 회수하는 제3장 생각해 볼 문제의 답이 산업화되었다. 블록을 호스트 RAM으로 스왑하는 방법도 있지만 §10.2의 PCIe 계산을 보면 보통 재계산이 유리한 이유를 알 수 있다. 실습 스윕(§11.6)은 이 패널을 그대로 훑는다. 결과물은 9주 차 이후 모든 스윕과 같다. 의도적으로 찾은 굴곡점을 기록한다.

선점은 실제 모습으로 한 번 볼 가치가 있다. 개념을 배우기 전에 로그 줄로 먼저 만나게 되며, 이 문단을 읽지 않으면 오류로 오해할 수 있기 때문이다.

\begin{terminalbox}
\begin{Verbatim}[breaklines=true,breakanywhere=true,fontsize=\small,formatcom=\color{tbbtermfg},codes=\tbbcodeactive]
# push admission past the pool on purpose (c=48, long answers) and
# watch the engine's overdraft plan execute:
WARNING  Sequence group chatcmpl-8f21 is preempted by
         PreemptionMode.RECOMPUTE because there is not enough KV
         cache space. total_num_cumulative_preemption=1
INFO     Avg generation throughput: 691.2 tokens/s,
         Running: 31 reqs, Waiting: 17 reqs, GPU KV cache usage: 99.8%
# not a crash: one conversation's blocks were evicted; it will
# re-prefill when space frees. the victim's next tokens stall ~400 ms.
# occasional = healthy elasticity. CONSTANT = your pool is undersized:
# quantize the weights, shrink max-model-len, or buy the bigger card.
\end{Verbatim}
\end{terminalbox}
\figcaption{터미널 기록 11-4  실제 선점. WARNING은 설계가 작동한다는 뜻이다. 제3장의 OOM killer에는 더 부드러운 선택지가 없었지만 블록 관리자는 퇴거, 대기, 재프리필을 할 수 있다. 카운터를 활력 징후처럼 읽어라. 0에 가까우면 탄력적인 상태이고 계속 오르면 용량이 부족하다.}

정직하게 한계를 짚어 보자. 연속 배치는 \textit{길이가 제각각인 상황의 처리량}을 최적화한다. §11.3의 대역폭 계산을 폐기하지 않는다. 승객들이 메모리 버스 하나를 공유하므로 배치 1에서 32로 가면 TPOT는 여전히 \textasciitilde{}3× 나빠진다. \textbf{동승자가 없는} 워크로드에도 도움이 되지 않는다. 사용자 한 명의 에이전트 루프가 영원히 배치 1이라면 메모리 벽의 포로로 남는다. §11.5의 추측 디코딩이 바로 이 경우를 위해 존재한다. 엔진이 주는 선물은 더 좁지만 엄청나다. \textit{누군가 기다리는 동안 비용을 지불한 좌석이 절대 비지 않게} 보장한다. 좌석이 충분히 저렴한지는 여전히 제10장과 제12장의 몫이다.

\begin{calloutbox}{tbbpurple}{tbbpurplefill}
{\normalsize\bfseries\color{tbbpurple}생각해 보기}\par\vspace{3pt}
\begin{tbbitemize}
  \item 답변 길이가 32/180/60/210 tokens인 정적 배치(static batching) 4를 생각하라(그림 11-4). 정확한 좌석-단계 낭비와 10단계에 도착한 요청의 TTFT를 계산하고, 연속 배치에서 두 수치를 다시 계산하라. 이어 입력 길이별 버킷팅(제10장 규칙 ③)이 이 문제를 전혀 해결하지 못하는 이유를 설명하라.
  \item 2,000토큰 프롬프트에 chunk = 512 tokens인 청크 프리필을 적용한다. decode step ≈ 21 ms, prefill chunk ≈ 55 ms를 함께 스케줄링할 때 새 승객에게 추가되는 TTFT와 네 혼합 단계 동안 기존 승객의 TPOT를 추정하라. 이제 반대편을 주장하라. 언제 의도적으로 전체 프롬프트 프리필을 실행하겠는가? 힌트: 전용 프리필 레플리카는 상자 11-1의 DistServe 아이디어다.
  \item 터미널 기록 11-3의 구간 3은 running이 32에 고정되고 waiting=26인 상태다. W ≈ 4.4 s에서 L = λW를 사용해 도착률, 대기열 증가율, 26번째 대기자의 TTFT를 추정하라. 실패를 보고하는 SLO 조항은 무엇이며, 제12장은 어떤 해결책을 판매하는가?
  \item vLLM은 퇴거 후 재계산으로 선점한다. 3,000토큰 대화를 선점할 때 (a) 낭비되는 프리필 연산, (b) 피해자의 다음 턴 TTFT, (c) 공정성 측면에서 비용을 계산하라. 가장 최근, 가장 긴, 재계산 비용이 가장 싼 시퀀스 중 누구를 선점할지 정책을 설계하고 피해자에게 정당화하라.
\end{tbbitemize}
\end{calloutbox}

\section{PagedAttention과 vLLM: 운영체제가 안으로 들어가다}

지금까지 좌석이 귀중하다는 사실을 거듭 확인했다. §11.3은 KV 바이트로 가격을 매겼고 §11.4는 어느 좌석도 쉬지 않게 하겠다고 약속했다. 그러니 2023년 이전 엔진이 캐시를 \textit{저장}한 방식을 알면 아플 수밖에 없다. 입장할 때 시퀀스마다 \textbf{가능한 최대 길이에 맞춘 하나의 연속 슬랩}을 할당했다. 어텐션 커널이 깔끔한 스트라이드를 원해 연속 공간을 썼고, 입장 시 답변이 30 tokens일지 3,000일지 알 수 없어 최대 크기로 잡았다. 이미 본 영화다. malloc의 가장 오래된 비극이 예상대로 되풀이되었다. \textbf{내부 낭비}는 슬랩마다 예약했지만 쓰지 않는 꼬리다. 2,048을 예약해 180을 쓰면 91\%를 낭비한다. \textbf{외부 단편화}는 해제된 슬랩이 새 최대 크기 요청을 넣을 수 없는 틈을 남기는 현상이다. vLLM 저자들이 SOSP 2023 논문을 위해 프로덕션급 시스템을 측정했을 때 결과는 충격적이었다. \textbf{KV 메모리의 60–80\%가 유용한 것을 전혀 담지 않았다.} 애써 확보한 13 GB 풀이 4 GB처럼 동작했다. 버스 좌석의 절반 이상을 뜯어내 공기를 보관한 셈이다.

\begin{tbbfigure}
  \adjustbox{max width=0.94\linewidth,max totalheight=0.46\textheight}{\includegraphics{figures/fig11-5_paged.png}}\\[3pt]
  \figcaption{그림 11-5  malloc 시대 대 페이징 시대. 연속 최대 길이 슬랩은 예약 꼬리와 쓸 수 없는 틈 때문에 풀의 60–80\%를 낭비한다. 고정 16토큰 블록과 시퀀스별 블록 테이블은 낭비를 4\% 아래로 줄이고, 복사 대신 참조 횟수로 공유할 수 있게 한다.}
\end{tbbfigure}

\subsection*{60년 된 해결책을 300 GB/s에 적용하다}

해결책은 컴퓨팅 역사상 가장 성공적인 간접 참조 기법이며, 이 책에서 이미 두 번 보았다. 운영체제는 1960년대부터 프로세스에 연속 물리 메모리를 주지 않았다. 고정 크기 \textbf{페이지}를 나눠 주고 \textbf{페이지 테이블}이 프로세스의 정돈된 가상 주소를 흩어진 물리 프레임에 대응시키게 했다. 그러면 단편화는 "마지막 페이지의 일부"로 줄어든다. 3주 차의 커널은 지금까지 실행한 모든 실습 아래에서 이를 수행한다. 2주 차에는 이 기법을 \textit{쌓아 올린} 모습도 보았다. EPT가 게스트 전체의 "물리" 메모리를 또 하나의 프로세스처럼 페이징했다. \textbf{PagedAttention}(Kwon et al., SOSP 2023)은 이를 통째로 어텐션에 이식한다. 제1장에서 필독으로 지정했고 이 책이 vLLM을 가르치는 이유가 된 논문이다. KV 풀을 \textbf{16 tokens}의 고정 블록으로 나누고, 각 시퀀스에 논리 위치 → 물리 블록을 대응시키는 \textbf{블록 테이블(block table)}을 주며, 어텐션 커널이 테이블을 따라가게 한다. 이제 시퀀스는 한 번에 \textit{블록} 하나씩 자라고 블록은 어디에나 놓일 수 있다. 낭비는 각 시퀀스의 덜 찬 마지막 블록으로 줄어 논문에서 \textbf{4\% 미만}으로 측정되었다. 모델도 정밀도도 바꾸지 않고 장부 관리만 바꿨다. 좌석의 효과가 누적되므로 처리량 개선은 제10장만큼 크다. 낭비를 회수하면 같은 풀에 2–3× 더 많은 동시 시퀀스가 들어가고, 그만큼 많은 승객이 매 20 ms 가중치 스트림을 상각한다. 종단 간 측정에서 \textbf{같은 지연 시간으로 당시 가장 강력한 시스템보다 2–4×, 단순 HF 서빙보다 14–24× 높은 처리량}을 냈다. 해당 세대 최고의 서빙 시스템 논문은 본질적으로 메모리 할당기에 관한 논문이다. cgroup을 다룬 제3장부터 이 책이 강조해 온 감각과 정확히 같다.

\subsection*{공유: 블록 테이블의 두 번째 묘기 — 제6장의 빚을 갚다}

페이지 테이블은 단편화를 없애는 것보다 더 많은 일을 운영체제에 해 주었다. 두 프로세스가 같은 물리 프레임을 매핑한 덕분에 Unix 공유 라이브러리와 쓰기 시 복사 fork가 가능해졌다. 블록 테이블도 이를 재현한다. 제품의 모든 대화가 같은 \textbf{900-token 시스템 프롬프트}로 시작한다고 하자. 연속 슬랩은 이를 32번 저장한다. 32 × 900 × 112 KB ≈ \textbf{동일한 바이트 3.2 GB}로 4k 좌석 거의 일곱 개가 순수 중복에 쓰인다. 페이징하면 900토큰은 \textbf{참조 횟수}가 붙은 \textasciitilde{}57개 블록이 된다. 모든 대화의 테이블이 \textit{같은} 물리 블록을 가리키므로 \textbf{한 번만} 저장한다(≈ 101 MB). 갈라지는 시퀀스에는 블록 단위 \textbf{쓰기 시 복사}를 적용한다. 이것이 \textbf{접두사 캐싱}이며 두 번 보상한다. 먼저 메모리 좌석 6개 이상을 회수한다. 그리고 \textit{시계}도 줄인다. 접두사가 이미 상주한 새 대화는 \textbf{그 접두사의 프리필을 완전히 건너뛰어} TTFT에서 연산 \textasciitilde{}100 ms를 없앤다. 돌아온 대화의 5번째 턴은 3,000토큰 기록 전체가 아니라 새 접미사만 다시 프리필한다. §11.3의 다중 턴 세금을 대부분 돌려받는다.

이로써 제6장에서 보류해 둔 고정 세션의 긴장을 해결할 수 있다. 캐시된 블록은 \textbf{한 레플리카의 VRAM}에 산다. 그렇다면 라우터가 대화를 레플리카에 고정해야 할까? 업계가 수렴한 성숙한 답은 \textbf{정확성은 상태 비저장으로 유지하고 어피니티는 성능 힌트로만 쓴다}는 것이다. 요청은 여전히 전체 기록을 담으므로 \textit{어떤} 레플리카도 5번째 턴을 정확히 처리할 수 있다. 캐시 미스는 접미사 하나가 아니라 \textit{전체} 프리필(\textasciitilde{}350 ms)의 비용을 내게 하지만 잘못된 답을 만들지는 않는다. 따라서 대화 id에 일관된 해시를 적용하는 세션 해시 라우팅은 편할 때 순수한 이점이고, 그렇지 않아도 해 없이 성능만 저하된다. 롤아웃, 스케일 다운, 레플리카 사망이라는 제6장 고정 상태의 세 재앙도 접두사 적중률을 잠시 낮출 뿐이다. 제6장에서 경고한 설계, 즉 잃으면 대화가 깨지는 세션 \textit{저장소}와 비교해 보자. "고정"이라는 말은 같지만 엔지니어링은 반대다. \textbf{캐시를 계약으로 만들지 마라.}

\subsection*{기관별로 살펴보는 vLLM — 추측이라는 수단}

이 장의 내용을 조립하면 vLLM의 블록 다이어그램(그림 11-6)이 나온다. 제9장의 여섯 기관 중 두 기관을 이식한 구조다. \textbf{배처}는 §11.4의 반복 스케줄러다. \textbf{모델 저장소}의 핫 경로에는 §11.5의 테이블, 참조 횟수, 선점을 실행하는 VRAM 속 "운영체제"인 블록 관리자가 생겼다. 주변은 익숙하다. OpenAI 호환 프런트 도어가 SSE로 통신하므로 제6장의 클라이언트, LB, 타임아웃 교훈이 그대로 적용된다. 제6장의 safetensors 로딩과 터미널 기록 11-3에서 이미 이름을 읽은 Prometheus 메트릭도 있다. 실행기 안에서는 제10장의 두 기법이 제값을 한다. 디코드 단계는 \textbf{CUDA Graphs}로 캡처한다. 3B 디코드 단계는 21 ms 예산 안에서 수백 개의 작은 커널을 실행하므로 실행 오버헤드가 중요한 적이다. 그래프가 전체 단계를 제출 한 번으로 재생해 제10장의 \textit{a}를 빌드가 아니라 런타임에서 줄인다. 가중치 슬롯은 제10장에서 그룹별 스케일 용어를 배운 \textbf{GPTQ/AWQ INT4} 산출물을 받는다. 책의 모델은 INT4에서 \textasciitilde{}1.9 GB이므로 배치 1 단계 트래픽이 \textasciitilde{}2.4 GB → \textbf{\textasciitilde{}8 ms → tokens/s 2.5×}로 줄고, KV 풀은 4 GB 넓어져 4k 좌석이 ≈ 9개 늘어난다. 게이트를 거친 산출물 하나가 두 시계와 좌석 수를 모두 개선한다. 퍼플렉시티나 작업 지표를 명시한 허용 오차와 비교하는 §10.3의 게이트는 그대로 적용된다.

실행기의 한 줄은 별도 문단을 쓸 가치가 있다. 제10장의 마지막 약속을 갚고 §11.2의 "본질적으로 직렬"이라는 규칙을 조금 휘게 하는 \textbf{추측 디코딩}이다. 핵심은 메모리 벽의 비대칭에서 차익을 얻는 것이다. 디코드 단계의 20 ms는 토큰 하나의 \textit{대역폭}을 사지만 수백 토큰 분량의 \textit{FLOP}도 함께 산다. AI ≈ 1이고 능선이 400이므로 연산 비용은 이미 냈지만 유휴 상태다. 아주 작은 \textbf{초안 모델(draft model)}(예: \textasciitilde{}160M, \textasciitilde{}1 ms/token)이 k = 4 tokens를 저렴하게 추측하게 한 뒤, 목표 모델이 \textbf{네 개를 한 패스에서 모두 검증}하게 한다. 토큰 \textit{하나}를 생성할 때와 같은 6 GB를 스트리밍한다. 가장 긴 올바른 접두사를 받아들이고 표본 추출을 보정해 출력 분포가 \textit{증명 가능하게 동일}하도록 만든다. 품질 절충이 아니라 시스템 기법이다. 일반 텍스트에서는 네 개 중 \textasciitilde{}3개가 살아남는다. 유효 비용은 (4 × 1 ms 초안) + (21.5 ms 검증 한 번) ≈ \textasciitilde{}3 tokens당 25.5 ms이므로 \textbf{\textasciitilde{}8.5 ms/token, 배치 1에서 TPOT 2.5× 개선}이다. 전부 낭비될 뻔한 FLOP으로 샀다. 정직한 작은 글씨도 있다. 어렵거나 특이한 텍스트에서는 수락률이 무너져 초안 비용은 내면서 속도 향상 → 1×가 된다. 초안 모델도 VRAM을 차지하며, 배치가 높을수록 유휴 FLOP 보조금이 줄어든다. 그래서 추측은 \textit{단일 사용자 에이전트 루프}의 수단이고, 다수 사용자를 위한 연속 배치와 대체로 상호 보완적이다.

\begin{calloutbox}{tbbpurple}{tbbxF5F3FF}
{\normalsize\bfseries\color{tbbx5B21B6}상자 11-1  더 넓은 LLM 서빙 선택지 — 이 책에서 선반에 남겨 둔 수단}\par\vspace{3pt}
남은 주에는 제12장이 필요하므로 여기서는 각 기법을 한 줄로 정직하게 설명한다. 나중에 만났을 때 아이디어와 대가를 알아볼 정도면 충분하다.

\begin{tbbitemize}
  \item \textbf{분리형 프리필(DistServe 방식):} 프리필과 디코드를 \textit{별도의 레플리카 풀}에서 실행한다. §11.2의 두 장치를 물리적으로 분리해 서로 간섭하지 않게 한다. 그 대가로 프리필한 KV 캐시를 네트워크 너머로 전달해야 한다. §10.2의 PCIe 교훈을 클러스터에 적용한 셈이다.
  \item \textbf{Multi-LoRA 서빙:} 하나의 기본 모델 가중치와 KV 풀을 수십 개의 미세 조정 어댑터가 공유하는 S-LoRA 기법이다. 어댑터는 수 메가바이트이므로 "고객마다 모델 하나"가 더 이상 "고객마다 GPU 하나"를 뜻하지 않는다.
  \item \textbf{구조화/유도 디코딩:} 문법이나 JSON 스키마로 표본 추출을 제한한다. 단계별 토큰 마스크에 작은 비용을 내고 에이전트 파이프라인의 출력 유효성을 얻는다. SGLang의 RadixAttention은 프로그램 형태 호출 사이에서 적극적인 접두사 재사용과 이를 결합한다.
  \item \textbf{KV 오프로딩과 계층화:} 차가운 대화의 블록을 호스트 RAM이나 NVMe로 내보내고 다음 턴에 복원한다. §10.2의 PCIe 바이트와 §11.3의 재프리필 연산을 절충한다. 블록 테이블 덕분에 말 그대로 페이징 정책이 된다.
\end{tbbitemize}

모든 기법은 이 장의 세 기본 요소인 두 단계 루프, 블록 테이블, 반복 스케줄러를 다시 조합한 것이다. 기본 요소를 익혀라. 제품은 낡아도 기본 요소는 낡지 않는다.
\end{calloutbox}

\vspace{3.0pt}

\begin{tbbfigure}
  \adjustbox{max width=0.94\linewidth,max totalheight=0.46\textheight}{\includegraphics{figures/fig11-6_vllm.png}}\\[3pt]
  \figcaption{그림 11-6  제9장의 구조를 다시 만든 vLLM. 배처는 토큰마다 스케줄링하고 블록 관리자는 페이지 테이블을 실행하며 메트릭은 12주 차 오토스케일러에 공급된다. 오른쪽은 2026년의 엔진 선택지다. §9.4의 같은 규칙에 따라 자신의 모델, SLO, 하드웨어에서 후보를 벤치마크해야 한다.}
\end{tbbfigure}

vLLM 주변의 선택지(그림 11-6 오른쪽)를 제9장의 방식으로 간단히 살펴보자. \textbf{TGI}는 Hugging Face의 접착제로, 그 생태계에서 허브부터 엔드포인트까지 가장 짧은 경로다. \textbf{TensorRT-LLM}은 §10.4의 절충을 LLM 규모에 적용한다. GPU마다 엔진을 빌드해 최대 성능의 커널을 얻고 좁은 범위도 함께 받아들인다. Hopper/Ada 실리콘에서 \textbf{FP8}이 제값을 하는 곳이다. 가중치와 KV 캐시에 모두 적용한다. 표 10-1의 "11주 차 FP8은 여기서 동작한다"는 예고가 실현되었다. 플릿의 L4에도 회로가 있으며 같은 §10.3 게이트가 적용된다. \textbf{SGLang}은 호출이 프로그램처럼 보일 때 앞선다. 에이전트, 많이 공유하는 접두사, 스키마에 묶인 출력이 해당한다. \textbf{llama.cpp/GGUF}는 제7장 이후 GPU가 없을 때 쓰는 대안으로 남아 있다. 선택 규칙은 §9.4 이후 변하지 않았고 앞으로도 변하지 않는다. \textit{자신의 모델, SLO, 하드웨어에서 후보를 벤치마크하라. 형태는 옮겨 가지만 수치는 그렇지 않다.} 이 책이 vLLM으로 실습하는 이유는 개방형 기본값이고, 메트릭이 개념을 가르치며, 논문 자체가 이 장의 내용이기 때문이다.

\begin{calloutbox}{tbbpurple}{tbbpurplefill}
{\normalsize\bfseries\color{tbbpurple}생각해 보기}\par\vspace{3pt}
\begin{tbbitemize}
  \item max\_len 2,048 슬랩을 사용하는 연속 엔진이 [30, 2,000] tokens에서 균일한 길이의 답변을 서빙한다. 시퀀스당 기대 내부 낭비와 L4의 13 GB 풀에서 잃는 좌석 수를 16토큰 페이징 블록과 비교하라. 이어 블록 크기가 1이나 512가 아니라 16인 이유와 각 극단의 비용을 설명하라.
  \item 제품 A/B 테스트에서 접두사 캐싱은 평균 TTFT를 38\% 줄였지만 p99 TTFT는 6\%만 줄였다. §11.4의 구간과 레플리카별 적중률 분포로 두 수치를 설명하라. LB 설정 한 줄로 p99를 가장 크게 개선할 라우팅 변경도 제시하라.
  \item 추측 디코딩은 "출력 분포를 보존함이 증명"되었지만 §10.3은 모든 손실 변경에 게이트를 요구했다. 추측은 손실 변환인가? 배포 전에 여전히 검증해야 할 사항을 제시하라. 지연 시간 분포, 토큰당 비용, 수락률 텔레메트리를 고려하고 조용한 실패 유형을 감지할 대시보드 항목도 말하라.
  \item L4 한 대의 Multi-LoRA 입장 정책을 설계하라. 기본 모델은 6 GB, 어댑터는 각각 40 MB이며 KV 풀을 공유한다. 고객 어댑터에는 캐시할 수 있는 300토큰 시스템 접두사가 있다. 어댑터 수가 늘수록 좌석은 어디로 가는가? "기본 모델 하나와 여러 어댑터"가 "기본 모델을 양자화해 좌석을 더 확보"하는 방식보다 유리해지는 지점도 계산으로 개략적으로 나타내라.
\end{tbbitemize}
\end{calloutbox}

\section{실습: 실제 LLM 서빙, 스트리밍, 측정}

이번 실습은 9주 차와 10주 차와 결이 다르다. vLLM은 \code{pip install} 한 번이면 되므로 빌드할 것이 없다. 따라서 \textbf{모든 시간을 측정에 쓴다.} 결과물은 수치로 뒷받침한 프로젝트의 새 SLO 문장이다. 수업용 T4에서 작은 오픈 모델을 서빙한다. Qwen2.5-1.5B-Instruct는 게이트가 없고 FP16에서 3.1 GB라 16 GB 안에서 KV 풀과 함께 여유롭게 들어간다. 직접 눈으로 스트리밍한 뒤 9주 차에 시계 하나를 측정했던 것처럼 두 시계를 측정한다. GPU 시간 예산은 \textasciitilde{}\$2–3이며 언제나처럼 결제 내역 스크린샷을 제출한다.

\begin{tbbfigure}
  \adjustbox{max width=0.94\linewidth,max totalheight=0.46\textheight}{\includegraphics{figures/fig11-7_lab.png}}\\[3pt]
  \figcaption{그림 11-7  실습 파이프라인과 예상되는 결과의 형태. 서빙하고, 하나를 손으로 스트리밍하고, 토큰 인식 부하 테스트를 실행하며, 엔진의 자체 집계를 관찰한다. 마지막에는 표, SLO 문장, 두 가지 통과 의례를 남긴다.}
\end{tbbfigure}

\subsection*{일곱 단계}

\begin{tbbitemize}
  \item \textbf{① 서빙한다.} \code{vllm serve Qwen/\allowbreak Qwen2.5-1.5B-Instruct --max-model-len 4096 --gpu-memory-utilization 0.90}을 실행한다. 시스템 엔지니어처럼 부팅 로그를 읽는다. 가중치 로드(제6장의 범위 GET), CUDA 그래프 캡처(제10장의 a), \textbf{풀이 보유한 KV 블록 수}를 출력하는 줄을 확인한다. 터미널 기록 11-2의 계산을 엔진이 대신 해 준 결과다. 시작에는 1분 이상 걸리므로 제5장의 준비 상태 프로브 예산을 떠올려야 한다.
  \item \textbf{② 하나를 손으로 스트리밍한다.} 터미널 기록 11-5처럼 \code{"stream": true}와 \code{curl -N}을 사용하고 SSE 청크가 도착하는 모습을 \textit{지켜본다}. 잠깐 숨을 고른 뒤 첫 청크(TTFT)가 오고 이어 일정한 가랑비(TPOT)가 내린다. 어떤 하네스(harness)가 추상화하기 전에 두 시계를 직접 눈으로 본다.
  \item \textbf{③ 직접 캐시 크기를 계산한다.} 모델 카드에 맞춰 kv\_math.py를 실행한다. 28개 계층에 KV 헤드 2개 × 128이므로 작은 모델의 캐시는 \textasciitilde{}28 KB/token으로 3B의 4분의 1이다. ①단계에서 엔진이 출력한 블록 수를 예측하고 10\% 안에서 맞춘다. 틀리더라도 실습에서 가장 저렴한 교훈이다. 빠뜨린 항을 찾으면 되며 보통 활성값이다.
  \item \textbf{④ 토큰을 인식하는 방식으로 부하를 건다.} 제공된 비동기 하네스는 c = 1 / 8 / 32 / 64개의 동시 스트림을 연다. 500토큰 프롬프트와 \textasciitilde{}200토큰 답변을 사용하고 길이를 통제하기 위해 \code{ignore\_eos}를 적용한다. 이 사실을 보고서에 명시해야 한다. 요청 지연 시간이 아니라 \textbf{스트림별 TTFT와 TPOT 백분위수}를 기록한다. Locust의 요청 수준 관점은 구조적으로 이 상황을 볼 수 없다. 하네스는 aiohttp 50줄이고, 이를 읽는 것도 실습에 포함된다.
  \item \textbf{⑤ ④가 실행되는 동안 집계를 지켜본다.} 터미널 기록 11-3의 세 구간을 실시간으로 재현한다. \code{running}은 좌석 수까지 올라가고, 풀이 차면 \code{waiting}이 나타나며, \code{gpu\_cache\_usage\_perc}는 높은 값에 고정된다. \textbf{높은 값에 고정되는 것은 정상}이다. 제3장의 페이지 캐시 교훈을 GPU에 적용한 모습이다. TTFT p99가 느끼는 것은 캐시가 아니라 대기열이다.
  \item \textbf{⑥ 통과 의례를 수집한다.} 앞에 60초 유휴 LB나 프록시를 두거나 실습 자료의 nginx 설정을 사용해 \textbf{정확히 60.0 s에 끊기는 스트림}을 재현한다. SSE keep-alive 주석을 추가하고 유휴 타임아웃을 높여 고친다. 제6장의 빚을 직접 갚는 과정이다. 선택 사항이지만 같은 모델의 \textbf{AWQ INT4} 빌드로 스윕 한 행을 반복하고 두 TPOT를 나란히 비교하기를 권장한다. 보고서에서는 이 비교 옆에 §10.3의 게이트 논의를 넣어야 한다.
  \item \textbf{⑦ 보고서와 문장을 작성한다.} 그림 11-7의 표에는 각 c에서 TTFT p50/p99, 스트림별 tok/s, 집계 tok/s를 넣는다. 이어 형태의 원인을 밝히는 문단 하나를 쓴다. c가 증가할수록 스트림별 TPOT는 어디로 갔는가? TTFT의 굴곡점은 어디에 나타났으며 왜 대기열 때문인가? 각 조항을 뒷받침하는 측정 근거와 함께 프로젝트의 SLO 문장을 다시 쓴다. 추정 → 측정 → 최적화 → \textbf{재계약}으로 이어지며, 제안서의 비용 이야기가 이제 스트림 단위로 말한다.
\end{tbbitemize}

\begin{terminalbox}
\begin{Verbatim}[breaklines=true,breakanywhere=true,fontsize=\small,formatcom=\color{tbbtermfg},codes=\tbbcodeactive]
$ vllm serve Qwen/Qwen2.5-1.5B-Instruct --max-model-len 4096 \
      --gpu-memory-utilization 0.90
INFO  Loading weights: 3.09 GiB in 14.2 s (Ch. 6's shelf at work)
INFO  # GPU blocks: 21,856  (= 349,696 tokens of KV pool)   # step ③!
INFO  Capturing CUDA graphs for decode... done (Ch. 10's a, crushed)
INFO  Uvicorn running on http://0.0.0.0:8000

$ curl -N localhost:8000/v1/chat/completions -H "Content-Type: application/json" \
   -d '{"model":"Qwen/Qwen2.5-1.5B-Instruct","stream":true,
        "messages":[{"role":"user","content":"What is a KV cache?"}]}'
data: {"choices":[{"delta":{"content":"The"}}]}          <- ~0.18 s: TTFT
data: {"choices":[{"delta":{"content":" KV"}}]}          <- then a steady
data: {"choices":[{"delta":{"content":" cache"}}]}       <-   ~21 ms drizzle
  ...
data: [DONE]
# you have now SEEN both clocks. everything else is percentiles.
\end{Verbatim}
\end{terminalbox}
\figcaption{터미널 기록 11-5  손으로 서빙하고 스트리밍하기. 부팅 로그는 KV 블록 수를 알려 주므로 직접 계산한 결과와 맞춰야 한다. SSE 스트림은 TTFT와 TPOT를 물리적으로 느끼게 한다. 한 번 숨을 고른 뒤 읽기보다 빠른 가랑비가 이어진다.}

\begin{terminalbox}
\begin{Verbatim}[breaklines=true,breakanywhere=true,fontsize=\small,formatcom=\color{tbbtermfg},codes=\tbbcodeactive]
$ python sweep.py --url :8000 --conc 1 8 32 64 --prompt-tok 500 --gen-tok 200
 c    TTFT p50   TTFT p99   TPOT p50    per-stream   aggregate
  1     0.19 s     0.24 s    21.4 ms      47 tok/s      47 tok/s
  8     0.28 s     0.90 s    26.3 ms      38 tok/s     304 tok/s
 32     0.60 s     2.10 s    34.5 ms      29 tok/s     928 tok/s
 64     1.40 s     4.80 s    55.6 ms      18 tok/s   1,150 tok/s  (plateau)
# read it with the chapter: aggregate x20 for per-stream /1.6 (c=32);
# TTFT p99 dies FIRST (queueing - regime 3), TPOT never falls below
# reading speed; and the plateau is §11.3's bandwidth, not "load".
# now write the SLO sentence your own numbers can sign.
\end{Verbatim}
\end{terminalbox}
\figcaption{터미널 기록 11-6  토큰 인식 스윕. 9주 차 방법을 두 시계 너비로 확장했다. 굴곡점은 TTFT p99의 대기열에, 고원은 집계 tokens/s의 메모리 버스에 있다. SLO 문장은 홍보 자료가 아니라 표에서 나온다.}

\begin{calloutbox}{tbbred}{tbbxFEF2F2}
{\normalsize\bfseries\color{tbbx991B1B}상자 11-2  보고서의 위험 신호 — 스트리밍판 채점자 점검표}\par\vspace{3pt}
\begin{tbbitemize}
  \item \textbf{스트리밍 엔드포인트 주변에 요청 수준 p99 사용} — 자신의 보고서에서 사례 11-A의 실수를 반복한 것이다.
  \item \textbf{TTFT와 TPOT를 함께 평균 냄} 또는 수식어 없이 "지연 시간"이라 표현 — 두 시계와 두 분포가 있으며 하나의 수치로 합치지 않는다.
  \item \textbf{단일 스트림 tokens/s를 용량으로 제시} — 용량은 동시 부하에서 집계한 값이며 스트림별 하한도 명시해야 한다.
  \item \textbf{통제하지 않은 답변 길이를 밝히지 않고 스윕} — ignore\_eos나 길이 분포를 명시하라. 숨겨진 변동 위의 백분위수는 정장을 입은 잡음이다.
  \item \textbf{gpu\_cache\_usage를 경보로 취급} — 높은 값에 고정되는 것은 설계가 작동하는 상태다. 경보는 num\_requests\_waiting에 있다.
  \item \textbf{결제 내역 스크린샷 누락} — 이 통과 의례는 책이 앞으로 가르칠 모든 아키텍처에서도 살아남는다.
\end{tbbitemize}
\end{calloutbox}

\vspace{3.0pt}

제출물은 네 가지 동시성에서 두 시계를 담은 스윕 표, ①과 ③의 블록 수 대조, 두 통과 의례(60.0 s 스크린샷과 수정, 구간 3 메트릭 캡처), 시도했다면 §10.3 게이트가 검사할 내용을 한 문장으로 설명한 AWQ 비교 행, 다시 쓴 SLO 문장, 결제 내역 스크린샷이다. 수치 한 쪽과 글 한 쪽이면 프로젝트 제안서의 서빙 절이 조용히 완성된다.

\begin{calloutbox}{tbbpurple}{tbbpurplefill}
{\normalsize\bfseries\color{tbbpurple}생각해 보기}\par\vspace{3pt}
\begin{tbbitemize}
  \item ③단계의 예측이 엔진이 출력한 블록 수와 22\% 차이 났다. 가능성이 큰 원인 세 가지를 활성값 예비 공간, CUDA 그래프 메모리, 캐시 dtype 순으로 나열하고, 각각을 분리해 확인할 플래그 하나짜리 실험을 제시하라.
  \item 터미널 기록 11-6에서 c=1에서 c=32로 갈 때 TPOT p50은 21.4 → 34.5 ms로 늘었다. §11.3의 최악 예측인 3×보다 작다. 이번 실습의 컨텍스트가 4k가 아니라 500+200 tokens라는 점이 KV 항을 작게 만드는 이유를 설명하라. T4 수치가 L4 계산 예시와 같은 형태가 되려면 컨텍스트 길이가 얼마여야 하는가?
  \item 플릿으로 변환하라. T4는 c=32에서 집계 928 tok/s, TTFT p99 = 2.1 s로 측정되었다. 제품에는 TTFT p99 ≤ 500 ms에서 동시 스트림 100개가 필요하다. L = λW와 스윕을 사용해 레플리카 수를 구하고, 용량과 지연 시간 중 어느 조항이 그 수를 강제했는지 설명하라.
  \item AWQ 행에서 TPOT는 21.4 → 13.1 ms지만 TTFT p50은 0.19 → 0.17 s로 거의 움직이지 않았다. §11.2의 두 장치를 이용해 두 델타를 설명하라. 모델 스토어에서 AWQ 빌드와 함께 두어야 할 §10.3의 산출물 규율 항목도 제시하라.
\end{tbbitemize}
\end{calloutbox}

\namedsection{요약}{열여섯 가지 핵심 내용}

\qitem{1}{\textbf{요청은 작업 단위가 아니게 되었다.} 생성 답변은 사용자의 요구와 모델의 선택에 따라 40–4,000 tokens 범위에 걸친다. 요청 수준 p99는 시스템 상태가 아니라 답변 길이를 추적한다. 사례 11-A에서는 모든 하위 시스템이 녹색인데도 SLO를 20× 위반했다. 죽은 것은 장치가 아니라 메트릭이다.}

\qitem{2}{\textbf{스트리밍은 속도 향상이 아니라 계약 변경이다.} 1968년 텔레타이프를 재발견한 셈이다. Miller의 상수(0.1 s / 1 s / 10 s)와 읽기 속도(\textasciitilde{}5–7 tok/s)가 제품을 정의한다. 첫 토큰을 흐름 임곗값 아래에서 전달하고 이후 독자보다 앞서가야 한다. 같은 14.8 s 연산도 스트리밍하면 즉각적으로 느껴진다.}

\qitem{3}{\textbf{시계 하나를 둘로 바꾼다.} TTFT는 대기열 + 프리필로 반응성을 나타내고, TPOT/tokens-per-second는 생성 속도를 나타낸다. 프로젝트 SLO도 변환 후 살아남는다. "TTFT p99 ≤ 500 ms, 스트림당 ≥ 15 tok/s p50, 동시 스트림 100개, 월 \$900." 백분위수, 임곗값, 부하, 가격이라는 같은 규율에 새 단위를 쓴다.}

\qitem{4}{\textbf{한 번의 생성은 두 장치로 이루어진다.} 프리필은 전체 프롬프트를 한 패스에서 처리하며 AI ≈ 프롬프트 길이인 연산 병목이다. 500 tokens에서 \textasciitilde{}55 ms로 TTFT를 결정한다. 디코드는 토큰마다 전체 패스를 한 번 실행하며 AI ≈ 1이다. §10.2의 하한인 6 GB ÷ 300 GB/s = 20 ms에 묶인 메모리 병목으로 TPOT를 결정하고, 시퀀스 안에서는 본질적으로 직렬이다.}

\qitem{5}{\textbf{L = λW는 종이 바뀌어도 살아남았다.} 스트림 하나가 W ≈ TTFT + N·TPOT 동안 좌석을 차지한다. 200 tokens에서 \textasciitilde{}4.4 s이므로 좌석 32개가 \textasciitilde{}7.3 req/s를 흡수한다. 사례 11-A에서 측정한 6.8을 설명한다. 제2장의 플릿 법칙에 단위만 다시 부여했다.}

\qitem{6}{\textbf{KV 캐시는 과거 재계산을 피하기 위해 교환한 메모리다.} 2 × layers × kv\_heads × head\_dim × bytes로 계산하면 책의 3B (GQA)에서 112 KB/token, 4k 대화마다 448 MiB다. 사례 11-B의 실패한 할당량을 도출했다. 대화가 살아 있는 동안 토큰마다 행 하나씩 자란다.}

\qitem{7}{\textbf{동시성은 메모리 수치가 되었다.} L4 budget = 21.6 usable − 6 weights − 2.5 runtime = 13 GB pool → \textasciitilde{}115k tokens → 4k 좌석 28개다. 이 계산 어디에도 FLOP은 나오지 않는다.}

\qitem{8}{\textbf{캐시는 점유 공간일 뿐 아니라 트래픽이다.} step bytes = weights + Σ KV(context)다. batch 32 × 4k에서는 21 GB → 70 ms → 스트림당 14 tok/s, 집계 \textasciitilde{}460으로 스트림당 비용 3×를 대가로 처리량 10×를 얻는다. 이런 좌석이 \textasciitilde{}13개를 넘으면 캐시가 모델보다 무겁다. 긴 컨텍스트는 모두에게 세금을 부과하며 GQA/MQA, 게이트를 적용한 KV-FP8, 슬라이딩 윈도가 경량화 수단이다.}

\qitem{9}{\textbf{정적 배치는 제각각인 출력 때문에 실패한다.} 가장 긴 글이 끝날 때까지 기다려 좌석-단계의 \textasciitilde{}50\%가 죽고, TTFT는 낯선 사람에게 붙잡힌다. 상자 9-1의 작은 글씨가 폭발한 결과다.}

\qitem{10}{\textbf{연속 배치(Orca)는 요청이 아니라 반복을 스케줄링한다.} \textasciitilde{}20 ms 단계마다 끝난 시퀀스를 퇴거하고 KV가 허용하는 만큼 대기 시퀀스를 입장시킨 뒤 모두를 진행한다. 배치는 집단이 되며 누군가 기다리는 동안 좌석이 쉬지 않는다. 다음 출발까지 한 단계뿐이므로 max\_queue\_delay는 사라진다.}

\qitem{11}{\textbf{두 시계는 엔진 하나 안에서 싸운다.} 전체 프롬프트 프리필은 디코더 30개를 멈춰 TPOT를 끊기게 한다. 청크 프리필은 이를 단계에 나누며 max\_num\_batched\_tokens가 TTFT↔TPOT 다이얼이다. 부하 실패는 언제나 대기열부터 시작한다. 터미널 기록 11-3의 세 구간처럼 TPOT는 완만하게 나빠지는 동안 TTFT p99가 무너진다.}

\qitem{12}{\textbf{PagedAttention은 캐시용 가상 메모리다.} 연속 최대 길이 슬랩은 풀의 60–80\%를 낭비했다(SOSP 2023 측정). 16토큰 블록 + 시퀀스별 블록 테이블은 낭비를 4\% 아래로 줄였다. 회수한 낭비가 좌석과 처리량으로 누적되어 당시 가장 강한 시스템보다 2–4×, 단순 서빙보다 14–24× 높은 성능을 냈다. 한 세대 최고의 시스템 논문이 메모리 할당기에 관한 것이었다.}

\qitem{13}{\textbf{블록 테이블의 두 번째 묘기는 공유다.} 900토큰 시스템 프롬프트를 32× (3.2 GB) 저장하지 않고 참조 횟수 + 쓰기 시 복사로 한 번만 저장한다. 접두사 캐싱은 다중 턴 재프리필과 TTFT \textasciitilde{}100 ms를 돌려준다. 고정 세션 문제도 해결된다. 정확성은 상태 비저장으로 유지하고 어피니티는 캐시 힌트로만 쓴다. 캐시를 계약으로 만들지 않는다.}

\qitem{14}{\textbf{vLLM은 운영체제가 들어간 제9장의 구조다.} 배처용 반복 스케줄러, 메모리용 블록 관리자, CUDA 그래프 디코드(제10장의 a), AWQ/GPTQ INT4 가중치(두 시계 + 좌석 9개, 게이트를 거친 산출물 하나), OpenAI 호환 SSE 프런트 도어, 12주 차 오토스케일러가 사용할 메트릭(running / waiting / cache \%)을 갖춘다.}

\qitem{15}{\textbf{추측 디코딩은 메모리 벽의 유휴 FLOP에서 차익을 얻는다.} 작은 초안 모델이 k tokens를 추측하고 목표 모델이 토큰 하나와 같은 바이트의 한 가중치 스트림으로 검증한다. 네 개 중 \textasciitilde{}3개가 살아남아 배치 1 TPOT가 \textasciitilde{}2.5× 개선되고 분포는 증명 가능하게 동일하다. 다수 사용자용 연속 배치를 보완하는 단일 사용자 수단이며, 수락률이 떨어지면 속도 향상이 사라지는 부드러운 실패를 텔레메트리로 감시해야 한다.}

\qitem{16}{\textbf{실습의 결론은 한 줄로 정리된다.} \$0.53/hr T4에서 c = 32일 때 집계 처리량은 ×20, 스트림별 처리량은 ÷1.6이고 TTFT p99가 언제나 먼저 깨지지만 모든 스트림은 독자보다 빠르다. 프로젝트 제안서는 이제 스트림 단위로 추정 → 측정 → 최적화 → 재계약을 말한다.}

\namedsection{연습문제}{}

\subsection*{A. 객관식}

\qitem{1}{한 달 동안 채팅 서비스의 요청 수준 p99 지연 시간이 세 배가 되었지만 TTFT p99와 TPOT p50은 그대로다. 가장 가능성 높은 원인은 무엇인가?}

\begin{tbboptions}
  \item ① 메모리 누수        ② 사용자가 더 긴 답변을 요구한다. E2E는 상태가 아니라 길이를 추적한다
  \item ③ GPU가 스로틀링된다        ④ 로드 밸런서 설정이 잘못되었다
\end{tbboptions}

\qitem{2}{프리필은 보통 연산 병목이고 디코드는 메모리 병목인 이유는 무엇인가?}

\begin{tbboptions}
  \item ① 프리필은 CPU에서 실행된다        ② 프리필은 모든 프롬프트 토큰을 한 패스에서 처리하고(AI ≈ 토큰 수), 디코드는 토큰 하나를 위해 모든 가중치를 스트리밍한다(AI ≈ 1)
  \item ③ 디코드는 더 큰 행렬을 사용한다        ④ 프리필은 어텐션을 건너뛴다
\end{tbboptions}

\qitem{3}{책의 3B 모델(28개 계층, KV 헤드 8개 × 128, FP16)에서 4,096토큰 대화 하나의 KV 캐시 크기에 가장 가까운 것은 무엇인가?}

\begin{tbboptions}
  \item ① 45 MB        ② 448 MiB        ③ 4.5 GB        ④ 배치 크기에 따라 다르다
\end{tbboptions}

\qitem{4}{연속 배치가 정적 배치보다 나은 주된 이유는 무엇인가?}

\begin{tbboptions}
  \item ① 시퀀스를 같은 길이로 패딩한다        ② 배치 크기를 키운다
  \item ③ 디코드 단계마다 배치 구성원을 바꾼다. 끝난 시퀀스가 나가고 기다리던 시퀀스가 들어온다        ④ 더 빠른 GPU에서 프리필을 실행한다
\end{tbboptions}

\qitem{5}{지속적인 과부하에서 연속 배치 엔진이 거의 언제나 가장 먼저 위반하는 SLO 조항은 무엇인가?}

\begin{tbboptions}
  \item ① TPOT(토큰이 느려진다)        ② TTFT p99(입장 대기열 — 구간 3)
  \item ③ 집계 tokens/s        ④ 정확도
\end{tbboptions}

\qitem{6}{PagedAttention의 핵심 메커니즘은 무엇인가?}

\begin{tbboptions}
  \item ① KV 캐시 압축        ② 고정 크기 KV 블록 + 논리 위치를 물리 위치에 대응시키는 시퀀스별 블록 테이블. 어텐션용 운영체제 페이지 테이블이다
  \item ③ 캐시를 호스트 RAM으로 이동        ④ K/V를 저장하지 않고 다시 계산
\end{tbboptions}

\qitem{7}{접두사 캐싱에서 LLM 채팅의 세션 어피니티(고정 라우팅)는 무엇이 되는가?}

\begin{tbboptions}
  \item ① 정확성에 필수다        ② 성능 힌트다. 미스는 프리필 비용만 다시 낼 뿐 잘못된 답을 만들지 않으므로 롤아웃과 실패에서 점진적으로 성능만 저하된다
  \item ③ 불가능하다        ④ GGUF 모델에만 관련 있다
\end{tbboptions}

\qitem{8}{추측 디코딩은 어떻게 속도를 높이는가?}

\begin{tbboptions}
  \item ① 어려운 토큰을 건너뛴다        ② 정밀도를 낮춘다
  \item ③ 토큰 하나를 만들었을 동일한 가중치 스트림인 목표 모델 패스 한 번에서 초안 토큰 k개를 검증한다. 유휴 FLOP을 쓰며 분포는 변하지 않는다
  \item ④ 답변 전체를 캐시한다
\end{tbboptions}

\subsection*{B. 단답형}

\qitem{1}{TTFT와 TPOT를 기계적으로 정의하고 각 수치를 결정하는 단계를 밝힌 뒤, 예산의 근거가 되는 인간 상수 두 개를 제시하라.}

\qitem{2}{KV 헤드 8개 × 128차원을 갖는 40계층 FP8 모델의 KV bytes/token과 8k 대화 하나의 전체 크기를 계산하라.}

\qitem{3}{한 시퀀스 안의 디코드 단계를 병렬화할 수 없는 이유는 무엇인가? 부분적으로 빠져나오는 기법 하나와 이를 가능하게 하는 비대칭도 제시하라.}

\qitem{4}{반복 수준 스케줄링(iteration-level scheduling)의 세 박자 주기를 서술하고, 9주 차의 max\_queue\_delay를 무엇이 대체했는지 설명하라.}

\qitem{5}{청크 프리필은 누구를 위해 어떤 시계와 어떤 시계를 절충하는가? 이 절충을 표현하는 조절점을 한 문장으로 설명하라.}

\qitem{6}{엔진이 gpu\_cache\_usage\_perc = 97\%, num\_requests\_waiting = 0을 보고한다. 정상인가, 경보 상황인가? 제3장의 비유를 이용해 두 문장으로 정당화하라.}

\subsection*{C. 서술형}

\qitem{1}{두 장치에 관한 글을 작성하라. 한 스타트업이 A100-80GB (2,000 GB/s)에서 7B GQA 모델(32개 계층, KV 헤드 8개 × 128, FP16 가중치 14 GB)을 서빙한다. 디코드 하한, KV/token, 가중치 + 3 GB 런타임을 뺀 KV 풀, 8k 컨텍스트의 좌석 수, 8k에서 배치 32의 단계 비용, 캐시 트래픽이 가중치 트래픽을 넘어서는 교차점을 계산하라. 이어 각 조항의 측정 근거를 밝힌 SLO 문장을 작성하라.}

\qitem{2}{"PagedAttention은 이 책에서 가장 단순하면서 가장 가치 있는 아이디어다. 어떤 커널보다 가치 있는 장부 관리다"라는 주장을 옹호하거나 공격하라. 낭비 측정값, 좌석→승객→처리량으로 누적되는 효과, 공유/CoW의 이점, 페이징의 간접 참조가 비용을 만드는 지점 하나 이상을 활용하라. 커널 복잡성, 블록 크기 절충, TLB와 비슷한 조회를 고려할 수 있다.}

\qitem{3}{간섭에 관한 글을 작성하라. 서비스는 한 vLLM 플릿에서 30토큰 분류형 호출과 2,000토큰 문서 채팅을 섞는다. 프리필–디코드 간섭, 거대한 KV, 입장 정책을 이용해 배포를 설계하라. 할당량이 있는 단일 풀, 두 풀, 분리형 프리필 등을 고려하고 트래픽 클래스별 SLO와 설계 실패를 알려 줄 메트릭을 제시하라.}

\qitem{4}{가격 재계산에 관한 글을 작성하라. 제10장의 프로젝트 4×L4 플릿(\$2,381/mo, 요청 형태)을 가져온다. 가정한 답변 길이 혼합에서 L = λW로 스트림 용량을 다시 도출하고, AWQ INT4(두 시계 + 좌석)와 다섯째 레플리카 중 어느 쪽이 30\% TTFT-p99 격차를 더 잘 줄이는지 결정하라. 플릿 비용을 제안서에 반영하고, 의도적으로 남겨 둔 12주 차 수단인 대기열 깊이 기반 오토스케일링을 표시하라.}

\namedsection{정답과 해설}{}

\subsection*{A. 객관식}

\qitem{1}{\textbf{②} E2E ≈ TTFT + N·TPOT다. 두 시계가 그대로라면 움직인 것은 N뿐이다. 사례 11-A를 진단 패턴으로 적용해 인프라보다 먼저 길이 분포를 확인해야 한다. (§11.1–11.2)}

\qitem{2}{\textbf{②} 산술 집약도 관점에서 프리필은 스트리밍한 각 가중치를 모든 프롬프트 토큰에 재사용하고(AI ≈ 500 > 능선), 디코드의 행렬–벡터 패스는 각 가중치를 한 번만 쓴다(AI ≈ 1 « 능선). 같은 실리콘이 §10.2의 벽 반대편에 놓인다. (§11.2)}

\qitem{3}{\textbf{②} 2 × 28 × 8 × 128 × 2 B = 112 KB/token이고 × 4,096 = 448 MiB다. 사례 11-B에서 정확히 이 양을 할당하다 실패했다. ④는 함정이다. 시퀀스별 KV는 배치와 무관하다. (§11.3)}

\qitem{4}{\textbf{③} 반복 수준 스케줄링은 \textasciitilde{}20 ms 토큰 경계마다 퇴거/입장/진행한다. ①은 입력을 다룬 제10장의 패딩 이야기이고 여기서는 제각각인 출력이 문제다. (§11.4)}

\qitem{5}{\textbf{②} 풀의 좌석이 입장을 제한한다. 과부하는 num\_requests\_waiting에 쌓이고 TTFT가 대기열을 물려받는다(L = λW). TPOT는 배치와 함께 완만하게 나빠지지만 대기열은 폭발한다. 터미널 기록 11-3의 구간 3이다. (§11.4)}

\qitem{6}{\textbf{②} 고정 블록 + 변환 테이블이 페이징이며 단편화는 마지막 블록 하나로 줄어든다(60–80\%에서 <4\%). 압축도 오프로딩도 아닌 간접 참조다. (§11.5)}

\qitem{7}{\textbf{②} 요청이 전체 기록을 담으므로 어떤 레플리카도 정확히 답한다. 미스는 프리필 한 번(\textasciitilde{}350 ms)의 비용을 내고 적중은 이를 아낀다. 성능 힌트일 뿐 정확성 요구사항은 아니다. 제6장의 긴장을 해결했다. (§11.5)}

\qitem{8}{\textbf{③} k tokens를 검증하는 데 토큰 하나를 생성할 때와 같은 6 GB를 스트리밍한다. 메모리 벽의 유휴 FLOP이 검사를 수행하고 거부 표본 추출이 출력 분포를 정확히 유지한다. (§11.5)}

\subsection*{B. 단답형}

\qitem{1}{TTFT = 대기열 + 토큰화 + \textbf{프리필}이며, 연산 병목으로 반응성을 결정한다. TPOT = \textbf{디코드} 단계당 시간이며, 대역폭 병목으로 생성 속도를 결정한다. TTFT 예산은 Miller의 \textasciitilde{}1 s 흐름 임곗값, TPOT 예산은 여유를 둔 읽기 속도 \textasciitilde{}5–7 tok/s를 바탕으로 한다. (§11.1–11.2)}

\qitem{2}{2 × 40 × 8 × 128 × 1 B = 81,920 B ≈ \textbf{80 KB/token}이고 × 8,192 ≈ 대화당 \textbf{671 MB}다. FP8은 FP16의 바이트를 절반으로 줄이지만 게이트는 여전히 적용한다. (§11.3)}

\qitem{3}{token t+1의 \textbf{입력이} token t이므로 어떤 스케줄러도 없앨 수 없는 데이터 의존성이 있다. 추측 디코딩은 저렴한 초안이 제안하고 목표 모델이 한 단계 비용으로 k개를 검증해 부분적으로 빠져나온다. 디코드가 메모리 병목이라 FLOP이 유휴 상태이기 때문에 가능하다. (§11.2, §11.5)}

\qitem{4}{끝난 시퀀스 퇴거 → KV가 허용하는 만큼 대기 시퀀스 입장 → 모두 진행을 토큰마다 반복한다. 출발이 매 단계 있으므로 max\_queue\_delay가 사라지고 단계 예산 max\_num\_batched\_tokens, 즉 TTFT↔TPOT 다이얼이 대체했다. (§11.4)}

\qitem{5}{새 승객의 TTFT를 보호하기 위해 각 단계에 청크를 태워 기존 승객의 TPOT에 작은 비용을 부과한다. 동시에 전체 프롬프트 프리필이 만들 반대 방향의 끊김을 제한한다. 조절점은 반복당 프리필 + 디코드 토큰 수를 합한 max\_num\_batched\_tokens다. (§11.4)}

\qitem{6}{정상이다. 제3장의 페이지 캐시처럼 가득 찬 KV 풀은 회수 가능한 좌석이 제 역할을 하며 자원이 \textbf{동작}하는 상태다. 압박은 대기열(waiting > 0)로 나타나므로 이 항목이 경보다. 97\%이면서 대기열과 선점이 증가하면 경보 상황이다. (§11.4–11.6)}

\subsection*{C. 서술형(채점 기준)}

\qitem{1}{하한은 14 ÷ 2,000 = \textbf{7 ms/token}이며 상한은 143 tok/s다. KV/token은 2×32×8×128×2 = \textbf{128 KB}다. 풀은 80×0.9 − 14 − 3 = \textbf{55 GB} → \textasciitilde{}420k tokens → \textbf{8k 좌석 51개}다. 배치 32 단계는 14 + 32×1.07 GB ≈ 48 GB → 24 ms → 스트림당 41 tok/s, 집계 1,330이다. 교차점은 14 ÷ 1.07 ≈ \textbf{좌석 13개}다. 좋은 SLO는 TTFT를 프리필 + 대기열 측정에 연결하고, 풀이 실제로 보유한 좌석 수에서 TPOT ≥ 15–20 tok/s를 약속한다. (§§11.2–11.4)}

\qitem{2}{좋은 글은 누적 효과(낭비 → 좌석 → 승객 → 2–4×/14–24×)를 정량화하고 공유의 이점(시스템 프롬프트 중복 제거, CoW, 접두사 TTFT 환급)을 인정한 뒤 비용을 정직하게 계산한다. 어텐션 커널은 테이블을 통해 모아야 해 융합이 더 어렵고 §10.4와 긴장이 생긴다. 블록 크기는 큰 블록의 내부 낭비와 작은 블록의 테이블 오버헤드를 절충한다. "장부 관리일 뿐"이어도 300 GB/s에서 실행하려면 커널과 공동 설계해야 했다. 누적 효과와 비용을 모두 계산하면 어느 결론이든 점수를 얻는다. (§11.5)}

\qitem{3}{포함할 내용: 분류 호출의 TTFT가 문서 프리필 뒤에서 고통받는 간섭에는 먼저 청크 프리필을 적용한다. 거대한 KV에는 할당량이나 별도 풀을 둔다. 규모를 키우면 분리형 프리필을 쓰되 KV 전송 네트워크 비용을 명시한다. 클래스별 SLO를 둔다. 분류는 요청 형태 SLO도 괜찮다. 실패 메트릭은 거대한 요청이 부풀리는 집계 tok/s가 아니라 클래스별 TTFT p99와 선점 횟수다. (§§11.3–11.5, 상자 11-1)}

\qitem{4}{포함할 내용: 한 점이 아니라 답변 길이 혼합에서 정직한 W를 구하고, 가정한 컨텍스트에서 KV 풀의 좌석을 계산하며, L = λW를 동시 스트림 100개를 위한 레플리카 수로 바꾼다. AWQ와 레플리카를 두 시계 모두에서 비교하고 §10.3 게이트를 언급한다. INT4는 디코드 하한을 절반으로 줄이고 풀을 넓힌다. 좌석이 늘어 대기열을 비우므로 대기열 때문에 생긴 TTFT-p99 격차에서 다섯째 레플리카보다 나은 경우가 많다. 12주 차에는 이 장에서 읽는 법을 배운 num\_requests\_waiting을 기준으로 오토스케일링한다. (§§11.3–11.6)}

\namedsection{더 읽을거리}{}

\textbf{[이번 주]}로 표시한 자료는 실습을 직접 뒷받침하고, \textbf{[N주 차]}는 나중에 도움이 될 자료를 뜻한다.

\begin{tbbitemize}
  \item \textbf{[이번 주] W. Kwon et al., "Efficient Memory Management for Large Language Model Serving with PagedAttention," SOSP 2023} — 제1장에서 필독으로 지정했다. 오늘 읽으면 벤치마크가 포함된 이 장 자체처럼 보인다. 실습 전에는 §§1–4를, 실습 뒤에는 평가 절을 읽어 보자.
  \item \textbf{[이번 주] G.-I. Yu et al., "Orca: A Distributed Serving System for Transformer-Based Generative Models," OSDI 2022} — 문제에 이름을 붙인 팀이 반복 수준 스케줄링과 선택적 배치를 설명한다. 제9장에서 "KV 캐시 강의 뒤"까지 미뤄 두었으니 이제 읽어도 된다.
  \item \textbf{R. B. Miller, "Response Time in Man-Computer Conversational Transactions," AFIPS 1968} — 지금도 TTFT 예산을 지배하는 텔레타이프 시대의 17쪽짜리 논문이다. 즐겁고 약간 겸손해지는 글이다.
  \item \textbf{N. Shazeer, "Fast Transformer Decoding: One Write-Head Is All You Need" (2019) · J. Ainslie et al., "GQA" (2023)} — §11.3의 바이트를 아키텍처 측면에서 공격한다. MQA의 질문과 GQA의 절충을 통해 모델 카드에 "KV 헤드 8개"라고 적힌 이유를 설명한다.
  \item \textbf{Y. Leviathan et al., "Fast Inference from Transformers via Speculative Decoding" (2023) · C. Chen et al., "Accelerating LLM Decoding with Speculative Sampling" (2023)} — 초안 후 검증 논문으로, 거부 표본 추출을 적용해 분포가 유지된다는 증명도 포함한다.
  \item \textbf{[12주 차] Y. Zhong et al., "DistServe" (OSDI 2024)} — 프리필–디코드 분리를 플릿 설계로 발전시킨다. 상자 11-1의 첫 항목을 GPU당 유효 처리량 계산으로 확장한 연구다.
  \item \textbf{vLLM docs — "Engine Arguments" and "Production Metrics"} — 실습의 조절점(max\_num\_batched\_tokens, gpu\_memory\_utilization)과 터미널 기록 11-3에서 읽은 모든 /metrics 항목을 공식적으로 설명한다.
  \item \textbf{[이번 주, 재미] 좋아하는 채팅 UI의 브라우저 네트워크 탭} — SSE 프레임이 도착하는 모습을 보고 첫 프레임을 초시계로 재어 보자. 바깥에서 누군가의 TTFT SLO를 감사하는 셈이며, 이번 주가 끝나면 그들의 대시보드가 무엇을 말하는지도 정확히 알 수 있다.
\end{tbbitemize}

\vspace{5.0pt}

\begin{calloutbox}{tbbblue}{tbbbluefill}
{\normalsize\bfseries\color{tbbblue}다음 장 — 제12장: 스케일링과 비용 최적화}\par\vspace{3pt}
7주 차 이후 모든 장은 레플리카 하나를 개선했지만, 사용자가 대화하든 잠들어 있든 플릿 자체는 L4 네 대를 24/7 유지했다. 다음 주에는 플릿이 숨 쉬는 법을 배운다. 이번 장에서 읽는 법을 배운 메트릭, 즉 CPU\%가 아니라 대기열 깊이와 KV 압박으로 \textbf{HPA와 KEDA 오토스케일링}을 구동한다. \textbf{0으로 스케일링}과 그 과정에서 부딪히는 GPU 콜드 스타트 문제도 다룬다. 2주 차의 Firecracker와 제6장에서 임계 경로의 6 GB 가중치에 관해 배운 모든 내용이 돌아온다. \textbf{GPU 공유}에서는 MIG와 시간 분할로 제2장의 카드당 테넌트 하나 규칙을 마침내 굽힌다. \textbf{KServe}로 이 책의 서버를 플랫폼처럼 실행하고, 스팟 인스턴스, 약정, \$/M-token 원장을 포함한 \textbf{FinOps 도구 상자}로 제안서 비용표를 관리 가능한 수치로 바꾼다. 프로젝트 비용에는 큰 수단 하나가 남았다. 12주 차에 이를 사용한다.
\end{calloutbox}

%% Chapter 12 — 스케일링과 비용 최적화
\setcounter{chapter}{11}
\chapter{스케일링과 비용 최적화}
\chaptersubtitle{오토스케일링, 콜드 스타트, 공유 실리콘, FinOps 포트폴리오 — 숨 쉬는 법을 배운 플릿과 마침내 도착한 청구서}

\begin{calloutbox}{tbbblue}{tbbbluefill}
{\normalsize\bfseries\color{tbbblue}학습 목표}\par\vspace{3pt}
\begin{tbbitemize}
  \item 플릿의 하루를 \textbf{부하율 문제}(필요 레플리카 시간 ÷ 결제한 레플리카 시간)로 읽고 유휴 격차의 가격을 계산한다. 전력 회사가 한 세기 전에 서빙 경제성을 개념적으로 해결한 이유도 설명한다.
  \item 아무도 쓰지 않는 CPU\%가 아니라 사용자가 느끼는 대기열이라는 \textbf{올바른 지표}에 오토스케일러를 연결하고, 목표 레플리카 공식, 안정화 창, 예약 스케일링, 0으로 스케일링을 포함한 HPA/KEDA 제어 루프를 운영한다.
  \item \textbf{GPU 콜드 스타트}를 스케줄링, 이미지, 가중치, 엔진, 워밍업으로 세분화하고 노드까지 시작해야 하면 ×3이 되는 이유를 이해한다. Firecracker의 125 ms가 청구서에서 차지하는 위치를 정확히 파악하고 완화 사다리를 한 단계씩 오른다.
  \item 격리 요구에 따라 \textbf{시분할, MPS, MIG} 중 하나를 선택하고, 큰 카드의 일부가 작은 카드 전체보다 유리해지는 시점을 계산한다.
  \item 하나의 CRD로 제5–11장의 운영을 수행하는 플랫폼 계층인 \textbf{KServe}의 위치를 파악하고, 원시 Deployment + KEDA가 더 현명한 경우를 판단한다.
  \item \textbf{가격 포트폴리오}(약정 바닥, 온디맨드 어깨, 스팟 버스트, 비혼잡 시간 채우기)를 구성하고 영수증을 근거로 프로젝트 비용을 \$2,381에서 \$900 상한 아래로 낮춘다. 측정한 \$/M-token 수치로 제1장의 자체 구축 대 외부 도입 손익분기점을 마무리해 중간 점검 패키지를 완성한다.
\end{tbbitemize}
\end{calloutbox}

\vspace{2.5pt}

\begin{calloutbox}{tbbteal}{tbbtealfill}
{\normalsize\bfseries\color{tbbteal}핵심 용어}\par\vspace{3pt}
부하율 · 일중 곡선 · HPA / KEDA · ScaledObject · 목표 레플리카 공식 · 안정화 창 · 예약 스케일링 · 0으로 스케일링 · GPU 콜드 스타트 · 사전 워밍업 풀 · 클러스터 오토스케일러 · 시분할 · MPS · MIG · 시끄러운 이웃 · KServe / InferenceService · 약정 사용 · 온디맨드 · 스팟/선점형 · 드레인 훈련 · 단위 경제성(\$/M tokens) · 비용 공개
\end{calloutbox}

\vspace{2.5pt}

\begin{calloutbox}{tbblightgray}{tbbboxfill}
{\normalsize\bfseries\color{tbbink}장 구성}\par\vspace{3pt}
\textbf{12.1} 잠들지 않는 청구서 — 1902년에 이를 해결한 사람    \textbf{12.2} 오토스케일링: 올바른 지표와 초시계    \textbf{12.3} 콜드 스타트와 0으로 스케일링    \textbf{12.4} 실리콘 공유: 시분할, MPS, MIG    \textbf{12.5} KServe와 FinOps 포트폴리오    \textbf{12.6} 실습: 프로젝트 중간 점검 — 이어서 요약, 연습문제, 더 읽을거리
\end{calloutbox}

\section{잠들지 않는 청구서 — 1902년에 이를 해결한 사람}

7주 차 이후 모든 기술 장은 레플리카 하나를 개선했다. 9주 차에는 정직하게, 10주 차에는 토큰당 비용을 낮게, 11주 차에는 스트리밍하면서 좌석이 가득 차게 만들었다. 하지만 제1장 사례 B가 유령 GPU 한 대에 던진 무례한 질문을 플릿에는 아직 묻지 않았다. \textbf{새벽 3시에 정확히 누가 사용하고 있는가?} 프로젝트의 L4 레플리카 네 개는 \code{replicas: 4}가 의미하는 대로 24시간 실행된다. 제5장의 컨트롤러는 점심시간의 혼잡에도 한밤중의 정적에도 완벽하고 비싼 성실함으로 약속을 지킨다. 이 장은 그 성실함의 비용을 측정하는 데서 시작하며, 사례 12-A가 그 결과다.

\begin{terminalbox}
\begin{Verbatim}[breaklines=true,breakanywhere=true,fontsize=\small,formatcom=\color{tbbtermfg},codes=\tbbcodeactive]
$ python fleet_ledger.py --day yesterday       # measured, not modeled
 hour    streams   replicas needed   running   note
 03:00       8          2*              4      # *floor: rollout safety
 09:00      71          3               4      # morning climb
 13:00     104          4               4      # the only honest hour
 17:00      55          3               4
 22:00      19          2               4

 needed today:  4x4h + 3x5h + 2x15h = 61 replica-hours
 paid today:    4x24h               = 96 replica-hours
 load factor:   61/96 = 63.5%
 monthly:  paid $2,381  =  needed $1,513  +  IDLE $868
# the fleet is sized for 13:00 and billed for 03:00.
# Exhibit 1-B's ghost GPU, diurnal edition: $868/mo of night shift.
\end{Verbatim}
\end{terminalbox}
\figcaption{사례 12-A  측정한 하루. 플릿은 가장 바쁜 한 시간을 기준으로 프로비저닝되지만 모든 시간에 비용을 낸다. 하루 35레플리카 시간은 아무도 서빙하지 않는다. 필요한 시간을 결제한 시간으로 나눈 부하율이 이 장에서 높일 단 하나의 수치다.}

문제의 종류에 주목하자. \textbf{잘못 설정된 것은 아무것도 없다.} 모든 레플리카는 최적화되고 게이트를 통과했으며 정직하다. 낭비는 구조적이다. 수요는 시간에 따라 \textit{모양}을 지닌다. 제1장은 서빙 트래픽을 "사용자가 결정"한다고 했고 사용자는 잠을 잔다. 반면 플릿은 수평선이다. 제7장의 생각해 볼 문제에서는 이 격차를 종이에 추정하게 하며 "12주 차가 무엇을 갚아야 하는가?"라고 물었다. 사례 12-A는 그 추정치를 실제로 측정한 결과다. 해결책은 레플리카별 수단일 수 없다. 제10장은 요청을 저렴하게 만들고 제11장은 좌석을 더 채웠지만, \textit{프로비저닝된 채 유휴인} 레플리카는 효율과 관계없이 시간당 \$0.80를 태운다. 남은 수단은 \textbf{개수}를 움직인다.

\begin{tbbfigure}
  \adjustbox{max width=0.94\linewidth,max totalheight=0.46\textheight}{\includegraphics{figures/fig12-1_diurnal.png}}\\[3pt]
  \figcaption{그림 12-1  하루를 그린 모습. 계단은 롤아웃 안전을 위한 바닥 2개를 포함해 SLO가 실제로 요구하는 양이고, 평평한 붉은 선은 정적 플릿이 지불하는 양이다. 그 차이는 매일 야간 근무 35레플리카 시간, 월 \$868다.}
\end{tbbfigure}

\subsection*{1902년: 똑같은 청구서에 직면했던 마지막 산업}

위안이 되는 소식이 있다. 새로운 문제가 아니다. GPU 시간이 아니라 킬로와트시를 팔던 산업이 조부모 세대가 태어나기도 전에 해결했다. 1900년 무렵 전력 회사는 정확히 사례 12-A의 계산 때문에 파산하고 있었다. 발전소는 막대한 고정 비용이 드는데 \textit{저녁 조명 피크에 맞춰 규모를 정하고} 나머지 20시간에는 충분히 쓰이지 않았다. Thomas Edison의 전 비서였고 당시 시카고 전력 회사를 운영하던 \textbf{Samuel Insull}은 수익이 아닌 한 지표에 집착했다. 평균 수요를 최대 수요로 나눈 \textbf{부하율(load factor)}로 사례 12-A 아래쪽과 정확히 같은 비율이다. 전력 회사의 진짜 제품은 \textit{용량}이고 유휴 용량은 순수한 손실이라는 것이 그의 통찰이었다. 골짜기를 채우고 피크를 깎는 게임이다. 그의 회사는 전체 전략을 개척했다. \textbf{소비량뿐 아니라 수요를 계량}했다. 고객의 피크가 발전소 크기를 정하므로 청구서도 피크가 정해야 했으며, 클라우드 가격의 조상이다. \textbf{피크가 서로 맞물리는 고객을 유치}했다. 낮에는 전차, 저녁에는 조명, 밤에는 제빙소와 공장을 운영해 발전소 하나가 세 수요 곡선을 처리했다. 작업을 골짜기로 옮길 수 있는 사람에게 \textbf{비혼잡 시간 전력을 저렴하게 판매}했고, 피크에 공급을 끊어도 되는 산업 고객에게 \textbf{중단 가능 요금}을 제공했다. 전기는 20년 안에 킬로와트시당 약 20센트였던 사치품에서 몇 센트 수준으로 내려왔다. 주된 이유는 더 좋은 발전기가 아니라 더 나은 \textit{부하율}이었다. 이 장의 모든 메커니즘인 오토스케일링, 0으로 스케일링, GPU 공유, 스팟 인스턴스, 야간 배치 기법은 Kubernetes 티셔츠를 입은 Insull의 전략이다. §12.5에서는 당시 고객이 전력을 샀던 방식 그대로 용량을 구매한다.

\subsection*{뻔한 해결책이 두 매복을 만나다}

"수요에 맞춰 레플리카 수만 바꾸면 된다"는 해결책은 기계적으로 모욕적일 만큼 쉽다. 제5장도 이를 약속했다. 오토스케일러는 정수 하나를 수정하는 \textit{원하는 상태의 또 다른 작성자}이고 나머지는 컨트롤러가 모두 처리한다. 흥미로운 엔지니어링은 쉬운 버전이 실패하는 두 방식에 있다. 사례 12-B가 둘을 모두 보여 준다.

\begin{terminalbox}
\begin{Verbatim}[breaklines=true,breakanywhere=true,fontsize=\small,formatcom=\color{tbbtermfg},codes=\tbbcodeactive]
# attempt 1, monday: the tutorial default.
$ kubectl autoscale deploy/chat --cpu-percent=70 --min=2 --max=6
$ kubectl get hpa chat -w              # tuesday, 09:00 wave arrives
NAME   REFERENCE     TARGETS       MINPODS  MAXPODS  REPLICAS
chat   deploy/chat   cpu: 9%/70%      2        6        2
# GPUs drowning (queue growing, TTFT p99 6.1s). CPU: napping at 9%.
# verdict: "9% < 70%, all quiet." it will NEVER scale. wrong eye.

# attempt 2, wednesday: right metric (KEDA on the vLLM queue).
09:00:15  waiting=41 -> desired 2->4       # right eye reacts in 15 s
09:00:24  2 pods Pending -> Init -> ...
09:02:54  ...Ready                         # cold start: 2.5 minutes
# 09:00-09:03: TTFT p99 red. right metric, LATE muscle.
# lesson 1: scale on what users feel. lesson 2: reaction time
# is part of capacity. the wave does not wait for your pods.
\end{Verbatim}
\end{terminalbox}
\figcaption{사례 12-B  아침마다 한 번씩 일어난 두 실패. CPU 기반 오토스케일러는 GPU의 바쁨을 볼 수 없어 전혀 작동하지 않는다. 대기열 기반 오토스케일러는 15초 만에 작동하지만 새 레플리카를 만드는 데 2분 30초가 걸린다. §12.2는 눈을, §12.3은 초시계를 고친다.}

두 매복이 이 장의 뼈대를 만든다. \textbf{§12.2}는 제11장에서 읽는 법을 배운 입장 대기열 지표에 오토스케일러를 연결하고, 목표 공식, 비대칭 안정화 창(stabilization window), 예측 가능한 피크를 위한 예약 증가라는 제어 루프의 실제 조절점을 다룬다. \textbf{§12.3}은 2분 30초의 탄생을 이미지, 가중치, 엔진, 워밍업으로 세분화하고 2주 차 Firecracker를 다시 불러 업계가 이미 해결한 항목이 무엇인지 설명한다. 이어 완화 사다리를 올라 0으로 스케일링과 그 정직한 틈새에 도달한다. \textbf{§12.4}는 사례 1-B가 드러낸 \textit{또 다른} 낭비인 GPU 한 대를 온전히 쓰기에는 너무 작은 서비스를 공격한다. 시분할, MPS, MIG로 제2장의 카드당 테넌트 하나 규칙을 굽힌다. \textbf{§12.5}는 Insull의 포트폴리오를 클라우드 가격으로 조립한다. 약정 바닥, 온디맨드 어깨, 스팟 버스트, 비혼잡 시간 채우기를 사용해 프로젝트 비용을 \$2,381에서 \textbf{\$877}, 즉 영수증을 근거로 1주 차 상한 아래까지 낮춘다. 측정한 단위 경제성(unit economics)으로 제1장의 자체 구축 대 외부 도입 손익분기점도 마무리한다. \textbf{§12.6}은 강의 계획서가 계속 준비해 온 실습으로, 제시하는 모든 수치가 스윕이나 가격표에 근거하는 팀 프로젝트 중간 점검이다.

\begin{calloutbox}{tbbpurple}{tbbpurplefill}
{\normalsize\bfseries\color{tbbpurple}생각해 보기}\par\vspace{3pt}
\begin{tbbitemize}
  \item 사례 12-A의 바닥은 "롤아웃 안전"을 근거로 레플리카 2개다. 제5장의 롤링 업데이트 메커니즘(maxUnavailable/maxSurge)과 제11장의 스트림을 이용해 논거를 재구성하라. 반대편도 주장하라. 바닥을 1이나 0으로 두어도 정당하려면 트래픽과 사용자에 대해 무엇이 참이어야 하는가?
  \item Insull의 수요 계량기는 총소비량이 아니라 피크를 기준으로 고객에게 요금을 부과했다. 이를 클라우드 GPU 가격에 대응시켜라. 클라우드에서 수요 요금에 해당하는 것은 무엇이며, 시간당 온디맨드 가격은 왜 이미 Insull의 논리를 인코딩하는가? 힌트: 각 가격 모델에서 유휴 위험은 누가 부담하는가?
  \item 사례 12-B의 시도 1도 한 가지 트래픽 패턴에서는 실제로 스케일링한다. 그 패턴을 설명하고 GPU 채팅 서비스에서는 절대 나타나지 않는 이유를 말하라. 제9장의 어느 사례를 거울처럼 뒤집은 모습인지도 제시하라.
  \item 팀원이 max=6과 영구적으로 낮은 목표값, 즉 "하루 종일 더 많은 여유 용량"으로 09:00의 매복을 해결하자고 제안한다. 사례 12-A의 원장을 기준으로 가격을 계산하고 실제로 무엇을 사는지 설명하라. 힌트: 계단을 다시 평평하게 만든다.
\end{tbbitemize}
\end{calloutbox}

\section{오토스케일링: 올바른 지표와 초시계}

신비부터 걷어 내자. 어려운 부분은 이미 제5장에서 해결했다. \textbf{Horizontal Pod Autoscaler}는 15초마다 지표를 읽고, 그 지표를 목표에 맞추려면 레플리카가 몇 개 \textit{있어야 할지} 계산한 뒤 Deployment의 \code{replicas} 필드에 답을 쓰는 제어 루프다. 전체 장치는 \textbf{관찰 → 목표 계산 → 정수 하나 수정}이 전부다. 이후의 스케줄링, 이미지 풀, 가중치 로드, 프로브, Service 구성원 관리는 제5장의 조정 메커니즘이 5주 차부터 지켜 온 약속을 이행한다. §5.2의 생각해 볼 문제는 "오토스케일러 자체가 알아야 하는 것이 얼마나 적은지" 알아보라고 했다. 이제 그 보상을 받는다. 목표 계산은 한 줄이다. \textbf{desired = ceil(current × metric ÷ target)}. 레플리카 두 개가 41개의 대기 요청을 처리하고 레플리카당 목표가 10이면 desired = ceil(2 × 20.5 ÷ 10) = \textbf{5}지만 \code{maxReplicas: 4}에 의해 제한된다. 머신러닝도 예측도 없고 기상학자가 아니라 온도 조절기다. 엔지니어링의 핵심은 공식이 주어진 것으로 보는 두 질문에 있다. \textit{어느 지표}를 사용할 것인가, 그리고 \textit{새 용량이 실제로 얼마나 빨리 도착하는가}다.

\begin{tbbfigure}
  \adjustbox{max width=0.94\linewidth,max totalheight=0.46\textheight}{\includegraphics{figures/fig12-2_autoscale.png}}\\[3pt]
  \figcaption{그림 12-2  루프와 두 실패 유형. 왼쪽은 원하는 상태를 수정하는 네 박자 주기이고, 오른쪽은 GPU 서비스에서 CPU\%가 전혀 작동하지 않는 이유다. 아래의 초시계를 보면 감지는 몇 초지만 탄생에는 몇 분이 걸리며 그 차이를 대기열이 먹는다.}
\end{tbbfigure}

\subsection*{눈: 사용자가 느끼는 신호로 스케일링하라}

사례 12-B의 월요일 실패는 논리 오류가 모든 대시보드에서 \textit{보이지 않기} 때문에 분해해 볼 가치가 있다. 기본 HPA는 웹 서버 플릿에서 CPU가 실제 경합 자원이므로 \textbf{CPU\%}를 본다. 채팅 서비스는 구조가 반대다. 제9장 사례 A에서는 GPU가 7\%로 쉬는 동안 CPU가 100\%에 고정되었다. 부하가 걸린 GPU 서비스는 거울상으로, GPU는 94\%인데 토크나이저와 HTTP를 처리하는 CPU는 9\%만 조금 쓰며, 대기열이 끝없이 커지는 동안 CPU만 보는 오토스케일러는 \textit{조용하다}고 결론 낸다. \textbf{GPU 사용률}은 두 번째로 끌리는 선택이지만 두 가지 이유로 역시 틀리다. 먼저 \textit{포화}된다. 카드가 \textasciitilde{}100\%를 표시하면 "수요가 그보다 3× 더 많다"는 사실을 표현할 수 없어 격차가 가장 클 때 오토스케일러가 부족하게 늘린다. 또 \textit{부분 부하에서 거짓말}한다. 제9장에서 커널 실행 병목 서비스는 용량의 절반만 처리하면서도 90\%+를 표시할 수 있음을 배웠다. 올바른 눈은 제11장이 엔진에 넣은 입장 대기열 \textbf{vllm:num\_requests\_waiting}이다. 수요가 좌석을 넘는 순간, 지연 시간 백분위수가 움직이기 전에 0이 아니게 되는 \textit{선행} 지표다. 대기 요청 41개는 용량이 얼마나 틀렸는지 정확히 말하므로 \textit{상한이 없다}. 또한 L = λW에서 대기열은 TTFT의 분자이므로 \textit{사용자가 느끼는 것}이다. 앞으로 작성할 모든 런북(runbook)의 여백에 적을 일반 규칙은 다음과 같다. \textbf{공급 측 게이지(CPU\%, GPU\%, 메모리)가 아니라 수요 신호(대기열 깊이, 진행 중 동시성)로 스케일링하라.} 게이지는 장치가 얼마나 열심히 일하는지 말하고 대기열은 사용자가 얼마나 기다리는지 말한다. SLO는 사용자에 관해 쓴다.

기계적으로 사용자 정의 지표는 두 경로 중 하나로 오토스케일러에 도달한다. \textbf{HPA 경로}는 PromQL을 Kubernetes 사용자 정의 지표 API로 변환하는 지표 어댑터(prometheus-adapter)를 설치한다. 강력하지만 손이 많이 가고 최소 레플리카가 1에서 멈춘다. 실습에서 사용하는 \textbf{KEDA 경로}에서 KEDA는 70개가 넘는 \textit{스케일러}(Prometheus, SQS, Kafka, cron 등)를 갖춘 오퍼레이터이며 CRD 하나로 HPA를 대신 관리한다. HPA와 달리 \textbf{0으로} 스케일링할 수 있다. §12.3에서 0이 왜 전혀 다른 대상인지 설명한다. \code{ScaledObject} 하나가 이 절의 모든 주장을 표현한다.

\begin{terminalbox}
\begin{Verbatim}[breaklines=true,breakanywhere=true,fontsize=\small,formatcom=\color{tbbtermfg},codes=\tbbcodeactive]
apiVersion: keda.sh/v1alpha1
kind: ScaledObject
metadata: {name: chat-scaler}
spec:
  scaleTargetRef: {name: chat}          # the Deployment to edit
  minReplicaCount: 2                    # the floor (rollout safety)
  maxReplicaCount: 4                    # the budget ceiling
  cooldownPeriod: 300                   # scale DOWN slowly (5 min)
  triggers:
  - type: prometheus
    metadata:
      serverAddress: http://prometheus:9090
      query: sum(vllm:num_requests_waiting{app="chat"})
      threshold: "10"                   # tolerate 10 waiting/replica
# scale on the queue users feel - not the CPU nobody uses.
\end{Verbatim}
\end{terminalbox}
\figcaption{기록 12-1  실습의 ScaledObject. 15줄에 올바른 지표(vLLM 입장 대기열), 안전을 위한 바닥, 예산을 위한 상한, 올릴 때는 빠르고 내릴 때는 느린 비대칭 성향이 담긴다.}

\subsection*{성향: 올릴 때는 빠르게, 내릴 때는 느리게}

공식만 사용하면 플릿은 신경질적으로 변한다. 트래픽이 출렁일 때 모든 흔들림을 따라 하는 스케일러는 \textbf{플래핑}한다. 레플리카가 나타나 2분 30초의 탄생 비용을 내고 8분간 서빙한 뒤 죽었다가 다음 흔들림에 다시 태어난다. 청구서 전체를 콜드 스타트 오버헤드로 바꾼다. 제10장의 생각해 볼 문제에서 최악의 특수 사례를 만났다. 스케일 아웃하는 포드마다 도착 시 40초의 torch.compile 워밍업 비용을 내고 매번 p99가 치솟았다. 업계의 해결책은 \textbf{비대칭 안정화}다. 실제 대기열은 평활화를 기다리게 해서는 안 되고 몇 초가 중요하므로 순간 신호에 따라 \textbf{증가}시킨다. \textbf{감소}는 \textit{후행 창의 최악값}을 기준으로만 한다. HPA에서는 \code{scaleDown.stabilizationWindowSeconds: 300}, KEDA에서는 \code{cooldownPeriod}다. 이 비대칭은 CFO에게도 옹호할 수 있는 비대칭 손실 함수를 인코딩한다. 너무 빠르게 늘리면 레플리카 몇 분을 낭비하지만 너무 빠르게 줄이면 다섯 장에 걸쳐 측정한 SLO를 깨뜨린다. 감소 창은 트래픽의 흔들림 주기에 맞춘다. 일중 채팅에는 5분이 적합하고 배치형 트래픽에는 더 길게 둔다. \textbf{바닥}은 부하가 아니라 운영을 기준으로 정한다. 레플리카 두 개면 제5장의 롤링 업데이트 중 하나가 서빙하는 동안 다른 하나를 종료할 수 있다. 바닥이 하나면 모든 배포가 부분 장애가 된다.

\subsection*{초시계: 반응 시간도 용량의 일부다}

이제 수요일의 매복을 보자. 올바른 눈과 차분한 성향이 있어도 루프의 종단 간 반응 시간은 \textbf{T = 감지(≤15 s 지표 스크레이프) + 결정(≤15 s 루프 틱) + 프로비저닝(§12.3의 콜드 스타트: 따뜻한 노드의 포드 \textasciitilde{}150 s, GPU 노드까지 태어나야 하면 6–10 min)}이다. 정직한 경우에도 3분으로 보자. 그동안 \textit{기존} 플릿이 \textit{새} 수요를 흡수하고 대기열이 차이를 적분한다. 09:00의 파도는 좌석 64개에 제공되는 부하를 동시 스트림 64개에서 104개로 올린다. 초당 \textasciitilde{}9건의 턴 도착이 좌석을 찾지 못하고, Ready 시점에는 스트림 약 40개 분량의 수요가 대기열에 모인다. 기록 11-3의 구간 3과 정확히 같고 매일 아침 예약된 시각에 TTFT p99가 붉어진다. 세 가지 해결책을 적용할 순서대로 살펴보자. 볼 수 있는 피크에는 \textbf{예약 스케일링(scheduled scaling)}을 사용한다. 평일마다 09:00에 오르는 트래픽은 예측 문제가 아니라 달력 문제다. 08:40에 바닥을 4로 올리는 KEDA cron 트리거는 40분 × 레플리카 2개 ≈ \textbf{하루 \$1.07}의 비용으로 매복을 없앤다. 볼 수 없는 피크에는 \textbf{여유 용량}을 둔다. 대기 10개 대신 5개를 허용해 목표를 조금 넉넉히 잡거나 영업시간 동안 레플리카 하나를 더 유지하면(\textasciitilde{}\$50/mo) 대기열이 감지 + 탄생 시간을 버틸 시간을 산다. 마지막은 \textbf{탄생 자체를 줄이는 것}이며 다음 절 전체의 주제다. 효과 없는 방법은 루프가 더 빨라지기를 바라는 것이다. 150초 탄생 앞에서 1초 스크레이프(scrape) 간격은 매복을 더 정밀하게 측정할 뿐이다.

\begin{terminalbox}
\begin{Verbatim}[breaklines=true,breakanywhere=true,fontsize=\small,formatcom=\color{tbbtermfg},codes=\tbbcodeactive]
# wednesday again, with the 08:40 cron floor + KEDA queue trigger:
$ kubectl get hpa keda-hpa-chat-scaler -w
TIME      TARGETS        REPLICAS   EVENT
08:40:02  0/10 (cron:4)  2 -> 4     # calendar, not clairvoyance
08:42:33  0/10           4          # 2 births paid BEFORE the wave
09:00:00  wave arrives: streams 64 -> 104
09:00:15  3/10           4          # queue: single digits. absorbed.
09:14:00  peak passes; waiting=0 for 5 min (cooldown)...
09:19:05  0/10           4 -> 3     # reluctant down. no flapping.
# TTFT p99 through the whole event: 0.46 s. the ambush is deleted -
# for $1.07 of pre-warm. compare Exhibit 12-B's three red minutes.
\end{Verbatim}
\end{terminalbox}
\figcaption{기록 12-2  파도를 견딘 결과. cron 트리거가 피크 전에 두 탄생 비용을 내고, 대기열 트리거가 곡선을 처리하며, 쿨다운이 차분하게 감소시킨다. 약 1달러로 사건 내내 SLO를 지킨다. 이 교재가 판매하는 가장 저렴한 보험이다.}

경계를 정직하게 밝히자. 앞의 모든 방법은 \textbf{프로비저닝된 노드 풀 안의 포드}를 스케일링한다. 원하는 포드 수가 노드 용량을 넘으면 \textbf{클러스터 오토스케일러}가 GPU \textit{장치}를 탄생시켜야 한다. 용량 찾기, 부팅, 드라이버 설치에 3–7분이 걸리며 때로는 "no capacity in zone"으로 끝난다. GPU 부족은 실제 날씨와 같고 제2장의 할당량 실습은 온화한 예고편이었다. 수치에서 운영 규칙이 나온다. \textbf{포드는 파도를 쫓고 노드는 예측을 따른다.} 하루의 노드 범위, 즉 계단의 최댓값을 미리 프로비저닝하고 그 안에서 포드 오토스케일링이 숨 쉬게 한다. 노드 오토스케일링은 느린 하루의 조수로 다뤄야 하며 09:00 파도와 TTFT 조항 사이에 세워서는 안 된다.

\begin{calloutbox}{tbbpurple}{tbbpurplefill}
{\normalsize\bfseries\color{tbbpurple}생각해 보기}\par\vspace{3pt}
\begin{tbbitemize}
  \item 기록 11-3의 구간 3(running 32, waiting 26, 레플리카 하나, target 10)에 목표 공식을 적용하라. HPA는 몇 개를 요청하고 maxReplicas는 어디에서 제한하는가? §11.3의 좌석 계산으로 그 상한이 옳은지 판단하고, 이를 높이려면 어떤 예산 수치가 정당화해야 하는지 설명하라.
  \item 트래픽에는 예측 가능한 09:00 증가와 예측할 수 없는 소셜 미디어 스파이크(90초 안에 3×, 10분 뒤 사라짐)의 두 피크가 있다. 둘을 위한 트리거 집합(cron + 대기열 + 여유 용량)을 설계하라. 의도적으로 희생하는 스파이크와 그것이 CFO와 나눌 올바른 대화인 이유도 밝히라.
  \item 팀원이 "더 빨리 비용을 아끼기 위해" scaleDown 안정화를 30 s로 설정한다. 약 7분마다 피크가 오는 출렁이는 트래픽 한 시간을 모델링하라. 탄생 횟수를 세고 각각 150 s인 콜드 스타트 오버헤드 가격을 계산해 "절감"이 음수임을 보인 뒤, 팀원이 옳아지는 트래픽 패턴을 찾아라.
  \item 제5장의 용어인 롤링 업데이트, maxUnavailable, 프로브 게이팅만 사용해 비용 검토자에게 바닥 2개를 옹호하라. 이어 바닥의 월 비용을 계산하고 이를 1로 낮출 수 있는 제품 결정 하나를 제시하라.
\end{tbbitemize}
\end{calloutbox}

\section{콜드 스타트와 0으로 스케일링}

오토스케일러의 힘은 레플리카가 태어나는 속도만큼만 빠르다. 이 절에서는 제7장이 지연 시간에 했던 일을 콜드 스타트에 적용한다. \textbf{하나의 수치로 취급하지 않고 청구 항목으로 세분화한다.} 지난 아홉 장 동안 항목을 하나씩 모았다. 제4장은 이미지 풀의 가격을 계산했고, 제6장은 가중치 로딩 막대를 만들며 "12주 차가 이 막대를 두고 싸운다"고 경고했다. 제10장은 런타임 컴파일 비용을, 제11장은 엔진의 KV 풀 할당과 CUDA 그래프 캡처를 추가했다. 포드 하나가 이 모든 비용을 한꺼번에 내는 모습을 보자.

\begin{terminalbox}
\begin{Verbatim}[breaklines=true,breakanywhere=true,fontsize=\small,formatcom=\color{tbbtermfg},codes=\tbbcodeactive]
$ kubectl describe pod chat-7f9b4-x2kkj | grep -A2 Events; kubectl logs -f ...
09:00:24  Scheduled    assigned to node gpu-pool-b7   #     +2 s
09:00:26  Pulling      image "registry/chat:v9"
09:00:42  Pulled       3.1 GB in 16.1 s               # Ch. 4's slim image,
                                                      # warm node layer cache
09:00:44  vllm  INFO   Fetching weights: 16-way ranged GET
09:00:53  vllm  INFO   6.02 GiB in 8.9 s              # Ch. 6's loader
09:02:28  vllm  INFO   engine init: VRAM load, KV pool alloc,
                       CUDA graph capture... 95 s     # Ch. 11's engine
09:02:54  Ready        readiness probe passed          # after 26 s of
                                                      # scripted warm-up
# total: 150 s. the pod path. if a NODE had to be born first:
# +3-7 min (VM provision, boot, driver install) - and sometimes
# "no capacity in zone." never put a node birth on the SLO path.
\end{Verbatim}
\end{terminalbox}
\figcaption{기록 12-3  한 번의 탄생을 항목별로 나눈 결과. 모든 줄에 해당 장이 있다. 16초 풀은 제4장의 엔지니어링, 9초 가중치 가져오기는 제6장, 95초 엔진 초기화는 제11장이다. 합계 150초는 §12.2의 초시계 문제를 정량화한다.}

\begin{tbbfigure}
  \adjustbox{max width=0.94\linewidth,max totalheight=0.46\textheight}{\includegraphics{figures/fig12-3_coldstart.png}}\\[3pt]
  \figcaption{그림 12-3  청구서, 두 경로, 완화 사다리. 따뜻한 노드의 포드는 \textasciitilde{}2.5분이 걸리고 이제 엔진 초기화가 지배한다. 노드 경로는 이를 몇 배로 늘린다. 완화 사다리의 각 단은 이미 이 교재에서 배운 한 장에 해당한다.}
\end{tbbfigure}

\subsection*{다시 등장한 Firecracker — 해결했지만 엉뚱한 항목}

2주 차에 심어 둔 약속을 이제 지켜야 한다. Firecracker의 microVM은 \textbf{\textasciitilde{}125 milliseconds}에 부팅한다. 2005년에 몇 분 걸리던 격리 환경은 2018년에 사실상 공짜가 되었다. 제2장 §2.4는 12주 차의 교훈을 미리 말했다. \textit{"병목은 환경이 아니라 페이로드가 되었다."} 기록 12-3을 이 주장에 비추어 읽어 보자. VM/컨테이너 메커니즘은 2초짜리 \code{Scheduled} 줄에 기여하고 나머지는 거의 없다. 남은 148초는 \textbf{페이로드}다. 3.1 GB 이미지, 6 GB 가중치, CUDA 컨텍스트, 엔진 전체의 할당과 그래프 캡처가 차지한다. 서버리스 CPU는 영리한 문제였지만 \textbf{서버리스 GPU}가 진짜 어려운 이유다. Lambda의 기법인 미리 부팅한 microVM 풀 + 스냅샷/복원은 밀리초 규모 문제를 공격한다. GPU 버전은 \textit{기가바이트}를 미리 배치하고 초기화된 CUDA 프로세스를 되살려야 한다. 초 단위 GPU 요금을 파는 플랫폼(Modal, RunPod, Cloud Run GPU)은 따뜻한 풀, 공격적 이미지 스트리밍, 점차 늘어나는 \textbf{GPU 체크포인트/복원}으로 바로 이 전쟁을 치른다. 완전히 초기화한 엔진 프로세스를 스냅샷하고 몇 초 안에 복원한다. 이 장의 더 읽을거리에 있는 "1초 미만 콜드 스타트" 엔지니어링 글은 모두 이 기법의 변형이다. 가격표를 보면 전쟁이 공짜가 아님을 알 수 있다. 초 단위 요금은 원시 VM보다 프리미엄이 붙는다. 교재의 플릿에는 이런 플랫폼이 아니라 사다리가 필요하다.

\subsection*{완화 사다리 — 0이 차지하는 정직한 틈새}

그림 12-3의 사다리를 가장 저렴한 단부터 오르면 이미 모든 단을 보유했음을 알 수 있다. \textbf{① 이미지 위생}(제4장): 따뜻한 노드 캐시가 있는 다단계 3.1 GB 이미지는 16 s가 걸렸지만 단순한 8.9 GB 이미지는 1분이 걸렸다. 지연 풀 스냅샷터(eStargz, 이미지 스트리밍)는 이미지가 모두 도착하기 전에 컨테이너를 시작한다. \textbf{② 가중치 물류}(제6장): 16방향 범위 GET + safetensors mmap은 260 s를 9초로 줄였다. 노드 로컬 캐시는 같은 노드의 \textit{두 번째} 탄생을 디스크 읽기로 바꾼다. \textbf{③ 준비 상태 안의 워밍업}(제5/10장): 프로브가 스크립트로 추론을 실행해 \code{Ready}가 \textit{따뜻함}을 뜻하게 한다. 피할 수 없는 작업이지만 올바르게 게이트하면 LB가 차가운 포드로 라우팅하지 않고 제10장의 플래핑이 재발하지 않는다. \textbf{④ 사전 워밍업 여유 용량}(§12.2): 영업시간 동안 레플리카 하나를 더 두는 \textasciitilde{}\$50/mo는 알려진 환율로 탄생 지연 시간을 돈으로 바꾼다. \textbf{⑤ 예약 스케일링}: 08:40 cron으로 수요가 생기기 전에 탄생 비용을 낸다. \textbf{⑥ 0으로 스케일링}: 가장 먼 단이며 전혀 다른 대상이다. 0 이후 \textit{첫} 요청이 150 s 전체, 노드 경로라면 몇 분을 사용자에게 보이는 지연 시간으로 먹기 때문이다. KEDA는 탄생 중 요청을 세워 두는 HTTP 애드온으로 이를 지원하고, Knative는 이를 신념처럼 만들었다. 정직한 결정 규칙은 한 문장이다. \textbf{0으로 스케일링은 호출자가 첫 적중의 몇 분 지연을 감당할 수 있는 엔드포인트에 쓴다.} 주 2회 쓰는 롱테일 내부 모델은 유휴 L4 \$595/mo를 \textasciitilde{}\$3로 줄일 수 있고, 개발 환경과 배치 트리거 스코어러도 해당한다. 채팅 플릿의 바닥은 2로 남는다. 03:00의 불면증 사용자도 돈을 내고, 제5장의 롤아웃에는 짝이 필요하다. 엔드포인트마다 계산하면 플릿은 자연스럽게 나뉜다. \textbf{롱테일은 0, SLO를 지는 헤드는 바닥을 둔다.}

\begin{calloutbox}{tbbpurple}{tbbpurplefill}
{\normalsize\bfseries\color{tbbpurple}생각해 보기}\par\vspace{3pt}
\begin{tbbitemize}
  \item 기록 12-3에서 엔진 초기화(95 s)가 이제 지배적인 항목이다. 제10–11장을 이용해 모델 크기와 함께 증가하는 부분, KV 풀 크기와 함께 증가하는 부분, 고정된 부분을 구분하라. 가장 유망한 공격 두 가지를 제안하고 각각의 절감 시간을 추정하라. 힌트: AWQ 가중치는 1.9 GB이고 CUDA 그래프 캡처 범위를 줄일 수 있다.
  \item 세 엔드포인트에서 0으로 스케일링의 가격을 계산하라. (a) 채팅 플릿, (b) 야간 배치 스코어러, (c) 주 \textasciitilde{}5회 사용하는 내부 RAG 데모다. 각각 유휴 비용 절감액, 첫 적중 지연 시간 비용, 결론을 제시하고 (a)의 결론을 뒤집을 제품 문장을 말하라.
  \item "1초 미만 GPU 콜드 스타트" 블로그 글이 화제가 된다. 이 절을 바탕으로 제목 뒤에서 무엇이 참이어야 하는지 나열하라. 무엇을 스냅샷했는가? 가중치는 어디에 사는가? 따뜻한 풀 비용은 누가 내는가? 벤치마크가 거의 확실히 포함하지 않은 것도 밝히라.
  \item 클러스터 오토스케일러가 피크에 L4에 대해 가끔 "no capacity in zone"을 보고한다. 다른 영역 → 다른 GPU 유형 → 온디맨드→스팟 반전 → 대기열 후 성능 저하로 이어지는 폴백 연쇄를 설계하고, 제5장의 nodeSelector/affinity 메커니즘으로 지금 표현할 수 있는 단계를 표시하라.
\end{tbbitemize}
\end{calloutbox}

\section{실리콘 공유: 시분할, MPS, MIG}

오토스케일링은 \textit{너무 많은 전체 GPU}의 낭비를 고친다. 이 절에서는 반대편의 부끄러운 문제를 해결한다. 제1장이 교재를 시작하며 보여 준 \textbf{사례 1-B의 유령}이다. 사용률 0.2\%로 한 달에 \$779.66가 청구된 A10G는 GPU 한 대를 쓰기에는 너무 작은 서비스가 카드 전체를 소유한 경우였다. 제2장 §2.6은 낭비가 구조적인 이유를 설명했다. GPU는 가상화에 저항하고 VT-x에 해당하는 기능이 없으며 커널이 드라이버 대기열의 작업을 선점할 수 없다. 그래서 클라우드 대여 단위는 \textbf{테넌트 하나, 카드 하나}가 되었다. 패스스루의 깔끔한 격리를 잔혹한 세분성으로 제공했다. 또한 12주 차에 이 규칙을 세 가지로 굽히겠다고 약속했다. 얼마나 많은 격리를 포기해 얼마나 많은 사용률을 회수하는지에 따라 세 방법을 순서대로 살펴보자.

\begin{tbbfigure}
  \adjustbox{max width=0.94\linewidth,max totalheight=0.46\textheight}{\includegraphics{figures/fig12-4_sharing.png}}\\[3pt]
  \figcaption{그림 12-4  공유의 스펙트럼. 시분할은 전체 다이를 번갈아 쓰며 격리가 없다. MPS는 다이를 동시에 공유하고 장애 반경이 하나다. MIG는 하드웨어에서 다이를 분할해 실제 벽을 만들지만 세분성이 거칠다. 제2장의 규칙을 실리콘 자체가 굽힌다.}
\end{tbbfigure}

\subsection*{시분할: 번갈아 쓰고 정직하게 표시하기}

\textbf{시분할}은 가장 단순하게 규칙을 굽힌다. 1970년대 시분할 운영체제가 CPU 하나를 공유했듯 드라이버가 몇 밀리초씩 테넌트 사이에서 \textit{전체 GPU}의 컨텍스트를 전환한다. Kubernetes는 장치 플러그인 설정으로 이를 연결한다. \code{replicas: 4}를 광고하면 물리 카드 한 장이 \code{nvidia.com/\allowbreak gpu: 4}로 나타나고, 각각 자신이 소유한다고 믿는 포드 네 개를 스케줄링할 수 있다. 어떤 GPU에서도 동작하고 실습의 T4도 포함되며, ConfigMap 하나로 설정할 수 있고 사용률 10\%인 서비스 네 개가 카드 3.6장을 낭비하지 않게 된다. 그 대신 제2장이 패스스루에서 중요하게 본 모든 것을 포기하며 실습에서 이를 \textit{측정}한다. \textbf{메모리 격리가 없다.} VRAM은 먼저 온 사람이 먼저 쓰고 한 테넌트의 누수가 이웃을 OOM으로 만든다. 사례 11-B의 증가가 이제 \textit{공유} 증가가 된다. \textbf{성능 격리가 없다.} 자신의 p99에 이웃의 배치 크기가 들어오고 컨텍스트 전환(context switch) 자체가 모든 테넌트에 비용을 부과한다. \textbf{장애 격리가 없다.} Xid 오류 하나가 모두의 오후를 초기화한다. 정직한 사용 사례는 자연스럽다. 개발 장치, 노트북, 대화형 실험, 모든 이웃이 자기 팀인 사용률 낮은 \textit{내부} 서비스다. 고객 대상 SLO 아래에서 시분할을 쓰면 제14장의 설명할 수 없는 p99 전쟁 이야기가 탄생한다.

\begin{terminalbox}
\begin{Verbatim}[breaklines=true,breakanywhere=true,fontsize=\small,formatcom=\color{tbbtermfg},codes=\tbbcodeactive]
# lab part A-④: slice the T4 four ways, then MEASURE the rent.
$ kubectl apply -f timeslice-config.yaml   # device-plugin: replicas: 4
$ kubectl describe node gpu-node | grep nvidia.com/gpu
  nvidia.com/gpu:  4        # one card, four schedulable "GPUs"
$ kubectl scale deploy/clf --replicas=2    # two classifiers co-tenant

# same sweep as Week 9, alone vs. with a busy neighbor:
            p50      p99      req/s@SLO
alone      2.9ms    9.8ms      507
neighbor   3.1ms   38.4ms      214       # p99 x3.9, capacity /2.4
# nothing is "broken." the card is just honest: two tenants, one die.
# noisy neighbor is not folklore - it is a number you now own.
\end{Verbatim}
\end{terminalbox}
\figcaption{기록 12-4  측정한 공유. 분류기 포드 두 개가 실습 T4를 시분할한다. 각자 실행할 때는 문제가 없지만 함께 실행하면 p99가 네 배가 된다. 제2장이 패스스루로 샀다고 한 격리가 사라지자 그 부재가 보인다. 공유는 실제 할인이지만 실제 대가도 있다.}

\subsection*{MPS와 MIG: 중간 지점과 하드웨어의 답}

\textbf{MPS}(Multi-Process Service)는 중간 단이다. 번갈아 쓰는 대신 테넌트 프로세스의 커널이 공유 CUDA 컨텍스트 하나를 통해 \textbf{동시에} 실행되며, 선택적으로 클라이언트별 SM 비율 상한을 둘 수 있다. 컨텍스트 전환 비용이 없고 유휴 SM이 실제로 이웃 작업을 흡수하므로 사용률은 시분할보다 높다. Volta+부터 주소 공간 격리도 생겼다. 그러나 메모리 대역폭은 여전히 공유하고 장애 반경도 같다. 그래서 MPS는 \textbf{한 팀의 여러 작은 모델}에 가장 적합하다. 자신이 자기 자신의 시끄러운 이웃이므로 스스로 협상할 수 있다. \textbf{MIG}는 제2장 §2.6이 약속한 답이다. Ampere 이후의 큰 다이(A100, A30, H100이며 실습의 T4나 플릿의 L4는 \textit{해당하지 않는다})는 \textbf{하드웨어에서} 최대 일곱 인스턴스로 분할된다. 각각 자체 SM 조각, L2 조각, DRAM 컨트롤러를 가진다. MIG 조각은 진짜 작은 GPU다. 자체 OOM과 장애가 있고, 대역폭은 경합하지 않고 \textit{프로비저닝}된다. 바로 제2장이 표현한 대로 격리가 실리콘 안으로 이동했다. Kubernetes는 각 프로파일을 확장 자원(\code{nvidia.com/\allowbreak mig-1g.10gb})으로 본다. 플릿 설계 질문은 계산 문제가 된다. 9주 차 분류기는 \textasciitilde{}10 GB와 적당한 연산 성능이 필요하다. 전용 T4는 사용률 \textasciitilde{}30\%에서 \$0.53/hr다. \textbf{A100의 1g.10gb 조각 하나는 ≈ \$3.60 ÷ 7 ≈ \$0.51/hr}다. 같은 서비스를 운영하면서 형제용 조각 여섯 개가 남고 유령 GPU 세금을 돌려받는다. 함정은 운영과 §10.2의 모양을 닮았다. 1/7 조각은 \textit{대역폭}도 \textasciitilde{}1/7이므로 조각마다 두 줄 진단을 해야 한다. 얇은 조각의 메모리 병목 모델은 여전히 메모리 병목이다. MIG 구성을 바꾸려면 카드를 드레인해야 하므로 노드 풀별 프로파일을 계획하고 시간마다 바꾸지 않는다. 통합은 \textit{노드} 수준 장애 반경을 집중시킨다. 이제 호스트 하나가 실패하면 서비스 일곱 개가 사라지므로 제5장의 안티어피니티(anti-affinity)가 제값을 한다.

\begin{tbbtablewrap}
  \begin{tabular}{@{}>{\raggedright\arraybackslash}m{\dimexpr 0.1862\tbltextwidth-2\tabcolsep\relax} >{\raggedright\arraybackslash}m{\dimexpr 0.2553\tbltextwidth-2\tabcolsep\relax} >{\raggedright\arraybackslash}m{\dimexpr 0.2606\tbltextwidth-2\tabcolsep\relax} >{\raggedright\arraybackslash}m{\dimexpr 0.2979\tbltextwidth-2\tabcolsep\relax}@{}}
  \arrayrulecolor{tbbnavy}\specialrule{1pt}{0pt}{0pt}
  \rowcolor{tbbtablehead}{\centering \bfseries\color{tbbnavy}모드} & {\centering \bfseries\color{tbbnavy}격리} & {\centering \bfseries\color{tbbnavy}동작 대상 / 세분성} & {\centering \bfseries\color{tbbnavy}정직한 용도} \\
  \arrayrulecolor{tbbnavy}\specialrule{0.7pt}{0pt}{2pt}
  {\textbf{시분할}} & {없음 — VRAM과 장애 공유, 전체 다이를 번갈아 사용} & {모든 GPU (T4 ✓) · 카드당 N개 가상 장치} & {개발, 노트북, 사용률 낮은 내부 서비스 — 고객 SLO에는 사용하지 않음} \\
  \arrayrulecolor{tbbtableborder}\hline
  {\textbf{MPS}} & {동시 커널, SM 상한. 대역폭과 장애 반경 공유} & {Volta+ · 부분 SM 공유} & {한 팀의 여러 작은 모델 — 자신이 자기 이웃} \\
  \arrayrulecolor{tbbtableborder}\hline
  {\textbf{MIG}} & {하드웨어: 자체 SM, L2, DRAM 채널, 자체 OOM} & {A100/A30/H100+ · 최대 7개 조각} & {대형 카드에서 SLO를 지는 멀티테넌트 서빙. 조각 ≈ 작은 GPU} \\
  \arrayrulecolor{tbbtableborder}\hline
  {\textbf{(전체 카드)}} & {패스스루 — 제2장의 깔끔한 기준} & {모든 GPU · 테넌트당 1개} & {플릿의 L4: 카드를 실제로 채우는 인기 모델 하나} \\
  \arrayrulecolor{tbbnavy}\specialrule{1pt}{2pt}{0pt}
  \end{tabular}\par
  \figcaption{표 12-1  공유 결정 요약. 위에서 아래로 읽으며 워크로드가 감당할 수 있는 격리를 제공하는 첫 행을 선택한다.}
\end{tbbtablewrap}

마지막 관점으로 이 절을 장의 주제와 연결해 보자. 공유는 \textbf{Insull의 다양성 기법을 다이 규모에 적용한 것}이다. 그의 발전소는 피크가 맞물렸기 때문에 전차, 조명, 제빙소를 함께 처리했다. MIG로 나눈 A100도 같은 기대를 바탕으로 작은 서비스 일곱 개를 처리한다. 실패 유형도 같다. 피크가 \textit{같은} 고객 일곱을 카드 하나에 넣으면 한 계층 아래에서 13:00 문제를 다시 만든다. 수요의 다양성이 자산이므로 통합하기 전에 측정해야 한다.

\begin{calloutbox}{tbbpurple}{tbbpurplefill}
{\normalsize\bfseries\color{tbbpurple}생각해 보기}\par\vspace{3pt}
\begin{tbbitemize}
  \item 기록 12-4에서 이웃 세금으로 p99는 9.8 → 38.4 ms가 되었지만 p50은 2.9 → 3.1로 거의 움직이지 않았다. 시분할 메커니즘으로 두 수치를 모두 설명하라. 컨텍스트 전환이 운 나쁜 요청과 평균적인 요청에 각각 어떤 영향을 주는가? 평균과 꼬리 중 어느 SLO 조항 형식이 공유를 기만적으로 안전해 보이게 하는가?
  \item 가상의 학과를 위한 MIG 구성을 설계하라. 9주 차 분류기(10 GB, 스파이크형), RAG 임베더(6 GB, 안정적), 내부 요약기(20 GB, 영업시간), 개발 노트북 두 개를 A100-80GB 한 대에 배치한다. 프로파일(1g.10gb/2g.20gb/3g.40gb 혼합)을 고르고 사용 주기로 정당화한 뒤 카드에 넣지 않을 워크로드 하나를 제시하라.
  \item 1g.10gb 조각은 \$0.51로 T4의 \$0.53보다 "유리"하다. 하지만 둘 모두에 §10.2의 같은 두 줄 진단을 적용해 정직하게 다시 비교하라. T4는 320 GB/s, A100의 1/7 조각은 \textasciitilde{}250 GB/s다. 시간당 가격이 아니라 SLO로 결정하라. 언제 T4가 여전히 유리한가?
  \item MPS는 "자신이 자기 자신의 시끄러운 이웃"이라고 설명한다. 자신의 프로젝트에서 피크가 완전히 상관된 같은 트리거의 서비스 두 개로 반례를 구성하고, 피크에는 MPS가 아무것도 절감하지 못하는 이유를 보이라. Insull의 다양성 요구를 CUDA로 다시 표현한 문제다.
\end{tbbitemize}
\end{calloutbox}

\section{KServe와 FinOps 포트폴리오}

두 가지 통합으로 기술 과정이 마무리된다. 하나는 \textit{YAML}, 다른 하나는 \textit{돈}을 위한 통합이다. YAML부터 보자. 이 장에서 설정한 모든 요소를 모델 하나에 적용하면 이제 줄마다 이해하는 약 300줄이 된다. Deployment, 느린 로드를 감싼 프로브(제5장), 저장소 URI에서 모델 가져오기(제6장), 서빙 런타임(제9/11장), Service와 인그레스(제5–6장), 대기열 지표 오토스케일러(§12.2), 적절한 곳에서 0으로 스케일링(§12.3)이 들어간다. 이를 모델 열 개와 팀 세 개로 늘리면 YAML이 스스로 잘못된 내용을 쓰기 시작한다. \textbf{KServe}는 대신 작성해 주는 플랫폼이다. \code{InferenceService} CRD 하나(\textasciitilde{}15줄)에 \textit{모델, 런타임, 오토스케일링 방침, 트래픽 분할}을 선언하면 오퍼레이터가 나머지를 생성한다. 제9장이 KServe를 프레임워크 지도에서 제외하고 다른 종이라고 부른 이유다. \textbf{프레임워크는 "레플리카 하나가 어떻게 잘 서빙하는가?"에 답하고 플랫폼은 "여러 레플리카가 어떻게 살아가는가?"에 답한다.}

\begin{terminalbox}
\begin{Verbatim}[breaklines=true,breakanywhere=true,fontsize=\small,formatcom=\color{tbbtermfg},codes=\tbbcodeactive]
$ kubectl apply -f - <<EOF
apiVersion: serving.kserve.io/v1beta1
kind: InferenceService
metadata: {name: chat}
spec:
  predictor:
    model:
      modelFormat: {name: vLLM}
      storageUri: s3://models/chat/v9-int4/   # Ch. 6's shelf, Ch. 10's gate
    minReplicas: 2
    maxReplicas: 4
    scaleMetric: concurrency                  # §12.2's lesson, built in
    scaleTarget: 28
  canaryTrafficPercent: 10                    # 90/10 split, old/new -
EOF                                           # Chapter 13 drives this knob
inferenceservice.serving.kserve.io/chat created
$ kubectl get isvc chat
NAME   URL                        READY   LATEST   PREV   AGE
chat   http://chat.team.svc...    True    90       10     61s
\end{Verbatim}
\end{terminalbox}
\figcaption{기록 12-5  15줄로 표현한 플랫폼 하나. InferenceService는 이 교재의 제5–12장을 필드로 인코딩한다. canaryTrafficPercent 줄은 제13장을 예고한다. CI 파이프라인이 제어하기를 기다리는 트래픽 분할이다.}

\begin{tbbfigure}
  \adjustbox{max width=0.94\linewidth,max totalheight=0.46\textheight}{\includegraphics{figures/fig12-5_kserve.png}}\\[3pt]
  \figcaption{그림 12-5  CRD와 CRD가 만들어 내는 것. 모델과 팀이 늘어나면 플랫폼을 도입한다. 인기 모델 하나가 제품 전체라면 원시 Deployment + KEDA를 유지한다. 이제 플랫폼이 대체할 모든 줄을 이해한다.}
\end{tbbfigure}

도입 규칙은 제9장의 프레임워크 조언과 닮아 익숙할 것이다. 여러 모델, 릴리스마다 카나리를 적용하려는 계획(제13장), 잠들어야 할 롱테일 엔드포인트(§12.3)처럼 \textbf{플랫폼의 견해가 잘못 쓰던 YAML에서 구해 줄 때 플랫폼을 도입한다.} 인기 모델 하나가 제품 전체라면 \textbf{도입을 미룬다.} 프로젝트의 원시 Deployment + KEDA가 더 작고 완전히 이해할 수 있으며 의존성도 하나 적다. 관리형 엔드포인트인 SageMaker와 Vertex는 오퍼레이터까지 외주화한 같은 절충이며 제13장에서 다시 다룬다. 어느 쪽이든 KServe가 제도화한 내용에 주목하자. 기본 스케일링 지표는 CPU가 아니라 요청의 \textbf{레플리카당 동시성}이다. §12.2의 주장이 교훈에서 플랫폼 기본값으로 승격되었다.

\subsection*{포트폴리오: Insull이 팔던 방식으로 용량 구매하기}

이제 돈을 보자. 지금까지는 소비하는 레플리카 시간이 \textit{몇 시간}이어야 하는지 적정 규모를 정했다. 마지막 수단은 \textit{각 시간}의 가격을 정하며, 선택지는 Insull의 요금표가 다시 태어난 모습이다. 클라우드가 같은 GPU 시간을 세 가격으로 파는 것은 판매자 관점에서 그와 같은 문제인 유휴 용량을 관리하기 때문이다. \textbf{약정 사용}(1–3년 약정, 30–60\% 할인)은 기본 부하 요금이다. 실행할 것이 확실한 바닥의 유휴 위험을 구매자가 맡아 양쪽 모두 이익을 얻는다. \textbf{온디맨드}(정가)는 부하 추종 요금이다. 약속이 전혀 없고 프리미엄 자체가 \textit{유연성}이다. \textbf{스팟/선점형(preemptible)}(60–90\% 할인)은 말 그대로 중단 가능 산업 요금이다. 클라우드는 피크에 GCP 30초, AWS 2분 전 통보 후 회수할 수 있으며, 할인은 이를 견디도록 엔지니어링한 대가다. 사례 12-A의 계단에서 포트폴리오 규칙이 자연스럽게 나온다. \textbf{바닥은 약정하고, 어깨는 온디맨드로, 버스트는 스팟으로 처리한다. 바닥을 스팟에 두지 않는다.} 03:00의 불면증 사용자에게는 약정한 실리콘이 필요하다. Insull의 기법 하나를 더하면 구성이 완성된다. \textbf{비혼잡 시간 채우기}다. 그의 야간 전력은 제빙소로 갔고 야간 GPU는 제7장의 배치 모드로 간다. 재임베딩 작업, 평가 스윕, 보정 실행처럼 설계상 중단 가능한 작업이 스팟에 가장 잘 맞는다. SLO를 건드리지 않고 부하율을 높인다.

\begin{tbbtablewrap}
  \begin{tabular}{@{}>{\raggedright\arraybackslash}m{\dimexpr 0.2181\tbltextwidth-2\tabcolsep\relax} >{\raggedright\arraybackslash}m{\dimexpr 0.1862\tbltextwidth-2\tabcolsep\relax} >{\raggedright\arraybackslash}m{\dimexpr 0.2766\tbltextwidth-2\tabcolsep\relax} >{\raggedright\arraybackslash}m{\dimexpr 0.3191\tbltextwidth-2\tabcolsep\relax}@{}}
  \arrayrulecolor{tbbnavy}\specialrule{1pt}{0pt}{0pt}
  \rowcolor{tbbtablehead}{\centering \bfseries\color{tbbnavy}계층(전력망의 조상)} & {\centering \bfseries\color{tbbnavy}정가 대비 가격} & {\centering \bfseries\color{tbbnavy}거래 조건} & {\centering \bfseries\color{tbbnavy}교재 플릿에서의 용도} \\
  \arrayrulecolor{tbbnavy}\specialrule{0.7pt}{0pt}{2pt}
  {\textbf{약정(기본 부하)}} & {−30–60\%} & {1–3 yrs의 유휴 위험을 구매자가 부담} & {레플리카 2개의 바닥: 2 × \$595 × 0.60 = \$714/mo} \\
  \arrayrulecolor{tbbtableborder}\hline
  {\textbf{온디맨드(부하 추종)}} & {정가} & {약속 없음. 프리미엄 = 유연성} & {오토스케일러가 추가하고 제거하는 어깨} \\
  \arrayrulecolor{tbbtableborder}\hline
  {\textbf{스팟(중단 가능 요금)}} & {−60–90\%} & {30 s–2 min 전 통보. 드레인을 설계해야 함} & {피크 버스트: 13 rep-h/day × \$0.28 ≈ \$113/mo} \\
  \arrayrulecolor{tbbtableborder}\hline
  {\textbf{비혼잡 시간 채우기(제빙소)}} & {야간 스팟} & {골짜기의 중단 가능한 작업} & {03:00의 배치 재임베딩, 평가, 보정} \\
  \arrayrulecolor{tbbnavy}\specialrule{1pt}{2pt}{0pt}
  \end{tabular}\par
  \figcaption{표 12-2  변환한 요금표. 같은 GPU 시간에 세 가격이 붙으며 유휴 위험을 누가 지고 중단을 누가 견디는지가 다르다.}
\end{tbbtablewrap}

\subsection*{눈물 없이 스팟에서 서빙하기 — 드레인 훈련}

배치뿐 아니라 \textit{서빙}에 스팟을 쓰려면 자격을 얻어야 하며, 그 자격은 이미 가진 조각을 안무하는 데서 나온다. 메타데이터 엔드포인트나 노드 테인트로 중단 통보가 도착한다. 포드의 \textbf{preStop hook + terminationGracePeriodSeconds}가 준비 상태를 바꿔 새 입장을 막고 진행 중 스트림이 끝나게 한 뒤 LB에서 등록을 해제한다. 제5장의 우아한 종료 메커니즘이 마침내 시스템을 지탱한다. 제11장에서 측정한 가장 긴 정직한 스트림 시간이 바로 유예 기간이다. \textbf{상태 비저장은 마지막 배당금을 지급한다}(제6/11장). 끝내지 못한 스트림은 다른 레플리카에서 다시 프리필한다. 짧은 지연 시간 증가일 뿐 잘못된 답은 나오지 않는다. 롤아웃과 접두사 캐시(prefix cache) 미스를 안전하게 만든 같은 성질이다. 스팟 풀은 영역과 인스턴스 유형에 걸쳐 다양화하고 용량 최적화 할당을 사용한다. 폴백도 설정한다. 스팟이 회수되고 대체되지 않으면 오토스케일러가 온디맨드로 격차를 채운다. 한 시간 동안 할인을 잃을 뿐 SLO를 잃지는 않는다. 중단이 지루하게 느껴질 때까지 스테이징에서 훈련한다. \textbf{절대 규칙은 스팟이 바닥이 아니라 버스트 계층이라는 것}이다. 누군가는 언제나 시도하므로 점검표에 위험 신호로 넣었다.

\subsection*{원장을 마감하다 — 제1장의 손익분기점도 함께}

\begin{terminalbox}
\begin{Verbatim}[breaklines=true,breakanywhere=true,fontsize=\small,formatcom=\color{tbbtermfg},codes=\tbbcodeactive]
$ python fleet_ledger.py --final          # the course's bill, walked
 Wk 7  estimated, 8 x L4, 24/7             $4,762   # the syllabus bill
 Wk 10 INT8: 2x goodput -> 4 replicas      $2,381   # same SLO
 Wk 12 autoscale: 61 rep-h/day             $1,513   # same SLO
 Wk 12 portfolio:
   commit floor  2 x $595.20 x 0.60      $  714.24
   spot burst    403 rep-h x $0.28       $  112.84
   pre-warm + cron cushion               $   49.60
   TOTAL                                 $  876.68   < $900 ceiling

 unit economics (measured, Wk 11 sweeps):
   capacity  ~3.1B tok/mo   served ~2.0B tok/mo  (load factor again!)
   $/M tokens: 876.68 / 2000 = $0.44   vs API $0.50: self-host wins
   ...at THIS utilization. at half the traffic, the API wins. Ch. 1's
   break-even curve: closed, with measured numbers, 11 weeks later.
\end{Verbatim}
\end{terminalbox}
\figcaption{기록 12-6  결승선. \$4,762 → \$2,381(최적화) → \$1,513(오토스케일링) → \$877(포트폴리오)로 줄어 1주 차 상한보다 \$23 낮아졌다. 모든 단계가 측정한 수단이다. \$/M-token 항목은 교재 첫머리의 자체 구축 대 외부 도입 질문을 마무리한다. 답은 구호가 아니라 사용률이었다.}

세 가지 습관이 일회성 승리를 영웅적인 오후가 아닌 지속적 규율인 \textbf{FinOps}로 바꾼다. \textbf{총액이 아니라 단위를 계량한다.} 중요한 청구서는 \$/M tokens 또는 \$/K requests다. 주, 모델, API 대안과 비교할 수 있기 때문이다. 총지출은 제품이 성장했다는 사실만 알려 준다. \textbf{태그를 붙이고 비용을 공개한다.} 모든 자원에 team/model/env 레이블을 붙여 원장을 분해한다. 제6장의 NAT 함정(\$5k/yr, 조용함)은 태그 없는 배관에 숨은 지출의 대표 사례다. \textbf{이상에 경보를 건다.} 예측의 80\%/100\%/120\%에서 예산 경보를 울린다. 제2장 실습의 가드레일(guardrail)이 프로덕션으로 승격되어 누수 GPU, 잊힌 개발 플릿, 트래픽으로 청구되는 재시도 폭풍을 잡는다. 화려한 일은 아니지만 \$877이 그대로 유지되는 것과 유령 GPU가 조용히 돌아오는 것의 차이다. 한 세기 뒤에도 Insull의 진짜 제품은 같다. \textbf{감시하는 용량}이다.

\begin{calloutbox}{tbbred}{tbbxFEF2F2}
{\normalsize\bfseries\color{tbbx991B1B}상자 12-1  비용 검토의 위험 신호 — 공개 점검표}\par\vspace{3pt}
\begin{tbbitemize}
  \item \textbf{단위 비용 없음} — \$/M tokens 없이 총지출만 있으면 비용 검토가 아니라 성장 차트다.
  \item \textbf{피크에 맞춘 약정} — 유휴 시간을 다시 산 셈이다. 바닥만 약정하고 나머지는 포트폴리오로 구성하라.
  \item \textbf{바닥에 스팟 사용} — 03:00에 호출하는 할인이다. 버스트 계층에만 쓰고 드레인 훈련을 반복한다.
  \item \textbf{트래픽 그래프의 평평한 레플리카 선} — 사례 12-A는 한눈에 진단할 수 있다. 부하율을 대시보드에 넣어야 한다.
  \item \textbf{태그 없는 지출 / 조용한 배관} — NAT 함정, 잊힌 개발 플릿, 좀비 볼륨이다. 귀속할 수 없으면 검토할 수도 없다.
  \item \textbf{SLO 인용 없이 주장하는 절감} — 줄인 모든 비용은 깨뜨리지 않은 조항을 밝혀야 한다. 이 장은 \textit{같은} SLO에서 비용을 절반으로 줄였으며 그 문장이 규율의 전부다.
\end{tbbitemize}
\end{calloutbox}

\vspace{3.0pt}

\begin{calloutbox}{tbbpurple}{tbbpurplefill}
{\normalsize\bfseries\color{tbbpurple}생각해 보기}\par\vspace{3pt}
\begin{tbbitemize}
  \item 트래픽이 두 배이고(피크 동시 스트림 208개, 같은 형태) 기록 12-6을 다시 작성하라. 계단, rep-h/day, 포트폴리오 분할, 최종 청구서, \$/M tokens를 계산하라. 어느 항목이 개선되고 어느 항목이 나빠지는가? 제품이 성장할수록 단위 비용이 \textit{낮아지는} 이유도 설명하라. Insull이라면 알고 있을 것이다.
  \item 약정은 1년간 레플리카 2개를 포함한다. 4개월째에 INT4 빌드(§11.5)가 도입되어 레플리카마다 스트림을 1.6× 더 처리한다. 과도하게 약정했는가? 성장해 채우기, 바닥 워크로드 이동, 재판매/수정 선택지를 검토하고 최적화 로드맵 아래에서 약정할 계획 규칙을 도출하라.
  \item KServe 카나리 필드는 트래픽을 90/10으로 나눈다. 하지만 §10.3의 게이트는 INT4 모델이 안전하다고 이미 결정했다. 카나리는 게이트가 확인할 수 없는 무엇을 검사하는가? 실제 트래픽 분포, 실제 인프라, 실제 SLO를 고려하라. 제13장의 첫 논거를 세 문장으로 작성해 보라.
  \item 실습 원장에서 자신의 프로젝트 손익분기 API 가격을 계산하라. 어느 \$/M tokens에서 자체 호스팅의 우위가 사라지며, 사용률, GPU 가격, 레플리카당 토큰 중 어느 단일 입력이 답을 가장 크게 움직이는가? "왜 그냥 API를 쓰지 않는가?"에 답하는 한 장짜리 슬라이드로 제시하라.
\end{tbbitemize}
\end{calloutbox}

\section{실습: 프로젝트 중간 점검 — 숨 쉬게 만든 뒤 원장 발표하기}

강의 계획서는 이번 주 제출물을 \textbf{팀 프로젝트 중간 점검}이라고 분명히 부른다. 따라서 실습은 두 부분으로 나뉜다. A부는 약 90분의 실습으로 눈앞에서 플릿이 숨 쉬게 하고 반사 작용의 가격을 계산한다. B부는 7주 차 이후 모든 결과를 팀이 발표할 8분짜리 패키지로 조립한다. 우연이 아니라 스태프 엔지니어가 용량 검토에 가져갈 것과 같은 패키지다. GPU 시간 예산은 \textasciitilde{}\$3이고, 이 장에서 마침내 내부 의미를 설명한 결제 내역 스크린샷도 제출한다.

\begin{tbbfigure}
  \adjustbox{max width=0.94\linewidth,max totalheight=0.46\textheight}{\includegraphics{figures/fig12-7_midpoint.png}}\\[3pt]
  \figcaption{그림 12-7  실습의 두 부분. A에서는 스케일러를 연결하고 파도를 보내며 탄생 시간을 재고 T4를 분할해 네 가지를 측정한다. B는 모든 수치의 출처를 추적할 수 있는 산출물 여섯 개로 구성한 8분짜리 중간 점검 패키지다.}
\end{tbbfigure}

\subsection*{A부 — 네 가지 측정(①–⑤)}

\begin{tbbitemize}
  \item \textbf{① KEDA를 대기열에 연결한다.} 실습 자료의 데모 환경을 배포한다. GPU 네 대 없이도 스케일링이 보이도록 CPU 모드 레플리카를 사용한다. Prometheus + KEDA + 기록 12-1의 ScaledObject를 \code{vllm:num\_requests\_waiting}에 연결하고 최소 1 / 최대 4로 설정한다. KEDA가 생성한 HPA 객체를 확인한다. 한 단계 위에서 관리자를 관리하는 제5장의 패턴이다.
  \item \textbf{② 09:00 파도를 보낸다.} 제11장의 토큰 인식 하네스로 동시성을 높이고 세 터미널(\code{kubectl get hpa -w}, \code{get pods -w}, 하네스의 TTFT 로그)에서 \code{waiting → desired → pods}를 지켜본다. 사건 전체의 TTFT p99 그래프를 그리고 감지, 탄생, 드레인을 표시한다. 이어 08:40 cron 트리거를 추가하고 같은 파도를 보낸다. GPU 비용 \textasciitilde{}\$0로 그래프 두 개와 교훈 하나를 얻는다.
  \item \textbf{③ 실제 탄생 하나의 시간을 잰다.} GPU 노드에서 vLLM 포드를 삭제하고 \code{kubectl describe pod} + 엔진 로그를 사용해 기록 12-3의 항목을 초시계로 잰다. 각 줄을 담당 장과 ±20\% 안에서 맞추고 플릿의 지배 항목과 이를 공격할 사다리 단을 제시한다. 이 수치는 \textit{자신의} 반응 시간 예산으로 중간 점검 패키지에 들어간다.
  \item \textbf{④ T4를 분할하고 대가를 측정한다.} 시분할 ConfigMap(replicas: 4)을 적용하고 분류기 포드 두 개를 함께 스케줄링한 뒤, 9주 차 스윕을 단독 대 이웃 있음으로 다시 실행한다(기록 12-4). 두 경우의 p50, p99, req/s@SLO를 보고하고 SLO 아래에서 시분할에 관한 한 문장 결론을 작성한다.
  \item \textbf{⑤ 원장을 마감한다.} 기록 12-6의 빈칸을 \textit{자신이} 측정한 수치로 채운다. 실제 하루의 지표나 실습 자료 재생에서 얻은 계단, rep-h/day, 실시간 가격표에서 정한 포트폴리오 분할, API 대안과 비교한 \$/M tokens를 넣는다. 모든 빈칸의 출처를 측정값이나 가격표로 추적할 수 있어야 한다. 채점자가 가장 먼저 확인하는 항목이다.
\end{tbbitemize}

\subsection*{B부 — 중간 점검 패키지(⑥–⑦)}

산출물 여섯 개를 조립한 뒤 8분에 맞춰 연습한다. \textbf{근거가 있는 SLO 문장}은 7주 차 계약을 11주 차 조항으로 다시 쓰고 각 조항 옆에 이를 인증하는 측정값을 둔다. \textbf{최적화 이야기}에는 10주 차 전후 표, 산출물 옆의 게이트 보고서, 한 문단으로 옹호한 INT4/INT8 결정을 넣는다. \textbf{스케일링 정책}은 ScaledObject, ③의 반응 시간 예산, cron/사전 워밍업 계획, 파도를 견딘 모습을 보여 주는 ②의 그래프 하나로 구성한다. \textbf{원장}은 자신의 수치를 넣은 그림 12-6의 폭포이며 상한 아래에 도달하거나 격차와 계획을 정직하게 설명한다. \textbf{위험과 요청 사항}에는 완화책이 있는 구체적인 위험 세 가지와 13–15주 차 전에 필요한 것을 넣는다. "안전 슬라이스에서 AWQ 게이트 여유가 0.1 pt다"가 "양자화가 정확도를 해칠 수 있다"보다 낫다. \textbf{데모}에서는 스트리밍 요청 하나와 실시간 스케일 아웃 하나 또는 녹화를 보여 준다. 발표하고 피드백을 수집한다. 15주 차의 드레스 리허설이다.

\begin{calloutbox}{tbbteal}{tbbxF0FDFA}
{\normalsize\bfseries\color{tbbx115E59}상자 12-2  중간 점검 평가 기준 — 8분을 채점하는 방법}\par\vspace{3pt}
\begin{tbbitemize}
  \item \textbf{추적 가능성(40\%)} — 모든 수치가 스윕, 기록, 가격표를 인용해야 한다. 출처 없는 "약 500"은 이 교재에서 유일하게 실패하는 답이다.
  \item \textbf{두 시계 제시(20\%)} — 백분위수와 부하를 명시한 TTFT와 TPOT를 모두 보여 준다. "지연 시간" 하나만 쓰면 해당 행 점수를 잃는다. 상자 11-2의 규칙은 계속 살아 있다.
  \item \textbf{원장 마감(20\%)} — 고정 → 오토스케일링 → 포트폴리오 → 단위 비용의 계산을 실시간으로 확인할 수 있어야 한다. 상한 아래라고 추가 점수를 주지 않고, \textit{설명하지 않은} 격차는 감점한다.
  \item \textbf{위험을 엔지니어링으로 표현(10\%)} — 구체적이고 측정했으며 완화해야 한다. 하나 이상은 약정 대 로드맵, 스팟 가용성, 이그레스 같은 비용 위험이어야 한다.
  \item \textbf{시간 규율(10\%)} — 8분 발표 후 질문을 받는다. 말을 멈추지 못하는 스태프 엔지니어는 용량 검토 일정이 다시 잡히며 학생도 마찬가지다.
\end{tbbitemize}
\end{calloutbox}

\vspace{3.0pt}

제출물은 A부의 측정 산출물 네 개(파도 그래프 두 개, 항목별 탄생, 이웃 세금 표), 완성한 원장, 슬라이드나 문서 형태의 산출물 여섯 개 패키지, 그리고 이 교재에서 기술 실습으로는 마지막인 결제 내역 스크린샷이다. 이 장을 마치면 스크린샷은 더 이상 의례가 아니라 \textbf{읽을 수 있는 명세}가 된다. 모든 항목에 담당 장, 수단, 소유자가 생겼다. 처음부터 바로 그것이 목적이었다.

\begin{calloutbox}{tbbpurple}{tbbpurplefill}
{\normalsize\bfseries\color{tbbpurple}생각해 보기}\par\vspace{3pt}
\begin{tbbitemize}
  \item ②의 그래프에서 cron 실행은 +\$1.07/day의 비용으로 TTFT p99를 0.46 s에 평평하게 유지하지만 반응형 실행은 추가 비용 없이 3 min 동안 위반한다. 자신의 제품을 위한 두 문장 권장안을 작성하고, 결정을 좌우하는 인프라 사실이 아닌 제품 사실을 제시하라.
  \item ③에서 탄생은 150초가 아니라 190 s였다. 추가 40 s는 노드 캐시가 없는 새 노드의 첫 포드가 이미지를 풀하는 시간이다. 어느 사다리 단으로 해결할 수 있고 각각의 월 비용은 얼마인가? 플릿이 주 2회 노드를 교체한다면 실제로 무엇을 권장할지 설명하라.
  \item 중간 점검 검토자가 "원장은 일중 곡선이 유지된다고 가정한다. B2B 고객이 02:00에 50k requests의 야간 배치를 추가하면 무엇이 깨지는가?"라고 묻는다. 계단과 포트폴리오에서 어느 계층이 처리하는지, 부하율에 미치는 영향으로 답하라. 이것이 단위 비용에는 좋은 소식인 이유도 설명하라.
  \item 팀의 위험 세 줄을 초안으로 지금 작성하라. 9/11주 차 측정에서 SLO 위험 하나, 이 장의 원장에서 비용 위험 하나, 13–15주 차의 전달 위험 하나를 택한다. 평가 기준의 구체성 잣대를 자신의 문장에 적용해 한 번 수정하라.
\end{tbbitemize}
\end{calloutbox}

\namedsection{요약}{열여섯 가지 핵심 내용}

\qitem{1}{\textbf{낭비는 잘못된 설정이 아니라 구조에서 나온다.} 수요는 시간에 따라 모양이 있지만 정적 플릿은 수평선이다. 측정한 사례 12-A에서는 하루 61레플리카 시간이 필요하지만 96시간에 비용을 냈다. 부하율은 63.5\%, 야간 근무 비용은 \$868/month다. 사례 1-B의 유령 GPU를 일중 곡선으로 확장한 모습이다.}

\qitem{2}{\textbf{1902년의 문제를 1910년까지 해결했다.} Samuel Insull은 부하율(평균 ÷ 피크)로 시카고 전력망을 운영했다. 수요를 계량하고 고객 피크를 맞물리게 하며 골짜기를 저렴하게 팔고 피크에는 중단 가능한 부하를 끊었다. 이 장의 오토스케일링, 공유, 스팟, 야간 배치는 모두 그의 기법이다.}

\qitem{3}{\textbf{오토스케일러는 기상학자가 아니라 온도 조절기다.} 15 s마다 desired = ceil(current × metric ÷ target)을 계산해 제5장의 메커니즘이 따르는 정수 하나에 쓴다. 엔지니어링은 공식이 아니라 지표와 초시계에 있다.}

\qitem{4}{\textbf{공급 게이지가 아니라 수요 신호로 스케일링한다.} GPU가 빠지는 동안 CPU는 9\%라 CPU\%는 GPU 서비스에서 작동하지 않는다. 제9장 사례의 거울상이다. GPU 사용률은 포화되고 부분 부하에서 거짓말한다. 올바른 눈은 선행하고 상한이 없으며 사용자가 느끼는 입장 대기열(vllm:num\_requests\_waiting)이다(L = λW).}

\qitem{5}{\textbf{루프의 성향은 비대칭이어야 한다.} 순간 신호로 증가시키고 후행 창에서만 감소시킨다(안정화/쿨다운 \textasciitilde{}300 s). 손실 함수가 비대칭이기 때문이다. 빠른 증가는 레플리카 몇 분을 낭비하지만 빠른 감소는 SLO를 깨뜨린다. 플래핑은 청구서를 제10장의 torch.compile 악순환 같은 콜드 스타트 오버헤드로 바꾼다.}

\qitem{6}{\textbf{반응 시간도 용량의 일부다.} T = 감지(≤15 s) + 결정(≤15 s) + 탄생(따뜻한 노드 \textasciitilde{}150 s)이다. 대기열은 T 전체에 걸쳐 초과분을 적분한다. 반응형 경로에서 09:00 파도는 3분 동안 TTFT를 위반하지만, 08:40 cron은 \$1.07로 없앤다. 볼 수 있는 피크에는 예지력보다 달력이 낫다.}

\qitem{7}{\textbf{포드는 파도를 쫓고 노드는 예측을 따른다.} GPU 노드 탄생은 3–7 min이고 때로는 "no capacity"로 끝난다. SLO 경로에 두지 않는다. 하루의 노드 범위를 미리 프로비저닝하고 그 안에서 포드를 오토스케일링한다.}

\qitem{8}{\textbf{콜드 스타트는 항목별 청구서다.} 스케줄링 2 s + 이미지 16 s(제4장) + 가중치 9 s(제6장) + 엔진 초기화 95 s(제11장) + 워밍업 28 s(제5/10장) = 150 s다. Firecracker의 125 ms는 환경 항목을 해결했고 페이로드 항목은 사용자의 몫이다. 서버리스 GPU는 프리미엄을 내고 따뜻한 풀 + 초기화한 프로세스의 스냅샷/복원을 사용하는 것이다.}

\qitem{9}{\textbf{완화 사다리는 이미 배운 장으로 이루어진다.} 작고 캐시되며 지연 로드하는 이미지 → 병렬 가중치 가져오기 + 노드 캐시 → 준비 상태 안의 워밍업 → 여유 용량 +1 (\textasciitilde{}\$50/mo) → 예약 스케일링 → 0으로 스케일링 순이다. 0은 호출자가 첫 적중의 몇 분 지연을 감당하는 엔드포인트용이다. 롱테일은 잠들고 SLO를 지는 헤드는 롤아웃과 03:00 사용자를 위해 바닥 2개를 유지한다.}

\qitem{10}{\textbf{공유는 카드당 테넌트 하나 규칙을 세 방식으로 굽힌다.} 시분할은 전체 다이를 번갈아 쓰며 격리가 없어 개발용으로만 적합하다. 실습에서 이웃 하나로 p99가 9.8 → 38.4 ms가 되었다. MPS는 동시에 실행하지만 장애 반경이 하나라 한 팀의 작은 모델에 적합하다. MIG는 큰 다이를 최대 7개로 하드웨어 분할해 자체 SM/L2/DRAM/OOM을 제공하므로 SLO가 있는 멀티테넌트에 안전하지만 T4/L4에는 없다.}

\qitem{11}{\textbf{통합 계산은 Insull의 다양성 기법이다.} A100의 1g.10gb 조각(≈\$0.51/hr)은 조각 여섯 개를 남기면서 T4 전체(\$0.53)와 비슷하다. 단, 테넌트 피크가 맞물리고 조각이 자체 §10.2 진단을 통과해야 한다. 대역폭도 1/7이다. 상관된 피크는 한 계층 아래에서 13:00 문제를 다시 만든다.}

\qitem{12}{\textbf{KServe는 제9장이 이름 붙인 플랫폼 종이다.} 15줄 InferenceService 하나가 Deployment, 프로브, 모델 가져오기, 런타임, 인그레스(ingress), 동시성 기반 오토스케일러, 0으로 스케일링, 카나리 트래픽 분할을 만든다. 모델과 팀이 늘면 도입하고 인기 모델 하나가 제품이면 원시 Deployment + KEDA를 유지한다.}

\qitem{13}{\textbf{용량을 포트폴리오로 구매한다.} 바닥은 약정(−30–60\%), 어깨는 온디맨드, 버스트는 스팟(−60–90\%, 30 s–2 min 통보), 골짜기는 배치로 채운다. 바닥을 스팟에 두지 않는다. 드레인 훈련은 제5장의 preStop + 제11장의 상태 비저장이다. 스트림을 끝내고 등록을 해제한 뒤 다른 곳에서 다시 프리필한다. 짧은 지연일 뿐 잘못된 답은 아니다.}

\qitem{14}{\textbf{원장을 상한 아래에서 마감했다.} \$4,762(7주 차 추정) → \$2,381(10주 차 INT8) → \$1,513(오토스케일링, 61 rep-h/day) → \textbf{\$877}(= \$714 약정 + \$113 스팟 + \$50 사전 워밍업) < \$900이다. 모든 단계가 \textit{같은} SLO에서 측정한 수단이다.}

\qitem{15}{\textbf{단위 경제성이 제1장의 손익분기점을 마무리한다.} \$877 ÷ 서빙한 2.0B tokens = \$0.44/M으로 API의 \$0.50보다 낮다. 자체 호스팅은 \textit{이} 사용률에서 이기지만 절반이면 진다. 자체 구축 대 외부 도입은 구호가 아니라 부하율 문제였다.}

\qitem{16}{\textbf{FinOps는 감시하는 용량이다.} 총액이 아니라 \$/M tokens를 계량하고, 태그와 비용 공개로 NAT 함정이 태그 없는 배관에 숨지 않게 하며, 예측의 80/100/120\%에서 이상 경보를 건다. 중간 점검 패키지는 모든 수치의 출처를 추적할 수 있게 8분에 이를 발표한다. 15주 차의 드레스 리허설이다.}

\namedsection{연습문제}{}

\subsection*{A. 객관식}

\qitem{1}{GPU 채팅 서비스가 대기열 증가와 붉은 TTFT로 고전하지만 CPU 기반 HPA는 전혀 스케일링하지 않는다. 근본 원인은 무엇인가?}

\begin{tbboptions}
  \item ① HPA 루프가 너무 느리다        ② 경합 자원인 GPU/대기열이 관찰 지표에 보이지 않는다. CPU는 9\%로 쉰다
  \item ③ maxReplicas가 너무 낮다        ④ 이미지가 너무 크다
\end{tbboptions}

\qitem{2}{vLLM 플릿의 최적 오토스케일링 지표가 num\_requests\_waiting인 이유는 무엇인가?}

\begin{tbboptions}
  \item ① 수집 비용이 가장 저렴하다        ② 지연 시간이 움직이기 전 작동하는 선행 지표이고, 용량 오차의 크기를 나타내는 상한 없는 값이며, TTFT의 분자인 대기열로 사용자가 느낀다
  \item ③ Kubernetes에 내장되어 있다        ④ 정상 상태에서 언제나 0이다
\end{tbboptions}

\qitem{3}{스케일 업 창은 즉시 반응하지만 스케일 다운 창은 \textasciitilde{}5분인 이유는 무엇인가?}

\begin{tbboptions}
  \item ① 스케일 다운이 기계적으로 더 느리다        ② 손실 함수가 비대칭이다. 빠른 증가는 레플리카 몇 분을 낭비하지만 빠른 감소는 SLO를 깨뜨리고 콜드 스타트 비용을 다시 내며 플래핑한다
  \item ③ Kubernetes가 요구한다        ④ 5분 단위로 요금을 부과한다
\end{tbboptions}

\qitem{4}{레플리카 2개, 전체 대기 41개, 레플리카당 목표 10일 때 HPA가 계산하는 desired는 얼마인가?}

\begin{tbboptions}
  \item ① 3        ② 4        ③ ceil(2 × 20.5 ÷ 10) = 5 (이후 maxReplicas로 제한)        ④ 41 ÷ 10 = 4.1 → 4
\end{tbboptions}

\qitem{5}{Firecracker의 125 ms microVM이 GPU 콜드 스타트를 해결하지 못한 이유는 무엇인가?}

\begin{tbboptions}
  \item ① CPU에서만 실행된다        ② 청구서의 대부분이 환경이 아니라 페이로드다. 수 기가바이트의 이미지와 가중치, CUDA 컨텍스트, 엔진 초기화가 150초 중 148초를 차지한다
  \item ③ Kubernetes는 microVM을 사용할 수 없다        ④ GPU는 부팅이 느리다
\end{tbboptions}

\qitem{6}{0으로 스케일링이 적합한 대상은 무엇인가?}

\begin{tbboptions}
  \item ① 유휴 시간이 있는 모든 서비스        ② 야간의 채팅 플릿
  \item ③ 호출자가 첫 적중의 몇 분 지연을 감당할 수 있는 엔드포인트. 롱테일 내부 모델, 개발 엔드포인트, 배치 트리거 스코어러        ④ GPU에서는 아무것도 해당하지 않는다
\end{tbboptions}

\qitem{7}{한 팀이 SLO를 지는 서비스 두 개를 시분할 GPU 한 대에 둔다. 예상되는 결과는 무엇인가?}

\begin{tbboptions}
  \item ① 처리량 2×        ② 아무것도 변하지 않는다
  \item ③ p50은 거의 유지되지만 운 나쁜 요청이 이웃의 시간 조각을 먹어 p99가 몇 배가 된다. 테넌트 하나의 OOM이 둘 다 죽인다        ④ MIG가 자동 활성화된다
\end{tbboptions}

\qitem{8}{일중 계단의 포트폴리오 규칙은 무엇인가?}

\begin{tbboptions}
  \item ① 할인을 위해 모두 스팟에 둔다        ② 할인을 위해 모두 약정한다
  \item ③ 바닥은 약정하고 어깨는 온디맨드, 버스트는 스팟으로 처리하며 바닥을 스팟에 두지 않는다        ④ 온디맨드만 쓴다. 가격 전략은 엔지니어링이 아니다
\end{tbboptions}

\subsection*{B. 단답형}

\qitem{1}{GPU 플릿의 부하율을 정의하고 사례 12-A(필요 61, 결제 96)에서 계산한 뒤, 이 지표로 전력 회사를 운영한 1900년대 인물을 제시하라.}

\qitem{2}{목표 레플리카 공식을 쓰고 다음 조건에서 계산하라. 레플리카 3개, 전체 대기 66개, 레플리카당 목표 15, maxReplicas 6이다.}

\qitem{3}{따뜻한 노드의 150초 콜드 스타트를 담당 장과 함께 다섯 항목으로 세분화하고, vLLM 플릿에서 지배적인 항목을 밝히라.}

\qitem{4}{시분할, MPS, MIG를 각각 무엇을 격리하거나 격리하지 않는지 중심으로 한 문장씩 설명하라.}

\qitem{5}{스팟 드레인 훈련의 네 박자를 서술하고, 중단된 스트림이 잘못된 답이 아니라 짧은 지연만 겪게 하는 제11장의 성질을 제시하라.}

\qitem{6}{중간 점검 원장은 왜 월 지출 대신 \$/M tokens를 보고하는가? 비교 가능성과 제1장의 어느 질문에 답하는지를 모두 설명하라.}

\subsection*{C. 서술형}

\qitem{1}{용량 검토에 관한 글을 작성하라. 서비스는 평일 11:00–15:00에 동시 스트림 340개로 피크를 찍고, 어깨는 120개, 야간과 주말 바닥은 25개다. 레플리카당 좌석 32개, \$1.10/hr, 콜드 스타트 3 min이다. 계단, 반응 계획(cron + 대기열 트리거 + 여유 용량), 포트폴리오(1년 약정 −38\%, 스팟 −64\%를 사용하는 약정/온디맨드/스팟), 최종 월 청구서를 구성하라. Insull 방식으로 전후 부하율도 제시하라.}

\qitem{2}{"오토스케일링은 이 장에서 가장 덜 흥미로운 수단이다. 실제 비용 절감은 제10장과 제11장에서 일어났고 실제 위험은 여기 있다"라는 주장을 옹호하거나 공격하라. 교재 원장(\$4,762 → \$2,381 → \$1,513 → \$877), 각 수단이 만드는 SLO 위반 유형, 각 단계의 엔지니어링 비용을 활용하라. 이 프로젝트에서 엔지니어링 비용의 다음 한계 달러를 어디에 써야 하는가?}

\qitem{3}{통합에 관한 글을 작성하라. 한 부서가 작은 모델 아홉 개를 각각 T4 한 대(\textasciitilde{}\$0.53/hr, 사용 주기 5–40\%, 팀 세 개, SLO를 지는 모델 두 개)에서 실행한다. 시분할/MPS/MIG/전체 카드로 목표 플릿을 설계하라. 모델 클래스별 배치, 각 선택의 격리 논거, 장애 반경과 재구성 드레인 위험, 월 전후 비용을 포함하라. 승인 전에 요구할 측정 하나도 제시하라. 힌트는 Insull의 요구다.}

\qitem{4}{손익분기점에 관한 글을 작성하라. 최종 원장(\$877, 서빙한 2.0B tokens, 용량 3.1B)을 사용해 월 토큰의 함수로 자체 호스팅 대 API 결정을 도출한다. \$0.50/M API와의 교차점을 계산하고 INT4 (1.6×)와 트래픽 두 배가 각각 이를 어떻게 움직이는지 보인 뒤, 0.4B tokens/month인 스타트업에 보내는 한 문단 메모로 마무리하라.}

\namedsection{정답과 해설}{}

\subsection*{A. 객관식}

\qitem{1}{\textbf{②} 오토스케일러의 눈이 경합하지 않는 자원을 본다. 제9장 사례 A의 거울상이다. 그때는 CPU가 고정되고 GPU가 쉬었지만 여기서는 GPU가 고전하는 동안 CPU가 쉰다. 수요 신호로 스케일링해야 한다. (§12.2)}

\qitem{2}{\textbf{②} 한 지표가 세 속성을 모두 갖고 있으며 제11장에서 읽는 법을 배우고 기록 11-3에서 관찰한 바로 그 항목이다. 공급 게이지인 CPU\%, GPU\%는 세 속성 중 하나 이상에서 실패한다. (§12.2)}

\qitem{3}{\textbf{②} 비대칭은 레플리카 몇 분과 SLO 위반 + 다시 내는 탄생 비용을 비교하는 손실 함수를 인코딩한다. 플래핑의 최악 사례는 제10장의 스케일 아웃 시 컴파일 생각해 볼 문제다. (§12.2)}

\qitem{4}{\textbf{③} 레플리카당 지표는 20.5다. desired = ceil(2 × 20.5/10) = ceil(4.1) = 5를 계산한 뒤 상한을 적용한다. ④는 고전적 오류다. 공식은 현재 수를 스케일링하며 원시 합계를 나누는 것이 아니다. (§12.2)}

\qitem{5}{\textbf{②} 기록 12-3에서 환경은 150초 중 \textasciitilde{}2 s만 차지한다. Firecracker는 2005년 문제를 훌륭하게 해결했지만 2026년 문제는 수 기가바이트와 CUDA 컨텍스트다. 더 빠른 부팅이 아니라 따뜻한 풀과 스냅샷이 필요하다. (§12.3)}

\qitem{6}{\textbf{③} 결정 규칙은 유휴 시간이 아니라 호출자의 지연 시간 허용 범위다. (a)는 이를 충족하지 못한다. 03:00 사용자도 실제이고 바닥은 롤아웃에도 필요하다. (§12.3)}

\qitem{7}{\textbf{③} 기록 12-4의 측정값은 p50 3.1 ms, p99 38.4 ms다. 꼬리 요청이 컨텍스트 전환과 이웃의 배치를 먹으며 VRAM은 공유한다. SLO 아래의 시분할은 제14장의 전쟁 이야기를 미리 구매하는 셈이다. (§12.4)}

\qitem{8}{\textbf{③} 계단은 계층별로 요금표에 대응한다. 바닥을 스팟에 두는 것은 점검표에서 가장 붉은 위험 신호다. 03:00에 회수되면 대체 용량이 없다. (§12.5)}

\subsection*{B. 단답형}

\qitem{1}{부하율 = 필요한(또는 평균) 용량 시간 ÷ 결제한(또는 피크) 용량 시간이다. 61/96 = \textbf{63.5\%}다. Chicago Edison의 Samuel Insull은 수요 곡선으로 전력 가격을 정하고 골짜기를 채워 바로 이 수치를 높였다. (§12.1)}

\qitem{2}{desired = ceil(current × metric ÷ target) = ceil(3 × 22 ÷ 15) = ceil(4.4) = \textbf{5}다. 상한 6 이하여서 5가 유지된다. 레플리카당 지표는 66/3 = 22다. (§12.2)}

\qitem{3}{스케줄링 2 s(제5장) + 이미지 풀 16 s(제4장) + 가중치 9 s(제6장) + \textbf{엔진 초기화 95 s(제11장 — VRAM 로드, KV 풀, CUDA 그래프. 지배 항목)} + 프로브 워밍업 28 s(제5/10장) = 150 s다. (§12.3)}

\qitem{4}{시분할은 테넌트가 전체 다이를 번갈아 써 아무것도 격리하지 않는다. MPS는 SM 상한을 두고 동시에 실행해 연산 몫은 격리하지만 대역폭과 장애는 격리하지 않는다. MIG는 자체 SM/L2/DRAM/OOM을 지닌 하드웨어 파티션으로 거의 모든 것을 격리하지만 큰 다이에서만 1/7 세분성을 제공한다. (§12.4)}

\qitem{5}{통보 → 새 입장 중단(준비 상태 해제) → 유예 기간 안에 진행 중 스트림 완료 → 등록 해제 순이다. 각 요청에 전체 기록을 담고 KV가 계약이 아닌 캐시인 상태 비저장 덕분에 회수된 스트림은 다른 곳에서 다시 프리필한다. 짧은 지연이 있지만 답은 정확하다. (§12.5, 제5/11장)}

\qitem{6}{단위 비용은 주, 모델, 플릿 크기, API의 공개 \$/M과 비교할 수 있지만 총지출은 성장만 측정한다. 또한 이 사용률에서 \$0.44 대 \$0.50이라는 제1장의 자체 구축 대 외부 도입 질문에 직접 답한다. (§12.5)}

\subsection*{C. 서술형(채점 기준)}

\qitem{1}{계단은 피크 ceil(340/32)=11 (4 h), 어깨 ceil(120/32)=4 (평일 \textasciitilde{}8 h), 바닥 ceil(25/32)=1→\textbf{2}(운영 바닥)이므로 평일 \textasciitilde{}76 rep-h, 주말 \textasciitilde{}48이다. 반응 계획은 3분 탄생을 고려해 10:40에 cron으로 11개까지 늘리고, 형태에는 대기열 트리거를 쓰며, 11:00–15:00에 여유 용량 +1을 둔다. 포트폴리오는 2개 약정(−38\%), 어깨 온디맨드, 7레플리카 버스트 스팟(−64\%)이다. 좋은 답은 고정 11×24 = 264 rep-h/day와 평균 \textasciitilde{}72를 비교해 부하율 \textasciitilde{}27\% → 이후 \textasciitilde{}64\%에 가까운 계산을 보인다. 일관된 청구서는 인정하지만 바닥을 스팟에 두거나 파도 안에 노드 탄생을 두면 감점한다. (§§12.1–12.5)}

\qitem{2}{좋은 글은 "흥미로움"과 "가치"를 구분한다. 오토스케일링 절감은 \$868/mo이며 2,381→1,513이 스케일러, 1,513→877이 포트폴리오다. 제10장의 \$2,381 절감이 더 컸다. 하지만 오토스케일링은 플래핑, 콜드 스타트 매복, 잘못된 지표로 위반을 \textit{만들 수 있는} 고유한 위험이 있어 엔지니어링의 대부분은 안전에 관한 것이다. 다음 한계 비용은 이 프로젝트에서 §12.3의 95 s 엔진 초기화나 원장의 서빙/용량 격차에 투입할 가능성이 크다. 후자는 인프라가 아니라 마케팅 문제다. 수단이 곱셈으로 합성되고 SLO 위반 인용 규율(상자 12-1)이 실제 논지임을 알아보는 글을 높이 평가한다. (장 전체)}

\qitem{3}{포함할 내용: SLO를 지는 두 모델은 MIG 조각에 두되 피크가 상관되면 전체 카드를 고려한다. 한 팀의 작은 모델 묶음은 MPS, 개발/노트북은 시분할, 카드를 채우는 모델은 전체 카드를 쓴다. 위험은 노드 수준 장애 반경(두 SLO 모델을 안티어피니티로 분산), MIG 재구성 드레인(노드 풀별 프로파일), 조각 대역폭과 모델별 §10.2 진단이다. 요구할 측정은 \textbf{아홉 사용 주기의 피크 상관관계}, 즉 Insull의 다양성이다. 상관된 피크를 통합하면 문제를 재현한다. 다양성이 있을 때만 전후가 약 9 × \$394 → A100 한 대(\textasciitilde{}\$2,678) + 전체 카드 1–2대가 된다. (§12.4)}

\qitem{4}{자체 호스팅 비용은 대략 고정 바닥 + 가변 rep-h다. 바닥 위에서는 토큰에 거의 선형이고 아래에서는 평평하므로 규모가 늘수록 단위 비용이 떨어진다. Insull의 원리다. 교재 수치로 \$0.50/M과 교차하는 지점은 \textasciitilde{}1.7–1.8B tokens/mo다. 아래는 바닥이 지배하고 위는 포트폴리오 가격이 지배한다. INT4는 rep-h당 토큰을 늘려 교차점을 \textasciitilde{}1.6× 낮추고, 트래픽 두 배는 같은 플릿을 활용해 \$0.28/M로 낮춘다. 0.4B/mo에서는 API가 \textasciitilde{}2× 유리하다. 메모에는 "지금은 구매하고 \textasciitilde{}1.5B/mo 또는 데이터 거버넌스가 자체 호스팅을 강제할 때 다시 검토"라고 쓰며 지배 입력인 사용률과 tokens/replica를 인용한다. 일관된 계산은 인정하고 구호는 인정하지 않는다. (§12.5, 제1장 §1.5)}

\namedsection{더 읽을거리}{}

\textbf{[이번 주]}로 표시한 자료는 실습을 직접 뒷받침하고, \textbf{[N주 차]}는 나중에 도움이 될 자료를 뜻한다.

\begin{tbbitemize}
  \item \textbf{[이번 주] Kubernetes HPA documentation — "Horizontal Pod Autoscaling" (algorithm + behavior/stabilization)} 및 \textbf{KEDA docs — ScaledObject, Prometheus scaler, cron scaler} — 기록 12-1/12-2의 권위 있는 출처다. behavior 절은 §12.2의 성향을 그대로 설명한다.
  \item \textbf{[이번 주] NVIDIA MIG User Guide · GPU Operator time-slicing docs} — 제2장 §2.6이 이번 주 직접 다룬다고 약속한 프로파일, K8s 장치 플러그인 전략, 재구성 드레인 주의점을 운영 관점에서 설명한다.
  \item \textbf{[이번 주] KServe documentation — InferenceService \& autoscaling} — 제9장의 "CRD를 훑고 12주 차에 배포하라"는 과제를 수행할 때가 되었다. 동시성 기반 기본값은 §12.2가 제도화된 모습이다.
  \item \textbf{A. Agache et al., "Firecracker" (NSDI 2020), §§1–2 다시 읽기} — 2주 차 필독 자료를 두 번째로 읽는다. 이번에는 "내 150초 중 무엇을 해결하는가?"라고 물어보자. 답이 "2초"라는 사실이 §12.3의 교훈이다.
  \item \textbf{클라우드 공급자 문서: Spot best practices (AWS) / Preemptible \& Spot VMs (GCP) · committed-use discounts / Savings Plans} — 통보 창, 용량 최적화 할당, 약정 범위 등 포트폴리오의 작은 글씨다. §12.5 형태의 계약에 서명하기 전에 읽어야 한다.
  \item \textbf{J. F. Wasik, "The Merchant of Power" (2006)} — Insull의 전기로 부하율, 수요 계량기, 맞물리는 피크, 전력 경제학의 탄생을 다룬다. 이 장의 중심 이야기를 책 한 권으로 확장하며 이 장에서 생략한 경고성 결말도 있다.
  \item \textbf{"Sub-second GPU cold start" 계열 엔지니어링 글(Modal, Fly.io, Cloud Run GPU) + containerd 지연 풀(eStargz) 문서} — 제4장이 콜드 스타트 최전선으로 예고한 스냅샷/복원, 이미지 스트리밍, 따뜻한 풀 경제성을 다룬다.
  \item \textbf{[이번 주, 재미] 내림차순으로 정렬한 클라우드 콘솔 결제 페이지} — NAT gateway(제6장), 유휴 개발 GPU, 태그 없는 수수께끼를 찾아보자. 이 장을 마치면 의례가 아니라 명세처럼 읽힌다. 모든 줄에 수단과 소유자가 있다.
\end{tbbitemize}

\vspace{5.0pt}

\begin{calloutbox}{tbbblue}{tbbbluefill}
{\normalsize\bfseries\color{tbbblue}다음 장 — 제13장: MLOps와 배포 자동화}\par\vspace{3pt}
플릿은 이제 숨 쉬고 원장도 마감되었다. 하지만 이 교재에서 배포한 모든 개선은 \textbf{사람이 kubectl을 입력해} 배포했다. 기록 12-5에는 \code{canaryTrafficPercent: 10}이라는 조절점이 매달린 채 남았다. 아무도 제어하지 않는 트래픽 분할이다. 제13장은 그 운전자를 만든다. 제4장의 다이제스트와 제6장의 저장소를 승격 워크플로로 발전시킨 \textbf{모델 레지스트리와 버전 관리}, 모든 후보에 제10장의 게이트와 제9장의 스윕을 실행하는 \textbf{CI/CD 파이프라인}(실습에서는 GitHub Actions), 오프라인 게이트가 확인할 수 없는 것을 10\% 분할로 검사하는 \textbf{카나리, A/B, 섀도 배포(shadow deployment)}, 그리고 오퍼레이터까지 외주화한 KServe와 같은 플랫폼 절충인 \textbf{관리형 엔드포인트}(SageMaker, Vertex AI)를 냉정하게 살펴본다. 목표는 모든 당직 엔지니어가 참이기를 바라는 한 문장이다. \textit{배포는 지루하다.}
\end{calloutbox}


\part{운영}
%% Chapter 13 — MLOps와 배포 자동화
\setcounter{chapter}{12}
\chapter{MLOps와 배포 자동화}
\chaptersubtitle{레지스트리, 파이프라인, 카나리, 섀도 — 배포가 엔지니어링에서 가장 지루한 문장이 되는 법}

\begin{calloutbox}{tbbblue}{tbbbluefill}
{\normalsize\bfseries\color{tbbblue}학습 목표}\par\vspace{3pt}
\begin{tbbitemize}
  \item 잘못된 배포를 개인의 부주의가 아니라 \textbf{빠진 제도}(버전이 있는 릴리스, 파이프라인, 점진적 노출)의 문제로 진단하고 피해 사용자-분 수치로 차이의 가격을 계산한다.
  \item \textbf{모델 레지스트리}를 운영한다. 릴리스 삼중항(가중치 + 설정 + 이미지 다이제스트), 게이트 증거, 계보, 단계 별칭을 관리하고 롤백을 연습된 포인터 전환으로 수행한다.
  \item \textbf{모델 CI/CD}를 구축한다. 세 경로(PR 검사, 병합 후 레지스트리 등록, 카나리와 함께 승격), 모델 CI와 앱 CI의 차이, 자체 호스팅 GPU 러너를 포함한 GitHub Actions 메커니즘을 이해한다.
  \item 각각이 \textbf{답하는 질문}에 따라 배포 전략을 선택한다. 롤링은 "부팅되는가?", 블루/그린은 "즉시 되돌릴 수 있는가?", 카나리는 "고장 났는가?", A/B는 "더 나은가?", 섀도는 "무엇을 했을 것인가?"에 답한다.
  \item 생성 모델의 \textbf{카나리 분석}을 설계한다. SLO 조항과 가드레일 프록시, 세션 일관적 분할, 자동 승격/자동 후퇴를 포함하고 오프라인 게이트가 확인할 수 없는 무엇을 검사하는지 설명한다.
  \item \textbf{관리형 엔드포인트}(SageMaker, Vertex AI)를 오퍼레이터 사다리 위에 배치하고, 무엇을 실행하고 어떤 비용을 청구하며 무엇은 언제나 사용자 몫인지 파악한다.
\end{tbbitemize}
\end{calloutbox}

\vspace{2.5pt}

\begin{calloutbox}{tbbteal}{tbbtealfill}
{\normalsize\bfseries\color{tbbteal}핵심 용어}\par\vspace{3pt}
피해 사용자-분 · 모델 레지스트리 · 릴리스 삼중항 · 계보 · 단계 별칭 · 승격 / 롤백 · CI/CD · GitHub Actions · 자체 호스팅 러너 · OIDC · 골든 요청 · 빠른 게이트 대 전체 게이트 · 점진적 전달 · 롤링 · 블루/그린 · 카나리 · 자동 분석 · 가드레일 메트릭 · A/B 테스트 · 섀도 / 다크 론치 · Argo Rollouts · 관리형 엔드포인트 · 배포 가드레일 · DORA 메트릭
\end{calloutbox}

\vspace{2.5pt}

\begin{calloutbox}{tbblightgray}{tbbboxfill}
{\normalsize\bfseries\color{tbbink}장 구성}\par\vspace{3pt}
\textbf{13.1} 금요일 배포 — 역사상 가장 비싼 누락 서버    \textbf{13.2} 레지스트리: 의미 있는 버전    \textbf{13.3} 파이프라인: 모든 약속을 장치가 검사한다    \textbf{13.4} 실제 트래픽과 만나기: 롤링, 블루/그린, 카나리, A/B, 섀도    \textbf{13.5} 관리형 엔드포인트: 완성된 오퍼레이터 사다리    \textbf{13.6} 실습: 파이프라인을 만들고 프로덕션을 깨뜨리는 데 실패하기 — 이어서 요약, 연습문제, 더 읽을거리
\end{calloutbox}

\section{금요일 배포 — 역사상 가장 비싼 누락 서버}

제12장은 정직하게 부끄러운 장면으로 끝났다. 플릿은 숨 쉬고 청구서는 마감되었지만, 모든 개선을 \textbf{사람이 노트북에서 kubectl을 입력해} 배포했다. 이 교재는 몇 주 동안 다이제스트, 저장소, 게이트, 프로브, 카나리 필드라는 대안을 조용히 쌓아 왔다. 하지만 아직 조각을 조립하지 않았고 조립되지 않은 조각은 누구도 보호하지 못한다. 현재 팀이 세부까지 그대로 재현할 수 있는 금요일 저녁 6시의 모습은 다음과 같다.

\begin{terminalbox}
\begin{Verbatim}[breaklines=true,breakanywhere=true,fontsize=\small,formatcom=\color{tbbtermfg},codes=\tbbcodeactive]
# friday 18:07. v10-awq passed the offline gate this afternoon.
$ kubectl set env deploy/chat MODEL_URI=s3://models/chat/v10-awq/
# 18:12  rolling update complete. probes green. weekend, here we come.
# 18:19  support: "answers keep cutting off mid-sentence??"
$ kubectl get pods; curl :8000/metrics | grep -E "ttft|error"
#        pods Ready. TTFT p99 0.44s. errors 0.0%. ...all GREEN.
# 18:31  more tickets. decision: roll back. ...roll back to WHAT?
$ history | grep MODEL_URI      # archaeology. Slack. more grep.
# 18:43  found it: v9-int4. re-deploy. 18:56: healthy.
#
# the bug: v10's bundled generation_config.json shipped
# max_tokens: 256 (v9: 1024). the model was FINE. the gate never
# noticed - it scores teacher-forced eval, it never GENERATES.
# damage: 37 minutes x 100% of users. mean answer: 211 -> 54 tokens.
\end{Verbatim}
\end{terminalbox}
\figcaption{사례 13-A  금요일 배포. 충돌한 것은 없고 모든 대시보드는 녹색이며 모델 자체도 완벽했다. 릴리스 밖에서 이동한 설정 파일 하나가 37분 동안 모든 답변 길이를 4분의 1로 잘랐다. 아무도 "이전"이 무엇인지 말할 수 없어 롤백에는 셸 기록을 발굴하는 12분이 걸렸다.}

조직이 가장 먼저 제시할 \textit{"엔지니어가 max\_tokens를 확인했어야 한다"}는 진단을 경계하라. 사실이지만 쓸모가 없다. 개인은 대략 일정한 비율로 실수하고 각 실수의 비용은 시스템이 결정한다. 판사가 아니라 감사인의 눈으로 사례를 다시 읽으면 세 가지 \textbf{빠진 제도}가 보인다. \textit{"이전에 무엇이 실행 중이었는가?"}에 12분 안에 답하지 못한 것은 \textbf{버전 관리} 실패다. 저장소에는 불변 바이트가 있었지만 어떤 바이트와 설정이 실행 릴리스를 구성했는지 기록하지 않았다. 게이트는 모델 점수를 계산했지만 \textit{배포 설정으로 생성해 보지 않았다.} 이는 \textbf{파이프라인} 실패다. 배포된 산출물이 아니라 확인하기 편한 대상을 검사했다. v10이 첫 1초부터 사용자 100\%를 만난 것은 \textbf{노출} 실패다. 현실과 점진적으로 접촉하지 않았고 자동 후퇴도 없었다. 엔지니어만 고치면 다른 엔지니어와 같은 금요일이 반복된다. 세 제도를 고치면 금요일은 그림 13-1의 두 번째 시간선이 된다. 같은 버그를 10\% 카나리의 답변 길이 가드레일이 5분 안에 잡고 자동 후퇴하며 PR에 증거를 기록한다. \textbf{피해 사용자-분이 37에서 0.5로} 74× 줄고 누구의 주말도 희생하지 않는다.

\begin{tbbfigure}
  \adjustbox{max width=0.94\linewidth,max totalheight=0.46\textheight}{\includegraphics{figures/fig13-1_friday.png}}\\[3pt]
  \figcaption{그림 13-1  배포 하나와 두 시간선. 수동/전체/연습 없음은 사용자 전체에 49분간 영향을 준다. 파이프라인/카나리/레지스트리는 사용자 열 명 중 한 명에게 4m47s 동안 노출한 뒤 자동 후퇴하고 증거를 첨부한다. 이 장은 두 행 사이의 거리를 다룬다.}
\end{tbbfigure}

\subsection*{2012년 8월 1일: 45분과 4억 6천만 달러}

사례 13-A가 견딜 만해 보인다면 실제로 그랬다. 산업 역사는 견딜 수 없었던 버전을 제공하며, 이 장의 규율을 대표하는 사례이므로 두 문단을 할애할 가치가 있다. \textbf{Knight Capital}은 미국 주식 시장 최대 마켓 메이커 중 하나였고 평소 미국 주식 약 열 주 중 한 주를 거래했다. NYSE 출시 기한을 맞추려고 엔지니어가 새 주문 라우팅 코드를 회사의 \textbf{프로덕션 서버 여덟 대에 수동으로 배포하다가 한 대를 빠뜨렸다.} 새 릴리스는 오래된 기능 플래그를 재사용했다. 업데이트하지 않은 여덟째 서버에서 이 플래그는 여전히 \textbf{2003년의 죽은 코드}를 가리켰다. 다른 시스템을 실습 환경에서 시험하려고 의도적으로 비싸게 사고 싸게 파는 Power Peg라는 테스트 모듈이었다. 2012년 8월 1일 시장이 열리자 플래그가 활성화되었다. 서버 일곱 대는 새 코드를 실행했고 여덟째 서버는 좀비를 깨웠다. \textbf{45분} 동안 154개 종목에 의도하지 않은 주문 수백만 건을 냈다. 손실은 약 \textbf{\$460 million}, 초당 약 \$170,000였고 회사 시가총액은 \textasciitilde{}\$1 billion이었다. 이후 SEC 명령은 이 장의 절 제목이 빠져 있는 목록과 같다. 자동화되고 검증된 배포가 없었다. \textit{사람}이 서버 여덟 대에 푸시했고 \textit{장치}는 8/8을 확인하지 않았다. 배포의 두 번째 검토도 없었다. 재사용 플래그가 새 동작을 버전 없는 죽은 코드에 연결했고 킬 스위치(kill switch)도 없었다. 반드시 기억할 세부는 \textbf{시스템이 피를 흘리는 동안 어디에서 무엇이 실행되는지 아무도 확신하지 못했다}는 것이다. 그래서 첫 "수정"은 플래그를 켜 둔 채 정상 서버 일곱 대의 \textit{새} 코드를 롤백해 상황을 \textit{악화}시켰다. Knight는 1년 안에 사실상 사라져 경쟁사에 흡수되었다.

업계의 긍정적인 극은 재난보다 먼저 등장해 대비가 더 선명하다. 2009년 Flickr의 두 엔지니어 John Allspaw와 Paul Hammond는 청중 대부분에게 오타처럼 들리는 제목으로 발표했다. \textbf{"10+ Deploys per Day."} 전체 DevOps 운동의 씨앗이 된 주장은 배포가 위험하므로 \textit{드물어야 한다}는 당시 직관을 뒤집었다. 배포가 위험하므로 \textbf{작고 자동화되며 관찰 가능하고 지속적이어야 한다.} 드문 배포는 큰 배포이고, 큰 배포는 연습하지 않은 중대한 사건이며, 연습하지 않은 중대한 사건이 Knight Capital이다. 이후 15년간의 업계 측정인 \textbf{DORA} 연구 프로그램이 이 역전을 정량화했다. \textit{가장 자주} 배포하는 팀이 변경 실패율도 \textit{가장 낮고} 복구도 가장 빠르다. 빈도가 이 장에서 구축할 메커니즘을 강제하기 때문이다. 목표 상태는 네 수치로 요약되며 그림 13-1 아래 띠가 교재 프로젝트에 맞게 가격을 매긴다. 병합부터 프로덕션까지 \textbf{리드 타임}은 며칠간 독촉하던 상태에서 22분 파이프라인으로, \textbf{배포 빈도}는 "누군가 용감할 때"에서 모든 병합으로, \textbf{변경 실패 장애 반경}은 사용자 100\%에서 자동 후퇴하는 10\% 카나리로, \textbf{복구 시간}은 49분간 발굴하던 상태에서 2분 별칭 전환으로 바뀐다. 이 장의 논지는 한 문장이다. \textbf{배포는 자동화되고 버전이 있으며 점진적이고 되돌릴 수 있는 만큼 정확히 지루해진다.} 지루함은 운영이 줄 수 있는 최고의 칭찬이다.

앞으로 90분의 지도를 보자. 각 절이 하나의 제도를 만든다. \textbf{§13.2}는 \textit{기억}을 만든다. 제6장의 저장소 위에 레지스트리를 두고, 릴리스를 게이트 보고서가 붙은 다이제스트 삼중항으로 정의하며, 승격은 별칭 전환이 된다. "이전에 무엇이 실행 중이었는가?"는 0.4초 쿼리로 답한다. \textbf{§13.3}은 \textit{반사 작용}을 만든다. 세 경로의 GitHub Actions 파이프라인이 모든 후보에서 교재의 모든 게이트인 제7장의 동등성, 제10장의 허용 오차, 제9장의 스윕을 동일하게 실행하고 큰 소리로 거부한다. \textbf{§13.4}는 \textit{신중함}을 만든다. 롤링부터 카나리, 섀도까지 점진적 전달을 다루고 카나리 자동 분석으로 제12장이 남긴 질문에 답한다. 실제 트래픽은 게이트가 확인할 수 없는 무엇을 검사하는가? \textbf{§13.5}에서는 이 메커니즘 대부분을 서비스로 판매하는 \textbf{관리형 엔드포인트}(SageMaker, Vertex)를 살펴본다. 개념은 같고 YAML은 다르며 마진이 붙는다. \textbf{§13.6}은 파이프라인을 만든 뒤 게이트와 카나리에서 한 번씩, \textit{프로덕션을 두 번 깨뜨리는 데 실패}해 점수를 얻는 실습이다. 배포는 지루하지만 거기까지 가는 과정은 그렇지 않다.

\begin{calloutbox}{tbbpurple}{tbbpurplefill}
{\normalsize\bfseries\color{tbbpurple}생각해 보기}\par\vspace{3pt}
\begin{tbbitemize}
  \item 사례 13-A의 대시보드는 사고 내내 녹색이었다. TTFT, TPOT, 오류가 모두 정상이었다. 제7/11장을 이용해 교재의 전체 메트릭 용어가 이 버그를 놓치는 이유와 이를 잡을 수 있는 가장 작은 새 메트릭을 설명하라. 방금 가드레일 메트릭을 발명했으니 §13.4까지 기억해 두자.
  \item Knight의 첫 대응자는 \textit{새} 코드를 롤백해 상황을 악화시켰다. 실제 트리거는 여덟째 서버의 재사용 \textit{플래그}였기 때문이다. 이 장의 용어로 실패를 다시 설명하라. "릴리스"는 무엇이었고 무엇에 버전이 있었으며 무엇에는 없었는가? "어디에서 무엇이 실행되는가"를 추측이 아니라 쿼리로 만들었을 §13.2의 산출물도 제시하라.
  \item 반대 문화를 가장 강한 형태로 논증하라. 한 은행은 2주 검토 게이트를 거쳐 위험 모델을 분기마다 배포한다. DORA의 논리로 드물고 무거운 방식이 올바른 체제인 경우가 있는지 논하고, 은행과 Flickr를 화해시키는 숨은 변수인 배포당 변경 배치 크기를 찾아라.
  \item 사례 13-A의 세 세계에서 피해 사용자-분을 계산하라. (a) 실제 상황, (b) 10\% 카나리로 5분 안에 감지, (c) 출시 전에 잡는 섀도 모드(노출 0\%)다. 각 세계는 추가 메커니즘에 어떤 비용을 냈으며 교재 프로젝트에는 무엇을 가장 먼저 구축할 것인가?
\end{tbbitemize}
\end{calloutbox}

\section{레지스트리: 의미 있는 버전}

제6장은 멋진 선반을 만들고 13주 차에 사서를 세우겠다고 했다. 저장소는 \textit{바이트} 문제를 해결했다. 깔끔한 접두사 아래에 불변 버전을 두고, 다이제스트로 신원을 식별하며, 병렬로 로드하고, 데이터셋과 골든 픽스처도 나란히 놓았다. 하지만 객체 더미라는 구조로는 새벽 2시에 중요한 유일한 질문에 답할 수 없다. \textbf{어느 버전이 실행 중인가? 어느 버전으로 안전하게 롤백할 수 있는가? 누가 무엇을 기준으로 검사했고 누가 승인했는가?} 금요일에 12분간 발굴한 것은 파일 시스템에 거버넌스 질문을 던진 대가였다. \textbf{모델 레지스트리}는 답을 주는 데이터베이스다. 작고 지루하지만 시스템을 지탱하는 \textit{바이트에 대한 약속}의 카탈로그다. 특정 공급업체 화면이 아니라 개념으로 스키마를 배워야 한다. MLflow, Weights \& Biases, SageMaker와 Vertex의 내장 기능 등 여러 로고를 달고 만나게 되며, JSON만으로 오후 한 번에 쓸 만한 것을 직접 만들 수도 있다.

\begin{tbbfigure}
  \adjustbox{max width=0.94\linewidth,max totalheight=0.46\textheight}{\includegraphics{figures/fig13-2_registry.png}}\\[3pt]
  \figcaption{그림 13-2  저장소 대 레지스트리. 저장소는 불변 바이트를 보관하고 레지스트리는 신원(릴리스 삼중항), 증거(게이트 보고서), 계보, 수명 주기(별칭), 감사를 보관한다. 승격과 롤백은 포인터를 옮기는 같은 작업이다.}
\end{tbbfigure}

\subsection*{릴리스 삼중항 — 금요일 버그를 구조적으로 불가능하게 만들기}

레지스트리의 원자는 "모델 하나"가 아니라 \textbf{릴리스}다. \textit{함께 요청에 답하는 완전하고 불변인 산출물 집합}을 뜻한다. 교재 서비스에서는 \textbf{다이제스트 삼중항}이다. 가중치(제6장의 접두사, 다이제스트 \code{a17b…}), max\_tokens, 템플릿, 표본 추출 기본값이 있는 서빙/생성 \textbf{설정}(\code{c4d2…}), 런타임 \textbf{이미지}(\code{i-9f3c…}, 제4장)다. 하나의 버전 번호 아래 셋을 함께 고정하면 금요일의 실패 유형이 구조적으로 사라진다. \textit{릴리스란 함께 이동하는 것이라고 정의}했으므로 설정이 "릴리스 밖에서 이동"할 수 없다. 물려받은 두 규율이 원자를 신뢰할 수 있게 한다. 제4장의 \textbf{다이제스트를 승격하고 다시 태그하지 않는다}는 규칙이 첫째다. 그 장의 생각해 볼 문제는 이유를 물었고 이제 답을 체감할 수 있다. 다시 푸시한 \code{prod} 같은 태그는 발밑에서 의미가 바뀌는 이름이다. Knight의 여덟째 서버는 돈이 걸린 태그 드리프트의 모습이다. 다이제스트는 드리프트할 수 없다. 둘째는 제10장의 \textbf{게이트 보고서를 버전에 붙인다}는 규칙이다. PASS/FAIL, 허용 오차, 슬라이스 표, 평가 세트 자체의 다이제스트를 포함한다. 그러면 "v10은 안전한가?"는 회의가 아니라 조회가 된다.

\subsection*{별칭: 포인터로 승격하고 연습한 대로 롤백하기}

두 번째 아이디어는 컴퓨팅보다 오래된 \textbf{간접 참조}다. 배포는 버전을 직접 참조하지 않는다. 제6장의 \code{MODEL\_URI}는 \textbf{별칭}(\code{…/\allowbreak chat/\allowbreak prod})을 가리키고 레지스트리가 별칭을 버전에 대응시킨다. 승격은 포인터 이동(\code{prod: v9 → v10})이고 롤백은 같은 이동의 반대 방향이다. 둘 다 원자적이고 로그에 기록되며 1초보다 훨씬 짧게 걸린다. 이후 플릿은 5주 차부터 신뢰해 온 제5장의 일반 롤아웃, 즉 추가 파드, 프로브, 준비 상태 게이팅으로 수렴하거나 다음 파드 탄생 시 별칭을 다시 해석한다. 새벽 2시 이야기가 어떻게 달라지는지 보자. \textit{결정}("v9으로 돌아간다"), \textit{지식}("v9은 게이트를 통과하고 프로덕션에서 11일간 서빙한 마지막 버전이다"), \textit{메커니즘}(명령 하나)이 이제 연습된 상태로 한곳에 있다. 실습에서 시간을 재면 \textbf{별칭 전환 0.4 s, 플릿 수렴 \textasciitilde{}2 minutes}다. §12.3의 탄생 계산에 따라 롤백 하한은 콜드 스타트이며 이를 줄이는 연습을 한 이유다. 수명 주기 단계(staging → production → archived)는 규칙이 붙은 구별된 별칭일 뿐이다. 게이트 보고서, 승인 기록, §13.4 이후에는 카나리 판정이 없으면 아무것도 \code{production}에 도달하지 못한다.

\begin{terminalbox}
\begin{Verbatim}[breaklines=true,breakanywhere=true,fontsize=\small,formatcom=\color{tbbtermfg},codes=\tbbcodeactive]
$ python registry.py list chat
 ver  stage       weights        config   image     gate        by
 v8   archived    a03f…/int8     c11e…    i-77aa…   PASS 09-28  jslee
 v9   production  a17b…/int4     c4d2…    i-9f3c…   PASS 10-05  minj
 v10  staging     b920…/awq      c9aa…    i-2be1…   PASS 10-12  CI

$ python registry.py show chat/v10 --lineage
 commit 8d31f2  base Qwen2.5-3B  calib calib-2026-10a (digest 77e1…)
 eval-set golden-v4 (digest 3c9d…)  gate: delta -0.18pt (tol 0.5) PASS

$ python registry.py rollback chat --to v9        # the 2 a.m. one-liner
alias prod: v10 -> v9   (flip 0.4s; fleet converges via rollout)
audit: 2026-10-12T02:13Z rollback by oncall-kim, reason="ans-len"
\end{Verbatim}
\end{terminalbox}
\figcaption{기록 13-1  직접 사용한 레지스트리. 금요일에 답하지 못한 세 질문인 무엇이 실행 중이고 무엇이 안전하며 무엇이 바뀌었는지를 각각 한 줄로 답한다. 롤백은 승격과 같은 메커니즘이므로 지루해질 때까지 연습할 수 있다.}

\subsection*{무엇이 버전을 올리는가 — 변경 분류 체계}

팀은 겉보기에는 작은 질문에서 멈추곤 한다. \textit{어떤 변경에 새 버전이 필요한가?} 레지스트리의 답은 기계적이다. \textbf{삼중항에서 어떤 다이제스트든 바뀌면 무조건 새 릴리스다.} 하지만 변경 클래스마다 \textit{결과}가 다르므로 한 번 적어 두면 반복되는 논쟁을 끝낼 수 있다. 두 가지 미묘한 경우를 따로 짚자. \textbf{설정 전용 릴리스}는 금요일의 수정처럼 같은 가중치에 \code{max\_tokens}만 고친 경우다. 실제 릴리스이므로 새 버전과 새 게이트 실행이 필요하지만 게이트의 \textit{범위}를 줄일 수 있다. 가중치가 같으면 동등성 검사는 의미가 없지만, 금요일 문제가 설정에 있었으므로 골든 요청과 가드레일 기준선 등 생성에 관련된 모든 검사를 다시 실행해야 한다. \textbf{평가 세트 변경}은 더 교묘하다. 릴리스를 전혀 만들지 않으면서 이후 모든 판정의 의미를 조용히 다시 정의한다. 그래서 평가 세트도 코드처럼 버전 관리(golden-v4 → golden-v5, 다이제스트 고정)한다. 기존 게이트 보고서는 사용한 세트의 이름을 계속 남기고, 세트 변경 후 첫 실행은 비교하지 않고 \textit{기준선을 다시 설정}한다. 어느 시험인지 모르면 성적에는 의미가 없다.

\begin{tbbtablewrap}
  \begin{tabular}{@{}>{\raggedright\arraybackslash}m{\dimexpr 0.2500\tbltextwidth-2\tabcolsep\relax} >{\raggedright\arraybackslash}m{\dimexpr 0.2287\tbltextwidth-2\tabcolsep\relax} >{\raggedright\arraybackslash}m{\dimexpr 0.2660\tbltextwidth-2\tabcolsep\relax} >{\raggedright\arraybackslash}m{\dimexpr 0.2553\tbltextwidth-2\tabcolsep\relax}@{}}
  \arrayrulecolor{tbbnavy}\specialrule{1pt}{0pt}{0pt}
  \rowcolor{tbbtablehead}{\centering \bfseries\color{tbbnavy}변경} & {\centering \bfseries\color{tbbnavy}새 릴리스?} & {\centering \bfseries\color{tbbnavy}다시 실행할 항목} & {\centering \bfseries\color{tbbnavy}일반적인 전략(§13.4)} \\
  \arrayrulecolor{tbbnavy}\specialrule{0.7pt}{0pt}{2pt}
  {\textbf{새 가중치(재학습/양자화)}} & {예 — 새 삼중항} & {전체: 동등성, 게이트, 스윕, 카나리} & {카나리. 모델 계열 경계에서는 섀도} \\
  \arrayrulecolor{tbbtableborder}\hline
  {\textbf{설정만(표본 추출, 토큰)}} & {예 — 설정 다이제스트} & {생성 경로 검사: 골든 요청, 가드레일 기준선, 카나리} & {카나리(금요일 유형이 여기 있다!)} \\
  \arrayrulecolor{tbbtableborder}\hline
  {\textbf{이미지만(의존성, 서버 코드)}} & {예 — 이미지 다이제스트} & {단위 + 골든 요청 + 스모크 스윕. 가중치/설정이 같으면 게이트 생략 가능} & {롤링. 지연 시간 관련이면 카나리} \\
  \arrayrulecolor{tbbtableborder}\hline
  {\textbf{평가 세트 / 골든 픽스처}} & {릴리스 없음 — 새 세트 버전} & {배포 없음. 다음 게이트에서 명시적으로 기준선 재설정} & {— (배포가 아닌 측정 변경)} \\
  \arrayrulecolor{tbbtableborder}\hline
  {\textbf{별칭 전환(승격/롤백)}} & {아니요 — 포인터 이동} & {없음. 참조하는 릴리스는 이미 검사됨} & {0.4 s 작업 자체} \\
  \arrayrulecolor{tbbnavy}\specialrule{1pt}{2pt}{0pt}
  \end{tabular}\par
  \figcaption{표 13-1  변경 분류 체계. 다이제스트 변경은 모두 릴리스라는 한 줄 규칙이지만 무엇을 다시 실행할지 알면 파이프라인은 빠르고 논쟁은 짧아진다.}
\end{tbbtablewrap}

로고보다 개념이 중요하므로 자체 구축 대 외부 도입을 짚고 넘어가자. \textbf{MLflow의 Model Registry}는 단계, 별칭, 계보, UI를 제공하는 일반적인 개방형 기본값이고 클라우드는 자체 기능을 묶어 제공한다. 실습의 \code{registry.py}는 저장소의 JSON 파일 위에 만든 200줄짜리 프로그램이다. 소프트웨어가 아니라 \textit{스키마와 규율}에 가치가 있음을 보여 주려고 의도적으로 단순하게 만들었다. 무엇을 도입하든 금요일 테스트를 적용하라. 아무 준비 없는 터미널에서 온콜(on-call) 엔지니어가 1분 안에 \textit{실행 중 / 안전 / 변경 / 담당자}에 답할 수 있는가? 그렇다면 레지스트리다. 위키 페이지가 필요하다면 의견을 가진 바이트 더미일 뿐이다. 구체적으로 실습 레지스트리의 전체 스키마를 보자. 이 절의 모든 주장을 항목 하나와 JSON 20줄에 담았고, 당연히 저장소 안에서 버전 관리한다.

\begin{terminalbox}
\begin{Verbatim}[breaklines=true,breakanywhere=true,fontsize=\small,formatcom=\color{tbbtermfg},codes=\tbbcodeactive]
{ "model": "chat",
  "aliases": { "prod": "v9", "staging": "v10" },
  "versions": {
    "v10": {
      "triple": { "weights": "sha256:b920…", "config": "sha256:c9aa…",
                  "image":   "sha256:2be1…" },
      "gate":   { "verdict": "PASS", "delta_pt": -0.18, "tol": 0.5,
                  "worst_slice": ["streetcar", -0.6],
                  "eval_set": "golden-v4@sha256:3c9d…" },
      "lineage":{ "commit": "8d31f2", "base": "Qwen2.5-3B",
                  "calib": "calib-2026-10a@sha256:77e1…" },
      "history":[ {"t":"2026-10-12T09:14Z","event":"registered","by":"CI"},
                  {"t":"2026-10-12T14:03Z","event":"canary-retreat",
                   "by":"deploy.yaml","reason":"ans-len guardrail"} ]
  } } }
// promotion = edit "aliases" (atomically, via a conditional write).
// every field above is one 2 a.m. question, pre-answered.
\end{Verbatim}
\end{terminalbox}
\figcaption{기록 13-2  registry-lite의 전체 스키마. 삼중항, 증거, 계보, 수명 주기, 감사라는 다섯 아이디어를 JSON 20줄로 표현했다. 공급업체는 UI와 확장성을 추가하지만 제도는 스키마에 담긴다.}

\begin{calloutbox}{tbbpurple}{tbbpurplefill}
{\normalsize\bfseries\color{tbbpurple}생각해 보기}\par\vspace{3pt}
\begin{tbbitemize}
  \item 금요일 세계에서는 표현조차 할 수 없던 릴리스의 레지스트리 항목을 설계하라. v9과 같은 가중치에 새 설정(max\_tokens 수정)을 쓴다. 버전은 무엇이고 계보에는 무엇을 기록하는가? "모델은 바뀌지 않았는데"도 게이트의 어느 부분을 다시 실행해야 하며 그 이유는 무엇인가?
  \item 제4장은 파이프라인이 다시 태그하는 대신 다이제스트를 승격하는 이유를 물었다. 두 실패 이야기를 바탕으로 답하라. 두 엔지니어가 태그 하나를 두고 벌이는 경합과, v9이 누군가 다시 푸시한 태그라면 "v9 재배포"가 무엇을 뜻하는지 불분명한 롤백을 각각 한 문단으로 설명하라.
  \item PM이 "MLflow도 운영할 서비스가 하나 더 늘어나는 것인데 스프레드시트를 쓰면 안 되는가?"라고 묻는다. 전제를 인정하고 깨지는 지점을 찾아라. 원자적 별칭 전환, 감사, CI용 API 중 어느 레지스트리 속성이 스프레드시트에서 먼저 죽으며 어떤 사고를 만드는가?
  \item 레지스트리는 각 승격을 누가 승인했는지 기록한다. 누군가 게이트를 2× 여유로 통과한 모든 것을 자동 승인하자고 제안한다. 게이트가 설정을 보지 못한 사례 13-A와 자동 승인자인 §13.4 카나리를 이용해 양쪽을 논증하라. 사람은 정확히 어디에 있어야 하는가?
\end{tbbitemize}
\end{calloutbox}

\section{파이프라인: 모든 약속을 장치가 검사한다}

레지스트리는 기억하고 \textbf{파이프라인}은 강제한다. 7주 차 이후 실제로 쌓은 것을 돌아보면 놀라운 사실이 드러난다. 약속한 거의 모든 품질이 \textit{이미 실행 가능}하다. §7.6의 동등성 검사는 스크립트고 §10.3의 게이트는 허용 오차가 있는 스크립트다. §9.5의 스윕은 locustfile이고 골든 요청은 저장소의 파일이다. 제6장은 이를 "13주 차 CI 게이트를 위해" 심어 두었다. 빠진 것은 검사가 아니라 \textbf{절차}다. 모든 후보에서 모든 검사를 같은 방식으로 순서대로 실행하고, 금요일이라고 지루한 항목을 건너뛸 수 없게 하는 장치가 필요하다. 그 장치가 CI/CD다. \textit{지속적 통합}은 모든 변경을 빌드하고 검사하며, \textit{지속적 전달}은 통과한 모든 변경을 패키징하고 등록해 승격할 준비를 한다. 올바른 제5장의 용어로 표현하면 파이프라인은 \textbf{원하는 상태의 또 다른 작성자}일 뿐이다. 레지스트리 별칭과 Deployment env 변수를 수정하고 조정 루프가 나머지를 처리한다. 프로덕션에서 새로 실행되는 것은 없지만, 새롭게도 검사하지 않은 것은 아무것도 프로덕션에 도달하지 못한다.

\begin{tbbfigure}
  \adjustbox{max width=0.94\linewidth,max totalheight=0.46\textheight}{\includegraphics{figures/fig13-3_pipeline.png}}\\[3pt]
  \figcaption{그림 13-3  세 경로. 경로 1은 모든 PR에서 CPU로 빠르게 위생을 검사한다. 경로 2는 모든 병합에서 GPU 러너로 무거운 산출물과 게이트를 처리하고 증거가 있는 등록 릴리스로 끝난다. 경로 3은 사람이 승인하고 카나리가 판정하며 장치가 실행하는 승격이다. 붉은 경로가 제품이다.}
\end{tbbfigure}

\subsection*{모델 CI가 앱 CI와 다른 이유 — 설계를 좌우하는 네 가지 차이}

앱 파이프라인에 모델을 덧붙인 팀은 네 가지 불일치를 발견한다. 그림 13-3의 형태는 각각을 수용한 결과다. \textbf{첫째, 판정은 코드의 불리언이 아니라 데이터에 대한 허용 오차다.} "테스트 통과"는 "\textit{평가 세트 golden-v4에서} Δ ≤ 0.5 pt top-1"이 된다. 따라서 평가 세트 자체도 버전과 다이제스트가 고정된 픽스처다. 시험을 바꾸면 모든 과거 성적의 의미가 조용히 변하므로 제6장에서 이를 저장했다. \textbf{둘째, 산출물은 기가바이트다.} 앱 CI는 환경마다 저렴하게 다시 빌드하지만 모델 CI는 \textbf{한 번 빌드하고 다이제스트를 승격}한다. 같은 바이트가 dev → staging → prod로 흐른다. 제4장의 레지스트리는 빠진 계층만 옮기고 제6장의 저장소는 아무것도 옮기지 않는다. 다른 방식은 "스테이징에서는 동작했다"는 거짓말을 되살린다. \textbf{셋째, 일부 단계에는 실리콘이 필요하다.} 보정, 게이트, 엔진 빌드에는 GPU가 필요하지만 GitHub 호스팅 러너에는 적합한 장치가 없다. 저장소에 등록된 스팟 T4 한 대인 \textbf{자체 호스팅 러너}가 \code{gpu} 태그 작업을 가져간다. 실습 자료 스크립트로 20분이면 설정한다. §12.5의 포트폴리오 사고를 CI에 적용한 것이다. 러너는 중단 가능한 배치 작업으로 스팟의 진정한 용도다. \textbf{넷째, 최종 검사에는 실제 트래픽이 필요하다.} 어떤 오프라인 단계도 프로덕션 분포를 보지 못한다. 따라서 파이프라인의 마지막 단계는 테스트가 아니라 \textit{전략}이다. §13.4의 카나리를 경로 3으로 연결한다.

\begin{terminalbox}
\begin{Verbatim}[breaklines=true,breakanywhere=true,fontsize=\small,formatcom=\color{tbbtermfg},codes=\tbbcodeactive]
# .github/workflows/ci.yaml (condensed - the handout has the full file)
on:
  pull_request:                      # LANE 1: minutes, CPU
  push: {branches: [main]}          # LANE 2: the heavy lane
jobs:
  checks:                           # lane 1
    runs-on: ubuntu-latest
    steps: [lint, unit, docker build --push (by digest),
            quick-gate  # golden requests + 200-sample eval, CPU]
  release-candidate:                # lane 2
    if: github.ref == 'refs/heads/main'
    runs-on: [self-hosted, gpu]     # the spot T4 runner
    steps: [export + equivalence (§7.6),
            quantize + calibrate (Ch. 10),
            gate.py --tol 0.5       # FAIL -> the job fails. loudly.
            smoke-sweep (§9.5, 3 min),
            registry.py register --stage staging  # triple + report]
    permissions: {id-token: write}  # OIDC: no long-lived cloud keys
\end{Verbatim}
\end{terminalbox}
\figcaption{기록 13-3  요약한 워크플로. 경로 1은 몇 분 안에 PR을 정직하게 유지하고, 경로 2는 병합마다 무거운 비용을 한 번 낸 뒤 등록 릴리스로 끝난다. OIDC 줄은 제4장의 인증 질문에 올바르게 답한다. 파이프라인은 몇 분 동안 클라우드 역할을 맡고 유출할 키가 없다.}

\subsection*{중요한 두 실행 — 팀을 바꾸는 문장}

\begin{terminalbox}
\begin{Verbatim}[breaklines=true,breakanywhere=true,fontsize=\small,formatcom=\color{tbbtermfg},codes=\tbbcodeactive]
# run 1: merge 8d31f2 ("switch calib set to 2026-10a")   -- GREEN
  export+equivalence   PASS   max|dz|=3.1e-3, top1-agree 100%   4m
  quantize+calibrate   OK     awq int4, 1.91 GiB                 6m
  gate                 PASS   delta -0.18pt (tol 0.5);
                              worst slice "streetcar" -0.6pt     7m
  smoke sweep          PASS   468 tok/s agg (-2.1% vs base)      3m
  register             OK     chat/v10 -> staging (evidence attached)
  total 22m 14s        bot comment on PR: report + artifact links

# run 2: merge a90c11 ("improve prompt template")        -- RED
  quick-gate           FAIL   golden request 07: expected JSON,
                              got prose; 09: missing citation;
                              12: language mismatch (ko->en)     1m31s
  >> merge blocked. no build, no registry entry, no debate.
# the check that refused has passed 41 times. it is not an opinion.
\end{Verbatim}
\end{terminalbox}
\figcaption{기록 13-4  중요한 두 실행. 녹색은 증거가 있는 레지스트리 항목으로 끝나는 기분 좋은 관료 절차다. 붉은색이 바로 제도다. 91초 안에 이름 붙은 회귀 세 개와 비활성화된 병합 버튼이 나타나 품질 이견이 논쟁으로 번지지 않는다.}

구매한 가치는 붉은 실행이며 문화적 메커니즘도 한 문단을 쓸 만하다. 빠른 게이트가 프롬프트 템플릿 회귀를 잡으면 \textit{병합 버튼이 회색으로 바뀐다.} \textbf{아무도 논쟁하지 않는다.} 거부한 검사는 의견을 가진 동료가 아니라 41번 통과시킨 바로 그 검사다. 파이프라인은 품질 이견을 사람 사이 문제에서 인프라 문제로 바꾼다. 두 규율이 신뢰성을 지킨다. \textbf{건너뛰지도 분기하지도 않는다.} \code{[skip-gate]} 커밋 태그가 생기는 순간 금요일이 돌아온다. 비상 경로는 게이트 생략이 아니라 \textit{더 빠른 카나리}다. \textbf{파이프라인은 자신이 만든 산출물을 먹는다.} 게이트를 거친 것과 배포되는 것은 바이트 단위로 동일해야 한다. 다시 다이제스트다. 검사와 배포 사이에 다시 빌드하면 모든 검사가 무효가 된다. 보안 팀이 언젠가는 보낼 공급망 메모도 미리 반영하자. 산출물 서명(cosign)과 출처 기록(SLSA)은 제4장의 "다이제스트를 신뢰하라"가 성숙한 모습이며 처음부터 경로 2에 넣어야 한다. 지금은 워크플로 두 단계지만 나중에는 발굴 프로젝트가 된다.

실습은 보안 문단 하나를 강제한다. 너무 흔한 실수라 별도 사고 장르가 되었기 때문이다. \textbf{자체 호스팅 러너는 워크플로가 시키는 모든 것을 실행한다.} \textit{공개} 저장소에 붙고 \textit{포크} 풀 리퀘스트로 트리거되는 러너는 GPU, 네트워크 위치, 부주의했다면 클라우드 자격 증명까지 가진 임의 코드 실행 서비스다. 규율은 다음과 같다. 자체 호스팅 러너는 \textbf{비공개 저장소나 포크가 아닌 트리거만} 처리한다. 경로 1에서 포크를 상대하는 검사는 GitHub 호스팅 러너에 남겨 둔다. 러너의 클라우드 역할은 작업에 꼭 필요한 범위로 제한하고 OIDC 신뢰 정책에 저장소와 브랜치를 명시한다. \code{pull\_request\_target}은 문서가 경고하는 대로 위험한 장치로 다룬다. 이국적인 이야기가 아니다. 제2장의 격리 직관을 프로덕션에서 실행할 것을 빌드하는 장치에 적용한 것이다. 그래서 프로덕션보다 덜 민감한 것이 아니라 \textit{더} 민감하다. 파이프라인이 침해되면 이후 모든 녹색 검사는 거짓말이 된다.

\begin{calloutbox}{tbbpurple}{tbbpurplefill}
{\normalsize\bfseries\color{tbbpurple}생각해 보기}\par\vspace{3pt}
\begin{tbbitemize}
  \item 경로 1의 빠른 게이트는 속도를 위해 평가 표본 200개를 쓰고 경로 2의 전체 게이트는 10,000개를 쓴다. 통계적 대가를 계산하라. 실제 회귀가 1.0 pt일 때 각각 어떤 Δ를 신뢰성 있게 감지할 수 있는가? 잡음 ≈ 1/√n이라는 경험 법칙을 사용하라. 경로 1이 정확도 드리프트가 아니라 어떤 버그를 잡기 위해 존재하는지도 설명하라.
  \item 자체 호스팅 GPU 러너는 스팟 인스턴스다. 게이트 중간에 회수되는 실패를 따라가 보자. 워크플로에는 무엇이 표시되는가? 이것이 단지 성가신 일에 그치려면 작업은 어떤 성질을 가져야 하는가? 멱등 단계와 재개 가능한 산출물을 고려하라. CI 러너가 §12.5의 표현대로 스팟의 진정한 용도인 이유도 설명하라.
  \item 팀원이 스윕 처리량에 "req/s@SLO가 >5\% 떨어지면 거부"라는 하드 임곗값 게이트를 적용하자고 한다. 옹호하고 공격하라. 제9장의 어떤 인프라 잡음이 하드 처리량 게이트를 불안정하게 만드는가? 02:00 오경보 없이 신호를 유지할 경고 대 차단, 다시 설정한 기준선 대비 비교의 2단 설계를 제안하라.
  \item "게이트를 절대 건너뛰지 않는다"는 전제에서 팀의 비상 배포 정책을 세 문장으로 작성하라. 비상시 카나리 속도, 승인 정족수, 롤백 사전 스테이징 중 무엇이 바뀌고 무엇은 절대 바뀌지 않는가? Knight에 적용해 정책이 그들을 구했을지 검토하라.
\end{tbbitemize}
\end{calloutbox}

\section{실제 트래픽과 만나기: 롤링, 블루/그린, 카나리, A/B, 섀도}

파이프라인의 마지막 단계는 테스트일 수 없다. 검사할 대상인 \textbf{프로덕션 트래픽}이 오프라인에 존재하지 않기 때문이다. 평가 세트는 지난달 사용자를 비추는 백미러이고 스테이징 클러스터는 축소 모형이다. 오늘 밤의 신조어, 오늘 아침 장애로 생긴 다급한 쿼리, 30k-token 계약서를 붙여 넣는 거대 사용자를 포함한 실제 분포는 정확히 한곳에만 존재한다. 이 절의 각 전략은 새 릴리스의 목줄을 잡은 채 \textit{현실을 만나게} 하는 서로 다른 프로토콜이며, 중요한 규율은 \textbf{서로 다른 질문}에 답한다는 점이다. 유행이 아니라 질문으로 선택한다. 롤링은 \textit{부팅되는가?}, 블루/그린(blue/green)은 \textit{1초 안에 되돌릴 수 있는가?}, 카나리는 \textit{실제 트래픽으로 판단했을 때 고장 났는가?}, A/B는 \textit{사용자가 판단할 때 더 나은가?}, 섀도는 누구도 위험을 감수하기 전에 \textit{무엇을 했을 것인가?}를 묻는다.

\begin{tbbfigure}
  \adjustbox{max width=0.94\linewidth,max totalheight=0.46\textheight}{\includegraphics{figures/fig13-4_strategies.png}}\\[3pt]
  \figcaption{그림 13-4  다섯 전략과 다섯 질문. 롤링은 시작을 보호하고, 블루/그린은 2× 플릿 비용으로 즉시 되돌리기를 사며, 카나리는 실제 트래픽을 판정하는 핵심 도구다. A/B는 안전이 아니라 제품 도구이고 섀도는 사용자 위험 없이 연습한다. 아래에는 Haldane의 새와 이를 은퇴시킨 감지기가 있다.}
\end{tbbfigure}

\subsection*{롤링과 블루/그린 — 물려받은 두 단}

\textbf{롤링}은 제5장부터 사용했다. 파드를 몇 개씩 교체하고 준비 상태 프로브가 각각의 입장을 허용한다. 정확히 무엇을 보호하는지 알아야 한다. \textit{시작하지 못하거나 프로브에 실패하는 파드}를 막는다. 반대로 무엇을 보지 못하는지도 알아야 한다. 금요일 파드는 아름답게 시작했고 따뜻함을 확인하는 프로브를 통과했으며 완벽한 상태로 짧은 답을 서빙했다. 롤링만으로 충분한 경우는 의존성 업데이트나 자원 조정처럼 실패 유형이 시작 자체인 변경뿐이다. 모델 릴리스에서는 더 영리한 전략이 올라타는 \textit{기반}이다. \textbf{블루/그린}은 전체 플릿 두 개를 실행한다. 블루가 서빙하고 그린은 후보를 올린 채 쉬며, 그린을 스모크 테스트와 따뜻한 캐시로 검증한 뒤 라우터를 전환한다. 후회하면 몇 초 만에 다시 전환한다. GPU 플릿에서의 정직한 문제는 제12장의 계산이다. 둘째 플릿은 둘째 청구서다. 릴리스 창마다 교재 플릿에 \textasciitilde{} \$1,500/mo가 들며 몇 분 동안은 견딜 수 있지만 습관이 되면 터무니없다. 전환 전 검증도 여전히 \textit{합성} 트래픽이므로 롤링처럼 금요일 유형 버그를 보지 못한다. 블루/그린은 엔진 업그레이드, CUDA 업데이트처럼 즉시 되돌리기가 필요한 \textit{설정/인프라} 릴리스에서 빛난다. 라우터 전환 메커니즘은 우연이 아니라 카나리가 더 세밀하게 재사용하는 힘이다.

\subsection*{카나리 — Haldane의 새와 새를 은퇴시킨 감지기}

핵심 도구의 어원은 비유 자체가 시스템을 지탱하므로 살펴볼 가치가 있다. 1890년대 광산 폭발을 조사한 생리학자 \textbf{John Scott Haldane}은 작업자가 작고 온혈인 보초인 카나리를 데리고 다니자고 제안했다. 대사가 더 빨라 광부가 아무것도 느끼기 \textit{몇 분 전}에 일산화탄소에 쓰러졌다. 작업자 \textbf{앞에 노출}한 민감하고 희생적인 감지기로, 새가 노래를 멈추고 휘청이는 모습을 저렴하게 관찰하고 밖으로 나가는 저렴한 조치를 취할 수 있었다. 영국 광산은 전자 감지기가 마침내 은퇴시킨 \textbf{1986년}까지 새를 데리고 다녔다. 모든 조항을 대응시켜 보자. \textbf{카나리 릴리스}는 실제 트래픽의 5–10\%를 받아 새가 실제 공기를 마시게 하고, 90\%는 기존 릴리스에 두어 작업자가 입구에서 기다리게 한다. \textbf{자동 분석}이 카나리를 지켜본다. 1911년보다 발전한 전자 감지기라 누구도 새를 봐야 한다는 사실을 기억할 필요가 없다. 나쁜 판정은 몇 초 안에 밖으로 나가는 \textbf{자동 후퇴}를 일으킨다. 이 교재 영역에서는 세 설계 선택이 성공을 좌우한다. \textit{요청이 아니라 세션으로 분할}한다. 제11장의 대화는 한쪽에 온전히 살아야 한다. 대화 ID의 일관된 해시로 §11.5의 어피니티 메커니즘을 다른 용도에 쓴다. 그렇지 않으면 사용자는 인격이 깜박이는 것을 보고 메트릭은 뒤섞인 것을 비교한다. \textit{카나리와 기준선을 동시에 비교}하고 어제와 비교하지 않는다. 트래픽 자체에 제12장의 일중 성격이 있어 같은 시각의 대조군만 릴리스를 분리한다. \textit{SLO뿐 아니라 가드레일을 관찰}해야 한다. 여기서 금요일 버그가 수업료를 낸다.

\begin{tbbfigure}
  \adjustbox{max width=0.94\linewidth,max totalheight=0.46\textheight}{\includegraphics{figures/fig13-5_canary.png}}\\[3pt]
  \figcaption{그림 13-5  카나리의 구조. 세션 해시 분할, 동시 기준선, SLO 조항과 가드레일 프록시를 확인하는 60초 분석 루프, 두 자동 출구로 구성된다. FAIL 행은 금요일 버그를 10\% 노출에서 5분 안에 잡은 모습이다.}
\end{tbbfigure}

\textbf{가드레일 메트릭}은 생성 시대가 카나리 전통에 더한 요소다. §13.1의 생각해 볼 문제에서 이미 발명했다. 진짜 품질 평가가 아직 잠든 동안 \textit{"모델이 이상해졌다"}를 잡는 저렴하고 모델에 구애받지 않는 프록시다. 교재에서 사용하는 집합은 모두 60초 루프에서 로그로 계산할 수 있다. \textbf{답변 길이 분포}는 금요일의 211 → 54 tokens를 한 창 안에서 외친다. \textbf{거부/사과 비율}은 과도한 정렬 회귀를 잡고, \textbf{형식 유효성}은 JSON이 파싱되는지 인용이 해석되는지 확인한다. \textbf{반복 점수}와 \textbf{언어 일치}도 본다. 한국어 질문에는 한국어 답이 나와야 한다. 어느 것도 "좋음"을 측정하지 않고 모두 \textit{동작의 동일성}을 측정한다. 바로 이것이 배포의 귀무가설이다. 제11장의 TTFT/TPOT/오류라는 SLO 조항이 시스템을 지키고 가드레일은 \textit{동작}을 지킨다. 더 미묘한 것은 A/B를 기다린다. 판정 논리는 의도적으로 단순하게 둔다. 고정 창, 기준선 대비 명시적 임곗값, 하나라도 실패하면 후퇴한다. 원시 Deployment에서는 제5장의 더 읽을거리에 등장한 Argo Rollouts가 단계별 가중치와 Prometheus 분석을 자동화한다. InferenceService에서는 제12장이 남긴 KServe의 \code{canaryTrafficPercent}가 이를 수행한다. 실습은 마법이 없음을 보여 주려고 후자를 40줄 루프로 제어한다.

\begin{terminalbox}
\begin{Verbatim}[breaklines=true,breakanywhere=true,fontsize=\small,formatcom=\color{tbbtermfg},codes=\tbbcodeactive]
$ gh workflow run deploy.yaml -f version=v10      # lane 3, replayed
14:00:12  canary: chat canaryTrafficPercent 0 -> 10 (session-hashed)
14:01:12  window 1/5:  TTFT p99 0.46s vs 0.44s OK | err 0.0% OK
          | ans-len 208t vs 211t OK | refusal 1.2% vs 1.1% OK
14:02:12  window 2/5:  all OK
14:03:12  window 3/5:  ans-len 55t vs 210t
          guardrail: |delta| 74% > 30%  -> FAIL
14:03:14  verdict: RETREAT. canary 10% -> 0%. alias prod stays v9.
14:03:20  posted to PR #241: verdict, metric graphs, sample diffs
          (12 canary answers, all ending mid-sentence at 54 tokens)
# exposure: ~10% of users for 3m02s = 0.3 bad user-minutes.
# friday, replayed with institutions: the bird went quiet,
# the detector noticed, the crew never entered the mine.
\end{Verbatim}
\end{terminalbox}
\figcaption{기록 13-5  금요일 버그 대 카나리. 세 창에서 가드레일 하나가 작동하고 자동 후퇴했으며 수정이 일어날 곳에 증거를 기록했다. 사람의 개입이 필요하지 않아 아무도 호출받지 않았다. 사고 전체가 PR 댓글 하나다.}

\subsection*{얼마나 길어야 충분한가 — 냅킨 위의 카나리 통계}

창 길이는 속설이 아니라 머릿속으로 계산할 수 있는 표본 크기 문제다. 교재 플릿은 피크에 \textasciitilde{}23 turn-arrivals/s를 흡수하므로(§11.2) 10\% 카나리는 \textbf{\textasciitilde{}2.3 requests/s ≈ 5분에 답변 690개}를 본다. 690개 표본이 감지할 수 있는 것은 효과 크기에 전적으로 달려 있다. 답변 길이가 211 → 54 tokens로 74\% 떨어지고 분포 폭이 \textasciitilde{}±30\%인 금요일의 절벽은 \textit{첫} 60초 창 안에서도 명백하다. 절벽에는 수십 개 표본이면 충분하므로 기록 13-5의 판정은 여유 있게 세 번째 창에 나온다. 하지만 형식 유효성이 97\%에서 3\% 떨어지는 \textit{미묘한} 회귀는 잡음과 구분하기 위해 각 집단에 수천 개 표본이 필요하다. 이항 경험 법칙은 감지 가능한 Δ ≈ 2√(p(1−p)/n)이며 n = 690에서는 잘해야 약 ±1.3 percentage points다. 기억할 설계 태도는 \textbf{카나리는 기울기가 아니라 절벽을 잡는다}는 것이다. 충돌, SLO 위반, 가드레일 급변이라는 절벽 유형에 짧은 체류와 단단한 임곗값을 설정한다. 기울기는 이를 위해 만든 도구에 맡긴다. 제14장의 드리프트 모니터는 몇 주의 프로덕션 데이터를 보고, §13.4의 A/B는 균형 잡힌 집단과 사전 등록한 검정력을 사용한다. 단계 증가 일정도 노출 제어와 같은 논리를 인코딩한다. 1\% → 10\% → 50\% → 100\%로 각 단계에서 머무르면 가장 저렴한 절벽은 1\%에서 잡고, 충돌에는 \textasciitilde{}10개 표본이면 충분하다. 가드레일 절벽은 10\%에서 잡고 "더 지켜보자"는 롱테일은 마지막 체류로 제한한다. 각 단계는 사용자 노출로 감지력을 구매하며, 일정을 적어 두면 판단이 설정 필드가 된다.

\subsection*{A/B와 섀도 — 나머지 두 질문}

\textbf{A/B 테스트}는 카나리의 트래픽 분할 배관만 공유하고 \textit{나머지는 전혀 공유하지 않는다.} 둘을 혼동하는 것이 이 절의 고전적 오류다. 카나리는 몇 분에서 몇 시간 동안 작은 비율을 노출하고 "같이 동작한다"는 귀무가설을 자동 판정하는 \textbf{안전} 도구다. A/B는 며칠에서 몇 주 동안 크고 균형 잡힌 버킷(흔히 50/50)을 두고 "v10의 답이 \textit{더 낫다}"는 가설을 실제 통계로 검증하는 \textbf{제품} 도구다. 사전 등록한 메트릭, 검정력 분석, 신규성 효과를 기다리는 인내와 함께 좋아요 비율, 작업 완료, 유지율 같은 제품 결과를 측정한다. LLM 제품에서 A/B는 품질 주장을 정직하게 시험하는 곳이다. 실제 세션 승률이 어떤 오프라인 벤치마크보다 낫다. 하지만 후보가 이미 게이트와 카나리를 통과했다고 \textit{전제}한다. 안전한 모델 두 개 중 무엇이 존재해야 하는지 배우려고 A/B를 하지, 어느 것이 고장인지 찾으려고 하지 않는다. \textbf{섀도(다크 론치)}는 사용자 위험 0으로 스펙트럼을 완성한다. 실제 트래픽 복사본을 후보에 미러링하고 사용자는 기존 릴리스의 답만 본다. 섀도의 동작을 오프라인에서 비교해 노출 전에 충돌, 실제 혼합에서의 지연 시간, 출력 드리프트를 확인한다. Facebook의 유명한 2008년 Chat \textbf{다크 론치(dark launch)}는 이를 규모 있게 실행했다. 몇 주 동안 UI 없이 페이지가 조용히 연결을 열고 실제 트래픽 형태로 메시지 경로를 모사해 출시일 부하는 백엔드에 이미 익숙했다. 비용도 실제다. 미러링한 부분에 추론 연산이 \textasciitilde{}2× 들며 제12장의 가격으로 트래픽 10\% 섀도는 10시간마다 레플리카 시간 하나 정도다. \textbf{부작용은 스텁 처리}해야 한다. 결제에 쓰고 이메일을 보내거나 공유 캐시의 KV를 변경하는 섀도는 연습이 아니라 둘째 프로덕션이다. 첫 vLLM 배포, 새 기반 모델, 엔진 메이저 버전처럼 \textit{계열 경계}에서 섀도는 비용만큼 가치가 있다. 카나리 10\%도 진정으로 알 수 없는 동작에 실제 사용자 10\%를 노출하기 때문이다.

교재의 기본 사다리를 그림 13-4 아래 띠에 따라 한 줄씩 정리하자. 실패가 시작 형태인 코드/설정 변경에는 \textbf{롤링}만 쓴다. 예외 없이 \textit{모든} 모델 릴리스에는 \textbf{자동 분석이 있는 카나리}를 적용한다. 이 교재 전체에서 절감한 피해 사용자-분당 가장 저렴한 제도다. 릴리스가 계열 경계를 넘으면 \textbf{섀도를 먼저} 쓰고, 게이트와 카나리를 통과한 후보 두 개가 성능으로 겨룰 때 \textbf{A/B}를 쓴다. 압박에서도 사다리가 동작하게 하는 상위 규칙은 \textbf{전략이 영웅적 선택이 아니라 파이프라인 단계}라는 것이다. 화요일 오후든 금요일 18:07이든 경로 3은 같은 YAML을 실행한다. 바로 그것이 목적이다.

\begin{calloutbox}{tbbpurple}{tbbpurplefill}
{\normalsize\bfseries\color{tbbpurple}생각해 보기}\par\vspace{3pt}
\begin{tbbitemize}
  \item 카나리는 통과했지만 일주일 뒤 사용자가 불만을 제기한다. 30분 카나리가 구조적으로 잡을 수 없는 세 실패 유형인 느리게 진행되는 품질 드리프트, 저빈도 슬라이스, 신규성 효과를 나열하고 각각을 잡는 §13.4 또는 제14장의 도구를 제시하라.
  \item 세션 해시는 대화의 10\%를 카나리로 보낸다. 하루 대화가 200개인 헤비 사용자는 "항상 이상한 버전이 걸린다"고 불평한다. 세션별 해시와 사용자별 해시에서 확률을 계산하고 분할 키를 다시 설계하라. 수정이 분석의 표본 독립성에 어떤 비용을 부과하는가?
  \item 교재 프로젝트의 섀도 하네스를 자세히 설계하라. 트래픽은 어디에서 분기하는가(LB, 사이드카, 로그의 비동기 탭)? 무엇을 스텁 처리하고 무엇을 비교하며, "카나리 준비 완료"를 결정하는 단 하나의 수치는 무엇인가? §12 비용 항목도 포함하라.
  \item PM은 "작은 프롬프트 수정"에 카나리를 건너뛰고 싶어 한다. "작은 설정" 변경인 사례 13-A, 가드레일 표, 한계 비용이 거의 0인 카나리 비용을 이용해 두 문장으로 답하라. 파이프라인은 어차피 카나리를 실행한다. 카나리를 생략해도 실제로 정당한 변경 클래스 하나도 제시하라. 힌트는 §13.4의 첫 단이다.
\end{tbbitemize}
\end{calloutbox}

\section{관리형 엔드포인트: 완성된 오퍼레이터 사다리}

이 장에서 구축한 모든 것을 누군가는 판매한다. AWS의 \textbf{SageMaker 엔드포인트}와 GCP의 \textbf{Vertex AI 엔드포인트}, 그리고 모든 클라우드의 유사 서비스는 제12장에서 시작한 오퍼레이터 사다리의 마지막 단이다. 직접 구축 방식에서는 사용자가 모든 것을 운영하고, 플랫폼에서는 KServe 오퍼레이터가 \textit{사용자} 클러스터의 서빙 배관을 운영하며, \textbf{관리형}에서는 공급자의 오퍼레이터와 클러스터와 호출기를 쓴다. 모델 산출물을 업로드하고 인스턴스 유형을 고르면 공급자가 제5–12장에서 구축한 것을 만들어 낸다. 프로비저닝, 오토스케일링, 상태 기반 라우팅, 로깅, 패칭에 더해 이 장과 관련된 \textbf{프로덕션 변형 / 트래픽 분할}을 제공한다. 카나리가 필드가 되어 SageMaker는 \code{InitialVariantWeight}, Vertex는 \code{trafficSplit \{v9: 90, v10: 10\}}을 쓴다. 승인 워크플로가 있는 \textbf{모델 레지스트리}와 경보 시 자동 롤백하는 \textbf{배포 가드레일}, 즉 §13.4의 분석을 체크박스로 판매한다. 용어가 일대일로 대응하므로 이 절은 짧다. 조절점을 직접 만들었으니 각 조절점의 역할을 이미 안다.

\begin{tbbfigure}
  \adjustbox{max width=0.94\linewidth,max totalheight=0.46\textheight}{\includegraphics{figures/fig13-6_managed.png}}\\[3pt]
  \figcaption{그림 13-6  사다리의 마지막 단. 관리형 엔드포인트는 배관을 운영하고 이 장의 전략을 필드로 판매한다. 하지만 게이트, 평가 세트, 레지스트리 계보, SLO 문장, 단위 경제성 원장은 어느 단에서도 영원히 사용자 몫이다.}
\end{tbbfigure}

이제 가진 도구로 정직한 절충의 가격을 매겨 보자. \textbf{얻는 것}은 운영 팀이 0명이어도 오늘 오후에 프로덕션 모델을 실행하고, 규정 준수 점검표가 공급자 인증을 물려받으며, 차별화되지 않는 무거운 작업인 노드 풀, 드라이버, CVE 패칭, §12.3의 콜드 스타트 엔지니어링이 다른 사람의 3분기 일이 된다는 점이다. \textbf{지불하는 것}은 원시 VM보다 인스턴스 시간당 붙는 \textbf{마진}이다. 흔히 수십 퍼센트이므로 어느 방향의 주장도 믿기 전에 자신의 사용률로 §12.5 원장을 계산해야 한다. 서빙 스택 제어도 줄어든다. 엔진 버전, §9/§11 조절점, 사용자 정의 커널이 공급자 일정에 맞춰 도착하고, 자신의 정책이 아니라 공급자의 콜드 스타트와 0으로 스케일링 정책을 따른다. 이그레스, IAM 연결망, 재학습 역량의 퇴화라는 중력도 생긴다. 결정은 정장을 입은 제9장의 벤치마크 규칙이다. \textit{홍보 자료가 아니라 자신의 사용률과 팀을 기준으로 사다리 가격을 계산하라.} 작은 팀 + 적은 모델 + 스파이크 트래픽은 관리형 쪽이고, 안정적 규모 + 비용 압박 + 엔진 수준 튜닝은 아래 단 쪽이다. 이 교재의 프로젝트는 의도적으로 후자다. 어느 단에서도 절대 이전되지 않는 자산이 있다. \textbf{게이트와 허용 오차, 평가 픽스처, 레지스트리의 계보와 감사, 근거가 있는 SLO 문장, \$/M-token 원장}이다. 공급자는 모델을 서빙하지만 \textit{이} 모델이 안전하고 가격이 적절하며 정직하게 측정되었다고 약속할 수 있는 것은 사용자뿐이다. 눈치챘겠지만 이 목록이 바로 이 교재다.

"원장을 계산하라"는 과제로만 남기지 않고 한 번 보여 줄 가치가 있다. §12.5의 계산으로 교재 플릿을 관리형 단에서 다시 가격 책정하자. L4급 인스턴스에 대표적인 \textbf{35\% 마진}이 붙고 공급자 오토스케일링이 같은 61 rep-h/day 계단을 따른다고 가정한다. 공급자의 0으로 스케일링 정책이 바닥 2개 논리를 이기는 경우는 드물어 그대로 유지한다. 61 × 31 × (\$0.80 × 1.35) ≈ \textbf{\$2,042/mo}이며 자체 호스팅 \$1,513보다 \$529/mo 비싸다. 관리형 단에는 약정을 적용할 수 있지만 사용자의 스팟 버스트 기법은 쓸 수 없어 현실적인 격차는 \textbf{\textasciitilde{}\$1,100/mo 대 \$877} 쪽으로 벌어진다. 프리미엄이 사는 것을 유일하게 정직한 단위로 가격을 매기면 §12.2–12.3을 구축하지 않고 §13.4를 체크박스로 사용하며 패칭 반복 작업을 피하는 \textbf{분기당 엔지니어 약 1주}다. 모델 세 개를 가진 2인 팀에는 이 거래가 흔히 \textit{옳다.} 인기 모델 하나에 엔진 수준 튜닝을 하고 모든 조절점을 이제 소유한 이 교재 프로젝트에는 그렇지 않다. 중요한 것은 판정 자체가 아니라 홍보 자료 없이 계산 네 줄로 판정했다는 점이다. 14주간 측정한 수치는 바로 이를 위한 것이다.

\begin{calloutbox}{tbbpurple}{tbbpurplefill}
{\normalsize\bfseries\color{tbbpurple}생각해 보기}\par\vspace{3pt}
\begin{tbbitemize}
  \item 롱테일에 0으로 스케일링을 포함하고 인스턴스 마진 35\%인 관리형 엔드포인트에서 기록 12-6 원장을 다시 작성하라. 어느 플릿 크기에서 마진이 연간 SRE 한 달을 넘으며, GPU 가격이 아니라 이것이 실제 손익분기 질문인 이유는 무엇인가?
  \item 규제를 받는 고객은 "하이퍼바이저 아래에 공유 인프라가 없어야 한다"고 요구한다. 제2장의 격리 용어인 패스스루, MIG, 멀티테넌트 컨트롤 플레인으로 사다리의 각 단을 요구사항과 비교하라. 어느 단이 살아남고 비용은 얼마인가?
  \item 현재 파이프라인 경로 1–3은 kubectl을 대상으로 한다. 배포 대상이 Vertex 엔드포인트로 바뀌면 정확히 어느 단계가 바뀌는지 나열하라. 힌트: 경로 3의 동사만 바뀐다. 검사한 것을 무효로 만들지 않으려면 삼중항, 게이트 보고서, 레지스트리 항목 중 무엇이 바뀌지 않아야 하는가?
  \item 관리형 가드레일은 인프라 메트릭인 CloudWatch 경보에 따라 자동 롤백한다. 금요일 버그를 이용해 공급자 가드레일과 §13.4의 가드레일 \textit{메트릭} 사이의 격차를 설명하고 이를 메울 가장 작은 Lambda/Cloud Function을 설계하라. 무엇을 읽고 계산하고 호출하는가?
\end{tbbitemize}
\end{calloutbox}

\section{실습: 파이프라인을 만들고 프로덕션을 깨뜨리는 데 실패하기}

강의 계획서는 제출물을 \textbf{GitHub Actions로 모델 자동 배포 파이프라인 구축}이라고 명시한다. 이 실습의 채점 철학은 §13.3에서 가져온다. 파이프라인은 무엇을 \textit{거부하는지}로 증명한다. 실습 자료 뼈대에서 프로젝트 저장소에 세 경로를 연결하고 종단 간 녹색 경로 하나를 만든다. 이후 게이트와 카나리에서 한 번씩 자신의 프로덕션을 \textbf{두 번 공격하고 두 번 실패}해 점수 대부분을 얻는다. GPU 러너 시간 예산은 \textasciitilde{}\$2이고 T4는 자체 호스팅 러너와 서빙 노드 역할을 겸한다. CPU 폴백 경로의 팀은 작은 모델로 모든 단계를 실행한다. 파이프라인은 이를 알지도 신경 쓰지도 않으며 그것이 바로 핵심이다.

\begin{tbbfigure}
  \adjustbox{max width=0.94\linewidth,max totalheight=0.46\textheight}{\includegraphics{figures/fig13-7_lab.png}}\\[3pt]
  \figcaption{그림 13-7  실습. 왼쪽은 교재에서 작성한 모든 것을 트리 하나에 모은 저장소 뼈대다. 오른쪽은 첫 PR부터 연습한 롤백까지 여섯 단계다. 아래의 제출물은 실행 링크이며 붉은 실행이 녹색 실행보다 중요하다.}
\end{tbbfigure}

\subsection*{여섯 단계}

\begin{tbbitemize}
  \item \textbf{① 경로 1을 연결한다.} 뼈대(app/, 골든 요청이 있는 eval/, gate.py, registry.py, Dockerfile, 워크플로 파일 두 개)를 포크한다. 간단한 PR을 열고 린트, 단위 테스트, \textbf{다이제스트로} 푸시하는 이미지 빌드, CPU에서 실행하는 빠른 게이트(골든 요청 + 200표본 평가)를 지켜본다. 기록 12-3처럼 로그를 읽는다. 이제 모든 단계가 이해되어야 한다.
  \item \textbf{② 러너를 등록하고 경로 2를 연결한다.} 실습 자료 스크립트가 T4를 \code{gpu} 태그 자체 호스팅 러너로 등록한다. systemd 서비스 하나이며 스팟에 적합하다. PR을 병합하고 내보내기 → 동등성 → 양자화+보정 → \textbf{게이트} → 스모크 스윕 → \code{registry.py register v10 --stage staging}을 지켜본다. 레지스트리 항목에 삼중항, 게이트 보고서, 계보 다이제스트가 있는지 확인한다.
  \item \textbf{③ OIDC로 스테이징에 배포한다.} 어디에도 장기 키를 두지 않는다. 워크플로는 OIDC(\code{permissions: id-token: write})로 클라우드 역할을 맡은 뒤 \code{kubectl set env … MODEL\_URI=\$(registry.py resolve chat/\allowbreak staging)}를 실행한다. 제6장의 환경 변수 배포(기록 6-7)를 이제 범위가 제한되고 몇 분만 유효한 자격 증명을 가진 로봇이 입력한다.
  \item \textbf{④ 경로 3인 카나리를 연결한다.} \code{deploy.yaml}이 \code{canaryTrafficPercent: 10} 또는 실습 자료의 원시 Deployment 가중 Service 변형을 패치한 뒤 40줄 분석 루프를 실행한다. 60초 창 다섯 개에서 SLO 조항 + 가드레일 네 개(답변 길이, 거부, 형식, 반복)를 동시 기준선과 비교한다. 판정에 따라 별칭을 전환해 승격하거나 카나리 0\%로 후퇴하고 그래프가 있는 PR 댓글을 남긴다. 무해한 v10으로 녹색 실행을 한 번 만든다.
  \item \textbf{⑤ 의도적으로 두 번 깨뜨린다.} (a) 골든 요청 하나의 기대 형식을 오염시키고 푸시한다. \textbf{경로 1에서 게이트가 거부}하며 회색 병합 버튼을 스크린샷으로 남긴 뒤 되돌린다. (b) 금요일을 재현한다. \code{max\_tokens: 256} 설정을 배포한다. 게이트는 긴 글을 생성하지 않아 통과한다. 승격하면 답변 길이 가드레일 때문에 \textbf{카나리가 약 3번째 창에서 후퇴}한다. PR 댓글을 저장한다. 자신이 만든 메커니즘으로 사례 13-A를 물리친 결과다.
  \item \textbf{⑥ 새벽 2시 롤백을 연습한다.} 좋은 v10을 승격한 상태에서 \code{registry.py rollback chat --to v9}를 실행하고 별칭 전환부터 v9이 처음 요청을 서빙할 때까지 \textbf{시간을 잰다.} 목표는 3분 아래다. 하한은 §12.3의 탄생이며 사전 워밍업이 이를 어떻게 바꾸는지 런북에서 논한다. 배포, 관찰, 후퇴, 롤백을 담은 한 쪽짜리 런북을 작성해 파이프라인을 \textit{만들지 않은} 팀원도 운전할 수 있게 한다. 그 팀원이 서명한다.
\end{tbbitemize}

\begin{calloutbox}{tbbred}{tbbxFEF2F2}
{\normalsize\bfseries\color{tbbx991B1B}상자 13-1  배포 위험 신호 — 공개 점검표}\par\vspace{3pt}
\begin{tbbitemize}
  \item \textbf{노트북에서 프로덕션으로 kubectl 실행} — 파이프라인만 작성자이고 사람은 PR과 승인을 작성한다.
  \item \textbf{배포 경로 어디서든 `:latest` 또는 다시 태그할 수 있는 참조 사용} — 다이제스트가 아니면 배포하지 않은 것이다. 제4장의 규칙은 계속 적용된다.
  \item \textbf{"비상용" 게이트 생략 플래그} — 비상 경로는 더 빠른 카나리이지 검사를 건너뛰는 것이 아니다.
  \item \textbf{자동 분석 없는 카나리} — 아무도 판정하지 않는 10\% 트래픽은 카나리가 아니라 사용자의 10\%가 베타 테스트하는 것이다.
  \item \textbf{연습하지 않은 롤백} — 금요일의 12분은 프로덕션에서 첫 연습을 한 대가였다.
  \item \textbf{릴리스 밖의 설정} — 변경해도 버전이 오르지 않는다면 반드시 자신의 사례 13-A가 된다.
  \item \textbf{저장소 시크릿의 장기 클라우드 키} — OIDC가 있다. 유출된 키는 회사 로고를 붙인 Knight의 여덟째 서버다.
\end{tbbitemize}
\end{calloutbox}

\vspace{3.0pt}

\begin{calloutbox}{tbbteal}{tbbxF0FDFA}
{\normalsize\bfseries\color{tbbx115E59}상자 13-2  런북 템플릿 — 이 페이지를 그대로 사용하라}\par\vspace{3pt}
\textbf{배포} — main에 병합하면 경로 2가 vN을 staging에 자동 등록한다. 승격하려면 승격 PR(\code{registry.py promote chat vN --request})을 열고 승인 하나를 받아 병합한다. 경로 3이 카나리를 실행하고 직접 별칭을 전환한다. \textit{kubectl은 입력하지 않는다.}

\textbf{관찰} — 카나리 중에는 Actions 실행의 실시간 로그나 \code{watch registry.py status chat}을 본다. 후자는 별칭, 카나리 비율, 현재 창 판정을 보여 준다. 승격 후에는 제14장 대시보드를 본다. 그전까지는 /metrics + \code{eval/\allowbreak watch.py}의 가드레일 스크립트를 사용한다.

\textbf{후퇴}(카나리가 이상하지만 아직 실패 판정 전) — \code{gh run cancel <id>}로 분석을 중단하고 카나리를 0으로 만든다. 별칭은 움직이지 않았으므로 언제든 안전하다.

\textbf{롤백}(나쁜 릴리스가 이미 승격됨) — \code{registry.py rollback chat --to <last-good>}를 실행한다. \code{registry.py list chat}의 이전 \code{production} 행에서 대상을 찾는다. 전환 ≈ 0.4 s, 플릿 수렴 ≈ 2–3 min이다. §12.3의 탄생이 하한이고 사전 워밍업이 줄인다. 이어 릴리스의 레지스트리 항목에 사고 메모를 기록하고 수정 PR을 연다. 프로덕션을 손으로 패치하지 않는다.

\textbf{담당자} — 승격 승인자는 [팀 목록] 중 아무 두 명이다. 02:00 규칙에 따라 온콜은 혼자 롤백할 수 있다. 이미 모든 검사를 통과한 릴리스를 다시 가리키므로 롤백은 구조적으로 사전 승인되었다. 승격은 낮까지 기다린다.

\textit{파이프라인을 만들지 않은 팀원의 서명: \_\_\_\_\_\_\_\_\_\_\_ (서명이 테스트다).}
\end{calloutbox}

\vspace{3.0pt}

제출물은 채점자가 클릭할 수 있는 \textbf{스크린샷이 아닌 링크}다. 승격 릴리스로 끝나는 경로 1→2→3의 녹색 실행 하나, 증거가 있는 붉은 실행 두 개(게이트 거부, 카나리 후퇴), 프로덕션의 v9과 증거와 함께 거부된 오염 후보를 보여 주는 레지스트리 목록, 시간을 잰 롤백 그림, 서명된 런북이다. 일찍 끝난 팀은 30분 미러링 트래픽을 위한 섀도 단계를 추가하고 출력을 비교한다. 이어 제14장 첫머리를 읽어 보자. 카나리의 분석 창이 곧 영구적이 되고 \textbf{모니터링}이라는 이름을 얻는다.

\begin{calloutbox}{tbbpurple}{tbbpurplefill}
{\normalsize\bfseries\color{tbbpurple}생각해 보기}\par\vspace{3pt}
\begin{tbbitemize}
  \item ⑤(b)의 카나리는 3번째 창(\textasciitilde{}3 minutes)에서 금요일 버그를 잡았다. 10\% 트래픽과 측정한 요청률에서 노출을 계산하라. 이어 신중한 팀을 위해 첫 창 1\%, 지수 단계 증가(1→5→25→100), 각 단계 10창 체류로 다시 설계하라. 전체 승격까지 시간 비용은 얼마이며 언제 이런 신중함이 가치 있는가?
  \item 런북의 롤백 하한은 §12.3의 탄생(\textasciitilde{}150 s)이다. 팀원이 승격 후 한 시간 동안 이전 릴리스를 블루/그린 방식으로 따뜻하게 유지하자고 제안한다. §12의 계산으로 가격을 매기고 얻는 롤백 시간을 제시하라. 한 시간이 올바른 창인가? 제14장의 어떤 데이터로 조정할 수 있는가?
  \item 실습 레지스트리는 저장소에 있는 JSON 200줄이다. 팀 규모에서 먼저 나타날 실패 유형 세 가지인 동시 쓰기, 인증, 감사 변조 등을 나열하고 무엇이 가장 먼저 실제 레지스트리 서비스로 이동하게 만드는지 설명하라.
  \item 프로젝트의 경로 3 승인 정책을 작성하라. 누가 승격을 승인할 수 있고 PR에 게이트 보고서, 카나리 계획, 롤백 담당자 중 어떤 증거를 인라인으로 보여 주어야 하는가? 02:00에 승인자가 잠들어 있으면 어떻게 하는가? §13.1 생각해 볼 문제의 2주 릴리스 위원회를 다시 만들지 않으면서 Knight의 교훈을 인코딩하라.
\end{tbbitemize}
\end{calloutbox}

\namedsection{요약}{열다섯 가지로 정리한 이 장}

\qitem{1}{\textbf{잘못된 배포는 부주의한 사람보다 빠진 제도에서 비롯된다.} 사례 13-A에는 완벽한 모델, 릴리스 밖의 설정, 초록색 대시보드, 전체 사용자에게 영향을 준 37분, 발굴 작업이 필요했던 롤백이 있었다. 엔지니어 한 사람을 고치면 재발하지만 버전 관리 + 파이프라인 + 점진적 노출을 고치면 피해 사용자-분은 0.5로 줄어든다(74배).}

\qitem{2}{\textbf{Knight Capital은 이 문제의 대표 사례다.} 서버 여덟 대를 손으로 배포하다 한 대를 빠뜨렸고, 재사용된 플래그가 죽은 코드를 깨웠다. 45분 만에 4억 6천만 달러를 잃었으며 사고 중에는 어디에서 무엇이 실행되는지 아무도 몰라 첫 조치가 오히려 사태를 악화했다. §13의 모든 제목은 SEC 명령서에서 빠져 있던 제도 하나씩을 가리킨다.}

\qitem{3}{\textbf{배포 빈도는 허세가 아니라 안전이다(Flickr 2009 → DORA).} 드문 배포는 크고, 큰 배포는 연습되지 않았으며, 연습되지 않은 배포는 Knight가 된다. 네 숫자는 리드 타임(→22분), 배포 빈도(→병합마다), 장애 반경(→자동 후퇴를 포함한 10\%), 평균 복구 시간(MTTR, →2분짜리 별칭 전환)이다. 지루한 배포란 자동화되고, 버전이 붙고, 점진적이며, 되돌릴 수 있는 배포다.}

\qitem{4}{\textbf{릴리스는 증거를 갖춘 세 요소의 묶음이다.} 가중치 다이제스트 + 설정 다이제스트 + 이미지 다이제스트를 함께 고정하고, 허용 오차·슬라이스·평가 세트 다이제스트를 담은 게이트 보고서를 붙이며, 계보를 기록한다. 설정이 릴리스 밖에서 따로 이동했던 금요일의 버그는 이제 표현 자체가 불가능하다.}

\qitem{5}{\textbf{레지스트리는 바이트에 관한 약속을 담은 데이터베이스다.} 무엇이 실행 중이고 무엇이 안전하며 무엇이 바뀌었고 누가 만들었는지를 아무 정보 없는 터미널에서도 몇 초 안에 답할 수 있어야 한다. 제6장이 바이트를 보관했다면 레지스트리는 의미를 보관한다. MLflow든 클라우드 내장 기능이든 200줄짜리 JSON이든, 가치는 제품명이 아니라 스키마에 있다.}

\qitem{6}{\textbf{승격과 롤백은 같은 동작, 즉 포인터 이동이다.} 배포는 버전이 아니라 별칭(…/prod)을 참조한다. 전환은 0.4초, 플릿 수렴은 제5장의 롤아웃을 따라 약 2분이며 하한은 §12.3의 탄생 시간이다. 제4장 생각 상자의 답처럼 다시 태그하지 말고 다이제스트를 쓴다. 태그는 움직이지만 다이제스트는 움직이지 않는다.}

\qitem{7}{\textbf{파이프라인은 이미 가진 검사를 빠짐없이 치르기 위한 의식이다.} §7.6의 동등성, §10.3의 게이트, §9.5의 스윕, 제6장의 골든 픽스처를 모든 후보에 똑같이 실행하며 금요일이라고 건너뛰는 길은 없다. 제5장의 표현을 빌리면 원하는 상태를 작성하는 주체가 하나 더 생긴 셈이다.}

\qitem{8}{\textbf{모델 CI는 앱 CI와 네 가지가 다르다.} 판정은 버전이 붙은 평가 픽스처에 대한 허용 오차이고, 아티팩트는 GB 단위이므로 한 번 빌드한 다이제스트를 승격하며, 일부 단계에는 실리콘이 필요해 자체 호스팅 스팟 T4 러너를 쓴다. 마지막 검사는 실제 트래픽을 필요로 하는데, 이는 테스트가 아니라 배포 전략이다.}

\qitem{9}{\textbf{빨간 실행 결과가 제품이다.} 병합 버튼을 비활성화하는 게이트는 동료의 의견이 아니라 마흔 번의 초록 선례를 쌓은 검사다. 품질 논쟁이 인프라 문제가 되는 순간이다. 건너뛰지 말고, 검사와 배포 사이에 다시 빌드하지 말며, 첫날부터 cosign/SLSA로 아티팩트에 서명한다.}

\qitem{10}{\textbf{배포 전략은 답하려는 질문으로 고른다.} 롤링은 "부팅되는가?", 블루/그린은 플릿 비용 2배를 감수하고 "1초 안에 되돌릴 수 있는가?", 카나리는 "현실에서 망가졌는가?", A/B는 "사용자에게 더 나은가?", 섀도는 "실제로 내보냈다면 어떻게 행동했을까?"에 답한다. 섀도는 사용자 위험이 없지만 계산 비용이 2배이고 모든 부작용을 막아야 한다.}

\qitem{11}{\textbf{카나리는 Haldane의 새에 그 새를 은퇴시킨 감지기를 더한 것이다.} 실제 트래픽의 10\%를 대화 전체가 같은 모델로 가도록 세션 해시하고, 동시에 실행되는 기준선과 비교하며, 판정과 수초 안의 후퇴를 자동화한다. 게이트가 확인할 수 없는 실제 입력, 실제 통합(금요일의 문제), 실제 인프라, 실제 SLO를 검사한다.}

\qitem{12}{\textbf{가드레일 메트릭은 우수성이 아니라 행동을 지킨다.} 답변 길이, 거부율, 형식 유효성, 반복, 언어 일치는 "모델이 이상해졌다"를 60초 루프에서 저렴하게 감지하는 프록시다. SLO 조항은 시스템을, 가드레일은 행동을 지키고, 품질은 A/B가 판단한다.}

\qitem{13}{\textbf{A/B는 안전을 이미 확보했다고 전제하는 제품 도구다.} 며칠에서 몇 주 동안 균형 잡힌 버킷과 미리 등록한 제품 메트릭, 통계를 사용하며 게이트와 카나리를 통과한 후보에만 실행한다. A/B를 카나리와 혼동하면 안전은 느려지고 실험은 타당성을 잃는다.}

\qitem{14}{\textbf{관리형 엔드포인트는 사다리의 마지막 칸이다.} 변형과 트래픽 분할, 레지스트리, 가드레일을 체크박스로 제공한다. 같은 개념에 더 높은 가격을 내고 스택 통제권은 줄어든다. 자기 사용률로 가격을 계산해야 하며(§12.5의 원장), 어느 단계에서도 게이트, 픽스처, 계보, SLO 문장, 단위 경제성은 공급자에게 넘길 수 없다.}

\qitem{15}{\textbf{실습의 채점 기준이 이 장의 주장이다.} 초록 실행 한 번, 빨간 실행 두 번(게이트 거부와 카나리 후퇴), 시간을 잰 롤백 연습, 작성자가 아닌 사람도 실행할 수 있는 런북을 요구한다. 파이프라인은 무엇을 허용하느냐보다 무엇을 거부하느냐로 입증된다. 빨간 실행은 일어나지 않은 금요일 사고다.}

\namedsection{연습문제}{}

\subsection*{A. 객관식}

\qitem{1}{사례 13-A의 버그(max\_tokens=256)가 오프라인 게이트를 통과한 이유는?}

\begin{tbboptions}
  \item ① 게이트의 허용 오차가 너무 느슨했다.        ② 게이트가 교사 강제 평가만 채점하고 실제 배포한 설정으로 생성하지 않아 릴리스가 아니라 모델만 검사했다.
  \item ③ 평가 세트가 너무 작았다.        ④ 게이트가 잘못된 GPU에서 실행되었다.
\end{tbboptions}

\qitem{2}{이 장에서 말하는 "릴리스"는 무엇인가?}

\begin{tbboptions}
  \item ① 모델 가중치        ② Docker 이미지
  \item ③ 게이트 보고서와 계보를 포함해 고정한 가중치 + 설정 + 이미지 다이제스트의 세 요소        ④ prod 별칭이 가리키는 무엇이든
\end{tbboptions}

\qitem{3}{레지스트리 별칭을 이용한 롤백이 빠르고 안전한 주된 이유는?}

\begin{tbboptions}
  \item ① 준비 상태 프로브를 건너뛴다.        ② 결정, 무엇이 안전한지에 대한 지식, 실행 수단이 연습된 포인터 이동 하나로 합쳐지고, 수렴은 평소 제5장 롤아웃과 같기 때문이다.
  \item ③ 이미지 가져오기를 피한다.        ④ 별칭이 로드 밸런서를 우회한다.
\end{tbboptions}

\qitem{4}{Knight Capital의 첫 조치, 즉 정상 서버 일곱 대에서 새 코드를 롤백한 일이 상황을 악화한 이유는?}

\begin{tbboptions}
  \item ① 롤백이 너무 느렸다.        ② 재사용된 FLAG와 업데이트되지 않은 여덟 번째 서버의 죽은 코드가 트리거였고, 어디에서 무엇이 실행되는지 보여 주는 권위 있는 기록이 없었다.
  \item ③ 시장이 불리하게 움직였다.        ④ 새 코드는 사실 문제가 없었다.
\end{tbboptions}

\qitem{5}{모델 CI가 아티팩트를 한 번만 빌드한 뒤 환경 사이에서 다이제스트를 승격하는 이유는?}

\begin{tbboptions}
  \item ① 다시 빌드하면 느리다.        ② 검사와 배포 사이의 재구축은 모든 검사를 무효화하므로 게이트를 통과한 바이트가 실제 배포되는 바이트와 같아야 한다.
  \item ③ 레지스트리가 빌드마다 요금을 받는다.        ④ 다이제스트의 압축률이 더 좋다.
\end{tbboptions}

\qitem{6}{카나리가 어제의 메트릭이 아니라 동시에 실행되는 기준선과 비교하는 이유는?}

\begin{tbboptions}
  \item ① 어제의 데이터는 삭제되었다.        ② 트래픽에는 하루 주기의 성격이 있어(제12장) 같은 시간의 대조군만 릴리스 변화와 시각 효과를 분리할 수 있다.
  \item ③ Prometheus 보존 기간이 짧다.        ④ 오토스케일링 때문에 기준선이 드리프트한다.
\end{tbboptions}

\qitem{7}{카나리와 A/B의 차이를 올바르게 설명한 것은?}

\begin{tbboptions}
  \item ① 둘은 같은 말이다.        ② 카나리는 안전을 묻는 도구(몇 분, "망가졌는가?", 자동 판정)이고 A/B는 제품을 묻는 도구(몇 주, "더 나은가?", 통계)이며, A/B는 게이트와 카나리를 이미 통과했다고 전제한다.
  \item ③ A/B는 인프라용이고 카나리는 모델용이다.        ④ 카나리가 A/B보다 더 많은 트래픽을 필요로 한다.
\end{tbboptions}

\qitem{8}{섀도 배포에서 피할 수 없는 비용 두 가지는?}

\begin{tbboptions}
  \item ① 느린 롤아웃과 큰 이미지        ② 사용자가 체감하는 위험과 호출기 부하
  \item ③ 미러링한 트래픽 조각에 드는 약 2배의 계산 비용과 쓰기·청구·이메일을 비롯한 모든 부작용 차단        ④ 추가 레지스트리와 더 긴 게이트
\end{tbboptions}

\subsection*{B. 단답형}

\qitem{1}{사례 13-A에서 빠진 세 가지 제도와 각각을 구축하는 절을 말하자.}

\qitem{2}{평가 세트 자체에 버전과 다이제스트를 고정해야 하는 이유는 무엇인가? 누군가 그 자리에 "어려운 예제 몇 개만" 더하면 무엇이 조용히 깨지는가?}

\qitem{3}{모델 파이프라인의 세 경로를 트리거, 러너, 최종 상태와 함께 나열하자.}

\qitem{4}{생성형 카나리의 가드레일 메트릭 네 가지를 들고, 가드레일이라는 부류의 한 문장짜리 근거("행동의 동일성")를 설명하자.}

\qitem{5}{채팅 워크로드의 카나리 분할에 세션 해시를 사용하는 이유는 무엇인가? 요청마다 나눌 때 생기는 UX 실패와 메트릭 실패를 각각 말하자.}

\qitem{6}{긴급 배포 원칙("게이트를 건너뛰지 않는다. 긴급 경로는 더 빠른 카나리다")을 말하고, Knight 사고가 이 원칙을 어떻게 뒷받침하는지 설명하자.}

\subsection*{C. 논술형}

\qitem{1}{사례 13-A의 비난 없는 사후 분석을 작성하자. 시간선, 영향(피해 사용자-분), 사람 이름이 아니라 빠진 제도로 표현한 근본 원인, 담당자와 날짜를 포함해 §§13.2–13.4에 연결한 개선 조치를 담는다. 대부분의 사후 분석이 빼먹는 문단도 추가하자. 하지 않을 개선은 무엇이며, 재발 가능성에 비해 비용이 큰 이유는 무엇인가?}

\qitem{2}{다음 주장에 찬성하거나 반박하자. "레지스트리가 하중을 받는 핵심 제도이고 파이프라인과 카나리는 그 주위의 자동화다." 금요일 사고(레지스트리만으로 막을 수 있었는가?), Knight(가장 먼저 없었던 것은 무엇인가?), 새벽 2시 테스트를 근거로 쓰고, 스타트업에서 가장 먼저 만들 제도 하나와 나머지를 미뤄 받아들일 실패를 밝힌다.}

\qitem{3}{배포 전략을 설계하자. 팀은 (a) 매주 기반 모델 업그레이드, (b) 매일 프롬프트/설정 수정, (c) 분기마다 엔진 메이저 버전을 배포한다. 각각에 대해 경로, 전략 단계, 대기 시간, 가드레일을 포함한 전체 경로, 장애 반경 산술, 하한을 포함한 롤백 계획을 제시한다. 세 경로가 달라지는 모든 지점을 정당화한다.}

\qitem{4}{배포 관점의 자체 구축 대 외부 도입 논술이다. 자기 프로젝트 팀의 규모, 모델, 중간 점검에서 확인한 트래픽을 기준으로 KServe를 대상으로 하는 경로 1–3 파이프라인과 레지스트리 + 가드레일을 갖춘 SageMaker 엔드포인트를 비교하자. §12.5 원장으로 마진을 계산하고, 그대로 이전되는 요소(세 요소 묶음, 게이트, 픽스처)를 나열한 뒤 도입·보류·혼합 가운데 하나를 권고한다. 결론을 뒤집을 메트릭 하나로 끝맺는다.}

\namedsection{정답과 해설}{}

\subsection*{A. 객관식}

\qitem{1}{\textbf{②} 게이트의 사각지대는 구조적이었다. 게이트는 가중치를 평가했지만 실패는 실제 배포된 가중치 + 설정으로 이루어진 릴리스에 있었다. "배포하는 것을 검사하라"는 원칙 때문에 경로 2가 조립된 릴리스에 게이트를 적용하고 §13.4가 다시 카나리로 검사한다. (§13.1, §13.3)}

\qitem{2}{\textbf{③} ④는 순환 정의다. 별칭은 릴리스를 가리킬 뿐 릴리스를 정의하지 않는다. 이동하는 단위는 세 요소 + 증거이며, 이를 고정했기 때문에 금요일과 같은 버그를 표현할 수 없게 된다. (§13.2)}

\qitem{3}{\textbf{②} 속도보다 중요한 것은 결정 + 지식 + 수단이 한곳에 모인 연습된 일체성이다. 플릿 쪽 하한은 §12.3의 탄생 시간이므로 런북에서 이 시간을 측정한다. (§13.2)}

\qitem{4}{\textbf{②} 교훈은 절차보다 지식에 관한 것이다. "어디에서 무엇이 실행되는가"에 대한 권위 있는 기록이 없으면 조치는 추측이 되고, 실제로 그 추측이 트리거를 다시 작동시켰다. 레지스트리가 바로 그 권위다. (§13.1–13.2)}

\qitem{5}{\textbf{②} 파이프라인을 통과하는 동안 신원을 유지하는 것이 신뢰 사슬의 전부다. 게이트 → 레지스트리 → 배포가 같은 다이제스트를 참조하지 않으면 초록 검사는 실제로 배포되지 않는 바이트를 인증한다. (§13.3)}

\qitem{6}{\textbf{②} 제12장의 일간 곡선은 교란 변수다. 14:00과 03:00의 비교는 후보가 아니라 시계를 측정한다. 동시 대조군은 가장 저렴하면서 타당한 실험 설계다. (§13.4)}

\qitem{7}{\textbf{②} 서로 다른 질문에는 서로 다른 도구를 쓴다. 게이트 → 카나리 → A/B 순서는 과학보다 안전이 먼저라는 뜻이다. (§13.4)}

\qitem{8}{\textbf{③} 미러링한 조각은 실제 추론 자원을 소비하고(제12장에서 가격을 계산했다), 살아 있는 부작용을 가진 섀도는 연습이 아니라 두 번째 프로덕션이다. (§13.4)}

\subsection*{B. 단답형}

\qitem{1}{버전이 붙은 릴리스(§13.2, 레지스트리와 세 요소), 배포할 것을 검사하는 파이프라인(§13.3, 세 경로), 자동 후퇴를 포함한 점진적 노출(§13.4, 카나리)이다. (§13.1)}

\qitem{2}{게이트 판정은 "이 세트에서 Δ와 허용 오차의 비교"이므로 세트는 측정의 절반이다. 제자리에서 바꾸면 과거 모든 PASS/FAIL의 의미가 말없이 달라져 서로 다른 시험의 점수가 되고, 회귀 비교와 계보 고정이 깨진다. 새 예제에는 새 평가 세트 버전을 부여하고 기준선을 명시적으로 다시 정한다. (§13.3)}

\qitem{3}{경로 1: 모든 PR / CPU 러너 / 위생 검사 + 빠른 게이트 → 초록 검사. 경로 2: main 병합 / 자체 호스팅 GPU 러너 / 내보내기→양자화→게이트→스모크 테스트 → 증거가 붙은 등록된 스테이징 릴리스. 경로 3: 사람의 승인 / OIDC로 클러스터 접근 / 카나리 분석 → 자동 승격(별칭 전환) 또는 자동 후퇴. (§13.3)}

\qitem{4}{답변 길이 분포, 거부/사과 비율, 형식 유효성(JSON 파싱), 반복 점수에 언어 일치를 더할 수 있다. 근거는 배포의 귀무가설이 행동의 \textbf{동일성}이라는 데 있다. 가드레일은 "이상해짐"을 저렴하게 감지하고 실제 품질 판정은 A/B까지 기다린다. (§13.4)}

\qitem{5}{UX 측면에서는 한 대화가 여러 모델 사이를 오가며 중간에 성격이 흔들린다. 메트릭 측면에서는 같은 세션의 턴이 양쪽 집단에 들어가 두 모집단을 오염시키고 비교의 독립성을 무너뜨린다. 대화 ID를 해시한다(제11장의 어피니티 장치). (§13.4)}

\qitem{6}{긴급 상황은 \textbf{속도}만 바꾼다. 카나리 창을 줄이고 롤백을 미리 준비하며 호출 범위를 넓힐 수 있지만 \textbf{검사}는 바꾸지 않는다. 압박받을 때 검증되지 않은 변경이 가장 위험하기 때문이다. Knight의 배포는 마감에 쫓겨 수동으로 진행했고 검증도 롤백 준비도 없었다. (§13.3, §13.1)}

\subsection*{C. 논술형 채점 기준}

\qitem{1}{다음을 확인한다. 프로브와 대시보드가 무해하게 초록색이었던 두 지점을 포함한 분 단위 시간선, 막연한 표현이 아니라 피해 사용자-분 37로 계산한 영향, 릴리스 밖 설정을 정확히 지목하면서 사람 대신 제도로 표현한 원인, 각 절에 연결하고 이 실패 부류를 잡는 카나리 가드레일을 명시한 개선 조치, 그리고 납득할 수 있는 "하지 않을 일"이다. 예를 들어 전체 블루/그린은 §12의 비용을 계산한 뒤 한계 비용이 거의 없는 카나리에 비해 과하다고 기각할 수 있다. 비난하지 않으면서도 개선 조치의 \textbf{책임자}를 분명히 한 답이 최고 점수를 받는다. (§§13.1–13.4)}

\qitem{2}{좋은 글은 각 도구가 서로 다른 실패 부류에 답함을 알아본다. 레지스트리는 Knight에게 없었던 "어디에 무엇이 있는가"를 제공하는 지식의 바닥이지만, 금요일에는 카나리가 필요했다. 레지스트리는 나쁜 릴리스도 완벽하게 기록할 뿐이다. 무엇을 먼저 만들지에 대해서는 레지스트리(가장 저렴하고 연습된 롤백을 가능하게 함)와 카나리(투입 노력당 피해 사용자-분을 가장 크게 줄임) 모두 타당하다. 다만 미뤄서 감수하는 실패를 솔직히 밝혀야 한다. 레지스트리 우선이면 금요일 사고를 한 번 더, 카나리 우선이면 느린 롤백을 감수한다. (장 전체)}

\qitem{3}{세 경로는 실제로 달라야 한다. (a) 모델 계열 경계에서 섀도 + 긴 카나리 + 품질 주장을 위한 전체 A/B, (b) 경로 1 + 롤링 또는 빠른 카나리(설정 변경이지만 사례 13-A 때문에 가드레일 카나리는 유지), (c) 엔진에는 즉시 되돌릴 수 있고 인프라 위험에 맞지만 §12 기준 창 비용이 2배인 블루/그린 + 수치 검사용 섀도(§7.6/§10.4의 순서 변경 주의)다. 전략 이름만 나열하기보다 구체적인 체류 시간, 장애 반경 수치, 롤백 하한(§12.3)을 제시한 답을 높이 평가한다. (§13.4)}

\qitem{4}{원장에는 마진 × 인스턴스 시간과 절약한 SRE 시간을 함께 적어야 한다. 이전 가능한 아티팩트, 즉 세 요소, 게이트, 픽스처, 레지스트리 스키마와 그대로인 파이프라인 경로 1–2를 밝힌다. 바뀌는 것은 경로 3의 동사가 공급자 API 호출이 된다는 점뿐이다. 판정은 팀의 형태에 연결한다. 결론을 뒤집는 좋은 메트릭에는 모델×팀 증가(→관리형을 더 일찍 선택), 규모가 커졌을 때의 사용률(→자체 호스팅), 규정 준수 요구가 있다. 산술이 일관되면 어느 결론도 만점을 받을 수 있지만 홍보 문구를 옮겨 적으면 점수를 받지 못한다. (§13.5, §12.5)}

\namedsection{더 읽을거리}{}

\textbf{[이번 주]} 표시는 실습에 직접 도움이 되는 자료이고, \textbf{[N주 차]} 표시는 뒤에서 다시 가치가 커지는 자료다.

\begin{tbbitemize}
  \item \textbf{[이번 주] GitHub Actions 문서 — 워크플로 문법, 자체 호스팅 러너, 클라우드 공급자용 OIDC} — 실습을 떠받치는 세 페이지다. OIDC 안내서를 적용하면 프로젝트에 남은 마지막 장기 키를 없앨 수 있다.
  \item \textbf{[이번 주] MLflow Model Registry 문서(별칭, 단계, 계보)} — 이제 읽을 때가 된 제6장의 더 읽을거리다. UI를 갖춘 터미널 기록 13-1의 레지스트리 스키마라고 생각하자. 제품이 아니라 스키마로 읽는다.
  \item \textbf{[이번 주] Argo Rollouts 문서(카나리, 분석 템플릿)} — 제5장에서 예고한 후속 내용이다. 단계별 트래픽 가중치, Prometheus 기반 분석, 자동 롤백으로 §13.4를 컨트롤러로 구현한다.
  \item \textbf{Knight Capital 관련 SEC Administrative Order(2013)와 "Knightmare" 사후 분석 자료} — §13.1의 1차 자료다. 명령서의 조사 결과를 그림 13-3과 나란히 읽고 빠진 경로를 세어 보자.
  \item \textbf{J. Allspaw \& P. Hammond, "10+ Deploys per Day: Dev and Ops Cooperation at Flickr"(Velocity 2009)} — DevOps의 씨앗이 된 발표다. 지금도 온라인에서 볼 수 있다. 2009년의 고통 가운데 이 과정의 각 장이 제도로 바꿔 놓은 것이 얼마나 많은지 살펴보자.
  \item \textbf{Google SRE Workbook의 "Canarying Releases"} — §13.4를 성숙한 수준으로 다룬다. 모집단 선택, 메트릭 선택, 작은 카나리의 통계, 가드레일과 SLO의 차이를 프로덕션 수준의 깊이로 설명한다.
  \item \textbf{DORA의 "Accelerate: State of DevOps" 보고서} — §13.1의 네 숫자를 뒷받침하는 측정 프로그램이다. 거울처럼 사용해 과정 프로젝트를 솔직하게 채점하고 가장 약한 지표를 찾자.
  \item \textbf{[이번 주, 가볍게] 대형 OSS 모델 저장소의 Actions 탭(vLLM이 좋은 예다)} — 초록 실행 하나와 빨간 실행 하나를 처음부터 끝까지 읽어 보자. 이제 그 언어를 이해할 수 있다. 무엇에 게이트를 적용하고 무엇은 아직 손으로 배포하는지 확인한다.
\end{tbbitemize}

\vspace{5.0pt}

\begin{calloutbox}{tbbblue}{tbbbluefill}
{\normalsize\bfseries\color{tbbblue}다음 장 — 제14장: 모니터링과 신뢰성}\par\vspace{3pt}
카나리의 분석 루프는 다섯 창을 실행한 뒤 멈췄다. 하지만 그 루프가 물었던 "TTFT는 정직한가? 행동이 드리프트했는가?"라는 질문은 승격과 함께 만료되지 않는다. 프로덕션은 아무도 듣지 않는 청중을 향해 영원히 답을 내놓는다. 제14장은 그 청중을 상시 체제로 만든다. \textbf{Prometheus와 Grafana}는 제11장의 /metrics 줄과 이 장의 가드레일을 강의계획서에서 약속한 서빙 대시보드로 바꾼다. \textbf{경보}는 SLO 문장을 규칙으로 만들어 막연한 느낌이 아니라 소진율에 따라 호출기가 울리게 한다. \textbf{드리프트 탐지}는 어느 배포 게이트도 얼려 둘 수 없는 것, 즉 사용자 집단 하나씩 평가 세트에서 멀어지는 입력 분포를 감시한다. §12.4에서 예고한 전쟁 이야기도 돌아온다. 시간 분할 GPU의 설명되지 않는 p99를 실제 대시보드에서 해결한다. 실습에서는 대시보드와 경보를 구축하며, 파이프라인을 물려받은 온콜 담당자에게 마침내 계기판이 생긴다.
\end{calloutbox}

%% Chapter 14 — 모니터링과 신뢰성
\setcounter{chapter}{13}
\chapter{모니터링과 신뢰성}
\chaptersubtitle{서빙 지표, Prometheus와 Grafana, 소진율 경보, 드리프트 — ssh로 감당할 수 없게 된 플릿의 관제 센터}

\begin{calloutbox}{tbbblue}{tbbbluefill}
{\normalsize\bfseries\color{tbbblue}학습 목표}\par\vspace{3pt}
\begin{tbbitemize}
  \item 겉으로 건강해 보이는 플릿이 사용자에게는 실패할 수 있는 이유를 설명하고, probe·log·평균·\code{nvidia-smi}가 각각 SLO와는 \textit{다른} 질문에 답한다는 사실을 이해한다.
  \item 세 가지 지표 형태(counter, gauge, histogram), node·cAdvisor·kube-state·DCGM·engine \code{/\allowbreak metrics}로 이어지는 exporter 생태계, percentile을 \textit{올바르게} 집계하는 PromQL을 사용해 서빙 플릿 전체를 계측한다.
  \item 약속, 누구, 이유, 기반의 네 행으로 서빙 dashboard를 설계해 49분짜리 수색을 두 번의 시선으로 끝내는 진단으로 바꾼다.
  \item SLO 문장을 \textbf{소진율 경보}(여러 window, page와 ticket 구분)로 바꾸고, 올바른 이유로 올바른 사람만 깨우는 경보 목록을 설계한다.
  \item Canary도 page도 볼 수 없는 느린 실패인 \textbf{드리프트}를 분포 검사(PSI), 플릿이 이미 내보내는 sentinel metric, 관찰에서 재학습까지 이어지는 대응 사다리로 감지한다.
  \item Dashboard가 답하는 runbook, registry 항목에 연결한 비난 없는 postmortem, 모든 개선의 가치를 수치화하는 MTTD/MTTR로 전체 신뢰성 순환 고리를 운영한다.
\end{tbbitemize}
\end{calloutbox}

\vspace{2.5pt}

\begin{calloutbox}{tbbteal}{tbbtealfill}
{\normalsize\bfseries\color{tbbteal}핵심 용어}\par\vspace{3pt}
모니터링과 관측 가능성 · 세 기둥(metric · log · trace) · exporter · scrape · 시계열 · cardinality · counter / gauge / histogram · le bucket · PromQL · rate() · histogram\_quantile() · recording rule · ServiceMonitor · DCGM · kube-state-metrics · RED / USE 방법 · 네 행 dashboard · 증상과 원인 · 오류 예산 · 소진율 · 여러 window 경보 · 경보 피로 · dead man’s switch · synthetic probe · data / concept drift · PSI · sentinel metric · MTTD / MTTR · 비난 없는 postmortem
\end{calloutbox}

\vspace{2.5pt}

\begin{calloutbox}{tbblightgray}{tbbboxfill}
{\normalsize\bfseries\color{tbbink}이 장의 구성}\par\vspace{3pt}
\textbf{14.1} 두 번의 밤, 하나의 플릿, 그리고 아무도 몰랐던 스위치    \textbf{14.2} 지표: 플릿을 정직하게 계측하기    \textbf{14.3} Prometheus와 Grafana: 실제 네 행 dashboard    \textbf{14.4} 경보: page, 소진율, 거짓 경보를 피하는 기술    \textbf{14.5} 드리프트: 절대 page를 울리지 않는 실패    \textbf{14.6} 실습: 관제 센터를 만들고 두 번 깨뜨리기. 이어서 요약, 연습문제, 더 읽을거리를 다룬다.
\end{calloutbox}

\section{두 번의 밤, 하나의 플릿, 그리고 아무도 몰랐던 스위치}

제13장은 promotion과 조용한 경고로 끝났다. Canary의 analysis loop는 30분 동안 "TTFT가 정직한가? 행동이 drift했는가?"라는 정확한 질문을 던진 뒤 꺼졌지만, production은 \textit{아무도 듣지 않는 곳에서 영원히} 그 질문에 답하고 있었다. 이 장에서는 그 답을 들을 청중을 만든다. 제12장이 외상으로 남겨 둔 전쟁 이야기도 있다(§12.4: "SLO 아래의 time-slicing은 p99 미스터리가 태어나는 방식이다. 제14장에서 그중 하나를 만난다"). 이제 그 빚을 갚는 밤에서 시작하자. v9-int4의 canary가 끝난 지 3주 된 어느 화요일이다. 오늘 배포된 것은 없고, 아침과 똑같은 L4 레플리카 네 대가 실행 중이다. 21:30에 cluster 어딘가에서 예약 job이 시작되고, 22:06에 pager가 아니라 support 전화가 on-call을 깨운다. 아직 pager가 없기 때문이다.

\begin{terminalbox}
\begin{Verbatim}[breaklines=true,breakanywhere=true,fontsize=\small,formatcom=\color{tbbtermfg},codes=\tbbcodeactive]
# tuesday 22:06. "chat is molasses." uptime checker: /health 200 OK,
# green all night. no deploy since friday. so... what changed?
$ kubectl get pods -l app=chat
#   chat-1..chat-4  Running  READY 1/1  RESTARTS 0   <- reflexes: fine
$ kubectl logs deploy/chat --since=30m | grep -ci error
#   0                                                <- nothing confessed
$ curl -s chat-1:8000/metrics | grep ttft   # spot-check... ONE pod
#   p99 0.44 s. fine?!            <- you sampled a healthy replica
$ python mean_latency.py --since 21:00
#   mean e2e 15.2 s -> 21.3 s (+40%). "users exaggerate."
# 22:29  ssh gpu-node-2; nvidia-smi: util 93%, mem 21.9/24 GiB
#        "busy = evening peak, right?"      <- util counts ANYONE's kernels
# 22:33  wrong fix: scale to 6 replicas (150 s births). p99 -> 2.9 s.
#        better! not fixed. (the sick replica still serves 1 in 6.)
# 22:47  nvidia-smi once more, slower this time: TWO python processes
#        on one GPU. the 12.4-lab time-slice config, still on the node.
#        eval-nightly has been sharing chat-2's card since 21:30.
# 22:55  eval job killed. p99 0.46 s. total: 85 bad user-minutes.
\end{Verbatim}
\end{terminalbox}
\figcaption{사례 14-A  앞을 보지 못한 첫째 밤. 모든 도구가 정직하게 답했지만 질문이 모두 틀렸다. Probe는 약속이 아니라 생존을 확인하고, log는 침묵이 아니라 고백만 기록한다. 평균은 고통받은 사용자 4분의 1을 숨기며, pod 하나의 /metrics는 일화에 불과하다. GPU “utilization”은 누구의 것이든 kernel을 센다. 전화가 울리고 49분 뒤, 누군가 §12.4 실험을 기억해 낸 덕분에 마침내 원인을 찾았다.}

이 사례의 잔인함에는 배울 점이 있으니 잠시 곱씹어 보자. \textbf{거짓말한 것은 하나도 없었다.} Pod는 정말 Ready였다. Readiness는 반사 작용이고(제5장), process는 느린 batch 사이사이에 probe에 응답했다. Log도 정말 깨끗했다. Engine은 error를 내는 대신 \textit{기다리고} 있었고, 기다림은 log를 남기지 않는다. 평균도 실제로 40\%만 늘었다. 제7장에서 \textbf{평균은 아무도 묘사하지 않는다}고 경고했듯, 여기서는 건강한 레플리카 세 대와 물에 빠진 한 대를 섞어 버렸다. Curl로 확인한 \code{/\allowbreak metrics} endpoint의 p99도 정말 0.44초였다. 우연히 살아남은 레플리카를 골랐을 뿐이며, 어차피 레플리카별 percentile의 평균으로 플릿 전체를 설명할 수 없다(§14.2에서 정확히 살펴본다). \code{nvidia-smi}도 실제로 utilization 93\%를 표시했다. 다만 card에 상주하는 \textit{누군가의} kernel을 셀 뿐 \textit{누구의 것인지}는 말하지 않는다. 참인 문장 다섯 개, 진단은 0개다. 부족했던 것은 정직함이 아니라 \textbf{조립}이었다. SLO가 정의한 "정상"을 계속 계산하고 레플리카별로 분해해, 그 이유를 설명하는 계층 바로 위에 한눈에 놓는 장소가 필요했다. 그 장소에는 이름이 있다. 같은 화요일을 그 장소와 함께 다시 재생하면 처음부터 끝까지 20분이면 된다.

\begin{terminalbox}
\begin{Verbatim}[breaklines=true,breakanywhere=true,fontsize=\small,formatcom=\color{tbbtermfg},codes=\tbbcodeactive]
# the same tuesday, replayed with mission control (this week's lab).
# 21:31  slack #serving-alerts  [severity: ticket]
#        DuplicatePodsOnProdGPU: uuid GPU-...3d9f hosts 2 pods
#        (chat-2, eval-nightly-28471)          runbook R7
# 21:42  phone  [severity: page]
#        QueueSustained: num_requests_waiting > 8 for 5m on chat-2
#        (chat-1/3/4: 0)                       runbook R2 -> row 3
# 21:44  the walk (Figure 14-4): row 2 -> chat-2 only, LB fair;
#        row 3 -> decode step 2.1x, queue local, KV normal;
#        row 4 -> GPU ...3d9f: two pods. there is the neighbor.
# 21:47  kubectl delete job eval-nightly    # + incident note -> registry
# 21:50  queue drained. p99 0.46 s. page auto-resolves.
#        two-line diagnosis: "chat-2 only; shared GPU."
#        20 bad user-minutes, 8 of them yours. same on-call, same brain.
\end{Verbatim}
\end{terminalbox}
\figcaption{사례 14-B  둘째 밤. 이 장의 장치를 갖춘 뒤 같은 플릿에서 똑같은 사고가 일어났다. 21:31의 ticket은 사용자가 증상을 느끼기 전에 원인을 지목했고, 21:42의 page에는 runbook이 붙어 있었다. Dashboard는 진단을 번호가 붙은 순서로 바꾸었다. 두 밤 사이에 사람이 더 똑똑해진 것은 아니다. 질문을 미리 해 두었을 뿐이다.}

\begin{tbbfigure}
  \adjustbox{max width=0.94\linewidth,max totalheight=0.46\textheight}{\includegraphics{figures/fig14-1_twonights.png}}\\[3pt]
  \figcaption{그림 14-1  한 사고를 겪은 두 번의 밤. 첫째 밤에는 support를 통해 알아차리는 데 36분, 도구를 하나씩 발굴하는 데 다시 49분이 걸렸다. 나쁜 시간 85분 뒤 기억에 의존해 원인을 찾았다. 둘째 밤에는 21:31에 원인을 ticket으로 알리고, 21:42에 증상을 page로 보내, 21:50에 고쳤다. 나쁜 시간은 20분이었다. 두 행 사이의 거리가 바로 이 장이다.}
\end{tbbfigure}

\subsection*{1969년 11월 14일: telemetry가 엉망이 되었고, 한 관제사는 그 모습을 본 적이 있었다}

Apollo 12가 비 내리는 Florida에서 이륙한 지 36.5초 뒤, 번개가 Saturn V를 때렸다. Capsule 안에서는 fuel cell offline과 bus 저전압을 비롯한 panel의 모든 경고등이 한꺼번에 켜졌고, Houston 관제사의 화면은 잡음이 섞인 의미 없는 숫자로 무너졌다. Commander는 그것이 무엇을 뜻하는지 모른 채 침착하고 전문적인 목소리로 panel을 지상에 읽어 주었다. 16초 뒤 두 번째 번개가 상황을 더 악화했다. 이미 시야 밖으로 사라진 vehicle을 볼 수 있는 유일한 수단인 telemetry가 알아볼 수 없는 상태에서, 관제 센터는 2분 안에 달 착륙 임무를 중단할지 결정해야 했다. 그때 24세 비행 관제사 \textbf{John Aaron}이 우주 비행 역사에서 가장 유명한 문장 가운데 하나를 말했다. \textit{“Flight, try SCE to Aux.”} Flight director도 astronaut도 그 방의 어느 누구도 이 switch를 알지 못했다. 1년 전 화려할 것 없는 pad test에서 Aaron은 telemetry가 \textit{정확히 같은 형태}로 무너지는 모습을 보았고, 근무 외 시간에 원인을 추적해 Signal Conditioning Equipment라는 잘 알려지지 않은 box를 찾아냈다. Auxiliary power mode로 바꾸면 문제가 해결된다는 것도 알아냈다. Lunar module pilot Alan Bean이 switch를 찾았고 telemetry가 즉시 돌아왔다. Fuel cell을 reset한 뒤 임무는 달을 향해 계속되었으며, Aaron은 모든 운영 엔지니어가 남몰래 부러워할 별명인 \textit{the steely-eyed missile man}을 얻었다.

이 장의 모든 가르침이 그 90초 안에 있다. \textbf{첫째, telemetry는 미리 설계하지 않으면 존재하지 않는다.} Saturn V가 수천 channel을 지상으로 전송한 이유는 rocket에 ssh로 접속할 수 없음을 엔지니어들이 받아들였기 때문이다. 자기 플릿도 autoscaling node의 레플리카 네 대가 되는 날 같은 선을 넘었다. 사례 14-A는 그 규모에서 "직접 걸어가 확인하는 일"의 비용이다. \textbf{둘째, 원시 signal은 시야가 아니다.} 그 2분 동안 Houston에는 data가 있었다. 임무를 구한 것은 전에 본 적 있고 이름이 붙었으며 행동과 연결된 \textit{알아볼 수 있는 signature}였다. Dashboard가 정확히 그런 역할을 한다. Aaron이 SCE의 엉망인 pattern을 알아본 것처럼 화요일의 on-call이 "레플리카 하나는 빨강, GPU 하나에 pod 둘"을 알아보도록 signature를 미리 그려 둔다. \textbf{셋째, 알아보는 능력은 drill에서 산다.} Aaron은 위기 때가 아니라 \textit{test}에서 실패를 보았다. 그래서 이 장의 실습은 dashboard를 열어 둔 채 자기 플릿에서 사례 14-A를 일부러 재현하며 끝난다. NASA는 simulation, 제13장은 두 번 깨뜨리기, SRE 문헌은 game day라고 부른다. 이름보다 중요한 것은 일정에 넣는 일이다. Apollo에는 1202 alarm, 손으로 쓴 cheat sheet, 좋은 paging 정책의 기원이라는 두 번째 교훈도 있지만 §14.4에서 다룬다.

\subsection*{모니터링, 관측 가능성, 세 기둥 — 비용으로 이해하는 용어}

도구에 앞서 용어부터 정리하자. \textbf{모니터링}은 미리 정한 signal을 계속 수집하고 기대와 비교하는 활동이다. p99가 500 ms 아래인가, queue가 얕은가, budget이 소진되는가처럼 \textit{물어야 한다고 알고 있던} 질문에 답한다. \textbf{관측 가능성}은 밖으로 드러나는 signal을 이용해 미리 정의하지 \textit{않은} 질문에 답할 수 있는 system의 성질이다. 둘째 밤의 "레플리카 하나만 느린 이유는 무엇인가?"는 어느 runbook 제목에도 없었지만 레플리카별 queue, decode step time, GPU별 pod를 \textit{조합}해 두 번의 시선으로 답을 만들 수 있었다. 서로 독립적이고 label이 붙은 signal을 충분히 내보내 관측 가능성을 사고, 이미 아는 signal을 dashboard와 alarm에 연결해 모니터링을 수행한다. 업계는 둘을 \textbf{세 기둥}으로 제공하며, 무엇을 고를지는 비용의 형태로 정직하게 판단해야 한다(그림 14-2). \textbf{Metric}은 source에서 집계한 label이 붙은 숫자다. Token counter, queue depth gauge, TTFT histogram이 예다. Event마다 increment 한 번이면 되므로 비용이 거의 없고 traffic volume이 커져도 그대로 버티며, \textit{상시} 평가할 만큼 저렴한 유일한 기둥이다. 그래서 alert와 이 장이 metric에 자리 잡는다. \textbf{Log}는 event마다 남기는 기록, 즉 pod의 전기다. 이상하게 동작한 요청 하나를 찾는 데 꼭 필요하고 volume에 비례해 비용이 든다. 제5장의 entourage가 log shipper를 약속했고 실습에서 request\_id를 담은 구조화 JSON을 보내도록 설치한다. \textbf{Trace}는 요청 하나가 service를 지나는 경로를 따라 hop별 시간을 기록한다. Microservice debugger의 기둥이지만 event당 가장 비싸므로 sampling한다. 40개 service를 지나는 결제 경로와 달리 네 hop짜리 serving 경로는 그만큼 필요하지 않아 이 과정에서는 §14.2의 OpenTelemetry 한 문단으로 제한한다. 세 기둥을 묶는 규칙은 제6장에서 이미 비용을 냈다. \textbf{Telemetry는 pod 밖으로 계속 보낸다.} Pod가 죽기 직전 몇 분은 바로 그 pod가 직접 이야기할 수 없는 시간이기 때문이다. 제3장의 OOM kill은 app log가 아니라 cgroup counter와 event에 증거를 남겼다.

제5장은 이 장과 probe를 구분하는, 새겨 둘 만한 문장을 주었다. \textbf{Probe는 반사 작용이고 모니터링은 기억이다.} Liveness probe는 hang이 나면 restart하는 행동만 하고 설명하지 않는다. Readiness는 traffic을 차단하는 행동만 하고 설명하지 않는다. 화요일 내내 process가 살아서 응답했으므로 둘 다 초록색이었다. 모니터링 자체는 설명만 할 뿐 행동하지 않는다. System의 기억(21:35의 queue는 어땠는가?), 주의력(기억을 전화 진동으로 바꾸는 alert), §14.5의 drift detection을 더하면 예지력(현실이 어느 방향으로 벗어나는가?)이 된다. 성숙한 플릿에는 두 신경계가 모두 필요하다. 제5장의 논술은 관측 가능성 없는 reconciliation이 자기가 고친 일을 감추는 \textit{사고 억압}이라고 경고했다. 이번 주에는 Kubernetes가 조용히 모으던 restart counter, OOM event, endpoint 공백이 모두 graph와 alert가 되는 선으로 바뀐다. Drift correction의 \textit{빈도}를 chart에 그려야 90초마다 한 번씩 restart를 일으킬 만큼 작게 잡힌 limit를 발견할 수 있다.

앞으로 90분 동안 갈 길을 살펴보자. \textbf{§14.2}에서는 세 가지 metric 형태, node부터 DCGM과 engine의 \code{/\allowbreak metrics}까지 이어지는 exporter 생태계(마침내 graph가 되는 제11장의 줄), 플릿 규모에서 percentile을 올바르게 계산하는 두 가지 PromQL 주문으로 signal 계층을 만든다. \textbf{§14.3}에서는 약속, 누구, 이유, 기반으로 이루어진 \textbf{네 행 dashboard}를 조립하고, 그 위에서 panel을 하나씩 따라 둘째 밤을 재생해 §12.4의 빚을 모두 갚는다. \textbf{§14.4}에서는 \textbf{alarm}을 연결한다. SLO 문장을 burn rate로 compile하고, 빠른 소진은 page로, 느린 소진은 ticket으로 보내는 여러 window 정책을 만들며, 증상과 원인을 구분해 routing한다. 거부할 수 있는 모든 계층의 바깥에서 synthetic probe를 실행하고(마침내 감시받는 제6장의 문), 감시자를 감시하는 dead man’s switch를 더한다. Runbook, on-call, 비난 없는 postmortem으로 이어지는 사람의 순환 고리도 다룬다. \textbf{§14.5}에서는 절대 page를 울리지 않는 실패인 \textbf{드리프트}를 살펴본다. 제13장의 canary가 구조적으로 잡을 수 없는 경사면을 손으로 계산할 수 있는 PSI와 플릿이 이미 무료로 내보내는 sentinel로 찾는다. \textbf{§14.6}은 강의계획서가 약속한 serving dashboard와 alarm 구축 실습에, 강의계획서가 좀처럼 약속하지 않는 부분을 더한다. 두 번 깨뜨리고 자기 panel이 그것을 잡게 한다. 관제 센터를 세운 뒤 임무를 수행한다.

\begin{calloutbox}{tbbpurple}{tbbpurplefill}
{\normalsize\bfseries\color{tbbpurple}생각해 보기}\par\vspace{3pt}
\begin{tbbitemize}
  \item 사례 14-A의 on-call은 다섯 가지를 확인했고 모두 “정상”이라는 답을 얻었다. Probe, log, 평균 latency, pod 하나의 /metrics, nvidia-smi util이 \textbf{실제로} 답하는 질문과 SLO 문장이 답해야 했던 질문을 각각 말하자. 그림 14-4의 panel 하나로 다섯 가지를 모두 대신한다면 어느 것인가?
  \item 22:33의 “잘못된 해결책”, 즉 레플리카를 6대로 늘린 조치는 p99를 실제로 4.6초에서 2.9초로 개선했다. 제12장의 산술을 이용해 레플리카 추가가 아픈 레플리카 하나의 영향을 희석할 뿐 치료하지는 못하는 이유를 설명하자. 진단을 가리는 동안 추가 레플리카 두 대의 시간당 비용은 얼마였는가?
  \item John Aaron의 해결에는 (a) 이상 중에도 계속 흐른 telemetry, (b) drill에서 본 signature, (c) 실제로 존재한 switch가 필요했다. 세 가지를 둘째 밤에 대응시키고, 자기 프로젝트 플릿에서 가장 가능성 높은 실패에 대비할 때 지금 빠진 것이 무엇인지 밝히자.
  \item 둘째 밤에는 원인 alert(21:31, ticket)가 증상 alert(21:42, page)보다 먼저 왔다. 심각도를 반대로 배정하지 않은 이유를 설명하자. 원인처럼 보이는 모든 이상이 page를 울리면 on-call의 삶은 어떻게 되는가? 답은 잠시 보류하자. §14.4에서 병원 통계로 비용을 계산한다.
\end{tbbitemize}
\end{calloutbox}

\section{지표: 플릿을 정직하게 계측하기}

\textbf{지표(metric)}는 시스템이 자신에 관해 말할 수 있는 가장 작고 정직한 문장이다. \textit{이름}, \textit{레이블} 집합, \textit{숫자}, \textit{시각}으로 이루어진다. \code{vllm:num\_requests\_waiting\{pod="chat-2"\} 19 @21:41:30}에는 주어, 한정어, 서술어, 타임스탬프가 모두 들어 있다. 이 장의 모든 내용은 이런 문장을 15초마다 수집하고 몇 주 동안 보관한 뒤, 집계해서 질의하는 데서 출발한다. 이름과 레이블의 조합이 하나씩 달라질 때마다 별도의 \textbf{시계열(time series)}이 생긴다. 이 분야의 첫 번째 법칙은 관측값이 아니라 \textit{시계열} 단위로 비용을 치른다는 것이다. 값이 네 개인 레이블(파드)은 지표 하나를 네 시계열로 늘릴 뿐이라 괜찮다. 값이 900만 개인 레이블(사용자 ID나 요청 ID)은 이를 900만 시계열로 늘린다. 이것은 모니터링의 탈을 쓴 메모리 장벽이다. 비용을 감당하려면 다음 규칙을 지켜야 한다. \textbf{레이블은 한정된 목록에서 “어느 것인가?”에 답해야 한다.} 파드, 모델, GPU UUID, 경로가 그 예다. 값의 수가 무한히 늘어날 수 있는 정보는 로그에 두어야 한다. 로그는 이벤트별 저장이 재앙이 아니라 애초의 설계이기 때문이다. (이 강좌의 플릿에는 약 14,000개의 활성 시계열이 있다. 그림 14-2에서 대략 계산한 값으로, 초당 표본은 약 1,000개이고 15일 보존 기간에 필요한 공간은 2GB 미만이다. 부주의하게 \code{user\_id} 레이블 하나만 추가해도 이 수가 천 배로 늘어난다.)

\begin{tbbfigure}
  \adjustbox{max width=0.94\linewidth,max totalheight=0.46\textheight}{\includegraphics{figures/fig14-2_pipeline.png}}\\[3pt]
  \figcaption{그림 14-2  텔레메트리 파이프라인. Prometheus가 모든 계층에서 지표를 가져와 Grafana(질문)와 Alertmanager(호출)로 나누어 보낸다. 로그는 기록되는 즉시 DaemonSet을 통해 전송하고, 이 강좌에서는 트레이스를 최소한으로 유지한다. 아래쪽에는 여섯 관측 지점의 익스포터가 있다. 하나의 스크레이프 루프가 거부할 수 있는 모든 계층을 수집 가능한 대상으로 만든다.}
\end{tbbfigure}

\subsection*{세 가지 형태: 카운터, 게이지, 히스토그램 — 세 번째가 이 강좌를 움직이는 이유}

내보내는 지표는 모두 세 형태 가운데 하나이며, 처음 배울 때 가장 흔한 실수는 형태를 잘못 고르는 것이다. \textbf{카운터}는 생성된 토큰 수, 부팅 이후의 선점 횟수, 완료된 요청 수처럼 오르기만 한다. 순간값 자체는 별 의미가 없다(선점 1,841회라고 해도 \textit{언제부터} 센 값인가?). 신호가 되는 것은 \textbf{변화율}이다. 그래서 질의 언어에서 가장 자주 쓰는 함수가 \code{rate(x[5m])}이다. 일정 구간의 초당 기울기를 계산하며, 재시작에도 안전하다. 파드가 다시 태어나면서 카운터가 0으로 초기화되면 이를 감지해 이어 계산한다. 직접 만든 “이번 분의 요청 수” 게이지보다 카운터를 선호해야 하는 이유가 하나 더 생긴다. \textbf{게이지}는 대기열 깊이, KV 캐시 점유율, 온도처럼 오르내리는 값이다. 게이지는 그대로 읽고, 파드나 시간에 걸쳐 \code{max}/\code{avg}로 집계하며, 포화 신호를 표현하기에 알맞다. 가장 흥미로운 것은 \textbf{히스토그램}이다. 히스토그램은 임곗값 이하의 관측값 개수를 세는 \textit{누적 버킷}(\code{le="0.25"}, \code{le="0.5"}, ...) 묶음과 누적 합계 및 개수로 구성된다. 이는 7장에서 남긴 문제를 정확히 해결하려고 존재한다. \textbf{백분위수는 평균을 내거나 더할 수 없다.} 레플리카 열 개의 p99를 어떤 방식으로 계산해도 플릿의 p99가 되지 않는다. 반면 버킷 \textit{카운터}는 다른 카운터처럼 더할 수 있다. 레플리카 네 개의 \code{le="0.5"} 버킷을 합치면 플릿 전체 분포를 얻고, 여기에서 원하는 분위수를 구할 수 있다. 히스토그램은 7장의 문법 규칙을 데이터 구조로 구현한 것이다.

실제 시스템에서 세 형태를 살펴보자. 11주차부터 실행해 온 엔진은 줄곧 세 형태를 모두 내보내고 있었다.

\begin{terminalbox}
\begin{Verbatim}[breaklines=true,breakanywhere=true,fontsize=\small,formatcom=\color{tbbtermfg},codes=\tbbcodeactive]
$ curl -s chat-2:8000/metrics | grep -E 'vllm:' | head -14
vllm:num_requests_running{model="chat-v9"} 32.0     <- gauge: the seats,
vllm:num_requests_waiting{model="chat-v9"} 19.0        full (11.3); and THE
vllm:gpu_cache_usage_perc{model="chat-v9"} 0.96        alarm line (11.4)
vllm:num_preemptions_total{model="chat-v9"} 1841.0  <- counter: rate() it
vllm:spec_decode_draft_acceptance_rate{...} 0.74    <- gauge: 11.5's bet,
                                                       still paying
vllm:time_to_first_token_seconds_bucket{le="0.25"} 811021.0   <- histogram:
vllm:time_to_first_token_seconds_bucket{le="0.5"}  947210.0      cumulative
vllm:time_to_first_token_seconds_bucket{le="1.0"}  951003.0      buckets
vllm:time_to_first_token_seconds_bucket{le="+Inf"} 953882.0      (le = "at
vllm:time_to_first_token_seconds_count 953882.0                  or below")

$ curl -s gpu-node-2:9400/metrics | grep -E 'DCGM' | head -3
DCGM_FI_DEV_GPU_UTIL{UUID="GPU-...3d9f",pod="chat-2"} 93
DCGM_FI_DEV_FB_USED{UUID="GPU-...3d9f",pod="chat-2"} 22429
DCGM_FI_DEV_XID_ERRORS{UUID="GPU-...3d9f"} 0
\end{Verbatim}
\end{terminalbox}
\figcaption{실행 기록 14-1  주석을 단 두 번의 스크레이프. 엔진은 세 형태를 모두 내보낸다. 전체 누적 요청 953,882개 가운데 947,210개의 TTFT가 0.5 s 이하였다는 점에 주목하자. 99.3\%로, SLO 자체의 SLI가 두 버킷 카운터에 그대로 드러나 있다. DCGM은 실리콘의 상태를 말해 준다. GPU UUID를 레이블로 삼고, 장치 플러그인 매핑을 거쳐 파드 레이블도 붙인다. 바로 이 조인이 화요일의 문제를 푼다.}

왜 \textbf{풀(pull)} 방식일까? Prometheus는 푸시를 받지 않고 대상을 스크레이프한다. 이는 배관 방식이 아니라 철학의 선택이다. 풀 방식에서는 수집 대상을 \textit{열거}할 수 있다. 서버는 Kubernetes에서 실시간으로 발견한 대상 목록을 보유한다. 따라서 “어떤 레플리카가 존재하는가?”와 “어떤 레플리카가 응답을 멈췄는가?”가 같은 사실이 된다. 스크레이프할 때마다 \code{up\{pod=...\}} 시계열을 기록하므로, 사고 도중 죽은 파드는 아무도 알아채지 못하는 침묵 대신 \code{up == 0}이라는 묘비를 남긴다. 푸시 시스템에서는 부재를 별도로 모델링해야만 알아낼 수 있다. 풀 방식은 계측 로직도 단순하게 유지한다. 엔진은 HTTP로 텍스트만 제공하면 되고, 지연 시간에 민감한 프로세스 안에 버퍼링, 배치 처리, 역압 로직을 넣지 않아도 된다. 그 대가로 Prometheus가 대상에 \textit{도달}할 수 있어야 하지만, 단일 클러스터 안에서는 문제가 되지 않는다. 이 설계에는 계보가 있다. Google 내부의 \textbf{Borgmon}은 5장에 나온 Borg와 비슷한 2003년 무렵 탄생해 화이트박스 \code{/\allowbreak varz} 엔드포인트, 중앙 풀러, 레이블 시계열용 질의 언어, 질의로 정의한 경보라는 패턴을 세웠다. 전 Google 엔지니어 Matt T. Proud와 Julius Volz는 2012년 SoundCloud에 합류한 뒤 이와 견줄 만한 공개 시스템이 없다는 사실을 알고 이를 오픈소스로 다시 만들었다. \textbf{Prometheus}는 2015년에 공개되었고 Cloud Native Computing Foundation의 두 번째 프로젝트가 되었다. 첫 번째는 Kubernetes였다. 두 프로젝트가 나란히 선택되었다는 사실은 업계가 이 둘을 한 설계의 두 절반으로 받아들였음을 보여 준다. \textbf{Grafana}는 반대편에서 출발했다. Torkel Ödegaard는 2014년에 로그 검색 UI인 Kibana 3를 포크해 범용 시계열 대시보드로 바꾸었다. 실습에서 설치하는 스택은 이 두 계보가 합류한 결과다.

\subsection*{네 가지 질의로 익히는 PromQL — 화요일의 문제를 푼 두 질의까지}

PromQL은 레이블이 붙은 시계열에 적용하는 집합대수처럼 읽힌다. 플릿 운영에 필요한 관용 패턴은 열두 개 남짓이면 충분하다. 다음 네 가지는 이 장을 떠받치는 핵심 질의다. 그중 가운데 두 개에는 그림 14-3의 교훈이 담겨 있으며, 사고가 난 구간에 실행한 결과를 보여 준다.

\begin{terminalbox}
\begin{Verbatim}[breaklines=true,breakanywhere=true,fontsize=\small,formatcom=\color{tbbtermfg},codes=\tbbcodeactive]
# 1. the WRONG fleet p99 (tuesday, 21:44) -- averaging per-replica p99s:
> avg( job:vllm_ttft:p99_5m )
2.76                    # experienced by nobody (Figure 14-3)

# 2. the RIGHT fleet p99 -- merge bucket counters, THEN take the quantile:
> histogram_quantile(0.99,
    sum by (le) (rate(vllm:time_to_first_token_seconds_bucket[5m])))
9.08                    # the promise, actually broken

# 3. decompose: waiting or working? (Ch. 9's split, per replica)
> max by (pod) (vllm:num_requests_waiting)
{pod="chat-1"} 0   {pod="chat-2"} 19   {pod="chat-3"} 0   {pod="chat-4"} 0
> histogram_quantile(0.5, sum by (pod, le)
    (rate(vllm:time_per_output_token_seconds_bucket[5m])))
{pod="chat-2"} 0.150    # vs 0.070 fleet norm: decode itself is 2.1x slower

# 4. the SLI as a ratio -- the number 14.4 will set on fire:
> sum(rate(vllm:time_to_first_token_seconds_bucket{le="0.5"}[5m]))
  / sum(rate(vllm:time_to_first_token_seconds_count[5m]))
0.762                   # 76.2% good this window; target 0.990
\end{Verbatim}
\end{terminalbox}
\figcaption{실행 기록 14-2  사고 구간에 실행한 네 가지 질의. 질의 1은 그럴듯하지만 틀린 답이다. 질의 2는 7장의 규칙을 기계로 실행한다. 질의 3은 “대기 중”과 “처리 중”을 분리한다. 9장에서 이 진단 정보만으로 대시보드 설계의 절반이 완성된다고 약속했다. 질의 4는 §14.4에서 소진율 호출로 변환할 SLI 비율이다.}

질의 2의 구조를 안쪽부터 한 번만 읽어 보면 완전히 이해할 수 있다. \code{rate(...bucket[5m])}는 각 누적 버킷 카운터를 초당 흐름으로 바꾼다. \code{sum by (le)}는 \textit{버킷 경계를 보존하면서} 레플리카별 흐름을 합친다. 분포를 여러 레플리카에 걸쳐 집계할 때 수학적으로 허용되는 유일한 방법이다. 이어서 \code{histogram\_quantile(0.99, ...)}가 합쳐진 버킷을 따라가 99번째 백분위수를 찾고, 해당 버킷 안에서 보간한다. 그래서 결과의 정밀도는 버킷 너비에 좌우된다. 실습에서 엔진 버킷 경계를 SLO \textit{바로 그 값}에 두는 이유도 여기에 있다. 경계를 정확히 0.5 s에 두면 p99 추정값은 여전히 근사치여도 SLI 비율은 정확해진다. 이제 한 가지 관용 패턴을 더 이름 붙일 차례다. \textbf{기록 규칙(recording rule)}은 Prometheus가 일정에 따라 질의를 다시 실행하고 그 결과를 새 시계열로 저장하는 규칙이다. 질의 4에는 \code{job:slo\_ttft\_good:ratio\_rate5m}를 만들고, 대시보드용으로는 파드마다 \code{job:vllm\_ttft:p99\_5m}를 만든다. 패널과 경보는 새로 고칠 때마다 원시 버킷을 다시 집계하지 않고 미리 계산된 결과를 읽는다. 규칙은 저장소의 YAML로 관리하며, 13장 이후의 모든 것과 마찬가지로 PR 검토를 받는다.

\subsection*{익스포터 모음 — 여섯 관측 지점과 각각의 역할}

Prometheus 자체는 아무것도 측정하지 않는다. 온도계를 들고 있는 것은 \textbf{익스포터}다. 그림 14-2 아래쪽의 익스포터 모음은 이 강좌의 스택을 마지막으로 한 번 더 관통한다. DaemonSet으로 실행되는 \textbf{node-exporter}는 CPU, 메모리, 디스크, 네트워크처럼 2장에서 다룬 머신 하드웨어를 시계열로 말해 준다. kubelet에 내장된 \textbf{cAdvisor}는 \textit{3장에서 직접 읽었던 바로 그 파일}, 즉 cgroup의 \code{memory.current}, \code{memory.events}, CPU 사용량 정보를 읽어 모든 컨테이너의 상태를 말한다. 여기서 3장의 두 교훈이 결실을 맺는다. 첫째는 정직한 마운트에 관한 주석이다. 컨테이너 안의 프로세스가 \code{/\allowbreak proc}를 조회하면 컨테이너가 아니라 \textit{호스트의} CPU와 메모리를 본다. 단순한 컨테이너 내부 에이전트가 거짓말하는 이유이며, cAdvisor가 대신 바깥에서 cgroup 트리를 읽는 이유다. 둘째는 경보다. \code{memory.events}는 \code{high}와 \code{oom\_kill} 경계를 넘은 횟수를 센다. 따라서 3장 실습에서 \code{memory.high}로 설계한 벽 직전의 경고 단계는 이제 Grafana가 그래프로 그리고 §14.4가 종료 코드 137이 나타나기 \textit{전에} 티켓을 만들 수 있는 카운터가 된다. 이 종료 코드는 마지막으로 한 번 더 레이블 값으로 등장한다. \textbf{kube-state-metrics}는 컨트롤 플레인의 \textit{의도}를 말한다. 재시작 횟수, 마지막 종료 원인인 \code{OOMKilled}, 배포에서 원하는 레플리카 수와 사용 가능한 레플리카 수의 차이, 비어 있는 엔드포인트, HPA 상태를 제공한다. “가용성만이 아니라 드리프트 교정 빈도에도 경보를 걸라”던 5장의 목록이 이제 데이터 원천을 얻는다. \textbf{DCGM-exporter}는 \code{nvidia-smi}를 산업 규모로 만든 도구다. 이는 2장의 약속대로 GPU UUID별 사용률, 프레임버퍼, 전력, 클록, Xid 오류 코드를 제공하고, 장치 플러그인 매핑으로 파드 이름과 조인한다. 덕분에 화요일의 문제는 표를 한 번 조회하는 일로 바뀌었다. 또한 “사용률”이 회피하는 질문에 답하는 프로파일링 계열(\code{DCGM\_FI\_PROF\_SM\_ACTIVE}, 텐서 파이프 활동)도 제공한다. 이 거짓말은 잠시 뒤 자세히 살펴본다. \textbf{엔진 자체의 /metrics}는 이미 확인했다. 여기에 9장에서 분류기 서비스가 제공한 Triton 카운터도 더하자. 이 카운터는 대기열 시간과 계산 시간을 이미 분리해 준다. \textbf{블랙박스 익스포터}는 \textit{바깥에서} 조사한다. 게이트웨이의 공개 출입구로 HTTP GET을 보내고, TCP 핸드셰이크를 수행하고, 스트리밍 토큰을 확인한다. 6장에서 사용자와 파드 사이의 모든 계층이 저마다의 방식으로 거부할 수 있음을 배웠기 때문이다. 플릿 내부에서만 자신을 측정하면 로드 밸런서 상태 확인이 실패하는 동안에도 플릿은 결백하다고 주장한다. 실행 기록 6-11의 교훈처럼 증거는 클러스터 밖에 있었다. 합성 검사는 §14.4에서 다루며, 익스포터는 그 결과를 시계열로 바꾸는 수단이다.

이제 22시 29분에 마주한 거짓말을 살펴보자. \code{nvidia-smi}가 출력하는 값과 같은 \code{DCGM\_FI\_DEV\_GPU\_UTIL}은 단 하나의 질문에만 답한다. \textit{최근 시간 가운데 누구의 것이든 커널 하나 이상이 장치에 올라와 있던 비율은 얼마인가?} 이 값은 \textbf{이진 점유 듀티 사이클}이지, 여유 용량이나 효율, 소유권을 측정한 값이 아니다. 작은 커널 하나를 계속 올려 두는 메모리 병목 디코드 루프는 ALU의 10분의 1만 사용하면서도 90\% 이상으로 표시된다. 10장의 루프라인이 대시보드의 함정으로 다시 나타난 것이다. 두 테넌트가 카드 하나를 시분할하면 각자 절반의 작업만 처리하면서도 \textit{둘 다} 값을 100\% 가까이 끌어올린다. 그 숫자만으로는 \textit{누구의} 커널인지도 알 수 없다. 첫날 밤 당직자는 93\%를 보고 “바쁘다, 그러므로 부하가 걸렸다, 그러므로 정상이다”라고 판단했다. 근거 없는 추론을 세 번 연달아 한 셈이다. GPU에 관해 정직한 질문을 던지고 정직한 데이터 원천에서 답을 구해야 한다. \textit{실리콘이 제값을 하고 있는가?}는 10장에서 예상한 값과 \code{PROF\_SM\_ACTIVE}, 텐서 파이프 활동을 비교해 답한다. \textit{메모리가 장벽인가?}는 디코드 중 \code{PROF\_DRAM\_ACTIVE}가 포화에 가까우면 메모리 장벽이 눈앞에 드러난다. \textit{카드가 정상인가?}는 Xid 오류, 온도, 전력을 본다. \textit{누구의 것인가?}는 4행의 결정적 증거인 UUID당 파드 조인으로 답한다. 사용률은 “무언가 달라졌다”는 신호로 쓸 만하므로 대시보드에 남긴다. 그러나 이 값만으로는 아무 호출도 보내지 않는다. 실습의 JSON에는 이 패널 아래에 늘 다음 캡션이 붙는다. \code{busy ≠ yours ≠ efficient}

이런 익스포터 모음을 정리하는 두 약어가 있으며 면접에서도 자주 등장한다. \textbf{RED}는 Rate, Errors, Duration의 약자로, 사용자가 호출하는 대상마다 행 하나를 두는 \textit{서비스 측} 관점이다. 초당 요청 수, 실패 비율, 지연 시간 분포를 본다. \textbf{USE}는 Utilization, Saturation, Errors의 약자로, 서비스가 소비하는 자원마다 행 하나를 두는 \textit{자원 측} 관점이다. 얼마나 바쁜지, 얼마나 밀려 있는지, 무엇이 고장 났는지를 본다. 이 강좌의 플릿에 정확히 대응시켜 익혀 둘 가치가 있다. TTFT 히스토그램은 RED의 Duration이고, \code{num\_requests\_waiting}은 좌석 풀에 대한 USE의 Saturation이다. 이 책 전체에서 사용률이 아니라 포화가 선행 지표다. \code{memory.high} 압력이 OOM 종료보다 먼저 나타나고 7장의 무릎점이 하키 스틱 곡선보다 먼저 나타나듯, 대기열은 지연 시간 위반보다 \textit{먼저} 커진다. KV 캐시 점유율은 \textit{높게 유지되도록 설계된} USE의 Utilization이다. 11장에서 높게 고정된 상태가 건강하다고 설명했지만, 플릿을 처음 접하는 사람은 이 게이지를 비상 상황으로 오해하기 쉽다. 선점률은 같은 풀의 포화를 퇴출이라는 형태로 표현한 USE의 Saturation이다. 분류기의 대기열 시간과 계산 시간 분리는 RED의 Duration을 원인별로 나눈 것이다. 표 14-1에 필요한 지표를 모았다.

\begin{tbbtablewrap}
  \begin{tabular}{@{}>{\raggedright\arraybackslash}m{\dimexpr 0.2181\tbltextwidth-2\tabcolsep\relax} >{\raggedright\arraybackslash}m{\dimexpr 0.2766\tbltextwidth-2\tabcolsep\relax} >{\raggedright\arraybackslash}m{\dimexpr 0.1117\tbltextwidth-2\tabcolsep\relax} >{\raggedright\arraybackslash}m{\dimexpr 0.2606\tbltextwidth-2\tabcolsep\relax} >{\raggedright\arraybackslash}m{\dimexpr 0.1330\tbltextwidth-2\tabcolsep\relax}@{}}
  \arrayrulecolor{tbbnavy}\specialrule{1pt}{0pt}{0pt}
  \rowcolor{tbbtablehead}{\centering \bfseries\color{tbbnavy}신호} & {\centering \bfseries\color{tbbnavy}지표 · 데이터 원천} & {\centering \bfseries\color{tbbnavy}형태} & {\centering \bfseries\color{tbbnavy}답하는 질문} & {\centering \bfseries\color{tbbnavy}경보?} \\
  \arrayrulecolor{tbbnavy}\specialrule{0.7pt}{0pt}{2pt}
  {\textbf{TTFT 분포}} & {vllm:time\_to\_first\_token\_seconds · 엔진} & {히스토그램} & {약속을 지키고 있는가? (1행, SLI)} & {소진율 쌍 — 호출/티켓} \\
  \arrayrulecolor{tbbtableborder}\hline
  {\textbf{디코드 속도}} & {vllm:time\_per\_output\_token\_seconds · 엔진} & {히스토그램} & {생성 자체가 느려지는가? (두 번째 조항)} & {관찰, 추세 시 티켓} \\
  \arrayrulecolor{tbbtableborder}\hline
  {\textbf{대기열 깊이}} & {vllm:num\_requests\_waiting · 엔진} & {게이지} & {수요와 좌석의 차이 — 선행 증상(11장)} & {호출(지속 시)} \\
  \arrayrulecolor{tbbtableborder}\hline
  {\textbf{KV 캐시 점유율}} & {vllm:gpu\_cache\_usage\_perc · 엔진} & {게이지} & {풀은 건강한가? — 높게 고정되면 설계가 작동하는 상태} & {없음(의도된 설계)} \\
  \arrayrulecolor{tbbtableborder}\hline
  {\textbf{선점}} & {vllm:num\_preemptions\_total · 엔진} & {카운터} & {풀이 요동치는가? (눈에 보이는 KV 압력)} & {발생률에 따라 티켓} \\
  \arrayrulecolor{tbbtableborder}\hline
  {\textbf{초안 수락률}} & {vllm:spec\_decode\_draft\_acceptance\_rate · 엔진} & {게이지} & {추측이 여전히 효과적인가? (드리프트 감시 지표, §14.5)} & {주간 검토} \\
  \arrayrulecolor{tbbtableborder}\hline
  {\textbf{재시작 / OOM}} & {restarts\_total + OOMKilled 원인 · kube-state} & {카운터} & {자가 치유가 누수를 감추는가? (5장, 3장)} & {티켓, 플릿 전체면 호출} \\
  \arrayrulecolor{tbbtableborder}\hline
  {\textbf{준비된 엔드포인트}} & {엔드포인트 주소 준비 개수 · kube-state} & {게이지} & {Service가 어느 곳으로든 라우팅할 수 있는가? (6장의 503)} & {0이면 호출} \\
  \arrayrulecolor{tbbtableborder}\hline
  {\textbf{GPU 상태}} & {DCGM 사용률 / FB / Xid / PROF\_* · DCGM} & {혼합} & {실리콘이 무엇을 하는가? 누구의 것인가?} & {Xid → 티켓} \\
  \arrayrulecolor{tbbtableborder}\hline
  {\textbf{출입구 검사}} & {probe\_success · 블랙박스} & {게이지} & {바깥에서도 작동하는가? (6장의 출입구)} & {3회 연속이면 호출} \\
  \arrayrulecolor{tbbtableborder}\hline
  {\textbf{플릿 태세}} & {원하는 레플리카와 준비된 레플리카 · kube-state/HPA} & {게이지} & {파드가 물결을 따라가고 있는가? (12장)} & {멈추면 티켓} \\
  \arrayrulecolor{tbbnavy}\specialrule{1pt}{2pt}{0pt}
  \end{tabular}\par
  \figcaption{표 14-1  서빙 플릿에 필요한 지표 목록. 엔진 히스토그램 두 개, 엔진 게이지 세 개, 엔진 카운터 한 개에 플랫폼의 기억(kube-state), 실리콘의 활력 징후(DCGM), 외부 관점(블랙박스)을 더한다. 경보 열은 §14.4를 미리 보여 준다. 두 호출은 증상을 다루고, 나머지는 모두 기록을 남긴다.}
\end{tbbtablewrap}

신호 계층을 떠나기 전에 세 번째 기둥을 정직하게 짚고 넘어가자. 게이트웨이, 대기열, 엔진에 걸친 OpenTelemetry 스팬을 이어 붙인 \textbf{트레이스}는 요청 시간이 여러 서비스 사이의 \textit{인계 지점}에 숨어 있을 때 제값을 한다. 이 강좌의 경로에는 홉이 세 개뿐이고, 인계 과정의 수수께끼인 “대기 중인가, 처리 중인가?”는 질의 3의 두 지표로 이미 풀린다. 따라서 OTel SDK의 기본 HTTP 계측으로 트레이스의 1\%만 표본 추출하고, 복잡한 스팬 분석은 마이크로서비스 강좌에 맡긴다. \textbf{서비스 메시}도 같은 이유로 미룬다. 5장의 세 번째 사이드카이자 6장의 L7 친척인 서비스 메시는 모든 홉의 RED 지표를 \textit{무료로} 제공한다. 실제로 가장 설득력 있는 장점이다. 하지만 모든 파드에 프록시를 넣어야 하고, 재시도를 수행하는 계층이 하나 더 생긴다. 관심사마다 소유자를 하나만 두라는 6장의 규칙은 여전히 유효하다. 이 시스템에서는 게이트웨이가 재시도를 소유하므로 메시를 보류한다. 14주차에는 5장에서 예고한 수행원 가운데 지표 익스포터와 로그 전송기를 받아들여 둘 다 실습에서 설치한다. 메시 프록시는 “나중”으로 남으며, 그 나중은 영원히 오지 않을 수도 있다.

\begin{calloutbox}{tbbpurple}{tbbpurplefill}
{\normalsize\bfseries\color{tbbpurple}생각해 보기}\par\vspace{3pt}
\begin{tbbitemize}
  \item 실행 기록 14-1의 전체 누적 버킷에 따르면 chat-2 요청의 99.3\%가 0.5 s 안에 처리되었다. 그런데 사고 도중 병든 레플리카도 chat-2다. 이 두 사실을 조화롭게 설명하라. 전체 누적 카운터와 rate() 구간의 차이를 생각하고, 원시 카운터 값을 언제 확인해야 하는지 일반 규칙을 제시하라.
  \item 팀원이 “사용자별로 디버깅할 수 있게” 모든 엔진 지표에 session\_id 레이블을 붙이자고 제안한다. 시계열 약 14,000개라는 계산과 월간 턴 약 950만 개를 이용해 변경 후 시계열 수를 추정하고, 발생할 실패 형태를 설명하라. 요청별 디버깅 정보가 속해야 할 올바른 기둥과 키도 제시하라.
  \item 새벽 3시에 살아남은 야간 레플리카의 GPU 사용률은 41\%다. 비용에 민감한 PM이 “그렇다면 59\%를 낭비하고 있다는 뜻인가?”라고 묻는다. 듀티 사이클과 여유 용량을 구분하고, 메모리 병목 디코드 플릿의 낭비를 실제로 측정할 지표를 제시한 뒤, 질문의 초점을 12장의 레플리카 시간 계산으로 돌리는 세 문장의 답을 작성하라.
  \item 표 14-1의 각 행을 RED와 USE로 분류하라. 경보 열의 패턴에도 주목하라. 어느 방법의 행은 호출을 보내고 어느 행은 티켓을 만드는가? 이 구분이 어떻게 §14.4의 증상 대 원인 규칙을 자동으로 재현하는가?
\end{tbbitemize}
\end{calloutbox}

\section{Prometheus와 Grafana: 네 행 대시보드 실전 가동}

설치는 일부러 싱거울 만큼 단순하다. 실습 1단계에서 \textbf{kube-prometheus-stack} Helm 차트 명령 하나로 Prometheus, Grafana, Alertmanager, node-exporter, kube-state-metrics를 설치한다. 풀 방식을 가능하게 하는 Kubernetes 서비스 발견도 미리 연결되어 있다. Prometheus는 API 서버를 감시하고, 스크레이프 주석이 붙은 Service는 한 번의 스크레이프 주기 안에 수집 대상이 된다. 스택에 \textit{직접 만든} 서비스를 알려 주려면 서비스마다 작은 사용자 정의 리소스인 \textbf{ServiceMonitor} 하나와 §14.2의 기록 규칙을 추가하면 된다. 둘 다 저장소에 있는 평범한 YAML이며, 다른 릴리스 산출물과 마찬가지로 13장의 파이프라인을 거쳐 배포된다.

\begin{terminalbox}
\begin{Verbatim}[breaklines=true,breakanywhere=true,fontsize=\small,formatcom=\color{tbbtermfg},codes=\tbbcodeactive]
# monitoring/servicemonitor-chat.yaml -- "scrape every chat pod's :8000"
apiVersion: monitoring.coreos.com/v1
kind: ServiceMonitor
metadata: { name: chat-engine, labels: { release: kube-prometheus } }
spec:
  selector: { matchLabels: { app: chat } }
  endpoints: [ { port: http, path: /metrics, interval: 15s } ]
---
# monitoring/rules-slo.yaml -- the SLI, precomputed every 30 s
groups:
- name: slo.rules
  rules:
  - record: job:slo_ttft_good:ratio_rate5m
    expr: sum(rate(vllm:time_to_first_token_seconds_bucket{le="0.5"}[5m]))
          / sum(rate(vllm:time_to_first_token_seconds_count[5m]))

$ promtool check rules monitoring/rules-slo.yaml     # CI lane 1, extended
  SUCCESS: 1 rules found
$ kubectl -n monitoring get servicemonitors; open :9090/targets
  chat-engine ... 4/4 targets UP   dcgm ... 2/2 UP   blackbox ... 3/3 UP
\end{Verbatim}
\end{terminalbox}
\figcaption{실행 기록 14-3  스택 연결하기. 서비스마다 ServiceMonitor 하나를 두고, SLI용 기록 규칙을 만들며, promtool을 파이프라인의 새 게이트로 추가한다. 규칙도 릴리스 산출물이므로 13장의 규율은 모델 가중치에서 끝나지 않는다. /targets 페이지는 “모든 대상을 스크레이프하고 있는가?”를 확인하는 출석부다. up == 0 자체도 경보를 걸 수 있는 시계열이다.}

Grafana의 사용법은 오후 한나절이면 배울 수 있다. \textbf{패널}은 PromQL 질의와 시각화를 합친 것이고, \textbf{대시보드}는 패널에 시간 범위와 템플릿 변수(\code{\$pod}, \code{\$model})를 더한 것이다. 전체 구성을 JSON으로 직렬화할 수 있으며, 실습에서는 이 JSON을 \textit{저장소 안에} 보관한다. 대시보드는 운영 코드이기 때문이다. 변경할 때 검토하고, 개선할 때 버전을 남기며, Grafana 파드가 다시 스케줄되어도 복구할 수 있어야 한다. UI에만 존재하는 대시보드는 13장의 금요일 배포 사고를 기다리는 것과 같다. 이전 모습이 어땠는지 아무도 말할 수 없다. 사용법만 익혀서는 배울 수 없는 것이 \textbf{설계}다. 서빙 대시보드는 흔히 \textit{모든 것을 쌓은 벽}이 되어 실패한다. 알파벳순으로 배열한 패널 마흔 개는 하나씩 보면 모두 사실이지만, 새벽 2시에 한꺼번에 읽을 수는 없다. 이 강좌에서 가르치는 설계 원칙은 그림 14-4와 같다. \textbf{질문마다 행 하나를 두고, 당직자의 런북이 묻는 순서 그대로 배열한다.}

\begin{tbbfigure}
  \adjustbox{max width=0.94\linewidth,max totalheight=0.46\textheight}{\includegraphics{figures/fig14-4_dashboard.png}}\\[3pt]
  \figcaption{그림 14-4  화요일 21시 44분에 멈춰 세운 네 행 서빙 대시보드. 1행은 “약속을 지키고 있는가?”에 답한다. SLO 문장을 예산, 소진율, 플릿 p99, 스트림별 속도 패널로 옮겼다. 2행은 “누가 고통받는가?”에 답하며 레플리카별 RED를 보여 준다. 3행은 “왜 그런가?”에 답하며 대기열, KV, 디코드 속도, 선점 등 엔진 내부를 보여 준다. 4행은 “그 아래에는 무엇이 있는가?”에 답하며 UUID별 DCGM, GPU당 파드, 오토스케일러 태세를 보여 준다. 둘째 날 밤의 조사 순서는 ①→④로 표시했다.}
\end{tbbfigure}

\textbf{1행은 “약속을 지키고 있는가?”에만 답한다.} 패널은 SLO 문장의 각 항목을 그대로 옮긴다. 30일 오류 예산 게이지에는 61\%가 남아 있다. 7장에서 제품 관리자와의 논쟁을 끝내 줄 것이라고 약속한 바로 그 그래프다. 14.4× 호출 임계선을 그린 소진율, 올바른 방식으로 계산한 플릿 TTFT p99(질의 2를 기록 규칙으로 만든 것), 15 이상이라는 조항과 비교한 스트림별 tokens/s p50도 있다. §14.4의 용어로 이곳의 모든 값은 사용자가 볼 수 있는 약속 위반, 즉 \textit{증상}이다. 따라서 1행만 사람을 깨운다. \textbf{2행은 “누가 고통받는가?”에 답한다.} 같은 RED 신호를 레플리카별, 경로별로 나눈다. 모든 용량 사고에서 첫 진단 갈림길은 \textit{모두가 느린가?}와 \textit{일부만 느린가?}이기 때문이다. 전자는 무릎점, 물결, 폭풍 같은 수요 문제이고 후자는 이번 화요일과 같은 머신 문제다. 파드별 TTFT p99와 로드 밸런서의 요청률을 나란히 둔 패널이 한눈에 갈림길을 결정한다. 트래픽은 네 곳에 공평하게 나뉘지만 고통은 한 곳에 불공평하게 몰린다. \textbf{3행은 엔진의 언어로 “왜 그런가?”에 답한다.} 11장의 선들이 마침내 그래프가 된다. 파드별 \code{num\_requests\_waiting}은 선행 지표이자 11장에서 “바로 그 경보”라고 부른 선이다. KV 캐시 점유율은 플릿 전체에서 높게 고정되어 정상이다. 새벽 2시의 뇌가 같은 오해를 되풀이하지 않도록 패널에 정상이라고 주석을 달아 둔다. 디코드 단계 속도는 원인을 구분한다. 대기열이 생기면서 디코드도 느리면 머신 자체가 느려진 것이고, 디코드는 정상인데 대기열이 생기면 수요가 좌석을 추월한 것이다. 선점률과 초안 수락률은 각각 풀 압력과 추측의 건전성을 보여 준다. 9장에서 대기열 시간과 계산 시간을 별도 카운터로 받아 “14주차 대시보드 설계의 절반이 이미 끝났다”고 약속했다. 3행은 분류기에서 그 약속을 지키고 LLM으로 일반화한다. \textbf{4행은 “그 아래에는 무엇이 있는가?”에 답한다.} UUID별 DCGM을 GPU당 파드와 조인한다. 프로덕션의 모든 행이 1이어야 하는 표에서 화요일의 결정적 증거가 드러난다. 프레임버퍼, Xid, 노드 활력 징후, 오토스케일러가 원하는 파드 수와 실행 중인 파드 수도 보여 준다. 파드는 물결을 뒤쫓는다. 파드가 \textit{뒤쫓지 않는} 물결은 별도의 사고 유형이며, 12장의 안정화 구간이 평탄한 지연으로 드러난다.

이제 이 설계가 어떻게 값을 치르는지 보자. §12.4에서 예고한 사고를 이 네 행으로 실시간 해결한다. 13장 마지막 상자에서 \textit{라이브로} 보여 주겠다고 약속한 시연이다. \textbf{21시 42분, 호출}(①): chat-2에서 QueueSustained가 발생한다. 호출 본문에는 런북 링크 R2가 있고 대시보드는 이미 경보 레이블로 필터링되어 있다. 1행에서 소진율이 7.2×로 계속 오르고 플릿 p99가 3.8 s임을 확인한다. 30초면 호출이 진짜임을 알 수 있고 질의는 하나도 입력하지 않는다. \textbf{②, 2행의 갈림길:} 레플리카별 TTFT p99는 0.42 / 9.8 / 0.39 / 0.44다. 모두가 아니라 \textit{일부}만 느리다. 레플리카별 req/s는 모두 5.7이므로 로드 밸런서는 공정하며 고통의 원인은 chat-2 자체에 있다. 모두와 일부를 가른 순간 “무릎점인가? 폭풍인가? 배포 때문인가?”라는 불안의 나무 전체가 사라진다. 배포는 없었고 수요 문제도 아니다. \textbf{③, 3행의 메커니즘:} chat-2의 대기열은 19지만 KV는 다른 파드와 같은 96\%이므로 풀 문제가 아니다. 선점도 늘지 않아 요동치는 상태가 아니다. 하지만 디코드 단계 p50은 플릿의 70에 비해 150 ms다. \textit{머신 자체}가 2.1× 느리다. 10장의 루프라인과 11장의 20 ms 하한을 떠올리면, 메모리 병목 루프의 경우 다른 무언가가 이 카드의 메모리 버스를 마시고 있다는 뜻이다. \textbf{④, 4행의 기반:} UUID당 파드 표에는 GPU …3d9f: 2(chat-2, eval-nightly-28471)라고 나온다. 이웃이 여기 있다. \code{kubectl delete job eval-nightly}를 실행하고, 노드의 오래된 시분할 ConfigMap을 조사하도록 격리 메모를 남긴다. 21시 50분이 되면 대기열이 비워지고 소진율은 1× 아래로 내려가 호출이 자동으로 해제된다. 사람이 들인 시간은 8분이고, 그중 입력에 쓴 시간은 2분이다. \textit{“chat-2만 문제, GPU 공유”}라는 두 줄 진단은 6장의 기념 문구가 염원에서 워크플로로 발전한 결과다. 행 자체가 두 줄을 작성한다. 2행이 첫 줄을, 4행이 둘째 줄을 쓴다.

재현 과정을 정직하게 만드는 세 가지 주석이 있다. 첫째, 21시 31분의 \textbf{원인 티켓}(GPU당 파드 = 2)은 호출보다 11분 먼저 범인을 지목했다. 둘째 날 밤이라면 사용자가 고통을 느끼기 전에 당직자가 Slack에서 읽을 수도 있었다. 그래도 호출 대신 티켓을 만드는 이유는 §14.4에서 설명한다. 원인은 셀 수 없이 많고 저렴하며 해가 없는 경우도 흔하다. 증상은 적고 비싸며 정의상 실제로 발생한 일이다. 약속 위반에 호출을 보내고 원인은 주석으로 붙인다. 둘째, 대시보드에 \textit{없었던} 것을 살펴보자. “시분할 충돌”이라는 패널은 없었다. 아무도 이 실패를 예측하지 못했지만 네 행은 답을 찾아냈다. 각 행이 알려진 실패를 열거하는 대신 \textit{질문}에 답하도록 설계되었기 때문이다. 이것이 관측 가능성의 역할이다. Aaron은 훈련에서만 본 잡음을 알아보았다. 셋째, 수정은 \textbf{기록}을 남겼다. 13장의 런북이 있는 레지스트리 항목에 사고 메모를 등록했고, 이제 WATCH 단계에는 “/metrics를 curl로 받아 눈을 가늘게 뜨고 본다” 대신 “대시보드 1\textasciitilde{}2행”이라고 적혀 있다. §14.4의 사후 분석 조치로 새 경보도 하나 추가했다. 실습 문제였던 GPU당 파드 규칙이 플릿 정책이 되었고 패널 캡션 하나도 수정했다. 사고 후에도 바뀌지 않은 대시보드는 아무것도 배우지 못한 대시보드다. 사고가 계측을 개선하는 순환 덕분에 모니터링의 가치가 \textit{복리로} 쌓인다. 이 순환은 10장 실습에서 사고 검토자가 읽을 수 있도록 게이트 보고서를 저장하게 했을 때 조용히 시작되었다.

\begin{calloutbox}{tbbred}{tbbxFEF2F2}
{\normalsize\bfseries\color{tbbx991B1B}상자 14-1  서빙 대시보드의 일곱 가지 치명적 죄악}\par\vspace{3pt}
\begin{tbbitemize}
  \item \textbf{전광판을 차지한 평균값.} 치우친 분포 위에 평균을 제목처럼 내세우면 사례 14-A의 “40\% 상승, 사용자가 과장한다”라는 결론에 이른다. 언제나 백분위수나 히스토그램을 써야 한다. 평균은 누구도 설명하지 못한다.
  \item \textbf{백분위수의 평균.} 레플리카별 p99에 avg()를 적용하면 실제 9.08 s 대신 2.76 s가 나온다. 버킷을 합친 뒤(sum by le) 분위수를 구해야 한다. 압박 속에서 즉석 질의를 만들 때도 예외는 없다.
  \item \textbf{사용률을 여유 용량으로 해석하기.} GPU 사용률은 점유 듀티 사이클이다. 바쁨 ≠ 내 것 ≠ 효율적임이다. 사용률이 의미하는 척하는 정보는 포화 신호인 대기열, 선점, DRAM 활동에 담겨 있다.
  \item \textbf{그래프에 SLO 선이 없음.} 0.5 s 임계선을 그리지 않은 p99 패널은 새벽 2시의 기억력에 판단을 떠넘긴다. 약속을 측정하는 모든 패널에 그 약속을 선으로 그려라.
  \item \textbf{런북 없는 경보, 질문 없는 패널.} 모든 호출은 런북으로 연결되고, 모든 런북 단계는 특정 행을 지목하며, 모든 패널은 런북이 그 질문을 하기 때문에 존재해야 한다. 그 밖의 것은 장식일 뿐이다.
  \item \textbf{모든 것을 쌓은 벽.} 알파벳순 패널 마흔 개는 어떤 질문에도 답하지 못한다. 약속 → 누구 → 원인 → 기반의 순서로 배치한 네 행은 위에서 아래로 1분 안에 읽을 수 있다. 지표마다 패널 하나가 아니라 독자마다 대시보드 하나를 만든다.
  \item \textbf{버전 관리하지 않는 대시보드.} JSON을 저장소에 두고 코드처럼 검토해야 한다. 그렇지 않으면 파드가 한 번 다시 스케줄되는 순간 계기판이 기억을 잃는다. 모니터링판 금요일 배포 사고다.
\end{tbbitemize}
\end{calloutbox}

\begin{calloutbox}{tbbpurple}{tbbpurplefill}
{\normalsize\bfseries\color{tbbpurple}생각해 보기}\par\vspace{3pt}
\begin{tbbitemize}
  \item 2행이 없다고 가정하고 둘째 날 밤의 조사 과정을 재구성하라. 어떤 질문에 답하기가 어려워지고, 당직자는 어떤 순서로 확인하게 되며, 첫째 날 밤에 걸린 49분 가운데 얼마나 되돌아오는가? 이는 “어느 행을 없앨 것인가?”라는 채점 문제의 본뜻이다. 답은 취향이 아니라 논증이어야 한다.
  \item 플릿 p99 패널(질의 2)과 SLI 비율 패널(질의 4)은 서로 다른 결론을 낼 수 있다. p99는 0.5 s 이하인데 SLI 비율은 목표보다 낮은 5분 구간과, 그 반대인 구간을 각각 구성하라. 호출기와 연결된 1행에는 어느 값을 넣어야 하는가? §14.4가 비율을 선택하는 이유는 무엇인가?
  \item 팀원이 “비즈니스 — 가입자 수, 판매한 토큰 수, 분당 매출”이라는 다섯 번째 행을 원한다. 독자, 갱신 주기, 조치하는 사람을 고려해 \textit{이} 대시보드에 넣어야 한다는 주장과 넣지 말아야 한다는 주장을 모두 제시하라. 12장의 단위 경제성 선이 실제로 속해야 할 위치도 설명하라.
  \item GPU당 파드 패널에는 UUID 레이블이 붙은 DCGM과 장치 플러그인의 파드 매핑을 조인해야 했다. 클러스터의 DCGM 익스포터에 파드 레이블이 없다면 당직자가 21시 44분에 수행할 두 단계 수동 조인을 설계하고, 첫날 밤의 49분과 비교해 몇 분이 더 드는지 추정하라.
\end{tbbitemize}
\end{calloutbox}

\section{경보: Page, 소진율, 거짓 경보를 피하는 기술}

Dashboard가 기억이라면 \textbf{alert}는 주의력이다. 어느 시간이든 사람을 지금 당장 깨워야 한다고 대신 결정하는 장치다. 이 결정에도 serving path와 같은 수준의 엔지니어링 엄밀함이 필요하다. 역사상 가장 좋은 사례는 가장 가까운 pager에서 40만 km 떨어진 곳에서 일어났다. 1969년 7월 20일, 최초의 달 착륙을 위한 powered descent가 시작된 지 5분 뒤 guidance computer에 \textbf{1202}가 떴다. Capsule 안에서 아무도 외우지 않은 alarm code였다. 연료는 타고, 표면은 다가오는데 crew의 화면은 뜻을 알 수 없는 네 자리 숫자만 내놓았다. 30초도 지나지 않아 지상에서 GO라는 판단이 왔고, 근거는 \textbf{손으로 쓴 cheat sheet}였다. 3주 전 훈련 simulation에서 관제팀은 자매 alarm인 1201을 보고 당황해 abort를 선언했다. 나중에 \textit{잘못된 판단}으로 드러났고 simulation supervisor는 교훈을 뼈아프게 남겼다. Backroom의 24세 엔지니어 Jack Garman은 집에 돌아가 모든 program alarm code를 손으로 쓰고, 각각이 문제가 되는 조건을 옆에 적었다. Executive overflow 계열에 붙인 메모의 뜻은 이랬다. \textit{너무 자주 반복되지 않고 guidance가 계속 작동하면 GO.} 실제로 1202가 떴을 때 front room의 Steve Bales가 Garman의 판정을 전달했고 Armstrong은 하강을 계속했다. Alarm은 이후 세 번 더 울렸지만 아무 문제도 생기지 않았다. Computer가 overload 상태임을 알리면서도 \textit{설계대로 우선순위가 낮은 작업을 버리고} 있었기 때문이다. Margaret Hamilton의 priority executive는 4 KB RAM 위에서 제9장의 부하 차단 원칙을 실행하고 있었다.

이 이야기에는 이 절의 모든 규칙이 들어 있다. Alarm의 의미는 \textbf{미리 글로 정했다.} \textit{실패한 drill}에서 태어난 runbook이다. 실습의 두 번 깨뜨리기에는 반세기가 넘는 전통이 있는 셈이다. 판정은 \textbf{존재가 아니라 비율과 지속성}을 기준으로 했다. "너무 자주 반복되지 않는다"는 1969년판 \code{for: 5m}과 rate threshold다. 무엇이든 spike 한 번만으로 page를 울릴 일은 거의 없다. Alarm이 울린 system은 \textbf{설계대로 우아하게 성능을 낮추고 있었다.} 어떤 alarm은 설계가 \textit{제대로 작동함}을 나타낸다. KV cache가 높게 고정되어 있거나 executive가 task를 버리는 경우다. 이런 상태마다 page를 보내면 사람은 page를 무시하도록 훈련된다. 이 원칙을 어길 때 사람이 치르는 대가는 우주 비행 밖에서 실제 사망자 수로 나타난다. 지금껏 만들어진 가장 빽빽한 alarm 환경인 병원 telemetry 병동은 업계에 \textbf{경보 피로(alert fatigue)}라는 용어를 가르쳤다. Joint Commission의 2013년 sentinel-event alert는 3년 반 동안 alarm 관련 사고 98건, 그중 사망 사고 80건을 집계했다. 임상 alarm의 \textbf{85–99\%는 아무 조치도 필요하지 않다}는 연구도 인용했다. Beep에 빠진 간호사는 합리적으로 beep가 아무 의미도 없다고 학습하고, 그러다 어느 하나가 실제 의미를 갖는다. 엔지니어링으로 옮기면 잔인하면서도 명쾌하다. \textit{행동할 수 있고, 긴급하며, 실제인} 조건이 아닌 모든 alert는 팀이 정말 중요한 하나를 놓치도록 적극적으로 훈련한다. Page는 정보가 아니라 잠에서 깨라는 명령이다. 적게 발행해야 한다.

\subsection*{증상은 page로, 원인은 ticket으로}

"적게 발행하라"를 운영 규율로 만드는 것은 Rob Ewaschuk의 \textit{My Philosophy on Alerting}와 뒤이은 SRE 책이 정착시킨 두 단어짜리 routing 원칙이다. \textbf{증상은 page로, 원인은 ticket으로 보낸다.} \textit{증상}은 사용자가 볼 수 있는 약속 위반이다. SLI가 피를 흘리고, queue가 요청을 빠뜨리며, 바깥에서 문이 열리지 않는 상태다. \textit{원인}은 언젠가 증상을 만들 수 있는 조건이다. 새는 레플리카, 오래된 ConfigMap, GPU 하나에 올라간 pod 두 개가 예다. 이 규칙을 정당화하는 비대칭이 있다. 증상은 \textbf{적고, 정의상 실제이며, 완전하다.} 상상하지 못한 것을 포함해 사용자를 해치는 모든 실패는 결국 SLI에 나타나야 한다. 화요일의 time-slice 충돌은 alert rule을 작성한 뒤 생겨났지만 queue page가 잡아냈다. 반면 원인은 \textbf{무수하고, 저렴하며, 흔히 무해하다.} \textit{Dev} pool에서 pod 둘이 같은 GPU를 쓰는 일은 화요일에도 정상이다. §12.4가 그곳에서는 time-slicing을 정직하게 권했다. 원인마다 page를 울리면 새벽 3시의 무의미한 사건에 무감각해진 병원 병동 같은 rotation이 된다. 그래서 화요일의 GPU별 pod alert는 \textbf{ticket}을 만들었다. 21:31에 Slack에서 볼 수 있었고, \textit{증상이 관련성을 부여한 순간} 엄청난 가치를 발휘했지만 pager는 울리지 않았다. 그날 밤 유일한 page는 queue였다. 사용자가 이미 기다리고 있던 선행 증상이라 24분 늦게 울릴 burn-rate page라는 후행 증상을 앞섰다. 제3장의 memory.high처럼 벽 앞에 tripwire를 놓는 방식이다. 표 14-1이 지속 gauge 가운데 유일한 PAGE를 제11장에서 바로 이 역할을 맡긴 \code{num\_requests\_waiting}에 배정한 이유다.

\subsection*{SLO 문장 컴파일하기: 오류 예산과 소진율 계산}

제13장은 이 장에서 "느낌이 아니라 소진율에 따라 page가 울리게" 만들겠다고 약속했다. 이제 과정 SLO의 availability 조항을 처음부터 끝까지 compile해 보자. 11주 차 계약은 \textbf{TTFT p99 ≤ 500 ms}라고 말한다. 이는 \textit{turn의 99\% 이상이 0.5초 안에 첫 token을 받는다}는 문장과 같다. 비율로 다시 쓰면 \textbf{SLI}가 되고(터미널 기록 14-2의 query 4), 나머지 1\%는 \textbf{오류 예산}이 된다. 플릿은 한 달에 약 950만 turn, 즉 평균 답변 약 211 token으로 20억 token을 서빙하므로 이 약속은 \textbf{30일마다 나쁜 turn 약 95,000개}를 허용한다. 제7장이 시간 기준 99.9\%를 43분이라는 자원으로 표현한 것과 같다. 단위가 달라졌음에 주목하자. 요청 기준 budget은 분이 아니라 turn을 세므로 하루 주기의 traffic 파동에서도 의미가 유지된다. 새벽 3시의 나쁜 \textit{1분}은 거의 쓰지 않지만 21:30의 나쁜 1분은 큰 비용을 쓰며, 사용자 고통을 이렇게 가격화하는 것이 옳다. \textbf{소진율}은 소비 속도다. Burn = 나쁜 turn 비율 ÷ (1 − SLO)다. 1×로 소진하면 정확히 30일 만에 budget을 모두 쓴다. 손익분기 속도인 셈이다. 화요일 21:44에 chat-2가 도착 요청의 4분의 1을 빠뜨렸으므로 나쁜 비율은 약 24\%, burn은 약 24×였다. 한 달 budget을 30시간 만에 쓰는 속도다. 그림 14-5는 tank와 drain을 그린다.

\begin{tbbfigure}
  \adjustbox{max width=0.94\linewidth,max totalheight=0.46\textheight}{\includegraphics{figures/fig14-5_burnrate.png}}\\[3pt]
  \figcaption{그림 14-5  Budget tank(나쁜 turn 95,000개), 네 가지 drain 속도, 여러 window 정책, 1시간 burn panel에 그린 화요일의 모습. Queue page가 burn page보다 24분 먼저 울렸고, 점선은 완화하지 않았을 때 burn page가 울릴 시점인 22:06을 나타낸다.}
\end{tbbfigure}

Burn 경보는 window를 다루는 게임이다. 실습이 그대로 채택한 SRE Workbook의 표준 정책은 심각도마다 \textbf{window 두 개를 짝지어} 사용한다. Burn이 \textbf{1시간과 5분 모두에서 14.4× 이상}이면 page를 울린다. 이 속도는 한 시간마다 월간 budget의 2\%를 쓴다. \textbf{6시간과 30분에서 6×}면 page를 울린다. 6시간마다 5\%를 쓰는 속도로, 빠른 rule이 놓치는 느린 출혈을 잡는다. \textbf{하루 동안 3×}, \textbf{3일 동안 1×}면 낮에 처리할 trajectory 경고인 ticket을 만든다. 각 짝은 기질의 절충을 하나 담는다. \textbf{긴 window는 의심}이다. 깊이가 아무리 커도 90초짜리 blip은 1시간 평균을 14.4× 너머로 움직일 수 없어 page를 울리지 않는다. Apollo의 "너무 자주 반복되지 않는다"를 산술로 옮긴 것이다. \textbf{짧은 window는 용서}다. 사고가 끝나면 5분 평균이 즉시 내려가 page가 해제된다. 긴 window가 빠지는 동안 나머지 한 시간 내내 울리지 않는다. 의심의 정직한 비용은 latency다. 화요일의 24× burn에서 1시간 평균이 14.4×를 넘는 데 약 35분이 필요했다(24 × 35/60 = 14). 그래서 alert 목록은 burn에만 의존하지 않는다. Burn page는 \textbf{완전성 보장}이다. 이미 알든 새롭든 사용자가 볼 수 있는 모든 실패가 결국 이를 울린다. Queue page는 이 플릿이 실제로 겪는 실패의 \textbf{빠른 경로}다. Capacity 형태의 고통은 제7장의 knee와 제11장의 admission line처럼 queue에 먼저 나타난다. 빠른 경로는 허용하지만 증상에만 준다.

\begin{terminalbox}
\begin{Verbatim}[breaklines=true,breakanywhere=true,fontsize=\small,formatcom=\color{tbbtermfg},codes=\tbbcodeactive]
# monitoring/rules-alerts.yaml (excerpt) -- the roster's teeth
- alert: SLOBurnFast                              # completeness guarantee
  expr: >
    (1 - job:slo_ttft_good:ratio_rate1h) / 0.01 > 14.4
    and (1 - job:slo_ttft_good:ratio_rate5m) / 0.01 > 14.4
  labels: { severity: page }
  annotations: { runbook: runbooks/R1-slo-burn.md }

- alert: QueueSustained                           # fast path (symptom)
  expr: avg_over_time(vllm:num_requests_waiting[5m]) > 8
  for: 5m
  labels: { severity: page }
  annotations: { runbook: runbooks/R2-queue.md, dashboard: row 3 }

- alert: DuplicatePodsOnProdGPU                   # cause (tuesday's vaccine)
  expr: count by (UUID) (DCGM_FI_DEV_GPU_UTIL{pool="prod"}) > 1
  for: 2m
  labels: { severity: ticket }

- alert: Watchdog                                 # dead man's switch
  expr: vector(1)                                 # fires ALWAYS, on purpose
  labels: { severity: heartbeat }

# alertmanager.yaml (excerpt) -- routing, grouping, silence
route:
  group_by: [alertname, model]      # one incident = one page, not 40
  routes:
  - match: { severity: page }    receiver: oncall-phone
  - match: { severity: ticket }  receiver: slack-serving-alerts
  - match: { severity: heartbeat } receiver: deadman-external
$ amtool silence add alertname=QueueSustained --duration=2h \
    --comment "drill A in progress (lab 14)"      # planned breakage, quiet
\end{Verbatim}
\end{terminalbox}
\figcaption{터미널 기록 14-4  Alert는 심각도를 가진 query이고 routing은 정책이다. 형태에 주목하자. Burn 짝은 recording rule을 읽고, queue page에는 지속 조건(for: 5m)이 필요하며, 원인 rule은 ticket을 만든다. Watchdog는 영원히 울리므로 heartbeat를 기다리는 외부 receiver가 감지한 \textbf{침묵}은 monitoring stack 자체가 멈췄다는 뜻이다. 감시자는 누가 감시하는가? 사람이 쓰러지면 놓이는 switch가 감시한다.}

\begin{tbbtablewrap}
  \begin{tabular}{@{}>{\raggedright\arraybackslash}m{\dimexpr 0.2500\tbltextwidth-2\tabcolsep\relax} >{\raggedright\arraybackslash}m{\dimexpr 0.3085\tbltextwidth-2\tabcolsep\relax} >{\raggedright\arraybackslash}m{\dimexpr 0.1223\tbltextwidth-2\tabcolsep\relax} >{\raggedright\arraybackslash}m{\dimexpr 0.3191\tbltextwidth-2\tabcolsep\relax}@{}}
  \arrayrulecolor{tbbnavy}\specialrule{1pt}{0pt}{0pt}
  \rowcolor{tbbtablehead}{\centering \bfseries\color{tbbnavy}경보} & {\centering \bfseries\color{tbbnavy}조건(개요)} & {\centering \bfseries\color{tbbnavy}심각도} & {\centering \bfseries\color{tbbnavy}존재 이유 · 계보} \\
  \arrayrulecolor{tbbnavy}\specialrule{0.7pt}{0pt}{2pt}
  {\textbf{SLOBurnFast}} & {1시간 ∧ 5분 동안 burn ≥ 14.4×} & {PAGE} & {약속이 빠르게 무너짐. 상상하지 못한 실패도 잡는다(제7장 → 여기).} \\
  \arrayrulecolor{tbbtableborder}\hline
  {\textbf{SLOBurnSlow}} & {하루 동안 burn ≥ 3×(그리고 3일간 1×)} & {ticket} & {궤적 경고. 낮에 의도적으로 budget을 사용한다.} \\
  \arrayrulecolor{tbbtableborder}\hline
  {\textbf{QueueSustained}} & {어느 pod든 5분 동안 waiting > 8} & {PAGE} & {빠른 경로. Capacity 고통은 queue에 먼저 온다(제11장의 alarm line).} \\
  \arrayrulecolor{tbbtableborder}\hline
  {\textbf{EndpointsEmpty}} & {chat Service의 ready address = 0} & {PAGE} & {Routing 대상이 없으면 지금 모든 사용자가 제6장의 503을 본다.} \\
  \arrayrulecolor{tbbtableborder}\hline
  {\textbf{ProbeFail}} & {어느 probe든 probe\_success = 0이 3회 연속} & {PAGE} & {밖에서 문이 열리지 않음. 플릿이 볼 수 없는 계층이다(제6장).} \\
  \arrayrulecolor{tbbtableborder}\hline
  {\textbf{PreemptionSurge}} & {10분 동안 rate(preemptions) > 0.5/s} & {ticket} & {KV pool thrashing. Prompt-length drift가 도착할 때가 많다(§14.5).} \\
  \arrayrulecolor{tbbtableborder}\hline
  {\textbf{RestartsOOM}} & {Δrestarts ≥ 3/h 또는 OOMKilled 발견} & {ticket → 플릿 전체면 page} & {Self-healing이 leak을 감춘다(제5장 목록, 제3장 exit 137).} \\
  \arrayrulecolor{tbbtableborder}\hline
  {\textbf{DuplicatePodsOnProdGPU}} & {2분 동안 prod UUID별 pod > 1} & {ticket} & {화요일을 막는 백신. 원인이므로 깨우지 않고 주석을 남긴다.} \\
  \arrayrulecolor{tbbtableborder}\hline
  {\textbf{Watchdog}} & {vector(1), 항상 firing} & {heartbeat} & {Dead man’s switch. Cluster 밖 receiver가 그 침묵을 page로 보낸다.} \\
  \arrayrulecolor{tbbnavy}\specialrule{1pt}{2pt}{0pt}
  \end{tabular}\par
  \figcaption{표 14-2  경보 목록. Page 네 개는 모두 증상이고 runbook을 갖췄다. Ticket 네 개는 모두 원인이며 heartbeat가 하나 있다. 수백 개가 울리는 병동과 비교해 보자. 자기 pager를 믿는 rotation이 이 표의 결과물이다.}
\end{tbbtablewrap}

\subsection*{Synthetic probe: 거부할 수 있는 모든 계층의 바깥에 서기}

앞의 모든 장치는 플릿의 벽 안에서 플릿을 바라본다. 제6장은 그 시야가 얼마나 많이 놓치는지를 실습 하나로 증명했다. LB가 구현되지 않은 path를 health check하고, TLS가 만료되며, perimeter timeout이 자기 것보다 짧은 실패에서는 사용자가 503을 보는데 kubectl은 결백하다고 말한다. 증거가 \textit{cluster 바깥} console에 있기 때문이다. 모니터링의 답은 \textbf{synthetic probe}이며, 세 가지를 설계하라고 했던 제6장 연습문제의 공식 해답이 여기에 있다. 첫째는 \textbf{문}이다. Blackbox exporter가 public gateway URL에 30초마다 HTTPS \code{GET /\allowbreak v1/\allowbreak models}를 보낸다. DNS, TLS, LB target health, routing을 모두 거치며, \textit{LB 바깥에서 실행하지 않으면} 아무것도 검증하지 못한다. 검사 대상과 운명을 공유하는 probe는 증인이 아니라 거울이다. 매우 신중한 팀은 두 번째 region에서 실행한다. 둘째는 \textbf{답변}이다. 1분마다 10초 timeout으로 token 하나짜리 completion을 요청한다. Admission과 engine을 포함한 전체 경로를 거치며 문 probe가 그대로 통과하는 504/hang 부류를 잡는다. 셋째는 \textbf{stream}이다. 5분마다 짧은 streaming completion을 실행해 2초 안에 첫 token을 받고 마지막 chunk가 온전해야 한다. 제6장의 표 6-3이 답변 중간의 침묵 원인으로 지목한 \textit{idle-timeout} 계층 바깥에 서는 stream-cut detector다. Probe 세 개가 서로 다른 세 계층을 맡고 각각 책임을 물을 계층 바깥에 선다. \code{probe\_success}가 3회 연속 실패하면 page를 울린다. 한 번의 존재만으로 울리면 모든 network 재채기마다 page가 온다. Garman의 원칙이 다시 등장한다. 이들은 과정에서 가장 오래된 관측 가능성의 빈틈을 함께 닫는다. 사례 14-A의 uptime checker가 밤새 초록색이었던 이유는 \textit{첫 번째 probe만 있었기} 때문이다. 문이 살아 있다는 사실만 알렸고 뒤의 queue에 대해서는 아무 말도 하지 않았다. Synthetic은 과정에서 처음 LB target-health console로 보냈을 때의 약속을 완성한다. \textbf{관측 가능성은 거부할 수 있는 모든 계층을, 그 계층이 조작할 수 없는 시점에서 포괄해야 한다.}

\subsection*{사람이 맡는 절반: On-call, runbook, postmortem과 숫자의 가치}

Page는 사람에게 도착하며, 그 사람 주위의 장치도 PromQL만큼 세심하게 설계해야 한다. \textbf{On-call rotation}은 교대로 pager를 들고 명시적인 계약을 따른다. 몇 분 안에 응답하고, 부담 없이 escalate하며, 인수인계를 글로 남긴다. 표 14-2가 매일 밤이 아니라 한 달에 한 번쯤 page를 울리기 때문에 지속 가능한 구조다. \textbf{Page에는 runbook이 붙는다.} 제13장의 template에 이제 이가 생겼다. 증상 → 답을 주는 dashboard 행 → 되돌리기 쉬운 순서로 나열한 안전한 완화책 → escalation 연락처가 이어진다. 제13장에서는 \code{/\allowbreak metrics}와 script만 가리킬 수 있었던 runbook의 WATCH 단계가 이제 \textit{"1–2행, burn과 레플리카별 지표"}라고 적힌다. 해결 뒤에는 항공 분야와 John Allspaw의 노력을 통해 업계가 받아들인 \textbf{비난 없는 postmortem}이 이어진다. Blameless는 예의가 아니라 지식의 문제다. 개인은 일정한 비율로 실수하고 system이 그 가격을 정한다는 제13장의 논증처럼, 사람 이름으로 끝난 postmortem은 아무것도 배우지 못하지만 \textit{작동 원리}로 끝나면 action item이 생긴다. 화요일 사고에서는 하루 안에 세 가지를 기록했다. GPU별 pod rule을 실습 문제에서 플릿 정책으로 승격해 원인 alert답게 ticket으로 만들고, 오래된 time-slice ConfigMap에 owner와 CI expiry check를 부여했으며, dashboard caption 하나를 "busy ≠ yours"로 고쳤다. 보고서 자체는 v9-int4의 \textbf{registry 항목에 연결했다.} 제10장 실습이 언젠가 사고 검토자가 저장된 gate report를 읽을 것이라 했고, 바로 이 사람이 파일을 닫는다. 사고 → 계측이라는 순환이 모니터링에 \textit{복리}를 만든다. 불이 날 때마다 smoke detector 하나가 남는다.

장치 전체를 평가하는 숫자는 둘이며 몇 주 전부터 이 책이 열어 둔 장부를 닫는다. \textbf{MTTD}는 mean time to detect, 즉 사고 시작부터 사람이 알 때까지이고, \textbf{MTTR}은 mean time to restore, 즉 사고 시작부터 사용자가 정상 상태를 되찾을 때까지다. 그림 14-1의 36+49와 12+8 막대가 이를 보여 준다. MTTR은 DORA의 네 번째 key이며 제13장에서 나머지 셋의 가격을 계산했다. 과정의 배포 이야기는 네 key를 추정이 아니라 계측한 상태로 끝난다. 이 분포는 \textit{설계 입력}이기도 하다. 제13장 논술은 blue/green warm standby 시간을 어떤 data로 조정할지 물었다. 이제 답이 dashboard에 있다. 배포 형태 regression의 MTTD가 12분이면 비용이 2배인 warm standby를 한 시간 유지해도 약 20분 이후에는 얻는 것이 없으므로 window를 줄일 수 있다. 반대로 support 전화로 36분 뒤에야 감지하는 플릿은 예전 color를 일찍 없앨 수 없다. 제7장이 논술로 냈던 새벽 2시 상황, 즉 \textit{SLO를 위반했고 demo까지 6일 남았을 때 어느 조항을 적용할 것인가?}도 마침내 읽기 문제가 된다. 1행은 budget이 61\%이고 burn이 증가함을 보여 주며, 7주 차 결과물인 SLO 문서는 미리 합의한 완화와 owner를 밝힌다. 결정은 논술이 약속한 형태, 즉 논쟁이 아니라 산술과 서명으로 내려진다. 마지막 원장 항목도 확정하자. Prometheus, Grafana, Alertmanager, exporter, 50 GiB TSDB로 이루어진 monitoring stack은 cluster의 남는 CPU를 이용해 한 달 약 \textbf{\$23}가 든다. 합계는 \$877 + \$23 = \textbf{정확히 상한인 \$900}다. 과정 청구서의 마지막 항목은 다른 모든 항목을 감시하고, 그림 14-1 caption이 이미 비용 대비 효과를 계산했다. 이 stack은 첫 화요일에 한 달 값을 모두 갚았다.

\begin{calloutbox}{tbbpurple}{tbbpurplefill}
{\normalsize\bfseries\color{tbbpurple}생각해 보기}\par\vspace{3pt}
\begin{tbbitemize}
  \item Garman의 원칙은 “너무 자주 반복되지 않으면 GO”였다. 이 절에서 “너무 자주”를 표현하는 세 가지 서로 다른 장치(for:, avg\_over\_time, 짝지은 window)를 alert-rule 문법으로 옮기자. 각 장치는 어떤 실패 부류 때문에 page가 울리는 일을 막는가?
  \item “Queue depth는 원인이고 오직 SLI만 증상이다”라고 주장하는 엄격한 증상 우선주의자에게 QueueSustained page를 변호하자. 그 입장은 어느 지점에서 사용자에게 비용을 떠넘기는가? 빠른 경로가 정직하려면 queue threshold가 어떤 성질을 가져야 하는가? 힌트: SLI가 울리지 않을 때도 이 threshold가 울릴 수 있는가? 12주 차 autoscaler가 seat 수를 바꾸면 이 결합에는 어떤 유지보수가 필요한가?
  \item 화요일 사고의 inhibition rule을 설계하자. Chat-2에서 QueueSustained가 page를 울리는 동안 Alertmanager가 어떤 ticket을 억제해야 하며, 어떤 것은 절대 억제해서는 안 되는가? Watchdog 억제가 용납할 수 없는 configuration인 이유는 무엇인가?
  \item Dead man’s switch는 monitoring이 침묵하면 cluster \textbf{바깥} receiver를 통해 page를 울린다. Receiver가 (a) Slack webhook, (b) 만료된 card로 결제하는 phone service, (c) 같은 cloud region에서 hosting되는 경우에도 말없이 실패하는 것은 무엇인가? Probe와 heartbeat가 모두 지켜야 하는 shared fate의 일반 원칙을 도출하자.
\end{tbbitemize}
\end{calloutbox}

\section{드리프트: 절대 Page를 울리지 않는 실패}

지금까지 이 장에서는 \textit{시작점}이 있는 실패를 추적했다. 무언가 깨지는 순간이 있고 산술이 붙잡을 수 있는 window가 있었다. 하지만 registry가 몇 주 전 digest \code{a17b…}에서 model을 얼린 뒤에도, 어느 digest도 고정할 수 없는 input 하나가 남는다. 바로 세상이다. 사용자의 집단, 언어, 계절, 의도가 바뀌고, 뉴스는 존재하는 질문 자체를 바꾸며, 자기 성공도 찾아오는 사람을 바꾼다. 제12장 논술의 LMS integration처럼 학기 중간에 학생 1만 명이 플릿에 들어올 수 있다. Model은 고정된 어느 순간의 weights로 이 모든 변화를 맞는다. \textbf{Model은 바뀌지 않지만 세상은 바뀐다.} 둘 사이의 간격은 강줄기가 움직이듯 넓어진다. 절벽도 시작점도 alarm도 없다. 제13장은 blind spot을 정확히 이름 붙였다. \textbf{Canary는 경사면이 아니라 절벽을 잡는다.} 경사면은 여러 \textit{주} 동안 쌓은 production data로 만든 도구의 몫이라고 약속했다. 그 도구가 \textbf{drift detection}이다. 강의계획서가 이를 dashboard와 alarm 옆에 둔 것은 옳다. 아무것도 beep하지 않는 시간 규모에 기억과 주의력이라는 pattern을 적용하기 때문이다.

분류 체계는 대응할 때 가치를 발휘하므로 세 질문으로 익히자. \textbf{Data drift}(covariate shift)는 \textit{input} 분포 P(x)가 움직였는지를 묻는다. Prompt가 길어지고 한국어 비중이 두 배가 되며 새로운 주제 cluster가 나타난 경우다. 3B chat model에게 긴 prompt가 분포 밖 입력인 것은 아니므로 새 조합에서도 model 자체는 완벽히 유능할 수 있다. 보통 먼저 깨지는 것은 \textit{서빙 산술}과 \textit{평가 범위}다. §11의 prefill/KV budget은 이제 맞지 않는 input census를 기준으로 정했고, gate의 eval set은 예전 세상에서 표본을 뽑았기 때문이다. \textbf{Concept drift}는 같은 input에서 \textit{올바른 답} P(y|x)가 달라졌는지를 묻는다. "대통령은 누구인가?", "도서관 이용료는 얼마인가?"처럼 ground truth에 날짜가 붙는 질문이 예다. 어떤 serving metric도 concept drift를 볼 수 없다. Model은 유창하고 자신 있게 오래된 답을 내놓는다. 오직 지금 세상으로 새로 평가해야 잡을 수 있으며 대응 사다리의 둘째 칸이 존재하는 이유다. \textbf{Prior shift}는 \textit{과제의 구성} P(y)가 움직였는지를 묻는다. 시험 기간이 되면 일반 assistant가 proof checker가 되고, answer length 211과 refusal 0.8\% 같은 guardrail baseline이 조용히 중심을 옮긴다. 9월에 맞춘 threshold는 12월에 거짓 경보를 낸다. Drift detection은 \textit{다른} monitor를 정직하게 유지하는 장치이기도 하다. LLM에서는 세 부류를 잇는 결론이 하나 더 있다. \textbf{Eval set 자체도 drift한다.} 영원히 94.2\%를 기록하는 gate는 자기가 태어난 날의 세상을 측정한다. 성숙한 팀은 분기마다 production에서 eval slice를 다시 뽑고 baseline을 새로 잡으며, 그 diff 자체를 drift report로 남긴다.

\begin{tbbfigure}
  \adjustbox{max width=0.94\linewidth,max totalheight=0.46\textheight}{\includegraphics{figures/fig14-6_drift.png}}\\[3pt]
  \figcaption{그림 14-6  드리프트 계측. 왼쪽 위는 canary가 잡는 절벽과 여러 주가 지나야 보이는 경사면, 오른쪽 위는 플릿이 이미 내보내며 함께 움직이는 sentinel metric이다. 왼쪽 아래는 prompt length의 PSI를 손으로 계산한 결과인 0.144, 관찰 구간이다. 오른쪽 아래는 저렴한 대응에서 비싼 대응으로 올라가는 사다리다. 넷째 칸은 제13장의 pipeline으로 이어져 제1장의 순환 고리를 닫는다.}
\end{tbbfigure}

\subsection*{이미 내보내는 Sentinel과 손으로 계산할 수 있는 숫자 하나}

새로 만들기 전에 그림 14-6 오른쪽 위 panel부터 읽자. 11주 차부터 플릿이 내보낸 다섯 선은 어느 것도 page를 울리지 않지만 \textit{함께} 움직여 input이 변했음을 말한다. 3주 동안 0.74에서 0.61로 내려간 \textbf{draft acceptance rate}는 제11장의 열린 질문, 즉 \textit{speculative decoding의 조용한 실패를 어느 dashboard 선이 감지하는가?}에 대한 답이다. Acceptance는 작은 model이 실제 traffic을 얼마나 잘 예측하는지 실시간으로 추정하므로 무료로 상시 작동하는 분포 온도계가 된다. Input이 두 model 모두에게 편안한 영역에서 벗어나면 draft가 먼저 틀리고, speedup이 1×로 줄며, 선이 내려간다. 품질 비용보다 TPOT 비용이 먼저 오는 조기 경보 순서를 공짜로 얻는다. \textbf{Refusal rate}는 0.8\%에서 1.9\%로 서서히 오른다. Refusal 자체가 제품이지만, 그 비중이 \textit{늘면} 예전에 서빙하던 사람을 거부하거나 사람들의 구성이 달라진 것이다. \textbf{평균 answer length}는 제13장의 211-token baseline에서 새 중심으로 이동한다. \textbf{Format validity}는 평평하다. Input drift가 보이는 동안 capability drift를 배제하는 유용한 \textit{음성} sentinel이다. 마지막으로 \textbf{prefill token p50}이 612에서 918로 올라 다섯 선 가운데 가장 크게 외친다. §11의 결과도 즉시 나타난다. 긴 prompt는 admission마다 KV를 더 쓰고 preemption을 늘린다. 표 14-2의 PreemptionSurge ticket은 drift가 capacity plane을 통해 자신을 알리는 경우가 많다. 순전히 산술적인 이유로 TTFT가 SLO 선을 향한다. Sentinel은 저렴한 상관관계이지 증명이 아니다. 여러 선이 함께 움직이면 실제 분포 검사로 넘어간다.

업계에서 널리 쓰는 도구는 \textbf{Population Stability Index}다. 냅킨 위에서도 계산할 수 있어 vendor dashboard가 보여 줄 때 직접 \textit{감사}할 수 있다는 장점이 있다. \textbf{Reference window}, 즉 마지막 gate 실행 전후 한 달로 feature를 bin으로 나눈다. 이 기간은 eval set이 믿는 세상이다. 이번 주의 bin별 비중을 측정한 뒤 PSI = Σ (qᵢ − pᵢ) · ln(qᵢ/pᵢ)를 계산한다. Bin이 늘거나 줄면 각 항은 양수가 되며 상대 변화에 절대 이동량으로 가중치를 준다. 그림 14-6의 표는 prompt length에 적용했다. 2k+ bin이 5\%에서 12\%로 두 배 넘게 늘어 혼자 0.061을 기여하고, 합계는 \textbf{PSI = 0.144}다. 신용 평가 분야가 대출 신청자의 drift를 수십 년 동안 추적하며 물려준 관례적 구간(< 0.1 안정, 0.1–0.25 관찰, > 0.25 행동)에서 관찰에 해당한다. PSI를 정직하게 지키는 규율은 둘이다. 먼저 \textit{reference} quantile을 이용해 같은 질량의 bin을 만든다. 예전 세상에 질량이 있던 모든 곳에서 검사가 힘을 갖게 된다. 그리고 length, language, route, hour-of-day처럼 feature마다 별도로 실행한다. 조용한 합계가 반대 방향으로 drift하는 두 feature를 숨길 수 있기 때문이다. 손으로 bin을 나누기 어려운 의미에는 같은 방법을 \textbf{embedding}에 적용한다. 이번 주 prompt-embedding centroid와 reference centroid의 평균 cosine distance를 밤마다 계산하고, reference vector를 eval set 옆 registry에 저장한다. Model뿐 아니라 \textit{측정자}의 lineage도 남긴다. 통계적 엄밀함이 더 필요하면 Kolmogorov–Smirnov test가 bin 대신 CDF 사이의 supremum을 사용한다. p-value 보정은 더 좋지만 bin별 이야기를 해석하기는 어렵다. 이 과정이 PSI를 가르치는 이유는 "어느 사용자 집단이 document를 붙여 넣기 시작했다" 같은 \textit{bin 속 이야기}가 엔지니어가 실제로 행동할 수 있는 부분이기 때문이다.

\subsection*{대응 사다리와 판정 비용을 내는 사람}

대응을 정하지 않은 detection은 취미에 불과하다. 그래서 실습의 drift job은 모든 report를 그림 14-6 사다리의 어느 칸에서 끝낸다. \textbf{첫째 칸, 관찰}은 PSI가 관찰 구간에 있고 sentinel이 조용할 때다. 주석을 남기고 매주 실행하되 아무것도 하지 않는다. 대부분의 주가 여기에 해당하며 drift 업무는 주로 과잉 반응하지 않는 규율이다. \textbf{둘째 칸, 재평가}는 \textit{이번 주} traffic에서 새 slice를 뽑아 label을 붙이기 어려운 품질을 채점하는 LLM-as-judge까지 포함한 gate rubric을 실행한 뒤, 고정된 eval set의 판정과 비교한다. Concept drift와 오래된 eval을 잡는 칸이다. 경제성은 제12장의 순환 고리를 닫는다. 한 달 약 950만 turn의 0.2\%를 turn당 약 1,600 token으로 판정하면 매달 약 3천만 token의 judge inference가 필요하다. API 가격이라면 실제 돈이 들지만 committed-use 하한으로 31억 token 용량을 샀고 traffic은 20억 token만 사용한다. 따라서 judge는 이미 지불한 \textit{골짜기 안} 용량에서 새벽 3시에 제13장 Lane 3 nightly eval 옆으로 달린다. Samuel Insull이 밤의 골짜기를 ice house로 채웠다면 우리는 품질 관리로 채운다. 한계 비용은 거의 0이고 원장은 \$900에 머문다. \textbf{셋째 칸, serving 재조정}은 긴 prompt처럼 drift가 산술적일 때 쓴다. Lever는 model이 아니라 §11/§12의 chunked-prefill knob, seat 수, context cap, 필요하다면 autoscaler target이다. Retraining이 아니라 가격을 다시 매기는 일이다. \textbf{넷째 칸, 새 candidate}는 둘째 칸이 새 세상에서 실제 품질 저하를 보일 때 쓴다. 먼저 새 eval slice로 세상을 재고, fine-tune 또는 model 재선택, registry, gate, canary, promotion 순으로 제13장 pipeline의 맨 위에 다시 들어간다. 이 문장은 operations data가 development로 돌아가는 과정, 즉 마지막 장의 장치가 닫은 제1장의 네 단계 순환이다. Trigger가 된 drift report는 물러나는 registry 항목에 남겨 다음 사고 검토자가 v9가 은퇴한 이유를 민간전승이 아니라 증거로 읽게 한다.

\begin{terminalbox}
\begin{Verbatim}[breaklines=true,breakanywhere=true,fontsize=\small,formatcom=\color{tbbtermfg},codes=\tbbcodeactive]
$ kubectl logs job/drift-nightly | tail -18
drift-report 2026-XX-XX  ref=2026-XX (gate window for v9-int4)
  inputs:
    prompt_len   PSI 0.144  WATCH   (2k+ bin: 5% -> 12%)
    language     PSI 0.031  ok
    route_mix    PSI 0.008  ok
    embedding    d_cos 0.041 vs ref  ok (threshold 0.10)
  sentinels (21-day trend):
    accept_rate  0.74 -> 0.61   FALLING    refusal  0.8% -> 1.9%  RISING
    ans_len p50  211  -> 196    drifting   fmt_valid 97.1 -> 97.0  ok
    prefill p50  612  -> 918    RISING
  verdict: WATCH -> escalate rung 2 (fresh-slice eval tonight)
  actions:
    - queued: eval-fresh-slice (n=2000, stratified by length bin)
    - note: filed on registry entry chat/v9-int4 (drift/2026-XX-XX)
    - capacity: KV headroom check vs prefill p50 918 -> §11 sizing
  next run: tomorrow 03:10 (rides the valley, after lane-3 evals)
\end{Verbatim}
\end{terminalbox}
\figcaption{터미널 기록 14-5  Nightly drift job의 report. Dashboard에서 받은 느낌이 아니라 날짜와 version이 붙은 artifact다. Feature별 PSI에는 이야기가 있는 bin을 밝히고, sentinel은 추세로 표시하며, 판정은 \textbf{사다리의 칸}이고, 행동은 queue 항목과 registry note가 된다. 경사면에는 page가 아니라 review를 보낸다. 여기에는 아무도 깨우는 것이 없고, 모든 항목을 일부러 매주 낮에 읽는다.}

마지막 구분 하나가 오래된 함정에서 팀을 지킨다. \textbf{Drift monitoring은 anomaly detection이 아니다.} 이 절의 도구는 일정에 따라 \textit{window와 reference}를 비교한다. 느리고 신중하며 사람이 검토한다. 찾는 실패 자체가 느리고 신중하게 진행되고, 아직은 자동으로 고칠 수 없기 때문이다. 자기 ML system을 감시하는 또 다른 ML system이 “anomaly”마다 page를 울린다는 vendor의 주장을 경계하자. False positive의 조사 비용을 방금 보호하는 법을 배운 바로 그 on-call에게 떠넘기면서 병원 병동을 뒷문으로 다시 들인다. 효과적인 역할 분담은 \textbf{절벽에는 page}(§14.4, 초에서 분 단위, 완전 자동화), \textbf{경사면에는 review}(이 절, 주 단위, human-in-the-loop)를 주는 것이다. 주간 drift report는 alert가 아니라 calendar 일정으로 전달한다. 지금까지 배포된 가장 믿을 만한 drift detector는 여전히 매주 월요일 coffee를 들고 터미널 기록 14-5를 읽으며 제7장의 가장 오래된 질문, 즉 \textit{약속은 아직도 올바른 약속인가?}를 묻는 senior engineer다. SLO 자체도 drift하기 때문이다. 경쟁자가 5초 만에 답하던 때는 학생 플릿의 500 ms가 중요했지만 기준선은 움직인다. 어떤 PSI도 이를 계산하지 못한다.

\begin{calloutbox}{tbbpurple}{tbbpurplefill}
{\normalsize\bfseries\color{tbbpurple}생각해 보기}\par\vspace{3pt}
\begin{tbbitemize}
  \item (a) Prompt가 질문에서 붙여 넣은 document로 바뀜, (b) 대통령 선거가 “시사” 답변을 바꿈, (c) 시험 기간에 proof-checking 요청이 두 배로 늘어남을 각각 분류하자. 이 절의 도구 가운데 \textbf{가장 먼저} 움직일 것을 말하고 대응 사다리의 칸을 고르자. 완벽한 canary라면 세 가지 중 무엇을 배포 시점에 잡을 수 있었는가? 함정이 있는 질문이니 이유까지 말하자.
  \item 같은 질량을 가진 5개 bin에서 PSI = 0.144다. 같은 변화를 2개 bin(0–300, 300+)으로 다시 계산해 index가 작아짐을 보이자. 이어 binning의 검출력에 관한 일반 교훈과 이번 주가 아니라 reference quantile이 bin을 정의해야 하는 이유를 설명하자.
  \item Acceptance rate는 “무료” drift 온도계지만 input shift와 draft-model update라는 두 원인을 혼동한다. 둘을 구분하는 한 줄짜리 registry check를 설계하고, 이런 교란이 있어도 sentinel panel이 자리를 차지할 가치가 있는 이유를 설명하자.
  \item Judge는 한계 비용이 거의 0인 idle valley를 이용한다. 동료가 “무료니까” sample을 10배, 즉 turn의 2\%로 늘리자고 제안한다. 공짜 점심이 끝나는 두 지점, 즉 token으로 측정한 valley 깊이와 19:00 어깨에서의 judge 대 serving 경합(제12장의 계단)을 찾고 대략적인 상한을 계산하자.
\end{tbbitemize}
\end{calloutbox}

\section{실습: 관제 센터를 만들고 두 번 깨뜨리기}

강의계획서는 “서빙 대시보드 구축과 알람 설정”을 약속했다. 이 실습은 serving dashboard와 alarm을 모두 만들고, 계측 도구의 주인과 관광객을 가르는 단계를 더한다. 방금 만든 panel을 실시간으로 보며 두 가지 실패를 일부러 일으킨다. 과정에서 계속 사용한 T4 card와 아직 warm 상태인 11주 차 vLLM 위 Qwen2.5-1.5B server, 9–10주 차 classifier 한 쌍에서 모두 실행한다. 절대적인 숫자는 플릿과 다르지만 과정의 오래된 원칙에 따르면 아무 문제가 없다. \textbf{책의 상수와 숫자를 맞추는 것이 아니라 형태와 추론으로 채점한다.} 그림 14-7이 지도이며, 여섯 build 단계와 두 drill, 다섯 결과물이 있다.

\begin{tbbfigure}
  \adjustbox{max width=0.94\linewidth,max totalheight=0.46\textheight}{\includegraphics{figures/fig14-7_lab.png}}\\[3pt]
  \figcaption{그림 14-7  실습 14. 1–3단계에서는 stack을 세우고 네 행을 그리며, 4–5단계에서는 경보 목록을 작동시키고 page를 routing한다. 6단계에서는 drift job의 일정을 잡는다. 이어 drill A가 자기 hardware에서 화요일 밤을 재현해 §12.4의 외상을 갚고, drill B가 이번에는 계측 도구를 갖춘 채 제9장의 retry storm을 재생한다. 결과물은 link로 제출하며, 채점자가 받는 image는 주석을 붙인 drill screenshot 두 장뿐이다.}
\end{tbbfigure}

\textbf{1–3단계: Stack 세우기.} Helm으로 \code{kube-prometheus-stack}을 설치한다. Chart 하나에 Prometheus, Grafana, Alertmanager, node-exporter, kube-state-metrics가 들어 있다. 다른 일을 시작하기 전에 \code{/\allowbreak targets}에서 모든 target이 UP인지 확인한다. Scrape 점호가 실습의 첫 번째 정직한 결과물이며 \code{up == 0} 행이 첫 debugging 문제다. 흔한 원인은 ServiceMonitor가 찾는 port name이 Service에 없거나 selector가 namespace를 놓친 경우다. Engine \code{:8000/\allowbreak metrics}와 classifier를 위한 ServiceMonitor 두 개, GPU node의 DCGM DaemonSet, §14.4의 세 probe를 실행하는 blackbox exporter를 연결한다. 이어 \textbf{네 행 dashboard를 repository의 JSON으로} 만든다. 표 14-1의 panel과 약속을 측정하는 모든 panel의 SLO 선, 4행의 UUID별 pod 표, "높게 고정되는 것이 건강함", "busy ≠ yours"처럼 가르치는 caption을 넣는다. 11주 차 token-aware load test를 concurrency 세 단계로 실행하면서 행이 \textit{인과 순서대로} 켜지는지 확인한다. 먼저 3행의 queue, 다음 1행의 p99, knee를 넘으면 마지막으로 2행이 누군가 한 명이 아니라 모두의 문제임을 보여야 한다. 순서가 다르면 query가 틀렸다. 이제 dashboard가 스스로를 검사한다.

\textbf{4–6단계: 경보 작동시키기.} SLI recording rule과 표 14-2의 경보 목록을 code로 만든다. 사용자 한 명인 실습 환경에서는 플릿의 8이 아니라 터미널 기록 11-5에서 측정한 seat 수에 맞춰 queue threshold를 조정한다. CI에서 \code{promtool check rules}가 통과하는지 확인한다. Rule도 release artifact이므로 pipeline은 weights처럼 gate한다. Page는 전화와 가까운 receiver로 보낸다. 자기 device를 울리는 webhook이면 충분하다. Ticket은 channel로 보낸다. \code{amtool}로 synthetic alert를 발생시켜 grouping되고 silence가 유지되는지 확인해 연결을 입증한다. Dead man’s receiver는 \textit{cluster 바깥}에 세운다. Heartbeat service의 무료 tier나 lab VM의 cron이면 충분하다. 핵심은 운명을 공유하지 않는 데 있다. 10주 차 gate 실행으로 reference window를 고정하고 03:10에 drift CronJob을 예약한다. WATCH 판정을 포함한 첫 실제 report를 결과물과 함께 제출한다. 이제 채점하는 부분, 즉 \textbf{두 번 깨뜨리기}로 넘어간다.

\begin{terminalbox}
\begin{Verbatim}[breaklines=true,breakanywhere=true,fontsize=\small,formatcom=\color{tbbtermfg},codes=\tbbcodeactive]
# drill A -- stage tuesday night against your own server (lab GPU, qwen-1.5b)
$ kubectl patch cm nvidia-device-plugin -p '{"data":{"config":"replicas: 2"}}'
$ locust -f loadtest.py --headless -u 24 &          # steady load, seats full
$ kubectl create job eval-burn --image=... -- python soak.py --busy
# watch, in order (screenshot each transition):
#   row 3: decode step p50 0.021 -> 0.044   <- the neighbor's tax
#   row 3: num_requests_waiting 0 -> 11     <- the queue is local
#   row 1: TTFT p99 0.28 -> 2.9s; SLI ratio dives
#   21:0X  PAGE QueueSustained (your phone buzzes; note the clock)
#   row 4: pods-per-UUID: your-gpu: 2 (qwen-srv, eval-burn)
$ kubectl delete job eval-burn && date              # note the clock again
#   queue drains; page auto-resolves. MTTD = page - patch. MTTR = delete - patch.

# drill B -- week 9's retry storm, now with instruments (classifier pair)
$ locust -f storm.py --headless -u 60 &   # client retries ON (the bad config)
$ kubectl delete pod classifier-1         # kill one of two under load
#   rows: req rate x3 (retries multiplying), goodput /2, queue sawtooth --
#   the waveform from week 9, on YOUR panels. now apply ch. 9's fixes
#   (retry connect-only, one layer, jitter, shed early) and re-run:
#   same kill, flat goodput, no storm. file both panels side by side.
\end{Verbatim}
\end{terminalbox}
\figcaption{터미널 기록 14-6  두 가지 drill. Drill A는 자기 hardware에서 time-slice ConfigMap과 바쁜 이웃, 그림 14-4의 정확한 panel 순서를 재현해 §12.4의 외상을 갚는다. 자기 전화가 pager이고 벽시계가 채점자다. Drill B는 제9장의 storm을 재생해 waveform과 그것을 그리는 panel을 시각적 기억 속에 나란히 남긴다. 그 장에서 새벽 2시에 중요해질 것이라고 약속한 “Grafana screenshot 두 장”을 마침내 찍는다.}

\textbf{제출할 것}은 주석을 붙인 drill panel 두 장을 제외하면 screenshot이 아니라 link다. PR history가 있는 repository의 dashboard JSON과 rule YAML, \code{/\allowbreak targets} 점호, 그림 14-1의 12분과 8분에 견준 \textit{자기 MTTD와 MTTR}을 적은 drill A timeline, 한 문단 해설을 곁들인 drill B의 storm 전후 panel, drift job의 첫 report를 제출한다. 학기 전체를 묶는 결과물도 있다. 제13장 runbook을 사용해 \textbf{drill A의 incident note를 자기 registry 항목에 연결한다.} 이제 WATCH 줄에는 자기 dashboard 행의 이름이 들어간다. 한 페이지짜리 비난 없는 postmortem도 작성하며 action item에는 \textit{monitoring 자체의 변경}을 하나 이상 넣어야 한다. 새 alert, 수정한 threshold, 다음 새벽 2시의 두뇌에서 30초를 아껴 줄 caption이 예다. 마지막 요구 사항이 실습의 주장을 축소해 담고 있다. Monitoring stack은 정원과 같은 방식으로 완성된다. 일부러 절대 끝내지 않고 사고를 하나씩 겪으며 가꾼다.

\begin{calloutbox}{tbbamber}{tbbxFFFBEB}
{\normalsize\bfseries\color{tbbx92400E}상자 14-2  채점 기준 — 실습 14(감점되는 위험 신호 포함)}\par\vspace{3pt}
\begin{tbbitemize}
  \item \textbf{Stack + scrape(20\%).} 증거와 함께 모든 target이 UP이다. ServiceMonitor와 exporter는 손으로 적용하지 않고 repository에 둔다. 위험 신호: target이 0/N인데 그 metric을 query하는 dashboard, 즉 확신에 찬 빈 panel.
  \item \textbf{Dashboard 설계(25\%).} Runbook 순서로 네 질문에 답하는 네 행, 그려진 SLO 선, 모든 곳의 올바른 percentile 집계가 필요하다. Board 어디든 avg-of-p99s가 하나 있으면 이 항목은 절반이 상한이다. Caption은 사용법을 가르쳐야 한다.
  \item \textbf{경보 목록(20\%).} Recording rule을 읽는 burn 짝, 지속 조건을 둔 queue page, page가 아니라 ticket으로 보내는 원인, 처음부터 끝까지 검증한 deadman이 필요하다. 위험 신호: 정상적인 실습 traffic에서 매주 한 번보다 자주 page를 울릴 목록. 횟수를 직접 세고 밝힌다.
  \item \textbf{Drill(25\%).} 둘 다 직접 재현하고, 지켜보다가 알아차리는 대신 \textbf{자기} alert가 잡아야 한다. MTTD/MTTR을 측정해 정직하게 비교한다. Drill B는 실패가 일어나는 모습뿐 아니라 수정이 작동하는 모습도 보여야 한다. 위험 신호: 시계가 없는 screenshot.
  \item \textbf{순환 고리(10\%).} Registry 항목에 incident note를 연결하고 monitoring을 바꾸는 action item을 담은 postmortem을 PR로 배포한다. 실제 플릿을 잘 운영할 사람을 가장 잘 예측하는 기준이다.
\end{tbbitemize}
\end{calloutbox}

\begin{calloutbox}{tbbpurple}{tbbpurplefill}
{\normalsize\bfseries\color{tbbpurple}생각해 보기}\par\vspace{3pt}
\begin{tbbitemize}
  \item Time-sliced T4에서 drill A는 decode step을 약 2.1×로 만들지만 CPU의 classifier drill은 p50이 거의 움직이지 않는 동안 p99가 세 배가 된다(터미널 기록 12-4의 형태). 두 현상을 기본 원리로 설명하자. 각 workload는 어떤 resource를 포화시키며 time-slicing은 왜 median보다 tail에 먼저 세금을 부과하는가?
  \item Queue page threshold는 터미널 기록 11-5의 seat census에서 나왔다. 12주 차 autoscaler는 이제 매시간 레플리카 수를 바꾼다. Autoscaling에도 살아남는 threshold 식을 제안하자. 힌트: pod별 gauge와 플릿 합계 중 표 14-2는 무엇을 골랐으며 그 이유는 무엇인가? 반대로 선택할 때의 실패 형태도 말하자.
  \item Drill B의 나쁜 실행에서 처음 60초 동안 1행 SLI는 거의 움직이지 않았지만 3행은 비명을 질렀다. §14.4의 “증상은 page로, 원인은 ticket으로”와 조화시키자. Retry storm에서 queue sawtooth는 증상인가 원인인가? 답은 QueueSustained의 지속 window에 무엇을 요구하는가?
  \item Postmortem은 monitoring 변경을 요구한다. Drill A에서 나온 GPU UUID별 pod ticket, baseline 대비 decode-step 비율 panel, pool별 time-slice config expiry를 확인하는 CI 가운데 프로젝트 플릿에 가장 잘 일반화되는 변경 하나를 고르고 논증하자. 각 선택은 어떤 미래 사고 부류를 막지 못하는가?
\end{tbbitemize}
\end{calloutbox}

\namedsection{요약}{열여섯 가지로 정리한 이 장}

\qitem{1}{\textbf{두 번의 밤이 하나의 주장을 입증한다. 오후 2시에 그래프로 그린 것만 새벽 2시에 볼 수 있다.} 첫날 밤의 다섯 가지 확인은 모두 잘못된 질문에 정확히 답했다. 프로브는 살아 있는지 재고, 로그는 자백을 기록하며, 평균은 누구도 설명하지 못하고, 파드 하나는 일화일 뿐이며, GPU 사용률에는 누구의 작업이든 포함된다. 나쁜 사용자 시간은 85분이었다. 둘째 날 밤에는 같은 사람이 미리 준비한 질문으로 20분 만에 해결했다.}

\qitem{2}{\textbf{Apollo는 이 강의 계획서의 축소판이다.} 텔레메트리는 미리 설계하지 않으면 존재하지 않는다. 로켓이나 플릿에 ssh로 접속할 수는 없다. 원시 신호 자체가 아니라 \textbf{알아본 특징}이 사람을 구한다. SCE-to-AUX는 훈련에서 익힌 기억이었다. 착륙을 구한 요약표도 실패한 시뮬레이션 뒤에 작성되었다. 계측하고, 연습하고, 주석을 남겨라.}

\qitem{3}{\textbf{프로브는 반사 작용이고 모니터링은 기억이다.} 5장의 문장이 결실을 맺는다. 대시보드는 기억, 경보는 주의, 드리프트 모니터는 예견이다. 관측 가능성은 미리 정의하지 않은 질문에도 답할 수 있게 하는 속성이다. 화요일의 충돌을 위한 전용 패널은 없었지만 범용적인 네 행이 답을 찾아냈다.}

\qitem{4}{\textbf{세 기둥은 비용이 생기는 방식이 다르다.} 지표는 시계열마다 비용을 내며 저렴하고 항상 켜져 있다. 경보는 여기에 둔다. 로그는 이벤트마다 비용을 내는 파드의 전기이며, 6장의 대가를 치른 덕분에 파드가 죽기 \textit{전에 파드 밖으로} 보낸다. 트레이스는 요청마다 비용을 내고 표본을 추출하며, 홉이 세 개인 이 강좌에서는 최소한으로 쓴다. 5장에서 예고한 수행원 가운데 익스포터와 로그 전송기는 설치하지만, 관심사마다 소유자를 하나만 두기 위해 메시는 보류한다.}

\qitem{5}{\textbf{지표는 이름 + 레이블 + 값 + 시각이며, 비용은 시계열 단위로 낸다.} 레이블은 파드, UUID, 경로처럼 한정된 목록에서 “어느 것인가?”에 답한다. 사용자나 요청처럼 값이 무한히 늘어나는 정보는 로그에 둔다. 카디널리티(cardinality)는 모니터링의 메모리 장벽이다. 부주의한 레이블 하나가 시계열 14,000개를 천 배로 늘린다.}

\qitem{6}{\textbf{형태는 셋이지만 명령은 하나다.} 카운터는 오르기만 하므로 rate()로 변화율을 구하며 재시작에도 안전하다. 게이지는 지금 값을 읽으며 포화를 표현한다. 히스토그램은 \textbf{백분위수의 평균을 낼 수 없기 때문에} 존재한다. 버킷 카운터는 레플리카에 걸쳐 더할 수 있고, 합친 버킷의 백분위수만 올바른 플릿 p99다. p99의 평균은 2.76 s라고 했지만 실제 값은 9.08 s였다.}

\qitem{7}{\textbf{부재도 사실이어야 하므로 풀 방식을 쓴다.} 스크레이프마다 up\{pod\}를 기록하고, 죽은 대상은 침묵 대신 묘비를 남긴다. 계보는 Borgmon에서 Prometheus로 이어졌다. Prometheus는 2012년 SoundCloud에서 시작해 Kubernetes에 이어 CNCF의 두 번째 프로젝트가 되었다. Grafana는 2014년에 Kibana를 포크해 출발했다. 기록 규칙은 SLI를 미리 계산하며 13장의 파이프라인으로 배포한다. 규칙도 릴리스 산출물이다.}

\qitem{8}{\textbf{익스포터 모음은 강좌 전체를 다시 관통한다.} node-exporter는 2장의 머신을, cAdvisor는 바깥에서 cgroup 파일을 읽어 3장의 정직한 마운트를 구현한다. 이제 memory.events의 oom\_kill과 memory.high 경보도 그래프로 볼 수 있다. kube-state는 재시작, OOMKilled, 엔드포인트, 롤아웃 상태 등 5장의 경보 목록을 제공한다. DCGM은 nvidia-smi를 산업 규모로 확장해 UUID별 값을 파드와 조인한다. 엔진의 /metrics에는 11장의 선들이 있고, 블랙박스는 6장의 출입구를 바깥에서 조사한다.}

\qitem{9}{\textbf{GPU “사용률”은 점유 듀티 사이클일 뿐이다. 바쁨 ≠ 내 것 ≠ 효율적임이다.} 메모리 병목 디코드는 ALU의 10분의 1만 쓰면서도 90\%+로 표시될 수 있고, 테넌트 둘이 함께 쓰면 둘 다 \textasciitilde{}100\%로 보인다. 그 값만으로는 누구의 작업인지 알 수 없다. 대기열, 선점, DRAM 활동 같은 포화 신호가 먼저 나타나고, 사용률은 뒤처지며 오해를 부른다. 자원에는 USE, 서비스에는 RED를 적용하고, 대기열을 핵심 포화 게이지로 삼는다.}

\qitem{10}{\textbf{네 행 대시보드는 약속 → 누구 → 원인 → 기반의 순서다.} 1행은 오류 예산, 소진율, 올바른 p99, 스트림별 속도로 SLO 문장을 패널에 옮기며 호출을 보내는 유일한 행이다. 2행은 레플리카별 RED로 모두가 느린지 일부만 느린지 가른다. 3행은 대기열, 높게 고정되어 정상인 KV, 디코드 속도, 선점으로 엔진 내부의 원인을 보여 주며 9장의 대기열/계산 분리를 지킨다. 4행은 UUID별 DCGM, GPU당 파드, 오토스케일러 태세로 기반을 보여 준다.}

\qitem{11}{\textbf{증상에는 호출을, 원인에는 티켓을 보낸다.} 증상은 적고 실제이며 완전하다. 처음 보는 실패도 SLI에 드러난다는 사실을 충돌 사고가 입증했다. 원인은 셀 수 없이 많고 값싸며 해가 없는 경우도 흔하다. 병원에서는 경보의 85–99\%에 조치가 필요 없었고, 경보 피로가 사람을 죽인다는 대가를 배웠다. 모든 호출은 조치할 수 있고, 긴급하며, 실제여야 한다.}

\qitem{12}{\textbf{SLO 문장은 소진율로 컴파일된다.} p99 ≤ 500 ms ≡ 99\% of turns good → SLI ratio → budget = 1\% × \textasciitilde{}9.5M turns ≈ 95k bad turns/30 d라는 식으로 이어진다. Burn = bad-fraction ÷ 0.01이다. 다중 구간 정책은 14.4×/1 h∧5 m이면 호출, 6×/6 h이면 호출, 3×/1×이면 티켓을 만든다. 긴 구간은 의심하고 짧은 구간은 용서한다. 대기열 같은 선행 증상에는 빠른 경로를 허용한다.}

\qitem{13}{\textbf{거부할 수 있는 모든 계층의 바깥에 서라.} 출입구, 응답, 스트림을 검사하는 세 합성 프로브(synthetic probe)는 각각 자신이 지목할 계층 바깥에 있어야 하고, 실패의 존재가 아니라 지속성에 호출을 보낸다. 사례 14-A의 초록색 /health는 첫 번째 프로브 하나뿐이었다. Watchdog는 항상 발화하므로 그 침묵을 경보로 삼는다. 클러스터와 운명을 공유하지 않는 수신기가 침묵을 감지해 호출을 보낸다. 데드맨 스위치(dead man's switch)가 감시자를 감시한다.}

\qitem{14}{\textbf{사람이 참여하는 순환은 모니터링의 가치를 복리로 키운다.} 호출에는 WATCH 단계에서 대시보드 행을 지목하는 런북이 붙는다. 사고는 사람 이름이 아니라 메커니즘을 다루는 비난 없는 사후 분석으로 끝나며, 조치 항목은 반드시 \textbf{모니터링을 바꾼다}. 화재가 날 때마다 연기 감지기가 하나씩 생긴다. 메모는 레지스트리 항목에 등록되어 마침내 10장의 검토자가 읽는다. MTTD와 MTTR은 이 순환을 채점하고 DORA의 네 번째 핵심 지표를 완성한다. 스택 비용은 \textasciitilde{}\$23/mo이며 원장은 정확히 \$900에 마감된다.}

\qitem{15}{\textbf{드리프트는 호출을 보내지 않는 실패다. 카나리는 절벽을 잡고 드리프트 모니터는 경사를 잡는다.} 대응 방식에 따라 데이터 드리프트(data drift), 개념 드리프트(concept drift), 사전 확률 변화, 평가 세트 노후화로 나눈다. 데이터 드리프트에서는 입력이 움직여 서빙 계산과 평가 범위가 먼저 깨진다. 개념 드리프트에서는 정답이 움직여 최신 현실을 반영한 평가만 이를 발견한다. 사전 확률 변화에서는 작업 구성이 달라져 가드레일 기준선을 다시 잡아야 한다. 평가 세트도 낡으므로 분기마다 다시 표본을 뽑는다. 수락률, 거부율, 길이, 프리필 p50 같은 감시 지표(sentinel metric)는 추가 비용 없이 얻는다. PSI는 손으로 계산할 수 있으며 0.144는 관찰 구간이다. 경사에는 호출이 아니라 \textbf{검토}를 배정한다.}

\qitem{16}{\textbf{대응 사다리는 반응 비용을 정한다.} 관찰 → 새로운 표본으로 재평가 → 서빙 다시 조정 → 13장의 레지스트리→게이트→카나리를 거치는 새 후보 순으로 올라간다. 판정기는 이미 약정한 비용 범위에서 유휴 골짜기를 활용하므로 Insull의 얼음 창고처럼 한계비용이 0달러다. 운영 데이터가 다시 개발로 흘러가며 1장의 순환이 닫힌다. 어떤 PSI도 계산하지 못하는 월요일의 질문도 남는다. 지금도 이 약속이 올바른 약속인가?}

\namedsection{연습문제}{}

\subsection*{A. 객관식}

\qitem{1}{첫날 밤의 다섯 가지 확인(프로브, 로그, 평균, 파드 하나의 p99, nvidia-smi)이 실제 장애 중에도 모두 “통과”한 이유는?}

\begin{tbboptions}
  \item ① 도구 설정이 잘못되었다.        ② 각 도구는 생존, 자백, 중심 경향, 일화, 점유처럼 SLO와 다른 질문에 답하며 약속 자체를 계산한 도구는 하나도 없었다.
  \item ③ 사고가 너무 짧아 관측할 수 없었다.        ④ 플릿에 레플리카가 충분하지 않았다.
\end{tbboptions}

\qitem{2}{레플리카 네 개의 p99를 평균 내면 2.76 s지만 실제 플릿 p99는 9.08 s다. 근본적인 오류는?}

\begin{tbboptions}
  \item ① 구간이 너무 짧었다.        ② 백분위수는 선형이 아니므로 합친 분포의 분위수만 의미가 있으며, 레플리카에 걸쳐 더할 수 있는 것은 버킷 카운터뿐이다.
  \item ③ rate()를 빠뜨렸다.        ④ 버킷 경계가 SLO와 맞지 않았다.
\end{tbboptions}

\qitem{3}{DCGM\_FI\_DEV\_GPU\_UTIL 값이 93\%다. 이 값으로 실제로 알 수 있는 사실은?}

\begin{tbboptions}
  \item ① GPU가 포화 상태의 93\%에 도달했다.        ② 내 파드가 FLOP의 93\%를 사용한다.
  \item ③ 최근 시간의 93\% 동안 누구의 것이든 커널 하나 이상이 올라와 있었다. 점유 듀티 사이클일 뿐 여유 용량, 효율, 소유권은 알려 주지 않는다.        ④ 이웃이 쓸 수 있는 용량이 7\% 남았다.
\end{tbboptions}

\qitem{4}{Garman이 말한 “너무 자주 반복되지 않으면 진행한다”는 판단은 이 장의 체계에서 무엇에 해당하는가?}

\begin{tbboptions}
  \item ① 침묵        ② 심각도 사다리        ③ 데드맨 스위치        ④ 발생률과 지속성에 따른 경보. for:, 지속 구간, 짝을 이룬 소진 구간을 이용해 존재 자체가 아니라 반복에 호출을 보낸다.
\end{tbboptions}

\qitem{5}{99\% / 30-day SLO에서 소진율이 24×라는 뜻은?}

\begin{tbboptions}
  \item ① 요청의 24\%가 SLO 임곗값을 위반하며, 이 속도면 월간 예산이 \textasciitilde{}30 h (30 d ÷ 24) 만에 소진된다.        ② SLO가 24× 지나치게 엄격하다.
  \item ③ p99가 임곗값의 24×다.        ④ 호출기가 24번 울린다.
\end{tbboptions}

\qitem{6}{소진율 호출이 완전성을 보장하는데도 QueueSustained 호출을 따로 두는 이유는?}

\begin{tbboptions}
  \item ① 소진율 규칙은 신뢰할 수 없다.        ② 대기열은 원인이며 원인에도 호출을 보내야 한다.
  \item ③ 소진율이 24×여도 1시간 구간이 14.4×를 넘는 데 약 35분이 걸린다. 대기열은 선행 증상이므로 증상에는 빠른 경로를 허용한다.        ④ Alertmanager는 두 호출을 라우팅할 수 없다.
\end{tbboptions}

\qitem{7}{Watchdog 경보는 vector(1)을 평가해 계속 발화한다. 그 이유는?}

\begin{tbboptions}
  \item ① 호출기를 예열한다.        ② 외부 수신기에서 이 경보가 사라지는 것 자체가 경보다. 심장박동이 멈추면 모니터링 스택이 죽었다는 뜻이며, 이것이 데드맨 스위치다.
  \item ③ Prometheus에는 활성 경보가 하나 필요하다.        ④ 경보 지연 시간을 측정한다.
\end{tbboptions}

\qitem{8}{야간 작업이 프롬프트 길이 PSI = 0.144를 보고했고 여러 감시 지표도 같은 방향으로 움직인다. 올바른 대응은?}

\begin{tbboptions}
  \item ① 당직자를 호출한다.        ② 즉시 다시 학습한다.        ③ 관찰 구간(0.1\textasciitilde{}0.25)이므로 주간 검토를 하고 새로운 표본 평가인 사다리 2단계로 올린다. 경사에는 호출이 아니라 검토를 배정한다.        ④ SLO 임곗값을 높인다.
\end{tbboptions}

\subsection*{B. 단답형}

\qitem{1}{플릿의 TTFT p99를 구하는 PromQL을 작성하고 세 계층(rate, sum by le, histogram\_quantile)이 각각 기여하는 바를 설명하라. 가장 안쪽 계층이 재시작에도 안전한 이유도 포함하라.}

\qitem{2}{KV 캐시 사용률은 96\%에 고정되어도 경보가 울리지 않지만 num\_requests\_waiting이 6분 동안 9를 유지하면 사람을 깨운다. 11장의 아키텍처를 근거로 두 설계를 정당화하라. 각 게이지가 무엇을 측정하며, 어느 것이 수요가 좌석을 초과했음을 뜻하는가?}

\qitem{3}{세 합성 프로브의 경로, 실행 빈도, 경보 조건을 명시하라. 각 프로브가 거울이 아니라 목격자가 되려면 어느 계층 바깥에 있어야 하는지도 설명하고, 가장 먼저 발견하는 6장의 실패 유형(503 / 504 / 스트림 끊김)에 대응시켜라.}

\qitem{4}{7장의 99.9\%는 43분을 허용했고 이 장의 99\%는 95,000턴을 허용한다. 요청 기반 단위가 시간 기반 단위와 달리 하루 주기 물결의 비용을 올바르게 매기는 이유를 설명하라. 새벽 3시와 21시 30분의 나쁜 1분은 두 회계 방식에서 각각 얼마의 비용이 되는가?}

\qitem{5}{각 항목을 호출 또는 티켓으로 분류하고 한 절씩 근거를 제시하라. (a) 준비된 엔드포인트 = 0, (b) 프로덕션 UUID당 파드 = 2, (c) probe\_success가 3× 연속 실패, (d) 레플리카 하나에서 OOMKilled 한 번 관측, (e) 6 h 동안 소진율 6×, (f) 10 min 동안 선점률 0.8/s.}

\qitem{6}{수락률 감시 지표를 설명하라. 초안 수락률 하락이 입력 분포의 온도계인 이유, 어떤 비용이 다른 어떤 비용보다 먼저 나타나는지, 신호를 가장할 수 있는 교란 요인, 이를 구분하는 레지스트리 질의를 모두 다루어라.}

\subsection*{C. 논술형}

\qitem{1}{화요일 사고의 비난 없는 사후 분석을 작성하라. 두 경보, 대시보드 조사, 퇴출을 분 단위 연표로 정리하고, 첫날 밤의 반사실(20분 대 85분)을 이용해 나쁜 사용자 시간으로 영향을 계산하라. 사람 이름이 아니라 메커니즘, 즉 오래된 시분할 ConfigMap과 빠진 UUID당 파드 정책으로 근본 원인을 서술하라. 담당자가 있는 조치 항목 세 가지를 제시하되 적어도 하나는 모니터링 자체를 바꾸어야 한다. 대부분의 사후 분석이 빠뜨리는 문단, 즉 이 사고가 새로 만들 근거가 \textit{되지 못하는} 것은 무엇이며 왜 그런지도 덧붙여라.}

\qitem{2}{정상 상태에서 월 2회라는 명시적 경보 피로 예산 아래 프로젝트 플릿의 전체 호출 정책을 설계하라. 직접 측정한 트래픽, 좌석 수, 도착률을 근거로 심각도와 구간을 정한 경보 목록, 데드맨의 외부 수신기, 그룹화와 억제 규칙을 포함하라. 감사자가 가장 먼저 읽을 부분도 작성하라. 어떤 실패 유형을 의도적으로 티켓으로 낮췄으며, 각 결정으로 어느 정도의 사용자 고통을 받아들이는가?}

\qitem{3}{“모니터링은 카나리의 분석 구간을 영구화한 것이다.” 13장의 금요일 버그에서 답변 길이를 따라 카나리 가드레일, 기록 규칙, 드리프트 감시 지표까지 추적하라. 각 시간 척도에 대해 구간, 임계값 방식(고정값, 기준선 대비, 기준 분포 대비), 독자, 대응을 제시하라. 같은 숫자가 세 계측 도구에 모두 쓰이는 이유와, 어떤 지표만 존재할 가치가 있는지에 관한 함의도 논증하라.}

\qitem{4}{관측 비용에 관한 논술이다. 23달러 항목, TSDB 대략 계산, 판정기의 유휴 골짜기 활용, 가장 희소한 자원인 당직자의 주의를 포함해 모니터링 스택의 비용을 정직하게 계산하라. 이 규모의 플릿에서 관측 가능성의 한계가치와 한계비용이 만나는 지점을 찾는다. 그런 다음 오류 예산 패널, 사전 합의한 완화 조항, 런북을 이용해 7장의 새벽 2시 논술 문제를 해결하고 다음 주장에 찬성하거나 반박하라. “지루한 운영이 이 강좌 전체의 올바른 최종 상태이며, 23달러는 청구서에서 가장 저렴한 항목이었다.”}

\namedsection{정답과 해설}{}

\subsection*{A. 객관식}

\qitem{1}{\textbf{②} 정직한 계측 도구는 다섯 개였지만 이를 조합한 곳은 없었다. SLO가 정의한 “정상”, 즉 SLI를 어디에서도 계산하거나 분해하지 않았다. 1행이 약속을 확인하는 다섯 도구를 대신하고 2\textasciitilde{}4행이 원인 탐색을 대신한다. (§14.1)}

\qitem{2}{\textbf{②} 꼬리 값의 평균은 허구다. 9.8 s인 레플리카 하나와 0.4 s인 레플리카 세 개를 평균 내면 아무도 실제로 겪지 않은 지연 시간이 나온다. 합친 버킷 자체가 플릿 분포다. 이는 7장의 네 번째 문법 규칙을 도구로 강제한 것이다. (§14.2)}

\qitem{3}{\textbf{③} 누구의 것이든 커널 하나가 올라와 있던 시간의 듀티 사이클이다. 메모리 병목 디코드 루프는 ALU를 굶기면서도 이 값을 높게 고정할 수 있고, 이웃의 커널도 계산에 포함된다. 여유 용량은 포화와 PROF\_* 지표에서, 소유권은 UUID당 파드 조인에서 찾는다. (§14.2)}

\qitem{4}{\textbf{④} 1202 경보는 간헐적이었고 유도 장치는 계속 작동했다. 하강 전에 문서로 정한 발생률 조건과 증상 확인을 함께 적용한 것이다. for:, 지속 평균, 짝을 이룬 구간은 이 판단을 자동화한 형태다. (§14.4)}

\qitem{5}{\textbf{①} Burn = bad-fraction ÷ (1 − SLO) = 0.24 ÷ 0.01이고, time-to-empty = 30 d ÷ burn이다. 핵심 관점은 소진율이 예산의 \textbf{지출 속도}라는 것이다. 그래서 여러 속도의 정책을 일관되게 구성할 수 있다. (§14.4)}

\qitem{6}{\textbf{③} 긴 구간이 의심하는 동안 지연이 생긴다. 소진율 24×에서도 35 min이 든다. 용량 문제로 생긴 고통은 7장의 무릎점과 11장의 허용 선처럼 대기열에서 먼저 드러나므로, 증상 빠른 경로는 \textasciitilde{}7분 만에 호출한다. 둘 다 증상이며, 원칙주의적 반론은 생각해 보기 상자에서 다룬다. (§14.4)}

\qitem{7}{\textbf{②} 침묵을 경보로 삼으면 실패 방식이 뒤집힌다. 죽은 Prometheus는 자신의 죽음을 호출할 수 없으므로, Prometheus가 죽일 수 \textbf{없는} 대상인 외부 수신기가 호출해야 한다. 리전, 결제, 웹훅(webhook) 체인을 별도로 감사해 클러스터와 운명을 공유하지 않게 한다. (§14.4)}

\qitem{8}{\textbf{③} 드리프트는 경사다. 경사에 호출을 보내면 경보 피로가 다시 생긴다. 관찰 판정은 새로운 표본을 유휴 골짜기에서 평가하는 2단계를 예약하고 보고서를 레지스트리 항목에 등록하며, 정기 검토 주기에 사람이 참여하게 한다. (§14.5)}

\subsection*{B. 단답형}

\qitem{1}{histogram\_quantile(0.99, sum by (le) (rate(vllm:time\_to\_first\_token\_seconds\_bucket[5m]))). rate()는 누적 버킷 카운터를 초당 흐름으로 바꾸고 카운터 초기화를 이어 계산한다. 파드가 다시 태어나면 초기화로 나타나지만 이를 감지하고 보정한다. 그렇지 않으면 5장의 조정 과정이 모든 그래프를 오염시킨다. sum by (le)는 버킷 경계를 보존하면서 레플리카를 합치는 유일하게 올바른 교차 레플리카 집계다. histogram\_quantile은 해당 버킷 안에서 분위수를 보간하므로 오차는 버킷 너비 정도다. (§14.2)}

\qitem{2}{KV 풀은 캐시 형태의 \textbf{할당}이다. PagedAttention의 설계 목표는 풀을 가득 채우는 것이다. 좌석이 가중치 스트림 비용을 분산하므로 높게 고정된 상태는 머신이 제값을 하는 상태다. 3장의 페이지 캐시 교훈을 GPU에 적용한 것이다. 대기열은 \textbf{허용}을 측정한다. 대기가 지속되면 도착률이 좌석 처리량을 넘어 용량을 초과했다는 뜻이고, 사용자는 이미 TTFT로 비용을 치르고 있다. 한 게이지가 가득 차면 작동 중이고, 다른 게이지가 0보다 큰 상태로 지속되면 실패 중이다. (§14.2, 11장)}

\qitem{3}{출입구: 로드 밸런서 밖에서 30 s마다 HTTPS GET /v1/models를 보내고 3× 연속 실패하면 호출한다. 플릿이 볼 수 없는 경계 밖에 있으므로 TLS, 대상 상태, 라우팅과 같은 503 유형을 잡는다. 응답: 게이트웨이 밖에서 60 s마다 토큰 1개 완성을 요청하고 시간 제한은 10 s로 둔다. 출입구를 지난 뒤 허용이나 엔진에서 생기는 504/멈춤을 잡는다. 스트림: 5 min마다 짧은 스트리밍 완성을 요청해 첫 토큰이 ≤ 2 s이고 마지막 청크가 온전한지 확인한다. 유휴 시간 제한 계층 바깥에서 표 6-3 셋째 행의 스트림 끊김을 잡는다. 공통 규칙은 검사하는 계층 안의 프로브는 목격자가 아니라 거울이라는 것이다. (§14.4, 6장)}

\qitem{4}{시간 기반 예산은 모든 분에 같은 가격을 매기지만 계단형 수요는 그렇지 않다. 12장의 물결에서 새벽 3시의 1분에는 21시 30분의 약 3분의 1에 해당하는 턴만 들어온다. 요청 기반 예산은 사용자 단위로 지출한다. 화요일 피크 시간 20분의 비용은 \textasciitilde{}24\% × 23/s × 1,200 s ≈ 6.6k bad turns (7\% of the month)다. 시간 기반 회계는 고통받는 사람이 3분의 1뿐인 새벽 3시에도 같은 비용을 기록한다. 호출의 긴급도는 고통받는 사람 수를 따라야 하므로 턴을 사용한다. (§14.4)}

\qitem{5}{(a) \textbf{호출} — 라우팅할 곳이 없어 모든 사용자가 지금 503을 본다. (b) \textbf{티켓} — 원인이며 개발 환경에서는 해가 없지만 증상이 발생하면 결정적인 맥락이 된다. (c) \textbf{호출} — 이를 알 수 있는 유일한 외부 관점에서 출입구가 실제 사용자를 거부한다. (d) \textbf{티켓} — OOM 한 번은 5장의 완충된 사고이며, 플릿 전체에서 나타나 시스템 한도 오류가 되었을 때만 호출한다. (e) \textbf{호출} — 6 h 동안 6×인 느린 누출은 빠른 규칙에 걸리지 않으므로 이를 위해 짝 구간을 둔다. (f) \textbf{티켓} — 흔히 드리프트가 다가오며 생기는 풀 압력이므로 근무 시간에 조사한다(§14.5). (§14.4)}

\qitem{6}{수락률은 작은 초안 모델이 실제 트래픽을 얼마나 잘 예측하는지 추정한다. 요청마다 무료로 갱신되는 분포 일치 측정기다. 입력이 움직이면 초안이 먼저 틀리고, 품질 변화가 눈에 보이기 전에 TPOT가 초안 없는 하한으로 나빠져 \textbf{비용} 회귀가 먼저 나타난다. 올바른 순서로 도착하는 조기 경보다. 초안 모델이나 엔진 업데이트도 수락률을 움직일 수 있다는 교란 요인이 있다. 추세가 시작된 날 릴리스 세 요소의 다이제스트가 바뀌었는지 레지스트리 계보로 구분한다. 교란 여부와 관계없이 감시 패널에 둘 가치가 있지만, 판정이 아니라 합창의 한 목소리로 사용한다. (§14.5, 11장)}

\subsection*{C. 논술형(채점 기준)}

\qitem{1}{확인할 내용: 티켓이 11분 먼저 발생했음을 짚은 두 경보의 연표, 반사실인 85분을 단순 주장하지 않고 논증해 턴과 사용자 시간으로 계산한 영향, 소유자와 만료가 없는 오래된 ConfigMap 및 빠진 배치 정책이라는 메커니즘, 엔지니어 개인을 명시적으로 비난에서 제외한 서술, 서로 다른 종류의 조치 세 가지다. 경보를 정책으로 승격하고 CI에 만료 검사를 추가하고 캡션을 수정하는 식으로, 각각 담당자와 날짜가 있어야 한다. “만들지 않을 것” 문단이 답안의 수준을 가른다. 예를 들어 모든 곳에 파드별 반친화성을 강제하면 §12.4에서 정당하게 공유하는 개발 풀이 깨지고, 전용 이상 감지기는 §14.5의 함정이다. 느낌이 아니라 비용으로 판단해야 한다. (§§14.1\textasciitilde{}14.4)}

\qitem{2}{예산이 실제 선택을 강제한다. 측정한 트래픽을 바탕으로 규칙마다 예상 월간 발화 횟수를 \textbf{세고}, 이를 합해 2와 비교한 경보 목록을 높이 평가한다. 호출의 핵심에는 소진율 쌍, 선행 증상 하나, 출입구 프로브가 있어야 한다. 데드맨은 외부에 두고, 사고 하나가 호출 하나가 되도록 그룹화하며, 억제 규칙은 Watchdog에 절대 적용하지 않는다. 실패 유형을 낮춘 부분이 논술의 핵심이다. 레플리카 하나의 OOM을 티켓으로 낮추면 재시작 한 바퀴만큼의 지연을 받아들이게 되고(§3.6/§5.6 계산), 선점 급증을 티켓으로 낮추면 검토 전까지 TTFT가 천천히 악화되는 것을 받아들인다. 받아들인 고통을 \textbf{수치화}해야 만점을 받을 수 있다. (§14.4)}

\qitem{3}{흐름은 다음과 같다. 금요일에는 답변 길이가 몇 분 만에 211→54로 바뀌었고, 30\% 고정 가드레일을 카나리 컨트롤러가 읽어 자동 후퇴했다(13장). 기록 규칙은 영구적으로 5분 구간을 계산하고 이동 기준선 대비 임계값을 2행과 당직자가 읽어 조사한다. 드리프트 감시 지표는 게이트 구간의 기준과 21일 추세를 PSI 방식 구간으로 비교하며 월요일 검토자가 보고 대응 사다리를 밟는다. 좋은 답안은 불변 요소를 짚는다. 숫자 하나를 세 \textbf{구간}과 세 대응 속도에서 사용한다. 여러 시간 척도에서 쓸 수 있는 지표만 내보낼 가치가 있다는 결론으로 이어진다. 가드레일은 영속성을 얻지만 허영 지표는 두 번째 계측 도구를 얻지 못한다. (§§14.3\textasciitilde{}14.5, 13장)}

\qitem{4}{정직한 원장은 스택 23달러 + 유휴 골짜기를 쓰는 판정기 약 0달러 + 보존 기간으로 제한된 TSDB 증가분을 기록하고, 제약 자원인 당직자의 주의와 비교한다. 클라우드 청구서에 없어도 경보 피로는 \textbf{비용}이다. 추가 패널과 경보의 한계 신호로 조치할 수 없게 되는 순간 더 이상 비용을 정당화하지 못한다. 상자 14-1의 죄악은 그 지점을 넘었다는 증상이다. 새벽 2시 재현에는 1행의 예산 그래프와 7주차 문서에서 미리 합의한 완화 조항을 인용하고 “계산과 서명”에 도달해야 한다. 주장을 옹호하는 좋은 답안은 지루함 = 자동화 + 버전 관리 + 점진성 + 가역성(13장) + \textbf{관측됨}(14장)을 추적하고, 23달러로 첫 사고부터 둘째 날 밤의 차이인 사용자 시간 65분을 샀음을 지적한다. 반박하는 좋은 답안은 지루함이 작업이 아니라 운영의 속성이며, 흥미로운 문제는 드리프트, 평가, 제품처럼 스택 위로 이동했음을 지적한다. 어느 입장이든 계산이 일관되면 최고 등급을 받는다. (§14.4, 장 전체)}

\namedsection{더 읽을거리}{}

\textbf{[이번 주]} 표시는 실습에 직접 도움이 되는 자료이고, \textbf{[N주차]} 표시는 뒤에서 활용할 자료다.

\begin{tbbitemize}
  \item \textbf{[이번 주] Google SRE Book, 6장 “Monitoring Distributed Systems” + SRE Workbook, 5장 “Alerting on SLOs”} — 네 가지 황금 신호, 증상/원인 원칙, 표 14-2가 그대로 채택한 다중 구간·다중 소진율 정책을 다룬다. 그림 14-5를 옆에 두고 Workbook의 해당 장을 읽어 보라. 숫자가 의도적으로 일치한다.
  \item \textbf{[이번 주] Rob Ewaschuk, “My Philosophy on Alerting”} — “호출은 조치할 수 있고, 긴급하며, 실제여야 한다”라는 원칙의 씨앗이 된 짧은 문서다. 증상에는 호출을, 원인에는 티켓을 보낸다는 현대적 원칙도 여기서 시작되었다. 20분을 투자하면 경력 내내 보상을 얻는다.
  \item \textbf{[이번 주] Prometheus 문서의 “Histograms and summaries”, “Querying basics”와 Robust Perception 블로그} — 버킷 설계, 요약을 집계할 수 없는 이유, §14.2에서 다룬 날카로운 모서리, rate()의 미묘한 동작을 설명한다. Brian Brazil의 블로그에는 이 분야의 구전 지식이 글로 축적되어 있다.
  \item \textbf{[이번 주] kube-prometheus-stack 차트 문서 + Grafana 대시보드 모범 사례 + NVIDIA DCGM-exporter README} — 실습에서 설치할 세 구성 요소의 문서다. DCGM README의 지표 표에는 사용률의 세부 의미와 PROF\_* 계열이 있으며, 파드 매핑 절은 4행의 결정적 증거를 가능하게 한다.
  \item \textbf{D. Sculley 외, “Hidden Technical Debt in Machine Learning Systems” (NeurIPS 2015)} — §14.5에서 모니터링하는 얽힘, 피드백 루프, 고정된 모델 아래에서 변하는 세상에 이름을 붙인 논문이다. “monitoring and testing” 절은 마지막에 읽어 보라. 이 장의 내용을 10년 먼저 네 문단으로 압축해 두었다.
  \item \textbf{Chip Huyen, Designing Machine Learning Systems, 8장(“Data Distribution Shifts and Monitoring”)} — 이 강좌가 PSI 설명에 압축한 구간, 기준, 계절성과 같은 추정의 미묘함을 포함해 드리프트 분류를 책 한 장 깊이로 다룬다.
  \item \textbf{The Joint Commission, Sentinel Event Alert \#50: “Medical device alarm safety in hospitals” (2013)} — 병원 안전 공지를 엔지니어의 눈으로 읽고 §14.4의 모든 원칙이 가용성의 숫자 9 대신 사람의 생명으로 표현된 부분을 찾아보라. “85–99\%”의 근거이며, 시끄러운 경보 목록에 반대할 때 제시할 수 있는 가장 강력한 논거다.
  \item \textbf{[흥미 읽을거리] “13 Minutes to the Moon” (BBC World Service 팟캐스트, 시즌 1)} — 1202 경보 에피소드에서 Garman, Bales, Hamilton이 §14.4 도입부의 이야기를 직접 들려준다. 이어서 Apollo 12의 SCE-to-AUX를 설명하는 Sy Liebergot의 이야기를 찾아 §14.1과 비교해 보라. 관제실 음성은 역사의 모습을 한 운영 교육으로 지금도 훌륭하다.
\end{tbbitemize}

\vspace{5.0pt}

\begin{calloutbox}{tbbblue}{tbbbluefill}
{\normalsize\bfseries\color{tbbblue}다음 주 — 15주차: 시연의 날 — 비행을 완수하다}\par\vspace{3pt}
15주 전 이 강좌는 30초 동안의 침묵과 빙글빙글 도는 커서로 시작했다. 모델은 작동하지만 서비스는 작동하지 않았던 사례 A다. 15주차는 그 반대편 책갈피인 팀 프로젝트 발표와 시연이다. 각 팀은 이 책이 계층별로 하나씩 조립해 온 시스템을 실시간으로 보여 준다. 기계를 다루는 장은 모두 끝났고 이제 남은 것은 증거다. 가장 뛰어난 시연은 가장 지루하게 흘러갈 것이다. 키와 할당량이 있는 게이트웨이를 통해 트래픽이 흐르고(6장), 팀이 11주차에 다시 합의한 SLO 문장 안에서 응답하며, 의도적으로 선택한 숫자에 청구서를 마감한 플릿(\$4,762 → \$877 → \$900, 모든 단계를 설명할 수 있다)에서 실행된다. 목격자 앞에서 나쁜 릴리스를 거부한 파이프라인이 배포를 맡는다(13장). 시연 중에는 이번 주에 만든 네 행 대시보드가 모든 것을 지켜본다. 발표 도중 문제가 생기면 실시간으로 진단하라. 그것은 실패가 아니라 만점에 이르는 가장 높은 수준의 시연이다. 레지스트리 계보, 게이트 보고서, 두 번의 빨간 실행, 훈련 사후 분석, 드리프트 보고서, 최종 원장을 준비하라. 시연은 30분이지만 산출물의 흔적이 성적을 결정한다. 임무를 완수하라.
\end{calloutbox}


%% Epilogue — 지루한 서비스
\frontchapter{지루한 서비스}{에필로그}
\chaptersubtitle{15주, 하나의 대화 모델, 아무 일도 일어나지 않았다는 조용한 승리}

\begin{calloutbox}{tbbblue}{tbbbluefill}
{\normalsize\bfseries\color{tbbblue}이 맺음말에 있는 것과 없는 것}\par\vspace{3pt}
이 맺음말에는 학습 목표도, 실습도, 연습문제도 없으며 시험에 나올 내용도 없다. 실제로 다녀온 여행을 다룬 박물관을 걷듯 읽으면 된다. 데모를 마친 뒤 커피 한 잔을 들고, 즐거운 속도로 천천히 둘러보자. 여기에는 여러분이 만든 서비스의 \textbf{일대기}, 직접 얻은 한 줄 문장을 모은 \textbf{기념품 진열장}, 이제 참석할 자격이 생긴 \textbf{저녁 만찬}, 이 책의 어떤 내용이 곧 만료되고 무엇이 10년 동안 보증되는지 적은 \textbf{보증서}, 아직 겪지 않았지만 언젠가는 찾아올 밤을 위한 짧은 \textbf{편지}가 들어 있다.
\end{calloutbox}

\vspace{4.0pt}

\namedsection{E.1}{한 서비스의 일대기}

이 서비스는 1주 차에 부끄러운 모습으로 태어났다. 완벽한 모델 실행에는 0.31초가 걸렸지만, 그 앞뒤를 \textbf{30초의 침묵과 빙글빙글 도는 커서}가 감쌌다. 사례 1-A의 진짜 주제는 버그 자체가 아니었다. 첫 요청에서 차가운 모델을 불러오고, 아무도 그리지 않은 대기열이 생긴 일이 전부가 아니었다. 핵심은 한 학기 내내 배운 끝에 두 줄로 내릴 수 있게 된 진단, 곧 \textit{어느 구간이 병목인가?}라는 질문이었다. 그 진단 속에는 더 큰 선언이 숨어 있었다. 학습된 모델과 작동하는 서비스는 서로 다른 결과물이며, 그 둘 사이에는 이 책의 모든 내용이 놓여 있다. 이 서비스의 첫 주는 여느 서비스와 같았다. 완벽한 두뇌는 있었지만 몸도, 예의도, 약속도 없었고 무엇에 얼마가 드는지도 몰랐다.

2주 차부터 5주 차까지는 몸을 만드는 시간이었다. 서비스는 \textbf{머신}을 얻었고, 곧 그 머신이 공손한 허구임을 배웠다. 하이퍼바이저 위의 손님일 뿐이며, GPU만은 패스스루가 허구의 틈으로 실리콘을 손대지 않고 건네기 때문에 실제였다. 빌린 것 가운데 가상화되지 않은 유일한 부분이 바로 돈을 내는 부분이다. 이어 \textbf{울타리와 계량기}를 얻었다. 네임스페이스는 자기 자신만 보게 했고, cgroup은 자기 몫만 먹게 했다. 그리고 처음 죽었다. \code{exit 137}, 커널의 작별 인사다. 이 책에서 영원히 외우라고 요구하는 유일한 계산이기도 하다(128 + signal 9). 다음에는 \textbf{운반할 수 있는 몸}을 얻었다. 계층으로 쌓고 다이제스트로 고정한 이미지 덕분에 "내 머신에서는 되는데"가 "어디서나 바이트 단위로 똑같은 내 머신 그 자체"로 바뀌었다. 마지막으로 \textbf{반사 작용}을 얻었다. 프로브는 트래픽을 통제하고, 재시작은 회의를 기다리지 않는다. 플릿 관리자는 사용자의 의도를 데이터베이스로, 현실을 고쳐야 할 버그로 취급한다. 5주 차가 되자 서비스는 자신의 죽음에서도 살아남을 수 있었다. 그러나 아직 \textit{신뢰할} 수는 없었다. 반사 작용은 행동하지만 설명하지 못한다. 그 빚은 14주 차까지 조용히 남아 있었다.

6주 차에는 \textbf{문}이 생겼다. 로드 밸런서, TLS, 키, 할당량으로 이루어진 문이었다. 이 문과 함께 강의에서 처음으로 진짜 상업적인 문장이 등장했다. \textit{curl에 답하는 컨테이너}와 \textit{누군가에게 요금을 청구할 수 있는 서비스}의 차이는 문이다. 문은 외부 세계가 거절을 표현하는 여러 어휘도 알려 주었다. 503과 504, 문장 중간에서 끊기는 스트림이 그것이다. 또한 자기 벽 바깥에 서서 안쪽을 점검하는 규율을 가르쳤다. 6주 차에는 편집증처럼 보였던 습관이 14주 차에는 호출 정책이 되었다.

7주 차는 사춘기였다. 서비스는 \textbf{약속하는 법}을 배웠다. 막연히 "빨라야 한다"가 아니라 형태가 있는 약속이었다. \textit{이 위치에서, 이 시간 구간 동안, 이 부하와 이 예산 아래에서 측정한 p99가 500ms 이하여야 한다.} 백분위수의 문법도 배웠다. 평균은 아무도 설명하지 못하고, 백분위수는 더할 수도 평균 낼 수도 없다. 이 규칙은 이후 히스토그램을 입고 두 번 더 돌아왔다. 서비스는 용돈, 즉 오류 예산도 받았다. 덕분에 신뢰성은 미덕을 겨루는 일이 아니라 장부를 관리하는 일이 되었다. 그러자 첫 청구서가 도착했다. 단순하게 구성한 플릿의 월 비용은 \textbf{4,762달러}였다. 이 숫자는 부모가 치켜뜬 눈썹처럼 5주 동안 강의 위에 걸려 있었다. 9주 차에는 군중이 몰려왔다. 서비스는 대부분의 폭풍이 스스로 만든 것임을 배웠다. 재시도 폭풍은 이 책에서 피해자가 직접 조직하는 유일한 기상 현상이다. 그리고 일찍, 싸게 거절하는 법을 익혔다. \textit{거절이 곧 제품이다.} 10주 차에는 절차를 갖춰 다이어트를 시작했다. 비용을 절반으로 줄이는 INT8은 측정 \textit{전에} 허용 오차를 적어 둔 관문을 통과해야만 허용되었다. 판정 관문 없는 손실 최적화는 규모만 커진 희망 사항일 뿐이기 때문이다. 청구서는 \textbf{2,381달러}로 내려갔다.

11주 차에 서비스는 \textbf{말하는 법}을 배웠고, 말하기는 서비스의 물리 법칙을 바꾸었다. 하나였던 지연 시간은 두 개의 시계가 되었다. 첫 단어를 기다리는 시간과 그 뒤 단어가 나오는 속도다. 메모리는 부동산이 되었다. KV 캐시는 토큰당 112KB, 긴 대화 하나에는 448MiB가 필요했다. 테이블의 좌석은 메가바이트 단위로 값을 매겼고 연속 배칭으로 채웠다. 위대한 할당 기법인 PagedAttention은 malloc의 오래된 비극, 즉 최댓값만큼 예약하고 차이는 버리는 문제를 1962년 운영체제가 해결한 방식으로 고쳤다. 페이지 테이블이다. 스트리밍에 맞춰 SLO를 공개적으로, 계산을 보이며 다시 협상했다. 12주 차에는 \textbf{호흡하는 법}을 배웠다. 하루의 계단을 따라 레플리카가 움직였다. 파드는 파도를 쫓고, 노드는 예보를 따랐다. 콜드 스타트의 150초도 숨기지 않고 비용에 넣었다. 새뮤얼 인설이 전기를 팔던 방식으로 용량을 샀다. 기본 부하는 약정 용량으로, 피크는 스팟으로 감당하고 골짜기에는 유용한 일을 찾아 넣었다. 청구서는 \textbf{1,513달러}, 이어 \textbf{877달러}로 줄었다. 그리고 강의에서 가장 조용한 교훈이 남았다. \textit{자체 구축 대 외부 도입은 구호였던 적이 없다. 언제나 부하율의 문제였다.}

13주 차에는 \textbf{소란 없이 바꾸는 법}을 배웠다. 금요일 배포에서 완벽한 모델이 홀로 이동한 설정 파일 하나 때문에 망가졌다. 이 사건은 레지스트리, 파이프라인, 카나리를 낳았다. 릴리스는 증거를 붙인 다이제스트 \textit{세 개}의 묶음이 되었고, 파이프라인의 빨간 실행 결과는 일어나지 않은 금요일이라는 제품이 되었다. 홀데인의 새는 자동 후퇴 반사 작용을 지닌 카나리로 다시 태어났다. 배포는 엔지니어링에서 가장 지루한 문장으로 바뀌었다. \textit{자동화되고, 버전이 있으며, 점진적이고, 되돌릴 수 있다.} 14주 차에는 반사 작용에 없던 세 가지를 더해 신경계를 완성했다. \textbf{기억}은 새벽 2시의 두뇌가 질문하는 순서대로 배치한 네 줄의 대시보드였다. \textbf{주의}는 감이 아니라 소진율을 보고, 그것도 드물게 울리는 호출이었다. 경보 피로에는 실제 희생자가 있기 때문이다. \textbf{예지}는 매주 드리프트를 지켜보는 일이었다. 카나리는 절벽을 잡지만 세상은 비탈을 더 좋아한다. 첫 실제 장애는 12주 차에 예고한 바로 그 수수께끼, 같은 GPU를 쓰는 이웃 때문에 일어났다. 장애는 20분 만에 끝났고, 세 개의 조치 항목과 비난 0건을 남겼다. 마지막 청구 금액은 \textbf{정확히 900달러}였다. 마지막 23달러는 나머지 모든 부분을 감시하는 부분의 비용이었다.

그리고 15주 차, 서비스는 낯선 사람들 앞에 서서 시연되었다. 모든 것이 설계대로였다면 데모는, 자랑스럽게 말하자, \textbf{지루했다.} 트래픽은 키가 달린 문을 통과했고, 답은 숫자로 적은 약속 안에서 나왔다. 청구서는 스스로 선택한 금액으로 끝났고, 프로젝터 한쪽에서는 대시보드가 잔잔히 숨 쉬었다. 운영에서 지루함은 흥분이 없다는 뜻이 아니다. 관문과 리허설, 오후 2시에 미리 그린 그래프로 흥분의 값을 \textit{선불로 치렀다}는 뜻이다. 이 서비스의 일대기는 훌륭한 모든 인프라의 일대기와 같은 방식으로 끝난다. 공개된 자리에서, 의도한 대로, 아무 일도 일어나지 않는다.

\begin{tbbfigure}
  \adjustbox{max width=0.94\linewidth,max totalheight=0.46\textheight}{\includegraphics{figures/figE-1_line15.png}}\\[3pt]
  \figcaption{그림 E-1  15호선. 하나의 서비스, 15개 정거장, 세 개의 운임 구간. 엔지니어링이 개선될수록 운임은 4,762달러에서 스스로 선택한 900달러까지 내려간다. 열차는 8주 차에 정차하지 않는다. 종착역은 시연을 마친 지루한 서비스다. 현재 위치는 여기다.}
\end{tbbfigure}

\namedsection{E.2}{기념품 진열장}

이 책의 각 장에는 일부러 들고 다닐 수 있게 만든 문장을 적어도 하나씩 심었다. 지친 머릿속에서도 살아남을 만큼 짧고, 결정을 내릴 수 있을 만큼 날카로운 문장이다. 이는 문체상의 버릇이 아니라 핵심 설계였다. 프로덕션 엔지니어링은 대개 사람이 가장 온전하게 생각하기 어려울 때 수행한다. 새벽 2시, 데모 당일, 분기 9주 차 같은 때다. 좋은 격언은 \textbf{성능이 떨어진 하드웨어를 위한 압축}이다. 전체 논증을 미리 계산해 잠이 부족한 사람도 닿을 수 있는 곳에 캐시해 둔다. 이제 그 문장들을 얻게 된 사건과 함께 박물관 명패에 올려 두었다. 가끔 먼지를 털어 주자.

\begin{tbbtablewrap}
  \begin{tabular}{@{}>{\raggedright\arraybackslash}m{\dimexpr 0.3564\tbltextwidth-2\tabcolsep\relax} >{\raggedright\arraybackslash}m{\dimexpr 0.1596\tbltextwidth-2\tabcolsep\relax} >{\raggedright\arraybackslash}m{\dimexpr 0.4840\tbltextwidth-2\tabcolsep\relax}@{}}
  \arrayrulecolor{tbbnavy}\specialrule{1pt}{0pt}{0pt}
  \rowcolor{tbbtablehead}{\centering \bfseries\color{tbbnavy}기념품} & {\centering \bfseries\color{tbbnavy}얻은 때} & {\centering \bfseries\color{tbbnavy}오래 남는 이유} \\
  \arrayrulecolor{tbbnavy}\specialrule{0.7pt}{0pt}{2pt}
  {\textbf{exit 137: 커널의 작별 인사}} & {3주 차} & {어떤 오케스트레이터에서든 앞으로 만나게 될 모든 OOM을 설명한다. 128 + 9. 상한은 실제였고 계량기는 정직했다.} \\
  \arrayrulecolor{tbbtableborder}\hline
  {\textbf{시스템을 보는 시야는 마운트만큼 정직하다}} & {3주 차} & {네임스페이스 안의 ps부터 cgroup 파일을 읽는 익스포터까지, 무엇을 읽는지 확인하지 않으면 계측도 거짓말한다.} \\
  \arrayrulecolor{tbbtableborder}\hline
  {\textbf{태그는 떠돌지만 다이제스트는 변하지 않는다}} & {4주 차} & {13주 차의 레지스트리 전체는 이 문장이 자란 모습이다. 레이블이 아니라 내용으로 버전을 정하면 롤백은 포인터 하나가 된다.} \\
  \arrayrulecolor{tbbtableborder}\hline
  {\textbf{프로브는 반사 작용이고, 모니터링은 기억이다}} & {5→14주 차} & {관측 가능성 없는 자가 치유는 장애를 억눌러 숨기는 일이다. 두 신경계가 모두 필요하며 서로 경쟁하지 않는다.} \\
  \arrayrulecolor{tbbtableborder}\hline
  {\textbf{하나의 상태 엔드포인트가 유일한 진실 공급원이다}} & {6주 차} & {"정상"의 정의가 둘이면 분기마다 정확히 한 번, 언제나 가장 나쁜 장애에서 서로 다른 답을 낸다.} \\
  \arrayrulecolor{tbbtableborder}\hline
  {\textbf{평균은 아무도 설명하지 못한다}} & {7주 차} & {이 책에서 가장 많이 반복한 문장이다. 현장에서 가장 많이 되풀이하는 실수를 가리키기 때문이다.} \\
  \arrayrulecolor{tbbtableborder}\hline
  {\textbf{백분위수는 더할 수도, 평균 낼 수도 없다}} & {7→14주 차} & {진지한 시스템이 히스토그램을 제공하는 이유다. avg-of-p99s는 실제가 9.08초일 때 2.76초라고 말한 적이 있다.} \\
  \arrayrulecolor{tbbtableborder}\hline
  {\textbf{L = λW}} & {7→11주 차} & {회의실까지 따라오는 유일한 공식이다. 좌석, 도착률, 대기 시간을 대기열과 KV 풀 모두에 같은 대수로 적용한다.} \\
  \arrayrulecolor{tbbtableborder}\hline
  {\textbf{빠른 거절이 느린 가능성보다 낫다(거절이 곧 제품이다)}} & {9주 차} & {일찍, 싸게 덜어 내라. 헛된 대기열을 만드는 시스템은 자신의 폭풍에 스스로 먹이를 준다.} \\
  \arrayrulecolor{tbbtableborder}\hline
  {\textbf{허용 오차는 측정값보다 먼저 정한다}} & {10주 차} & {측정하기 전에 "충분히 좋다"의 뜻을 정하지 않으면 측정값이 대신 정한다. 손실 최적화에는 절차가 필요하다.} \\
  \arrayrulecolor{tbbtableborder}\hline
  {\textbf{높은 캐시 점유는 정상이고, 경보해야 할 것은 대기열이다}} & {11주 차} & {가득 찬 캐시는 설계가 작동하는 모습이다. 커지는 대기열은 수요가 좌석을 앞지른다는 뜻이다. 이 계기를 잘못 읽는 통과 의례는 건너뛰자.} \\
  \arrayrulecolor{tbbtableborder}\hline
  {\textbf{파드는 파도를 쫓고, 노드는 예보를 따른다}} & {12주 차} & {두 시간 척도, 두 메커니즘, 하나의 청구서. 여섯 단어로 요약한 오토스케일링이다.} \\
  \arrayrulecolor{tbbtableborder}\hline
  {\textbf{지루함 = 자동화 + 버전 관리 + 점진성 + 가역성(+ 관측)}} & {13→14주 차} & {14주 차의 보완을 더한 배포 문장이다. 지루하다는 말은 운영이 줄 수 있는 최고의 찬사다.} \\
  \arrayrulecolor{tbbtableborder}\hline
  {\textbf{카나리는 절벽을 잡지만 비탈은 잡지 못한다}} & {13주 차} & {모든 계측기에는 보지 못하는 시간 척도가 있다. 분 단위 계측기와 주 단위 계측기를 짝짓지 않으면 세상은 둘 사이로 흘러 지나간다.} \\
  \arrayrulecolor{tbbtableborder}\hline
  {\textbf{증상에는 호출, 원인에는 티켓}} & {14주 차} & {당직 순환을 온전하게 유지하고 호출기를 믿게 만드는 짧은 라우팅 규칙이다. 이 교훈의 수업료는 병원이 치렀다.} \\
  \arrayrulecolor{tbbtableborder}\hline
  {\textbf{새벽 2시에는 오후 2시에 그린 그래프만 보인다}} & {14주 차} & {이 책 후반부의 운영 내용을 한 문장으로 요약한다. 밤이 질문하기 전에 그 질문을 그래프로 그려 두어라.} \\
  \arrayrulecolor{tbbnavy}\specialrule{1pt}{2pt}{0pt}
  \end{tabular}\par
  \figcaption{기념품 진열장. 15주 동안 모은 16개의 명패다. 하나만 간직한다면 마지막 문장을 고르자. 나머지를 간직하는 방법까지 그 안에 들어 있다.}
\end{tbbtablewrap}

\namedsection{E.3}{저녁 만찬}

이 책에는 불안할 때 드러나는 버릇이 하나 있었다. 문제가 정말 어려워지기만 하면 방을 나섰다. 다음 장으로 간 것이 아니라 다른 \textit{세기}로 도망쳤다. 7주 차에 청구서가 도착하자 \textbf{1902년}으로 달려가 전차와 제빙 공장에 킬로와트시를 팔던 사람을 만났다. 플릿의 규모를 정해야 하자 \textbf{1961년}으로 가서 아무 가정도 하지 않았기에 어떤 상황에서도 살아남은 대기열 증명을 찾았다. KV 풀이 쓸 수 없는 조각으로 산산이 흩어지자 \textbf{1962년}의 페이지 테이블로 향했다. 배포를 판정해야 할 때는 \textbf{1890년대}로 가서 온혈의 보초를 탄광에 데려간 생리학자를 만났다. 호출기에 예절이 필요하자 \textbf{1969년} 휴스턴의 한 방으로 가서 손으로 쓴 종이에 적힌 네 자리 경보 코드를 읽었다. 이는 결코 장식이 아니었다. 눈치채지 못할 만큼 반복해서 전한 이 강의의 가장 깊은 주장이었다. \textbf{서빙의 어려운 문제는 새로운 문제가 아니다.} 단위만 바뀐 오래된 문제다. 킬로와트는 GPU 시간으로, 간선 회선은 좌석으로, 일산화탄소는 나쁜 백분위수로 바뀌었다. 먼저 그 문제를 해결한 사람들은 이제 직업적으로 여러분의 동료다.

그래서 저녁 만찬을 연다. 이 맺음말의 유일한 시험이며 스스로 치르는 시험이다. 지금의 기술 스택이 발명되기도 전에 세상을 떠난 사람들과 한 식탁에 앉아 자기 자리를 지킬 수 있는가? 15주 전의 정직한 답은 아니었다. 지능이 부족해서가 아니라 \textit{같은 청구서}를 받아 본 적이 없었기 때문이다. 유휴 용량에 돈을 내거나, 백분위수에 서명하거나, 보초 하나를 믿고 플릿을 걸거나, 깨워야 할 이유가 충분한 머신에게 잠을 깨워진 적이 없었다. 이제는 모두 경험했다. 이 사람들의 학문을 배웠기 때문이 아니라 같은 우주에서 청구서를 받았기 때문에 이 식탁에 앉을 자격이 있다. 피크와 골짜기, 좌석과 기다림, 절벽과 비탈, 작업자와 새가 적힌 같은 청구서다. 좌석도 그에 맞춰 배치했다. 각 의자의 번호는 \textbf{초대장을 보낸 주 차}다.

\begin{tbbtablewrap}
  \begin{tabular}{@{}>{\raggedright\arraybackslash}m{\dimexpr 0.0957\tbltextwidth-2\tabcolsep\relax} >{\raggedright\arraybackslash}m{\dimexpr 0.2074\tbltextwidth-2\tabcolsep\relax} >{\raggedright\arraybackslash}m{\dimexpr 0.3138\tbltextwidth-2\tabcolsep\relax} >{\raggedright\arraybackslash}m{\dimexpr 0.3830\tbltextwidth-2\tabcolsep\relax}@{}}
  \arrayrulecolor{tbbnavy}\specialrule{1pt}{0pt}{0pt}
  \rowcolor{tbbtablehead}{\centering \bfseries\color{tbbnavy}좌석} & {\centering \bfseries\color{tbbnavy}손님} & {\centering \bfseries\color{tbbnavy}같은 문제의 옛 모습} & {\centering \bfseries\color{tbbnavy}식탁에서 나눌 이야기} \\
  \arrayrulecolor{tbbnavy}\specialrule{0.7pt}{0pt}{2pt}
  {\centering \textcolor{tbbteal}{\textbf{7}}} & {\textbf{존 리틀}\newline \textit{대기행렬, 1961}} & {고객이 얼마나 빨리 도착하고 얼마나 오래 머무는지 알 때 가게 \textit{안}에 몇 명이 있는가. 두 값에 아무런 가정도 두지 않고 증명했다.} & {냅킨에 KV 풀을 그려 보이면 곧바로 알아볼 것이다. 좌석, 도착, 기다림이다. 그의 대수는 7주 차 레플리카와 11주 차 캐시에 두 번 사용했다. 그는 둘이 같은 용법이었다고 조용히 짚어 줄 것이다.} \\
  \arrayrulecolor{tbbtableborder}\hline
  {\centering \textcolor{tbbteal}{\textbf{9}}} & {\textbf{마거릿 해밀턴}\newline \textit{아폴로 소프트웨어, 1969}} & {하강 중 과부하가 걸리자 설계대로 우선순위가 낮은 작업을 버리고도 착륙에 성공한 4KB 실행 시스템.} & {군중이 몰릴 때 서비스가 무엇을 버리는지, 그 대상을 미리 정했는지 물을 것이다. 9주 차의 답을 준비하자. \textit{거절이 곧 제품이다.} 그의 소프트웨어는 달까지 날아갔고, 여러분의 서비스는 429를 반환한다. 원리는 같고 고도만 다르다.} \\
  \arrayrulecolor{tbbtableborder}\hline
  {\centering \textcolor{tbbteal}{\textbf{11}}} & {\textbf{Atlas 페이징 팀}\newline \textit{맨체스터, 1962}} & {프로그램은 연속된 메모리를 요구했지만 머신에는 흩어진 조각만 남았다. 고정 크기 페이지와 주소를 변환하는 테이블로 해결했다.} & {그들의 기술이 이제 \textit{어텐션 내부에서} 실행된다고 알려 주자. 16토큰 블록, 4\% 미만의 낭비, 2\textasciitilde{}4배의 처리량이다. 누구도 그 생각을 개선하지 않았고 목표만 바꾸어 적용했다. 놀라지 않을 것이다. 간접 참조에 익숙한 사람은 좀처럼 놀라지 않는다.} \\
  \arrayrulecolor{tbbtableborder}\hline
  {\centering \textcolor{tbbteal}{\textbf{12}}} & {\textbf{새뮤얼 인설}\newline \textit{시카고 에디슨, 1902}} & {저녁 조명 피크에 맞춰 지은 발전소는 하루 20시간 놀았다. 그의 집착은 \textbf{부하율}, 즉 피크 대비 평균 수요였다.} & {수프가 나오기도 전에 새벽 4시 GPU가 무엇을 하는지 물을 것이다. \textit{드리프트 평가가 유휴 골짜기에 올라탄다}고 답할 수 있다면, 12주 차의 약정 바닥 위에서 14주 차의 판정 작업이 돈다는 뜻이다. 디저트가 나오기 전에 채용하려 들 것이다.} \\
  \arrayrulecolor{tbbtableborder}\hline
  {\centering \textcolor{tbbpurple}{\textbf{13}}} & {\textbf{존 스콧 홀데인}\newline \textit{광산 생리학, 1896}} & {대사 속도가 더 빠른 보초를 작업자보다 먼저 노출했다. 관찰 비용과 대응 비용이 싸고 몇 분의 경고 시간을 벌었다.} & {트래픽 5\% 분할을 보고 만족할 것이다. 그의 새를 은퇴시킨 탐지기가 나오기까지 90년이 걸렸지만, 여러분의 것은 자동으로 지켜보다 몇 초 만에 후퇴한다고 담담히 말할 것이다. 그러고는 몸을 기울여 묻는다. \textit{카나리는 절벽을 잡는다. 비탈은 무엇이 지켜보는가?}} \\
  \arrayrulecolor{tbbtableborder}\hline
  {\centering \textcolor{tbbpurple}{\textbf{14}}} & {\textbf{존 에런과 잭 가먼}\newline \textit{미션 컨트롤, 1969}} & {훈련에서 본 이상한 텔레메트리 패턴 하나를 기억해 낸 SCE to AUX, \textit{실패한} 시뮬레이션 뒤 손으로 쓴 참고표로 해결한 1202.} & {새벽 2시에는 타고난 천재성보다 알아볼 수 있는 패턴이 낫다는 사실을 보여 준 최고의 두 사례다. 참고표가 어디에 있는지(런북), 어떤 실패가 그것을 쓰게 했는지(일부러 두 번 망가뜨린 실습), 장애 뒤 누가 갱신하는지 물을 것이다. 정직하게 답하자. 두 사람은 거짓말을 알아챈다.} \\
  \arrayrulecolor{tbbtableborder}\hline
  {\centering \textcolor{tbbred}{\textbf{15}}} & {\textbf{여러분}\newline \textit{15주, 하나의 서비스}} & {30초의 침묵 뒤에 있던 완벽한 모델을 약속이 글로 적힌 지루한 서비스로 바꾸었다.} & {대부분은 듣기만 하자. 청구서가 오면 비용이 부하율에 따라 나뉘고, 여러분의 금액은 스스로 고른 숫자, 곧 \textbf{정확히 900달러}로 끝난다고 말하자. 이 식탁의 모든 사람은 그 사실이 왜 그토록 만족스러운지 정확히 이해한다.} \\
  \arrayrulecolor{tbbnavy}\specialrule{1pt}{2pt}{0pt}
  \end{tabular}\par
  \figcaption{좌석 순서대로 적은 손님 명단. 1주 차부터 6주 차를 대표해 커널도 초대했지만 참석하지 못한다는 답을 보내왔다. 누군가는 불을 계속 켜 두어야 한다.}
\end{tbbtablewrap}

\begin{tbbfigure}
  \adjustbox{max width=0.94\linewidth,max totalheight=0.46\textheight}{\includegraphics{figures/figE-2_dinner.png}}\\[3pt]
  \figcaption{그림 E-2  좌석 배치도. 의자 번호는 초대장을 보낸 주 차다. 커널은 불참을 알렸고, 상석은 사례 1-A 때부터 줄곧 여러분의 자리였다. 비용은 세 번째 식사 규칙에 따라 나눈다.}
\end{tbbfigure}

\begin{calloutbox}{tbbteal}{tbbtealfill}
{\normalsize\bfseries\color{tbbteal}이 만찬을 여는 까닭}\par\vspace{3pt}
이 만찬은 보상이 아니라 방법이다. 에너지든, 토큰이든, 에이전트든, 2030년대가 무엇에 계량기를 달든 다음번의 터무니없는 청구서 역시 새로운 문제는 아닐 것이다. 그 방에서 가장 빠른 엔지니어는 누구보다 먼저 이렇게 묻는 사람이다. \textit{컴퓨팅이 생기기 전에 누가 같은 청구서를 냈으며, 무엇을 측정했는가?} 더 오래된 것에서 훔쳐라. 이 식탁에는 10년에 한 번쯤 의자가 하나씩 늘어나지만 집안의 격식은 변하지 않는다. 문제만 들여보낸다. 도구는 문 앞에 맡기며, 바로 그곳에서 유효기간이 끝난다.
\end{calloutbox}

\vspace{4.0pt}

\namedsection{E.4}{보증서}

모든 기술서는 유효기간과 함께 출고되지만, 그 날짜를 인쇄한 책은 거의 없다. 독자는 면접에서 폐기된 플래그를 말하고, 사라진 가격표로 예산을 잡고, 이름이 바뀐 YAML 키 하나에 오후를 통째로 쓰면서 어렵게 사실을 깨닫는다. 그러고는 책 전체가 상했다고 잘못 결론 내린다. 그렇지 않다. 시스템 책은 혼합 화물이다. 어떤 내용은 \textbf{우유}이고, 어떤 내용은 \textbf{통조림}이며, 또 어떤 내용은 \textbf{무쇠}다. 정직한 방법은 상자마다 표지를 붙이는 것이다. 여기 눈에 보이는 잉크로 인쇄한 보증서가 있다.

유효기간을 판별하는 데는 1초면 충분하며, 이 책은 물론 어떤 책의 문장에도 적용할 수 있다. 문장에 \textit{제품명, 가격, 버전, 순위}가 들어 있는가? 우유다. 특정 실습실, 지역, 연도 안에서만 참이다. 신선할 때 소비하고 다시 인용하기 전에는 재측정해야 한다. \textit{인터페이스나 메커니즘}, 예를 들어 다이제스트, cgroup 파일, 조정 루프, 히스토그램을 말하는가? 통조림이다. 인터페이스는 조약이고 조약을 깨는 데는 큰 비용이 들기 때문에 약 10년 동안 유행의 변화를 막아 낸다. \textit{항등식, 비율, 스트레스 상황의 인간에 관한 사실}을 말하는가? 무쇠다. L = λW는 1961년에도 특혜를 요구하지 않았고 지금도 그렇다. 부하율은 진공관보다 오래되었으며 경보 피로에는 실제 희생자가 있다. 무쇠 등급은 산업 자체가 바뀌어도 살아남는다. 의미 있는 내구성 시험은 오직 이것뿐이다.

\begin{tbbtablewrap}
  \begin{tabular}{@{}>{\raggedright\arraybackslash}m{\dimexpr 0.3191\tbltextwidth-2\tabcolsep\relax} >{\raggedright\arraybackslash}m{\dimexpr 0.1596\tbltextwidth-2\tabcolsep\relax} >{\raggedright\arraybackslash}m{\dimexpr 0.5213\tbltextwidth-2\tabcolsep\relax}@{}}
  \arrayrulecolor{tbbnavy}\specialrule{1pt}{0pt}{0pt}
  \rowcolor{tbbtablehead}{\centering \bfseries\color{tbbnavy}보증 대상} & {\centering \bfseries\color{tbbnavy}등급} & {\centering \bfseries\color{tbbnavy}보증 조건} \\
  \arrayrulecolor{tbbnavy}\specialrule{0.7pt}{0pt}{2pt}
  \cellcolor{tbbxFFF7ED}{이름, 가격, 플래그: vLLM 스위치, KEDA 스케일러, GPU 시간당 비용, 4,762→900달러 장부, T4와 H100, "소형 = 1\textasciitilde{}3B", 이번 철의 양자화 방법} & \cellcolor{tbbxFFF7ED}{\centering \textcolor{tbbamberdark}{\textbf{우유}\newline \textit{90일\textasciitilde{}3년}}} & \cellcolor{tbbxFFF7ED}{숫자는 특정 실습실, 지역, 연도에서 참이었다. 실제로 얻은 것은 그 숫자를 만든 \textit{방법}, 곧 측정하고, 관문을 적용하고, 다시 측정하는 절차다. 다시 실행해 자신의 숫자를 얻자. 어차피 더 좋아져 있을 것이다.} \\
  \arrayrulecolor{tbbtableborder}\hline
  \cellcolor{tbbxFFF7ED}{순위표: Triton 대 TorchServe 대 BentoML 대 Ray Serve, SageMaker와 Vertex가 관리해 주는 것과 그렇지 않은 것} & \cellcolor{tbbxFFF7ED}{\centering \textcolor{tbbamberdark}{\textbf{우유}\newline \textit{90일\textasciitilde{}3년}}} & \cellcolor{tbbxFFF7ED}{비교표는 상하지만 \textit{머리글에 적힌 질문}은 상하지 않는다. 다음에 나오는 제품에도 같은 질문을 던지고, 승자가 달라져도 미련 없이 받아들이자.} \\
  \arrayrulecolor{tbbtableborder}\hline
  \cellcolor{tbbxEFF6FF}{인터페이스: OCI 이미지와 다이제스트, cgroup v2 파일, 조정 루프, Prometheus의 시계열 모델, HTTP 스트림} & \cellcolor{tbbxEFF6FF}{\centering \textcolor{tbbnavy}{\textbf{통조림}\newline \textit{약 10년}}} & \cellcolor{tbbxEFF6FF}{유행이 아니라 조약이다. 그 위의 대시보드는 해마다 새로 칠하지만, 아래에서는 다이제스트가 여전히 해시이고, 상한은 여전히 커널이 읽는 파일이며, 원하는 상태는 계속 수렴한다. 한 번 배워 직장을 옮겨도 가져갈 수 있다.} \\
  \arrayrulecolor{tbbtableborder}\hline
  \cellcolor{tbbxEFF6FF}{메커니즘: 연속 배칭, 페이지드 KV, 레지스트리→관문→카나리, 두 시간 척도의 오토스케일링, 플릿 백분위수용 히스토그램} & \cellcolor{tbbxEFF6FF}{\centering \textcolor{tbbnavy}{\textbf{통조림}\newline \textit{약 10년}}} & \cellcolor{tbbxEFF6FF}{가중치 스트림과 메모리 벽이라는 물리 법칙, 밤낮에 따라 움직이는 사람 때문에 존재한다. 어느 쪽에도 지원 종료 일정은 없다. 이름은 자주 바뀌고 인수되겠지만 형태는 유지된다.} \\
  \arrayrulecolor{tbbtableborder}\hline
  \cellcolor{tbbxF0FDFA}{계산: L = λW, 백분위수는 더할 수도 평균 낼 수도 없음, 평균은 아무도 설명하지 못함, 128 + 9 = 137, 부하율 = 평균 ÷ 피크} & \cellcolor{tbbxF0FDFA}{\centering \textcolor{tbbteal}{\textbf{무쇠}\newline \textit{서면 보증 10년}}} & \cellcolor{tbbxF0FDFA}{리틀은 아무 가정 없이 법칙을 증명했으므로 시간이 깨뜨릴 것이 없다. 인설의 비율은 진공관보다 오래되었다. POSIX가 137을 약속한 기간은 웹의 역사보다 길다.} \\
  \arrayrulecolor{tbbtableborder}\hline
  \cellcolor{tbbxF0FDFA}{절차: 허용 오차는 측정값보다 먼저 정함, 증상에는 호출하고 원인에는 티켓 발행, 비난 없는 사후 분석, 지루함 = 자동화 + 버전 관리 + 점진성 + 가역성(+ 관측)} & \cellcolor{tbbxF0FDFA}{\centering \textcolor{tbbteal}{\textbf{무쇠}\newline \textit{서면 보증 10년}}} & \cellcolor{tbbxF0FDFA}{스트레스 상황의 인간을 다루는 절차이며 인간에게는 업데이트가 배포되지 않는다. 병원 데이터에는 실제 희생자가 있고, 참고표에는 달 착륙이 있다. 새벽 2시가 존재하는 한 유효하다.} \\
  \arrayrulecolor{tbbtableborder}\hline
  \cellcolor{tbbxF0FDFA}{\textbf{법칙: 새벽 2시에는 오후 2시에 그린 그래프만 보인다}} & \cellcolor{tbbxF0FDFA}{\centering \textcolor{tbbteal}{\textbf{무쇠}\newline \textit{평생}}} & \cellcolor{tbbxF0FDFA}{이 법칙이 실패한다면 그 실패는 누군가 미리 그린 그래프에서만 보일 것이다. 결국 법칙이 다시 이긴 셈이다. 이 조항에 대한 보증 청구는 보통 \textit{연구 결과}라고 부른다.} \\
  \arrayrulecolor{tbbnavy}\specialrule{1pt}{2pt}{0pt}
  \end{tabular}\par
  \figcaption{보증서. 측정하지 않은 경우 무효.}
\end{tbbtablewrap}

\begin{calloutbox}{tbbamber}{tbbamberfill}
{\normalsize\bfseries\color{tbbamber}보증을 청구하는 방법}\par\vspace{3pt}
부하를 걸어 재현하고, 위 조건에 따라 대낮에 그린 그래프를 첨부한다. \textbf{상한 우유}는 결함이 아니다. 날짜를 확인하고 벤치마크를 다시 실행하되 방법은 간직하자. \textbf{10년 안에 통조림이 상했다면} 분노한 블로그 글 한 편과 업계의 진심 어린 사과를 받을 자격이 있다. \textbf{무쇠가 깨졌다면}, 예를 들어 L = λW를 이기는 대기열, 정직하게 평균 낼 수 있는 백분위수, 순수한 기억만으로 해결한 새벽 2시의 디버깅을 발견했다면 보증 문제가 아니다. 출판할 수 있는 연구 결과이며, 이 분야는 인용으로 보상한다. 지금까지 승인된 청구는 한 건도 없다.
\end{calloutbox}

\vspace{4.0pt}

\namedsection{E.5}{아직 겪지 않은 밤을 위한 편지}

이제 봉인된 물건 하나만 남았다. 제14장에서는 네 줄의 대시보드, 소진율, 같은 GPU를 쓰는 이웃, 20분, 비난 0건으로 이루어진 밤을 함께 리허설했다. 그러나 리허설은 리허설일 뿐이다. 앞으로 몇 주가 지났는지도 세지 못할 어느 날, 진짜 밤이 찾아온다. 호출기가 처음 울리고 그 반대편 서비스가 \textit{여러분의 것}인 밤이다. 실습도, 강사도, 정답지도 없다. 이 편지는 그 밤을 위한 것이다. 지금 한 번 읽어 내용을 기억해 두자. 그런 다음 미래의 자신이 비상 물품을 보관할 곳에 넣어 두자. 런북에 붙이거나, 당직 인수인계 문서에 고정하거나, 존재가 떠오르는 그 순간 읽고 있을 대시보드 뒤에 접어 두어도 좋다.

\begin{calloutbox}{tbbpurple}{tbbpurplefill}
{\normalsize\bfseries\color{tbbpurple}새벽 2시에 개봉할 것}\par\vspace{3pt}
엔지니어에게,

예정대로 이 편지를 읽고 있다면 밖은 어둡고, 무언가 잘못되었으며, 호출기가 여러분을 골랐을 것이다. 무엇이든 타이핑하기 전에 의식적으로 한 번 숨을 쉬자. 서비스는 이미 몇 분 동안 망가져 있었다. 5초쯤 더 기다릴 수 있다. 실습에서는 매주 자신의 죽음에서 살아남았다. 그만큼 죽이기 어렵도록 만든 서비스다.

이 밤은 처음이지만 이곳에는 전에 와 본 적이 있다. 지금 벌어지는 일은 전례가 없는 듯 느껴질 것이다. 이 시간의 모든 장애가 그렇다. 그러나 거의 틀림없이 다섯 명의 오랜 지인 가운데 하나다. 무언가 \textbf{가득 찼거나}(대기열, 디스크, 좌석 풀), 무언가 \textbf{종료되었거나}(종료 코드를 읽자. 137은 계량기가 정직했다는 뜻이다), 무언가 \textbf{바뀌었거나}(배포를 찾자. 설정도 변경이다. 설정은 \textit{언제나} 변경이다), 무언가 \textbf{들어왔거나}(아무것도 요청하지 않으면서 모든 것을 가져가는 실리콘 위의 이웃), 무언가 \textbf{드리프트했다}(카나리가 절벽을 지켜보는 동안 세상이 천천히 움직였다). 다섯 가족이다. 이 책에서 모두 대낮에 커피를 마시며 만나 보았다.

대시보드의 줄을 순서대로 읽자. 오후 2시의 두뇌가 아니라 바로 지금의 두뇌를 위해 그 순서로 그려 두었다. 약속부터 본다. 예산이 실제로 타고 있는가, 아니면 원인이 증상처럼 꾸미고 나타났는가? 아드레날린보다 소진율을 믿자. 둘 다 이 밤을 측정하지만 보정된 것은 하나뿐이다. 실제 장애라면 \textbf{이해하기 전에 완화하라.} 롤백은 포인터 하나를 뒤집는 일이고 맑은 정신일 때 연습해 두었다. 사용하는 것은 패배가 아니라 설계가 작동하는 모습이다. 특히 오늘 밤에는 빠른 거절이 느린 가능성보다 낫다.

진행하면서 타임스탬프를 채널에 곧바로 적자. 보기 흉하고 문장 부호가 없어도 된다. 내일의 자신은 그것을 경전처럼 읽을 것이다. 오늘 밤의 목표는 작게 유지하자. 진실이 아니라 \textit{안정}이다. 진실은 대낮에 비난 없이, 사람이 아니라 메커니즘의 이름을 적은 문서로 찾는다. 진지한 사람이라면 왜 더 일찍 보지 못했느냐고 묻지 않는다. 사후 분석은 그저 미리 보았을 패널 하나를 추가한다. 모든 화재는 연기 감지기를 남긴다. 그러면 다음 엔지니어의 밤은 20분 짧아진다. 손으로 쓴 경보 코드 한 장이 달 하강을 30초 더 평온하게 만든 것과 같다. 누군가 여러분을 위해 그렇게 했다. 이제 다음 사람을 위해 그렇게 하자.

장애가 끝나면, 그리고 체감보다 일찍 끝날 것이다, 그 결말의 모습을 눈여겨보자. 승리감이 넘치지 않는다. \textbf{지루하다.} 그래프는 가라앉고, 호출기는 어깨를 으쓱하며, 대시보드는 다시 숨을 쉰다. 김빠지는 결말이 아니라 그것이 바로 제품이다. 15주 동안 만든 것은 절대 망가지지 않는 서비스가 아니다. 누구도 그런 서비스를 만든 적은 없다. 대신 일찍 끝나는 밤, 의미가 있었던 호출, 포인터 한 번으로 적용한 수정, 조정된 경보 임계값 하나만 희생된 아침을 만들었다.

다시 잠자리에 들자. 서비스는 괜찮다. 여러분이 지루하게 만들었다. 처음 30초의 침묵에서 시작했을 때부터 바로 그것이 전부였다.

{\raggedleft \textcolor{tbbgray}{\textit{오후 2시, 불을 켜고 씀}}\par}
\end{calloutbox}

\vspace{4.0pt}

이것이 마지막 사례다. 15주 전에는 완벽한 모델이 30초의 침묵 뒤에 있었고, 방 안의 누구도 그 이유를 설명하지 못했다. 언젠가 찾아올 오늘 밤에는 지루한 서비스가 대시보드 위에서 숨을 쉬며 밀리초로 적은 약속과 스스로 선택한 숫자로 끝나는 예산을 지킨다. 아무도 지켜보지 않는다. 그럴 필요가 없기 때문이다. 이 두 문장 사이에 이 책이 하고 싶었던 모든 이야기가 놓여 있다. 15호선을 타고 종점까지 함께해 주어 고맙다.

\begin{center}
{\large \textcolor{tbbink}{\textbf{이제 지루한 무언가를 만들어 보자.}}}
\end{center}

\begin{center}
\textcolor{tbbgray}{∎}
\end{center}


\end{document}
